% Options for packages loaded elsewhere
\PassOptionsToPackage{unicode}{hyperref}
\PassOptionsToPackage{hyphens}{url}
\PassOptionsToPackage{dvipsnames,svgnames,x11names}{xcolor}
\documentclass[
  10pt,
]{article}
\usepackage{xcolor}
\usepackage[margin=1in]{geometry}
\usepackage{amsmath,amssymb}
\setcounter{secnumdepth}{-\maxdimen} % remove section numbering
\usepackage{iftex}
\ifPDFTeX
  \usepackage[T1]{fontenc}
  \usepackage[utf8]{inputenc}
  \usepackage{textcomp} % provide euro and other symbols
\else % if luatex or xetex
  \usepackage{unicode-math} % this also loads fontspec
  \defaultfontfeatures{Scale=MatchLowercase}
  \defaultfontfeatures[\rmfamily]{Ligatures=TeX,Scale=1}
\fi
\usepackage{lmodern}
\ifPDFTeX\else
  % xetex/luatex font selection
\fi
% Use upquote if available, for straight quotes in verbatim environments
\IfFileExists{upquote.sty}{\usepackage{upquote}}{}
\IfFileExists{microtype.sty}{% use microtype if available
  \usepackage[]{microtype}
  \UseMicrotypeSet[protrusion]{basicmath} % disable protrusion for tt fonts
}{}
\makeatletter
\@ifundefined{KOMAClassName}{% if non-KOMA class
  \IfFileExists{parskip.sty}{%
    \usepackage{parskip}
  }{% else
    \setlength{\parindent}{0pt}
    \setlength{\parskip}{6pt plus 2pt minus 1pt}}
}{% if KOMA class
  \KOMAoptions{parskip=half}}
\makeatother
\setlength{\emergencystretch}{3em} % prevent overfull lines
\providecommand{\tightlist}{%
  \setlength{\itemsep}{0pt}\setlength{\parskip}{0pt}}
\usepackage{bookmark}
\IfFileExists{xurl.sty}{\usepackage{xurl}}{} % add URL line breaks if available
\urlstyle{same}
\hypersetup{
  colorlinks=true,
  linkcolor={Maroon},
  filecolor={Maroon},
  citecolor={Blue},
  urlcolor={Blue},
  pdfcreator={LaTeX via pandoc}}

\author{}
\date{}

\begin{document}

\section{\texorpdfstring{Unconditional One-Way Functions \(\to\) P
\(\neq\) NP via Semantic Conservation
Law}{Unconditional One-Way Functions \textbackslash to P \textbackslash neq NP via Semantic Conservation Law}}\label{unconditional-one-way-functions--p--np-via-semantic-conservation-law}

\textbf{Author}: Prasanth Marreddy \textbf{Email}:
\href{mailto:pmarreddy@gmail.com}{pmarreddy@gmail.com} \textbf{Date}:
December 2025

\subsection{Abstract}\label{abstract}

We prove P \(\neq\) NP for uniform probabilistic polynomial-time Turing
machines by constructing an NP-complete language L* whose exponential
configuration space provably cannot be compressed into polynomially many
computational states by any uniform algorithm. This
configuration-resource inequivalence yields a one-way function under
standard execution semantics; via the classical bridge (OWF
\(\Rightarrow\) FP \(\neq\) FNP \(\Rightarrow\) P \(\neq\) NP), we
establish P \(\neq\) NP for uniform PPT. This resolves
Impagliazzo\textquotesingle{}s Five Worlds conjecture (1995) by
unconditionally placing us in Cryptomania---where both private-key and
public-key cryptography are possible.

\textbf{Key Innovation: The Semantic Conservation Law.} The
incompressibility derives from a structural correctness constraint q +
\(\Phi\) \(\geq\) R, where \(\Phi\) measures simultaneously
distinguishable computational artifacts (log\(_2\) of state count),
while R and q denote required and resolved information. This constraint
yields exactly three operational routes: (1) Storage---maintain
2\^{}(R\(-\)q) distinguishable states in parallel; (2)
Resolution---learn correct values through sequential reads; (3)
Elimination---prune wrong values through testing. L*\textquotesingle{}s
structural properties (Keyedness; Emergence + Bandwidth; Per-Node
Antagonism) block all three routes simultaneously, forcing exponential
cost. Maintaining fewer than 2\^{}(R-q) artifacts causes collisions
between inequivalent seeds, leading to incorrect designated addresses
and verification failure---a correctness requirement, not a performance
heuristic.

\textbf{Construction and One-Wayness.} We define L* as a dependency DAG
with seed-locked problem access: distinct computational histories induce
distinct content-addressed seeds selecting designated memory locations.
The one-way function f: \{0,1\}\^{}m \(\to\) L*\_\{yes\} maps random
string r to planted instance x* = Plant(\(\varphi\), r). Every output
admits a per-instance, deterministic witness-finding lower bound: on any
fixed run, producing the canonical witness requires time \(\geq\)
n\^{}(\(\Omega\)(log n)) or \(\geq\) 2\^{}(\(\Omega\)(n)) (Theorem 8.A).

Any uniform PPT inverter \ensuremath{\mathcal{I}} succeeding with non-negligible probability
can be coin-fixed (Yao\textquotesingle{}s principle) to deterministic
algorithm \ensuremath{\mathcal{I}}\_\{c-\} succeeding on some instance x*. Composing with
polynomial-time extractor Ext yields witness W in polynomial time,
contradicting the per-instance lower bound. Hence f is one-way against
uniform PPT.

\textbf{Information-Theoretic Barrier.} L* achieves incompressibility
through seed-locked encoding: accessing ANY structural property of the
CNF formula \(\varphi\) (unit clauses, literal polarities, clause
structure, variable frequencies) requires computing the correct seed
chain, which requires knowing the solution \(\alpha\). This circular
dependency creates a perfect information barrier---\textbf{algorithms
need information to make progress; L* provides no information.} Every
algorithmic technique enabling shortcuts (unit propagation, clause
learning, branching heuristics, symmetry breaking) requires structural
information that L* systematically removes. The exponential barrier is
information-theoretic, not algorithmic.

\textbf{Significance and Barrier Circumvention.} The approach shifts
from analyzing algorithm behavior within specific computational models
to analyzing what problem structure requires for correctness. The
Semantic Conservation Law \textbf{articulates a common pattern} across
diverse lower bound techniques---decision trees, communication
complexity, pebbling games, branching programs/OBDDs, resolution,
backtracking, dynamic programming, streaming---as structural parallels
exhibiting configuration-space incompressibility. This work
\textbf{formalizes the TM observation paradigm} (bits observed = q,
configs visited = 2\^{}\(\Phi\)) in Lean, bridging SCL to information
theory and enabling unconditional complexity bounds. The other paradigm
correspondences are conceptual (see §11.4 for precise status). The
framework derives bounds from structural correctness requirements rather
than algorithm-specific analyses, using pigeonhole counting and
information-theoretic conservation laws---not barrier-sensitive
techniques like relativization or natural proofs. (See §11.4 for
complete paradigm catalog and formalization status.)

\textbf{Machine Verification.} Complete formalization in Lean 4:
approximately 90,000 lines across 90+ publication-ready files with
minimal trust boundary (two operational axioms). The axioms are: (1)
\texttt{algspec\_has\_tm}---Church-Turing bridge asserting any
polynomial-time algorithmic specification has a TM implementation; and
(2)
\texttt{collision\_indistinguishability\_under\_incomplete\_observation}---semantic
bound from A2 injectivity asserting correctness on planted instances
requires visiting all 2\^{}R configurations. Both are standard
CS/information-theory principles with low trust risk. The Lean
formalization uses a direct OWF-based proof path (OWF \(\to\)
FP\(\neq\)FNP \(\to\) P\(\neq\)NP), which differs from but is equivalent
to the paper\textquotesingle{}s 3-SAT reduction exposition. Axiom audits
via \texttt{\#print\ axioms} provide complete transparency. The Lean
code is the authoritative proof; this paper provides mathematical
intuition and proof narrative.

\textbf{Model Scope.} Classical uniform model: deterministic k-tape
Turing machines with constant tapes and alphabet. Randomized PPT
adversaries handled by coin-fixing (Yao); all bounds apply per fixed
run. Results apply to uniform classical models.

\textbf{Keywords:} one-way functions, P versus NP, computational
complexity, lower bounds, NP-completeness, configuration-space
compression, representation-invariance

\begin{center}\rule{0.5\linewidth}{0.5pt}\end{center}

\subsection{For Reviewers}\label{for-reviewers}

\textbf{Current focus: conceptual soundness.} We invite peer review to
validate the proof architecture before polishing details. Is the
direction correct? Are definitions sound? Do critical lemmas hold?
Let\textquotesingle{}s first confirm we\textquotesingle{}re heading in
the right direction; once validated, we make it airtight. Detailed
polish (missing cases, supporting lemmas, clarity, typos) is deferred
until the core argument is validated.

\textbf{Quick start}: For a 5-minute overview, see
\texttt{docs/TRAPDOOR\_OWF\_MECHANISM.md}.

\textbf{Three foundational shifts from current approaches}:

\begin{itemize}
\tightlist
\item
  \emph{Structure, not algorithms}: Traditional asks "Can we invent a
  faster algorithm?" --- we ask "What does problem structure require?"
\item
  \emph{Correctness, not speed}: Traditional proves "algorithms take too
  many steps" --- we prove "correctness requires too many states"
\item
  \emph{Engineered, not natural}: Traditional analyzes messy natural
  problems (3-SAT, TSP) --- we engineer L* to guarantee required
  properties by construction
\end{itemize}

\textbf{Falsification criteria} (see §3.6 for details):

\begin{itemize}
\tightlist
\item
  \emph{Approach-ending}: Universal compression that works regardless of
  structural properties \(\to\) structural incompressibility cannot
  deliver P\(\neq\)NP
\item
  \emph{Patchable}: L*-specific compression \(\to\) may motivate
  additional properties (A6, A7, ...)
\item
  \emph{Technical gaps}: Flaws in A1-A5, SCL derivation, OWF, or bridge
  \(\to\) repairable through proof revision
\end{itemize}

See \texttt{docs/CONTRIBUTIONS.md} for full reviewer guidelines,
contribution recognition (Co-authorship / Special Thanks / Thanks), and
how to report issues.

\begin{center}\rule{0.5\linewidth}{0.5pt}\end{center}

\textbf{Navigation: Full Path}

\begin{itemize}
\tightlist
\item
  §1: Exponential Spaces and Structural Necessity - why SCL is a
  correctness requirement; universal manifestation across paradigms
\item
  §2: Technical Overview: The Semantic Conservation Law - formal SCL (q
  + \(\Phi\) \(\geq\) R) and proof architecture
\item
  §3: Main Results and Parametric Spectrum - P \(\neq\) NP and
  unconditional Structural OWF construction; complexity bounds for
  different \(\lambda\) regimes
\item
  §4: Model and Semantic Framework - deterministic k-tape TMs; Semantic
  Multiplication Principle (SMP); DAG min-cut framework
\item
  §5: Computational Models and Classes - paradigm-specific SCL
  manifestations (backtracking, DP, OBDD, resolution, circuits)
\item
  §6: Instance Construction \& Invariants - L* structure with overlay;
  structural properties A1-A5 (Hermeticity, Injectivity, Emergence,
  Closure, Dependency)
\item
  §7: The Semantic Conservation Law (SCL) and Semantic Necessity -
  abstract proof: A1-A5 \(\to\) per-node SCL (Theorem 7.A, §7.2.1)
\item
  §8: Per-Instance Deterministic Bounds - foundation for OWF: every
  FG-wired (Frontier-Gate) instance hard on any fixed run (Theorem 8.A);
  coin-fixing extends to randomized PPT
\item
  §9: Unconditional One-Way Function Construction - OWF from
  per-instance bounds (§8) via coin-fixing; OWF \(\Rightarrow\) FP
  \(\neq\) FNP \(\Rightarrow\) P \(\neq\) NP (classical bridge)
\item
  §10: NP-Completeness and Classical Bridge - L* \(\in\) NP; NP-hardness
  via 3-SAT reduction; Proposition 10.4 (classical bridge); Main Theorem
  10.5 (P \(\neq\) NP)
\item
  §11: Related Work - connections to width-based lower bounds,
  algorithm-to-hardness frameworks; §11.4 documents SCL as structural
  parallel across 9 paradigms + formalized TM observation bridge
\item
  §12: Discussion and Implications - scope, limitations, future
  directions
\end{itemize}

\begin{center}\rule{0.5\linewidth}{0.5pt}\end{center}

\textbf{Alternative Navigation} - Two complementary layers with distinct
roles:

\textbf{For Reviewers}: See \texttt{CONTRIBUTIONS.md} for expectations,
contribution recognition (Co-authorship / Special Thanks / Thanks), and
reporting guidelines. For AI-assisted review (\textasciitilde{}10x
faster), see \texttt{docs/AI\_REVIEW\_GUIDE.md}.

\subsubsection{Dual-Layer Architecture: Lean (Authoritative) vs Paper
(Conceptual)}\label{dual-layer-architecture-lean-authoritative-vs-paper-conceptual}

\textbf{Ground truth}: The Lean formalization is the authoritative
proof. This paper provides conceptual understanding and intuition for
what the Lean code proves. For rigorous verification, consult the Lean
code directly.

\begin{center}\rule{0.5\linewidth}{0.5pt}\end{center}

\textbf{PATH 1: Paper (Conceptual Guide)} \emph{Use this path for
intuition, motivation, and informal understanding. This is NOT the
formal proof.}

\textbf{Purpose}: Understand the mathematical ideas, proof strategy, and
why the approach works. All formal claims are proven in Lean (PATH 2).

\textbf{Recommended reading order} (flexible - sections can be read
non-linearly):

\begin{enumerate}
\def\labelenumi{\arabic{enumi}.}
\tightlist
\item
  §1.1 Configuration diversity \(\to\) §2 SCL framework \(\to\) §3
  Results overview (including §3.6 falsifiability)
\item
  §6 L* construction (A1--A5 properties) \(\to\) §7.2.1 SCL proof
  (A1--A5 \(\Rightarrow\) q + \(\Phi\) \(\geq\) R)
\item
  §8.A Per-instance bounds (Theorem 8.A) \(\to\) §9 Extractor +
  Structural OWF security
\item
  §10 NP-completeness + classical bridge \(\to\) §5 Paradigm
  manifestations
\end{enumerate}

\textbf{Alternative for skeptics}: Start with §3.6 (falsifiability),
§2.2-§2.4 (formal definitions), then §6-§10 (proof chain), Appendices
C/A/J (technical details), finally §1+§5 (context).

\textbf{Caveat}: Paper sections provide mathematical narrative, not
formal proofs. Each major claim references the corresponding Lean
theorem for verification.

\begin{center}\rule{0.5\linewidth}{0.5pt}\end{center}

\textbf{PATH 2: Lean (Authoritative Proof)} \emph{Use this path for
rigorous verification, trust assessment, and formal proof checking. This
is the ground truth.}

\textbf{Purpose}: Machine-verified proof with complete trust
transparency. All theorems are proven with 0 \texttt{sorry} statements
and minimal axioms (2 operational axioms, all standard).

\textbf{Proof Architecture}: 13 critical theorems + 1 key definition
compose to establish P\(\neq\)NP. See \texttt{PROOF\_CONTROL\_FLOW.md}
for complete dependency graph and theorem enumeration.

\textbf{Three Verification Approaches} (choose based on verification
goal):

\textbf{Approach A: Critical Theorem Verification} (follow proof spine):
{[}1{]} SCL\_node (SCLNode.lean) \(\to\) {[}2{]} SCL\_cut (SCLCut.lean)
\(\to\) {[}3{]} A2\_keyedness (A2\_Injectivity.lean) \(\to\) {[}4{]}
A3\_emergence (A3\_Emergence.lean) \(\to\) {[}5{]} R\_of\_flat
(RanksExponential.lean) \(\to\) {[}7{]} segment\_reduction
(SegmentReduction.lean) \(\to\) {[}9{]} TM time bound
(TMAdapterExponential.lean) \(\to\) {[}10{]} extractor (Extractor.lean)
\(\to\) {[}11{]} Structural OWF security (StructuralOWFExponential.lean)
\(\to\) {[}12{]} OWF\(\to\)FP\(\neq\)FNP bridge
(StructuralOWFBridge.lean) \(\to\) {[}13{]} FINAL:
fpnefnp\_implies\_not\_peqnp (ParametricBitstringBridge.lean) Execute
\texttt{\#print\ axioms\ \textless{}theorem\_name\textgreater{}} at each
step to verify trust boundary.

\textbf{Approach B: Trust Boundary Audit} (verify axiom usage):

\begin{enumerate}
\def\labelenumi{\arabic{enumi}.}
\tightlist
\item
  Start: \texttt{StructuralOWFExponential.lean} (uses 2 axioms)
\item
  Verify Axiom 1: \texttt{algspec\_has\_tm} (Church-Turing bridge---any
  AlgSpec has TM implementation)
\item
  Verify Axiom 2:
  \texttt{collision\_indistinguishability\_under\_incomplete\_observation}
  (semantic bound from A2 injectivity)
\item
  Conclusion: 2 axioms total, all standard CS/information-theory
  principles (low risk)
\end{enumerate}

\textbf{Proven Theorems} (formerly axioms, eliminated 2025-12-08):

\begin{itemize}
\tightlist
\item
  \texttt{fg\_lossless\_encoding} --- Now fully proven (145-line theorem
  in EncodingDiscipline.lean)
\item
  \texttt{plant\_flat\_wf\_transfer} --- Eliminated by including
  WellFormed in WellFormedRandomness\_flat
\end{itemize}

\textbf{Axiom Layer Note}: Both axioms operate at the
inversion/information layer (TM semantics, A2 injectivity)---neither
mentions P, NP, or complexity bounds. The separation emerges from the
construction, not the axioms.

\textbf{Approach C: Layer-by-Layer Verification} (systematic
architecture): Layer 0 (Foundations) \(\to\) Layer 1 (Construction)
\(\to\) Layer 3 (Information Bounds) \(\to\) Layer 4 (Operational
Bridge) \(\to\) Layer 2 (Security Proof) \(\to\) Layer 5 (Complexity
Bridge) \(\to\) P!=NP Each layer builds on previous layers. Verify
compilation: \texttt{lake\ build} at each layer.

\textbf{Build Verification}: Execute
\texttt{cd\ lean\ \&\&\ lake\ build} (all 90 files,
\textasciitilde{}90,000 lines compile successfully).

\textbf{Detailed Control Flow}: See \texttt{PROOF\_CONTROL\_FLOW.md}
for:

\begin{itemize}
\tightlist
\item
  Complete theorem dependency matrix (13×13)
\item
  Verification checklist (14 items)
\item
  Supporting branches (9 branches, 49 theorems)
\item
  Proof statistics and axiom summary
\end{itemize}

\textbf{Status}: \textasciitilde{}90,000 lines, 2 operational axioms
(all standard).

\textbf{Authoritative status}: This Lean formalization is the definitive
proof. The paper (PATH 1) explains the ideas informally; the Lean code
proves them rigorously. Where paper and Lean diverge (e.g., Lean uses
direct OWF \(\to\) FP\(\neq\)FNP \(\to\) P\(\neq\)NP path, bypassing
paper\textquotesingle{}s 3-SAT reduction §10.2), the Lean formalization
is authoritative.

\subsubsection{Question-Driven
Navigation}\label{question-driven-navigation}

\textbf{Navigate by Your Questions} - Organized by logical dependency
and proof-breaking potential.

\textbf{Reading strategies:}

\begin{itemize}
\tightlist
\item
  \textbf{Skeptics}: Foundation questions first (hunt for fatal flaws
  before building intuition)
\item
  \textbf{Learners}: Context questions first (build intuition before
  verification)
\item
  \textbf{Experts}: Barrier questions first (check impossibility-result
  evasion before foundations)
\end{itemize}

\begin{center}\rule{0.5\linewidth}{0.5pt}\end{center}

\paragraph{\texorpdfstring{\textbf{GROUP 1: PROOF
FOUNDATIONS}}{GROUP 1: PROOF FOUNDATIONS}}\label{group-1-proof-foundations}

\emph{If any question fails, P\(\neq\)NP claim collapses}

\begin{enumerate}
\def\labelenumi{\arabic{enumi}.}
\item
  \textbf{How do I falsify this? What\textquotesingle{}s the compression
  criterion?}

  \begin{itemize}
  \tightlist
  \item
    \textbf{CHECK}: §3.6 (falsification criteria: can you compress
    2\^{}(\(\lambda\)\_base) states to poly(n)?)
  \item
    **CRITICAL*\emph{: Exhibit uniform PPT algorithm A achieving
    \(\lambda\)(A,x}) = O(log n) on test instances

    \begin{itemize}
    \tightlist
    \item
      Type 1 = fundamental breakthrough, Type 2 = fixable design issue
    \end{itemize}
  \item
    \textbf{LEAN}: No mechanical falsification test (would require
    proving negation of P\(\neq\)NP)

    \begin{itemize}
    \tightlist
    \item
      Formalization provides positive proof only
    \end{itemize}
  \end{itemize}
\item
  \textbf{Is L* actually NP-complete? Proof details?}

  \begin{itemize}
  \tightlist
  \item
    \textbf{CHECK}: §10.1 (L*\(\in\)NP via verifier V), §10.2
    (NP-hardness via 3-SAT reduction), §10.3 (combines both)
  \item
    **CRITICAL**: If membership or hardness fails, classical bridge
    breaks
  \item
    \textbf{LEAN}: Formalization uses OWF-based path (bypasses
    NP-completeness)

    \begin{itemize}
    \tightlist
    \item
      See \texttt{ParametricBitstringBridge.lean} for direct OWF \(\to\)
      FP\(\neq\)FNP \(\to\) P\(\neq\)NP proof
    \item
      Paper\textquotesingle{}s 3-SAT reduction path not formalized (both
      approaches valid)
    \end{itemize}
  \end{itemize}
\item
  \textbf{Quantifier structure: \(\forall\)x\(\forall\)A or
  \(\exists\)x\(\forall\)A? Why can\textquotesingle{}t non-uniform
  circuits hardcode solutions?}

  \begin{itemize}
  \tightlist
  \item
    \textbf{CHECK}: Abstract (per-instance deterministic claim), §4.3
    (uniform restriction), Theorem 8.A
  \item
    **CRITICAL**: \(\forall\)x\(\forall\)A with uniformity prevents
    hardcoding

    \begin{itemize}
    \tightlist
    \item
      \(\exists\)x\(\forall\)A would fail (non-uniform circuits can
      hardcode "if input=x, output L(x)")
    \end{itemize}
  \item
    \textbf{LEAN}: \texttt{TMAdapterExponential.lean} - theorem
    \texttt{fg\_first\_commit\_time\_lower\_bound}

    \begin{itemize}
    \tightlist
    \item
      Applies to all FG-wired instances (\(\forall\)x* quantifier)
    \item
      Uniformity enforced via \texttt{PPTAdversary} structure requiring
      fixed C,k constants
    \end{itemize}
  \end{itemize}
\item
  \textbf{How does extractor Ext work and why is it polynomial-time?}

  \begin{itemize}
  \tightlist
  \item
    \textbf{CHECK}: §9.2 (Ext construction: r\textquotesingle{}
    \(\mapsto\) W), §9.3 (poly-time analysis), §9.4 (composition
    \ensuremath{\mathcal{I}}\(\circ\)Ext)
  \item
    **CRITICAL**: Structural OWF security requires Ext maps inversion
    \(\to\) witness in poly-time

    \begin{itemize}
    \tightlist
    \item
      If Ext fails or is exponential, contradiction fails
    \end{itemize}
  \item
    \textbf{LEAN}: \texttt{Extractor.lean} - witness extraction from
    preimage (polynomial-time parsing)

    \begin{itemize}
    \tightlist
    \item
      Used in \texttt{StructuralOWFExponential.lean} theorem
      \texttt{f\_is\_structural\_owf\_exponential\_flat} for
      contradiction
    \end{itemize}
  \end{itemize}
\item
  \textbf{What is Frontier-Gate (FG) and how does it provide
  per-instance deterministic bounds?}

  \begin{itemize}
  \tightlist
  \item
    \textbf{CHECK}: §8.1 (FG overview), Theorem 8.A, Appendix C.1.1
    (digest cost), Appendix C.2 (segment count)
  \item
    **CRITICAL**: Per-instance deterministic bounds enable OWF via
    coin-fixing

    \begin{itemize}
    \tightlist
    \item
      Distributional bounds insufficient
    \end{itemize}
  \item
    \textbf{LEAN}: \texttt{FrontierGate.lean} - FG mechanism
    formalization

    \begin{itemize}
    \tightlist
    \item
      \texttt{SegmentReduction.lean} - theorem
      \texttt{refutation\_count\_exponential\_bound}
    \item
      Lower bound: \(\geq\) 2\^{}(\(\rho\)-s)
    \item
      Verify with
      \texttt{\#print\ axioms\ refutation\_count\_exponential\_bound}
    \end{itemize}
  \end{itemize}
\item
  \textbf{What are A1-A5 properties and why necessary/sufficient for
  SCL?}

  \begin{itemize}
  \tightlist
  \item
    \textbf{CHECK}: §6 (A1-A5 definitions), §7.2.1 (A1-A5 \(\to\)
    q+\(\Phi\)\(\geq\)R proof), Lemma 7.I
  \item
    **CRITICAL**: Entire SCL proof depends on these axioms

    \begin{itemize}
    \tightlist
    \item
      If any unnecessary/insufficient, proof structure suspect
    \end{itemize}
  \item
    \textbf{LEAN}: A1 (\texttt{A1\_Hermeticity.lean}), A2
    (\texttt{A2\_Injectivity.lean}), A3 (\texttt{A3\_Emergence.lean})

    \begin{itemize}
    \tightlist
    \item
      \texttt{SCLNode.lean} - theorem \texttt{SCL\_node} derives bound
      from \texttt{keyed} property (combines A1+A2)
    \item
      Execute \texttt{\#print\ axioms\ SCL\_node} to verify 0 custom
      axioms
    \end{itemize}
  \end{itemize}
\item
  \textbf{Lane dichotomy: Does it cover ALL algorithms (no escape
  routes)?}

  \begin{itemize}
  \tightlist
  \item
    \textbf{CHECK}: §7.3 (restart vs single-run), Theorem 7.B (every
    computation falls into exactly one lane), Appendix C (exhaustiveness
    proof)
  \item
    **CRITICAL**: If hybrids escape both lanes, super-poly bound might
    fail
  \item
    \textbf{LEAN}: Not explicitly formalized as separate "lane
    dichotomy" theorem. Bounds apply to any TM execution trace via
    \texttt{buildRunFromTMTrace} in \texttt{TMToExecutionPrefix.lean}.
    Exhaustiveness follows from TM semantics completeness.
  \end{itemize}
\item
  \textbf{Can you bypass the overlay and solve the CNF \(\varphi\)
  directly (OAP)?}

  \begin{itemize}
  \tightlist
  \item
    \textbf{CHECK}: §10.1.1 (OAP: Overlay-as-Problem), §6 (seed-locked
    decode schema \(\Phi\)\textasciitilde{}), §10.4.1 (bypass discussion Theorem
    10.4.1-BYP)
  \item
    **CRITICAL**: If CNF can be solved directly without overlay
    engagement, entire construction fails; must verify decode schema
    \(\Phi\)\textasciitilde{} prevents standalone access (mask bits R at seed-dependent
    addresses)
  \item
    \textbf{LEAN}: Seed-locked encoding enforced via
    \texttt{SeedChain.lean} and \texttt{Pools.lean} structures. CNF
    access requires seed computation (dependency enforced by type
    system). No explicit "bypass impossibility" theorem (structural
    property of construction).
  \end{itemize}
\item
  \textbf{Is representation-invariance proven or assumed? What about
  different encodings?}

  \begin{itemize}
  \tightlist
  \item
    \textbf{CHECK}: §10 (NP-completeness under polynomial-time
    reductions), §1.1 (configuration-resource inequivalence), Abstract
    (standard many-one encodings)
  \item
    **CRITICAL**: Claim is configuration-space compression persists
    across encodings; verify this holds under standard polynomial-time
    many-one reductions (not just L*\textquotesingle{}s specific
    encoding). NP-completeness provides representation-invariance via
    reduction closure.
  \item
    \textbf{LEAN}: Artifact-encoding invariance proven (cardinality
    bounds via \texttt{Fintype.card} independent of encoding choice).
    Cross-reduction stability not formalized (future work, see paper
    §1.7 Scope). Hardness proven for specific L* encoding (sufficient
    for P\(\neq\)NP via OWF path).
  \end{itemize}
\end{enumerate}

\begin{center}\rule{0.5\linewidth}{0.5pt}\end{center}

\paragraph{\texorpdfstring{\textbf{GROUP 2: CORE
MECHANISMS}}{GROUP 2: CORE MECHANISMS}}\label{group-2-core-mechanisms}

\begin{enumerate}
\def\labelenumi{\arabic{enumi}.}
\item
  \textbf{How does SCL (q + \(\Phi\) \(\geq\) R) actually work
  mathematically?}

  \begin{itemize}
  \tightlist
  \item
    \textbf{PATH}: §1.1 (intuition: collision argument) \(\to\) §2.2
    (formal definitions: R\_v, q\_v, \(\Phi\)\_v) \(\to\) §7.2.1
    (rigorous proof)
  \item
    \textbf{ANSWER}: A1-A5 \(\to\) Injectivity+Keyedness \(\to\) Alt\_v
    \(\geq\) 2\^{}(R\_v-q\_v) \(\to\) \(\Phi\)\_v \(\geq\) R\_v-q\_v
    \(\to\) SCL
  \item
    \textbf{LEAN}: \texttt{SCLNode.lean} - theorem \texttt{SCL\_node}
    (per-node bound: Fintype.card v.State \(\geq\) 2 \^{} lambda v).
    Proof uses pigeonhole principle via \texttt{keyed} property (no
    cross-seed merges). \texttt{SCLCut.lean} - theorem \texttt{SCL\_cut}
    (cut composition via multiplicative principle). Both have 0 axioms.
  \end{itemize}
\item
  \textbf{What is planting function Plant(\(\varphi\),r) and why does
  every output have a witness?}

  \begin{itemize}
  \tightlist
  \item
    \textbf{PATH}: §9.1 (f: r \(\mapsto\) Plant(\(\varphi\),r)
    definition), §9.2 (planting algorithm), Appendix O (full details)
  \item
    \textbf{ANSWER}: r = (assignment, gateDigests, structuralSalt)
    \(\in\) D(\(\varphi\)) \(\to\) Plant constructs x* with
    identity-based digests (non-leaking); seed computation is
    assignment-independent; any valid domain element r\textquotesingle{}
    \(\in\) D(\(\varphi\)) with plant(r\textquotesingle{})=x* contains a
    satisfying assignment (Lemma 9.DOM)
  \item
    \textbf{LEAN}: \texttt{PlantExponential.lean} - planting algorithm
    formalization. Witness extraction via \texttt{Extractor.lean}. FG
    wiring via \texttt{FrontierGate.lean} (FG gates). All outputs have
    canonical witness by construction (structural property, not separate
    theorem).
  \end{itemize}
\item
  \textbf{Three routes (Storage/Resolution/Elimination): How does L*
  block all simultaneously?}

  \begin{itemize}
  \tightlist
  \item
    \textbf{PATH}: §1.1-§1.2.1 (three-way framework), §7.3+Appendix C
    (lane dichotomy)
  \item
    \textbf{ANSWER}: Keyedness blocks Storage; Emergence+Bandwidth block
    Resolution; Per-node+CDT block Elimination \(\to\) no polynomial
    escape
  \item
    \textbf{LEAN}:

    \begin{itemize}
    \tightlist
    \item
      Storage blocked: \texttt{SCLNode.lean} (exponential state
      requirement)
    \item
      Resolution blocked: \texttt{A3\_Emergence.lean} (fresh bits per
      node)
    \item
      Elimination blocked: \texttt{SegmentReduction.lean} (exponential
      refutation count)
    \item
      Three-way simultaneity: consequence of A1-A5 combination, not
      single theorem
    \end{itemize}
  \end{itemize}
\end{enumerate}

\begin{center}\rule{0.5\linewidth}{0.5pt}\end{center}

\paragraph{\texorpdfstring{\textbf{GROUP 3: COMPLETENESS COVERAGE}
{[}YES{]}}{GROUP 3: COMPLETENESS COVERAGE {[}YES{]}}}\label{group-3-completeness-coverage-yes}

\begin{enumerate}
\def\labelenumi{\arabic{enumi}.}
\tightlist
\item
  \textbf{Randomization / heuristics / SAT solvers / CDCL: Are these
  covered?}

  \begin{itemize}
  \tightlist
  \item
    \textbf{PATH}: §9.4 (coin-fixing via Yao), §5.4 (CDCL behavior and
    lanes), Appendix C.4.2 (expected tries \(\geq\) 2\^{}\(\lambda\)),
    Appendix J (Theorem J.1: elimination factoring)
  \item
    \textbf{ANSWER}: Coin-fixing \(\to\) deterministic run (SCL applies
    per fixed run). CDCL with aggressive restarts and bounded/purged
    clause memory behaves like the restart lane (expected tries \(\geq\)
    2\^{}(\(\Omega\)(\(\lambda\)))); when persistent learned clauses are
    retained, they contribute to \(\Phi\) and are priced under the
    single-run lane. Heuristics affect constants, not asymptotics.
  \end{itemize}
\end{enumerate}

\begin{center}\rule{0.5\linewidth}{0.5pt}\end{center}

\paragraph{\texorpdfstring{\textbf{GROUP 4: BARRIER EVASION}
}{GROUP 4: BARRIER EVASION }}\label{group-4-barrier-evasion-construction}

\begin{enumerate}
\def\labelenumi{\arabic{enumi}.}
\tightlist
\item
  \textbf{How does this avoid relativization, natural proofs,
  algebrization barriers?}

  \begin{itemize}
  \tightlist
  \item
    \textbf{PATH}: §1 (problem structure vs algorithm behavior), §7.2.1
    (A1-A5 \(\to\) SCL structural proof), §12.2--§12.3 (detailed barrier
    discussion)
  \item
    \textbf{ANSWER}: Analyzes what L* requires for correctness (A1-A5
    structural obligations \(\to\) SCL via pigeonhole), not what
    algorithms do relative to oracles/models; explicit construction (not
    generic large-circuit lower bound)
  \end{itemize}
\end{enumerate}

\begin{center}\rule{0.5\linewidth}{0.5pt}\end{center}

\paragraph{\texorpdfstring{\textbf{GROUP 5: MODEL SCOPE}
}{GROUP 5: MODEL SCOPE }}\label{group-5-model-scope-dart}

\begin{enumerate}
\def\labelenumi{\arabic{enumi}.}
\item
  \textbf{Per-instance deterministic vs distributional/average-case:
  What\textquotesingle{}s the difference?}

  \begin{itemize}
  \tightlist
  \item
    \textbf{PATH}: §§8--10 (per-instance deterministic), §9.4
    (coin-fixing preserves per-instance)
  \item
    \textbf{ANSWER}: \textbf{Per-instance}: \(\forall\)x*\(\in\)L*\_FG
    hard on any fixed run (stronger). \textbf{Distributional}: hard on
    average over distribution (typical crypto: factoring). Every
    FG-wired instance hard \(\to\) no easy-instance escape hatch \(\to\)
    unconditional Structural OWF
  \end{itemize}
\item
  \textbf{Model scope: Uniform PPT only? What about quantum,
  non-uniform, advice?}

  \begin{itemize}
  \tightlist
  \item
    \textbf{PATH}: §4.3 (model specification), Abstract (scope
    limitations), §12.2--§12.3 (open questions)
  \item
    \textbf{ANSWER}: {[}YES{]} Uniform classical PPT
    (deterministic+randomized via coin-fixing). {[}NO{]} Non-uniform
    circuits (hardcoding defeats \(\forall\)x), quantum (out of scope),
    advice/oracles (defeats uniformity)
  \end{itemize}
\item
  \textbf{Search (witness-finding) vs decision (membership): Which task
  has lower bound?}

  \begin{itemize}
  \tightlist
  \item
    \textbf{PATH}: §1 (technical note: F\_can vs decision), §10.4.1
    (decision poly-time with witness), Theorem 8.A (witness-finding
    super-poly)
  \item
    \textbf{ANSWER}: Lower bound applies to \textbf{search} (canonical
    witness W=(w,G\_\(\tau\),Dig\_\(\tau\)) finding). Decision with
    witness provided is poly-time (§10.1 verifier V). Gap: verification
    easy, search hard
  \end{itemize}
\end{enumerate}

\begin{center}\rule{0.5\linewidth}{0.5pt}\end{center}

\paragraph{\texorpdfstring{\textbf{GROUP 6: UNDERSTANDING \& CONTEXT}
}{GROUP 6: UNDERSTANDING \& CONTEXT }}\label{group-6-understanding--context-bulb}

\begin{enumerate}
\def\labelenumi{\arabic{enumi}.}
\item
  \textbf{Main claim and proof chain?}

  \begin{itemize}
  \tightlist
  \item
    \textbf{QUICK PATH}: §3 (results) \(\to\) §1.1 (intuition) \(\to\)
    §7.2.1 (SCL) \(\to\) §8.A (bounds) \(\to\) §9.4 (OWF) \(\to\) §10.4
    (bridge)
  \item
    \textbf{CHAIN}: L* NP-complete + OWF f(r)=Plant(\(\varphi\),r) with
    FG + per-instance bounds + Ext poly-time + coin-fixing \(\to\)
    contradiction \(\to\) OWF secure \(\to\) P\(\neq\)NP
  \end{itemize}
\item
  \textbf{How is this different from cryptography?}

  \begin{itemize}
  \tightlist
  \item
    \textbf{PATH}: §1.1.1 (configuration-diversity), §2.6 (worked
    example), §3.5 (Structural OWF construction)
  \item
    \textbf{ANSWER}: Both use configuration-diversity (exponential
    patterns over polynomial resources). Crypto: algebraic trap doors
    (assume hardness). L*: dependency geometry A1-A5 (prove hardness via
    pigeonhole)
  \end{itemize}
\item
  \textbf{Where\textquotesingle{}s the technical innovation?}

  \begin{itemize}
  \tightlist
  \item
    \textbf{PATH}: §6 (A1-A5 engineered incompressibility), §2.6
    (configuration-diversity example), §8.1 (per-instance bounds),
    §9.1-§9.3 (FG+Ext)
  \item
    \textbf{ANSWER}: Configuration-diversity barrier (pattern diversity
    not resource quantity) + per-instance deterministic bounds (every
    instance hard any run) + structural proof (problem requirements not
    algorithm behavior) + unconditional Structural OWF (no crypto
    assumptions)
  \end{itemize}
\item
  \textbf{Verification roadmap?}

  \begin{itemize}
  \tightlist
  \item
    \textbf{DEPENDENCY CHAIN}: A1-A5 (§6) \(\to\) SCL (§7.2.1) \(\to\)
    FG+lanes (§8.A, Theorem 7.B + Appendix C) \(\to\) Ext (§9.2-9.3)
    \(\to\) OWF (§9.4) \(\to\) NP-complete (§10.1-10.3) \(\to\)
    P\(\neq\)NP (§10.4)
  \end{itemize}
\end{enumerate}

\begin{itemize}
\tightlist
\item
  \textbf{APPENDICES}: C (lane exhaustiveness proof), A (encoding
  injectivity), J (elimination factoring)
\end{itemize}

\paragraph{\texorpdfstring{Bottleneck residual \(\lambda\): How does it
determine complexity across
models?}{Bottleneck residual \textbackslash lambda: How does it determine complexity across models?}}\label{bottleneck-residual-ux3bb-how-does-it-determine-complexity-across-models}

\begin{itemize}
\tightlist
\item
  \textbf{PATH}: §1.2 (complexity spectrum), §5 (paradigm
  manifestations), Theorem 7.B (SCL \(\to\) bounds)
\item
  \textbf{ANSWER}: Must maintain \(\geq\) 2\^{}\(\lambda\)
  distinguishable artifacts (correctness via SCL). Same \(\lambda\),
  different "currency": backtracking (tree size), DP (keys), OBDD
  (width), resolution (proof size), k-tape TM (time via bandwidth)
\end{itemize}

\paragraph{P problems succeed, NP-complete resist:
Why?}\label{p-problems-succeed-np-complete-resist-why}

\begin{itemize}
\tightlist
\item
  \textbf{PATH}: §1.3 (three-way framework), §1.4 (search-verification
  gap)
\item
  \textbf{ANSWER} (interpretive for general; rigorous for L*): P
  problems allow all three routes (Storage, Resolution, Elimination)
  polynomial \(\to\) \(\lambda\)=O(log n). L* blocks all three
  simultaneously (Keyedness, Emergence+Bandwidth, Per-node+CDT) \(\to\)
  \(\lambda\)=\(\Omega\)(log² n) or \(\Omega\)(n)
\end{itemize}

\begin{center}\rule{0.5\linewidth}{0.5pt}\end{center}

\textbf{VERIFICATION SEQUENCES:}

\paragraph{A. FASTEST FAILURE-HUNTING
(paper-based)}\label{a-fastest-failure-hunting-paper-based}

\emph{Prioritize proof-breaking questions}

\begin{enumerate}
\def\labelenumi{\arabic{enumi}.}
\tightlist
\item
  Group 1.0 (falsifiability: compression criterion clear? empirical
  tests?)
\item
  Group 1.2 (quantifiers \(\forall\)x\(\forall\)A? uniform sufficient?)
\item
  Group 1.3 (extractor Ext poly-time? correct?)
\item
  Group 1.4 (FG per-instance deterministic? evasion possible?)
\item
  Group 1.5 (A1-A5 necessary? sufficient?)
\item
  Group 1.6 (lanes exhaustive? hybrids escape?)
\item
  Group 1.7 (OAP: can overlay be bypassed?)
\item
  Group 1.8 (representation-invariance: proven or assumed?)
\item
  Group 1.1 (NP-completeness reduction sound?)
\end{enumerate}

\paragraph{B. SYSTEMATIC FOUNDATIONS
(paper-based)}\label{b-systematic-foundations-paper-based}

\emph{Follow logical dependency chain}

Definitions (§2.2-§2.4) \(\to\) A1-A5 (§6) \(\to\) SCL (§7.2.1) \(\to\)
FG (§8.A) \(\to\) lanes (Theorem 7.B + Appendix C exhaustiveness)
\(\to\) Ext (§9.2-9.3) \(\to\) NP-complete (§10.1-10.3) \(\to\) bridge
(§10.4) \(\to\) quantifiers \& representation-invariance (verify)

\paragraph{C. BARRIER-EVASION FIRST (experts,
paper-based)}\label{c-barrier-evasion-first-experts-paper-based}

\emph{Check impossibility-result evasion before foundations}

Group 4 (barriers) \(\to\) Group 1.8 (representation-invariance) \(\to\)
Group 1.2 (quantifiers) \(\to\) Group 5.1 (per-instance vs
distributional) \(\to\) Group 1.5 (A1-A5 structural) \(\to\) then Group
1 (remaining foundations)

\paragraph{D. LEAN MECHANICAL VERIFICATION
(code-based)}\label{d-lean-mechanical-verification-code-based}

\emph{Machine-checkable trust assessment}

\begin{enumerate}
\def\labelenumi{\arabic{enumi}.}
\item
  \textbf{Trust boundary} \(\to\) \texttt{AXIOM\_FINAL\_COUNT.md} (read
  axiom documentation)

  \begin{itemize}
  \tightlist
  \item
    Execute \texttt{\#print\ axioms\ fpnefnp\_implies\_not\_peqnp}
    (inspect full dependency chain)
  \end{itemize}
\item
  \textbf{Main theorem} \(\to\) \texttt{ParametricBitstringBridge.lean}

  \begin{itemize}
  \tightlist
  \item
    Verify P\(\neq\)NP proof compiles, 0 sorries
  \end{itemize}
\item
  \textbf{Structural OWF security} \(\to\)
  \texttt{StructuralOWFExponential.lean} (verify one-wayness, 0 sorries)

  \begin{itemize}
  \tightlist
  \item
    Execute
    \texttt{\#print\ axioms\ f\_is\_structural\_owf\_exponential\_flat}
  \end{itemize}
\item
  \textbf{Per-instance bounds} \(\to\)
  \texttt{TMAdapterExponential.lean} (verify TM time bound)

  \begin{itemize}
  \tightlist
  \item
    Also verify \texttt{SegmentReduction.lean} (segment counting)
  \end{itemize}
\item
  \textbf{SCL framework} \(\to\) \texttt{SCLNode.lean} (verify per-node
  bound, 0 axioms)

  \begin{itemize}
  \tightlist
  \item
    Execute \texttt{\#print\ axioms\ SCL\_node} (confirm pure pigeonhole
    counting)
  \end{itemize}
\item
  \textbf{Construction} \(\to\) Properties directory (verify A1, A2, A3
  properties)

  \begin{itemize}
  \tightlist
  \item
    Also verify \texttt{PlantExponential.lean} (planting algorithm)
  \end{itemize}
\item
  \textbf{Build verification} \(\to\) Execute
  \texttt{cd\ lean\ \&\&\ lake\ build}

  \begin{itemize}
  \tightlist
  \item
    Confirm all 90 files compile successfully
  \end{itemize}
\end{enumerate}

\textbf{Lean verification advantages}: 2 operational axioms (all
standard), executable proofs, axiom audits via \texttt{\#print\ axioms},
complete trust transparency. Use this path for highest confidence or
when verifying specific claims mechanically.

\begin{center}\rule{0.5\linewidth}{0.5pt}\end{center}

\subsection{Part I: Motivation \&
Overview}\label{part-i-motivation--overview}

\subsubsection{1. Exponential Spaces and Structural
Necessity}\label{1-exponential-spaces-and-structural-necessity}

\textbf{Purpose of Section 1:} Establish the intuitive foundation for
the Semantic Conservation Law (SCL) and explain why it imposes
unavoidable computational requirements for L*. This section introduces
the core framework informally before the rigorous mathematics of
§§6--10.

\textbf{Three-stage approach:}

\begin{enumerate}
\def\labelenumi{\arabic{enumi}.}
\tightlist
\item
  \textbf{Intuition} (§1.1--§1.2): Why algorithms must satisfy q +
  \(\Phi\) \(\geq\) R (a correctness condition from L*'s structure), and
  why the three routes---Storage, Resolution, Elimination---cannot all
  be kept polynomial for L* simultaneously
\item
  \textbf{Formalization} (§1.3): Precise definitions showing how each
  route maps to SCL\textquotesingle{}s terms
\item
  \textbf{Search-verification gap} (§1.4): How the framework explains
  the asymmetry between search and verification
\end{enumerate}

\textbf{Roadmap:}

\begin{itemize}
\tightlist
\item
  \textbf{Impagliazzo\textquotesingle{}s Five Worlds}: P vs NP through a
  cryptographic lens---why this proof places us in Cryptomania
\item
  \textbf{§1.1}: The real question---compressing exponential
  configuration spaces
\item
  \textbf{§1.2}: The Semantic Conservation Law---a correctness
  condition, not a heuristic
\item
  \textbf{§1.3}: How the three routes determine complexity
\item
  \textbf{§1.4}: The search-verification gap
\item
  \textbf{§1.5}: Building on existing theory---unifying the structural
  perspective
\item
  \textbf{§1.6}: Technical foundations---SCL vs. Shannon
\item
  \textbf{§1.7}: Scope
\end{itemize}

\textbf{Relation to rigorous proofs:} The formal results appear in
§§6-10: L*\textquotesingle{}s construction with properties A1-A5 (§6);
proof that A1-A5 mathematically imply SCL (§7); per-instance
deterministic bounds (§8); Structural OWF construction (§9);
NP-completeness and classical bridge (§10). Section 1 provides
intuition; subsequent sections provide proof.

\textbf{Key Terms Used in Section 1:}

\begin{itemize}
\tightlist
\item
  \textbf{R\_v (Required bits):} Information that must emerge at node v
  (determined by problem structure via Emergence axiom A3)
\item
  \textbf{q\_v (Resolved bits):} Information determined by the algorithm
  at node v (via designated reads; measured by RWA (Receiving-Window
  Attribution))
\item
  \textbf{\(\Phi\)\_v (Potential):} log\(_2\)(simultaneously
  distinguishable computational artifacts) = log\(_2\)(Alt\_v)
  maintained at node v
\item
  \textbf{\(\lambda\) (Bottleneck residual):} min\_C
  \(\Sigma\)\_\{v\(\in\)C\}(R\_v - q\_v) = minimum unresolved residual
  across all cuts; determines complexity
\item
  \textbf{World:} A consistent assignment of values to all node outputs;
  distinct seeds identify distinct worlds
\item
  \textbf{Artifact:} A computational state (DP table key, backtracking
  node, OBDD node, etc.) distinguished by its content
\item
  \textbf{Collision:} When distinct worlds map to the same artifact
  \(\to\) wrong seed used \(\to\) wrong addresses \(\to\) wrong output
  \(\to\) verifier detects error
\item
  \textbf{SCL (Semantic Conservation Law):} q + \(\Phi\) \(\geq\) R (a
  correctness condition derived from problem structure)
\end{itemize}

\textbf{Key insight:} Rather than asking "how do algorithms behave in
model M?" we ask "what does L*\textquotesingle{}s structure require for
correctness?" This configuration-diversity question---can exponentially
many configurations be represented using only polynomially many
simultaneously distinguishable computational states?---yields
representation-invariant incompressibility (within L*) via
information-theoretic conservation laws. (§1.1 develops this
perspective.)

\textbf{Impagliazzo\textquotesingle{}s Five Worlds: P vs NP Through a
Cryptographic Lens}

Russell Impagliazzo (1995) reframed the P vs NP question by identifying
five possible "worlds" based on which cryptographic primitives can
exist:

\begin{itemize}
\tightlist
\item
  \textbf{Algorithmica}: P = NP; cryptography is impossible (all secrets
  efficiently recoverable)
\item
  \textbf{Heuristica}: P \(\neq\) NP but no OWF; worst-case hardness
  exists but doesn\textquotesingle{}t translate to average-case (no
  useful cryptography)
\item
  \textbf{Pessiland}: Hard problems exist but yield no useful OWF;
  problems are hard but unusable for crypto
\item
  \textbf{Minicrypt}: OWF exist; private-key cryptography possible
  (symmetric encryption, MACs, commitments)
\item
  \textbf{Cryptomania}: Structural OWF exist; public-key cryptography
  possible (PKE, signatures, key exchange)
\end{itemize}

This framework captures a fundamental question: \emph{what is the
relationship between computational hardness and cryptographic
possibility?} Proving P \(\neq\) NP alone doesn\textquotesingle{}t
resolve which world we inhabit---Heuristica and Pessiland both have P
\(\neq\) NP but no useful cryptography.

\textbf{This proof places us unconditionally in Cryptomania.} L*
provides both OWF (enabling all of Minicrypt) and Structural OWF
(enabling Cryptomania). The resolution is unconditional---we do not
assume factoring is hard or discrete log is intractable. OWF existence
follows from L*\textquotesingle{}s information-theoretic structure (A2
injectivity on R-bit emergent configurations, with identity digest as
discriminator), not from conjectured computational hardness of
number-theoretic problems.

\textbf{Proof Strategy: State Compression Impossibility}

Historically, NP was defined via nondeterministic Turing machines and
polynomial-time verification: L \(\in\) NP if a witness can be verified
in polynomial time. This frames P vs NP as "verification easy, decision
hard"---an abstract formulation difficult to connect to lower bounds.

We take a different approach: prove \textbf{One-Way Function (OWF)
existence}, which implies P \(\neq\) NP via the classical bridge (OWF
\(\Rightarrow\) FP \(\neq\) FNP \(\Rightarrow\) P \(\neq\) NP). The OWF
framing asks: "Can I reverse a computation?"---directly connected to
\textbf{state space requirements}.

\begin{itemize}
\tightlist
\item
  \textbf{Verification view} (traditional): "There exists a language
  where verification is easy but decision is hard." Abstract;
  step-counting arguments hit barriers (relativization, natural proofs).
\item
  \textbf{Inversion view} (this work): "There exists a function easy to
  compute but hard to invert." Concrete; hardness derives from state
  compression impossibility---algorithms cannot represent 2\^{}R
  configurations with fewer than 2\^{}R states while maintaining
  correctness.
\end{itemize}

\textbf{Why inversion admits structural proofs.} The critical
distinction is between \emph{computational hiding} and \emph{structural
incompressibility}:

\begin{itemize}
\tightlist
\item
  \textbf{Computational hiding} (e.g., factoring): N = p × q encodes
  both factors completely. Information is present but computationally
  hard to extract. Algorithms might find shortcuts; quantum computers do
  (Shor\textquotesingle{}s algorithm).
\item
  \textbf{Structural incompressibility} (this work): L* has 2\^{}R
  configurations that must map to distinct computational states
  (keyedness from A2 injectivity). Attempting to compress---using fewer
  states---causes correctness failures. This is a counting argument via
  pigeonhole, not an entropy bound.
\end{itemize}

The Semantic Conservation Law (q + \(\Phi\) \(\geq\) R) formalizes this:
to correctly process R bits of structural requirement, an algorithm must
either resolve them (q) or maintain distinguishing state (\(\Phi\)).
Violating this constraint doesn\textquotesingle{}t slow computation---it
produces wrong answers.

\textbf{The strategic path:} State compression impossibility \(\to\)
inversion hardness \(\to\) OWF \(\to\) P\(\neq\)NP.
L*\textquotesingle{}s A2 injectivity ensures 2\^{}R configurations map
to distinct seeds, which must map to distinct algorithm states. By
pigeonhole, compression is impossible without correctness failure.

\textbf{The hardness mechanism.} The 2\^{}R lower bound comes from
keyedness (A2 injectivity):

\begin{itemize}
\tightlist
\item
  \textbf{Keyedness}: Different R-bit emergent configurations \(\to\)
  different seeds \(\to\) different designated addresses \(\to\)
  different required outputs
\item
  \textbf{Pigeonhole}: 2\^{}R configurations exist; each requires a
  distinct state to produce its correct output
\item
  \textbf{Compression fails}: Merging two configurations into one state
  \(\to\) same output for both \(\to\) wrong answer for at least one
\end{itemize}

\textbf{Lean formalization}: Theorem \texttt{SCL\_node} (SCLNode.lean)
proves: if keyedness holds, then \textbar{}State\textbar{} \(\geq\)
2\^{}\(\lambda\). The proof is pure counting---keyedness provides
injection from configurations to states; pigeonhole gives the bound. No
entropy formulas, no Shannon bounds.

\textbf{Why This Approach Circumvents Known Barriers:}

We shift from analyzing algorithm behavior within computational models
to analyzing what problem structure requires for correctness. This
structural perspective naturally avoids known barriers:

\textbf{Natural proofs (Razborov--Rudich, 1997):} Apply to sufficiently
general circuit lower bounds. Our approach is non-generic---we construct
a specific problem L* with explicit structural properties (A1--A5) that
enforce incompressibility via pigeonhole counting, not general circuit
analysis.

\textbf{Relativization (Baker--Gill--Solovay, 1975):} Constrains
oracle-based arguments. Our proof analyzes L*\textquotesingle{}s
structural requirements (seed-determined addressing, emergence
constraints, keyedness) which are intrinsic to the problem, not
oracle-dependent.

\textbf{Algebraization and related barriers:} Target model-dependent
techniques. We unify lower bounds across paradigms (backtracking, DP,
OBDDs, resolution) via a single principle---the Semantic Conservation
Law q + \(\Phi\) \(\geq\) R---deriving bounds from structural
correctness requirements rather than model-specific analyses.

\paragraph{1.1 The Real Question: Compressing Exponential
Configurations}\label{11-the-real-question-compressing-exponential-configurations}

\textbf{The Fundamental Challenge}: Exponential search spaces pervade
NP-complete problems. For example, Hamiltonian Path requires considering
up to n! orderings in the worst case. NP-complete problems universally
exhibit exponential solution spaces:

\textbf{Common NP-complete problems: exponential search spaces}

\begin{itemize}
\tightlist
\item
  \textbf{3-SAT}: 2\^{}n possible assignments to explore
\item
  \textbf{Graph Coloring}: k\^{}n possible colorings to consider
\item
  \textbf{Subset Sum}: 2\^{}n subsets to check
\item
  \textbf{Hamiltonian Path}: n! orderings to enumerate
\end{itemize}

Yet verification is polynomial: checking a 3-SAT assignment takes O(n)
clause evaluations, and verifying a graph coloring requires O(n) edge
checks. The P vs NP question asks whether algorithmic
techniques---pruning, inference, memoization, branching---can compress
these exponential configuration spaces into polynomially many
simultaneously distinguishable computational states, independent of
representation. (We formalize "state" below via \(\Phi\), the log\(_2\)
of simultaneously distinguishable maintained artifacts.)

This compression challenge breaks into two parts:

\textbf{Question 1: What kind of exponentiality must be compressed?}

Traditional lower bound techniques prove exponential resource
requirements within restricted computational models: how many circuit
gates, how many OBDD nodes, how much memory. Some prior work has
analyzed configurations within these models---pebbling games track
frontier configurations, OBDD bounds depend on variable orderings,
communication complexity partitions information across players---but
these configuration-based analyses remain paradigm-specific and often
encoding-sensitive.

L* lifts configuration-based hardness to a representation-invariant form
with respect to artifact/state encoding inside L*. Rather than requiring
exponential resources, L* forces algorithms to distinguish exponentially
many access patterns over polynomial resources---this
configuration-resource inequivalence is defined within the fixed L*
overlay/encoding. NP-completeness connects L* to standard reductions
(§10). L* does not invent new hardness; it abstracts the collision
principle underlying paradigm-specific bounds (pebbling frontier
configurations, OBDD variable orderings, communication partitions) into
a semantic correctness requirement.

\textbf{The exponentiality}: A polynomial-size set of resources (memory
addresses) admits exponentially many access patterns---different
permutations of how and when those resources are used. In L*, each
computational history produces a unique content-addressed seed that
determines which addresses to read at each step.

\textbf{Question 2: Can this exponential configuration space be
compressed into polynomial states?}

For L*, the answer is \textbf{no---representation-invariantly (within
L*)}. The incompressibility mechanism operates through
L*\textquotesingle{}s structural properties: different computational
histories generate different content-addressed seeds, which in turn
compute different designated memory addresses and retrieve different
values. Compression requires merging distinct histories into a single
computational state, but such merging forces the use of one seed for
multiple histories, retrieving values correct for one history but
incorrect for the others---errors the verifier detects. This
information-theoretic incompressibility (structural, via pigeonhole
principle) is established for the fixed L* overlay.

By the pigeonhole principle, collisions are unavoidable: compressing
2\^{}\(\lambda\) (where \(\lambda\) := R \(-\) q) seed-consistent
configurations into fewer states forces at least two distinct histories
to share a computational artifact, causing one seed to be used and the
wrong addresses to be read for the merged alternative; the verifier
rejects. Therefore any correct algorithm must maintain at least
2\^{}\(\lambda\) simultaneously distinguishable states (equivalently,
\(\Phi\) \(\geq\) R \(-\) q). This configuration-resource inequivalence
arises from structural properties (Injectivity, Emergence, Closure) that
persist across representations.

\textbf{What is L*?} A dependency-chained DAG with content-addressed
seeds---an NP-complete language built from 3-SAT---whose structural
constraints force exponential barriers in all three operational routes
simultaneously, while still admitting polynomial-time verification with
a witness.

\textbf{Compression perspective on P versus NP.} Does there exist an
NP-complete problem whose exponential configuration space cannot be
compressed into polynomially many simultaneously distinguishable states
under standard polynomial-time encodings? (Formalized via SCL in §1.2.)

L* answers affirmatively by construction. Its engineered structural
properties (Injectivity, Emergence, Seed-Locked Addressing) eliminate
all compression pathways: compressing 2\^{}\(\lambda\) seed-consistent
configurations into fewer states necessarily induces collisions
(misroutes) that the verifier detects (A1--A5 + RWA; see §7.2.1 and
Lemma 7.Misroute). This configuration-resource inequivalence is
representation-invariant with respect to artifact encoding for L*. Since
L* is NP-complete (§10) and this incompressibility yields an
unconditional one-way function (§9), the classical bridge (OWF
\(\Rightarrow\) FP \(\neq\) FNP \(\Rightarrow\) P \(\neq\) NP) implies P
\(\neq\) NP for uniform PPT (§10.4--§10.5).

Falsifiability (§3.6): Exhibit a representation-invariant
compression---a uniform PPT algorithm that achieves \(\lambda\)(A, x*) =
O(log n) across standard polynomial-time many-one encodings of a
canonical NP-complete problem, proving configuration space = resource
space representation-independently (Type 1 falsification \(\to\) P =
NP).

\textbf{The Mechanism: Why Algorithms Must Satisfy q + \(\Phi\) \(\geq\)
R}

We formalize the compression barrier through key quantities: R denotes
the structural information requirement; q denotes resolved information;
\(\Phi\) denotes log\(_2\) of simultaneously distinguishable maintained
artifacts ("state"); and \(\lambda\) := R \(-\) q denotes residual
uncertainty.

\textbf{Construction approach}: We design L* so that any algorithm faces
a structural correctness constraint. When computation requires
determining R bits but only q bits have been resolved, there remain
2\^{}(R-q) distinct computational possibilities. Our construction
ensures these are not equivalent---they lead to different correct
outputs (proven §6--§7):

\begin{itemize}
\tightlist
\item
  \textbf{Different possibilities have different seeds}
  (content-addressed identifiers; Injectivity property)
\item
  \textbf{Different seeds compute different designated memory addresses}
  (Keyedness property)
\item
  \textbf{Reading from wrong addresses produces wrong values}
  (detectable by the verifier)
\end{itemize}

Together these imply that maintaining fewer than 2\^{}(R-q)
simultaneously distinguishable states forces collisions that violate
correctness---formalized below as the Semantic Conservation Law (SCL).

The SCL captures the correctness obligations that L*\textquotesingle{}s
structure imposes on any algorithm; computational complexity emerges as
the necessary cost of fulfilling these obligations.

\textbf{Incompressibility Principle (informal).} No uniform algorithm
can map 2\^{}\(\lambda\) pairwise inequivalent seed-determined
computational paths to fewer than 2\^{}\(\lambda\) simultaneously
distinguishable states without producing an incorrect output. Here
``paths'' are seed-determined computational histories, and ``states''
are maintained artifacts (e.g., branches, table entries, OBDD nodes,
proof clauses). This is the pigeonhole principle applied to correctness;
violating it yields rejection rather than a speedup.

The requirement can be stated formally as the \textbf{Semantic
Conservation Law (SCL)}:

\textbf{q + \(\Phi\) \(\geq\) R} (where \(\Phi\) = log\(_2\) of
distinguishable artifacts maintained; formal proof: Theorem 7.A; q is
defined via functional determination and is operationally attributed by
RWA in an order-invariant way)

Violating this law does not yield slower computation; it yields
incorrect results. It is a mathematical necessity from L*'s structure,
not an algorithmic limitation.

\textbf{Terminology note:} The "conserved" quantity is \textbf{R} (the
structural information requirement), which is fixed by instance
properties (A3: Emergence). The "conservation law" states this
obligation must be accounted for via q + \(\Phi\) \(\geq\) R - an
inequality constraint analogous to thermodynamic laws (e.g., \(\Delta\)S
\(\geq\) 0). The term emphasizes that information obligations cannot be
eliminated, only redistributed between resolution (q) and potential
(\(\Phi\)).

L*\textquotesingle{}s structural properties establish the minimum
requirement (q + \(\Phi\) \(\geq\) R); algorithms must choose among
three operational routes to satisfy this constraint.

*\emph{Three Routes to Satisfy SCL*}

\textbf{Fundamental constraint}: Any correct algorithm must distinguish
2\^{}(R-q) computational paths (different seed-determined histories).
Compressing these into fewer states necessarily causes collisions and
errors. This yields exactly three ways to satisfy q + \(\Phi\) \(\geq\)
R:

\begin{enumerate}
\def\labelenumi{\arabic{enumi}.}
\item
  \textbf{Storage (space)}: Maintain 2\^{}(R\(-\)q) distinguishable
  states in parallel

  \begin{itemize}
  \tightlist
  \item
    Cost: \(\Phi\) \(\geq\) (R\(-\)q) bits of state information
  \item
    L*\textquotesingle{}s structural barrier: \textbf{Keyedness} ---
    different computational histories produce different designated
    addresses and cannot be merged without causing errors
  \end{itemize}
\item
  \textbf{Resolution (forward time)}: Learn which possibility is correct
  by reading

  \begin{itemize}
  \tightlist
  \item
    Cost: \(\geq\) (R\(-\)q)/B sequential time steps (B = O(1)
    bits/step)
  \item
    L*\textquotesingle{}s structural barrier: \textbf{Emergence +
    Bandwidth} --- R\_v fresh bits at each node must be explicitly read
    from designated addresses; no inference shortcuts exist
  \end{itemize}
\item
  \textbf{Elimination (backward time)}: Prune wrong possibilities by
  testing candidates
\end{enumerate}

\begin{itemize}
\tightlist
\item
  Cost: Must test exponentially many candidates (\ensuremath{\mathbb{E}}{[}tries{]} \(\geq\)
  2\^{}(\(\Omega\)(\(\lambda\))) in restart strategies; see Theorem J.1
  in Appendix J)
\item
  Blocked by: \textbf{Per-Node Antagonism + CDT} (Constraint-Digest
  Tagging; with world commitment; App. C.2.a) --- each test eliminates
  \(\leq\)1 bit; no bulk-pruning cascades
\end{itemize}

\textbf{No Fourth Way (Storage-Resolution-Elimination)}: SCL is a
two-dimensional constraint in (q, \(\Phi\)). Operationally, there are
three routes to satisfy q + \(\Phi\) \(\geq\) R: \textbf{Resolution}
(increase q by reading), \textbf{Elimination} (increase q by testing
candidates), and \textbf{Storage} (increase \(\Phi\) via simultaneously
distinguishable states).

\textbf{Why these three exhaust the possibilities:} The inequality q +
\(\Phi\) \(\geq\) R has only two variables (q, \(\Phi\)). To satisfy it,
an algorithm must either: (1) Increase q \(\to\) achieved via
\textbf{Resolution} (learning correct values through designated reads)
or \textbf{Elimination} (pruning wrong values through testing) (2)
Increase \(\Phi\) \(\to\) achieved via \textbf{Storage} (maintaining
more simultaneously distinguishable artifacts)

Any purported "fourth way" must increase q or \(\Phi\), thereby reducing
to cases (1) or (2). Attempting to avoid all three strategies would
require representing 2\^{}(R\(-\)q) distinguishable possibilities with
fewer than 2\^{}(R\(-\)q) artifacts. By the pigeonhole principle, this
forces collisions (different seeds map to the same artifact \(\to\)
wrong addresses \(\to\) incorrect outputs the verifier detects).

\textbf{Why L* Matters: Engineered Structural Incompressibility}

Other NP-complete problems may have exponential search spaces but often
allow shortcuts:

\begin{itemize}
\tightlist
\item
  \textbf{XOR-SAT}: Gaussian elimination (Resolution dimension bypasses
  exponential barrier)
\item
  \textbf{2-SAT}: Implication graphs (Elimination dimension allows
  polynomial pruning)
\item
  \textbf{Horn-SAT}: Unit propagation (partial Resolution shortcuts)
\end{itemize}

\textbf{L* is different}: It removes the polynomially \textbf{tractable
route} in each dimension. Algorithms cannot:

\begin{itemize}
\tightlist
\item
  Merge states efficiently (Storage blocked by Keyedness: merging
  different seeds causes errors)
\item
  Learn answers efficiently (Resolution blocked by Emergence +
  Bandwidth)
\item
  Eliminate candidates efficiently (Elimination blocked by Per-Node
  Antagonism + CDT)
\end{itemize}

This yields engineered incompressibility across the three routes: the
overall cost is exponential because all routes cannot be made cheap
simultaneously. The argument does not rely on the absence of clever
algorithms; it derives a correctness condition from
L*\textquotesingle{}s structure that precludes polynomial-time search
while preserving polynomial-time verification. For concrete
illustrations, see §2.7.

\textbf{The Information Barrier: Why Algorithmic Shortcuts Cannot Exist}

The fundamental insight is simple: \textbf{algorithms need information
to make progress; L* provides no information.} All efficient algorithmic
techniques---whether for SAT solving, constraint propagation, heuristic
search, or optimization---rely on extracting structural information from
the problem instance to guide search, prune candidates, or make informed
decisions. Standard 3-SAT instances provide such clues:

\begin{itemize}
\tightlist
\item
  \textbf{Unit clauses} (single-literal clauses) enable unit propagation
  (forced assignments)
\item
  \textbf{Pure literals} (variables appearing only positively or
  negatively) enable free assignments
\item
  \textbf{Clause structure} enables subsumption, resolution, and learned
  clauses (CDCL)
\item
  \textbf{Variable frequency} guides branching heuristics (VSIDS, VMTF)
\item
  \textbf{Conflict analysis} enables clause learning with exponential
  pruning
\item
  \textbf{Symmetries} enable symmetry breaking and search space
  reduction
\end{itemize}

L* systematically removes ALL these informational clues through
\textbf{seed-locked encoding} (§10.1.1 Overlay-as-Problem). The CNF
formula \(\varphi\) is not provided in plaintext but encoded as
E{[}i,p{]} = enc(lit{[}i,p{]}) \(\oplus\) R{[}i,p{]}, where mask bits
R{[}i,p{]} reside at seed-dependent addresses. Computing these addresses
requires the correct seed chain, which requires knowing the assignment
\(\alpha\), which IS the solution to \(\varphi\).

\textbf{The bootstrapping problem} (circular dependency with no entry
point):

\begin{itemize}
\tightlist
\item
  To get the information you need (decode \(\varphi\)) \(\to\) need the
  seed chain
\item
  To compute the seed chain \(\to\) need the assignment \(\alpha\)
\item
  To find \(\alpha\) \(\to\) need the information (the decoded
  \(\varphi\))
\end{itemize}

You cannot start anywhere without already having what
you\textquotesingle{}re trying to get. There is no "first step" that
doesn\textquotesingle{}t require the final answer. This is not "the
problem is hard"---this is "you must have the information before you can
get the information."

\textbf{No useful feedback from wrong guesses.} Even when trying
different assignments, L* provides no information to guide search.
Standard problems give useful feedback: trying x=5 in a search might
reveal "too high" (eliminating half the space), or a failed assignment
might reveal "clause 7 unsatisfied" (providing structural information
about variable interactions).

L*\textquotesingle{}s Frontier-Gate mechanism ensures failed attempts
yield only "digest mismatch"---telling you this exact guess is wrong but
providing zero information about which variables to change, which
direction to adjust, or which other candidates to eliminate. The WC-1
property (§6.2.8, Appendix C.2) formalizes this: each test eliminates
\(\leq\)1 candidate from 2\^{}\(\lambda\) possibilities.

No bulk pruning, no learned clauses, no conflict-driven shortcuts---just
exhaustive enumeration. Combined with the bootstrapping problem
(can\textquotesingle{}t get initial information) and the information
barrier (no structural clues), this creates perfect exponential
resistance: you must try exponentially many candidates, learning almost
nothing from each failure.

\textbf{Result}: Accessing ANY structural property of \(\varphi\)---unit
clauses, literal polarities, clause overlaps, variable
frequencies---requires computing the correct seed chain, which requires
already knowing the solution. Every algorithmic query is answered only
by "the solution would tell you that." This transforms the computational
barrier from algorithmic ("we haven\textquotesingle{}t found a fast
algorithm") to \textbf{information-theoretic} ("no information exists to
guide faster algorithms"). The exponential search requirement is not an
algorithmic limitation but a fundamental information-theoretic necessity
arising from L*\textquotesingle{}s construction (§6.1.1, §10.1.1).

\textbf{From correctness to time.} Correct algorithms on deterministic
k-tape TMs must explicitly create and maintain 2\^{}\(\lambda\)
distinguishable artifacts at bottlenecks.

\textbf{Intuition (Bandwidth).} Creating N distinct artifacts on a
k-tape TM with bandwidth B requires \(\Omega\)(N log N / B) steps; for N
= 2\^{}\(\lambda\) this is super-polynomial. See §7.3 (Lane Dichotomy)
for the formal proof and Frontier-Gate analysis.

\textbf{Application to L*:} The structural requirement for correctness
(must maintain 2\^{}\(\lambda\) distinguishable seed-consistent worlds)
translates directly into unavoidable computational cost. The rigorous
proof combines information-theoretic conservation (SCL, §7.2.1) with
per-segment work bounds (Frontier-Gate mechanism, Theorem 8.A).

\textbf{Universal manifestation.} This single structural requirement -
maintain 2\^{}\(\lambda\) distinguishable artifacts for correctness -
manifests differently across computational paradigms: as
2\^{}\(\lambda\) tree branches in backtracking, 2\^{}\(\lambda\) table
keys in dynamic programming, width-2\^{}\(\lambda\) in resolution
proofs, 2\^{}\(\lambda\) layer width in OBDDs. These are not separate
bounds demanding separate proofs; they are the same structural necessity
expressed in each paradigm\textquotesingle{}s native metric.

\textbf{Technical note:} Statements about "maintaining 2\^{}\(\lambda\)
artifacts" refer to the search/output task F\_can (producing canonical
witness W = (w, G\_\(\tau\), Dig\_\(\tau\))) or to runs whose
correctness uses overlay-bound artifacts; decision membership
\(\exists\)w: \(\varphi\)(w)=1 remains polynomial-time verifiable via
the overlay (§10.4.1). Formal bounds and model details in §1.7.

\paragraph{1.1.1 Summary: From Structural Hardness to Unconditional
Separation}\label{111-summary-from-structural-hardness-to-unconditional-separation}

What §1.1 established

\begin{itemize}
\tightlist
\item
  Where hardness comes from: configuration-diversity (2\^{}\(\lambda\)
  seed-determined access patterns over polynomial resources), not
  resource quantity.
\item
  Why algorithms cannot escape: information-theoretic incompressibility
  (pigeonhole principle); compressing below 2\^{}\(\lambda\) causes
  collisions. Equivalently, correctness requires \(\Phi\) \(\geq\) R
  \(-\) q (SCL).
\item
  How this creates exponential cost: the three operational routes
  (Storage, Resolution, Elimination) cannot be simultaneously polynomial
  on L* (cf. §1.3.3).
\item
  What makes this rigorous: Lane Dichotomy (Theorem 7.B) plus
  per-instance deterministic bounds (Theorem 8.A) yield time lower
  bounds via restart tries or rollback segments.
\end{itemize}

\textbf{Connection to main results (see abstract for full details):}

Designing L* to expose information-theoretic incompressibility (via
A1-A5 properties and pigeonhole principle) enables two results.

\textbf{First, per-instance deterministic bounds (Theorem 8.A):} Every
FG-wired (Frontier-Gate) instance requires super-polynomial
witness-finding time on any fixed run. These bounds hold per instance
and per fixed run; coin-fixing extends them to randomized settings
without distributional assumptions.

\textbf{Second, unconditional Structural OWF (§9):} Define f: r
\(\mapsto\) Plant(\(\varphi\), r) so every output satisfies the bound.
Any uniform PPT inverter with non-negligible success can be coin-fixed
to a deterministic run on some x* = f(r*); composing with the extractor
Ext produces a polynomial-time witness, contradicting the per-instance
bound on that (instance, run) pair. Therefore f is one-way.

Since L* is NP-complete (§10), the classical bridge (OWF \(\Rightarrow\)
FP \(\neq\) FNP \(\Rightarrow\) P \(\neq\) NP) completes the complexity
separation.

Broader perspective: one-way functions and computational asymmetry can
arise from structural correctness requirements (dependency geometry via
injectivity, emergence, closure). Both cryptography (algebraic trap
doors) and L* (configuration-diversity) derive hardness from unique-path
permutation spaces where compression causes collisions.

Related sections

\begin{itemize}
\tightlist
\item
  SCL proof: §7.2.1 (consolidated proof A1--A5 \(\Rightarrow\) q +
  \(\Phi\) \(\geq\) R)
\item
  Per-instance bounds: §8.A (Frontier-Gate mechanism with profile-tight
  timing)
\item
  Worked example: §2.6 (3-node toy instance with explicit seed/address
  calculations)
\item
  Philosophical context: §3.5 (structural hardness vs number-theoretic
  accidents)
\end{itemize}

\begin{center}\rule{0.5\linewidth}{0.5pt}\end{center}

\paragraph{\texorpdfstring{1.2 The Semantic Necessity:
L*\textquotesingle{}s Structural Obligations Drive
Complexity}{1.2 The Semantic Necessity: L*\textquotesingle s Structural Obligations Drive Complexity}}\label{12-the-semantic-necessity-ls-structural-obligations-drive-complexity}

§1.1 established why algorithms must satisfy the \textbf{Semantic
Conservation Law (SCL)}: q + \(\Phi\) \(\geq\) R. This is a correctness
requirement; violating it causes collisions between distinct seeds,
leading to wrong addresses and detectable errors.

Mathematically, SCL is a two-dimensional constraint in (q, \(\Phi\)),
yielding exactly three operational routes: \textbf{Storage (maintaining
2\^{}\(\lambda\) distinguishable states), Resolution} (learning correct
answers via reading), and \textbf{Elimination} (pruning wrong candidates
via testing).

L* blocks all three routes simultaneously. Keyedness blocks Storage,
emergence plus bandwidth constraints block Resolution, and per-node
antagonism plus CDT block Elimination; avoiding cost in one route forces
exponential payment in another.

We now provide rigorous mathematical definitions for SCL and show how
each way to satisfy it maps to the inequality's terms.

\textbf{What ``Semantic'' Means:} We analyze what correctness requires,
not what machines do. Semantic = structural correctness obligations¹ ---
the simultaneously distinguishable computational artifacts any correct
algorithm must maintain to avoid mistakes. This contrasts with syntactic
analysis (machine-specific rules like tape moves or circuit gates) or
Shannon-theoretic analysis (statistical average-case information flow).
Our focus is worst-case structural requirements for correctness. (For
technical comparison with Shannon theory, see §1.6.)

¹ Terminology note: This use of ``semantic'' concerns correctness
requirements (what must be distinguished), analogous to but distinct
from truth-theoretic semantics in logic. Both involve matching
representation to reality---truth values in logic; correctness
obligations in computation.

\textbf{The Framework in Three Parts:}

\textbf{1. The Semantic Conservation Law (SCL):} q + \(\Phi\) \(\geq\) R
(formal proof: Theorem 7.A; §7.2.1)

\textbf{"Resolved + Potential \(\geq\) Required"}

\begin{itemize}
\tightlist
\item
  q = information resolved so far (via Resolution or Elimination
  dimensions)
\item
  \(\Phi\) = log\(_2\)(states) tracking unresolved possibilities
  (Storage dimension: states = simultaneously distinguishable artifacts
  across the bottleneck)
\item
  R = total information that must be determined
\end{itemize}

Algorithms must increase q through \textbf{Resolution} (learning via
reading) or \textbf{Elimination} (pruning via testing), or increase
\(\Phi\) through \textbf{Storage} (maintaining distinguishable states).
These correspond to the three dimensions from §1.1. Maintaining fewer
states leads to collisions and incorrect answers. Recall from §1.1: when
distinct worlds (seeds) are merged, designated addresses are computed
incorrectly, leading to wrong reads and outputs; the verifier detects
these errors. This is a correctness requirement, not a complexity
constraint.

\textbf{2. Global bottleneck:} the incompressible residual at the worst
cut

\begin{itemize}
\tightlist
\item
  Computation can be viewed as flowing through a DAG from start to
  finish.
\item
  At any cut (a set of nodes separating start from finish), unresolved
  information accumulates.
\item
  The bottleneck cut with minimum residual \(\lambda\) determines
  overall complexity.
\item
  Any algorithm solving L* must account for at least 2\^{}\(\lambda\)
  simultaneously distinguishable artifacts for correctness; maintaining
  fewer forces collisions and wrong outputs (either within a single run
  or in expectation across restarts; see Theorem J.1, Appendix J).
\end{itemize}

\textbf{3. Complexity Spectrum:}

\begin{itemize}
\tightlist
\item
  \textbf{\(\lambda\) = 0} \(\to\) Polynomial
\item
  \textbf{\(\lambda\) = \(\Theta\)(log n\_core)} \(\to\) Polynomial
\item
  \textbf{\(\lambda\) = \(\Theta\)(log² n)} \(\to\) Quasi-polynomial
  (n\^{}(\(\Theta\)(log n)))
\item
  \textbf{\(\lambda\) = \(\Theta\)(n)} \(\to\) Exponential
  (2\^{}(\(\Theta\)(n)))
\end{itemize}

Key point: algorithms choose how much to resolve (q\_v), but L*'s
structure determines the consequences. Insufficient resolution within
polynomial time makes exponential state tracking structurally necessary
for correctness---algorithms must maintain at least 2\^{}(R-q)
distinguishable artifacts or produce wrong outputs (via the collision
mechanism). We refer informally to the bottleneck where the unresolved
requirement is largest; precise cut notation appears in §2.4.3 and §4.2.

\paragraph{1.2.1 How to Satisfy SCL: The Three
Ways}\label{121-how-to-satisfy-scl-the-three-ways}

The Semantic Conservation Law \textbf{q + \(\Phi\) \(\geq\) R} offers
exactly three ways to satisfy it. Each way maps precisely to terms in
the inequality:

\textbf{Dimension 1: Storage (space) \(\to\) \(\Phi\) term}

\begin{itemize}
\tightlist
\item
  \textbf{What it represents}: log\(_2\) of simultaneously
  distinguishable states maintained at the bottleneck
\item
  \textbf{Cost}: Must maintain \(\Phi\) \(\geq\) (R - q) bits of state
  information
\item
  \textbf{Manifestations}: dynamic-programming keys, OBDD layer-width
  units, backtracking-frontier elements
\item
  \textbf{L*'s block}: Keyedness (address injectivity across histories)
  ensures different computational histories produce different designated
  addresses; states cannot merge without causing errors.
\end{itemize}

\textbf{Dimension 2: Resolution (time-forward) \(\to\) +q contribution}

\begin{itemize}
\tightlist
\item
  \textbf{What it represents}: Bits resolved by learning correct answers
  through explicit reading
\item
  \textbf{Cost}: \(\geq\) (R - q)/B sequential time steps (B = O(1)
  bits/step bandwidth)
\item
  \textbf{Mechanism}: Reading designated information from the
  content-addressed overlay
\item
  \textbf{L*'s block}: emergence and bandwidth constraints ensure R\_v
  fresh bits must be explicitly read from designated addresses; no
  inference shortcuts exist.
\end{itemize}

\textbf{Dimension 3: Elimination (time-backward) \(\to\) +q
contribution}

\begin{itemize}
\tightlist
\item
  \textbf{What it represents}: Bits resolved by pruning wrong candidates
  through testing
\item
  \textbf{Cost}: Must test exponentially many candidates (\ensuremath{\mathbb{E}}{[}tries{]}
  \(\geq\) 2\^{}(\(\Omega\)(\(\lambda\))) in restart strategies; see
  Theorem J.1 in Appendix J)
\item
  \textbf{Mechanism}: Trial-and-error exploration across the exponential
  search space
\item
  \textbf{L*'s block}: per-node antagonism and CDT (Constraint-Digest
  Tagging with world commitment; App. C.2.a) ensure each test eliminates
  at most one bit; no bulk-pruning cascades.
\end{itemize}

\textbf{The Bottleneck Residual \(\lambda\):}

The residual \(\lambda\) = R \(-\) q at the bottleneck cut is the
incompressibility that determines complexity. This quantity is
operational: correctness requires maintaining at least 2\^{}\(\lambda\)
distinguishable artifacts across the bottleneck; maintaining fewer
causes collisions between inequivalent worlds, leading to wrong seeds,
wrong addresses, and incorrect outputs. This yields the runtime bound on
deterministic k-tape TMs:

time \(\geq\) 2\^{}(\(\Omega\)(\(\lambda\))) × (baseline cost per state)

where ``state'' denotes a simultaneously distinguishable artifact across
the bottleneck (see Theorem 7.B for hypotheses and constants).

Here, RWA bounds per-step legitimate inflow B = O(1) in our TM model,
and FG (Frontier-Gate) segment accounting determines the per-state
baseline (see §7.3/Appendix C).

Boundary: \(\lambda\) = O(log n) yields polynomial 2\^{}\(\lambda\);
\(\lambda\) = \(\omega\)(log n) yields super-polynomial
2\^{}\(\lambda\). For L*, \(\lambda\) = \(\Theta\)(log² n) (QP-sharp) or
\(\Theta\)(n) (flat), which are super-polynomial.

\textbf{Complexity Spectrum (rule of thumb):}

\begin{itemize}
\tightlist
\item
  \textbf{\(\lambda\) = 0} \(\to\) Polynomial (verification with
  witness)
\item
  \textbf{\(\lambda\) = O(log n\_core)} \(\to\) Polynomial
\item
  \textbf{\(\lambda\) = \(\Theta\)(log² n)} \(\to\) Quasi-polynomial
  (L*\textquotesingle{}s QP-sharp profile)

  \begin{itemize}
  \tightlist
  \item
    Bound: n\^{}(\(\Theta\)(log n))
  \end{itemize}
\item
  \textbf{\(\lambda\) = \(\Theta\)(\ensuremath{\sqrt{n}})} \(\to\) Subexponential

  \begin{itemize}
  \tightlist
  \item
    Bound: 2\^{}(\(\Theta\)(\ensuremath{\sqrt{n}}))
  \end{itemize}
\item
  \textbf{\(\lambda\) = \(\Theta\)(n)} \(\to\) Exponential
  (L*\textquotesingle{}s flat profile)

  \begin{itemize}
  \tightlist
  \item
    Bound: 2\^{}(\(\Theta\)(n))
  \end{itemize}
\end{itemize}

Profiles: ``QP-sharp'' refers to \(\lambda\) = \(\Theta\)(log² n\_core);
``flat'' to \(\lambda\) = \(\Theta\)(n).

\textbf{Why Algorithms Succeed or Fail:}

Modern solvers reduce \(\lambda\) through the three ways:

\begin{itemize}
\tightlist
\item
  \textbf{Storage techniques}: Exploit symmetry, decomposition to reduce
  branching
\item
  \textbf{Resolution techniques}: Propagation, learning to increase q
  faster
\item
  \textbf{Elimination techniques}: Smart branching, conflict analysis to
  prune efficiently
\end{itemize}

When problem structure permits all three ways to work in polynomial
time, \(\lambda\) stays at O(log n\_core), yielding polynomial time.
When the structure (as in L*) forces at least one way to be exponential,
\(\lambda\) remains large; SCL then enforces exponential growth:
algorithms must maintain 2\^{}\(\lambda\) distinguishable artifacts for
correctness, and maintaining fewer causes collisions and errors (§1.1).

\textbf{Incompressibility from Structure:}

We construct L* with specific structural properties (emergence,
completeness, dependency, injectivity; see §6) that ensure \(\lambda\) =
\(\Theta\)(log² n) or \(\Theta\)(n), depending on profile. These
properties create a mathematical necessity for correctness: any
algorithm solving L* must account for at least 2\^{}\(\lambda\)
simultaneously distinguishable artifacts across the bottleneck;
maintaining fewer causes collisions between worlds with different seeds
and thus incorrect outputs. The three dimensions (Storage, Resolution,
Elimination) face exponential barriers simultaneously, yielding a
provable gap: verification (with witness) achieves \(\lambda\) = 0,
whereas search faces \(\lambda\) = \(\Theta\)(log² n) or \(\Theta\)(n),
forcing super-polynomial complexity from structural correctness
requirements rather than algorithmic limitations.

\textbf{Connection to main results:} The per-instance deterministic
bounds (Theorem 8.A) enable the unconditional Structural OWF
construction (§9: every FG-wired instance hard for every uniform
algorithm). For quantifier structure and proof path details, see the
abstract and §1.1.1.

\paragraph{1.3 How the Three Ways Determine
Complexity}\label{13-how-the-three-ways-determine-complexity}

The three-way framework we develop rigorously for L* (§1.1--§1.2, proven
§6--9) also provides a conceptual lens for understanding P vs NP
complexity more broadly. We establish these results formally for L* and
explore how the framework illuminates complexity differences across
problem classes. For L* specifically, the residual \(\lambda\) = R \(-\)
q at the bottleneck (formal \(\lambda\)(A, x) in §2.4.3) determines
complexity:

\begin{itemize}
\tightlist
\item
  \(\lambda\) = O(log n) \(\to\) polynomial time (P)
\item
  \(\lambda\) = \(\Theta\)(log² n) \(\to\) quasi-polynomial time (as in
  L*)
\item
  \(\lambda\) = \(\Theta\)(n) \(\to\) exponential time
\end{itemize}

\subparagraph{1.3.1 P Problems --- All Three Ways
Polynomial}\label{131-p-problems--all-three-ways-polynomial}

P problems\textquotesingle{} polynomial-time solvability corresponds to
structures that allow all three ways to satisfy SCL---Storage,
Resolution, and Elimination---to execute in polynomial time. This drives
the residual \(\lambda\) = R \(-\) q to O(log n), making the required
2\^{}\(\lambda\) artifacts polynomial in size.

\textbf{How P Problems Achieve Small \(\lambda\):}

\textbf{Resolution-Friendly Structure (Polynomial Resolution):}

\textbf{Examples from known P problems:}

\begin{itemize}
\tightlist
\item
  \textbf{Maximum flow}: Resolves residual capacities along augmenting
  paths---each augmentation commits flow permanently.
\item
  \textbf{Linear programming}: Resolves feasible vertex/face
  decisions---each pivot locks in structure.
\item
  \textbf{Maximum matching}: Resolves match status via augmenting
  paths---each augmentation is locally verifiable and final.
\item
  \textbf{Shortest paths}: Resolves distance labels through
  relaxations---optimal substructure enables safe commitments.
\end{itemize}

Illustration---monotone incremental commitment: each local decision
(pivot, augmentation, relaxation) permanently resolves constraints
without introducing new conflicts. This enables incremental resolution:
q increases steadily toward R without branching, driving \(\lambda\) = R
\(-\) q to O(log n). For L* (proven §6--§8), the structure precludes
such incremental commitments.

\textbf{Why Storage and Elimination Are Also Polynomial:}

\begin{itemize}
\tightlist
\item
  \textbf{Storage}: Only poly(n) states are needed at the bottleneck
  when \(\lambda\) = O(log n).
\item
  \textbf{Elimination}: Few candidates must be tested; wrong choices are
  ruled out through local verification.
\end{itemize}

This perspective suggests that many P problems have cooperative
constraints---satisfying one constraint does not introduce
conflicts---so the three routes operate efficiently, keeping \(\lambda\)
small as a consequence of structure rather than algorithmic cleverness.

\subparagraph{1.3.2 NP-complete Problems --- At Least One Way
Exponential}\label{132-npcomplete-problems--at-least-one-way-exponential}

NP-complete problems\textquotesingle{} structural properties
(antagonistic constraints) force at least one of the three
routes---Storage, Resolution, or Elimination---to be exponential. This
keeps \(\lambda\) = \(\omega\)(log n), requiring 2\^{}\(\lambda\)
super-polynomially many artifacts. For L* specifically (proven §6--§9),
we establish this rigorously: L* is NP-complete and demonstrably
requires super-polynomial time.

\textbf{Why NP-complete Problems Have Large \(\lambda\):}

\textbf{Antagonistic Constraints:} Unlike P problems, NP-complete
problems have \textbf{constraints that conflict} - satisfying one
constraint can invalidate others. This conflict prevents incremental
resolution and maintains large residual:

\begin{itemize}
\tightlist
\item
  \textbf{3-SAT:} Clauses conflict---satisfying one can break others
  \(\to\) lacks incremental commitment \(\to\) must branch \(\to\)
  maintains many states (Storage exponential)
\item
  \textbf{Hamiltonian Path:} Global tour constraints conflict with local
  edge choices; valid segments can block completion, so many candidates
  must be tested (Elimination exponential).
\item
  \textbf{Graph coloring:} Each edge creates neighbor conflicts; choices
  cascade unpredictably, so many possibilities must be explored
  (Resolution exponential).
\end{itemize}

\textbf{The Key Difference:} Lack of incremental commitment---each local
decision can conflict with future requirements---forces algorithms to
either:

\begin{enumerate}
\def\labelenumi{\arabic{enumi}.}
\tightlist
\item
  \textbf{Maintain many states} (Storage exponential) to track
  conflicting possibilities
\item
  \textbf{Read extensively} (Resolution exponential) to determine which
  commitments are safe
\item
  \textbf{Test many candidates} (Elimination exponential) because wrong
  choices reveal little
\end{enumerate}

\textbf{Why \(\lambda\) Remains Large:} Antagonistic constraints prevent
q from approaching R within polynomial time; the three ways do not
cooperate efficiently. The residual \(\lambda\) = R \(-\) q remains
\(\omega\)(log n), so correctness requires 2\^{}\(\lambda\) =
super-polynomially many artifacts. This follows from structural
correctness, not from algorithmic limitations.

\subparagraph{1.3.3 L* --- No Simultaneous Polynomial
Bounds}\label{133-l--no-simultaneous-polynomial-bounds}

Unlike typical NP-complete problems that effectively block one or two
routes, L*'s structure precludes a polynomial-time solution along each
of the three routes. No algorithm can keep Storage, Resolution, and
Elimination simultaneously small; reducing the cost of one forces an
increase in others. With \(\lambda\) calibrated to \(\Theta\)(log² n)
(QP-sharp) or \(\Theta\)(n) (flat), the overall cost is exponential.

\textbf{How L* Blocks Each Way:}

\textbf{1. Storage Blocked (Keyedness):}

\begin{itemize}
\tightlist
\item
  \textbf{Property}: Seeds encode history injectively: Seed\_v = Enc(v
  \textbar{}\textbar{} parent\_tuples \textbar{}\textbar{}
  GateDigest\_v)
\item
  \textbf{Effect}: Different histories \(\to\) different seeds \(\to\)
  different designated addresses (§4.4)
\item
  \textbf{Consequence}: Cannot merge states---merging produces wrong
  addresses \(\to\) errors (§1.1)
\item
  \textbf{Result}: Must maintain \(\geq\) 2\^{}\(\lambda\)
  distinguishable states for correctness at bottleneck
\end{itemize}

\textbf{2. Resolution Blocked (Emergence + Bandwidth):}

\begin{itemize}
\tightlist
\item
  \textbf{Property}: Each node introduces R\_v fresh bits that cannot be
  inferred (A3: Emergence; §6.2.5)
\item
  \textbf{Bandwidth limit}: Can read only B = O(1) bits per step (Lemma
  5.5.1).
\item
  \textbf{Effect}: Cannot predict answers---must explicitly read;
  reading is slow.
\item
  \textbf{Result}: The single-run lane requires m\_seg \(\geq\)
  2\^{}(\(\rho\)-s) segments, each costing \(\Omega\)(n/W\_min) steps.
\end{itemize}

\textbf{3. Elimination Blocked (Per-node + CDT):}

\begin{itemize}
\tightlist
\item
  \textbf{Property}: Testing wrong candidates eliminates \(\leq\)1 bit
  per rejection (§6.1.1.A; factoring-style structure)
\item
  \textbf{CDT}: No ``free'' semantic progress from conflicts (Lemma
  CDT-1\textquotesingle{}; Appendix C).
\item
  \textbf{Effect}: Cannot prune efficiently---wrong guesses reveal
  little about right answers.
\item
  \textbf{Result}: The restart lane requires \ensuremath{\mathbb{E}}{[}tries{]} \(\geq\)
  2\^{}(\(\Delta\)(C*)) independent attempts.
\end{itemize}

\textbf{Why All Three Simultaneously:}

\textbf{Independence of Obstacles:} The three ways are
\textbf{orthogonal} - improving one does not help the others:

\begin{itemize}
\tightlist
\item
  Efficient storage does not make reading faster or pruning cheaper
\item
  Fast reading does not reduce storage needs or improve pruning
\item
  Efficient pruning does not reduce storage or speed up reading
\end{itemize}

\textbf{No fourth way:} These are the only mathematical avenues to
satisfy q + \(\Phi\) \(\geq\) R. Any attempt to bypass all three would
require maintaining fewer than 2\^{}(R-q) artifacts, which induces
collisions that violate verification (§1.1).

\textbf{L*\textquotesingle{}s Construction:} We design L* to block all
three routes simultaneously: it enforces exponential barriers in
Storage, Resolution, and Elimination, leaving no polynomial-time escape
route among these three operational dimensions.

\textbf{§1.3 Summary:} Three operational routes (Storage, Resolution,
Elimination) are the only ways to satisfy SCL\textquotesingle{}s q +
\(\Phi\) \(\geq\) R. P problems allow all three in polynomial time
\(\to\) \(\lambda\) = O(log n); NP-complete problems block at least one
\(\to\) \(\lambda\) = \(\omega\)(log n). For L* specifically (proven
§6--§9), all three are blocked simultaneously via Keyedness, Emergence +
bandwidth, and per-node antagonism + CDT \(\to\) \(\lambda\) \(\in\)
\{\(\Theta\)(log² n), \(\Theta\)(n)\} \(\to\) unavoidable
super-polynomial cost. See §7.2.1 (SCL proof), §8.A (per-instance
bounds), Appendix C (lane dichotomy).

\paragraph{1.4 The Search-Verification
Gap}\label{14-the-search-verification-gap}

This three-way framework explains the fundamental asymmetry between
verification and search:

\textbf{Verification (with witness W):}

\begin{itemize}
\tightlist
\item
  \textbf{Resolution}: Witness provides all answers directly \(\to\) q =
  R instantly
\item
  \textbf{Storage}: No branching needed \(\to\) only one state
  maintained (\(\Phi\) = 0)
\item
  \textbf{Elimination}: No candidates to test \(\to\) no pruning needed
\item
  \textbf{Result}: \(\lambda\) = R \(-\) q = 0; time is polynomial.
\end{itemize}

\textbf{Search (without witness):}

\begin{itemize}
\tightlist
\item
  \textbf{Resolution}: Must explicitly read to learn answers \(\to\)
  slow progress on q
\item
  \textbf{Storage}: Must maintain 2\^{}\(\lambda\) states for unresolved
  possibilities \(\to\) exponential space
\item
  \textbf{Elimination}: Must test candidates to rule out wrongs \(\to\)
  exponential attempts
\item
  \textbf{Result}: \(\lambda\) = R \(-\) q remains large; time is
  exponential.
\end{itemize}

\textbf{Key insight:} A witness collapses all three operational routes:
it provides resolution (answers), removes branching (no storage growth),
and obviates elimination (no candidates to test). Without a witness,
L*'s structure forces exponential cost across the three routes.

\textbf{Note on practical strategies:} While approximation, restriction,
parallelization, and heuristics may help in practice for other problems,
for L* these cannot reduce \(\lambda\) below the super-polynomial
threshold - the structural barrier is fundamental. (Discussion of
practical algorithmic approaches: §12.)

\textbf{P vs NP via three ways (interpretive framework):} In this
framework, P problems maintain \(\lambda\) = O(log n) because their
structure permits the three ways to succeed in polynomial time. For L*
specifically (proven §6--§9), its structure demands \(\lambda\) =
\(\omega\)(log n), forcing at least one way to be exponential; this is a
structural correctness requirement, not an algorithmic limitation.

\paragraph{1.5 Building on Existing Theory: Unifying the Structural
Perspective}\label{15-building-on-existing-theory-unifying-the-structural-perspective}

\textbf{A Unifying Perspective on Lower Bounds (via L*)}

Exponential lower bound results exist across paradigms: Resolution
requires width to handle certain formulas, OBDDs suffer from width
bottlenecks, and backtracking algorithms require exponential-sized
search trees for hard instances. While these results employ distinct
proof methodologies, they share a common underlying structure.

\textbf{This Work\textquotesingle{}s Two-Part Contribution} (see §11.4
for complete paradigm catalog):

\begin{enumerate}
\def\labelenumi{\arabic{enumi}.}
\tightlist
\item
  \textbf{Articulates common structure} across prior techniques via SCL
  (q + \(\Phi\) \(\geq\) R): 5 lower bound techniques (decision trees,
  communication complexity, pebbling, branching programs/OBDD,
  resolution) + 4 algorithmic paradigms (backtracking, DP, CDCL,
  streaming)
\item
  \textbf{Formalizes TM observation paradigm}: bits observed = q,
  configs visited = 2\^{}\(\Phi\) --- bridges SCL to information theory,
  enabling unconditional P\(\neq\)NP
\end{enumerate}

To illustrate this unity via L*, consider the pigeonhole principle:
placing n+1 objects into n containers necessitates that at least one
container holds multiple objects. This fundamental counting argument
manifests across computational paradigms under different terminology -
Resolution tracks "clause width," OBDDs measure "node count,"
backtracking algorithms maintain "branch count" - yet all quantify the
same abstract notion: the minimum number of distinguishable
computational states required for correctness. Our construction of L*
crystallizes this correspondence.

\textbf{Formal Correspondence via Projection Templates:} The language L*
enables precise mappings between paradigm-specific resources and our
unified measure of distinguishable artifacts (Alt):

\begin{itemize}
\tightlist
\item
  Backtracking branches correspond to distinct partial assignments that
  cannot coalesce without sacrificing correctness (§7.3.1)
\item
  Dynamic programming keys represent irreducible subproblem states
  requiring independent storage (§7.3.3)
\item
  OBDD nodes encode residual Boolean subfunctions necessitating distinct
  representations (Appendix B)
\item
  Resolution width captures the minimum variable set requiring
  simultaneous tracking (Appendix G)
\end{itemize}

Our technical framework provides explicit "adapters" that formalize
these correspondences. The central observation is that exponential lower
bounds all reflect the same phenomenon: an incompressible information
bottleneck of size 2\^{}\(\lambda\) arising from structural correctness
requirements. Through rigorous proof, we demonstrate that
L*\textquotesingle{}s structural properties (A1-A5) mathematically imply
the Semantic Conservation Law (SCL) - showing that q + \(\Phi\) \(\geq\)
R must hold for correctness as a consequence of these properties.

\textbf{The Core Structural Insight:}

At its heart, this paper explores one fundamental observation:
\textbf{L*\textquotesingle{}s mathematical structure imposes unavoidable
requirements for maintaining correctness during computation}. Just as a
sudoku puzzle\textquotesingle{}s constraints determine what information
must be tracked regardless of solving strategy, L*\textquotesingle{}s
structure dictates what any correct algorithm must maintain.

This single structural insight naturally manifests in multiple ways:

\begin{itemize}
\item
  \textbf{As a Conservation Law}: The requirement crystallizes
  mathematically as q + \(\Phi\) \(\geq\) R - information resolved plus
  potential maintained must meet structural requirements. This is not a
  new physical law but rather a formal accounting of
  L*\textquotesingle{}s correctness requirements.
\item
  \textbf{Across Computational Paradigms}: Whether using backtracking,
  dynamic programming, or resolution, any algorithm solving L* faces the
  same structural requirements. The conservation law is not about
  algorithmic limitations - it is about what L*\textquotesingle{}s
  structure demands for correctness.
\item
  \textbf{Through Instance Construction}: L* is explicitly designed to
  make these structural requirements observable and provable. Its DAG
  overlay with content-addressed seeds creates the precise dependencies
  that force the conservation law.
\item
  \textbf{Via Incompressibility}: The structural requirements create an
  irreducible information core - \(\geq\) 2\^{}\(\lambda\) states must
  remain distinguishable for correctness; maintaining fewer causes
  collisions \(\to\) errors (§1.1).
\end{itemize}

\textbf{Complementary perspectives:}

\begin{itemize}
\tightlist
\item
  Traditional view: algorithms cannot solve certain NP-complete problems
  efficiently (remains true).
\item
  Structural addition: for L*, this inefficiency arises from unavoidable
  structural correctness requirements (explains why).
\item
  Unified understanding: together, these perspectives provide both the
  ``what'' and the ``why'' of computational complexity.
\end{itemize}

\textbf{Dual perspective:}

Structure determines the floor; algorithms determine how close they get
to it. Both views are correct - the difference is perspective:

L*'s structural properties set a minimum requirement (any correct solver
with fewer than \(\lambda\) new bits resolved at the bottleneck must
account for \(\geq\) 2\^{}\(\lambda\) simultaneously distinguishable
artifacts), while different algorithms reach this floor via different
computational metrics---backtracking explores 2\^{}\(\lambda\) branches,
dynamic programming maintains 2\^{}\(\lambda\) keys, OBDDs have width
2\^{}\(\lambda\). The SCL captures the common constraint: q + \(\Phi\)
\(\geq\) R, with residual \(\lambda\) = R \(-\) q determining
complexity. (For technical details on cuts and expected tries, see
Appendix J: Theorem J.1 and Lemma J.1-Cart.)

\paragraph{1.5.1 How the Conservation Law Manifests Across Paradigms
(for
L*)}\label{151-how-the-conservation-law-manifests-across-paradigms-for-l}

\textbf{How the Conservation Law Manifests When Solving L*:}

When solving our constructed language L*, the Semantic Conservation Law
q + \(\Phi\) \(\geq\) R (where \(\Phi\) = log\(_2\)(Alt)) manifests
across all classical sequential models we analyze - each paradigm
measures the required distinguishable artifacts in its own units:

\textbf{How the conservation law manifests across paradigms (solving
L*)}

\begin{itemize}
\tightlist
\item
  \textbf{Backtracking}: Artifact: tree branches; counted unit: each
  branch = independent partial world; lower bound for L*: tree size
  \(\geq\) 2\^{}(\(\Omega\)(\(\lambda\)))
\item
  \textbf{Dynamic programming}: Artifact: table keys; counted unit: each
  key = distinct subproblem state; lower bound for L*: keys \(\geq\)
  2\^{}(\(\Omega\)(\(\lambda\)))
\item
  \textbf{OBDD/BDD}: Artifact: diagram width; counted unit: each node =
  residual subfunction; lower bound for L*: width \(\geq\)
  2\^{}(\(\Omega\)(\(\lambda\))) (order-robust with expander-parity)
\item
  **Resolution/CDCL**: Artifact: clause width \(\to\) size; counted
  unit: width = simultaneous variable tracking; lower bound for L*:
  proof size \(\geq\) 2\^{}(\(\Omega\)(\(\lambda\))) (via
  width\(\to\)size)
\item
  \textbf{Communication}: Artifact: monochromatic rectangles; counted
  unit: each rectangle = uniform protocol region; lower bound for L*:
  rectangles \(\geq\) 2\^{}(\(\Omega\)(\(\lambda\))) (discussion in §11;
  template/out of scope)
\item
  \textbf{Streaming}: Artifact: pass complexity; counted unit: passes
  accumulate limited bits per pass; lower bound for L*: passes \(\geq\)
  \(\Omega\)(\(\lambda\)/S) (precise tradeoffs in App. C.4; brief
  summary in §11)
\end{itemize}

Where \(\lambda\) denotes the run-dependent bottleneck residual that
captures the global difficulty for L* (formally defined in §2.4.3; we
write \(\lambda\) for brevity in the table).

Interpretation: "Artifact" means a simultaneously distinguishable,
seed-consistent equivalence class at the frontier (the alternatives that
must be kept apart at the bottleneck \textbf{for correctness} - merging
causes errors). The table lists canonical representatives per paradigm
(branches, keys, width, width\(\to\)size, rectangles, passes).

\textbf{Why the Same Bound Reappears for L*:} When solving L*, these
paradigms face the same correctness requirement - maintain
2\^{}(R\_v\(-\)q\_v) simultaneously distinguishable artifacts when R\_v
bits emerge but only q\_v are resolved. Whether these artifacts appear
as tree branches, table entries, or diagram nodes is a difference in
representation, not in the underlying necessity that
L*\textquotesingle{}s structure imposes (within a try, or accounted via
the expected number of tries across restarts; see Theorem J.1 in
Appendix J).

\textbf{The Critical Insight for L*:} Our constructed language enforces
that artifacts keyed by different seed histories cannot merge (different
seeds \(\to\) different designated addresses \(\to\) merging causes
errors; see Lemma 7.Misroute). This \textbf{Keyedness} property makes
the conservation law unavoidable across the paradigms we explicitly
model. NP verifiability: the verifier checks designated addresses and
seed parses, so illegal merges are detectable (see Lemma J.1-INJ).

\paragraph{1.6 Technical Foundations: SCL vs.
Shannon}\label{16-technical-foundations-scl-vs-shannon}

\textbf{Fundamental distinction.} Shannon's information theory analyzes
distributional/average-case properties---expected information H(·) and
mutual information through channels---appropriate for coding theory and
compression. The Semantic Conservation Law (SCL), as developed for L*,
analyzes worst-case structural requirements---how many feasible
computational worlds must remain simultaneously distinguishable for
correctness on a fixed instance. These address different questions.

\textbf{SCL (as manifested in L*):} q + \(\Phi\) \(\geq\) R (``Resolved
+ Potential \(\geq\) Required'')

where \(\Phi\) = log\(_2\)(simultaneously distinguishable artifacts).
Resolved bits plus the log of maintained states must meet the total
requirement at the bottleneck. Precise formulation in §4/§7.

\textbf{Rényi entropy spectrum.} Information theory provides a family of
entropy measures (Rényi entropy, parameterized by \(\alpha\) \(\in\)
{[}0, \(\infty\){]}), suited to different analytical needs:

\begin{itemize}
\tightlist
\item
  \textbf{Hartley entropy (Rényi-0)}: H\(_0\)(X) =
  log\(_2\)\textbar X\textbar{} --- counts uniform/worst-case
  distinguishability; zero-error scenarios
\item
  \textbf{Shannon entropy (Rényi-1)}: H\(_1\)(X) = \(-\)\(\Sigma\) p(x)
  log p(x) --- measures average uncertainty over distributions; coding
  theory, compression, channel capacity
\item
  \textbf{Min-entropy (Rényi-\(\infty\))}: H\(\infty\)(X) =
  \(-\)log\(_2\)(max p(x)) --- measures guessing hardness; cryptography
\end{itemize}

\textbf{For L*, SCL corresponds to Hartley (Rényi-0).} We analyze
worst-case structural correctness: how many computational worlds must
remain distinguishable to avoid collisions \(\to\) wrong addresses
\(\to\) detectable errors (§1.1). This is zero-error, uniform
counting---no probabilities, no averaging. Shannon (Rényi-1) averages
over distributions and thus does not capture worst-case guarantees.
Hartley is the natural measure for:

\begin{itemize}
\tightlist
\item
  Zero-error and uniform analyses: decision problems, per-instance
  deterministic bounds
\item
  Distribution-free and compositional reasoning: cut-additivity enables
  min-cut analysis
\item
  Computational pricing: \(\lambda\) converts to runtime via per-run
  pricing lemmas (§7.3/Appendix C)
\end{itemize}

\textbf{Derivation vs correspondence.} The SCL bound
\textbar{}State\textbar{} \(\geq\) 2\^{}\(\lambda\) is derived via
\textbf{pigeonhole counting} (the primary argument), then recognized as
equivalent to Hartley entropy (the explanatory correspondence). We do
not "apply" Rényi-0 formulas---we prove the bound from first principles
(keyedness \(\to\) injection \(\to\) pigeonhole), and note that
log\(_2\)\textbar State\textbar{} \(\geq\) \(\lambda\) is precisely
H\(_0\) = log\(_2\)\textbar support\textbar. The counting argument is
foundational; the Rényi-0 label is descriptive.

\textbf{Where Shannon appears (limited scope).} Shannon entropy is used
only to cap per-step fresh-bit inflow B via standard tape-bandwidth
arguments (Lemma 5.5.1). The core SCL inequalities q + \(\Phi\) \(\geq\)
R use Hartley throughout, proven via combinatorial partition-refinement
(§7.2.1), not information-theoretic channels.

\textbf{Operational differences} (Shannon vs SCL):

\begin{itemize}
\item
  \textbf{Measure:}

  \begin{itemize}
  \tightlist
  \item
    Shannon: entropy H (average uncertainty; Rényi-1)
  \item
    SCL: \(\Phi\) = log\(_2\)\textbar worlds\textbar{} (zero-error
    counting; Rényi-0)
  \end{itemize}
\item
  \textbf{Conservation:}

  \begin{itemize}
  \tightlist
  \item
    Shannon: data-processing inequality (mutual information \(\leq\)
    constant)
  \item
    SCL: partition-refinement (each legitimate read shrinks by \(\leq\)
    factor 2)
  \end{itemize}
\item
  \textbf{Composition:}

  \begin{itemize}
  \tightlist
  \item
    Shannon: rate cut-sets bound channel capacity
  \item
    SCL: worlds factor across disjoint pools \(\to\) requirements add
    across cuts \(\to\) min-cut residual \(\lambda\) prices into time
  \end{itemize}
\item
  \textbf{Credit rule:}

  \begin{itemize}
  \tightlist
  \item
    Shannon: any correlation (mutual information)
  \item
    SCL: only first-use, legitimate resolution (Keyedness + Hermeticity
    + RWA)
  \end{itemize}
\end{itemize}

\textbf{Extensions Beyond L* (Future Work).} While SCL analyzes
worst-case requirements for L* corresponding to Hartley (Rényi-0),
future frameworks extending these ideas to other settings could explore:

\begin{itemize}
\tightlist
\item
  \textbf{Average-case:} Layer Shannon/Rényi-1 measures to predict
  expected hardness on distributions (§12.12-F1)
\item
  \textbf{Cryptography:} Layer min-entropy and relativized credit rules
  for advice/oracle leakage (§12.12-F2/F5)
\item
  \textbf{Approximation:} Define \(\varepsilon\)-worlds and
  \(\Phi\)\^{}(\(\varepsilon\)) for approx-SCL with controlled slack
  (§12.12-F3)
\end{itemize}

\textbf{Scope.} Throughout this paper: \textbf{classical, uniform}
models (deterministic k-tape TMs; non-relativizing, no advice/oracles
unless explicitly noted). SCL corresponds to Hartley (Rényi-0) for
worst-case structural analysis.

\textbf{§1.6 Key Distinction:} SCL (Hartley/Rényi-0) analyzes worst-case
structural correctness (how many worlds must remain distinguishable);
Shannon (Rényi-1) analyzes average-case distributional uncertainty.
Different measures for different questions - SCL for zero-error
per-instance bounds, Shannon only for bandwidth caps (Lemma 5.5.1).
Formal SCL definition and proof: §4.2, §7.2.1.

\paragraph{1.7 Scope}\label{17-scope}

\begin{itemize}
\tightlist
\item
  Model: Deterministic k-tape TM; constant k,
  \textbar{}\(\Gamma\)\textbar; no advice; no oracles. Per-step
  legitimate inflow B =
  k\(\lceil\)log\(_2\)\textbar{}\(\Gamma\)\textbar{}\(\rceil\) (Lemma
  5.5.1). SCL corresponds to Hartley/Rényi-0 (\(\Phi\) = log\(_2\) Alt);
  Shannon appears only to cap per-step bandwidth B.
\item
  One-run, no-advice transcripts: inputs are only x*; solvers must
  produce W from x* alone in a single run (no external witness, advice,
  or oracles). For randomized solvers, we fix coins when applying
  worst-case lower bounds.
\item
  Language: Results are proved for our constructed NP-complete language
  L*, not arbitrary NP problems (see §10 for NP-completeness).
\item
  Quantifiers: \(\forall\)x*\(\in\)L*\_\{FG\}, \(\forall\) uniform
  algorithm A \(\to\) per-instance deterministic bounds (Theorem 8.A)
  \(\to\) f one-way \(\to\) P \(\neq\) NP.
\item
  Randomized solvers: Coin-fixing (Yao) extends per-instance bounds to
  randomized PPT (§9.4).
\item
  What \(\Phi\) counts: \(\Phi\) = log\(_2\) Alt counts simultaneously
  distinguishable artifacts at the bottleneck.
\item
  Polynomial baselines: all ``polynomial'' statements are in
  \textbar{}x*\textbar{}. With our overlays, \textbar{}x*\textbar{} =
  \(\Theta\)(n\_core · log² n\_core) and \(\Sigma\)\_v R\_v =
  \(\Theta\)(n\_core · log n\_core); in §§1--8 we write n for n\_core
  unless noted.
\item
  Out of scope: Pure randomness generation and quantum superposition
  (maintaining 2\^{}n amplitudes with n qubits) lie outside our
  classical, uniform framework; \(\lambda\)-accounting here applies only
  to classical models. (Appendix N discusses oracle variants; we assume
  no oracles.)
\end{itemize}

\begin{center}\rule{0.5\linewidth}{0.5pt}\end{center}

\textbf{Scope and Precision of Claims.} This work establishes
incompressibility for the specific seed-locked overlay encoding of L*
constructed in §6. The term "representation-invariant" appears
throughout with two distinct technical meanings: (1) \textbf{artifact
encoding invariance} (proven herein): the cardinality of distinguishable
computational artifacts Alt\_v is independent of how individual
artifacts are internally encoded---storing seeds as binary strings
versus hexadecimal versus structured tuples does not alter the count
(Lemma 4.2.2); compression from 2\^{}\(\lambda\) to poly(n) artifacts is
impossible regardless of such representation choices within
L*\textquotesingle{}s construction. (2) \textbf{cross-reduction
stability} (future work, §12 F6): whether incompressibility persists
across arbitrary polynomial-time many-one reductions to alternative
encodings of the same NP-complete problem.

The encoding-specific incompressibility proven here suffices for
P\(\neq\)NP via standard complexity-theoretic reasoning: L* (seed-locked
encoding) is NP-complete (§10); if L* \(\in\) P, then P=NP by definition
of NP-completeness; we prove L* \(\notin\) P for this encoding;
therefore P\(\neq\)NP. This argument requires demonstrating hardness for
one specific NP-complete language, not for all possible encodings.
Cross-reduction stability would strengthen the result but is not
required for the main theorem.

\textbf{Note}: All results, claims, and properties discussed in this
paper (including the Semantic Conservation Law, structural correctness
requirements, and complexity bounds) are specifically proven for our
constructed language L*, not general statements about all computational
problems or complexity classes. Our proofs rigorously establish these
properties for L*. While we believe similar principles may apply more
broadly, our formal results are confined to L*. This focused scope
enables concrete, verifiable proofs while acknowledging the broader P vs
NP question\textquotesingle{}s complexity. The P \(\neq\) NP result
follows via: (1) per-instance deterministic bounds for L* (Theorem 8.A),
(2) unconditional Structural OWF construction (§9), (3) classical bridge
OWF \(\Rightarrow\) FP \(\neq\) FNP \(\Rightarrow\) P \(\neq\) NP
(§10.4-10.5).

\subsubsection{2. Technical Overview: The Semantic Conservation
Law}\label{2-technical-overview-the-semantic-conservation-law}

Section 1 established the intuitive foundation: algorithms solving L*
must satisfy a \textbf{correctness requirement} q + \(\Phi\) \(\geq\) R
arising from L*\textquotesingle{}s structural properties. This
requirement offers exactly three ways to satisfy it---Storage
(maintaining 2\^{}\(\lambda\) distinguishable states), Resolution
(learning via reading), and Elimination (pruning via testing)---all
exponentially costly for L* simultaneously.

\textbf{Purpose of Section 2:} We now formalize these intuitions,
providing the precise mathematical statement of the \textbf{Semantic
Conservation Law (SCL)} and establishing the notation and framework used
throughout the paper.

\textbf{From intuition to rigor:} While §1 explained \emph{why} the law
must hold (collision mechanism, structural necessity), §2 provides the
\emph{formal machinery} needed for rigorous proofs.

\textbf{Roadmap:}

\begin{itemize}
\tightlist
\item
  §2.1: Formal statement of SCL and its dual forms (logarithmic vs
  multiplicative)
\item
  §2.2: The sharp phase transition at \(\lambda\) = \(\Theta\)(log n)
  --- why complexity jumps discontinuously
\item
  §2.3: Key distinctions addressing common misconceptions
\item
  §2.4: How dependencies create multiplicative growth and bottleneck
  residuals
\item
  §2.5: Concrete 3-node DAG example demonstrating the framework
\end{itemize}

\textbf{Relation to rigorous proofs:} The formal results appear in
§§6-10: L*\textquotesingle{}s construction and properties \textbf{A1-A5}
(§6); proof that A1-A5 mathematically imply SCL (§7.2.1);
\textbf{per-instance deterministic bounds} (§8.A) with FG details
(Appendix C); Structural OWF construction (§9); NP-completeness (§10).
Section 2 provides the conceptual bridge between §1\textquotesingle{}s
intuition and these technical results.

\paragraph{2.1 The Semantic Conservation Law: Formal Statement and
Structure}\label{21-the-semantic-conservation-law-formal-statement-and-structure}

\textbf{Symbols \& Notation Quick Reference}

\textbf{Core SCL Symbols:}

\begin{itemize}
\tightlist
\item
  SCL (node form): q\_v + \(\Phi\)\_v \(\geq\) R\_v
\item
  R\_v = required bits at node v (emergence rank)
\item
  q\_v = resolved bits at node v (via RWA)
\item
  Alt\_v = simultaneously distinguishable artifacts at v
\item
  \(\Phi\)\_v = log\(_2\)(Alt\_v) (potential in bits)
\end{itemize}

\textbf{Cut Aggregates:}

\begin{itemize}
\tightlist
\item
  \(\Lambda\)(C) = \(\Sigma\)\_\{v\(\in\)C\} R\_v (total emergence at
  cut)
\item
  Q(C) = \(\Sigma\)\_\{v\(\in\)C\} q\_v (total resolution)
\item
  \(\lambda\)(C) = \(\Sigma\)\_\{v\(\in\)C\}(R\_v\(-\)q\_v) (cut
  residual)
\item
  Alt(C) = \(\prod\)\_\{v\(\in\)C\} Alt\_v (multiplicative worlds)
\item
  \(\Phi\)(C) = log\(_2\) Alt(C) = \(\Sigma\)\_\{v\(\in\)C\} \(\Phi\)\_v
\item
  C*: designated bottleneck cut
\end{itemize}

\textbf{Residual Parameters:}

\begin{itemize}
\tightlist
\item
  \(\lambda\)(A, x) = min\_C \(\lambda\)(C) (bottleneck residual,
  run-dependent)
\item
  \(\lambda\)\_base (instance/profile baseline)
\item
  \(\rho\) (residual before final segment)
\item
  s (pre-final agreement bits)
\end{itemize}

\textbf{Model Parameters:}

\begin{itemize}
\tightlist
\item
  B = k\(\lceil\)log\(_2\)\textbar{}\(\Gamma\)\textbar{}\(\rceil\)
  (fresh bits/step on k-tape TM)
\item
  W\_min (profile width)
\item
  n\_core (core problem size)
\item
  \textbar{}x*\textbar{} = \(\Theta\)(n\_core · log² n\_core) (instance
  size, profile-dependent)
\end{itemize}

\textbf{Function Notation:}

\begin{itemize}
\tightlist
\item
  Enc(·) = seed encoder (injective, parseable)
\item
  enc(·) = literal packer (lowercase)
\item
  F\_overlay(Seed\_v; j,\(\ell\)) = designated address function
\item
  m\_seg = non-accepting rollback segment count
\end{itemize}

\textbf{Common Disambiguations:}

\begin{itemize}
\tightlist
\item
  \(\Phi\) (Phi) = log\(_2\) Alt (information measure)
\item
  \(\varphi\) (phi) = CNF formula (unrelated symbol)
\item
  w = witness, \(\omega\) = feasible world, \(\eta\) = Plant randomness
\item
  U\_v = per-node address pool, c(v) = R\_v \(-\) q\_v
\end{itemize}

See Appendix Z for complete cheat sheet and precise conversions.

\subparagraph{2.1.0 What Is L*? (Quick
Overview)}\label{210-what-is-l-quick-overview}

Before formalizing the conservation law, we briefly describe the
constructed language L* to which it applies:

\textbf{A1-A5 Properties at a Glance} (full construction §6; proof
§7.2.1):

\begin{itemize}
\tightlist
\item
  \textbf{A1 (Hermeticity):} No hidden channels - information flows only
  through designated addresses in disjoint per-node pools \{U\_v\}
\item
  \textbf{A2 (Injectivity):} Distinct parent tuples \(\to\) distinct
  seeds (Enc is injective on domain)
\item
  \textbf{A3 (Emergence):} R\_v fresh bits emerge at node v (rank(H\_v)
  = R\_v forces new information)
\item
  \textbf{A4 (Closure):} Seeds deterministically recover ancestry (Enc
  is parseable)
\item
  \textbf{A5 (Dependency):} Parents complete before children (DAG
  structure)
\end{itemize}

\textbf{SCL Derivation:} A1-A5 \(\to\) Keyedness (different seeds
\(\to\) different designated addresses) \(\to\) Collision necessity
\(\to\) q + \(\Phi\) \(\geq\) R

\textbf{L*} is a constructed language (§6) proven to be NP-complete
(§10) with these key properties:

\begin{itemize}
\tightlist
\item
  \textbf{Structure}: Blockchain-inspired dependency-chained DAG with
  seed-locked decode (content-addressed; no cryptography; deterministic
  construction)
\item
  \textbf{Decision problem}: Given instance x*, does there exist witness
  W such that the decoded CNF formula \(\varphi\) is satisfiable?
\item
  \textbf{Witness W}: Complete information (seeds \{Seed\_v\},
  assignments \{y\_v\}, gate proofs G\_\(\tau\), digests Dig\_\(\tau\))
  enabling polynomial-time verification
\item
  \textbf{Verification}: Algorithm V checks witness validity in O(n²)
  time (§10.1, Theorem 10.1)
\item
  \textbf{Canonical output task F\_can}: Produce canonical witness W
  from instance x* alone (without being given W)
\item
  \textbf{Hardness result}: F\_can requires super-polynomial time:
  \(\geq\) n\^{}(\(\Omega\)(log n)) (QP-sharp profile) or \(\geq\)
  2\^{}(\(\Omega\)(n)) (flat profile)
\item
  \textbf{NP-completeness}: Witness-preserving reduction from 3-SAT
  (§10.2-10.3, Theorem 10.2)
\item
  \textbf{Main result}: FP \(\neq\) FNP \(\to\) P \(\neq\) NP for
  deterministic k-tape TMs (§10.5)
\end{itemize}

\textbf{Key Innovation - Overlay-as-Problem (OAP):}

Unlike traditional NP reductions that provide the problem formula
\(\varphi\) directly in the instance, L* \textbf{applies SCL to problem
access itself} via a seed-locked decode schema (§10.1.1). The CNF
formula is encoded as E\_lit = enc(lit) \(\oplus\) R where mask bits R
reside at designated addresses requiring GREQ=1 (gate requirement flag)
seed computation with gate digests (GateDigest\_v: identity digest wired
into seeds).

\textbf{Why decoding is hard:} Computing these decode seeds requires
evaluating gate digests at GREQ=1 nodes---work that itself is subject to
SCL. To correctly compute seeds, you must either resolve information (q)
or maintain distinguishable artifacts (\(\Phi\)), both exponentially
costly.

\textbf{Avoiding circularity:} The decode seed computation is
assignment-independent (using only instance salts/metadata and
parity-based gateDigests; Lemma 10.1.1-Adj), avoiding circularity. The
planting transform constructs instances with identity-based digests
(non-leaking: no assignment bits exposed in x*); any valid domain
element r' \(\in\) D(\(\varphi\)) with plant(r') = x* contains a
satisfying assignment by domain constraint.

\textbf{The cost barrier:} Computing the seeds and designated addresses
costs \(\Omega\)(n/W\_min) operations per segment because SCL applies to
the decoding process. This makes \textbf{standard SAT search (finding w)
already super-polynomial} because solvers must first decode \(\varphi\)
through SCL-constrained overlay work.

\textbf{Integration with FP\(\neq\)FNP:} The canonical witness W makes
this unavoidable overlay work verifiable for the FP \(\neq\) FNP
formulation. OAP is not a separate mechanism---it is the key innovation
of applying the same structural correctness requirement to discovering
what you are solving, not just solving it.

\textbf{Key mechanisms}:

\begin{itemize}
\tightlist
\item
  \textbf{Seeds}: Each node v has Seed\_v = Enc(v \textbar{}\textbar{}
  parent\_data \textbar{}\textbar{} GateDigest\_v) when required
  (Injectivity A2)
\item
  \textbf{Designated addresses}: F\_overlay(Seed\_v, j, \(\ell\))
  computes content-addressed locations (Keyedness)
\item
  \textbf{Seed-locked decode}: Mask bits R{[}i,p{]}{[}t{]} at GREQ=1
  addresses; decoding \(\varphi\) requires computing gate digests
\item
  \textbf{Emergence}: R\_v fresh bits at each node via rank(H\_v) = R\_v
  (A3)
\item
  \textbf{Dependencies}: Children\textquotesingle{}s seeds depend on
  parents\textquotesingle{} outputs (A5), creating multiplicative state
  growth
\end{itemize}

Full construction: §6; axioms \textbf{A1-A5}: Lemma 4.A (verified in
§6); NP-completeness: §10; \textbf{per-instance bounds}: §8.A.

\subparagraph{2.1.1 The ConservationLaw}\label{211-the-conservationlaw}

\textbf{Definition (Semantic Conservation Law for L*).} For any node v
in L*\textquotesingle{}s computational DAG, any correct algorithm must
satisfy:

\textbf{q\_v + \(\Phi\)\_v \(\geq\) R\_v}

where:

\begin{itemize}
\tightlist
\item
  \textbf{R\_v} = bits that must emerge at node v (structural
  requirement; determined by rank(H\_v) = R\_v per Axiom A3)
\item
  \textbf{q\_v} = bits resolved by the algorithm at node v (via
  Receiving-Window Attribution; §2.1.3)
\item
  \textbf{\(\Phi\)\_v} = log\(_2\)(Alt\_v) where Alt\_v = number of
  simultaneously distinguishable computational artifacts maintained at v
\end{itemize}

\textbf{Intuition.} The SCL states a \textbf{correctness requirement},
not a complexity constraint. At node v, resolving q\_v bits reduces
2\^{}(R\_v) initial possibilities to 2\^{}(R\_v-q\_v) remaining worlds.
Each world has a distinct seed (by Injectivity, A2) which computes
distinct designated addresses (by Keyedness - address injectivity;
defined §2.1.2 below). If an algorithm maintains fewer than
2\^{}(R\_v-q\_v) distinguishable artifacts, collision occurs: different
worlds map to the same artifact \(\to\) algorithm uses wrong seed
\(\to\) computes wrong designated addresses \(\to\) reads wrong overlay
values \(\to\) produces incorrect output \(\to\) verifier detects error.

\textbf{Violating SCL produces incorrect results, not merely slow
computation.} The distinction is fundamental: collision causes errors
detectable by the verifier, not inefficiency (collision mechanism
detailed in §1.1).

q semantics vs pricing. q is defined semantically via functional
determination (what the computation has logically fixed).
Receiving-Window Attribution (RWA; §2.1.3, §4.2) is the accounting rule
used to price when those bits are first revealed in runs. This matches
the metric note in §7.2.1 and the exposition in §1.

\begin{itemize}
\tightlist
\item
  \textbf{Note}: We use subscript \_v only when referring to specific
  nodes; the law q + \(\Phi\) \(\geq\) R holds universally across all
  nodes in L*
\end{itemize}

\subparagraph{2.1.2 The Three Ways to Satisfy
SCL}\label{212-the-three-ways-to-satisfy-scl}

As established in §1.1-1.2, the conservation law q + \(\Phi\) \(\geq\) R
offers exactly three ways to satisfy it. These are not algorithmic
strategies - they are the \textbf{only mathematical possibilities} for
handling 2\^{}(R-q) unresolved computational worlds:

\textbf{1. Storage (Space Dimension) \(\to\) \(\Phi\) term}

\begin{itemize}
\tightlist
\item
  \textbf{What it represents}: \(\Phi\) = log\(_2\)(Alt) where Alt =
  simultaneously distinguishable states at bottleneck
\item
  \textbf{Cost}: Must maintain \(\Phi\) \(\geq\) (R - q) bits of state
  information \(\to\) 2\^{}(R-q) artifacts
\item
  \textbf{Manifestations}: DP table keys, OBDD layer-width nodes,
  backtracking-tree frontier elements
\item
  \textbf{L*\textquotesingle{}s block}: \textbf{Keyedness} (address
  injectivity from A1-A2): different seed histories \(\to\) different
  designated addresses \(\to\) cannot merge without errors
\end{itemize}

\textbf{2. Resolution (Time-Forward Dimension) \(\to\) +q contribution}

\begin{itemize}
\tightlist
\item
  \textbf{What it represents}: Bits resolved by learning correct answers
  through explicit reading
\item
  \textbf{Cost}: \(\geq\) (R - q)/B sequential time steps (B = O(1)
  bits/step; Lemma 5.5.1)
\item
  \textbf{Mechanism}: Reading designated information from
  content-addressed overlay
\item
  \textbf{L*\textquotesingle{}s block}: Emergence (A3) + Bandwidth limit
  ensure R\_v fresh bits must be explicitly read \(\to\) no inference
  shortcuts
\end{itemize}

\textbf{3. Elimination (Time-Backward Dimension) \(\to\) +q
contribution}

\begin{itemize}
\tightlist
\item
  \textbf{What it represents}: Bits resolved by pruning wrong candidates
  through testing
\item
  \textbf{Cost}: Must test exponentially many candidates (\ensuremath{\mathbb{E}}{[}tries{]}
  \(\geq\) 2\^{}(\(\Omega\)(\(\lambda\))) in restart lane; see Theorem
  J.1 in Appendix J)
\item
  \textbf{Mechanism}: Trial-and-error exploration across exponential
  search space
\item
  \textbf{L*\textquotesingle{}s block}: Per-node antagonism + CDT
  (Constraint-Digest Tagging; App. C.2.a) ensure each test eliminates
  \(\leq\)1 bit \(\to\) no bulk-pruning cascades
\end{itemize}

\textbf{Why no fourth way:} Attempting to bypass all three
simultaneously would require representing 2\^{}(R-q) distinguishable
possibilities with fewer than 2\^{}(R-q) artifacts. By the pigeonhole
principle, collision is unavoidable: different seeds \(\to\) wrong
addresses \(\to\) wrong outputs (§1.1). The structure does not permit
it.

\subparagraph{2.1.3 Receiving-Window Attribution
(RWA)}\label{213-receiving-window-attribution-rwa}

\textbf{Definition.} A bit of information is charged to node v (in q\_v)
at its \textbf{first valid use} that functionally determines it on the
current seed chain. Re-reads or later uses of already-revealed
information are not re-charged.

\textbf{Operational Definition (First Valid Use).} To determine whether
a designated read at step t of bit b from address u constitutes a first
valid use for node v:

\begin{enumerate}
\def\labelenumi{\arabic{enumi}.}
\tightlist
\item
  \textbf{Parse address}: Compute u = F\_overlay(Seed\_v; j,\(\ell\)) to
  identify node v (the address determines ownership by construction via
  disjoint pools A1)
\item
  \textbf{Check prior reads}: Search transcript \(\pi\)\_\{\textless t\}
  for any previous read of bit b from address u

  \begin{itemize}
  \tightlist
  \item
    If no prior read found: \textbf{This is the first valid use for v}
    \(\to\) Credit 1 bit to q\_v
  \item
    If prior read(s) found: \textbf{Not a first use} \(\to\) Credit 0
    bits (already charged)
  \end{itemize}
\item
  \textbf{Functional determination}: The bit b becomes functionally
  determined at v at the first step where:

  \begin{itemize}
  \tightlist
  \item
    The algorithm reads b from a designated address u in
    v\textquotesingle{}s pool, AND
  \item
    The value of b had not been previously observed in the transcript
  \end{itemize}
\item
  \textbf{Re-reads and caching}: Subsequent reads of the same bit b
  (whether from the original address u, from cache, or via
  recomputation) are not charged --- the information was already counted
  at first use.
\end{enumerate}

\textbf{Key Properties}:

\begin{itemize}
\tightlist
\item
  \textbf{Unique attribution}: Each bit is charged exactly once (to the
  node whose pool contains the read address)
\item
  \textbf{Schedule-independent}: Credits depend on which bits are read,
  not when (Lemma 6.1-RWA)
\item
  \textbf{Observable}: First-use events are deterministic functions of
  the transcript (Lemma I.2.3)
\item
  \textbf{No double counting}: Hermeticity ensures designated reads are
  the only information source; re-reads do not add new information.
\end{itemize}

\textbf{Purpose.} RWA provides a representation-independent accounting
rule that works across all computational paradigms. Combined with
Hermeticity (A1), RWA ensures each counted bit corresponds to at least
one designated read from an address pool, preventing hidden information
channels.

\textbf{Technical note.} While q\_v is defined via functional
determination (which bit patterns are possible given the transcript),
\textbf{pricing} q\_v into time/space resources uses RWA to ensure
first-use designation. This distinction matters: SCL (q + \(\Phi\)
\(\geq\) R) uses functional q; time bounds use RWA-priced q. For L*,
these align (proven in Theorem I.1, Appendix I; operational details in
§5.5.1 for time conversion).

\textbf{Lemma 2.1.3-SIM (RWA/Hermeticity Model Equivalence).} For
deterministic k-tape TMs, RWA and Hermeticity are analysis conventions
that do not restrict the computational model nor change asymptotic time.

\begin{itemize}
\tightlist
\item
  (i) RWA is accounting, not restriction: it credits information at the
  earliest first-use read; reordering/re-reads cannot reduce the total
  credited q\_v across a run.
\item
  (ii) Hermeticity captures standard inputs: for L*, all fresh instance
  information enters via designated addresses by construction; there are
  no hidden channels in the input beyond designated reads.
\item
  (iii) No simulation overhead: given any TM transcript, RWA credits can
  be computed post-hoc and per-step inflow bounds (Lemma 5.5.1) still
  hold; this does not alter asymptotic time complexity.
\end{itemize}

Consequence. Lower bounds derived under RWA/Hermeticity apply to all
standard deterministic k-tape TMs solving L*.

\textbf{Lemma 2.1.3-ADDR (Schedule-measurability of addresses).} At any
time t, the multiset of designated addresses issued by the algorithm is
a function of the transcript prefix up to t together with the public
overlay. In particular, observing addresses carries no fresh information
about unresolved bits beyond what is already contained in the transcript
prefix.

\emph{Proof sketch.} Designated addresses are computed by evaluating
public functions \texttt{F\_overlay(Seed\_v;\ j,l)} along the current
seed chain. Each \texttt{Seed\_v} is a deterministic function of
previously observed primitive values and public metadata. Thus addresses
depend only on the observed transcript (plus public overlay), not on
hidden instance bits directly. By the data processing inequality,
conditioning on the transcript ensures that addresses cannot increase
mutual information with unresolved bits. \(\blacksquare\)

\textbf{Definition Block (Operational Primitives \& Verifier Checks).}

\begin{itemize}
\tightlist
\item
  Designated address: a parseable name \texttt{(pool-id=v,\ key)}
  produced by F\_overlay(Seed\_v; j,\(\ell\)) for a published index
  (j,\(\ell\)); keys map via the seed-dependent permutation into the
  disjoint pool \texttt{U\_v}.
\item
  First valid use / credited read: a bit contributes to \texttt{q\_v}
  exactly at the first designated read that functionally determines it
  on the current seed chain (RWA; see above).
\item
  Seed-consistent artifact (Keyedness): an artifact for node \texttt{v}
  is correct iff from its content the verifier can deterministically
  derive exactly the designated addresses F\_overlay(Seed\_v;
  j,\(\ell\)) on the current seed chain. Cross-seed reuse is invalid.
\item
  Collision \(\to\) error: merging two inequivalent seed histories
  forces misaddressing; the verifier recomputation detects the mismatch
  (Keyedness + Injectivity).
\item
  Verifier check (Algorithm V): Steps 4a-4d recompute seeds, designated
  addresses, and any required digests on the current seed chain and
  reject on any mismatch (polytime; see §10.2 Algorithm V).
\end{itemize}

\subparagraph{2.1.4 Dual Forms of the Conservation
Law}\label{214-dual-forms-of-the-conservation-law}

The Semantic Conservation Law has two equivalent mathematical
formulations:

\textbf{Logarithmic Form (SCL):} \textbf{q\_v + \(\Phi\)\_v \(\geq\)
R\_v} where \(\Phi\)\_v = log\(_2\)(Alt\_v)

\textbf{Multiplicative Form (SMP - Semantic Multiplication Principle):}
\textbar{}\(\Pi\)\_\{after v\}\textbar{} \(\geq\)
\textbar{}\(\Pi\)\_\{before v\}\textbar{} · 2\^{}(R\_v-q\_v)

\textbf{Equivalence.} Taking log\(_2\) of SMP yields SCL. Both express
the same necessity: at node v, 2\^{}(R\_v-q\_v) remaining worlds (after
resolving q\_v bits) must map to at least 2\^{}(R\_v-q\_v)
simultaneously distinguishable artifacts, or collision causes errors.

\textbf{When to use which:}

\begin{itemize}
\tightlist
\item
  \textbf{SCL (logarithmic)}: Easier for additive composition across
  cuts (\(\Sigma\)\_v \(\Phi\)\_v \(\geq\) \(\Sigma\)\_v (R\_v-q\_v))
\item
  \textbf{SMP (multiplicative)}: More intuitive for understanding growth
  dynamics (worlds multiply across dependencies)
\end{itemize}

Both are rigorously proven for L* in §7.2.1 (Consolidated SCL Theorem).

\subparagraph{2.1.5 Structural Foundation: A1-A5
Axioms}\label{215-structural-foundation-a1-a5-axioms}

\textbf{Where SCL comes from.} The conservation law is not assumed - it
is \textbf{mathematically derived} from L*\textquotesingle{}s structural
properties. We construct L* (§6) with specific properties:

\begin{itemize}
\tightlist
\item
  \textbf{A1 (Hermeticity)}: No hidden information channels (disjoint
  designated address pools)
\item
  \textbf{A2 (Injectivity)}: Distinct parent tuples \(\to\) distinct
  seeds (Enc injective)
\item
  \textbf{A3 (Emergence)}: R\_v fresh bits at each node v (rank forcing
  via H\_v)
\item
  \textbf{A4 (Closure)}: Seeds deterministically recover ancestors (Enc
  parseable)
\item
  \textbf{A5 (Dependency)}: Parents complete before children (DAG
  structure)
\end{itemize}

\textbf{Note on A6 (Independence)}: The Structural OWF construction (§9)
uses \textbf{A1-A5 with per-instance deterministic bounds} (Theorem
8.A), not distributional arguments.

\textbf{The proof.} Section 7.4.1 provides the complete rigorous proof
that properties A1-A5 \textbf{mathematically imply} the conservation
bound q + \(\Phi\) \(\geq\) R via Keyedness (from A1-A2: disjoint pools
+ injectivity). The non-trivial contribution is showing this precise
quantitative relationship follows necessarily from these qualitative
structural properties.

\textbf{Dual-purpose design.} The A1-A5 axioms serve two roles:

\begin{enumerate}
\def\labelenumi{\arabic{enumi}.}
\tightlist
\item
  \textbf{Create computational requirements}: L*\textquotesingle{}s
  structure forces unavoidable obstacles (Keyedness blocks Storage,
  Emergence blocks Resolution, etc.)
\item
  \textbf{Enable measurement}: A1-A5 allow representation-independent
  tracking of resolved bits (q\_v) and distinguishable states (Alt\_v)
\end{enumerate}

This design makes SCL both \textbf{provable} (we show A1-A5 \(\to\) SCL
rigorously) and \textbf{paradigm-invariant} (applies to any algorithm
solving L*, regardless of implementation).

\textbf{§2.1 Summary:} SCL (q + \(\Phi\) \(\geq\) R) is a correctness
requirement derived from A1-A5 properties, not an algorithmic
limitation. Three ways to satisfy it: Storage (maintain 2\^{}\(\lambda\)
artifacts), Resolution (read to learn), Elimination (test to prune) -
all blocked by L*. RWA tracks first-use bits; multiplicative form shows
exponential growth. Formal proof: §7.2.1 (Consolidated SCL Theorem);
per-instance bounds: §8.A; A1-A5 construction: §6.

\paragraph{2.2 The Sharp Phase Transition: A Mathematical
Cliff}\label{22-the-sharp-phase-transition-a-mathematical-cliff}

The Semantic Conservation Law creates a \textbf{discontinuous complexity
jump} at the boundary \(\lambda\) = \(\Theta\)(log n). This is not a
gradual spectrum - it is a sharp mathematical cliff separating
polynomial from super-polynomial complexity.

\subparagraph{2.2.1 The Phase Diagram}\label{221-the-phase-diagram}

The residual \(\lambda\) = R - q at the bottleneck determines complexity
through the required number of distinguishable artifacts:
2\^{}\(\lambda\). The exponential function exhibits phase-transition
behavior:

\textbf{Polynomial regime:}

\begin{itemize}
\tightlist
\item
  \textbf{\(\lambda\) = 0} \(\to\) 2\^{}0 = 1 artifacts \(\to\)
  Polynomial
\item
  \textbf{\(\lambda\) = O(log n)} \(\to\) 2\^{}(O(log n)) = poly
  artifacts \(\to\) Polynomial
\end{itemize}

\textbf{{[}SHARP DISCONTINUITY AT \(\lambda\) = \(\omega\)(log n){]}}

\textbf{Super-polynomial regime:}

\begin{itemize}
\tightlist
\item
  \textbf{\(\lambda\) = \(\omega\)(log n)} \(\to\) 2\^{}(\(\omega\)(log
  n)) \textgreater{} poly artifacts \(\to\) Super-polynomial
\item
  \textbf{\(\lambda\) = \(\Theta\)(log² n)} \(\to\) n\^{}(\(\Theta\)(log
  n)) artifacts \(\to\) Quasi-polynomial
\item
  \textbf{\(\lambda\) = \(\Theta\)(\ensuremath{\sqrt{n}})} \(\to\) 2\^{}(\(\Theta\)(\ensuremath{\sqrt{n}}))
  artifacts \(\to\) Subexponential
\item
  \textbf{\(\lambda\) = \(\Theta\)(n)} \(\to\) 2\^{}(\(\Theta\)(n))
  artifacts \(\to\) Exponential
\end{itemize}

\textbf{The boundary.} The critical threshold occurs precisely at
\(\lambda\) = \(\Theta\)(log n):

\begin{itemize}
\tightlist
\item
  \textbf{Below threshold} (\(\lambda\) = O(log n)): 2\^{}\(\lambda\) =
  2\^{}(c log n) = 2\^{}(log n\^{}c) = n\^{}c = polynomial
\item
  \textbf{Above threshold} (\(\lambda\) = \(\omega\)(log n)):
  2\^{}\(\lambda\) grows faster than any polynomial in n
\end{itemize}

There is \textbf{no smooth transition} because 2\^{}\(\lambda\) is
exponential in \(\lambda\) - the moment \(\lambda\) crosses from O(log
n) to \(\omega\)(log n), complexity jumps discontinuously from
polynomial to super-polynomial.

\subparagraph{2.2.2 Verification vs. Search Across the
Cliff}\label{222-verification-vs-search-across-the-cliff}

\textbf{Verification (with witness W):}

\begin{itemize}
\tightlist
\item
  Witness provides all seeds directly \(\to\) designated addresses
  computable immediately
\item
  All bits resolved: q\_v = R\_v at every node v
\item
  Residual at bottleneck: \(\lambda\) = 0 (no unresolved bits)
\item
  Required artifacts: 2\^{}0 = 1 (single path through DAG)
\item
  \textbf{Time complexity}: O(n²) polynomial (Theorem 10.1)
\item
  \textbf{Position}: Left side of cliff (\(\lambda\) = 0)
\end{itemize}

\textbf{Search for canonical witness (F\_can, without being given W):}

\begin{itemize}
\tightlist
\item
  Must determine seeds through search (no direct access)
\item
  Limited resolution in polynomial time: q\_v \(\ll\) R\_v at bottleneck
  nodes
\item
  Residual at bottleneck: \(\lambda\) = \(\Theta\)(log² n) or
  \(\Theta\)(n) (depending on L* profile)
\item
  Required artifacts: n\^{}(\(\Theta\)(log n)) or 2\^{}(\(\Theta\)(n))
  (QP-sharp vs. flat profile; §8.A)
\item
  \textbf{Time complexity}: Super-polynomial (§9, Structural OWF
  construction)
\item
  \textbf{Position}: Right side of cliff (\(\lambda\) = \(\omega\)(log
  n))
\end{itemize}

\textbf{The gap location.} Verification sits at \(\lambda\) = 0 (left
side of cliff), while search sits at \(\lambda\) = \(\Theta\)(log² n) or
\(\Theta\)(n) (right side of cliff). They occupy \textbf{opposite sides}
of the phase transition - there is no continuous path between them that
stays polynomial.

\begin{center}\rule{0.5\linewidth}{0.5pt}\end{center}

\textbf{Note:} For why L* avoids the \(\lambda\) = \(\Theta\)(log n)
boundary, see §8. For connections to broader P vs NP, see Appendix C (FG
details). Technical details on tight bounds: §5.5 (arity/time pricing),
Appendix C (Frontier-Gate mechanism).

\paragraph{2.3 Key Clarifications}\label{23-key-clarifications}

Several common confusions arise when first encountering the SCL
framework:

\textbf{"SCL limits what algorithms can do"} \(\to\) \textbf{No.} SCL is
a \textbf{correctness requirement}, not an algorithmic limitation. It
states what any correct algorithm must satisfy to avoid producing wrong
outputs. Violating SCL \(\to\) incorrect outputs (collision \(\to\)
wrong addresses \(\to\) verifier rejects), not "slow but correct"
outputs. Like conservation of energy in physics: a mathematical
necessity any working system must respect, not a rule forbidding clever
techniques.

\textbf{"Storage and Time are interchangeable resources"} \(\to\)
\textbf{Partially.} Both space and time arise from the same underlying
\textbf{mechanism requirement} - handling 2\^{}(R-q) unresolved
possibilities. However, you cannot escape the exponential factor by
switching resources: the mechanism requirement (SCL) is upstream of
resource costs. For L*, this manifests as exponential barriers in all
dimensions simultaneously (§7.5; §8).

\textbf{"The three dimensions (Storage, Resolution, Elimination) can be
traded off"} \(\to\) \textbf{Not for L*.} The dimensions are
\textbf{orthogonally blocked}: Storage efficiency does not make reading
faster (bandwidth B = O(1) bits/step is model-fixed); fast reading does
not reduce storage needs (Keyedness: different seeds \(\to\) different
addresses \(\to\) cannot merge); efficient pruning does not help either
(per-node antagonism: each test eliminates \(\leq\)1 bit). Many
NP-complete problems allow partial workarounds; L* is \textbf{engineered
to block all three routes simultaneously}.

\textbf{Mechanisms vs. Resources:} To satisfy q + \(\Phi\) \(\geq\) R,
algorithms can (1) RESOLVE by increasing q (via Resolution: read
designated values; or Elimination: test/prune candidates), or (2)
MAINTAIN by increasing \(\Phi\) (via Storage: keep distinguishable
states). These \textbf{mechanisms} translate to \textbf{resources} (TIME
for sequential operations, SPACE for state storage) depending on
computational model. For L*, when \(\lambda\) = \(\omega\)(log n), each
dimension enforces exponential cost (§7.5).

\paragraph{2.4 Why Dependencies Matter: From Per-Node to Global
Bounds}\label{24-why-dependencies-matter-from-per-node-to-global-bounds}

Dependencies transform local requirements into global exponential
growth. Understanding this multiplicative composition is key to seeing
why \(\lambda\) determines overall complexity.

\subparagraph{2.4.1 The Multiplicative
Principle}\label{241-the-multiplicative-principle}

At independent nodes, requirements add linearly. At dependent nodes,
\textbf{worlds multiply}:

\textbf{Independent computation:}

\begin{itemize}
\tightlist
\item
  Nodes A and B compute separately (no shared state)
\item
  Total states needed: max(2\^{}(r\_A), 2\^{}(r\_B)) \(\approx\)
  2\^{}(max(r\_A, r\_B))
\item
  Largest bottleneck dominates
\end{itemize}

\textbf{Dependent computation (DAG structure):}

\begin{itemize}
\tightlist
\item
  Node C depends on both A and B (needs their outputs)
\item
  Each (y\_A, y\_B) combination produces different Seed\_C
\item
  Total states needed: 2\^{}(r\_A) × 2\^{}(r\_B) = 2\^{}(r\_A+r\_B)
\item
  Requirements \textbf{add in the exponent} (multiply in the base)
\end{itemize}

\textbf{Why multiplication occurs (Injectivity + Dependency):}

At node C with parents A and B: Seed\_C = Enc(C \textbar{}\textbar{} (A,
Seed\_A, y\_A) \textbar{}\textbar{} (B, Seed\_B, y\_B)
\textbar{}\textbar{} GateDigest\_C)

Different parent outputs yield different seeds (A2: Injectivity):

\begin{itemize}
\tightlist
\item
  (y\_A=0, y\_B=0) \(\to\) Seed\_C(0,0)
\item
  (y\_A=0, y\_B=1) \(\to\) Seed\_C(0,1)
\item
  (y\_A=1, y\_B=0) \(\to\) Seed\_C(1,0)
\item
  (y\_A=1, y\_B=1) \(\to\) Seed\_C(1,1)
\end{itemize}

All four are distinct \(\to\) cannot merge without collision \(\to\)
must maintain all 4 states.

Generalizing: 2\^{}(R\_A) possibilities for A × 2\^{}(R\_B)
possibilities for B = 2\^{}(R\_A+R\_B) total.

\subparagraph{2.4.2 Cut Composition: From Nodes to
Bottlenecks}\label{242-cut-composition-from-nodes-to-bottlenecks}

\textbf{Terminology (Vertex Cuts):} Throughout this paper, all cuts are
\textbf{vertex cuts} (sets of nodes), not edge cuts. A vertex cut C
\(\subseteq\) V is a set of nodes whose removal disconnects source from
sink. By convention, cuts exclude Start and Goal (equivalently, in
node-splitting set c(Start)=c(Goal)=+\(\infty\)).

\textbf{Computational DAG model:} Start \(\to\) {[}Layer 1 nodes{]}
\(\to\) {[}Layer 2 nodes{]} \(\to\) ... \(\to\) Goal

A \textbf{cut} C is any set of nodes separating Start from Goal (every
source\(\to\)sink path hits C, i.e., contains at least one node from C).

\textbf{Per-node SCL (§2.1):} At each node v: q\_v + \(\Phi\)\_v
\(\geq\) R\_v

\textbf{Cut composition (additive):} For any cut C:

\(\Sigma\)\_\{v\(\in\)C\} q\_v + \(\Sigma\)\_\{v\(\in\)C\} \(\Phi\)\_v
\(\geq\) \(\Sigma\)\_\{v\(\in\)C\} R\_v

\textbf{Define cut quantities:}

\begin{itemize}
\tightlist
\item
  \textbf{\(\Lambda\)(C)} = \(\Sigma\)\_\{v\(\in\)C\} R\_v

  \begin{itemize}
  \tightlist
  \item
    Total emergence at cut
  \end{itemize}
\item
  \textbf{Q(C)} = \(\Sigma\)\_\{v\(\in\)C\} q\_v

  \begin{itemize}
  \tightlist
  \item
    Total resolution at cut (first-use RWA)
  \end{itemize}
\item
  \textbf{\(\lambda\)(C)} = \(\Lambda\)(C) - Q(C) =
  \(\Sigma\)\_\{v\(\in\)C\}(R\_v - q\_v)

  \begin{itemize}
  \tightlist
  \item
    Cut residual
  \end{itemize}
\end{itemize}

\textbf{Theorem 7.A.1 (Cut Composition):} For any cut C in
L*\textquotesingle{}s DAG and any correct algorithm:

\(\Sigma\)\_\{v\(\in\)C\} \(\Phi\)\_v \(\geq\) \(\lambda\)(C) where
\(\lambda\)(C) = \(\Sigma\)\_\{v\(\in\)C\}(R\_v - q\_v)

\textbf{Equivalently:}

\(\Phi\)(C) := log\(_2\)(\(\prod\)\emph{\{v\(\in\)C\} Alt\_v) \(\geq\)
\(\Sigma\)}\{v\(\in\)C\}(R\_v - q\_v)

\textbf{Proof intuition:}

\begin{itemize}
\tightlist
\item
  By multiplicative principle, 2\^{}(R\_v-q\_v) worlds at each v compose
  multiplicatively across cut
\item
  Total distinguishable artifacts \(\geq\) 2\^{}(\(\Sigma\)(R\_v-q\_v))
  = 2\^{}(\(\lambda\)(C))
\item
  Taking log yields result
\end{itemize}

(Full proof: §7.2 Theorem 7.A.1; see also Appendix J)

\subparagraph{\texorpdfstring{2.4.3 The Bottleneck Residual
\(\lambda\)(A,x)}{2.4.3 The Bottleneck Residual \textbackslash lambda(A,x)}}\label{243-the-bottleneck-residual-ux3bbax}

Not all cuts are equally hard. The \textbf{minimum residual cut}
determines overall complexity.

\textbf{Definition:} For algorithm A running on instance x:

\(\lambda\)(A,x) := min\_\{cuts C\} \(\lambda\)(C) = min\_C
\(\Sigma\)\_\{v\(\in\)C\}(R\_v - q\_v)

\textbf{Terminology note:} This is a \textbf{vertex-capacitated s-t
min-cut} (also called \textbf{minimum s-t vertex cut} or \textbf{s-t
vertex separator}):

\begin{itemize}
\tightlist
\item
  A vertex cut C \(\subseteq\) V that hits every source\(\to\)sink path
\item
  For minimum cuts in DAGs, an antichain representative exists
\item
  Capacity: cap(C) = \(\Sigma\)\_\{v\(\in\)C\} c(v), where node
  capacities c(v) = R\_v - q\_v
\item
  Computed via node-splitting reduction + max-flow
\end{itemize}

\textbf{Intuition:} Think of information flow through a DAG like water
through pipes. The narrowest pipe (bottleneck) determines maximum flow
rate. Similarly, the cut with minimum unresolved residual determines
computational difficulty - algorithms cannot avoid their hardest
bottleneck.

\textbf{Why the minimum matters:}

\begin{itemize}
\tightlist
\item
  Algorithm must pass through EVERY cut to reach the goal
\item
  Each cut requires \(\geq\) 2\^{}(\(\lambda\)(C)) distinguishable
  artifacts
\item
  The cut with minimum \(\lambda\)(C) still requires 2\^{}(min\_C
  \(\lambda\)(C)) artifacts
\item
  Cannot "route around" bottleneck - every computation path hits it
  (contains nodes from it)
\end{itemize}

\textbf{Example (3-node chain A\(\to\)B\(\to\)C):}

\begin{itemize}
\tightlist
\item
  Cut \{A\}: \(\lambda\) = R\_A - q\_A
\item
  Cut \{B\}: \(\lambda\) = R\_B - q\_B
\item
  Cut \{C\}: \(\lambda\) = R\_C - q\_C
\item
  Cut \{A,B\}: \(\lambda\) = (R\_A-q\_A) + (R\_B-q\_B)
\item
  Bottleneck: \(\lambda\)(A,x) = min of all above
\end{itemize}

\textbf{Summary:} Dependencies cause multiplicative growth (worlds
multiply across nodes), creating bottleneck \(\lambda\)(A,x).
Verification (with witness) achieves \(\lambda\)=0 (costs add); search
(without witness) maintains \(\lambda\) \(\geq\) \(\lambda\)\_base
(states multiply exponentially). Formal DAG framework: §4.2; cut
composition proofs: §7.2 and Appendix J (Theorem J.1; Lemma J.1-Cart).

\paragraph{2.5 Preview: Complete Worked Example (See
§2.6)}\label{25-preview-complete-worked-example-see-26}

A detailed walkthrough with explicit verification and search traces on a
3-node DAG instance appears in §2.6. The example demonstrates: instance
structure with seed dependencies, step-by-step verification trace
(polynomial time), multiple search attempts with failures (exponential
tries), collision mechanism explanation (why maintaining fewer than
2\^{}\(\lambda\) states causes errors), and gap analysis showing how
\(\lambda\)=5 yields 32× separation that scales to n\^{}(\(\Theta\)(log
n)) or 2\^{}(\(\Theta\)(n)) for full L* instances.

\begin{center}\rule{0.5\linewidth}{0.5pt}\end{center}

\paragraph{2.6 Complete Worked Example: 3-Node
Instance}\label{26-complete-worked-example-3-node-instance}

To ground the abstract principles of §§2.1-2.4, we now present a
complete worked example on the smallest non-trivial instance: a 3-node
DAG with \(\lambda\) = 5 bits of bottleneck residual. This instance is
small enough to trace by hand yet large enough to demonstrate the
exponential gap between verification and search. All calculations below
are verifiable with a hex calculator.

\textbf{Purpose of this walkthrough:}

\begin{itemize}
\tightlist
\item
  See SCL (q + \(\Phi\) \(\geq\) R) in action with concrete numbers
\item
  Trace polynomial verification (51 operations) step-by-step
\item
  Watch exponential search fail (32 attempts × 30 ops \(\approx\) 960
  operations)
\item
  Understand why maintaining fewer than 2\^{}\(\lambda\) artifacts
  causes collisions \(\to\) errors
\item
  Preview how this gap scales to n\^{}(\(\Omega\)(log n)) or
  2\^{}(\(\Theta\)(n)) for full L* (32× artifact separation from 2\^{}5
  = 32 required states; \textasciitilde{}19× operational ratio detailed
  in §2.6.4)
\end{itemize}

\subparagraph{2.6.1 The Instance}\label{261-the-instance}

\textbf{DAG Structure:}

\begin{verbatim}
Node 1 (Root, depth 0)
  |
Node 2 (Middle, depth 1)
  |
Node 3 (Leaf, depth 2)

Cut: {Node 2, Node 3}
\end{verbatim}

This is a simple linear chain - the simplest dependency structure that
creates a non-trivial bottleneck.

\textbf{Emergence Requirements (A3: Emergence):}

\begin{itemize}
\tightlist
\item
  R\(_1\) = 3 bits (Node 1 must produce 3 independent bits of
  information)
\item
  R\(_2\) = 3 bits (Node 2 must produce 3 additional independent bits)
\item
  R\(_3\) = 2 bits (Node 3 must produce 2 additional independent bits)
\item
  \textbf{R\_total = 3 + 3 + 2 = 8 bits} (entire instance)
\end{itemize}

\textbf{Bottleneck Analysis:} The cut C = \{Node 2, Node 3\} separates
the root from the leaves. At this cut:

\begin{itemize}
\tightlist
\item
  \(\lambda\)(C) = (R\(_2\) \(-\) q\(_2\)) + (R\(_3\) \(-\) q\(_3\))
  where q\(_2\), q\(_3\) are resolved bits at the cut
\item
  \textbf{Maximum residual: \(\lambda\)\_max = R\(_2\) + R\(_3\) = 3 + 2
  = 5 bits}
\item
  This occurs when the algorithm reaches the cut without resolving any
  bits at Node 2 or Node 3 (pure search mode)
\item
  Required artifacts at bottleneck: 2\^{}\(\lambda\) = 2\^{}5 = 32
  simultaneously distinguishable states
\end{itemize}

\textbf{Why 32 states?} With q\(_2\) = q\(_3\) = 0 (no bits resolved),
there remain 2³ × 2² = 32 possible worlds for (y\(_2\), y\(_3\)). Each
world corresponds to a distinct seed at Node 2, which determines
different designated addresses. To correctly solve all 32 cases, the
algorithm must either:

\begin{enumerate}
\def\labelenumi{\arabic{enumi}.}
\tightlist
\item
  Resolve bits (increase q) to reduce the 32 possibilities, OR
\item
  Maintain 32 distinguishable artifacts (\(\Phi\) \(\geq\) log\(_2\)(32)
  = 5)
\end{enumerate}

Maintaining fewer than 32 artifacts \(\to\) collision between worlds
\(\to\) wrong designated addresses \(\to\) verification failure.

\textbf{Seed Encoding Function (A2: Injectivity, A4: Closure):}

For this toy instance we use a simplified encoding that is injective on
the restricted domain we exercise here (distinct y for a fixed parent
seed produce distinct child seeds). The full L* construction in §6 uses
a stronger, length-delimited injective scheme over parent tuples.

Concretely, in this example the seed at each node is computed via:
Enc(v, parent\_seed, y) = (parent\_seed \^{} (v \textless{}\textless{}
12) \^{} (y\_int \textless{}\textless{} 4) \^{} y\_int) \& 0xFFFF where:

\begin{itemize}
\tightlist
\item
  v: node ID (1, 2, or 3)
\item
  parent\_seed: 16-bit seed value from parent (in hex)
\item
  y: output bits from parent as binary string (e.g., "101" for y = 5)
\item
  y\_int: integer value of y (e.g., int("101", 2) = 5)
\item
  \(\oplus\): bitwise XOR
\item
  \(\ll\): left bit shift
\item
  \& 0xFFFF: mask to 16 bits
\end{itemize}

\textbf{Initial Seed:}

Seed\(_1\) = 0x4A2F (fixed starting point for this instance)

\textbf{Seed Chain (Dependency):} Seed\(_2\) = Enc(2, Seed\(_1\),
y\(_1\)) where y\(_1\) in \{0,1\}³ Seed\(_3\) = Enc(3, Seed\(_2\),
y\(_2\)) where y\(_2\) in \{0,1\}³

\textbf{Key Property (Injectivity):} Different y\(_1\) values produce
different Seed\(_2\) values. For example:

\begin{itemize}
\tightlist
\item
  y\(_1\) = "010" (binary 2) \(\to\) Seed\(_2\) = 0x6A0D
\item
  y\(_1\) = "101" (binary 5) \(\to\) Seed\(_2\) = 0x6A7A
\item
  y\(_1\) = "110" (binary 6) \(\to\) Seed\(_2\) = 0x6A49
\end{itemize}

Since Enc is injective, all 8 possible y\(_1\) values produce 8 distinct
Seed\(_2\) values. This keyedness property (§1.1) is what prevents
algorithms from "merging" different computational paths:

\begin{itemize}
\tightlist
\item
  Each seed determines a unique set of designated addresses
\item
  Using the wrong seed yields wrong values
\item
  Verification catches the error
\end{itemize}

\textbf{Designated Address Function (A1: Hermeticity):}

Each node\textquotesingle{}s bits are stored in designated addresses
determined by its seed, using node-specific address pools to enforce
Hermeticity (disjoint pools per node):

\textbf{Notation note:} In this toy example, we write F\_overlay(v,
seed, j, \(\ell\)) for readability. In the formal framework we write
F\_overlay(Seed\_v; j, \(\ell\)), with the node id v implicitly encoded
in Seed\_v. F\_overlay(v, seed, j, l) = ((v-1) \textless{}\textless{} 4)
\textbar{} (((seed \textgreater{}\textgreater{} 2) \^{} (j
\textless{}\textless{} 2) \^{} l) \& 0x0F) where:

\begin{itemize}
\tightlist
\item
  v: node ID (2 or 3 in our toy instance; Node 1 is root with no
  designated reads)
\item
  seed: 16-bit seed value for this node
\item
  j: bit index (0 to R\_v \(-\) 1)
\item
  \(\ell\): level selector (0 or 1) - each bit uses 2 addresses for
  redundancy/encoding
\item
  \(\gg\), \(\ll\): right/left bit shift
\item
  \textbar{}: bitwise OR (combines node base with position offset)
\item
  Result: 6-bit address (0x00 to 0x3F, i.e., 64 possible addresses)
\end{itemize}

\textbf{Address Space (Partitioned by Node):}

\begin{itemize}
\tightlist
\item
  Total available: 64 addresses (0x00 to 0x3F)
\item
  \textbf{Node 2 pool}: 0x10-0x1F (16 addresses; needs 6 for R\(_2\)=3
  bits × 2 levels)
\item
  \textbf{Node 3 pool}: 0x20-0x2F (16 addresses; needs 4 for R\(_3\)=2
  bits × 2 levels)
\item
  \textbf{Hermeticity}: Node-specific base offsets ((v\(-\)1) \(\ll\) 4)
  mathematically guarantee disjoint pools

  \begin{itemize}
  \tightlist
  \item
    Holds across all nodes and all seed values
  \item
    No overlap possible even under wrong guesses
  \end{itemize}
\end{itemize}

\textbf{Example: Address Sets for Different Seeds}

For Node 2, Seed\(_2\) = 0x6A7A (correct seed when y\(_1\) = "101"):

\begin{itemize}
\tightlist
\item
  Bit 0: F\_overlay(2, 0x6A7A, 0, 0) = 0x1E, F\_overlay(2, 0x6A7A, 0, 1)
  = 0x1F
\item
  Bit 1: F\_overlay(2, 0x6A7A, 1, 0) = 0x1A, F\_overlay(2, 0x6A7A, 1, 1)
  = 0x1B
\item
  Bit 2: F\_overlay(2, 0x6A7A, 2, 0) = 0x16, F\_overlay(2, 0x6A7A, 2, 1)
  = 0x17
\item
  Address set: \{0x1E, 0x1F, 0x1A, 0x1B, 0x16, 0x17\} (all in Node 2
  pool 0x10-0x1F)
\end{itemize}

For Node 2, Seed\(_2\)\textquotesingle{} = 0x6A0D (incorrect seed when
y\(_1\) = "010"):

\begin{itemize}
\tightlist
\item
  Bit 0: F\_overlay(2, 0x6A0D, 0, 0) = 0x13, F\_overlay(2, 0x6A0D, 0, 1)
  = 0x12
\item
  Bit 1: F\_overlay(2, 0x6A0D, 1, 0) = 0x17, F\_overlay(2, 0x6A0D, 1, 1)
  = 0x16
\item
  Bit 2: F\_overlay(2, 0x6A0D, 2, 0) = 0x1B, F\_overlay(2, 0x6A0D, 2, 1)
  = 0x1A
\item
  Address set: \{0x13, 0x12, 0x17, 0x16, 0x1B, 0x1A\} (all in Node 2
  pool 0x10-0x1F)
\end{itemize}

For Node 3, Seed\(_3\) = 0x5A1C (correct seed when y\(_2\) = "110"):

\begin{itemize}
\tightlist
\item
  Bit 0: F\_overlay(3, 0x5A1C, 0, 0) = 0x27, F\_overlay(3, 0x5A1C, 0, 1)
  = 0x26
\item
  Bit 1: F\_overlay(3, 0x5A1C, 1, 0) = 0x23, F\_overlay(3, 0x5A1C, 1, 1)
  = 0x22
\item
  Address set: \{0x27, 0x26, 0x23, 0x22\} (all in Node 3 pool 0x20-0x2F)
\end{itemize}

\textbf{Crucial observations:}

\begin{enumerate}
\def\labelenumi{\arabic{enumi}.}
\tightlist
\item
  \textbf{Node-specific pools}: Node 2 addresses are always 0x10-0x1F,
  Node 3 addresses are always 0x20-0x2F, regardless of seed value -
  mathematically guaranteed disjoint (Hermeticity A1).
\item
  \textbf{Seed-specific patterns}: Within a node\textquotesingle{}s
  pool, different seeds produce different address patterns. Wrong seed
  \(\to\) wrong addresses \(\to\) wrong values \(\to\) verification
  failure (Keyedness enforces SCL).
\end{enumerate}

\textbf{Correct Solution (for reference in traces below):}

\begin{itemize}
\tightlist
\item
  y\(_1\) = "101" (binary 5) \(\to\) Seed\(_2\) = 0x6A7A
\item
  y\(_2\) = "110" (binary 6) \(\to\) Seed\(_3\) = 0x5A1C
\item
  y\(_3\) = "10" (binary 2)
\end{itemize}

This is the satisfying assignment we\textquotesingle{}ll trace in
§2.6.2. The verifier with witness (y\(_1\), y\(_2\), y\(_3\)) can check
correctness in polynomial time. A searcher without the witness must
explore up to 2³ × 2³ × 2² = 256 possibilities (or 2³ × 2² = 32 at the
cut), demonstrating exponential cost.

\textbf{Summary of Instance Parameters:}

\begin{itemize}
\tightlist
\item
  Nodes: 3 (linear chain)
\item
  Emergence: R = (3, 3, 2) bits
\item
  Cut residual: \(\lambda\) = 5 bits \(\to\) 2\^{}5 = 32 required states
\item
  Seeds: Enc(v, parent, y) with InitSeed = 0x4A2F
\item
  Addresses: F\_overlay(v, seed, j, \(\ell\)) over per-node pools (Node
  2: 0x10-0x1F, Node 3: 0x20-0x2F)
\item
  Gap preview: Verification \(\approx\) 51 ops vs Search \(\approx\) 960
  ops (\textasciitilde{}19× operational ratio; 32× artifact ratio from
  2\^{}\(\lambda\))
\end{itemize}

\subparagraph{2.6.2 Verification Trace (Polynomial
Time)}\label{262-verification-trace-polynomial-time}

We now trace the verifier\textquotesingle{}s execution when given the
correct witness W = (y\(_1\), y\(_2\), y\(_3\)). The verifier checks
that the witness is consistent with the instance structure and satisfies
all constraints. This entire process takes \textbf{51 operations} -
polynomial in R\_total = 8 bits.

\textbf{Given Witness:}

W = (y\(_1\), y\(_2\), y\(_3\)) = ("101", "110", "10") These are the
correct output values for nodes 1, 2, and 3 respectively. The verifier
must confirm:

\begin{enumerate}
\def\labelenumi{\arabic{enumi}.}
\tightlist
\item
  Seeds are correctly computed from the witness
\item
  Designated addresses are correctly derived from seeds
\item
  All completeness constraints (H\_v · x\_v = y\_v) are satisfied
\item
  The decoded CNF formula is satisfied by the final assignment
\end{enumerate}

\textbf{Step 1: Parse Witness (3 operations)}

\begin{itemize}
\tightlist
\item
  Extract y\(_1\) = "101" (3 bits) \(\to\) 1 op
\item
  Extract y\(_2\) = "110" (3 bits) \(\to\) 1 op
\item
  Extract y\(_3\) = "10" (2 bits) \(\to\) 1 op
\item
  \textbf{Subtotal: 3 operations}
\end{itemize}

\textbf{Step 2: Verify Initial Seed (2 operations)}

\begin{itemize}
\tightlist
\item
  Check Seed\(_1\) = 0x4A2F (given in instance) \(\to\) 1 op
  (comparison)
\item
  Store Seed\(_1\) for chain verification \(\to\) 1 op (memory write)
\item
  \textbf{Subtotal: 2 operations}
\end{itemize}

\textbf{Step 3: Verify Node 1 Output (2 operations)}

\begin{itemize}
\tightlist
\item
  Node 1 is the root with fixed output y\(_1\) = "101" (from witness)
\item
  Verify y\(_1\) has correct length (3 bits = R\(_1\)) \(\to\) 1 op
\item
  Store y\(_1\) for seed chain \(\to\) 1 op
\item
  \textbf{Subtotal: 2 operations}
\end{itemize}

\textbf{Step 4: Verify Seed Chain - Node 2 (2 operations)}

\begin{itemize}
\tightlist
\item
  Compute Seed\(_2\) = Enc(2, Seed\(_1\), y\(_1\)) = Enc(2, 0x4A2F,
  "101")
\item
  Calculation:

  \begin{itemize}
  \tightlist
  \item
    v = 2, parent\_seed = 0x4A2F, y = "101" \(\to\) y\_int = 5
  \item
    Seed\(_2\) = (0x4A2F \(\oplus\) (2 \(\ll\) 12) \(\oplus\) (5 \(\ll\)
    4) \(\oplus\) 5) \& 0xFFFF
  \item
    Seed\(_2\) = (0x4A2F \(\oplus\) 0x2000 \(\oplus\) 0x50 \(\oplus\)
    0x5) \& 0xFFFF
  \item
    Seed\(_2\) = 0x6A7A
  \end{itemize}
\item
  Perform Enc computation \(\to\) 1 op (arithmetic)
\item
  Store Seed\(_2\) \(\to\) 1 op
\item
  \textbf{Subtotal: 2 operations}
\item
  \textbf{Result: Seed\(_2\) = 0x6A7A}
\end{itemize}

\textbf{Step 5: Verify Node 2 Designated Reads (15 operations)}

Node 2 has R\(_2\) = 3 bits, each requiring designated address
computations and reads. For each bit j \(\in\) \{0, 1, 2\}:

\textbf{Bit 0 (j=0):}

\begin{itemize}
\tightlist
\item
  Compute addr\(_0\)\textsuperscript{(}\textsuperscript{0}\textsuperscript{)} = F\_overlay(2, 0x6A7A, 0, 0) = 0x10 \textbar{}
  (((0x6A7A \(\gg\) 2) \(\oplus\) 0 \(\oplus\) 0) \& 0x0F) = 0x1E
  \(\to\) 1 op
\item
  Compute addr\(_0\)\textsuperscript{(}¹\textsuperscript{)} = F\_overlay(2, 0x6A7A, 0, 1) = 0x10 \textbar{}
  (((0x6A7A \(\gg\) 2) \(\oplus\) 0 \(\oplus\) 1) \& 0x0F) = 0x1F
  \(\to\) 1 op
\item
  Read value from addr 0x1E \(\to\) 1 op
\item
  Read value from addr 0x1F \(\to\) 1 op
\item
  Verify consistency (both addresses yield same bit value) \(\to\) 1 op
\item
  \textbf{Subtotal for bit 0: 5 operations}
\end{itemize}

\textbf{Bit 1 (j=1):}

\begin{itemize}
\tightlist
\item
  Compute addr\(_1\)\textsuperscript{(}\textsuperscript{0}\textsuperscript{)} = F\_overlay(2, 0x6A7A, 1, 0) = 0x1A \(\to\) 1
  op
\item
  Compute addr\(_1\)\textsuperscript{(}¹\textsuperscript{)} = F\_overlay(2, 0x6A7A, 1, 1) = 0x1B \(\to\) 1
  op
\item
  Read from 0x1A \(\to\) 1 op
\item
  Read from 0x1B \(\to\) 1 op
\item
  Verify consistency \(\to\) 1 op
\item
  \textbf{Subtotal for bit 1: 5 operations}
\end{itemize}

\textbf{Bit 2 (j=2):}

\begin{itemize}
\tightlist
\item
  Compute addr\(_2\)\textsuperscript{(}\textsuperscript{0}\textsuperscript{)} = F\_overlay(2, 0x6A7A, 2, 0) = 0x16 \(\to\) 1
  op
\item
  Compute addr\(_2\)\textsuperscript{(}¹\textsuperscript{)} = F\_overlay(2, 0x6A7A, 2, 1) = 0x17 \(\to\) 1
  op
\item
  Read from 0x16 \(\to\) 1 op
\item
  Read from 0x17 \(\to\) 1 op
\item
  Verify consistency \(\to\) 1 op
\item
  \textbf{Subtotal for bit 2: 5 operations}
\end{itemize}

\textbf{Total for Step 5: 5 + 5 + 5 = 15 operations}

\textbf{Designated addresses for Node 2:} \{0x1E, 0x1F, 0x1A, 0x1B,
0x16, 0x17\}

\textbf{Step 6: Verify Node 2 Completeness (9 operations)}

Check that H\(_2\) · x\(_2\) = y\(_2\) where:

\begin{itemize}
\tightlist
\item
  H\(_2\) is a 3×3 emergence matrix (full rank, R\(_2\) = 3)
\item
  x\(_2\) is the input vector at Node 2 (includes parent outputs and
  local variables)
\item
  y\(_2\) = "110" is the output claimed by the witness
\end{itemize}

This requires a linear system check:

\begin{itemize}
\tightlist
\item
  Set up 3 equations (one per row of H\(_2\)) \(\to\) 3 ops
\item
  Substitute x\(_2\) values and compute each equation \(\to\) 3 ops (1
  per equation)
\item
  Verify each equation equals corresponding bit of y\(_2\) \(\to\) 3 ops
  (1 per equation)
\item
  \textbf{Subtotal: 3 + 3 + 3 = 9 operations}
\end{itemize}

This is a simplified count; actual Gaussian elimination would be
O(R\(_2\)³) = O(27), but since we are just checking (not solving), we
can verify the linear constraints in O(R\(_2\)²) = O(9) operations.

\textbf{Step 7: Verify Seed Chain - Node 3 (2 operations)}

\begin{itemize}
\tightlist
\item
  Compute Seed\(_3\) = Enc(3, Seed\(_2\), y\(_2\)) = Enc(3, 0x6A7A,
  "110")
\item
  Calculation:

  \begin{itemize}
  \tightlist
  \item
    v = 3, parent\_seed = 0x6A7A, y = "110" \(\to\) y\_int = 6
  \item
    Seed\(_3\) = (0x6A7A \(\oplus\) (3 \(\ll\) 12) \(\oplus\) (6 \(\ll\)
    4) \(\oplus\) 6) \& 0xFFFF
  \item
    Seed\(_3\) = (0x6A7A \(\oplus\) 0x3000 \(\oplus\) 0x60 \(\oplus\)
    0x6) \& 0xFFFF
  \item
    Seed\(_3\) = 0x5A1C
  \end{itemize}
\item
  Perform Enc computation \(\to\) 1 op
\item
  Store Seed\(_3\) \(\to\) 1 op
\item
  \textbf{Subtotal: 2 operations}
\item
  \textbf{Result: Seed\(_3\) = 0x5A1C}
\end{itemize}

\textbf{Step 8: Verify Node 3 Designated Reads (10 operations)}

Node 3 has R\(_3\) = 2 bits:

\textbf{Bit 0 (j=0):}

\begin{itemize}
\tightlist
\item
  Compute addr\(_0\)\textsuperscript{(}\textsuperscript{0}\textsuperscript{)} = F\_overlay(3, 0x5A1C, 0, 0) = 0x20 \textbar{}
  (((0x5A1C \(\gg\) 2) \(\oplus\) 0 \(\oplus\) 0) \& 0x0F) = 0x27
  \(\to\) 1 op
\item
  Compute addr\(_0\)\textsuperscript{(}¹\textsuperscript{)} = F\_overlay(3, 0x5A1C, 0, 1) = 0x26 \(\to\) 1
  op
\item
  Read from 0x27 \(\to\) 1 op
\item
  Read from 0x26 \(\to\) 1 op
\item
  Verify consistency \(\to\) 1 op
\item
  \textbf{Subtotal for bit 0: 5 operations}
\end{itemize}

\textbf{Bit 1 (j=1):}

\begin{itemize}
\tightlist
\item
  Compute addr\(_1\)\textsuperscript{(}\textsuperscript{0}\textsuperscript{)} = F\_overlay(3, 0x5A1C, 1, 0) = 0x23 \(\to\) 1
  op
\item
  Compute addr\(_1\)\textsuperscript{(}¹\textsuperscript{)} = F\_overlay(3, 0x5A1C, 1, 1) = 0x22 \(\to\) 1
  op
\item
  Read from 0x23 \(\to\) 1 op
\item
  Read from 0x22 \(\to\) 1 op
\item
  Verify consistency \(\to\) 1 op
\item
  \textbf{Subtotal for bit 1: 5 operations}
\end{itemize}

\textbf{Total for Step 8: 5 + 5 = 10 operations}

\textbf{Designated addresses for Node 3:} \{0x27, 0x26, 0x23, 0x22\}
(all in Node 3 pool 0x20-0x2F)

\textbf{Step 9: Verify Node 3 Completeness (4 operations)}

Check H\(_3\) · x\(_3\) = y\(_3\) where:

\begin{itemize}
\tightlist
\item
  H\(_3\) is a 2×2 emergence matrix (R\(_3\) = 2)
\item
  x\(_3\) is the input at Node 3
\item
  y\(_3\) = "10" from witness
\end{itemize}

Linear system check:

\begin{itemize}
\tightlist
\item
  Set up 2 equations \(\to\) 2 ops
\item
  Verify both equations \(\to\) 2 ops
\item
  \textbf{Subtotal: 4 operations}
\end{itemize}

\textbf{Step 10: Final Acceptance (2 operations)}

All structural verification is now complete:

\begin{itemize}
\tightlist
\item
  Seed chain verified for all 3 nodes (correct Enc computation)
\item
  Designated reads performed at correct addresses (Hermeticity
  satisfied)
\item
  Emergence constraints H\_v · x\_v = y\_v satisfied at all nodes
  (Completeness)
\end{itemize}

Final acceptance check:

\begin{itemize}
\tightlist
\item
  Confirm all nodes processed successfully \(\to\) 1 op
\item
  Return ACCEPT \(\to\) 1 op
\item
  \textbf{Subtotal: 2 operations}
\end{itemize}

\emph{Note on full L*: The complete construction includes a seed-locked
decode schema \(\Phi\)\textasciitilde{} where even the CNF formula \(\varphi\) is
accessed through designated addresses (§10.1.1, OAP mechanism). For this
simplified toy example, the emergence constraints H\_v · x\_v = y\_v
serve the analogous structural role---they define what "correct" means
at each node and are what the wrong-seed detection mechanism checks
against.}

\begin{center}\rule{0.5\linewidth}{0.5pt}\end{center}

\textbf{Verification Summary:}

\begin{enumerate}
\def\labelenumi{\arabic{enumi}.}
\tightlist
\item
  \textbf{Parse witness} --- 3 operations
\item
  \textbf{Verify initial seed} --- 2 operations
\item
  \textbf{Verify Node 1 output} --- 2 operations
\item
  \textbf{Verify Seed\(_2\) chain} --- 2 operations
\item
  \textbf{Node 2 designated reads (3 bits × 5)} --- 15 operations
\item
  \textbf{Node 2 completeness check} --- 9 operations
\item
  \textbf{Verify Seed\(_3\) chain} --- 2 operations
\item
  \textbf{Node 3 designated reads (2 bits × 5)} --- 10 operations
\item
  \textbf{Node 3 completeness check} --- 4 operations
\item
  \textbf{Final acceptance} --- 2 operations
\end{enumerate}

\textbf{Total: 51 operations}

\textbf{Time Complexity Analysis:}

The verification algorithm\textquotesingle{}s running time is polynomial
in R\_total = 8:

\begin{itemize}
\tightlist
\item
  Seed computations: O(R\_total) = O(8) operations (constant per node)
\item
  Designated reads: O(\(\Sigma\) R\_v over nodes with designated reads).
  In this toy: 5 bits (nodes 2-3) × 2 addresses per bit; O(R\_total) in
  general.
\item
  Completeness checks: O(R\_total²) \(\approx\) O(64) operations (linear
  algebra on R\_v × R\_v matrices)
\item
  \textbf{Overall: O(R\_total²) = O(64) operations}
\end{itemize}

For the actual 51 operations counted above:

\begin{itemize}
\tightlist
\item
  51 \(\approx\) 0.8 × R\_total² = 0.8 × 64
\item
  This confirms polynomial scaling
\end{itemize}

For general L* instances with R\_total = \(\Theta\)(n log n):

\begin{itemize}
\tightlist
\item
  Verification time: O((n log n)²) = O(n² log² n) = polynomial in n
\end{itemize}

\textbf{Key Insight:} With the witness W, the verifier directly computes
the correct seeds at each step, accesses the correct designated
addresses, and verifies all constraints in polynomial time. The verifier
never needs to explore alternative possibilities - the witness provides
a certificate that eliminates all exponential branching. This is exactly
the P vs NP dichotomy: verification in polynomial time, but finding the
witness (as we\textquotesingle{}ll see in §2.6.3) requires exponential
search.

\textbf{SCL Satisfied:} During verification with witness:

\begin{itemize}
\tightlist
\item
  At Node 2: q\(_2\) = 3 (all 3 bits resolved), so \(\Phi\)\(_2\)
  \(\geq\) R\(_2\) \(-\) q\(_2\) = 0 (no branching needed)
\item
  At Node 3: q\(_3\) = 2 (all 2 bits resolved), so \(\Phi\)\(_3\)
  \(\geq\) R\(_3\) \(-\) q\(_3\) = 0
\item
  \textbf{Bottleneck residual: \(\lambda\) = 0} (verification mode)
\item
  Artifacts maintained: 2\^{}\(\lambda\) = 2\^{}0 = 1 state (single path
  through the computation)
\end{itemize}

This is why verification is polynomial: \(\lambda\) = 0 \(\to\) no
exponential cost.

\subparagraph{2.6.3 Search Trace (Exponential
Time)}\label{263-search-trace-exponential-time}

Now we examine what happens when an algorithm \emph{does not} have the
witness - it must search for the correct (y\(_1\), y\(_2\), y\(_3\))
values. Without prior knowledge, the algorithm faces 2³ × 2³ × 2² = 256
total possibilities (or 32 possibilities for (y\(_2\), y\(_3\)) at the
cut if y\(_1\) is known). This section traces two specific wrong guesses
to show how the collision mechanism detects errors and forces
backtracking.

\textbf{Search Space:}

\begin{itemize}
\tightlist
\item
  Node 1 outputs: 2³ = 8 possibilities for y\(_1\)
\item
  Node 2 outputs: 2³ = 8 possibilities for y\(_2\)
\item
  Node 3 outputs: 2² = 4 possibilities for y\(_3\)
\item
  \textbf{Total: 8 × 8 × 4 = 256 complete paths}
\item
  \textbf{At cut \{2,3\}: 8 × 4 = 32 seed-distinct states} (if y\(_1\)
  is resolved)
\end{itemize}

Without the witness, the algorithm must either:

\begin{enumerate}
\def\labelenumi{\arabic{enumi}.}
\tightlist
\item
  \textbf{Resolve bits} through propagation/learning (increase q to
  reduce search space), OR
\item
  Maintain 2\^{}\(\lambda\) distinguishable artifacts to track all
  possible worlds
\end{enumerate}

For this instance with \(\lambda\) = 5, that means maintaining 32
simultaneous states - which translates to 32× cost multiplier in most
paradigms (backtracking tries, DP keys, OBDD width, resolution proof
size).

\begin{center}\rule{0.5\linewidth}{0.5pt}\end{center}

\textbf{Failure Example 1: Wrong y\(_1\) Guess}

\textbf{Scenario:} The algorithm guesses y\(_1\) = "010" instead of the
correct y\(_1\) = "101".

\textbf{Step 1: Compute Wrong Seed at Node 2}

\begin{itemize}
\tightlist
\item
  Guess: y\(_1\) = "010" (binary 2) --- incorrect guess
\item
  Compute Seed\(_2\)\textquotesingle{} = Enc(2, 0x4A2F, "010")

  \begin{itemize}
  \tightlist
  \item
    Seed\(_2\)\textquotesingle{} = (0x4A2F \(\oplus\) (2 \(\ll\) 12)
    \(\oplus\) (2 \(\ll\) 4) \(\oplus\) 2) \& 0xFFFF
  \item
    Seed\(_2\)\textquotesingle{} = (0x4A2F \(\oplus\) 0x2000 \(\oplus\)
    0x20 \(\oplus\) 0x2) \& 0xFFFF
  \item
    \textbf{Seed\(_2\)\textquotesingle{} = 0x6A0D} (incorrect; correct
    value is 0x6A7A)
  \end{itemize}
\end{itemize}

\textbf{Step 2: Compute Wrong Addresses}

Using the wrong seed Seed\(_2\)\textquotesingle{} = 0x6A0D, the
algorithm computes designated addresses:

\begin{itemize}
\tightlist
\item
  Bit 0: F\_overlay(2, 0x6A0D, 0, 0) = 0x13, F\_overlay(2, 0x6A0D, 0, 1)
  = 0x12
\item
  Bit 1: F\_overlay(2, 0x6A0D, 1, 0) = 0x17, F\_overlay(2, 0x6A0D, 1, 1)
  = 0x16
\item
  Bit 2: F\_overlay(2, 0x6A0D, 2, 0) = 0x1B, F\_overlay(2, 0x6A0D, 2, 1)
  = 0x1A
\item
  \textbf{Wrong address set: \{0x13, 0x12, 0x17, 0x16, 0x1B, 0x1A\}}
  (all in Node 2 pool 0x10-0x1F)
\end{itemize}

Compare to correct addresses (from §2.6.2 with Seed\(_2\) = 0x6A7A):

\begin{itemize}
\tightlist
\item
  \textbf{Correct address set: \{0x1E, 0x1F, 0x1A, 0x1B, 0x16, 0x17\}}
  (all in Node 2 pool 0x10--0x1F)
\end{itemize}

\textbf{Crucial observation:} The two address sets share some addresses
(\{0x1A, 0x1B, 0x16, 0x17\}) but use them for \textbf{different bit
positions} - bit 1 and bit 2 in correct execution vs. different
positions in wrong execution. This mismapping \(\to\) wrong retrieved
values \(\to\) detectable failures. Both sets stay within Node
2\textquotesingle{}s pool 0x10-0x1F, demonstrating Hermeticity.

\textbf{Step 3: Read Wrong Values}

When the algorithm reads from addresses \{0x13, 0x12, ...\}, it
retrieves values at wrong positions or gets inconsistent data. Even
where addresses overlap (\{0x1A, 0x1B, 0x16, 0x17\}), they are assigned
to different bit indices, causing semantic errors.

\textbf{Step 4: Detect Failure}

The error is detected in one of two ways:

\textbf{(A) Bit Position Mismapping:} The wrong seed causes the
algorithm to read from addresses assigned to different bit positions.
For example, addresses \{0x17, 0x16\} store bit \textbf{2} data in the
correct execution (Seed\(_2\)=0x6A7A), but the wrong seed
(Seed\(_2\)\textquotesingle=0x6A0D) tries to use them for bit
\textbf{1}. This cross-bit mismapping \(\to\) wrong semantic context
\(\to\) inconsistent or wrong values.

\textbf{(B) Completeness Constraint Violation:} When the algorithm
checks H\(_2\) · x\(_2\) = y\(_2\), the values read (even from
overlapping addresses) represent the wrong seed\textquotesingle{}s
context and will not satisfy the linear emergence constraints \(\to\)
verification failure.

Overlap within Node 2\textquotesingle{}s pool does not save the try:
designated contents are seed-context dependent, and using an address
under the wrong seed context violates completeness; the try still fails.

Either way, the algorithm detects that y\(_1\) = "010" is incorrect and
must backtrack to try a different value.

\textbf{Operations for this failed attempt:} \textasciitilde{}30 ops
(seed computation, address computation, reads, consistency checks)

\begin{center}\rule{0.5\linewidth}{0.5pt}\end{center}

\textbf{Failure Example 2: Wrong y\(_2\) Guess (Correct y\(_1\))}

\textbf{Scenario:} The algorithm correctly guesses y\(_1\) = "101" but
then guesses y\(_2\) = "000" instead of the correct y\(_2\) = "110".

\textbf{Step 1: Correct Seed at Node 2}

\begin{itemize}
\tightlist
\item
  Correct guess: y\(_1\) = "101" \(\to\) Seed\(_2\) = 0x6A7A {[}YES{]}
\item
  Node 2 designated reads succeed (correct addresses accessed)
\end{itemize}

\textbf{Step 2: Compute Wrong Seed at Node 3}

\begin{itemize}
\tightlist
\item
  Guess: y\(_2\) = "000" (binary 0) --- incorrect guess
\item
  Compute Seed\(_3\)\textquotesingle{} = Enc(3, 0x6A7A, "000")

  \begin{itemize}
  \tightlist
  \item
    Seed\(_3\)\textquotesingle{} = (0x6A7A \(\oplus\) (3 \(\ll\) 12)
    \(\oplus\) (0 \(\ll\) 4) \(\oplus\) 0) \& 0xFFFF
  \item
    Seed\(_3\)\textquotesingle{} = (0x6A7A \(\oplus\) 0x3000 \(\oplus\)
    0x00 \(\oplus\) 0x0) \& 0xFFFF
  \item
    \textbf{Seed\(_3\)\textquotesingle{} = 0x5A7A} (incorrect; correct
    value is 0x5A1C)
  \end{itemize}
\end{itemize}

\textbf{Step 3: Compute Wrong Addresses at Node 3}

Using wrong seed Seed\(_3\)\textquotesingle{} = 0x5A7A:

\begin{itemize}
\tightlist
\item
  Bit 0: F\_overlay(3, 0x5A7A, 0, 0) = 0x2E, F\_overlay(3, 0x5A7A, 0, 1)
  = 0x2F
\item
  Bit 1: F\_overlay(3, 0x5A7A, 1, 0) = 0x2A, F\_overlay(3, 0x5A7A, 1, 1)
  = 0x2B
\item
  \textbf{Wrong address set: \{0x2E, 0x2F, 0x2A, 0x2B\}} (all in Node 3
  pool 0x20-0x2F)
\end{itemize}

Correct addresses (from §2.6.2 with Seed\(_3\) = 0x5A1C):

\begin{itemize}
\tightlist
\item
  \textbf{Correct address set: \{0x27, 0x26, 0x23, 0x22\}} (all in Node
  3 pool 0x20-0x2F)
\end{itemize}

\textbf{Disjoint address sets within the same node\textquotesingle{}s
pool} (for this pair of seeds) - the wrong seed points to entirely
different locations within the 0x20-0x2F space.

\textbf{Step 4: Detect Failure}

The algorithm reads from wrong addresses \(\to\) gets wrong bit values
\(\to\) either:

\begin{itemize}
\tightlist
\item
  (A) Completeness constraint H\(_3\) · x\(_3\) = y\(_3\) fails, OR
\item
  (B) Final CNF check fails (the decoded formula is not satisfied)
\end{itemize}

Disjointness within Node 3\textquotesingle{}s pool makes misreads
immediate: the wrong seed\textquotesingle{}s designated addresses are
disjoint from the correct set, so there is no accidental overlap to
exploit; reads return values for the wrong world, and verification
fails.

The algorithm detects y\(_2\) = "000" is incorrect and must backtrack.

\textbf{Operations for this failed attempt:} \textasciitilde{}35 ops
(includes Node 2 success + Node 3 failure)

\begin{center}\rule{0.5\linewidth}{0.5pt}\end{center}

\textbf{Search Complexity Analysis}

\textbf{Exhaustive Search Cost:}

At the bottleneck cut \{Node 2, Node 3\}, there are 2³ × 2² = 32
possible combinations of (y\(_2\), y\(_3\)) values (assuming y\(_1\) has
been resolved). In the worst case, an exhaustive search algorithm must
try all 32 combinations before finding the correct one.

\textbf{Operations per attempt:}

\begin{itemize}
\tightlist
\item
  Guess y values: \textasciitilde{}3 ops
\item
  Compute seeds: \textasciitilde{}6 ops (one per node in the cut)
\item
  Compute addresses and perform reads: \textasciitilde{}15-20 ops
\item
  Check constraints: \textasciitilde{}5 ops
\item
  \textbf{Total: \textasciitilde{}30-35 ops per attempt}
\end{itemize}

\textbf{Worst-case total:}

\begin{itemize}
\tightlist
\item
  32 attempts × 30 ops/attempt = \textbf{960 operations}
\end{itemize}

\textbf{Compare to verification:}

\begin{itemize}
\tightlist
\item
  Verification with witness: \textbf{51 operations}
\item
  Search without witness: \textbf{960 operations}
\item
  \textbf{Gap: 960 / 51 \(\approx\) 19× for \(\lambda\) = 5 bits}
\end{itemize}

\textbf{Why 19× and not 32×?} The factor of 32 = 2\^{}\(\lambda\)
appears in the \emph{number of attempts}, but each attempt is cheaper
than full verification (failures detected early). The 19× ratio is the
realized gap for this toy instance. For larger instances, early failure
detection becomes less effective, and the gap approaches
2\^{}\(\lambda\).

\textbf{SCL During Search:}

When searching without witness:

\begin{itemize}
\tightlist
\item
  At Node 2: q\(_2\) = 0 (no bits resolved yet), so \(\Phi\)\(_2\)
  \(\geq\) R\(_2\) \(-\) q\(_2\) = 3 bits
\item
  At Node 3: q\(_3\) = 0 (no bits resolved yet), so \(\Phi\)\(_3\)
  \(\geq\) R\(_3\) \(-\) q\(_3\) = 2 bits
\item
  \textbf{Bottleneck residual: \(\lambda\) = 3 + 2 = 5 bits}
\item
  Artifacts required: 2\^{}\(\lambda\) = 2\^{}5 = 32 states
\end{itemize}

The algorithm must maintain 32 distinguishable computational states
(backtracking nodes, DP table entries, OBDD nodes, or resolution
clauses) to correctly handle all possible worlds at the cut. Maintaining
fewer than 32 states \(\to\) collision between distinct seeds \(\to\)
wrong answers on some instances.

\textbf{Scaling to Full L*:}

This 3-node toy instance demonstrates the mechanism with \(\lambda\) = 5
bits \(\to\) 19× gap. For full L* instances:

\begin{itemize}
\tightlist
\item
  \textbf{QP-sharp profile:} \(\lambda\)\_base = \(\Theta\)(log² n)
  \(\to\) time \(\geq\) n\^{}(\(\Theta\)(log n))
\item
  \textbf{Flat profile:} \(\lambda\)\_base = \(\Theta\)(n) \(\to\) time
  \(\geq\) 2\^{}(\(\Theta\)(n))
\end{itemize}

The exponential relationship between \(\lambda\) and complexity is the
core of the SCL framework.

\subparagraph{2.6.4 Collision Mechanism: Why Keyedness Enforces
SCL}\label{264-collision-mechanism-why-keyedness-enforces-scl}

The two failure examples above demonstrate the core mechanism that makes
SCL mathematically necessary for L*: \textbf{keyedness via
seed-determined addresses}. This subsection explains \emph{why} wrong
guesses inevitably lead to detectable failures, and why maintaining
fewer than 2\^{}\(\lambda\) artifacts causes collisions.

\textbf{The Keyedness Property:}

Every distinct seed value Seed\_v determines a unique set of designated
addresses via F\_overlay. For our instance:

\begin{itemize}
\tightlist
\item
  Different y\(_1\) \(\to\) different Seed\(_2\) \(\to\) different
  address sets at Node 2
\item
  Different y\(_2\) \(\to\) different Seed\(_3\) \(\to\) different
  address sets at Node 3
\end{itemize}

This property is guaranteed by:

\begin{enumerate}
\def\labelenumi{\arabic{enumi}.}
\tightlist
\item
  \textbf{Injectivity of Enc (A2):} distinct input tuples \(\to\)
  distinct seeds
\item
  \textbf{Hermeticity (A1):} information flows only through designated
  addresses
\item
  \textbf{Emergence (A3):} R\_v independent bits required at each node
\end{enumerate}

\textbf{Why Collisions Are Unavoidable:}

Consider an algorithm that tries to "compress" the computation by
maintaining fewer than 2\^{}\(\lambda\) = 32 distinct artifacts at the
cut \{Node 2, Node 3\}. By the pigeonhole principle, at least two of the
32 possible seed combinations must map to the same artifact state.

\textbf{Collision scenario:}

\begin{itemize}
\tightlist
\item
  World \(\alpha\): (y\(_2\), y\(_3\)) = ("110", "10") \(\to\)
  (Seed\(_2\), Seed\(_3\)) = (0x6A7A, 0x5A1C) \(\to\) addresses \{0x1E,
  0x1F, ..., 0x27, 0x26, ...\}
\item
  World \(\beta\): (y\(_2\), y\(_3\)) = ("011", "01") \(\to\)
  (Seed\(_2\)\textquotesingle, Seed\(_3\)\textquotesingle) = (0x6A0D,
  0x...) \(\to\) addresses \{0x13, 0x12, ..., 0x2E, 0x2F, ...\}
\end{itemize}

If the algorithm merges worlds \(\alpha\) and \(\beta\) into a single
artifact state (to save space), it loses the distinction between
Seed\(_2\) = 0x6A7A and Seed\(_2\)\textquotesingle{} = 0x6A0D. When it
later needs to access designated addresses at Node 2:

\begin{itemize}
\tightlist
\item
  It cannot correctly access both address sets \{0x1E, ...\} and \{0x13,
  ...\}
\item
  It must choose one seed - say, 0x6A7A
\item
  World \(\beta\) then reads from wrong addresses \(\to\) gets wrong
  values \(\to\) fails verification
\end{itemize}

\textbf{This is not a probabilistic or heuristic argument - it is
deterministic:} Given the instance structure (\textbf{A1-A5}
properties), any algorithm that conflates distinct seeds \emph{will}
produce wrong answers on some inputs. The only ways to avoid collisions
are:

\begin{enumerate}
\def\labelenumi{\arabic{enumi}.}
\tightlist
\item
  Maintain \(\geq\) 2\^{}\(\lambda\) distinguishable artifacts (satisfy
  SCL: \(\Phi\) \(\geq\) \(\lambda\)), OR
\item
  \textbf{Resolve bits} to reduce \(\lambda\) (satisfy SCL: q + \(\Phi\)
  \(\geq\) R)
\end{enumerate}

There is no third option - these are the only ways to correctly solve
L*.

\textbf{Paradigm-Specific Manifestations:}

The collision mechanism appears differently in each computational
paradigm:

\begin{itemize}
\item
  \textbf{Backtracking}

  \begin{itemize}
  \tightlist
  \item
    Artifact Type: Tree nodes
  \item
    Collision Consequence: Must explore \(\geq\) 2\^{}\(\lambda\)
    branches; pruning incorrect paths \(\to\) exponential tree
  \end{itemize}
\item
  \textbf{Dynamic Programming}

  \begin{itemize}
  \tightlist
  \item
    Artifact Type: Table keys
  \item
    Collision Consequence: Must store \(\geq\) 2\^{}\(\lambda\) distinct
    keys; merging seeds \(\to\) wrong subproblem solutions
  \end{itemize}
\item
  \textbf{OBDD}

  \begin{itemize}
  \tightlist
  \item
    Artifact Type: Diagram width
  \item
    Collision Consequence: Must maintain \(\geq\) 2\^{}\(\lambda\) nodes
    at cut; reducing width \(\to\) lost distinctions
  \end{itemize}
\item
  \textbf{Resolution}

  \begin{itemize}
  \tightlist
  \item
    Artifact Type: Proof clauses
  \item
    Collision Consequence: Must derive \(\geq\) 2\^{}\(\lambda\) width;
    narrower proofs \(\to\) semantic gaps
  \end{itemize}
\item
  \textbf{TM}

  \begin{itemize}
  \tightlist
  \item
    Artifact Type: Time steps
  \item
    Collision Consequence: Must spend \(\geq\) 2\^{}\(\lambda\) / B time
    to distinguish via reads (B = per-step bandwidth)
  \end{itemize}
\end{itemize}

All paradigms face the same underlying constraint (q + \(\Phi\) \(\geq\)
R), but express it in their native metric.

\textbf{Mathematical Necessity:}

This is not a statement about algorithmic limitations - it is a
statement about problem structure. L*\textquotesingle{}s construction
(via \textbf{A1-A5} properties) creates a situation where:

\begin{itemize}
\tightlist
\item
  2\^{}(R-q) remaining possible worlds exist at the cut
\item
  Each world corresponds to a distinct seed
\item
  Each seed determines a disjoint address set (by F\_overlay
  injectivity)
\item
  Correctness requires accessing the right addresses for the right world
\end{itemize}

Representing 2\^{}(R-q) distinct computational states with fewer than
2\^{}(R-q) artifacts \(\to\) collision \(\to\) wrong answers. SCL (q +
\(\Phi\) \(\geq\) R) is the mathematical expression of this necessity.

\textbf{Gap Summary for This Instance:}

\textbf{Verification (\(\lambda\)=0) vs Search (\(\lambda\)=5)
Comparison:}

\begin{itemize}
\item
  \textbf{Bottleneck residual}

  \begin{itemize}
  \tightlist
  \item
    Verification: 0 bits
  \item
    Search: 5 bits
  \end{itemize}
\item
  \textbf{Required artifacts}

  \begin{itemize}
  \tightlist
  \item
    Verification: 2\^{}0 = 1
  \item
    Search: 2\^{}5 = 32
  \item
    Ratio: 32×
  \end{itemize}
\item
  \textbf{Operations}

  \begin{itemize}
  \tightlist
  \item
    Verification: 51
  \item
    Search: \textasciitilde{}960
  \item
    Ratio: \textasciitilde{}19×
  \end{itemize}
\item
  \textbf{Paradigm manifestation}

  \begin{itemize}
  \tightlist
  \item
    Verification: Single path
  \item
    Search: 32 DP keys / tree nodes / OBDD width
  \item
    Ratio: 32×
  \end{itemize}
\end{itemize}

The 19× operation ratio is smaller than the 32× artifact ratio because
failed attempts terminate early, but the fundamental gap is exponential
in \(\lambda\).

\subparagraph{2.6.5 Generalization: From Toy Instance to Full
L*}\label{265-generalization-from-toy-instance-to-full-l}

This 3-node, 8-bit instance demonstrates the SCL mechanism in miniature.
Scaling to full L* instances amplifies the same principles:

\textbf{Scaling Parameters:}

\textbf{Toy Instance vs QP-Sharp L* vs Flat L*:}

\begin{itemize}
\item
  \textbf{Nodes}

  \begin{itemize}
  \tightlist
  \item
    Toy Instance: 3
  \item
    QP-Sharp L*: O(log n) depth
  \item
    Flat L*: O(log n) depth
  \end{itemize}
\item
  \textbf{Total bits R\_total}

  \begin{itemize}
  \tightlist
  \item
    Toy Instance: 8
  \item
    QP-Sharp L*: \(\Theta\)(n log n)
  \item
    Flat L*: \(\Theta\)(n log n)
  \end{itemize}
\item
  \textbf{Bottleneck residual \(\lambda\)}

  \begin{itemize}
  \tightlist
  \item
    Toy Instance: 5
  \item
    QP-Sharp L*: \(\Theta\)(log² n)
  \item
    Flat L*: \(\Theta\)(n)
  \end{itemize}
\item
  Required artifacts 2\^{}\(\lambda\)

  \begin{itemize}
  \tightlist
  \item
    Toy Instance: 32
  \item
    QP-Sharp L*: n\^{}(\(\Theta\)(log n))
  \item
    Flat L*: 2\^{}(\(\Theta\)(n))
  \end{itemize}
\item
  \textbf{Verification time}

  \begin{itemize}
  \tightlist
  \item
    Toy Instance: 51 ops (polynomial)
  \item
    QP-Sharp L*: O(n² log² n)
  \item
    Flat L*: O(n² log² n)
  \end{itemize}
\item
  \textbf{Search time}

  \begin{itemize}
  \tightlist
  \item
    Toy Instance: \textasciitilde{}960 ops
  \item
    QP-Sharp L*: \(\geq\) n\^{}(\(\Omega\)(log n))
  \item
    Flat L*: \(\geq\) 2\^{}(\(\Omega\)(n))
  \end{itemize}
\item
  \textbf{Gap}

  \begin{itemize}
  \tightlist
  \item
    Toy Instance: \textasciitilde{}19×
  \item
    QP-Sharp L*: Quasi-poly
  \item
    Flat L*: Exponential
  \end{itemize}
\end{itemize}

\textbf{What Stays the Same:}

\begin{enumerate}
\def\labelenumi{\arabic{enumi}.}
\tightlist
\item
  \textbf{A1-A5 Properties:} Full L* has the same structural properties
  - Hermeticity (A1), Injectivity (A2), Emergence (A3), Closure (A4),
  Dependency (A5).
\item
  \textbf{Enc and F\_overlay:} The same seed encoding and address
  functions (or equivalent injective functions with similar properties).
\item
  \textbf{Keyedness:} For a fixed node v, different seed histories
  induce disjoint designated address sets within U\_v. Across nodes,
  pools U\_v are disjoint by construction in full L*. This enforces
  collision detection under seed merges.
\item
  \textbf{SCL:} q + \(\Phi\) \(\geq\) R holds at every node, with the
  same proof structure (Lemma 7.I: Injectivity + Keyedness \(\to\)
  Alt\_v \(\geq\) 2\^{}(R\_v - q\_v)).
\end{enumerate}

\textbf{What Scales:}

\begin{enumerate}
\def\labelenumi{\arabic{enumi}.}
\tightlist
\item
  \textbf{DAG Depth:} O(log n) instead of 3, creating longer dependency
  chains.
\item
  \textbf{Per-Node Emergence:} R\_v = \(\Theta\)(log n) bits per node
  instead of 2-3 bits.
\item
  \textbf{Bottleneck Residual:} \(\lambda\)\_base = \(\Theta\)(log² n)
  or \(\Theta\)(n) instead of 5 bits, depending on instance profile.
\item
  \textbf{Address Space:} Larger per-node pools U\_v that are disjoint
  by construction (still polynomial total size).
\item
  \textbf{CNF Size:} Polynomial number of clauses (seed-locked decode
  schema ensures NP membership).
\end{enumerate}

\textbf{Frontier-Gate (FG) Mechanism:}

Full L* includes an additional refinement not present in this toy
instance: \textbf{Frontier-Gate (FG)} wiring (§6.2.8, Appendix C). FG
ensures:

\begin{itemize}
\tightlist
\item
  Tight single-run bounds via digest cost \(\Omega\)(n/W\_min) per
  segment
\item
  High-probability hardness (not just expectation)
\item
  Profile-tight complexity: exactly n\^{}(\(\Omega\)(log n)) for
  QP-sharp, exactly 2\^{}(\(\Omega\)(n)) for flat
\end{itemize}

FG works by wiring GateDigest\_v into seeds at GREQ\_v=1 nodes beyond
the gate horizon, creating additional per-segment costs that compound
exponentially with rollback count. This refinement does not change the
core SCL mechanism - it sharpens the lower bounds to match the
parametric spectrum exactly.

\textbf{Key Insight: The Mechanism Scales Transparently}

The verification-vs-search gap demonstrated on this toy instance
(\(\lambda\)=5 \(\to\) 19× gap) scales exponentially with \(\lambda\):

\begin{itemize}
\tightlist
\item
  \(\lambda\) = O(log n) \(\to\) polynomial time (verification)
\item
  \(\lambda\) = \(\Theta\)(log² n) \(\to\) quasi-polynomial
  n\^{}(\(\Theta\)(log n))
\item
  \(\lambda\) = \(\Theta\)(n) \(\to\) exponential 2\^{}(\(\Theta\)(n))
\end{itemize}

Because the collision mechanism is \emph{structural} (guaranteed by
\textbf{A1-A5}), not heuristic or probabilistic, the same mathematical
necessity applies at all scales. This is why SCL provides
paradigm-invariant lower bounds: the framework captures what
L*\textquotesingle{}s problem structure requires for correctness,
regardless of algorithmic approach.

\textbf{What This Example Teaches:}

\begin{enumerate}
\def\labelenumi{\arabic{enumi}.}
\tightlist
\item
  \textbf{Concrete Grounding:} Abstract principles (q + \(\Phi\)
  \(\geq\) R, keyedness, collision) become tangible with specific
  numbers (0x6A7A, \{0x1E, 0x1F, ...\}, 51 ops).
\item
  \textbf{Verification vs. Search Dichotomy:} The same instance takes 51
  ops with witness, \textasciitilde{}960 ops without - pure search vs.
  resolution distinction.
\item
  \textbf{Collision Mechanism:} Wrong seeds \(\to\) wrong addresses
  \(\to\) detectable failures (no escape hatch).
\item
  \textbf{Exponential Scaling:} \(\lambda\) = 5 \(\to\) 32 artifacts;
  \(\lambda\) = 10 \(\to\) 1024 artifacts; \(\lambda\) = 100 \(\to\)
  2\^{}100 artifacts.
\item
  \textbf{Paradigm Invariance:} The 32-artifact requirement appears as
  backtracking nodes, DP keys, OBDD width, resolution width, or TM time
  - same bottleneck, different metrics.
\end{enumerate}

This walkthrough provides the intuition for the full L* construction
(§6), the abstract SCL theorem (§7), and the complexity separation
Appendix C (FG details). The core insight - that problem structure can
\emph{mathematically force} exponential costs for correct algorithms -
extends from this 3-node toy to the full P \(\neq\) NP result.

\begin{center}\rule{0.5\linewidth}{0.5pt}\end{center}

\paragraph{2.7 The Architectural Insight: Configuration Space vs
Resource
Space}\label{27-the-architectural-insight-configuration-space-vs-resource-space}

The toy example above demonstrates a profound architectural principle
that distinguishes L* from typical complexity-theoretic constructions.
This subsection makes explicit the key insight underlying exponential
complexity despite polynomial resources.

\textbf{The Central Question:}

With only \textbf{polynomial-sized address pools} (Node 2 has 16
addresses: 0x10-0x1F; Node 3 has 16 addresses: 0x20-0x2F; total 32
addresses for an 8-bit problem), how can the complexity be
\textbf{exponential}?

\textbf{The Answer: Exponentiality Lives in Configuration Space, Not
Resource Space}

The breakthrough is separating two distinct spaces:

\begin{enumerate}
\def\labelenumi{\arabic{enumi}.}
\item
  \textbf{Resource Space (Polynomial):} The number of available
  addresses/memory cells

  \begin{itemize}
  \tightlist
  \item
    Toy instance: 32 addresses total
  \item
    Full L*: poly(n\_core) addresses total
  \item
    Physical constraint: polynomial-sized address pools U\_v
  \end{itemize}
\item
  \textbf{Configuration Space (Exponential):} The number of distinct
  ways to access those resources

  \begin{itemize}
  \tightlist
  \item
    Toy instance: 2\^{}5 = 32 distinct seed values \(\to\) 32 distinct
    address-selection patterns
  \item
    Full L*: 2\^{}\(\lambda\) distinct seed chains \(\to\)
    2\^{}\(\lambda\) distinct permutations
  \item
    Structural property: each seed induces a different permutation over
    the same polynomial pool
  \end{itemize}
\end{enumerate}

\textbf{Piano analogy (worked example context):}

A piano has only 88 keys (polynomial resource space) yet admits
extremely many compositions (an exponential configuration space over
patterns). The keys are finite and fixed; complexity arises from the
diversity of patterns over those keys. Similarly:

\begin{itemize}
\tightlist
\item
  L* ``keys'': polynomial-sized address pools U\_v (e.g., the 32
  addresses 0x10--0x2F in the toy instance)
\item
  L* ``compositions'': 2\^{}\(\lambda\) distinct seed-dependent
  permutations determining which addresses to read
\item
  Consequence: simultaneous realization of all patterns is infeasible;
  with 2\^{}\(\lambda\) patterns, exhaustive search is required
\end{itemize}

\textbf{How It Works Technically:}

Each seed value Seed\_v induces a \textbf{different permutation}
\(\pi\)\_v : L\_v \(\to\) U\_v mapping logical coordinates to the
polynomial pool U\_v:

\begin{itemize}
\tightlist
\item
  Seed\(_1\) = 0x6A7A \(\to\) reads from addresses \{0x1E, 0x1F, 0x1A,
  0x1B, 0x16, 0x17\}
\item
  Seed\(_2\) = 0x6A0D \(\to\) reads from addresses \{0x13, 0x12, 0x17,
  0x16, 0x1B, 0x1A\}
\end{itemize}

Both seeds use addresses from the same pool (0x10--0x1F for Node 2) but
in different patterns. With 2\^{}R possible seed values, there are
2\^{}R distinct permutations over the same polynomial address pool.

\textbf{Why This Creates Exponential Complexity:}

\begin{enumerate}
\def\labelenumi{\arabic{enumi}.}
\item
  \textbf{Full-churn property} (Lemma A.1.F in full L*): Different seeds
  map coordinates to \textbf{substantially different address sets} - at
  least a constant fraction of addresses change between any two distinct
  seeds.
\item
  \textbf{Keyedness requirement}: Computational artifacts (DP keys, OBDD
  nodes, resolution clauses) must be \textbf{seed-consistent} for
  correctness. You cannot reuse artifacts computed for Seed\(_1\) when
  working with Seed\(_2\), because they depend on reading from mostly
  different addresses (which contain different salt values, yielding
  different outputs).
\item
  \textbf{No precomputation bypass}: Even reading all poly(n\_core)
  addresses does not help --- the designated subset for a given seed is
  unknown until the seed chain is known, and there are exponentially
  many possible seed chains. Computing correctly for all 2\^{}(R\(-\)q)
  possible seeds requires 2\^{}(R\(-\)q) · poly(n\_core) work ---
  exactly the exponential cost SCL predicts.
\end{enumerate}

\textbf{Where Exponentiality Lives:}

The exponential requirement comes from the number of \textbf{distinct
permutations (configuration space: 2\^{}R possibilities), not the number
of available addresses} (resource space: poly(n\_core) cells). This
sidesteps the apparent limitation "polynomial resources \(\to\)
polynomial complexity" by locating exponentiality in the
\textbf{diversity of arrangements} rather than the \textbf{quantity of
elements}.

\textbf{Key Formalization:}

\begin{itemize}
\tightlist
\item
  \textbf{Resource Space}: \textbar{}U\_v\textbar{} = poly(n\_core) per
  node; \(\Sigma\)\_v \textbar{}U\_v\textbar{} = poly(n\_core) total
\item
  \textbf{Configuration Space}: 2\^{}(R\_v) distinct seed values at node
  v \(\to\) 2\^{}(R\_v) distinct permutations \(\pi\)\_v
\item
  \textbf{Bottleneck}: min\_C \(\Sigma\)\_\{v\(\in\)C\} (R\_v - q\_v) =
  \(\lambda\) requires distinguishing \(\geq\) 2\^{}\(\lambda\)
  configurations
\item
  \textbf{Consequence}: Exponential complexity in \(\lambda\) despite
  polynomial \textbar{}U\_v\textbar{}
\end{itemize}

\textbf{Why This Matters for P \(\neq\) NP:}

Traditional complexity analyses focus on resource quantity (memory size,
circuit depth, communication bits). L* highlights a
configuration-diversity barrier: the problem structure forces
distinguishing exponentially many patterns of access to polynomial
resources. This separation:

\begin{itemize}
\tightlist
\item
  Explains why polynomial resources don\textquotesingle{}t imply
  polynomial time
\item
  Shows hardness can arise from \textbf{permutation-space
  incompressibility}, not just resource scarcity
\item
  Demonstrates that computational asymmetry (verification easy, search
  hard) can be \textbf{structurally engineered} via injectivity +
  emergence + closure
\end{itemize}

This configuration-vs-resource separation underpins L*'s one-way
function construction (§9) and the P \(\neq\) NP proof (§10).

\textbf{Cross-references:}

\begin{itemize}
\tightlist
\item
  Full mathematical treatment: §6.2.3 (seed-dependent permutations),
  Appendix A.1.1 (explicit construction)
\item
  Formal properties: Lemma A.1.F (full-churn), Lemma 6.5 (disjoint
  designated atoms)
\item
  SCL proof using this: §7.2.1 (Consolidated SCL Theorem)
\end{itemize}

\begin{center}\rule{0.5\linewidth}{0.5pt}\end{center}

This concludes the formal framework for the Semantic Conservation Law.
Section 2 has established: the precise SCL statement q + \(\Phi\)
\(\geq\) R (§2.1), the sharp phase transition at \(\lambda\) =
\(\Theta\)(log n) (§2.2), key clarifications on correctness vs.
complexity (§2.3), and how dependencies create multiplicative growth
(§2.4). The detailed example in §2.6 demonstrates these principles
concretely. We now turn to the main technical results and parametric
spectrum of complexity bounds.

References map (quick navigation)

\begin{itemize}
\tightlist
\item
  Per-node SCL: Theorem 7.A (§7.2)
\item
  Cut composition/min-cut: Theorem 7.A.1 (§7.2); Appendix J: Theorem J.1
  and Lemma J.1-Cart
\item
  Time lower bounds: Theorem 7.B (§7.2); Theorem 8.A (single-run FG
  bound); arity (§5.5.1); Frontier-Gate (§6.2.8; App. C.1.1); segment
  counting (App. C.2); FrontierPeak (Cor. 7.1.1); conversions (App. C.4)
\item
  Construction hooks: Disjoint pools (§6.2.3); Seed/Enc + Keyedness/RWA
  (§4.2, §6.2.7); gate horizon (§6.2.10); expander-parity gate
  {[}HOO06{]} (§6.2.10)
\end{itemize}

\subsubsection{3. Main Results and Parametric
Spectrum}\label{3-main-results-and-parametric-spectrum}

Sections 1 and 2 established the conceptual foundation: algorithms
solving L* must satisfy the Semantic Conservation Law q + \(\Phi\)
\(\geq\) R (a correctness requirement from L*\textquotesingle{}s
structure), and provided the formal mathematical framework. Section 2
showed how dependencies create multiplicative growth leading to
exponential bottlenecks measured by residual \(\lambda\) = min\_C
\(\Sigma\)\_\{v\(\in\)C\}(R\_v-q\_v).

\textbf{Purpose of Section 3:} We now state the main theorems and
results precisely, specifying exactly what is proven, for which
computational model, and under what quantifier structure. This section
serves as a reference point - the formal claims that subsequent sections
will prove rigorously.

\textbf{Falsifiability:} Our approach rests on a single structural
assertion---the bottleneck residual \(\lambda\) cannot be compressed to
O(log n) for L* instances with FG wiring (i.e., \(\lambda\)(A,x*)
\(\notin\) O(log n)). This claim is falsifiable: Does there exist a
polynomial-time algorithm compressing 2\^{}(\(\lambda\)\_base)
distinguishable states at the bottleneck cut to poly(n) states? We
provide explicit test instances and verification criteria (§3.6).

\textbf{Roadmap:}

\begin{itemize}
\tightlist
\item
  \textbf{§3.0}: Results at a Glance - main separation theorem,
  model/quantifiers/scope, complexity spectrum
\item
  \textbf{§3.1}: The Mathematics of Search vs Verification - why
  verification is polynomial but search is exponential
\item
  \textbf{§3.2}: The Computational Conservation Principle - formal
  conservation law setup
\item
  \textbf{§3.3}: Cross-Paradigm Manifestations and Formal Connection -
  universal manifestation across models
\item
  \textbf{§3.4}: Scope Summary - what is proven and what is not
\item
  \textbf{§3.5}: OWF Packaging via Deterministic FG (direct
  construction)
\item
  \textbf{§3.6}: Falsifiability: The Compression Barrier
\end{itemize}

\textbf{Relation to the proof:} The theorems stated here are proven
across §§6-10: L*\textquotesingle{}s construction (§6), SCL derivation
(§7), \textbf{per-instance deterministic bounds} (§8.A) with FG details
(Appendix C), Structural OWF construction (§9), and NP-completeness
(§10). Section 3 tells you \emph{what} we\textquotesingle{}ll prove; the
remaining sections show \emph{how}.

\paragraph{3.0 Results at a Glance¹}\label{30-results-at-a-glanceuxb9}

Convention: In this section, n := n\_core (see §1.7).

¹ \emph{Polynomials are measured in \textbar{}x}\textbar{}; notation and
sizing in §4 (Notation). Unless stated otherwise, §§1-8 use n :=
n\_core; §12 restates bounds with n := \textbar{}x*\textbar{}.*

\textbf{Main Separation (OWF Construction).} \textbf{Theorem (P \(\neq\)
NP via OWF).} We construct an unconditional one-way function f from
NP-complete language L*, proving \textbf{P \(\neq\) NP} via the
classical bridge OWF \(\Rightarrow\) FP \(\neq\) FNP \(\Rightarrow\) P
\(\neq\) NP.

\textbf{Construction}: f: r \(\mapsto\) Plant(\(\varphi\), r) with FG
wiring. Every output x*=f(r) has \textbf{per-instance deterministic
bound (Theorem 8.A): witness-finding requires time \(\geq\)
n\^{}(\(\Omega\)(log n)) (QP-sharp) or \(\geq\) 2\^{}(\(\Omega\)(n))
(flat) on any fixed computational run}.

\textbf{Key Innovation (Per-Instance Deterministic vs Distributional):}

\begin{itemize}
\tightlist
\item
  \textbf{Traditional OWFs} (e.g., factoring): Assume
  \emph{average-case} hardness---most instances hard over a distribution
\item
  \textbf{This construction}: Proves \emph{per-instance deterministic}
  hardness---every instance hard on every run (\(\forall\)x*
  \(\forall\)run \(\to\) hard)
\item
  \textbf{Why this matters}: Stronger property survives coin-fixing.
  When a randomized inverter is fixed to deterministic coins, the
  per-instance bound still applies to that specific (instance, run)
  pair, enabling the contradiction without distributional assumptions.
\end{itemize}

\textbf{OWF Security}: For any uniform PPT inverter \ensuremath{\mathcal{I}}, if
Pr\_\{r\(\leftarrow\)D\_n, coins\}{[}\ensuremath{\mathcal{I}}(f(r)) \(\in\) f\^{}(-1)(f(r)){]}
\(\geq\) 1/poly(n), coin-fixing yields coins c- with non-negligible
success over r\(\leftarrow\)D\_n. Composing \ensuremath{\mathcal{I}}\_\{c-\} with Ext (see
§9.3) maps inversion \(\to\) witness in poly-time on some x* = f(r*),
contradicting the per-instance bound. Therefore f is one-way against
classical PPT.

Since L* is NP-complete (§10) and f is one-way (§9), this establishes P
\(\neq\) NP for classical uniform PPT via Proposition 10.4.
\(\blacksquare\)

\textbf{Model, Quantifiers \& Scope (This Paper):}

\textbf{Computational Model:}

\begin{itemize}
\tightlist
\item
  \textbf{Primary (OWF path)}: Classical uniform PPT (probabilistic
  Turing machines) with constant k (tapes) and
  \textbar{}\(\Gamma\)\textbar{} (alphabet)

  \begin{itemize}
  \tightlist
  \item
    Per-run deterministic bounds apply to any fixed coin string
  \item
    \textbf{Randomized adversaries covered via coin-fixing (Yao)}: any
    PPT inverter reduces to deterministic run
  \end{itemize}
\item
  \textbf{Per-step bandwidth}: B =
  k\(\lceil\)log\(_2\)\textbar{}\(\Gamma\)\textbar{}\(\rceil\) = O(1)
  bits per step (Lemma 5.5.1); information-theoretic per-run bounds
\item
  \textbf{No advice, no oracles, no hidden channels}: Fresh information
  enters only via first-use designated reads (Hermeticity; Axiom A1)
\item
  \textbf{Quantum adversaries}: Out of scope (different combination
  rules; §12.4)
\end{itemize}

\textbf{Quantifier Structure:}

\textbf{Main (OWF Construction):}

\begin{itemize}
\tightlist
\item
  For FG-wired instances: \textbf{\(\forall\)x*\(\in\)L*\_\{FG\},
  \(\forall\) uniform algorithm A} \(\to\) time \(\geq\) super-poly
\item
  This quantifier form (every instance hard for every uniform algorithm)
  enables the Structural OWF construction (§9)
\item
  Uniform algorithms cannot hardcode instance-specific solutions, so the
  \(\forall\)x* claim is valid under uniform restriction
\item
  Yields: f one-way against uniform PPT \(\to\) FP \(\neq\) FNP \(\to\)
  P \(\neq\) NP via classical bridge
\end{itemize}

\textbf{Scope Boundaries:}

\begin{itemize}
\tightlist
\item
  \textbf{Proven for}: Classical uniform PPT (§9 Structural OWF
  construction)

  \begin{itemize}
  \tightlist
  \item
    Per-run deterministic bounds (Theorem 8.A) apply to any fixed
    computational trace
  \item
    Randomized PPT adversaries covered via \textbf{coin-fixing (Yao)}:
    any successful randomized inverter yields successful deterministic
    run
  \item
    Extractor maps inversion \(\to\) witness in poly-time, contradicting
    per-instance bound
  \end{itemize}
\item
  \textbf{Core mechanism}: Structural OWF construction (§9) from
  per-instance deterministic bounds (§8.A)

  \begin{itemize}
  \tightlist
  \item
    f: r \(\mapsto\) Plant(\(\varphi\), r) with FG wiring; every output
    has deterministic witness-finding lower bound
  \item
    Classical bridge: OWF \(\Rightarrow\) FP \(\neq\) FNP
    \(\Rightarrow\) P \(\neq\) NP (§10.3)
  \end{itemize}
\item
  \textbf{Implies}: P \(\neq\) NP unconditionally (§10.5)
\item
  \textbf{Adapters mentioned}: RAM, Boolean circuits, other paradigms
  (§5, §7.3)
\item
  \textbf{Not proven for}: Quantum computation (different combination
  rules; §12.4), non-uniform models with advice
\end{itemize}

Detailed model specification: §4 (Model \& Framework); quantifier
justification: §8 (per-instance bounds Theorem 8.A).

\begin{itemize}
\tightlist
\item
  Cross-paradigm artifacts (any correct solver of L*). At the
  bottleneck, artifacts \(\geq\) 2\^{}(\(\lambda\)(A,x)):

  \begin{itemize}
  \tightlist
  \item
    Backtracking tree size, DP distinct keys
  \item
    OBDD width and diagram size (order-robust only with expander-FG
    gates)
  \item
    Resolution width and proof size (via width\(\to\)size)
  \end{itemize}
\item
  Parameterization. Bounds scale with the run-dependent bottleneck
  residual \(\lambda\)(A,x) (precise definition deferred). For brevity
  we write \(\lambda\) when unambiguous, and \(\lambda\)\_base for the
  instance profile residual (§4.3; Part V notation recap).
\item
  Randomization. Distributional lower bounds imply worst-case via
  Yao\textquotesingle{}s principle (averaging and coin-fixing).
\item
  Tightness and NP. Time conversions use arity-bounded analysis;
  single-run tightness via instance-side calibrations (Appendix C) while
  preserving NP membership.
\end{itemize}

\textbf{Informal Main Statement.} We construct an NP-complete language
L* with specific structural properties (\textbf{A1-A5}: Hermeticity,
Injectivity, Emergence, Closure, Dependency), then prove these
qualitative properties mathematically imply a quantitative conservation
bound - any algorithm solving L* must satisfy q + \(\Phi\) \(\geq\) R.
The non-trivial contribution is deriving this precise quantitative
relationship from structural properties:

q + \(\Phi\) \(\geq\) R (where \(\Phi\) = log\(_2\)(distinct\_states))
This holds at every node throughout the computation, where q is
information resolved, distinct\_states are simultaneously
distinguishable artifacts maintained, and R is the
node\textquotesingle{}s information requirement.

\textbf{Key Principle:} At any node, resolving q\_v bits reduces
2\^{}R\_v initial possibilities to 2\^{}(R\_v-q\_v) remaining worlds.
Correctness requires at least 2\^{}(R\_v-q\_v) simultaneously
distinguishable artifacts. Maintaining fewer \(\Rightarrow\) collision
between inequivalent worlds (different seeds \(\to\) wrong designated
addresses) \(\Rightarrow\) wrong answer (§1.1).

This yields a parametric spectrum of lower bounds:

\textbf{Parametric Spectrum: Profile-Dependent Complexity}

\begin{itemize}
\tightlist
\item
  \textbf{Verification}: Bottleneck residual \(\lambda\)\_base = 0;
  Lower bound: Polynomial; Description: All bits resolved (q\_v = R\_v)
\item
  \textbf{QP-sharp}: Bottleneck residual \(\lambda\)\_base =
  \(\Theta\)(log² n); Lower bound: n\^{}\(\Theta\)(log n); Description:
  Optimal quasi-polynomial
\item
  \textbf{\ensuremath{\sqrt{n}}-profile}: Bottleneck residual \(\lambda\)\_base =
  \(\Theta\)(\ensuremath{\sqrt{n}}); Lower bound: 2\^{}\(\Theta\)(\ensuremath{\sqrt{n}}); Description:
  Sub-exponential
\item
  \textbf{Exponential}: Bottleneck residual \(\lambda\)\_base =
  \(\Theta\)(n); Lower bound: 2\^{}\(\Theta\)(n); Description: Full
  exponential
\end{itemize}

These bounds apply to all algorithms that correctly solve L*. The
paradigm-specific manifestations (backtracking trees, DP tables, OBDD
widths, Resolution proofs) all reflect the same underlying conservation
law - each paradigm must maintain sufficient distinguishable artifacts
to avoid conflating distinct computational paths.

\textbf{Instance construction:} L* was designed with specific structural
properties (\textbf{A1-A5}) that create \textbf{per-instance
deterministic witness-finding lower bounds} (Theorem 8.A).
\textbf{Frontier-Gate (FG) wiring} ensures every instance exhibits
super-polynomial hardness for any fixed computational run - this enables
the Structural OWF construction (§9) where f(r) outputs FG-wired
instances.

\paragraph{3.1 The Mathematics of Search vs
Verification}\label{31-the-mathematics-of-search-vs-verification}

\textbf{The fundamental dichotomy for L*:} Algorithms solving L* face a
stark choice at each node v with R\_v bits of information:

\begin{itemize}
\tightlist
\item
  \textbf{Verification mode:} Already knows the answer \(\to\) resolves
  all R\_v bits \(\to\) polynomial cost
\item
  \textbf{Search mode:} Does not know \(\to\) must track 2\^{}(R\_v)
  possibilities \(\to\) exponential cost
\end{itemize}

\textbf{Why multiplication occurs:} At node v with R\_v\(-\)q\_v
unresolved bits, there are 2\^{}(R\_v\(-\)q\_v) still-possible worlds
that must map to distinct, seed-consistent artifacts. Due to
dependencies (child nodes need parent outputs), these requirements
multiply as computation progresses toward the bottleneck: by
multiplicative composition (Theorem 7.A.1), the total number of
simultaneously distinguishable artifacts must reach at least 2\^{}((sum
of the unresolved contributions along the hardest route)).

\textbf{State multiplication across dependencies:} When nodes depend on
each other (as in our DAG), unresolved possibilities multiply:

\begin{itemize}
\tightlist
\item
  \textbf{Verification:} Costs add: R\(_1\) + R\(_2\) + R\(_3\) =
  polynomial
\item
  \textbf{Search:} Distinguishable artifacts multiply: 2\^{}(R\(_1\)) ×
  2\^{}(R\(_2\)) × 2\^{}(R\(_3\)) = exponential
\end{itemize}

This multiplication - where unknown information at multiple nodes forces
tracking exponentially many simultaneously distinguishable artifacts -
is precisely what SCL captures. Resolving bits (via designated reads)
reduces this burden exponentially: each bit resolved halves the
artifacts needed.

\paragraph{3.2 The Computational Conservation
Principle}\label{32-the-computational-conservation-principle}

\textbf{The Conservation Principle:} For L*, when R structurally
required bits must be satisfied (A3: rank-based emergence):

\textbf{q + log\(_2\)(distinct\_states) \(\geq\) R}

Resolving q bits reduces 2\^{}R possibilities to 2\^{}(R-q) remaining
worlds. Correctness requires representing these with \(\geq\) 2\^{}(R-q)
artifacts - algorithms must either resolve more or maintain more
simultaneously distinguishable artifacts.

This principle is mathematically necessary for L*: representing
2\^{}(R-q) distinct possibilities with fewer than 2\^{}(R-q) artifacts
causes collisions. When distinct worlds map to the same artifact, the
algorithm cannot distinguish between them \(\to\) wrong answers on some
L* instances. The requirement comes from L*\textquotesingle{}s problem
structure, not algorithmic limitations.

\textbf{Two operational categories to satisfy the conservation law:}

\begin{enumerate}
\def\labelenumi{\arabic{enumi}.}
\tightlist
\item
  \textbf{Resolve bits (q term):} Reduce uncertainty through:

  \begin{itemize}
  \tightlist
  \item
    \textbf{Resolution}: Sequential reads from designated addresses
  \item
    \textbf{Elimination}: Testing/pruning wrong candidates
  \end{itemize}
\item
  \textbf{Maintain distinguishable artifacts (\(\Phi\) term):}

  \begin{itemize}
  \tightlist
  \item
    \textbf{Spatial} (Storage): DP tables, OBDD nodes
  \item
    \textbf{Temporal} (branches): Backtracking, restarts
  \end{itemize}
\end{enumerate}

This maps to §1.1\textquotesingle{}s three ways: Resolution and
Elimination both increase q; Storage increases \(\Phi\).

Key insight: States (i.e., simultaneously distinguishable artifacts) and
branches count the same at the bottleneck cut C* - max(s,b), not s+b.
Rationale: we count simultaneously distinguishable families at the same
frontier (cut C*); parallel tracks do not add beyond the widest
frontier.

\textbf{The algorithmic spectrum:} Pure strategies at extremes, hybrids
in between:

\begin{itemize}
\tightlist
\item
  \textbf{Pure verification:} q = R, states (artifacts) = 1 \(\to\)
  polynomial
\item
  \textbf{Pure memoization:} q = 0, states (artifacts) = 2\^{}R \(\to\)
  exponential space
\item
  \textbf{Pure search:} q = 0, branches = 2\^{}R \(\to\) exponential
  time
\item
  \textbf{Optimal hybrid:} For L* (not a general claim), if the
  effective residual \(\rho\) on the bottleneck cut is \(\Theta\)(log²
  n), artifacts are 2\^{}\(\rho\) = n\^{}(\(\Theta\)(log n\_core)); if
  \(\rho\) = \(\Theta\)(log n\_core), they are polynomial. (\(\rho\) is
  the per-run effective residual on the bottleneck cut; see §7.3.3.)

  \begin{itemize}
  \tightlist
  \item
    Example hybrid (illustration only): If we resolve \ensuremath{\sqrt{R}} bits and
    maintain 2\^{}\ensuremath{\sqrt{R}} states: (R-\ensuremath{\sqrt{R}}) + log\(_2\)(2\^{}\ensuremath{\sqrt{R}}) = (R-\ensuremath{\sqrt{R}}) + \ensuremath{\sqrt{R}} =
    R {[}YES{]}
  \item
    Demonstrates how the principle governs the complexity landscape
  \end{itemize}
\end{itemize}

\paragraph{3.3 Cross-Paradigm Manifestations and Formal
Connection}\label{33-cross-paradigm-manifestations-and-formal-connection}

\textbf{How the conservation principle appears in different models:}

\begin{itemize}
\tightlist
\item
  \textbf{SAT solvers (§5.4):} Balance learning (unit propagation) vs
  states (branches/restarts)
\item
  \textbf{Dynamic programming (§5.3-§5.4):} Trade space (state tables)
  for time (recomputation)
\item
  \textbf{Resolution/CDCL (§5.4):} Width\(\to\)size; clauses reflect
  tracked variables
\item
  \textbf{OBDD/ROBDD (§5.3, §6.2.10):} Width at some level; order-robust
  with expander-FG gates
\item
  \textbf{Deterministic k-tape TM (Appendix C; §5.5):} Segment counting
  and per-segment pricing
\end{itemize}

\textbf{Analogous principles in other models (not proven for L*):}

\begin{itemize}
\tightlist
\item
  \textbf{Streaming:} Limited memory forces multiple passes or
  approximation
\item
  \textbf{Communication:} Bits sent = learning, protocol states =
  simultaneously distinguishable artifacts
\end{itemize}

\textbf{Formal connection to SCL:}

The informal Conservation Principle maps directly to our formal Semantic
Conservation Law:

\begin{itemize}
\tightlist
\item
  \textbf{Informal:} q + log\(_2\)(distinct\_states) \(\geq\) R
\item
  \textbf{Formal:} q + \(\Phi\) \(\geq\) R (where \(\Phi\) =
  log\(_2\)(Alt), the log of alternatives)
\end{itemize}

These are the same statement - Section 3 builds intuition, Section 4
provides proof machinery. For the Structural OWF construction (§9),
\textbf{every FG-wired instance} forces any algorithm to satisfy this
conservation via per-instance deterministic bounds (Theorem 8.A).

\textbf{Implications:}

\begin{itemize}
\tightlist
\item
  \textbf{For L*:} P-membership allows polynomial verification;
  NP-hardness of witness-finding forces exponential states when
  searching without witness
\item
  \textbf{Algorithm design:} Optimize allocation within the conservation
  constraint
\item
  \textbf{Lower bounds unity:} Different techniques measure the same
  underlying conservation
\end{itemize}

\textbf{Scope:} Classical computation only. Quantum computation has
different combination rules (see §12.4).

\begin{center}\rule{0.5\linewidth}{0.5pt}\end{center}

This concludes our presentation of the main results and the conservation
principle that underlies them. We\textquotesingle{}ve shown how
L*\textquotesingle{}s structure makes observable the fundamental
tradeoff between information acquisition and state maintenance that
manifests across the computational paradigms we analyze (backtracking,
DP, OBDD, resolution, TMs). In Part II, we develop the formal framework
that makes these intuitions mathematically precise.

\paragraph{3.4 Scope Summary}\label{34-scope-summary}

\textbf{Proven:} SCL (Thm 7.A), min-cut composition (Thm J.1: Cartesian
factoring), L* satisfies A1-A5 axioms (Lemma 4.A) \textbf{Derived:}
Per-paradigm bounds via adapters (§5) and templates (§7.3)
\textbf{Extensions:} Randomization (§7.7), time conversions (Appendix
C.4)

\paragraph{3.5 OWF Packaging via Deterministic
FG}\label{35-owf-packaging-via-deterministic-fg}

We establish P \(\neq\) NP via three complementary results:

\textbf{One-way function construction (§9).}

\textbf{The function:} §9 defines a total length-regular function f by
sampling r \(\in\) D(\(\varphi\)) uniformly from the domain and
outputting x* := f(r), where x* is FG-wired with identity-based digests
(non-leaking: no assignment bits in x*). The domain D(\(\varphi\)) = \{
r \textbar{} WellFormedRandomness(\(\varphi\), r) \} requires
\(\varphi\).satisfies(r.assignment).

\textbf{Security proof:} For any randomized PPT inverter \ensuremath{\mathcal{I}} with
non-negligible success probability: averaging over coins yields fixed
coins c- and instance r* where \ensuremath{\mathcal{I}}\_\{c-\}(f(r*)) = r' \(\in\)
D(\(\varphi\)) with f(r') = x*. By Lemma 9.DOM, r'.assignment satisfies
\(\varphi\). Ext then produces a canonical witness in poly-time,
contradicting the per-instance deterministic bound (Theorem 8.A) on that
fixed run. Note: r' need not equal r* (preimages are non-unique in the
non-leaking model).

\textbf{Result:} Therefore, f is one-way against classical PPT,
establishing \textbf{worst-case FP \(\neq\) FNP} via the Structural OWF
construction.

This OWF packaging does not rely on average-case hardness or {[}IL90{]};
it uses the deterministic FG-based structural lower bounds.

\textbf{Proposition 10.4 (classical bridge):} OWF \(\Rightarrow\) FP
\(\neq\) FNP \(\Rightarrow\) P \(\neq\) NP. Let f be a total,
polynomial-time computable function that is one-way against PPT
inverters. Since FP \(\subseteq\) PPT, f is also one-way against FP
inverters. Define the inversion relation R(y, x) := {[}f(x) = y{]}. Then
R \(\in\) P (verification runs f and checks equality), so inverting f is
an FNP search problem. If P = NP, then FP = FNP and there exists a
polynomial-time inverter Inv with f(Inv(y)) = y for all y in Range(f),
contradicting one-wayness against FP. Therefore FP \(\neq\) FNP and
hence P \(\neq\) NP.

\textbf{Why hardness arises - the permutation-space foundation:}

\textbf{Analogy (not formal proof):} Both cryptography and L* derive
hardness from \textbf{unique-path permutation spaces}:

\begin{itemize}
\tightlist
\item
  \textbf{RSA}: Unique prime factorization creates 2\^{}k factor pairs;
  trap-door structure (multiply easy, factor hard)
\item
  \textbf{Discrete log}: Each exponentiation path is unique; algebraic
  structure (exponentiate easy, discrete log hard)
\item
  \textbf{L*}: Injective encoding creates 2\^{}\(\lambda\) seed chains;
  dependency geometry (verify easy, search hard)
\end{itemize}

The difference is mechanism, not mathematics. Cryptography discovered
hardness through algebraic trap doors. L* constructs it through
\textbf{dependency geometry - the Conservation Law q + \(\Phi\) \(\geq\)
R ensures you cannot compress 2\^{}(R-q) unique paths into fewer
artifacts without collision. No modular arithmetic required; just
injectivity} (distinct paths stay distinct), \textbf{emergence} (new
choices expand the space), and \textbf{closure} (deterministic history
recovery). These three structural properties guarantee incompressibility
unconditionally.

\textbf{What this connects:}

The Conservation Law is a counting argument (pigeonhole:
2\^{}\(\lambda\) possibilities cannot compress into fewer artifacts),
showing how dependency geometry yields the required asymmetry without
number-theoretic assumptions.

\textbf{The structural insight}: Hardness transcends number theory - it
\textbf{can arise as a fundamental property of permutation spaces} with
the right structure. Both RSA (modular arithmetic) and L* (dependency
chains) derive security from unique-path permutation lattices.
Computational asymmetry can arise universally from structure
(injectivity + emergence + closure), not merely from algebraic
accidents. Cryptography discovered hardness through number theory; we
have proven it is structurally inevitable for L*.

\textbf{Architectural foundation:} The key insight is separating
\textbf{configuration space (2\^{}\(\lambda\) distinct seed-dependent
permutations over addresses) from resource space} (polynomial-sized
address pools). Exponential complexity arises from the diversity of
access patterns, not resource quantity - like a piano with 88 keys
enabling infinitely many songs. See §2.7 for the complete explanation
with concrete examples from the toy instance.

Packaging note (Appendix O). For the explicit OWF packaging, we use a
planted-instance sampler that maps domain elements r = (assignment,
gateDigests, structuralSalt) \(\in\) D(\(\varphi\)) to instances x* with
identity-based digests (non-leaking: no assignment bits in x*). The
domain D(\(\varphi\)) requires \(\varphi\).satisfies(r.assignment), so
any valid preimage contains a satisfying assignment (Lemma 9.DOM).
Preimages are non-unique in this model---multiple satisfying assignments
can produce valid domain elements---but security is preserved: finding
any r' \(\in\) D(\(\varphi\)) with f(r') = x* requires solving SAT. This
makes the inverter\(\Rightarrow\)witness-finder reduction immediate
under the deterministic single-run FG bounds (no distributional
assumption).

\paragraph{3.6 Falsifiability: The Compression
Barrier}\label{36-falsifiability-the-compression-barrier}

\emph{For a quick intuitive overview, see
\texttt{docs/TRAPDOOR\_OWF\_MECHANISM.md}.}

\subparagraph{Canonical Forms (definition and
scope)}\label{canonical-forms-definition-and-scope}

\begin{itemize}
\tightlist
\item
  Canonical witness \texttt{W(x*)}. A single deterministic serialization
  of the witness used by both the verifier and the extractor:
  \texttt{W\ =\ (w,\ G\_tau,\ Dig\_tau)} with parent lists sorted;
  within each path \texttt{P} gates are listed in topological order;
  \texttt{G\_tau} entries follow the published path order; fixed field
  layouts. Verifier \texttt{V} accepts only this form; \texttt{Ext}
  outputs exactly this form in polynomial time.
\item
  Canonical dependency/seed encoding. Parent lists are sorted under a
  fixed ordering; \texttt{Seed\_v} is content-addressed from a
  versioned, self-delimiting \texttt{Enc} with fixed endianness/widths.
  The encoding forbids equivalent multi-representations, ensuring
  parseable, unique reconstructions.
\item
  Canonical decode schema \texttt{Phi\_tilde}. A global order over
  decode slots \texttt{T\_dec} with predetermined sizes. The schema
  defines how identity-based digests (GateDigest values) are computed
  from emergent configurations. The domain D(\(\varphi\)) requires
  WellFormedRandomness: correct parities matching the assignment.
\item
  Poly-time normalizer. There is a polynomial-time procedure that
  converts any valid representation to its canonical form without
  changing semantics. Falsifiability (the compression test) is evaluated
  against these canonical forms to avoid spurious ``wins'' from
  alternative encodings.
\end{itemize}

Why this matters here. The compression criterion in §3.6 asks whether a
solver can collapse \(\geq\) 2\^{}\(\lambda\) distinguishable states to
polynomially many artifacts on a fixed run. Canonicalization pins the
target: counts, comparisons, and extractor outputs are measured against
one unique, checkable representation per instance, eliminating
ambiguity.

Our approach rests on a single structural claim about the L* language:

\textbf{Core Claim (Incompressibility of Bottleneck Residual).} There
exists an explicit infinite family \{x*\emph{n \(\in\) L*}\{FG\}(n)\}
constructed in §6 such that for every uniform probabilistic
polynomial-time (PPT) algorithm A there exist c \textgreater{} 0 and N
with the following property: for all n \(\geq\) N, for every
deterministic execution of A on x*\_n that outputs a valid canonical
witness W (including any specific coin-fixing for randomized A), the
min-cut residual satisfies \(\lambda\)(execution, x*\_n) \(\geq\) c ·
\(\lambda\)\_base(n). In particular, no successful execution achieves
\(\lambda\)(execution, x*\_n) \(\in\) O(log n) for all sufficiently
large n.

\textbf{The P\(\neq\)NP Question as a Compression Problem.} For our
construction, the complexity question "Is P=NP?" reduces to:

\textbf{Does there exist a polynomial-time algorithm that compresses the
exponential configuration space at L*\textquotesingle{}s bottleneck cut
from 2\^{}(\(\lambda\)\_base) distinguishable states down to poly(n)
states?}

Our proof asserts: \textbf{No such compression exists.}

This assertion has three components, each independently falsifiable:

\begin{enumerate}
\def\labelenumi{\arabic{enumi}.}
\tightlist
\item
  \textbf{Existence}: the bottleneck cut C* with residual
  \(\lambda\)\_base exists (by construction; §6.2).
\item
  \textbf{Necessity}: any correct solver must account for
  2\^{}(\(\lambda\)(A, x*)) distinguishable artifacts across C* (Theorem
  7.A (§7.2.1); Theorem J.1-PROD).
\item
  \textbf{Persistence}: no algorithm can compress \(\lambda\) to O(log
  n) (Theorem 8.A via Frontier-Gate).
\end{enumerate}

\textbf{Falsification Criterion.}

\textbf{Definition 3.6.1 (Compression Algorithm).} A compression
algorithm for L* is a uniform probabilistic polynomial-time (PPT)
algorithm A such that for the family \{x*\emph{n \(\in\) L*}\{FG\}(n)\}
constructed in §6, A (via some deterministic execution strategy,
including specific coin-fixing for randomized A) outputs a canonical
witness W = (w, G\_\(\tau\), Dig\_\(\tau\)) while achieving min-cut
residual:

\(\lambda\)(execution, x*\emph{n) = min\_C \(\Sigma\)}\{v\(\in\)C\}
max\{0, R\_v \(-\) q\_v(execution)\} = o(\(\lambda\)\_base(n))

\textbf{Where:}

\begin{itemize}
\tightlist
\item
  \textbf{execution}: a specific deterministic run of A on x*\_n (for
  randomized A, this means a specific coin-fixing)
\item
  \textbf{q\_v(execution)}: bits resolved at node v via first-use
  designated reads (RWA accounting; §4.2)
\item
  \textbf{little-o notation}: with respect to the sequence parameter n
  (i.e., the bound holds for all sufficiently large n)
\end{itemize}

An execution succeeds if it outputs a valid canonical witness W. For
randomized A with Pr\_coins{[}A(x*\_n) outputs valid W{]} \(\geq\) 2/3:

\begin{itemize}
\tightlist
\item
  By averaging over coin strings, there exists a coin-fixing where A
  succeeds
\item
  We measure \(\lambda\) on such a successful deterministic run
  (coin-fixing; §9.4)
\item
  Coins are sampled fresh per run
\item
  The \(\geq\) 2/3 requirement precludes embedding per-instance advice
  into rare coin strings
\end{itemize}

The accounting enforces:

0 \(\leq\) q\_v(execution) \(\leq\) R\_v

\(\lambda\)(execution, x*\emph{n) = min\_C \(\Sigma\)}\{v\(\in\)C\}
max\{0, R\_v \(-\) q\_v(execution)\}

\textbf{Theorem 3.6.1 (Falsification Criterion).} The approach is
falsified by exhibiting:

\begin{itemize}
\tightlist
\item
  A uniform probabilistic polynomial-time (PPT) algorithm A
\item
  A deterministic execution strategy (including a specific coin-fixing
  for randomized A)
\end{itemize}

Such that, for the family \{x*\_n\} constructed in §6, the execution
achieves:

\(\lambda\)(execution, x*\_n) = o(\(\lambda\)\_base(n)) on infinitely
many instances

That is: \(\forall\)c \textgreater{} 0 \(\exists\)N \(\forall\)n
\(\geq\) N: \(\lambda\)(execution on x*\_n) \textless{} c ·
\(\lambda\)\_base(n)

\textbf{Stronger sufficient condition:} \(\lambda\)(execution, x*\_n) =
O(log n) would also falsify the claim.

\textbf{Two paths to falsification.}

The P vs NP question from the compression perspective: \textbf{Can
exponential configuration space be compressed into polynomial resource
space, independent of representation?}

\textbf{Key terms:}

\begin{itemize}
\tightlist
\item
  \textbf{Configuration space}: The 2\^{}\(\lambda\) seed-consistent
  computational histories that must remain distinguishable for
  correctness
\item
  \textbf{Resource space}: The number of simultaneously distinguishable
  computational states (Alt = 2\^{}\(\Phi\)) the algorithm maintains
\item
  \textbf{Compression}: Achieving \(\Phi\) = O(log n), equivalently Alt
  = 2\^{}\(\Phi\) = poly(n)
\item
  \textbf{Representation-independence}: The compression principle
  applies across standard polynomial-time many-one encodings of the same
  NP-complete problem, not just L*\textquotesingle{}s specific encoding
\end{itemize}

Discovering a compression algorithm has fundamentally different
implications depending on whether it is representation-independent:

\textbf{Type 1 (Representation-Invariant Compression)}: A uniform PPT
algorithm proves that \textbf{configuration space = resource space}
representation-independently by establishing a correctness-preserving
compression principle that:

\begin{itemize}
\tightlist
\item
  Maps exponentially many configuration alternatives (2\^{}\(\lambda\)
  seed-consistent histories)
\item
  Into polynomially many simultaneously distinguishable computational
  states
\item
  Achieves Alt = poly(n), equivalently \(\Phi\) = O(log n) where
  \(\Phi\) = log\(_2\)(Alt)
\item
  Works across standard polynomial-time many-one encodings of a
  canonical NP-complete problem
\end{itemize}

This demonstrates that the configuration-resource equivalence holds
independent of representation. If such representation-invariant
compression exists, it falsifies our claim and (via the classical
bridge) implies P = NP in the classical uniform model (§4.0).

\textbf{Type 2 (Design-Specific Bypass)}: The algorithm exploits a
\textbf{loophole in THIS particular L* construction}:

\begin{itemize}
\tightlist
\item
  Examples: seed circularity shortcuts, FG gate budget bypass, gaps in
  A1--A5 allowing merges
\item
  \textbf{Implication}: L* construction needs refinement, NOT that P=NP
\item
  The construction can be patched by:

  \begin{itemize}
  \tightlist
  \item
    Closing the loophole (tightening axioms, adjusting FG parameters, or
    fixing the witness encoding)
  \item
    Re-running the proof with the improved construction
  \end{itemize}
\end{itemize}

\textbf{Distinguishing the types}: Test whether the compression
principle is \textbf{representation-invariant}.

\textbf{Type 1 indicators} (fundamental \(\to\) implies P = NP):

\begin{itemize}
\tightlist
\item
  Algorithm achieves \(\Phi\) = O(log n) (equivalently Alt = poly(n)
  states)
\item
  Works across standard polynomial-time many-one encodings of a
  canonical NP-complete problem
\item
  Proves that configuration space = resource space
  representation-independently
\item
  Falsifies our claim and implies P = NP (uniform)
\end{itemize}

\textbf{Type 2 indicators} (construction loophole \(\to\) needs
refinement):

\begin{itemize}
\tightlist
\item
  Compression depends on L*-specific machinery (seeds, FG, overlay
  addressing)
\item
  Fails when the same NP-complete problem is encoded differently
\item
  Requires construction refinement, not a fundamental breakthrough
\end{itemize}

\textbf{Out of scope}: Methods using non-uniform advice or
quantum/oracle resources are outside our scope (§4.0) and do not imply
classical P = NP.

\textbf{Note}: The Type 1 test is deliberately stronger than what we
formally prove. Our incompressibility is rigorously established for the
fixed L* overlay/encoding; representation-invariant compression across
reduction-stable overlays is discussed in §12 F6.

\textbf{Algorithmic checkability.} This criterion is verifiable by
execution:

\begin{enumerate}
\def\labelenumi{\arabic{enumi}.}
\tightlist
\item
  Run A (or a specific coin-fixing for randomized A) on a growing family
  \{x*\_n\} with n \(\to\) \(\infty\) from the test suite generator
\item
  Track first-use designated reads at each node v to compute
  q\_v(execution) (RWA accounting; requires algorithm instrumentation).
  The harness audits first-use reads from both input and constant pools;
  program constants, PRNG seeds, and preloaded tables are priced under
  \(\Phi\) (distinguishable states) or counted toward q per §4.2.
\item
  Calculate the vertex-capacitated s--t min-cut residual:
  \(\lambda\)(execution, x*\emph{n) = min\_C \(\Sigma\)}\{v\(\in\)C\}
  max\{0, R\_v \(-\) q\_v(execution)\} (in polynomial time via a
  node-splitting reduction: split each node v into v\_in\(\to\)v\_out
  with edge capacity w\_v = max\{0, R\_v \(-\) q\_v(execution)\},
  redirect edges u\_out\(\to\)v\_in, then run a standard s--t min-cut
  algorithm; treat s and t as uncuttable by excluding them or
  equivalently setting c(s) = c(t) = +\(\infty\) in the reduction).
\item
  Verify: Is \(\lambda\)(execution, x*\_n) = o(\(\lambda\)\_base(n))?
\item
  For randomized A: Either (a) apply coin-fixing (Yao\textquotesingle{}s
  principle; §9.4) to obtain deterministic runs and analyze each, or (b)
  empirically verify Pr{[}success{]} \(\geq\) 2/3 via sampling, then use
  averaging to identify a successful coin-fixing to analyze in steps
  1--4
\end{enumerate}

If yes, our Theorem 8.A (per-instance bounds) and Theorem 10.5 (P
\(\neq\) NP) are disproved.

\textbf{Invitation.} We explicitly invite researchers to implement
algorithms attempting to compress \(\lambda\) to O(log n). Methods might
include: algebraic techniques (Gaussian elimination, Gröbner bases),
constraint propagation (CSP, SAT solvers with learning), implicit
representations (BDDs, symmetries), or decode bypass strategies. A
successful compression algorithm---demonstrated empirically by showing
the asymptotic trend \(\lambda\)(execution, x*\_n)/\(\lambda\)\_base(n)
\(\to\) 0 as n grows, or the sufficient condition \(\lambda\)(execution,
x*\_n) = O(log n) across growing instance sizes---would constitute
immediate falsification.

\textbf{Why We Believe Compression is Impossible.}

The claim rests on \textbf{structural necessity} rather than algorithmic
limitations:

\textbf{Multi-Level Barrier.} Compression must simultaneously defeat
three structural layers that are \textbf{coupled by design}:

\textbf{(i) A1-A5 axioms} (Hermeticity, Injectivity, Emergence, Closure,
Dependency) force SCL q+\(\Phi\)\(\geq\)R at overlay level---Cartesian
factorization across cuts (Theorem J.1-PROD) requires 2\^{}\(\lambda\)
distinguishable artifacts for correctness.

\textbf{(ii) FG (Frontier-Gate)} provides per-instance deterministic
bounds (Theorem 8.A)---every FG-wired instance requires super-poly time
on any fixed run, enabling Structural OWF construction.

\textbf{(iii) Witness layer} (NP-complete 3-SAT reduction) requires
searching witness space.

\textbf{Crucially, these layers are interdependent}: OAP
(Overlay-as-Problem; §10.1.1) implements seed-locked decoding where
witness-solving requires overlay engagement---no standalone CNF bypass.
FG wiring (§6.2.8) embeds GateDigest\_v into seeds at GREQ\_v=1 nodes,
coupling gate constraints to seed chains. Breaking any single layer is
insufficient---must compress overlay AND bypass FG AND solve NP-complete
witness simultaneously, but the coupling prevents individual bypass.

\textbf{Tight Empirical Bounds.} Testing shows enumeration-based
algorithms achieve time n\^{}(O(log n)) (QP-sharp) or 2\^{}(O(n))
(flat), exactly matching proven lower bounds n\^{}(\(\Omega\)(log n))
and 2\^{}(\(\Omega\)(n)) (Theorem 8.A). When \(\Omega\)=\(\Theta\), no
gap remains for improvement.

\textbf{Conservation Law.} The inequality q + \(\Phi\) \(\geq\) R (where
\(\Phi\) = log\(_2\) of simultaneously distinguishable artifacts) is a
\textbf{counting argument}: representing 2\^{}(R-q) distinguishable
possibilities requires resolving R-q bits (via reading or testing) or
maintaining 2\^{}(R-q) states.

The conservation inequality (two quantities q and \(\Phi\) summing to
requirement R) offers no bypass---algorithms must pay via: \textbf{(i)
resolution} (increase q through designated reads/tests), \textbf{(ii)
maintained artifacts} (increase \(\Phi\) via simultaneous
distinguishable states), or \textbf{(iii) repeated trials} (affecting
expected time). There is no fourth dimension---the inequality is
absolute per run.

\textbf{Randomization via coin-fixing (Yao):} Any randomized algorithm
decomposes into deterministic runs indexed by coin strings. Each fixed
run satisfies q + \(\Phi\) \(\geq\) R per cut (Theorem 7.2.1), so
\(\lambda\)(A\_coins, x*) \(\geq\) c · \(\lambda\)\_base(n) for some
constant c \textgreater{} 0 on every successful run (as in the core
claim). Randomization affects \textbf{expected trials} (averaging over
coins) but cannot compress the per-run bottleneck---thus \(\lambda\)
incompressibility holds for both deterministic and randomized PPT.

\begin{center}\rule{0.5\linewidth}{0.5pt}\end{center}

\textbf{The central claim---incompressibility of \(\lambda\) for L*---is
falsifiable, algorithmically checkable, and empirically testable.} The
configuration space perspective makes the barrier concrete: can you
compress 2\^{}(\(\lambda\)\_base) states to poly(n)? We assert no, and
provide instances for verification.

\begin{center}\rule{0.5\linewidth}{0.5pt}\end{center}

\subsection{Part II: Building the
Framework}\label{part-ii-building-the-framework}

\textbf{Reader\textquotesingle{}s Guide:} Part II transforms the
intuitive Conservation Principle from §3.2 into the rigorous Semantic
Conservation Law. The informal inequality q\_learned +
log\(_2\)(distinct\_states) \(\geq\) R becomes the formal q + \(\Phi\)
\(\geq\) R, proven through the DAG-Semantic Framework. Section 4
provides the mathematical machinery that makes Section
3\textquotesingle{}s intuitions bulletproof.

\subsubsection{4. Model and Semantic
Framework}\label{4-model-and-semantic-framework}

\paragraph{4.0 Model Scope
(Consolidated)}\label{40-model-scope-consolidated}

\begin{itemize}
\tightlist
\item
  Included: deterministic k-tape TMs with constant k and alphabet
  \textbar{}\(\Gamma\)\textbar; classical uniform algorithms (no
  advice); randomized PPT via coin-fixing (analyzed per fixed run).
\item
  Excluded: non-uniform circuits/advice (hardcoding defeats
  \(\forall\)x); oracle/advice models; quantum algorithms;
  algebrized/algebraic oracle access (see Appendix N).
\item
  Assumptions: read-only input with designated payloads (Hermeticity);
  correctness checked by Algorithm V; instance-side universality (FG
  wiring) is constructed, not assumed.
\end{itemize}

Uniformity references. Our lower bounds apply to uniform algorithms (no
advice strings or hardcoded instances). This restriction is essential
for the \(\forall\)x* universality used by the OWF packaging (§9):
non-uniform families can hardcode per-instance solutions and thus
vacuously satisfy lower-bound statements that quantify only over inputs,
not over description size and construction uniformity. See standard
texts for uniform vs non-uniform models and their implications
{[}AB09{]}, {[}GOL01{]}.

Section 3 stated the main results: P \(\neq\) NP for deterministic
k-tape TMs via an explicit NP-complete language L* whose canonical
witness output task requires super-polynomial time. But what exactly is
this computational model, and how do we rigorously define the semantic
quantities (R\_v, q\_v, \(\Phi\)\_v, \(\lambda\)) that appear in the
Semantic Conservation Law?

\textbf{Purpose of Section 4:} We establish the \textbf{abstract
DAG-Semantic Framework} - the paradigm-invariant mathematical machinery
underlying all our results. This framework separates the
\textbf{semantic core} (what L* structurally requires for correctness)
from \textbf{paradigm adapters} (how these requirements manifest as
concrete resources in different computational models).

\textbf{Roadmap:}

\begin{itemize}
\tightlist
\item
  \textbf{§4.1}: Core concepts (LEARN vs KEEP TRACK, DAG dependencies,
  content-addressed seeds)
\item
  \textbf{§4.2}: The DAG-Semantic Framework (axioms \textbf{A1-A5},
  keyedness, RWA attribution)
\item
  \textbf{§4.3}: Parameters and Profiles (\(\lambda\)\_base, instance
  sizing, QP-sharp vs flat)
\item
  \textbf{§4.4}: Language Definition and Correctness (formal
  specification)
\end{itemize}

\textbf{Framework Scope.} The Structural OWF construction (§9) uses
axioms \textbf{A1-A5} with per-instance properties. Section 4
establishes the A1-A5 core framework.

\textbf{Why this matters for P \(\neq\) NP:} We prove the conservation
law (§7) for the abstract framework, then provide thin adapters showing
how the same residual \(\lambda\) translates into different metrics:
tree branches (backtracking), table keys (DP), width (OBDDs), clause
width (resolution). This architecture ensures that the exponential
bottleneck (\(\lambda\) \(\geq\) \(\omega\)(log n)) cannot be
circumvented by switching computational models - the residual is
paradigm-invariant; only its interpretation changes. This
model-independence grounds the Structural OWF construction (§9) and
enables unconditional separation results.

\textbf{Model extensions:} §12.4 discusses quantum computation and other
models. For TM adapter details: addresses are logical, self-delimiting
record names over designated tapes; disjoint pools via name prefixes;
first-use credit when head enters payload (Appendix D.5).

\textbf{Adapter scope note (post-hoc vs from-scratch).} From-scratch
solving inherits \(\lambda\)-driven blow-ups as branches/keys/width/size
in each model; however, once satisfying assignment w is fixed,
assembling \(\tau\)-budgeted canonical witness G\_\(\tau\) from (x*, w)
is polynomial by design across paradigms (matches §10.4.1) - this does
not weaken from-scratch lower bounds.

Notation (sizes). Throughout the paper:

\begin{itemize}
\tightlist
\item
  \textbf{n\_core} := \textbar{}Encode(\(\varphi\))\textbar{} denotes
  the CNF encoding length of the base predicate \(\varphi\) (number of
  bits to encode \(\varphi\))
\item
  \textbf{n\_tot} := \(\Sigma\)\_v R\_v denotes the total emergence mass
  = \(\Theta\)(n\_core · log n\_core) under the flat profile
\item
  Unless explicitly stated otherwise, \textbf{n} refers to n\_core
\item
  When we write "\(\Sigma\)\_v R\_v = \(\Theta\)(n log n)", this means
  \(\Theta\)(n\_core log n\_core)
\end{itemize}

Notation (symbols). Throughout §4-§10:

\begin{itemize}
\tightlist
\item
  \textbf{\(\Phi\)} (no subscript) denotes the Hartley counter in SCL:
  \(\Phi\) := log\(_2\)(Alt)
\item
  \textbf{\(\Phi\)\textasciitilde{}} denotes the published seed-locked CNF decode schema
  (no standalone CNF table)
\item
  \textbf{\(\varphi\) := Decode\_from\_seeds(\(\Phi\)\textasciitilde{}; Seed\_chain)} is
  the decoded 3-CNF; for any assignment w, \(\varphi\)(w)=1 defines
  acceptance of the base predicate
\end{itemize}

\paragraph{4.1 Core Concepts}\label{41-core-concepts}

We model computation as a DAG where nodes represent computational phases
with dependencies:

\begin{enumerate}
\def\labelenumi{\arabic{enumi}.}
\item
  \textbf{LEARN vs KEEP TRACK:} At each node, algorithms either resolve
  bits (LEARN, increasing q\_v) or maintain multiple possibilities (KEEP
  TRACK, increasing \(\Phi\)\_v)
\item
  \textbf{DAG Structure:} Nodes have dependencies - e.g., node v
  requires parents u\(_1\), u\(_2\) for computation
\item
  \textbf{Content-Addressed Seeds (Seed\_v encodes computation
  history):} Each node\textquotesingle{}s seed deterministically encodes
  its computation history, preventing illegal state merging
\end{enumerate}

This framework will capture (§7.2.1) why distinct computational states
cannot merge without resolution.

\textbf{Definition Block: Core Semantic Concepts}

The following concepts are formalized throughout §4-§7; this box
provides quick reference:

\begin{itemize}
\item
  \textbf{Feasible world (\(\omega\) \(\in\) \(\Pi\)\_v)}: A total
  assignment to problem variables consistent with the transcript prefix
  up to node v and all overlay constraints (see §7.2.1 for setup)
\item
  \textbf{Seed-consistent equivalence (\(\omega\) \(\approx\)\_v
  \(\omega\)')}: Two worlds \(\omega\), \(\omega\)' at node v are
  equivalent if they induce identical designated addresses and outcomes
  at all descendants along the current seed chain (i.e.,
  future-distinguishable behavior is identical). Equivalence classes
  partition feasible worlds: \(\Pi\)\_v / \(\approx\)\_v (full
  definition: §7.2.1 before Theorem 7.A)
\item
  \textbf{Alternatives (Alt\_v)}: The number of seed-consistent
  equivalence classes at node v: Alt\_v := \textbar{}\(\Pi\)\_v /
  \(\approx\)\_v\textbar{} These are the simultaneously distinguishable
  computational artifacts that a correct algorithm must realize to avoid
  collision between inequivalent worlds (representation-invariant; Lemma
  4.2.2)
\item
  \textbf{Potential (\(\Phi\)\_v := log\(_2\)(Alt\_v))}: Hartley entropy
  (Rényi-0; zero-error, worst-case); the logarithmic measure of
  maintained distinguishable states. Counts "bits of residual
  uncertainty" requiring explicit tracking (not Shannon entropy; see
  §1.6)
\item
  \textbf{Required bits (R\_v)}: Information that must emerge at node v,
  defined by rank(H\_v) = R\_v where H\_v is the emergence selector
  matrix (Axiom A3, Emergence; §6.2.5). Creates 2\^{}(R\_v) possible
  values for y\_v = H\_v x\_v
\item
  \textbf{Resolved bits (q\_v)}: Information functionally determined at
  node v through designated reads or testing, measured via
  Receiving-Window Attribution (RWA; defined in next bullet). Quantifies
  "what the algorithm has definitively learned" on the current seed
  chain
\item
  \textbf{Receiving-Window Attribution (RWA)}: Accounting rule that
  charges each bit of information to node v at its first valid use on
  the current seed chain (§4.2; before DAG-Semantic Framework). Re-reads
  or later uses of already-revealed information are not re-charged.
  Combined with Hermeticity (A1), ensures each counted bit corresponds
  to at least one designated read from an address pool
\item
  \textbf{Keyedness (seed-consistency correctness)}: Correctness
  requirement stating that computational artifacts at node v are correct
  if and only if they are seed-consistent - from the
  artifact\textquotesingle{}s content, a verifier can deterministically
  derive exactly the designated addresses F\_overlay(Seed\_v; j,
  \(\ell\)) (addressing function; see §6.2.7 and §2.1.2) for all
  (j,\(\ell\)) in the artifact\textquotesingle{}s scope on the current
  seed chain (§4.2; early in DAG-Semantic Framework). This is a
  correctness condition on reuse, not a storage mandate: any sufficient
  statistic that deterministically yields the correct addresses
  qualifies. All lower bounds count representation-invariant equivalence
  classes under this seed-consistency relation
\item
  \textbf{Min-cut residual (\(\lambda\)(A,x))}: The run-dependent
  bottleneck residual, defined as: \(\lambda\)(A,x) := min\_C
  (\(\Lambda\)(C) - Q(C)) = min\_C \(\Sigma\)\_\{v\(\in\)C\}(R\_v -
  q\_v) where C ranges over all s-t cuts in the DAG. Measures the
  incompressible information at the hardest bottleneck after algorithm A
  has resolved what it can (precise definition: Step 5 of Theorem 7.2.1;
  §7.2.1)
\item
  \textbf{Profile residual (\(\lambda\)\_base)}: Instance-structural
  residual before any algorithmic work, determined by the DAG profile
  (W\_min, r, L). For QP-sharp: \(\lambda\)\_base = \(\Theta\)(log²
  n\_core); for flat: \(\lambda\)\_base = \(\Theta\)(n\_core) (§4.3)
\end{itemize}

\textbf{Key relationships}:

\begin{itemize}
\tightlist
\item
  SCL per-node: q\_v + \(\Phi\)\_v \(\geq\) R\_v (Eq. (2.1); Theorem
  7.A; §7.2)
\item
  SCL across cut: Q(C) + \(\Phi\)(C) \(\geq\) \(\Lambda\)(C) where
  \(\Phi\)(C) = \(\Sigma\)\_\{v\(\in\)C\} \(\Phi\)\_v (Eq. (2.2);
  Theorem 7.A.1; §7.2)
\item
  Time bound: Complexity \(\geq\) 2\^{}(\(\Omega\)(\(\lambda\)(A,x))) ·
  baseline (Theorem 7.B; §7.2)
\end{itemize}

\emph{Detailed formalizations and proofs}: §4.2 (Framework), §6
(Instance Construction), §7 (SCL Theorems).

\paragraph{4.2 The DAG-Semantic
Framework}\label{42-the-dag-semantic-framework}

The \textbf{DAG-Semantic Framework} captures how computation progresses
through structured dependencies. This formalizes the intuition from §3.1
about "selection under constraints" - transforming the conceptual
understanding into precise mathematical machinery. At its core, the
framework models:

\begin{enumerate}
\def\labelenumi{\arabic{enumi}.}
\tightlist
\item
  \textbf{Semantic properties} that create conservation laws:

  \begin{itemize}
  \tightlist
  \item
    \textbf{Emergence}: Intermediate state vectors x\_v \(\in\)
    \{0,1\}\^{}(R\_v) must be discovered at node v
  \item
    \textbf{Completeness}: Full resolution y\_v = H\_v x\_v with
    rank(H\_v) = R\_v ensures complete states
  \end{itemize}
\end{enumerate}

\begin{itemize}
\tightlist
\item
  \textbf{Dependency}: DAG structure via Seed\_v = Enc(v
  \textbar{}\textbar{} sort\{(u, Seed\_u, y\_u) : u \(\in\) P(v)\}
  \textbar{}\textbar{} GateDigest\_v) (GateDigest\_v empty when
  GREQ\_v=0; GREQ\_v is the per-node gate-requirement flag; acyclic:
  digests use pre-horizon indices only - see §6.2.8)
\item
  \textbf{Injectivity}: Enc ensures distinct histories \(\to\) distinct
  seeds \(\to\) cannot merge without resolution
\end{itemize}

\begin{enumerate}
\def\labelenumi{\arabic{enumi}.}
\setcounter{enumi}{1}
\tightlist
\item
  \textbf{Conservation laws} that govern complexity:

  \begin{itemize}
  \tightlist
  \item
    \textbf{Semantic Conservation Law (SCL)}: q + \(\Phi\) \(\geq\) R
    (where \(\Phi\) = log\(_2\)(Alt); holds at each node)
  \item
    \textbf{Semantic Multiplication Principle (SMP)}:
    \textbar{}\(\Pi\)\_\{after v\}\textbar{} \(\geq\)
    \textbar{}\(\Pi\)\_\{before v\}\textbar{} × 2\^{}(R\_v-q\_v)
  \item
    \textbf{Min-cut composition}: \(\lambda\)(A,x) = min\_C
    \(\Sigma\)\_\{v\(\in\)C\}(R\_v - q\_v)
  \end{itemize}
\end{enumerate}

\textbf{Definition 4.2.1 (Semantic Multiplication Principle).}

SCL (engine reference): The per-node conservation law q\_v + \(\Phi\)\_v
\(\geq\) R\_v is stated and proved abstractly in §7 (Theorem 7.A). It is
the logarithmic form of the multiplicative refinement (SMP) below.

\textbf{SMP (Dynamic View):} At node v with requirement R\_v and q\_v
bits already resolved, the partition of feasible worlds refines:
\textbar{}\(\Pi\)\_\{after v\}\textbar{} \(\geq\)
\textbar{}\(\Pi\)\_\{before v\}\textbar{} × 2\^{}(R\_v-q\_v)

\textbf{Multiplicative Growth:} The SMP shows that at each node, after
resolving q\_v bits out of R\_v required, we have 2\^{}(R\_v-q\_v)
remaining distinguishable possibilities that must map to distinct
computational states within each existing partition element. This
multiplication is unavoidable - maintaining fewer than this bound causes
distinct worlds to collide, yielding incorrect computation.

\textbf{Theorem 4.2 (Chain Product Bound).} For any root\(\to\)leaf
chain P in the DAG, \(\Sigma\)\_\{v\(\in\)P\} log\(_2\)(Alt\_v) \(\geq\)
\(\Sigma\)\_\{v\(\in\)P\} (R\_v - q\_v) - (Eq. 4.2a) equivalently the
artifact size along P satisfies \(\Pi\)\_\{v\(\in\)P\} Alt\_v \(\geq\)
2\^{}(\(\Sigma\)\_\{v\(\in\)P\}(R\_v-q\_v)). - (Eq. 4.2b) \emph{Proof.}
Apply SMP iteratively along the chain, charging first-use revelations by
RWA. \(\blacksquare\)

\textbf{Cut Composition:} For any cut C through the DAG: distinguishable
artifacts(C) \(\geq\) 2\^{}(\(\Sigma\)\_\{v\(\in\)C\}(R\_v-q\_v))

\textbf{Per-Instance via H1-H5.} Multiplicativity across a cut uses
\textbf{disjoint designated atoms + injectivity} (H1-H5 properties); see
\textbf{Lemma J.1-Cart (App. J)}. This is realized concretely by the
disjoint address pools U\_v and Unique-Neighbor mapping (§6.2.3)
together with content-addressed seeds Enc (Seed definition §6.2.7).

\textbf{Equivalence (scope):} SMP implies SCL when \(\Phi\)\_v :=
log\(_2\) Alt\_v and Alt\_v counts within-try, simultaneously
distinguishable artifacts that differ in unresolved coordinates at v
(Keyedness forbids merging such worlds). Conversely, SCL is the
logarithmic form of the multiplicative refinement lower bound. Equality
holds when Alt\_v matches the per-node refinement factor. Both express
that unresolved requirements force multiplicative growth of
distinguishable possibilities.

\textbf{Axioms A1-A5 (Main Framework; formalized in §4 and Appendix A).}

These are not model restrictions. They ensure representation-neutral,
sound attribution of learned bits (q\_v) and distinct futures (Alt\_v),
so that SCL/min-cut accounting is meaningful across paradigms. They do
not constrain algorithm design (including efficient P algorithms):
caching, reordering, look-ahead, compression, or hybrid strategies are
all compatible, provided reuse is seed-consistent (Keyedness/Enc
parseability) and learning is charged at first valid use (RWA). In this
light, SCL "explains" polytime runs on L* (e.g., verification with
witness: residual kept at 0 by having all bits resolved) rather than
forbidding optimizations.

\textbf{Framework axioms A1-A5:} Structure L* to create obligations and
enable consistent measurement:

\begin{itemize}
\tightlist
\item
  \textbf{A1 (Hermeticity):} No hidden information channels (disjoint
  designated address pools)
\item
  \textbf{A2 (Injectivity):} Distinct parent tuples \(\to\) distinct
  seeds (Enc injective, parseable)
\item
  \textbf{A3 (Emergence):} Fresh R\_v bits at each node v (rank(H\_v) =
  R\_v; realizability)
\item
  \textbf{A4 (Closure):} Seeds deterministically recover ancestors (Enc
  parseable; verifier reconstructs)
\item
  \textbf{A5 (Dependency):} Parents complete before children (Seed\_v =
  Enc(v \textbar{}\textbar{} parents \textbar{}\textbar{} digest))
\end{itemize}

\textbf{Cut Composition (H1-H5 Properties).} For cut-level bounds,
multiplicativity follows from \textbf{per-run properties H1-H5}
(disjoint address pools, unique-neighbor, Enc injectivity,
realizability, no cross-coupling ensured by FG/horizon; full definitions
and verification in §7.2.1). See Theorem J.1-PROD, Lemma J.1-Cart
(Appendix J) for Cartesian factoring: \(\Pi\)\_C(\(\pi\)) \(\cong\)
\(\prod\)\_\{v\(\in\)C\} S\_v.

\textbf{Completeness (rank forcing; used alongside A1-A5).} In our
overlay, each node v uses a full-rank map with rank(H\_v)=R\_v (see
§6.2.5), ensuring that y\_v fully determines the R\_v emergent bits at
v. We refer to this instance property as "Completeness" and cite it
together with A1-A5 when required (e.g., Appendix J).

\textbf{Lemma 4.A (Instance Compliance).} The constructed language L*
satisfies axioms \textbf{A1-A5}. \emph{Proof (sketch; see Appendix A for
full proof):}

\begin{itemize}
\tightlist
\item
  \textbf{A1 (Hermeticity)}: Achieved through designated disjoint
  address pools U\_v (§6.2.3)
\item
  \textbf{A2 (Injectivity)}: Enc function is injective and
  length-delimited (§6.2.7)
\item
  \textbf{A3 (Emergence)}: Rank-R\_v check matrix H\_v ensures R\_v
  fresh bits; all 2\^{}(R\_v) values realizable (§6.2.5, Lemma 6.1)
\item
  \textbf{A4 (Closure)}: Seeds deterministically recover ancestors via
  Enc parseability (§6.2.7)
\item
  \textbf{A5 (Dependency)}: Seed chaining Seed\_v = Enc(v
  \textbar{}\textbar{} parents \textbar{}\textbar{} digest) enforces DAG
  structure (§6.2.7)
\end{itemize}

For cut composition, \textbf{H1-H5 properties} (disjoint pools,
unique-neighbor, Enc injectivity, realizability, no cross-coupling)
enable Cartesian factoring across cuts (Appendix J).

The SCL (Theorem 7.A) is then proved FROM these properties, not assumed.
\ensuremath{\square}

\textbf{Keyedness (seed-consistency).} Artifacts at node v are correct
iff they are seed-consistent: from the artifact\textquotesingle{}s
content, a verifier can deterministically derive exactly the designated
addresses F\_overlay(Seed\_v; j,\(\ell\)) for all (j,\(\ell\)) in the
artifact\textquotesingle{}s scope on the current seed chain. Including
Seed\_v (and GateDigest\_v when GREQ\_v=1) is sufficient but not
necessary; any injective encoding that deterministically yields the same
addresses qualifies. This is a correctness condition on reuse - not a
storage mandate: the artifact may carry any sufficient statistic that
yields the same designated addresses. All lower bounds count
representation-invariant equivalence classes under this seed-consistency
relation.

\textbf{Definition (Seed-consistent equivalence and Alt\_v).} Two worlds
\(\omega\), \(\omega\)' at node v are equivalent, written \(\omega\)
\(\approx\)\_v \(\omega\)', if along the current seed chain through v
they induce identical designated addresses and outcomes at all
descendants (i.e., future-distinguishable behavior is identical). Define
Alt\_v := \textbar{}\(\Pi\)\_v / \(\approx\)\_v\textbar. Under Keyedness
+ Injectivity, correct solvers must realize at least Alt\_v disjoint
computational states within a try (no cross-seed merges).

\textbf{RWA (Receiving-Window Attribution).} A bit of information is
charged to node v (in q\_v) at the first valid use that functionally
determines it on the current seed chain. Re-reads or later uses of
already revealed information are not re-charged. Combined with
Hermeticity, RWA ensures each counted bit corresponds to at least one
designated read of a cell in the appropriate address pool. See Theorem
I.1 (Ledger Soundness \& Schedule-Invariance) for the schedule-invariant
equivalence of semantic and RWA-credited q and the per-step B-bits
inflow pricing. For the complete operational definition of "first valid
use," see the 4-step procedure in §1.6 (Operational Definition
subsection).

\textbf{Lemma 4.2.I (Overlay independence; no side channels).} Let O
denote the published overlay tuple (G, Sel\_v, H\_v, Enc schema,
F\_overlay, GREQ, PathOf, S(P), salts, public parameters, \(\Phi\)\textasciitilde{}),
where \(\Phi\)\textasciitilde{} is a seed-locked decode schema for \(\varphi\). For any
run on input x* = O:

\begin{itemize}
\tightlist
\item
  O is fixed before the run and is a deterministic function of
  \(\varphi\) and public parameters (see §10 reduction f: \(\varphi\)
  \(\to\) x*).
\item
  O does not depend on any unresolved designated symbols \(\sigma\)\_u
  nor on any run-dependent outputs y\_v.
\item
  Consequently, conditioned on any transcript prefix, unrevealed
  designated symbols carry no additional information through O; fresh
  information can enter only via first-use designated reads (Hermeticity
  + RWA).
\end{itemize}

\emph{Proof (sketch).} By §10, the reduction constructs x*
deterministically from \(\varphi\); all overlay components are explicit
constants in the input. During a run, only computed values (Seeds,
GateDigest) change and are not part of O. Address functions and paths
are public and seed-dependent but do not encode hidden bits. Hence no
published field leaks information about unresolved designated symbols;
per Lemma 5.5.1/RWA, fresh bits are charged only at first-use reads.
\(\blacksquare\)

\textbf{Notation (\(\Pi\)\_v, \(\Phi\)\_v and Alt\_v).} \(\Pi\)\_v
denotes the set of feasible worlds (assignments to unresolved
coordinates consistent with the transcript up to node v). For
counting-type paradigms we write \(\Phi\)\_v := log\(_2\) Alt\_v, where
\texttt{Alt\_v} is the number of simultaneously distinguishable
artifacts (branches/keys/nodes/subfunctions) that differ on unresolved
bits at node v within a single try. For degree-type paradigms we use the
native degree/width measure \(\Phi\)\_v directly; when a standard
degree/width\(\to\)size bridge applies, this yields size lower bounds.

\textbf{Representation invariance.} Alt\_v counts equivalence classes
modulo internal encoding (compression/coordinate changes); see §4.2.x
(Representation Invariance) below.

\textbf{Formalization note (Alt\_v lower bound).} Under A2 (injectivity)
and Keyedness, two worlds that differ on unresolved coordinates at v
induce different seeds and cannot share an artifact without error. Hence
Alt\_v \(\geq\) 2\^{}(R\_v-q\_v) within a try, and with \(\Phi\)\_v :=
log\(_2\) Alt\_v we obtain SCL from SMP.

\subparagraph{4.2.x Representation
Invariance}\label{42x-representation-invariance}

\textbf{Lemma 4.2.2 (Representation-invariant lower bound):} Any
sufficient statistic S must have \textbar{}im(S)\textbar{} \(\geq\)
2\^{}(R\_v-q\_v) at node v with R\_v-q\_v unresolved bits.

\emph{Proof:} Let U index the t := R\_v \(-\) q\_v unresolved
coordinates. For each b \(\in\) \{0,1\}\^{}t, define world \(\omega\)\_b
with x\_v\textbar{}\_U = b. Since rank(H\_v) = R\_v:

\begin{itemize}
\tightlist
\item
  Different b \(\to\) different y\_v = H\_v x\_v (injectivity of H\_v)
\item
  Different y\_v \(\to\) different Seed\_c for children (injectivity of
  Enc)
\item
  Different seeds \(\to\) different designated addresses F\_overlay
\item
  Therefore \(\omega\)\_b and \(\omega\)\_\{b\textquotesingle\} require
  distinguishable artifacts in S
\end{itemize}

Hence \textbar{}im(S)\textbar{} \(\geq\) 2\^{}t. \(\blacksquare\)

\textbf{Key implications:}

\begin{itemize}
\tightlist
\item
  Compression/encoding changes do not reduce distinguishable classes
  (bound counts equivalence classes)
\item
  Probabilistic beliefs do not help worst-case (both possibilities must
  be tracked until resolved)
\item
  Only logical exclusion (achieving probability 0 or 1) reduces states
\end{itemize}

\textbf{Paradigm manifestations} of the conservation law:

\begin{itemize}
\tightlist
\item
  \textbf{Backtracking}: \(\Phi\)\_v = log\(_2\)(tree branches)
\item
  \textbf{Dynamic Programming}: \(\Phi\)\_v = log\(_2\)(distinct keys)
\item
  **Resolution/CDCL**: \(\Phi\)\_v = log\(_2\)(clause width)
\item
  \textbf{OBDD}: \(\Phi\)\_v = log\(_2\)(diagram nodes)
\end{itemize}

With the framework established, we now specify how to instantiate it
with concrete parameters.

\textbf{§4.2 Summary:} Established DAG-Semantic Framework: axioms A1-A5
create obligations (Hermeticity, Injectivity, Emergence, Closure,
Dependency); RWA tracks first-use resolution; SMP shows multiplicative
growth (\textbar{}\(\Pi\)\_\{after\}\textbar{} \(\geq\)
\textbar{}\(\Pi\)\_\{before\}\textbar{} × 2\^{}(R\_v-q\_v)); Keyedness
forbids seed-inconsistent merging. Proved in §7.2.1 (SCL Theorem); cut
composition in Appendix J (Cartesian factoring); used throughout §§6-10.

\paragraph{4.3 Parameters and
Profiles}\label{43-parameters-and-profiles}

\textbf{Primary size parameter.} The total emergence mass is n\_tot :=
\(\Sigma\)\_\{v\(\in\)V\} R\_v (in bits), where V denotes the set of DAG
nodes, and n\_tot = \(\Theta\)(n\_core · log n\_core) under the flat
profile.

Node count \textbar{}V\textbar{}, layer depth L, and per-layer widths
W\_\(\ell\) (nodes at layer \(\ell\)) with per-layer ranks r\_\(\ell\)
must satisfy the rank-budget invariant:
\(\Sigma\)\_\{\(\ell\)=1\}\^{}(L) W\_\(\ell\) r\_\(\ell\) =
\(\Theta\)(n\_tot).

\textbf{Min-cut capacity (run-dependent).} For any cut C:
\(\lambda\)(A,x; C) = \(\Sigma\)\_\{v\(\in\)C\}(R\_v - q\_v). The
bottleneck is \(\lambda\)(A,x) = min\_C \(\lambda\)(A,x; C). When
unambiguous we write \(\lambda\) for \(\lambda\)(A,x). We use
\(\lambda\)\_base for the instance-side profile residual (see §4.1
Definition Block).

\subparagraph{Profile Families}\label{profile-families}

A profile (W\_min, r, L) specifies uniform DAG parameters (constant
width W\_min and rank r across layers) satisfying the rank budget
L·W\_min·r = \(\Theta\)(n\_tot), yielding profile residual
\(\lambda\)\_base := W\_min·r.

\textbf{QP-sharp Profile (Our Main Construction):}

\begin{itemize}
\tightlist
\item
  Width: W\_min = \(\Theta\)(log n\_core)
\item
  Rank: r = \(\Theta\)(log n\_core)
\item
  Depth: L = \(\Theta\)(n\_tot/(W\_min·r)) = \(\Theta\)((n\_core · log
  n\_core)/(log² n\_core)) = \(\Theta\)(n\_core/log n\_core) to satisfy
  budget
\item
  Min-cut (profile): \(\lambda\)\_base = W\_min·r = \(\Theta\)(log²
  n\_core)
\item
  Complexity: distinguishable artifacts \(\geq\) 2\^{}(\(\Theta\)(log²
  n\_core)) = n\_core\^{}(\(\Theta\)(log n\_core)) (quasi-polynomial)
\end{itemize}

\textbf{Why QP-sharp:} This profile achieves the tightest bounds - only
polynomial gap between upper and lower bounds while maintaining clean
P/NP separation. The calculation:

\begin{enumerate}
\def\labelenumi{\arabic{enumi}.}
\tightlist
\item
  Rank budget: L·W\_min·r = \(\Theta\)(n\_core/log
  n\_core)·\(\Theta\)(log n\_core)·\(\Theta\)(log n\_core) =
  \(\Theta\)(n\_core · log n\_core) = \(\Theta\)(n\_tot) {[}YES{]}
\item
  Min-cut (profile): \(\lambda\)\_base = W\_min·r = \(\Theta\)(log²
  n\_core)
\item
  Lower bound: 2\^{}(\(\lambda\)\_base) = 2\^{}(\(\Theta\)(log²
  n\_core)) = n\_core\^{}(\(\Theta\)(log n\_core)) (quasi-polynomial)
\end{enumerate}

\textbf{Worked mini-calculation (toy scaling).} Let n\_core=2\^{}k and
choose W\_min=r=k, L=\(\Theta\)(n\_core/(k·k)). Then
\(\lambda\)\_base=W\_min·r=k² and
2\^{}(\(\lambda\)\_base)=2\^{}(k²)=n\_core\^{}(k)=n\_core\^{}(\(\Theta\)(log
n\_core)), matching the QP-sharp bound.

\textbf{Other Key Profiles:}

\begin{itemize}
\tightlist
\item
  \textbf{Exponential} (L = \(\Theta\)(1), W = \(\Theta\)(n\_core), r =
  \(\Theta\)(1)): \(\lambda\)\_base = \(\Theta\)(n\_core) \(\to\)
  2\^{}(\(\Theta\)(n\_core)) distinguishable artifacts
\item
  \textbf{Sub-exponential} (L = \(\Theta\)(\ensuremath{\sqrt{n}}\_core), W =
  \(\Theta\)(\ensuremath{\sqrt{n}}\_core), r = \(\Theta\)(1)): \(\lambda\)\_base =
  \(\Theta\)(\ensuremath{\sqrt{n}}\_core) \(\to\) 2\^{}(\(\Theta\)(\ensuremath{\sqrt{n}}\_core))
  distinguishable artifacts
\item
  \textbf{Sub-critical} (\(\lambda\)\_base = o(log n\_core)): No P/NP
  separation (since 2\^{}(\(\lambda\)\_base) = n\_core\^{}(o(1)) remains
  polynomial)
\end{itemize}

\textbf{Profile Reference (Quick Lookup):}

\textbf{1. QP-sharp Profile (Main Construction)}

\begin{itemize}
\tightlist
\item
  \(\lambda\)\_base = \(\Theta\)(log² n\_core)
\item
  W\_min = \(\Theta\)(log n\_core)
\item
  r = \(\Theta\)(log n\_core)
\item
  L = \(\Theta\)(n\_core/log n\_core)
\item
  \(\tau\)(n) = \(\Theta\)(log n\_core/log² n\_core) = \(\Theta\)(1/log
  n\_core)
\item
  Time lower bound: n\^{}(\(\Omega\)(log n\_core))
\item
  Distinguishable artifacts: 2\^{}(\(\lambda\)\_base) =
  n\_core\^{}(\(\Theta\)(log n\_core)) (quasi-polynomial)
\end{itemize}

\textbf{2. Flat Profile (Exponential)}

\begin{itemize}
\tightlist
\item
  \(\lambda\)\_base = \(\Theta\)(n\_core)
\item
  W\_min = \(\Theta\)(n\_core)
\item
  r = \(\Theta\)(1)
\item
  L = \(\Theta\)(1)
\item
  \(\tau\)(n) = \(\Theta\)(log n\_core/n\_core) = o(1)
\item
  Time lower bound: 2\^{}(\(\Omega\)(n\_core))
\item
  Distinguishable artifacts: 2\^{}(\(\lambda\)\_base) =
  2\^{}(\(\Theta\)(n\_core)) (exponential)
\end{itemize}

\textbf{3. Sub-exponential Profile}

\begin{itemize}
\tightlist
\item
  \(\lambda\)\_base = \(\Theta\)(\ensuremath{\sqrt{n}}\_core)
\item
  W\_min = \(\Theta\)(\ensuremath{\sqrt{n}}\_core)
\item
  r = \(\Theta\)(1)
\item
  L = \(\Theta\)(\ensuremath{\sqrt{n}}\_core)
\item
  \(\tau\)(n) = \(\Theta\)(log n\_core/\ensuremath{\sqrt{n}}\_core)
\item
  Time lower bound: 2\^{}(\(\Omega\)(\ensuremath{\sqrt{n}}\_core))
\item
  Distinguishable artifacts: 2\^{}(\(\lambda\)\_base) =
  2\^{}(\(\Theta\)(\ensuremath{\sqrt{n}}\_core)) (sub-exponential)
\end{itemize}

\textbf{4. Verification Mode (With Witness)}

\begin{itemize}
\tightlist
\item
  \(\lambda\)\_base = 0 (all bits resolved: q\_v = R\_v)
\item
  Time: poly(n\_core)
\item
  Distinguishable artifacts: 1 (unique witness path)
\end{itemize}

\textbf{Key Formulas:}

\begin{itemize}
\tightlist
\item
  Profile residual: \textbf{\(\lambda\)\_base = W\_min · r}
\item
  FG calibration: \textbf{\(\tau\)(n) = \(\Theta\)(log n /
  \(\lambda\)\_base)} for super-polynomial profiles
\item
  Time bound (single run): \(\geq\) 2\^{}(\(\rho\)-s) ·
  \(\Omega\)(n/W\_min) with \(\rho\) \(\geq\) \(\lambda\)\_base \(-\) s
  and s \(\leq\) \(\Theta\)(\(\tau\) · \(\lambda\)\_base)
\item
  Rank budget constraint: \textbf{L · W\_min · r = \(\Theta\)(n\_tot) =
  \(\Theta\)(n\_core · log n\_core)}
\end{itemize}

\emph{Note: The layer count L must be chosen to satisfy the rank budget,
not fixed arbitrarily.}

\subparagraph{Instance Design vs Algorithmic
Response}\label{instance-design-vs-algorithmic-response}

\textbf{Key Insight:} The profile residual \(\lambda\)\_base is a
structural property of the constructed instance (set by the profile and
wiring), not an algorithmic choice. For any correct run, algorithms can
reduce the residual by resolving bits (increasing q\_v), but
Frontier-Gate (§5.1.2/§6.2.8) limits this compression: the realized
residual satisfies \(\lambda\)(A,x) := min\_C
\(\Sigma\)\_\{v\(\in\)C\}(R\_v\(-\)q\_v) \(\geq\) \(\lambda\)\_base
\(-\) \(\tau\)(n), where \(\tau\)(n) = O(\(\tau\)·\(\lambda\)\_base)
with \(\tau\) = o(1) the gate-horizon calibration parameter (see §6.2.9
and Appendix C.2). Thus FG ensures algorithms cannot compress residual
below \(\Omega\)(\(\lambda\)\_base).

Clarification (seeds vs. lower bound). Theorem 8.A does not assume
hidden or undiscoverable seeds. In our construction the seed chain is
reconstructable from the public overlay, but reconstructing and binding
it to the final chain requires designated reads and digest computations
on seed-bound addresses. The exponential pressure arises from residual
credit accounting and FG's gate-horizon constraint (Appendix C.2), which
together force exponentially many non-accepting rollback segments or
exponentially large frontier --- independent of whether seeds are
discoverable in principle.

\textbf{Multiplicity note.} For any fixed transcript prefix, the
designated bottleneck cut C* admits 2\^{}(cap(C*)) potential seed
evolutions (one per assignment to the unresolved cut bits). Only those
consistent with a satisfying assignment w yield valid witnesses. Search
requires identifying the correct evolution among exponentially many;
given w, the correct chain is efficiently computable (cf. §9.3 Ext, Step
2).

\textbf{The Three-Route Trap (two-dimensional constraint):} At our
constructed instance, algorithms face:

\begin{itemize}
\tightlist
\item
  Cannot reduce \(\lambda\)\_base (instance property)
\item
  Cannot avoid the Semantic Conservation Law (mathematical necessity)
\item
  Must satisfy q + \(\Phi\) \(\geq\) R (a two-dimensional constraint in
  (q, \(\Phi\))) via three operational routes, all exponentially hard:
  Resolution (reading), Elimination (testing), or Storage (maintaining
  2\^{}(R\_v-q\_v) distinguishable artifacts)
\end{itemize}

\textbf{Bounds for QP-sharp:} With \(\lambda\)\_base = \(\Theta\)(log²
n), we achieve tight characterization:

n\^{}\(\Omega\)(log n\_core) \(\leq\) T(n) \(\leq\) n\^{}O(log n\_core)

The gap is only polynomial factors - a rare achievement in complexity
theory.

Having specified the parametric framework, we now define the concrete
language L* and its correctness requirements.

\textbf{§4.3 Summary:} Profile families determine complexity spectrum
via \(\lambda\)\_base = W\_min × r: QP-sharp (\(\lambda\)\_base =
\(\Theta\)(log² n) \(\to\) n\^{}(\(\Theta\)(log n))); Flat
(\(\lambda\)\_base = \(\Theta\)(n) \(\to\) 2\^{}(\(\Theta\)(n)));
Sub-exp (\(\lambda\)\_base = \(\Theta\)(\ensuremath{\sqrt{n}}) \(\to\)
2\^{}(\(\Theta\)(\ensuremath{\sqrt{n}}))); Verification (\(\lambda\)=0 \(\to\) poly). Rank
budget L·W\_min·r = \(\Theta\)(n\_tot) constrains design. FG calibration
\(\tau\)(n) = \(\Theta\)(log n/\(\lambda\)\_base) ensures residual
incompressibility. Used in §8 (per-instance bounds), §6 (instance
construction).

\paragraph{4.4 Language Definition and
Correctness}\label{44-language-definition-and-correctness}

\textbf{Language Definition.} L* overlays a 3-SAT core with
DAG-structured metadata engineered to exhibit SCL properties:

Seed\_v = Enc(v \textbar{}\textbar{}
sort(\{(u,Seed\_u,y\_u):u\(\in\)P(v)\}) \textbar{}\textbar{}
GateDigest\_v) (GateDigest\_v empty when GREQ\_v=0; see §6.2.8 for
GREQ/GateDigest)

where Enc is a canonical, length-delimited, domain-separated encoding
ensuring:

\begin{itemize}
\tightlist
\item
  \textbf{Injectivity}: Different inputs produce different seeds
\item
  \textbf{Parseability}: Can recover (v, parent tuples) from Seed\_v
\item
  \textbf{No cryptography required}: Only needs deterministic encoding
\end{itemize}

\emph{Concrete implementation: Length-delimited concatenation with
sorted parent tuples.}

\textbf{Correctness Principles:}

\begin{itemize}
\item
  \textbf{Hermeticity}: Computations depend only on published metadata
  and seeds (content-addressed discipline)
\item
  \textbf{Keyed reuse (representation-invariant)}: Artifacts are correct
  iff their addressing is \textbf{seed-consistent}: from the
  artifact\textquotesingle{}s content, the solver can compute exactly
  the designated addresses F\_overlay(Seed\_v; j,\(\ell\)) for all
  selector indices (j,\(\ell\)) at node v on the current seed chain.
  This can be achieved by including Seed\_v (and GateDigest\_v when
  GREQ\_v=1) or by any injective encoding that deterministically derives
  the same addresses. The lower bounds are
  \textbf{representation-invariant}: they count equivalence classes of
  artifacts that differ on unresolved bits, independent of how those
  classes are encoded. This is a correctness requirement on reuse, not a
  demand to literally store Seed\_v: any sufficient statistic that
  deterministically yields the same designated addresses qualifies.
\item
  \textbf{Anti-equivalence}: Different seed histories cannot legally
  merge (merging causes wrong designated addresses \(\to\) verifier
  detects error)
\end{itemize}

\emph{Formal proofs in Appendix B.}

\textbf{Key Properties:}

\begin{itemize}
\tightlist
\item
  States with different frontier tuples cannot merge (Anti-equivalence)
\item
  Seeds form a dependency DAG (blockchain-inspired but hashless - using
  content-addressed encoding instead)
\item
  Once seeded, parent states are immutable (No-Rebase)
\item
  Sink seeds contain parseable history of all ancestors (Closure)
\end{itemize}

\emph{Operational contract and real-system mappings in Appendix B.}

\textbf{Acceptance predicate (explicit; base CNF).} To verify x* \(\in\)
L*, the verifier receives instance x* and witness candidate w, where x*
publishes overlay objects for each node v (H\_v, Sel\_v, U\_v salts,
DAG, GREQ\_v flags, and path assignments), the Enc schema and addressing
function F\_overlay, the published path-gate index sets S(P) (Appendix
A.9/C.1.1), public parameters, and \(\Phi\)\textasciitilde{} (seed-locked decode
schema). The verifier accepts (x*,w) iff:

\begin{itemize}
\tightlist
\item
  Base NP predicate: decode \(\varphi\) :=
  Decode\_from\_seeds(\(\Phi\)\textasciitilde{}; Seed\_chain) and check \(\varphi\)(w)=1
\item
  Overlay consistency: for each node v in topological order,

  \begin{enumerate}
  \def\labelenumi{\arabic{enumi})}
  \tightlist
  \item
    compute designated addresses from selector indices (j,\(\ell\)) at
    node v: Addr\_v = \{F\_overlay(Seed\_v; j,\(\ell\))\} (addressing
    function; see §6.2.7/§2.1.2),
  \item
    read the corresponding salts and compute e\_\{v,j,\(\ell\)\},
    assemble x\_v := Sel\_v e\_v, compute y\_v := H\_v x\_v,
  \item
    if GREQ\_v=1, compute GateDigest\_v as in §6.2.8 using the canonical
    path P(v) and XOR over S(P(v)); else set
    GateDigest\_v:=\(\varepsilon\),
  \item
    recompute Seed\_v := Enc(v \textbar{}\textbar{}
    sort(\{(u,Seed\_u,y\_u)\}\_\{u\(\in\)P(v)\}) \textbar{}\textbar{}
    GateDigest\_v), and verify all published structural equalities.
  \end{enumerate}
\end{itemize}

Canonical witness check (completeness and correctness). A canonical
witness has the form W = (w, G\_\(\tau\), Dig\_\(\tau\)) with
G\_\(\tau\) := \{(P,S(P)) : P \(\in\) \ensuremath{\mathcal{P}}\} equal to the published
canonical gate list and Dig\_\(\tau\){[}P{]} the digest bit for each P.
The verifier additionally checks:

\begin{itemize}
\tightlist
\item
  Coverage: G\_\(\tau\) in the witness matches exactly the published
  \{(P,S(P)) : P \(\in\) \ensuremath{\mathcal{P}}\}.
\item
  Values: For each P \(\in\) \ensuremath{\mathcal{P}}, recompute the parity over S(P) on the
  current seed chain and require Dig\_\(\tau\){[}P{]} to match. This
  enforces engagement of all post-horizon gate paths and forbids
  post-hoc synthesis of Dig\_\(\tau\) from w alone (cf. §C.1.1; §10.4.1
  Step 1).
\end{itemize}

All checks are deterministic and polynomial-time (O(n²)). The key
insight: algorithms must either \textbf{resolve} missing bits (increase
q\_v) or \textbf{track distinguishable artifacts} (Alt\_v) for
unresolved seeds - this is SCL in action.

\textbf{Key conventions:}

\begin{itemize}
\tightlist
\item
  \textbf{Receiving-Window Attribution (RWA):} Information charged at
  first valid use
\item
  \textbf{Keyedness:} Artifacts must match actual Seed\_v
\item
  \textbf{Emergence:} 2\^{}(R\_v) possible x\_v values before resolution
\item
  \textbf{Dependencies:} Node v requires all parents complete
\end{itemize}

\subparagraph{Core Mathematical
Results}\label{core-mathematical-results}

\textbf{Key Theorem (Min-Cut Bound):} For any classical algorithm on L*,
the semantic artifact size satisfies:

artifacts \(\geq\) 2\^{}(\(\lambda\)(A,x)) where \(\lambda\)(A,x) =
min\_\{cuts C\} \(\Sigma\)\_\{v\(\in\)C\}(R\_v - q\_v)

\textbf{Proof Structure:}

\begin{enumerate}
\def\labelenumi{\arabic{enumi}.}
\tightlist
\item
  \textbf{Per-node bound}: At node v, after resolving q\_v bits: Alt\_v
  \(\geq\) 2\^{}(R\_v-q\_v) (by emergence; §4.1)
\item
  \textbf{Cut composition}: For cut C: \(\Sigma\)\_\{v\(\in\)C\}
  log\(_2\)(Alt\_v) \(\geq\) \(\Sigma\)\_\{v\(\in\)C\}(R\_v - q\_v) (by
  H1-H5 Cartesian factoring; §7.2.1)
\item
  \textbf{Anti-equivalence}: States with different frontier tuples
  cannot legally merge (by Enc injectivity)
\item
  \textbf{Min-cut}: Take minimum over all s-t cuts for tightest bound
\end{enumerate}

\textbf{Distributional Extension (for randomized algorithms):} Under
product-salts distribution, facts across cut C are mutually independent
given ancestors. This yields:

\begin{itemize}
\tightlist
\item
  Pr{[}guess correctly without reading{]} \(\leq\)
  \(\prod\)\_\{v\(\in\)C\} 2\^{}(-(R\_v-q\_v)) =
  2\^{}(-\(\lambda\)(A,x))
\item
  Expected complexity: Randomized algorithms face expected
  \(\Omega\)(2\^{}(\(\lambda\)(A,x))) trials
\item
  Per-instance deterministic bounds: Coin-fixing extends these bounds to
  per-instance deterministic time (§8; Theorem 8.A)
\end{itemize}

\begin{center}\rule{0.5\linewidth}{0.5pt}\end{center}

With the DAG-Semantic Framework and language L* fully specified, we now
analyze how specific computational paradigms manifest the Semantic
Conservation Law.

\subsubsection{5. Computational Models and
Classes}\label{5-computational-models-and-classes}

Section 4 established the abstract DAG-Semantic Framework with semantic
quantities R\_v (emergence), q\_v (resolution), \(\Phi\)\_v =
log\(_2\)(alternatives), and \(\lambda\) (min-cut residual). We defined
axioms \textbf{A1-A5} as the abstract properties that imply the Semantic
Conservation Law q + \(\Phi\) \(\geq\) R, and specified our
computational model (deterministic k-tape TMs).

\textbf{Purpose of Section 5:} We now show that the abstract
conservation law has \textbf{universal manifestation} - the requirement
to maintain 2\^{}(R-q) distinguishable artifacts appears across
\emph{all} computational paradigms, expressed in each
paradigm\textquotesingle{}s natural units. This demonstrates that our
framework captures a fundamental structural necessity, not a
model-specific artifact.

\textbf{Analysis Approach.} The Structural OWF construction (§9) uses
instance axioms \textbf{A1-A5} (Hermeticity, Injectivity, Emergence,
Closure, Dependency) with per-run properties \textbf{H1-H5} (disjoint
pools, unique-neighbor, Enc injectivity, realizability, no
cross-coupling) for cut composition. Section 5 focuses on
paradigm-invariant manifestations of the conservation law.

\textbf{Roadmap:}

\begin{itemize}
\tightlist
\item
  \textbf{§5.0}: Universal principle - same bottleneck measured in
  different metrics
\item
  \textbf{§5.1}: Min-cut characterization across DAG dependencies
\item
  \textbf{§5.2}: Algorithm classes and paradigm-specific manifestations
  (Backtracking, DP, OBDD, Resolution, Streaming)
\item
  \textbf{§5.3}: How correctness manifests in each paradigm
\item
  \textbf{§5.4-§5.7}: Extensions (technical bounds, practical
  considerations, parallelism)
\end{itemize}

\textbf{Why this matters for P \(\neq\) NP:} By showing the same
structural requirement appears in every paradigm, we establish that
\(\lambda\) is paradigm-invariant. An adversary cannot escape
exponential cost by switching computational models - the bottleneck
follows from L*\textquotesingle{}s structure, not from algorithmic
limitations in any particular model.

\textbf{Scope labels for §5.} Results marked \textbf{{[}Thm{]}} are
fully proved here or in appendices; \textbf{{[}Adapter{]}} are model
projections from Theorem 7.A via representation-invariance (Lemma 5.3.1
+ Corollary 5.3.2); \textbf{{[}Template{]}} are outlined connections
with complete details in cited appendices or §12.12.

\textbf{Preservation note (semantic vs simulation).} The SCL is proved
directly on runs via semantic counters (RWA-credited q, seed-consistent
Alt, min-cut \(\lambda\)). When we compile to other models (e.g.,
layered BPs), we do not assume the simulation preserves these counters
automatically. Instead, we use the SNF ledger (Appendix I.2) to
reconstruct first-use events and frontier classes from the transcript,
and then relate them to model-native measures (width/subfunctions).
Lemma I.2.4 (monotonicity) ensures compilation cannot reduce
ledger-defined counters; Lemma I.2.5 equates frontier classes with
width/subfunctions at cuts.

\paragraph{5.0 Universal State-Tracking Manifestations in
L*}\label{50-universal-state-tracking-manifestations-in-l}

\textbf{Core principle:} Through L*\textquotesingle{}s engineered
structure, we demonstrate that all paradigms must count the same
fundamental quantity - independent, simultaneously distinguishable
artifacts - expressed in their natural units. At node v with R\_v-q\_v
unresolved bits, algorithms need 2\^{}(R\_v-q\_v) distinguishable
artifacts to avoid collision (merging inequivalent worlds \(\to\) wrong
answers).

\textbf{Note on order-robustness (OBDD-specific).} For OBDD width bounds
(§5.2, §5.3), order-robustness requires \textbf{expander-parity} gates
(§6.2.10; {[}HOO06{]}); with identity gates H\_v=I, bounds are
order-specific. Other paradigms do not have this order-dependency.

\textbf{Formal basis:} By Keyedness + Injectivity (Lemma 7.I), states
with different seeds cannot merge. Across a cut C, requirements
multiply: total states \(\geq\)
2\^{}(\(\Sigma\)\_\{v\(\in\)C\}(R\_v-q\_v)) (Theorem 7.A.1 and Appendix
J (Theorem J.1; Lemma J.1-Cart)).

\textbf{Paradigm manifestations (detailed analysis in §5.2):} When
solving L*, the constraint 2\^{}(R\_v-q\_v) simultaneously
distinguishable artifacts appears as tree branches (backtracking), table
keys (DP), diagram width (OBDD), clause width (Resolution), or
monochromatic rectangles (Communication). These are not separate
phenomena but the same structural necessity for L* expressed in
different paradigm-specific units. Under Keyedness, artifacts with
different seeds cannot merge - merging would compute wrong addresses for
at least one seed history \(\to\) incorrect answers.

\paragraph{5.1 DAG dependency structure and min-cut
bounds}\label{51-dag-dependency-structure-and-min-cut-bounds}

Dependency requires that node v be settled (y\_v fixed) before any child
can be processed. For any correct algorithm, maintained distinguishable
artifacts grow exponentially with residual capacity \(\lambda\)(A,x) =
min\_C \(\Sigma\)\_\{v\(\in\)C\}(R\_v - q\_v).

\subparagraph{5.1.1 Min-cut
characterization}\label{511-min-cut-characterization}

\textbf{Key Insight:} In parallel/layered DAGs, \textbf{\(\lambda\)(A,x)
(min-cut capacity) - not path sums (\(\Sigma\)\_\{v\(\in\)P\}(R\_v -
q\_v) for root\(\to\)leaf chains) - controls the exponent}. On chains
they coincide, but for parallel structures min-cut gives tighter bounds
by capturing the true bottleneck.

Mini-example (two parallel branches). Let A and B be parallel parents of
C with R\_A\(-\)q\_A = R\_B\(-\)q\_B = 1. Path sums along A\(\to\)C or
B\(\to\)C each give 1, but the min-cut across \{A,B\} has residual 2,
forcing 2² simultaneously distinguishable artifacts at the cut.

\textbf{Lemma (Closure/Recoverability; see Appendix L).} From the
sink-reachable seeds recorded in the transcript, the verifier
deterministically reconstructs all ancestor tuples (Seed, y).
\emph{(Used to justify the cut bound; see also Appendix J for
compositional details.)}

Let G=(V,E) be a DAG with per-node ranks R\_v and resolutions q\_v. For
v with parent set P(v), define a \textbf{canonical, content-addressed}
seed

Seed\_v = Enc(v \textbar{}\textbar{} sort(\{(u, Seed\_u, y\_u) : u
\(\in\) P(v)\}) \textbar{}\textbar{} GateDigest\_v), where GateDigest\_v
is empty when GREQ\_v=0 (acyclic: digest ranges over pre-horizon indices
only; see §6.2.8),

where all fields are length-delimited and parents are lexicographically
sorted. \textbf{Keyedness}: artifacts at v must be seed-consistent - any
representation that deterministically yields the same designated
addresses F\_overlay(Seed\_v; j,\(\ell\)) qualifies
(representation-invariant; §4.2).

\textbf{Address pools.} Each node v has a designated address pool U\_v,
and \{U\_v\}\emph{v are pairwise disjoint. Designated cells at v are
computed as u}\{v,j,\(\ell\)\} := F\_overlay(Seed\_v; j,\(\ell\))
\(\in\) U\_v.

\textbf{Parent dependencies.} Since Seed\_v = Enc(v \textbar{}\textbar{}
sort(\{(u,Seed\_u,y\_u)\}\_\{u\(\in\)P(v)\}) \textbar{}\textbar{}
GateDigest\_v), addresses and artifacts at v are undefined until every
parent u fixes y\_u and (if GREQ\_v=1) the gate digest is computed on
the current seed chain. Any attempt to proceed with an incomplete parent
set necessarily maintains a frontier of distinguishable artifacts keyed
by multiple candidate seeds, counted in Alt\_v.

\textbf{Multiplicative composition.} Under H1-H5, across any s-t cut C,
distinguishable classes factor as Alt(C) = \(\prod\)\_\{v\(\in\)C\}
Alt\_v by Dependency (seed chaining), Keyedness, and Injectivity; see
Theorem J.1 (Appendix J) for the formal cut-product proof.

\textbf{Lemma 5.1.NC (No cross-coupling across a cut - main text).}
Under A1 (Hermeticity), A2 (Injective, parseable Enc), disjoint
designated pools \{U\_v\}\_v, and FG acyclicity (GateDigest\_v depends
only on pre-horizon indices; §6.2.8), fix any transcript prefix \(\pi\)
and any s-t cut C. Let S\_v(\(\pi\)) denote the set of admissible local
unresolved coordinates at node v consistent with \(\pi\). Then feasible
worlds across C factor as a Cartesian product:

\(\Pi\)\_C(\(\pi\)) \(\cong\) \(\prod\)\_\{v\(\in\)C\} S\_v(\(\pi\)).

In particular, Alt(C) =
\textbar{}\(\Pi\)\_C(\(\pi\))/\(\equiv\)\textbar{} \(\geq\)
\(\prod\)\emph{\{v\(\in\)C\} 2\^{}(R\_v-q\_v) =
2\^{}(\(\Sigma\)}\{v\(\in\)C\}(R\_v-q\_v)).

\emph{Proof sketch.} Hermeticity and disjoint pools ensure that local
designated addresses witnessing node-v constraints always lie in U\_v
and depend only on Seed\_v. By Dependency and Enc\textquotesingle{}s
parseable injectivity, Seed\_v is a function of ancestor (Seed,y) tuples
and (for GREQ\_v=1) GateDigest\_v, which in turn aggregates only
pre-horizon data; thus post-horizon digests do not introduce relations
among unresolved coordinates on C. Therefore no transcript-consistent
constraint simultaneously restricts two distinct nodes v\(\neq\)v' in C.
Choices at distinct cut nodes are independent under \(\pi\),
establishing the Cartesian factoring. See Appendix J (Theorem J.1, Lemma
J.1-Cart, Lemma J.1-NC) for the full proof. \ensuremath{\square}

\emph{(Interpretation note: Cartesian factoring applies to
overlay-feasible worlds - cut-coordinate tuples that produce well-formed
overlay structures. Whether resulting instances decode to satisfiable
\(\varphi\) is a separate filter that does not retroactively couple the
cut coordinates. See interpretation note following Lemma J.1-REAL in
Appendix J for detailed discussion.)}

\textbf{Example:} Diamond DAG with a,b \(\to\) c \(\to\) d where
(R\_a,q\_a)=(2,1), (R\_b,q\_b)=(3,1), (R\_c,q\_c)=(2,0),
(R\_d,q\_d)=(1,0). Min-cut analysis: Cut \{a,b\} has capacity 3, Cut
\{c\} has capacity 2, Cut \{d\} has capacity 1. Thus \(\lambda\)(A,x) =
1, requiring \(\geq\) 2\^{}1 = 2 distinguishable artifacts at the
bottleneck.

\textbf{Theorem (Min-Cut Characterization).} \emph{Framework refs:}
Axioms A2 (injective CA-Enc), A3 (Emergence), A4 (Closure), Hermeticity,
Keyedness, RWA; see \textbf{Theorem J.1} (Appendix J). We extend the
max-path bound to a \textbf{min-cut characterization}. With node
capacities c(v) = R\_v - q\_v, for any s-t cut C in the node-DAG:

\(\Sigma\)\_\{v\(\in\)C\} log\(_2\)(Alt\_v) \(\geq\)
\(\Sigma\)\_\{v\(\in\)C\} c(v)

Equivalently, defining the total across the cut as Alt(C) :=
\(\Pi\)\_\{v\(\in\)C\} Alt\_v, we have log\(_2\)(Alt(C)) \(\geq\)
\(\lambda\), where \(\lambda\) = min\_C \(\Sigma\)\_\{v\(\in\)C\} c(v)
is the min-cut value. The solver can reduce this only by paying
corresponding resolution q\_v at cut nodes. \textbf{See Appendix J for
the complete proof.}

\textbf{Prerequisites (A1-A5 framework).} The min-cut characterization
relies on the structural properties created by \textbf{A1-A5}: the
instance-side invariants (Emergence, Completeness), the canonical
content-addressed dependency encoder Enc that stores sorted
(u,Seed\_u,y\_u) tuples, Hermeticity (purity), Keyedness (correct reuse
iff keys include Seed\_v), and Receiving-Window Attribution (RWA). For
cut composition, we use \textbf{H1-H5} (disjoint pools, unique-neighbor,
Enc injectivity, realizability, no cross-coupling). Appendix J proves
closure/recoverability and Cartesian product factoring across cuts under
these properties.

\textbf{Lemma (Cut Cartesian Product).} See Appendix J, Theorem J.1-PROD
(Cut Product Theorem - minimal hypotheses) for the formal statement and
proof that \(\Pi\)\_C(\(\pi\)) \(\cong\) \(\prod\)\emph{\{v\(\in\)C\}
S\_v and \textbar{}\(\Pi\)\_C(\(\pi\))\textbar{} =
\(\prod\)}\{v\(\in\)C\} \textbar{}S\_v\textbar{}.

\textbf{Corollary 5.1.1-DT (Decision-tree lower bound; per-instance).}
For any deterministic decision tree that decides L* correctly on a given
instance x*, the number of queries satisfies D(L*, x*) \(\geq\)
\(\lambda\)(x*) = min\_C \(\Sigma\)\_\{v\(\in\)C\}(R\_v - q\_v).

\emph{Proof (Per-Instance).} By \textbf{A4 (Closure)} (Appendix J), sink
seeds uniquely determine all ancestor (Seed,y) along a
run\textquotesingle{}s transcript, so choices at cut nodes are decoded
from the sink. By \textbf{H4 (realizability) (Lemma 6.1), all
2\^{}(R\_v-q\_v) values of x\_v remain structurally valid at node v
before discovery reads. Together with H1 (disjoint designated pools)}
(Lemma 6.5), feasible choices across cut nodes C factor as a Cartesian
product via \textbf{H1-H5} (per-run properties; Appendix J), giving
\textbar{}\(\Pi\)\_C\textbar{} = \(\prod\)\emph{\{v\(\in\)C\}
2\^{}(R\_v-q\_v) = 2\^{}(\(\Sigma\)}\{v\(\in\)C\}(R\_v-q\_v)). Any
correct decision tree must distinguish these worlds, requiring at least
\(\Sigma\)\_\{v\(\in\)C\}(R\_v\(-\)q\_v) queries. Minimizing over cuts
gives the bound. \ensuremath{\square}

\textbf{Consequence for Randomized Algorithms.} For randomized decision
trees, \textbf{coin-fixing} (Yao\textquotesingle{}s principle) extends
the bound: any fixed coin string yields a deterministic tree subject to
the same per-instance bound. This applies to \textbf{every FG-wired
instance} (Theorem 8.A).

\textbf{Why min-cut?} Min-cut and max-path are incomparable in general
DAGs: min-cut excels with parallelism while max-path excels on long
chains (cf. the micro-example above). We use both: product bounds along
chains (Theorem 4.2) and min-cut bounds (Appendix J) for parallel
structure. This complements rather than strengthens the chain
characterization.

\textbf{Zero-\(\lambda\) sanity check.} If along some valid s-t cut C we
have q\_v=R\_v for all v\(\in\)C, then
\(\Sigma\)\_\{v\(\in\)C\}(R\_v-q\_v)=0 and the calculus certifies
\textbf{no forced distinguishable artifacts} past that cut. Intuitively,
verification "keeps up" with emergence, so no multiplicative uncertainty
survives - exactly what we expect for "easy" instances/languages.

With the min-cut characterization established, we now examine how
specific algorithmic paradigms manifest these bounds. (Technical
extensions such as machine-specific bounds and practical cut selection
guidance appear in §5.5-5.6.)

\subparagraph{5.1.2 Advanced: Frontier-Gate (FG) for Tight
Bounds}\label{512-advanced-frontier-gate-fg-for-tight-bounds}

\emph{Note: This subsection presents an advanced technique for achieving
tight bounds. Readers primarily interested in the paradigm
manifestations can skip to §5.2.}

Purpose (one line). FG converts artifact lower bounds (Alt) into
single-run time bounds by ensuring each non-accepting segment performs
priced designated work, so Alt-based multiplicative lower bounds yield
concrete step counts.

\textbf{Brief:} The Frontier-Gate (FG) is an instance-side mechanism,
implemented by wiring GateDigest into seeds, that controls \textbf{when}
information can be accumulated, preventing algorithms from
"front-loading" information gathering. This construction choice refines
the arity-bounded bounds from loose to tight (e.g., from
2\^{}\(\Theta\)(n) to n\^{}\(\Theta\)(log n\_core) for QP-sharp profile)
and is part of L*\textquotesingle{}s design.

\textbf{Key result:} FG caps pre-final agreement at s \(\leq\)
\(\tau\)·\(\rho\) where \(\rho\) is the effective residual (§4.3) and
\(\tau\) = \(\Theta\)(log n/\(\lambda\)\_base), ensuring tight bounds
unconditionally for single-run TMs.

\emph{Unpredictability without reads.} Skipping any designated primitive
in a cut-gate leaves two transcript-consistent completions that flip the
digest parity (Lemma C.1.2). Thus pre-scans cannot predict digests; a
fresh digest must be \emph{computed} on the current seed chain.

\textbf{Lemma 5.1.2-U (No-read unpredictability; restated, model-only).}
Fix a transcript prefix \(\pi\) and a gated node v with current Seed\_v.
Let S(P(v)) be the published index set for the gate on
v\textquotesingle{}s canonical path. If the run has not evaluated at
least one designated primitive e\_\{v,j,\(\ell\)\} in S(P(v)) on the
current seed chain, then there exist two completions of the instance,
both consistent with \(\pi\) and agreeing on all designated payloads
already read, that yield opposite values of GateDigest\_v.

\emph{Proof (combinatorial; no cryptography).} GateDigest\_v is the
parity of the designated terms indexed by S(P(v)) under the current
Seed\_v. By Unique-Neighbor addressing and disjoint pools, each term
refers to a distinct designated payload tied to v. If some term is
unevaluated, flip only that payload bit while holding all others (and
all previously read symbols) fixed. Both completions are consistent with
\(\pi\) by Completeness/realizability, and they yield opposite parity by
sensitivity of XOR. Hence the digest is not determined by \(\pi\);
computing it requires evaluating all designated terms on the current
seed chain. \(\blacksquare\)

\textbf{Critical Clarification - FG is Instance-Side, Preserves NP:}

{[}YES{]} \textbf{Instance-side mechanism}: FG is built into
L*\textquotesingle{}s construction via GateDigest wiring, not a model
restriction {[}YES{]} \textbf{Preserves NP membership}: The NP verifier
recomputes all GateDigest checks in polynomial time {[}YES{]} \textbf{No
model assumptions added}: Only requires Receiving-Window Attribution
(RWA) and keyedness for accounting {[}YES{]}
\textbf{Verifier-checkable}: All FG requirements are deterministic
functions of (x*,w) checkable in poly(n\_core) time

\textbf{Lemma (Per-segment baseline under FG).} In any \textbf{single
run}, once the \(\mu\) and \(\tau\) allowances are exhausted,
\textbf{every non-accepting rollback segment} that progresses the final
chain must either (i) reveal a new cut bit or (ii) evaluate a
\textbf{cut-gate} G\_C(y\_\{\(\leq\)C\}). By \textbf{FG-3 (Gate
requirement)} in Appendix C.1, each such gate evaluation requires
computing a digest over \textbar{}S(P)\textbar{} = \(\Theta\)(n/W\_min)
seed-dependent terms. By Lemma 5.5.1.c (TM digest cost) and Lemma
A.1.\(\Delta\) (address-churn), this requires \(\Omega\)(n/W\_min)
operations on a k-tape TM regardless of pre-scanning. Therefore each
non-accepting rollback segment costs \(\Omega\)(n/W\_min) steps.

\textbf{Key Technical Lemmas (stated here for immediate reference; full
proofs in §5.5.1):}

\textbf{Lemma 5.5.1.c (TM parity computation cost).} Computing
GateDigest\_v = \(\oplus\)\emph{\{(u,(j,\(\ell\))) \(\in\) S(P)\}
\(\sigma\)}\{F\_overlay(Seed\_u; j,\(\ell\))\} requires
\(\Omega\)(\textbar S(P)\textbar) tape operations even with all salts
pre-cached, since parity has full sensitivity and accessing
\textbar{}S(P)\textbar{} distinct cells requires
\(\Omega\)(\textbar S(P)\textbar) head moves.

\textbf{Lemma A.1.\(\Delta\) (Address churn).} When the seed chain
changes, \(\Theta\)(\textbar S(P)\textbar) designated addresses relocate
entirely, forcing fresh parity computation even with cached salts.
Formally, for two seeds Seed\_v \(\neq\) Seed'\_v and the seed-derived
permutations \(\pi\)\_v\^{}(Seed\_v), \(\pi\)\_v\^{}(Seed'\_v), there
exists a universal constant c \(\in\) (0,1) such that
\textbar{}\(\pi\)\_v\^{}(Seed\_v)(S(P)) \(\Delta\)
\(\pi\)\_v\^{}(Seed'\_v)(S(P))\textbar{} \(\geq\)
c·\textbar{}S(P)\textbar{} (see Lemma A.1.F and §A.1.1).

For complete technical details including the \(\mu\)-\(\tau\)
parameterization, two-tier gate mechanism, and full proofs, see
\textbf{Appendix C.1}.

See §8.1 (Theorem 8.A) for the consolidated single-run per-instance
bound and its proof using segment counting (Appendix C.2) and the
per-segment baseline under FG.

Soundness of FG time pricing. The gate digest GateDigest\_v =
\(\oplus\)\emph{\{(u,(j,\(\ell\))) \(\in\) S(P)\}
\(\sigma\)}\{F\_overlay(Seed\_u; j,\(\ell\))\} requires computing XOR
over \textbar{}S(P)\textbar{} = \(\Theta\)(n/W\_min) seed-dependent
addresses. Even with all salts pre-cached, by Lemma 5.5.1.c and Lemma
A.1.\(\Delta\) this computation costs \(\Omega\)(n/W\_min) TM
operations; pre-scanning does not reduce the work needed on a new seed
chain.

\subparagraph{5.1.3 Why Min-Cut?}\label{513-why-min-cut}

\textbf{Key question:} What measure captures the computational
bottleneck in a DAG? We prove it must be min-cut capacity
\(\lambda\)(A,x) = min\_C \(\Sigma\)\_\{v\(\in\)C\}(R\_v - q\_v).

\begin{quote}
\textbf{Terminology:} Throughout this paper, "min-cut" refers to
\textbf{vertex-capacitated s-t min-cut} (also called \textbf{minimum s-t
vertex cut} or \textbf{s-t vertex separator}) - a vertex cut C
\(\subseteq\) V hitting every source\(\to\)sink path; for minimum cuts
in DAGs, an antichain representative exists. Node capacities are c(v) =
R\_v - q\_v. Distinct from edge-weighted min-cut.
\end{quote}

\textbf{Relationship to Classical Min-Cut Theory}

Our use of min-cut builds upon well-established graph-theoretic
foundations while introducing a semantic interpretation for complexity
analysis. Understanding this relationship clarifies what is standard
mathematical machinery versus what constitutes our conceptual
contribution.

\emph{Standard graph-theoretic formulation.} In network flow theory, a
\textbf{vertex-capacitated min-cut} is computed as follows: Given a
directed acyclic graph G = (V,E) with designated source s and sink t, a
\textbf{vertex cut} C \(\subseteq\) V is a set of vertices whose removal
disconnects all paths from s to t. In DAGs, any minimum s-t vertex cut
admits an \textbf{antichain} representative hitting every s\(\to\)t
path. The capacity of such a cut is the sum of individual vertex
capacities: cap(C) = \(\Sigma\)\_\{v\(\in\)C\} c(v). The minimum vertex
cut is found by:

\begin{enumerate}
\def\labelenumi{\arabic{enumi}.}
\tightlist
\item
  \textbf{Node-splitting transformation}: Replace each vertex v with
  v\_in and v\_out, connected by an edge of capacity c(v)
\item
  \textbf{Standard max-flow algorithm}: Apply Ford-Fulkerson,
  Dinic\textquotesingle{}s algorithm, or push-relabel on the transformed
  graph
\item
  \textbf{Min-cut recovery}: By max-flow/min-cut theorem, the min-cut
  value equals max-flow
\end{enumerate}

This formulation has been standard since the 1950s (Ford \& Fulkerson)
and is covered in graph algorithms textbooks.

\emph{Our semantic interpretation.} The innovation lies not in the
min-cut algorithm itself, but in recognizing what to measure:

\textbf{Traditional network flow:} Vertex capacities c(v) represent
physical constraints - bandwidth limits, supply capacities, transmission
costs. The min-cut value answers: "What is the maximum throughput
possible?" This is a \textbf{fixed property} of the network topology.

\textbf{Our semantic framework:} Vertex capacities are defined as
\textbf{unresolved semantic obligations}: c(v) = R\_v - q\_v, where:

\begin{itemize}
\tightlist
\item
  R\_v: Bits required to emerge at node v (determined by instance
  structure, Axiom A3)
\item
  q\_v: Bits algorithm A has resolved at v via designated reads
  (run-dependent, credited by RWA)
\end{itemize}

The min-cut value \(\lambda\)(A,x) answers: "What is the irreducible
information bottleneck this algorithm faces on this instance?" This is a
\textbf{run-dependent measure} - different algorithms resolving
different information yield different min-cut values on the same
instance.

\emph{Key conceptual shift.} Where classical applications interpret
min-cut as a capacity limitation, we interpret it as a
\textbf{complexity predictor}. The mathematical structure is identical;
the semantic content is fundamentally different. This reinterpretation
enables \(\lambda\) to serve as a paradigm-invariant complexity measure:
the same min-cut value translates directly to tree size (backtracking),
table size (DP), proof width (resolution), diagram width (OBDDs), and
time (k-tape TMs).

\textbf{Analogy:} Just as Shannon recognized thermodynamic
entropy\textquotesingle{}s mathematical form applies to information
transmission (though measuring different quantities - heat vs
uncertainty), we recognize min-cut\textquotesingle{}s mathematical form
applies to computational complexity (though measuring different
quantities - flow capacity vs semantic obligation).

\begin{center}\rule{0.5\linewidth}{0.5pt}\end{center}

\textbf{Formal Characterization of Min-Cut as Canonical Measure}

Having established the conceptual connection between min-cut and
complexity, we now prove that this connection is mathematically
necessary - not merely a useful analogy, but the unique optimal
aggregator for DAG-structured semantic constraints. We characterize what
properties any reasonable complexity measure must satisfy, then show
min-cut is the unique function meeting these requirements.

\textbf{Definition 5.1.3.1 (Admissible Semantic Aggregator).} A function
A: (DAG, \{R\_v\}, \{q\_v\}) \(\to\) \(\mathbb{R}\)\(\geq\)0 is
\emph{admissible} if it satisfies:

\begin{enumerate}
\def\labelenumi{\arabic{enumi}.}
\item
  \textbf{Additivity (ADD):} On any cut C where the instance ensures
  node independence (no hidden correlations among unresolved
  coordinates; formally via disjoint designated pools and independence
  axioms), A aggregates residuals by \(\Sigma\)\_\{v\(\in\)C\}(R\_v -
  q\_v)
\item
  \textbf{Monotonicity (MON):} If q\textquotesingle{}\_v \(\geq\) q\_v
  for all v, then A(\{q\textquotesingle{}\_v\}) \(\leq\) A(\{q\_v\})
  (learning more never increases the bound)
\item
  \textbf{Representation-invariance (REP):} A depends only on semantic
  equivalence classes (not syntactic encodings or variable orderings)
\item
  \textbf{Cut-soundness (CUT):} Since any s\(\to\)t computation must
  traverse some cut, and algorithms can choose which cuts to make
  expensive (within DAG constraints), A must identify the bottleneck:
  the minimum unavoidable residual over all s\(\to\)t cuts. (This
  captures that the bound must be both necessary - some cut is hit - and
  tight - achieved at the minimum-cost cut.)
\item
  \textbf{Sharpness (TIGHT):} There exist calibrated instances/solvers
  achieving complexity 2\^{}(A+o(A)) (bound is asymptotically tight;
  rules out overly pessimistic measures)
\item
  \textbf{Normalization (NORM):} When R\_v = q\_v for all v (full
  resolution), A = 0 (no remaining uncertainty)
\end{enumerate}

\textbf{Lemma 5.1.3.2 (Cut Lower Bound).} For any s-t cut C in DAG G,
any correct algorithm for L* must maintain at least
2\^{}(\(\Sigma\)\_\{v\(\in\)C\}(R\_v-q\_v)) distinguishable artifacts
when passing through C.

\emph{Proof (Per-Instance via H1-H5).} By \textbf{disjoint designated
atoms} (H1, Lemma 6.5) + \textbf{realizability} (H4, Lemma 6.1),
feasible worlds across C factor as a \textbf{Cartesian product (Theorem
J.1, Lemma J.1-Cart). By SCL, each node v \(\in\) C requires Alt\_v
\(\geq\) 2\^{}(R\_v - q\_v) local artifact classes. By Keyedness} (H2),
artifacts keyed by different seeds cannot merge. Therefore, the total
across C is \(\prod\)\emph{\{v\(\in\)C\} Alt\_v \(\geq\)
\(\prod\)}\{v\(\in\)C\} 2\^{}(R\_v - q\_v) =
2\^{}(\(\Sigma\)\_\{v\(\in\)C\}(R\_v-q\_v)). \ensuremath{\square}

\textbf{Proposition 5.1.3.3 (Min-Cut is canonical and tight up to
subpolynomial factors).} The min-cut capacity \(\lambda\)(A,x) =
min\_\{C\} \(\Sigma\)\_\{v\(\in\)C\}(R\_v \(-\) q\_v) is the canonical
admissible semantic aggregator for our framework, and for our calibrated
families it is tight up to subpolynomial factors. (Per-node SCL is
Theorem 7.A.)

\emph{Proof.} We establish that \(\lambda\)(A,x) is admissible. For
tightness, Theorem C.2.T gives calibrated instances/solvers with
complexity 2\^{}(\(\lambda\)\_base+o(\(\lambda\)\_base)) (QP-sharp:
n\^{}(\(\Theta\)(log n\_core))). Any admissible aggregator
A\textquotesingle{} larger than \(\lambda\)(A,x) would contradict these
achievers, hence A\textquotesingle{} \(\leq\) \(\lambda\)(A,x) +
o(\(\lambda\)\_base). \(\blacksquare\)

\textbf{Part 1 (Admissibility):}

\begin{itemize}
\tightlist
\item
  ADD: \(\lambda\)(A,x) uses sum \(\Sigma\)\_\{v\(\in\)C\}(R\_v - q\_v)
  for each cut when instance ensures independence (A6/H1-H5) {[}YES{]}
\item
  MON: Increasing q\_v decreases R\_v - q\_v, hence decreases
  \(\lambda\)(A,x) {[}YES{]}
\item
  REP: Depends only on (R\_v - q\_v), not syntactic representations or
  orderings {[}YES{]}
\item
  CUT: Takes minimum over all s-t cuts (identifies bottleneck) {[}YES{]}
\item
  TIGHT: With FG (§5.1.2, Appendix C.2), algorithms achieve
  2\^{}(\(\lambda\)\_base+o(\(\lambda\)\_base)) {[}YES{]}
\item
  NORM: When R\_v = q\_v for all v, \(\lambda\)(A,x) = min\_C
  \(\Sigma\)\_v(0) = 0 {[}YES{]}
\end{itemize}

\textbf{Part 2 (Tightness):} Let A\textquotesingle{} be any admissible
aggregator.

\textbf{Necessity.} For any s\(\to\)t cut C, SCL + Cut Cartesian Product
(Lemma J.1-Cart) imply the run must maintain Alt(C) \(\geq\)
2\^{}(\(\Sigma\)\_\{v\(\in\)C\}(R\_v-q\_v)). Since some cut must be hit
by every source\(\to\)sink path, at the \textbf{bottleneck} (a min-cut
C*) we have Alt(C*) \(\geq\) 2\^{}(min\_C
\(\Sigma\)\_\{v\(\in\)C\}(R\_v-q\_v)) = 2\^{}(\(\lambda\)(A,x)). Hence
\(\lambda\)(A,x) lower-bounds any admissible aggregator:
A\textquotesingle{} \(\geq\) \(\lambda\)(A,x).

\textbf{Upper tightness.}

\begin{itemize}
\tightlist
\item
  If A\textquotesingle{} \textless{} \(\lambda\)(A,x),
  \textbf{Necessity} above is violated on the min-cut C*.
\item
  If A\textquotesingle{} \textgreater{} \(\lambda\)(A,x),
  \textbf{Sharpness (TIGHT)} plus \textbf{Theorem C.2.T} (FG tightness)
  give calibrated instances/solvers with complexity
  2\^{}(\(\lambda\)\_base+o(\(\lambda\)\_base)). Then
  2\^{}(A\textquotesingle{}+o(A\textquotesingle{})) \(\leq\)
  2\^{}(\(\lambda\)\_base+o(\(\lambda\)\_base)) forces
  A\textquotesingle{} \(\leq\) \(\lambda\)(A,x) + o(\(\lambda\)\_base).
\end{itemize}

Combining both, \(\lambda\)(A,x) \(\leq\) A\textquotesingle{} \(\leq\)
\(\lambda\)(A,x) + o(\(\lambda\)\_base), so \(\lambda\)(A,x) is
canonical and tight up to subpolynomial factors. \ensuremath{\square}

\textbf{Interpretation.} Proposition 5.1.3.3 establishes that
vertex-capacitated min-cut is not merely a convenient analogy for
complexity - it is the \textbf{unique canonical aggregator} (up to
subpolynomial factors) for our semantic framework. The admissibility
axioms (ADD, MON, REP, CUT, TIGHT, NORM) capture natural requirements
for any reasonable complexity measure in DAG-structured computation. The
proposition proves that any measure satisfying these axioms must
coincide with min-cut capacity, thus elevating the connection from
"useful perspective" to "mathematical necessity." This is why the same
\(\lambda\) value yields consistent, tight bounds across all paradigms:
the underlying mathematical structure uniquely determines the complexity
landscape.

\textbf{Axiom design notes:} ADD explicitly acknowledges that
independence comes from instance properties (A6/H1-H5), not the
aggregator itself. CUT now explicitly captures the "minimum over cuts"
(bottleneck) property, which previously emerged implicitly from the
combination of necessity and achievability. NORM ensures the measure
correctly handles the trivial case of full resolution. Together, these
six axioms form a complete, minimal characterization of semantic
complexity aggregators for DAG-structured problems.

\textbf{Why not alternative aggregators?}

While Proposition 5.1.3.3 proves min-cut is the unique canonical
aggregator, it is instructive to see explicitly why natural alternatives
fail the admissibility requirements:

\textbf{Why Alternative Measures Fail}

\begin{itemize}
\tightlist
\item
  \textbf{Maximum cut}: Claims 2\^{}(max\_C \(\Sigma\)(R\_v-q\_v))
  states needed; Why it fails: Violates CUT (takes maximum instead of
  minimum over cuts, not the bottleneck) and TIGHT (algorithms can avoid
  worst cut, making bound unachievable)
\item
  \textbf{Average cut}: Uses 2\^{}(avg\_C \(\Sigma\)(R\_v-q\_v)); Why it
  fails: Violates CUT (averaging obscures bottleneck; does not identify
  minimum unavoidable residual)
\item
  \textbf{Sum over all cuts}: Counts \(\Sigma\)\_C
  2\^{}(\(\Sigma\)(R\_v-q\_v)); Why it fails: Violates ADD (overcounts
  shared states; does not aggregate by simple sum of residuals on a cut)
\item
  \textbf{Treewidth}: Tree decomposition dependent; Why it fails:
  Violates REP (depends on decomposition choice, not purely on semantic
  residuals R\_v - q\_v)
\item
  \textbf{Spectral gap}: Eigenvalue based; Why it fails: Violates ADD
  (does not aggregate by \(\Sigma\)(R\_v-q\_v); uses linear algebra
  structure instead of residual sums) and REP (depends on graph matrix
  representation)
\item
  \textbf{Max-flow value} (as opposed to min-cut capacity): Using the
  maximum flow f\_max as complexity measure; Why it fails: Violates CUT
  (measures transmission capability, not minimum distinguishability
  requirement; does not compute minimum over cuts) and ADD (flow values
  do not correspond to residual sums)
\end{itemize}

Each failure maps to violation of specific admissibility axioms,
confirming that the axioms correctly capture the essential properties of
DAG-structured complexity measures. Notably, no alternative satisfies
all six axioms - even measures that seem plausible (max-cut, average)
fail on fundamental requirements (bottleneck structure, achievability).

\textbf{Summary:} Min-cut captures the unavoidable bottleneck - the
minimum uncertainty any algorithm must handle. With the Frontier-Gate
mechanism (§5.1.2), we achieve both necessity (\(\geq\)
2\^{}(\(\lambda\)(A,x))) and sufficiency (algorithms achieve
2\^{}(\(\lambda\)\_base+o(\(\lambda\)\_base)) via Theorem C.2.T),
confirming min-cut as the precise characterization.

\paragraph{5.2 Algorithm classes analyzed (semantic
only)}\label{52-algorithm-classes-analyzed-semantic-only}

\textbf{The Key Question:} How do different algorithmic paradigms handle
the distinguishable artifacts forced by the conservation law? Each
paradigm has its characteristic way of managing the 2\^{}(R\_v-q\_v)
possibilities that arise when avoiding resolution.

\textbf{Critical Insight for L*:} Our constructed language L* has a
special property - the DAG dependency structure (Seed\_v depends on
parent outputs) forces any correct algorithm to distinguish artifacts
(states) by their seed-history. This means paradigms cannot "cheat" by
merging inequivalent artifacts that should not be merged.

We analyze the following \textbf{semantics}, independent of any hardware
model:

\begin{itemize}
\item
  \textbf{EO / Backtracking.} Builds a \textbf{decision tree}. At node v
  it may resolve some facts (by reading designated addresses); if it
  will not resolve enough, it must \textbf{branch} and keep
  distinguishable artifacts alive. \emph{Artifact:} number of
  \textbf{tree nodes}.
\item
  \textbf{DP / Memoization.} Maintains a \textbf{state table} keyed by a
  declared \textbf{state key}. For the DAG-structured problem L*:

  \begin{itemize}
  \tightlist
  \item
    States at node v must be seed-consistent - any representation that
    deterministically yields the same designated addresses
    F\_overlay(Seed\_v; j,\(\ell\)) qualifies (Seed\_v itself is
    sufficient but not necessary; representation-invariant)
  \end{itemize}

  Merging states with different seeds is semantically incorrect on our
  instances. \emph{Artifact:} number of \textbf{states/keys}.

  \textbf{Why DP must track resolution history on L*:} Consider two DP
  states at node v:

  \begin{itemize}
  \tightlist
  \item
    State A has Seed\_v\^{}A from one set of ancestor resolutions
  \item
    State B has Seed\_v\^{}B from different ancestor resolutions
  \end{itemize}

  If we attempted to merge these states:

  \begin{itemize}
  \tightlist
  \item
    State A needs specific addresses derived from Seed\_v\^{}A
  \item
    State B needs specific addresses derived from Seed\_v\^{}B
  \item
    A merged state can compute only ONE Seed\_v value
  \end{itemize}

  Since Enc is injective, choosing either Seed\_v value makes the
  algorithm \textbf{compute wrong addresses} for the other history,
  yielding \textbf{incorrect answers}. Therefore, any correct DP
  implementation for L* must keep states with different seed histories
  distinct.

  \begin{quote}
  \textbf{DP Correctness Box (One-line proof for L*):} States with
  different Seed\_v must not merge: different seeds \(\Rightarrow\)
  different designated addresses F\_overlay(Seed\_v; j,\(\ell\))
  \(\Rightarrow\) any merged state computes wrong addresses for at least
  one seed \(\Rightarrow\) incorrect output. Hence Seed\_v (or an
  equivalent sufficient statistic) is mandatory in the DP state key.
  \end{quote}
\item
  \textbf{OBDD.} Represents simultaneously distinguishable artifacts as
  a binary decision diagram. Using the expander-parity gate ensures
  order-robust bounds: for any variable order, avoiding resolution
  forces \textbf{width} 2\^{}(\(\Omega\)(R\_v-q\_v)) at some cut. (Note:
  order-robustness requires expander-FG gates; with identity gates,
  bounds are order-specific.) \emph{Artifact:} \textbf{diagram size} via
  width layers.
\item
  \textbf{Resolution / CDCL (proof system).} Refutes the CNF that embeds
  the gate. For correct refutation of L*, skipping R\_v\(-\)q\_v forced
  literals at node v forces \textbf{clause width}
  \(\Omega\)(R\_v\(-\)q\_v); widths \textbf{add} across cuts to a global
  width parameter, and the Ben-Sasson-Wigderson width\(\to\)size bridge
  then yields exponential \textbf{size} (see Appendix G for the complete
  proof). \emph{Artifact:} \textbf{proof size}.
\item
  \textbf{Streaming / Oblivious (verifier).} Reads a fixed schedule per
  node with a small running state S. Each pass accumulates at most S
  resolved bits; \textbf{passes \(\geq\)
  \(\lceil\)(R\_v-q\_v)/S\(\rceil\)}. \emph{Artifact:} \textbf{passes}
  (they \textbf{add} across paths).
\item
  \textbf{One-path verification (any model).} Follows \textbf{one}
  candidate; resolves all R\_v at each node; costs \textbf{add} and are
  polynomial. This is the standard NP verifier.
\end{itemize}

\textbf{Uses (per paradigm, §5):}

\begin{itemize}
\tightlist
\item
  EO/Backtracking: Theorem 7.A + Corollary 5.3.2 part 2 {[}Adapter{]}
  (tree nodes \(\geq\) Alt\_v)
\item
  DP/Memoization: Theorem 7.A + Corollary 5.3.2 part 1 {[}Adapter{]}
  (state keys \(\geq\) Alt\_v)
\item
  OBDD: Theorem 7.A + Corollary 5.3.2 part 3 {[}Thm{]}; order-robust
  width bound via Appendix B.1 (requires expander-parity; §6.2.10;
  {[}WEG00{]}, {[}HOO06{]})
\item
  Resolution/CDCL: Theorem 7.A + Corollary 5.3.2 part 4 {[}Thm{]};
  width\(\to\)size via Appendix G ({[}BEN01{]})
\item
  Streaming/Oblivious: Theorem 7.A + passes accounting (per-path
  additive bound)
\end{itemize}

\textbf{Lemma (Fork-and-Merge Soundness).} Two partial states may merge
iff their frontiers maintain identical (Seed,y) tuples. Otherwise the
merge is semantically unsound on L*, since descendants derive distinct
addresses from distinct seeds.

\textbf{Lemma 5.3.1 (Representation-Invariant Alt\_v).} For any correct
solver and any node v, let two artifact instances \(\alpha\), \(\beta\)
be called equivalent if, on every transcript consistent with their use,
they compute the same designated addresses for all selector indices
(j,\(\ell\)) at node v on the current seed chain and lead to the same
downstream behavior. Then \(\alpha\) \(\equiv\) \(\beta\) iff their
frontiers share identical (Seed,y) tuples. In particular, if
Seed\_v(\(\alpha\)) \(\neq\) Seed\_v(\(\beta\)), then \(\alpha\) and
\(\beta\) are not equivalent: there exists (j,\(\ell\)) with
F\_overlay(Seed\_v(\(\alpha\)); j,\(\ell\)) \(\neq\)
F\_overlay(Seed\_v(\(\beta\)); j,\(\ell\)), so reusing one for both is
incorrect. Therefore Alt\_v counts equivalence classes of artifacts
modulo internal encoding; compression or alternative coordinate systems
do not reduce Alt\_v. (See Lemma 4.2.2 for a sufficient-statistic
formulation.)

\textbf{Corollary 5.3.2 (Paradigm Projection Lemmas - Explicit Bounds).}
The representation-invariant principle (Lemma 5.3.1) together with the
semantic lower bound (Lemma 4.2.2) directly yields paradigm-specific
artifact bounds:

\begin{enumerate}
\def\labelenumi{\arabic{enumi}.}
\item
  \textbf{{[}Adapter{]} DP/Memoization Projection:} Any correct DP
  algorithm for L* at node v maintains \(\geq\) Alt\_v \(\geq\)
  2\^{}(R\_v-q\_v) distinct state keys.

  \emph{Proof.} DP states are artifacts. By Lemma 5.3.1, two states are
  equivalent iff they share identical Seed\_v (else they compute
  different designated addresses \(\to\) incorrect outputs on descendant
  nodes). By Lemma 4.2.2, with R\_v-q\_v unresolved bits, there are
  \(\geq\) 2\^{}(R\_v-q\_v) inequivalent seed-consistent states.
  Therefore the DP table must maintain \(\geq\) 2\^{}(R\_v-q\_v)
  distinct keys. Across a cut C, keys multiply: \(\geq\)
  2\^{}(\(\Sigma\)\_\{v\(\in\)C\}(R\_v-q\_v)) = 2\^{}(\(\lambda\)(C)).
  \(\blacksquare\)
\item
  \textbf{{[}Adapter{]} Backtracking/Decision Tree Projection:} Any
  correct decision tree exploring L* at node v has \(\geq\) Alt\_v
  \(\geq\) 2\^{}(R\_v-q\_v) tree nodes.

  \emph{Proof.} Tree nodes (partial states) are artifacts. By Lemma
  5.3.1, two nodes are equivalent iff they have identical Seed\_v. By
  Lemma 4.2.2, \(\geq\) 2\^{}(R\_v-q\_v) inequivalent seed-consistent
  nodes exist. Therefore tree size \(\geq\) 2\^{}(R\_v-q\_v). Across a
  cut, tree nodes multiply: \(\geq\) 2\^{}(\(\lambda\)(C)).
  \(\blacksquare\)
\item
  \textbf{{[}Thm{]} OBDD Projection: Width \(\geq\)
  2\^{}(\(\Omega\)(R\_v-q\_v)) with expander-FG gates (order-robust).
  Proven in Appendix B.1} (Lemma B.1).
\item
  \textbf{{[}Thm{]} Resolution Projection: Clause width \(\geq\)
  R\_v-q\_v \(\to\) size \(\geq\) 2\^{}(\(\Omega\)(R\_v-q\_v)) via
  Ben-Sasson-Wigderson. Proven in Appendix G} (Theorem G.1).
\end{enumerate}

\textbf{Proof chain summary:} Lemma 4.2.2 (information-theoretic bound
on sufficient statistics) \(\to\) Lemma 5.3.1
(representation-invariance) \(\to\) Corollary 5.3.2 (paradigm-specific
projections). This establishes that the SCL bottleneck manifests
universally across computational models.

We make \textbf{no} additional machine assumptions (no bandwidth, head
speed, word size) for the results. Randomized algorithms: coin-fixing
establishes per-instance deterministic bounds (§8; Theorem 8.A).

\paragraph{5.3 How Correctness Manifests in Each
Paradigm}\label{53-how-correctness-manifests-in-each-paradigm}

Correctness requires maintaining distinctions between computational
paths that lead to different outputs. Merging distinct possibilities
without learning what distinguishes them produces incorrect answers.

\textbullet{} \textbf{Explicit search / backtracking (EO):} may branch on guesses;
may not merge distinct partial worlds before resolution. (Natural to
paradigm)

\textbullet{} \textbf{DP (Dynamic Programming):} States at node v must be
\textbf{seed-consistent} for correctness on L* (see §5.2 for
exposition). This follows from the Keyedness correctness rule (see §4.4;
formal proofs in Appendix B): merging states with different seeds would
lead to incorrect addresses at descendant nodes. Any representation that
deterministically yields the same designated addresses qualifies
(representation-invariant).

\textbf{Note on DP correctness for L*:} The dependency mechanism
(Seed\_v per DAG structure) means that any correct DP solver must
distinguish states by their seed history. Two states with different
seeds lead to different addresses and thus access different memory
locations at descendant nodes, so merging them would be incorrect. This
is a consequence of L*\textquotesingle{}s structural requirements, not
an artificial restriction. Under this correctness requirement, distinct
seed histories induce disjoint key spaces; when Alt\_v indistinguishable
worlds remain at node v (Alt\_v \(\geq\) 2\^{}(R\_v - q\_v) by SCL),
correctness forces \(\geq\) Alt\_v distinct keys.

\textbullet{} \textbf{OBDD:} nodes represent variable decisions; the expander parity
gate {[}HOO06{]} ensures width blow-up for any variable order

\textbullet{} \textbf{Resolution/CDCL:} only resolution inferences are allowed;
clause width tracks distinguishing information (a width-w clause
distinguishes \(\leq\) 2\^{}w worlds; ruling out 2\^{}\(\lambda\) worlds
requires width \(\geq\) \(\lambda\)). (Uses Tseitin CNF encoding)

\begin{center}\rule{0.5\linewidth}{0.5pt}\end{center}

\paragraph{5.4 Problem Statement and Main Results
Preview}\label{54-problem-statement-and-main-results-preview}

All paradigm-specific lower bounds in this subsection are parameterized
by the run-dependent bottleneck \(\lambda\)(A,x) := min\_C
\(\Sigma\)\_\{v\(\in\)C\}(R\_v\(-\)q\_v).

Given an overlaid instance x* (constructed in §6) and a solver:

\begin{itemize}
\item
  \textbf{Verification problem:} Given witness w, accept iff all checks
  pass. Polynomial time by resolving fully (q\_v = R\_v at each node).
\item
  \textbf{Decision/search problem:} Without witness, decide YES/NO. We
  prove super-polynomial lower bounds on artifacts (profile-dependent;
  quasi-polynomial for QP-sharp, exponential for flat):

  \begin{itemize}
  \tightlist
  \item
    \textbf{Backtracking:} Tree size \(\geq\) 2\^{}(\(\lambda\)(A,x))
    nodes
  \item
    \textbf{Dynamic Programming:} Table size \(\geq\)
    2\^{}(\(\lambda\)(A,x)) distinct keys
  \item
    \textbf{OBDD:} Diagram size \(\geq\)
    2\^{}(\(\Omega\)(\(\lambda\)(A,x))) nodes
  \item
    \textbf{Resolution:} Proof size \(\geq\)
    2\^{}(\(\Omega\)(\(\lambda\)(A,x))) clauses
  \end{itemize}
\end{itemize}

\subparagraph{\texorpdfstring{5.4.1 Bridges: Artifact Counts \(\to\)
Standard
Resources}{5.4.1 Bridges: Artifact Counts \textbackslash to Standard Resources}}\label{541-bridges-artifact-counts--standard-resources}

The lower bounds above connect \textbf{abstract artifact counts} (Alt\_v
from SCL) to \textbf{concrete resource measures} (time, space, proof
size) via paradigm-specific bridges:

\begin{itemize}
\item
  \textbf{Backtracking / Explicit Search}

  \begin{itemize}
  \tightlist
  \item
    \emph{Artifact}: Concurrent branches at frontier
  \item
    \emph{Bridge}: time \(\geq\) \#visited nodes; space \(\geq\) max
    concurrent branches
  \item
    \emph{Result for L*:} Tree size \(\geq\) 2\^{}(\(\lambda\)(A,x))
  \item
    \emph{Citation}: Direct (§7.3.3)
  \end{itemize}
\item
  \textbf{Dynamic Programming}

  \begin{itemize}
  \tightlist
  \item
    \emph{Artifact}: Distinct table keys (seed-consistent)
  \item
    \emph{Bridge}: space \(\geq\) \#keys at widest cut; time \(\geq\)
    \#distinct key computations
  \item
    \emph{Result for L*:} Keys \(\geq\) 2\^{}(\(\lambda\)(A,x))
  \item
    \emph{Citation}: Direct (§7.3.3)
  \end{itemize}
\item
  *\emph{Resolution / CDCL*}

  \begin{itemize}
  \tightlist
  \item
    \emph{Artifact}: Clause width (variables tracked)
  \item
    \emph{Bridge}: Width \(\Omega\)(\(\lambda\)) \(\to\) Size \(\geq\)
    2\^{}(\(\Omega\)(\(\lambda\))) via Ben-Sasson \& Wigderson
  \item
    \emph{Result for L*:} Proof size \(\geq\)
    2\^{}(\(\Omega\)(\(\lambda\)(A,x)))
  \item
    \emph{Citation}: {[}BEN01{]} (§7.3.3)
  \end{itemize}
\item
  \textbf{OBDD / ROBDD}

  \begin{itemize}
  \tightlist
  \item
    \emph{Artifact}: Diagram width at some level
  \item
    \emph{Bridge}: \(\exists\)level: width \(\geq\)
    2\^{}(\(\Omega\)(\(\lambda\))) (order-robust with expander gates)
  \item
    \emph{Result for L*:} Size \(\geq\)
    2\^{}(\(\Omega\)(\(\lambda\)(A,x)))
  \item
    \emph{Citation}: {[}WEG00{]} (§7.3.3)
  \end{itemize}
\item
  \textbf{Deterministic k-tape TM}

  \begin{itemize}
  \tightlist
  \item
    \emph{Artifact}: Rollback segments (single-run lane; §7.3, Appendix
    C)
  \item
    \emph{Bridge}: Segments m\_seg \(\geq\) 2\^{}(\(\rho\)-s) where
    \(\rho\) = effective residual, s = pre-final agreement (§4.3);
    per-segment cost \(\Omega\)(n/W\_min)
  \item
    \emph{Result for L*:} Time \(\geq\) 2\^{}(\(\rho\)-s) ·
    \(\Omega\)(n/W\_min)
  \item
    \emph{Citation}: FG+Segment Counting (Appendix C.1-C.2)
  \end{itemize}
\end{itemize}

\textbf{Key observations:}

\begin{enumerate}
\def\labelenumi{\arabic{enumi}.}
\item
  \textbf{Same bottleneck, different units}: All paradigms face the same
  structural requirement (exponentially many distinguishable artifacts
  at the min-cut bottleneck), but measure it in different metrics
  (branches vs keys vs width vs segments).
\item
  \textbf{Hypotheses matter}:

  \begin{itemize}
  \tightlist
  \item
    OBDD order-robustness requires expander-FG gates (§6.2.10;
    {[}HOO06{]})
  \item
    Resolution width\(\to\)size requires Tseitin CNF encoding and
    standard width-bottleneck theorem
  \item
    TM time bounds require Frontier-Gate wiring (§6.2.8) for tight
    single-run analysis
  \end{itemize}
\item
  \textbf{Adapters preserve bound}: Each bridge (artifact count \(\to\)
  resource) is proven independently, then composed with
  SCL\textquotesingle{}s cut-level bound (\(\Sigma\)\_\{v\(\in\)C\}
  log\(_2\)(Alt\_v) \(\geq\) \(\lambda\)(A,x) at min-cut C; §7.2.1) to
  yield paradigm-specific lower bounds.
\end{enumerate}

\emph{Full formalizations}: §7.3.3 (Artifact Extractors and Cost
Bridges) provides explicit extractor functions E\_v and bridge
mechanisms for each paradigm, with complete proofs of bounds.

\textbf{Cross-reference}: The abstract SCL (Theorem 7.A; §7.2) is
paradigm-agnostic; this table shows how to instantiate it in each
computational model.

These results establish the dichotomy: algorithms either resolve nearly
all information (polynomial verification) or face super-polynomial
artifacts (search; quasi-polynomial for QP-sharp, exponential for flat).

\begin{center}\rule{0.5\linewidth}{0.5pt}\end{center}

\paragraph{5.5 Technical Extensions: Machine-Specific
Bounds}\label{55-technical-extensions-machine-specific-bounds}

\subparagraph{5.5.1 Arity-Bounded Per-Step Information
(TMs)}\label{551-arity-bounded-per-step-information-tms}

(Worst-case analysis uses Lemma 5.5.1.b.)

We model a deterministic k-tape TM with tape alphabet \(\Gamma\). In
each step the machine \textbf{observes} at most the k-tuple of symbols
under its heads (plus finite control), then updates state and tapes.
This bounds how many \textbf{fresh} bits can be learned per step.

\textbf{Lemma 5.5.1 (Per-step information inflow, TM).} Let B :=
k\(\lceil\)log\(_2\)\textbar{}\(\Gamma\)\textbar{}\(\rceil\). In any
step, the verified transcript\textquotesingle{}s mutual information with
the unresolved fresh bits can increase by \textbf{at most B bits}.
Equivalently, a TM can learn \(\leq\) B fresh bits per step.

Formally, if X\_\{unres\} denotes the unresolved fresh bits of the
current node and T\_t the verified transcript up to step t, then
I(X\_\{unres\}; T\_\{t+1\}) - I(X\_\{unres\}; T\_t) \(\leq\) B.

\emph{Proof (outline; see also Lemma 5.5.1.d for a combinatorial
version).} The step\textquotesingle{}s new observable data from the
instance is a k-symbol tuple over \(\Gamma\), providing \(\leq\)
k\(\lceil\)log\(_2\)\textbar{}\(\Gamma\)\textbar{}\(\rceil\) bits.
\textbf{State}, \textbf{head moves}, and \textbf{writes} are
deterministic functions of the prior transcript and coins and, absent
new symbol reads, add \textbf{no mutual information about unresolved
bits}. By \textbf{Hermeticity}, no other channel leaks fresh bits; by
\textbf{RWA}, only the \textbf{first use} of input/witness symbols is
chargeable. Hence the per-step inflow is \(\leq\) B. \ensuremath{\square}

Parenthetical note. SCL corresponds to Hartley/Rényi-0 (\(\Phi\) :=
log\(_2\) Alt) for counting worst-case distinguishable artifacts, while
Lemma 5.5.1 phrases a per-step inflow bound in Shannon terms; these are
consistent because the inflow cap prices first-use revelations (q
growth), not \(\Phi\).

\textbf{Lemma 5.5.1.b (RWA bits imply read steps).} Under
Receiving-Window Attribution (RWA), every counted fresh bit corresponds
to the first valid use of some input/witness symbol, and thus to at
least one actual read of a designated cell. Consequently, any
(contiguous) run segment that accumulates Q fresh bits entails at least
\(\lceil\)Q/B\(\rceil\) TM steps. \emph{Proof (outline; see Lemma D.2.1
for the step accounting).} Each RWA-counted bit is attributable to a
unique first revealing read on a designated address; Hermeticity
excludes other channels. Grouping at most B fresh bits per step by Lemma
5.5.1 and applying Lemma D.2.1 yields the step lower bound. \ensuremath{\square}

\textbf{Lemma 5.5.1.c (TM parity computation cost).} For a k-tape TM
computing a gate digest GateDigest\_v =
\(\oplus\)\emph{\{(u,(j,\(\ell\))) \(\in\) S(P)\}
\(\sigma\)}\{F\_overlay(Seed\_u; j,\(\ell\))\} where
\textbar{}S(P)\textbar{} = M:

\begin{itemize}
\tightlist
\item
  Computing the XOR requires \(\Omega\)(M) tape operations even with all
  salts \(\sigma\) pre-cached
\item
  This holds regardless of when salts were read (pre-scan allowed)
\end{itemize}

\emph{Proof.} Let S(P) = \{(j\_1,\(\ell\)\_1),...,(j\_M,\(\ell\)\_M)\}
denote the M gate path indices. Under the current Seed\_v, these map to
designated addresses u\_i = F\_overlay(Seed\_v; j\_i,\(\ell\)\_i)
containing salts \(\sigma\)\_\{u\_1\},...,\(\sigma\)\_\{u\_M\}. Consider
any accepting TM that outputs the parity \(\oplus\)\emph{\{i=1\}\^{}M
\(\sigma\)}\{u\_i\}.

(i) Sensitivity: Parity has full sensitivity - flipping any input bit
flips the output. Therefore any correct algorithm must (logically)
incorporate the value of each \(\sigma\)\_\{u\_i\} at least once;
otherwise there exists an input assignment differing only at an
untouched \(\sigma\)\_\{u\_i\} that the algorithm cannot distinguish,
leading to an incorrect parity.

(ii) Operational cost on TMs: Even if \(\sigma\)\_\{u\_i\} are cached on
a work tape from an earlier pre-scan, in a fixed segment under a fixed
Seed\_v the machine must (a) compute each address u\_i =
F\_overlay(Seed\_v; j\_i,\(\ell\)\_i), (b) fetch the cached bit for
u\_i, and (c) XOR it into an accumulator. On a k-tape TM, fetching M
distinct cells on a work tape costs \(\Omega\)(M) head moves in the
worst case (cell-probe style: each distinct cell requires at least one
access). No re-use across different u\_i is possible because each
contributes independently to the parity.

A simple spacing argument makes this explicit: adversarially place the
cached values for u\_1,...,u\_M at pairwise distinct tape cells
separated by at least one blank cell. Any correct algorithm must access
each distinct cell at least once (by (i)), and each access incurs
\(\Omega\)(1) head motion. Summing over M distinct cells gives
\(\Omega\)(M) tape operations. Multiple tapes or pre-scanning only
change constant factors; the lower bound counts distinct cell probes.

Combining (i) and (ii) yields an \(\Omega\)(M) lower bound on the number
of tape operations to compute the digest, irrespective of when
\(\sigma\) were read. \(\blacksquare\)

\textbf{Corollary 5.5.1.a (Per-node baseline under arity).} If node v
requires learning R\_v-q\_v fresh bits, then any accepting run must
execute at least

\(\lceil\)(R\_v \(-\) q\_v)/B\(\rceil\)

steps inside the receiving window of node v.

\textbf{Lemma 5.5.1.d (Combinatorial per-step fresh-bit bound; no
information theory).} On a deterministic k-tape TM with alphabet
\(\Gamma\), in any step at most B :=
k\(\lceil\)log\(_2\)\textbar{}\(\Gamma\)\textbar{}\(\rceil\) fresh bits
(in the RWA sense) can be credited. Consequently, any run segment
accruing Q fresh bits contains at least \(\lceil\)Q/B\(\rceil\) steps.

\emph{Proof.} In one step the machine reads at most k tape symbols; each
symbol contributes at most
\(\lceil\)log\(_2\)\textbar{}\(\Gamma\)\textbar{}\(\rceil\) new bits. By
Hermeticity, the only chargeable fresh information comes from first
reads of designated cells; by RWA, each credited bit is tied to such
first-use reads within the segment. Group at most B credited bits per
step; summing yields the \(\lceil\)Q/B\(\rceil\) step bound.
\(\blacksquare\)

\textbf{Parameterization and scope.} In the standard complexity setting,
k,\textbar{}\(\Gamma\)\textbar{} are fixed constants, so B=O(1). If k or
\textbar{}\(\Gamma\)\textbar{} vary with input length n, replace B by
B(n)=k(n)\(\lceil\)log\(_2\)\textbar{}\(\Gamma\)(n)\textbar{}\(\rceil\)
throughout; all bounds then scale by 1/B(n).

\textbf{Randomness.} For randomized algorithms, the bound applies
\textbf{per fixed coin string} (via coin-fixing; Yao\textquotesingle{}s
principle). Algorithmic coins are independent of the instance and do not
increase mutual information with unresolved bits. RWA ensures that later
re-reads of symbols or of machine-written data do not double-charge.

\textbf{Adapters.} For Word-RAM/PRAM with word width w, the per-step
inflow bound becomes B := (words read per step) · w.

\textbf{Remark (Segments vs. steps).} The arity bound prices
\textbf{steps directly} and does \textbf{not} rely on segment
partitioning. Segments appear later only as an analytic tool in Appendix
C.2 (Segment Counting) to bound \textbf{pre-final agreement} s when we
seek \textbf{tight single-run} lower bounds (via FG). They are
\textbf{not} needed to obtain the per-try step baseline.

\textbf{Lemma 5.5.1.e (RWA monotonicity and subadditivity).} Let
q(C;\(\pi\)) denote the RWA-credited fresh bits across a cut C computed
from a transcript prefix \(\pi\). Then:

\begin{itemize}
\tightlist
\item
  (Monotonicity) If \(\pi\) \ensuremath{\preceq} \(\pi\)', then q(C;\(\pi\)) \(\leq\)
  q(C;\(\pi\)').
\item
  (Schedule invariance) Reordering non-first-use actions (re-reads,
  local writes, head moves) in \(\pi\) does not change q.
\item
  (Subadditivity) For concatenated segments \(\pi\)=\(\sigma\)\(\tau\)
  with disjoint sets of first-use events across C, q(C;\(\pi\)) =
  q(C;\(\sigma\))+q(C;\(\tau\)).
\item
  (Per-try minimum) In any try that changes the resolution prefix across
  C, at least one first-use event occurs across C.
\end{itemize}

\emph{Proof sketch.} By definition, RWA credits a bit exactly at its
first revealing designated read. Such events are determined by the
earliest step at which the symbol is both read and first needed for
correctness under Hermeticity; later schedule changes cannot move a
first-use earlier than the earliest read of its designated cell, so
credits are monotone under transcript extension and invariant under
permutations that do not introduce earlier first-reads. First-use events
are unique per designated payload, so disjoint segments add. A try that
changes the resolution prefix must import at least one new designated
bit across C to distinguish the accepting outcome from prior
commitments; otherwise acceptance could have been decided before the
try. \(\blacksquare\)

\subparagraph{5.5.2 Finality Depth (Advanced
Systems)}\label{552-finality-depth-advanced-systems}

Some systems treat nodes with distance \textgreater{} d from current
frontier as \textbf{final} (no rewrites). Model this by declaring nodes
at distance \textgreater{} d as final when processing node v.

\textbf{Effect.} The live frontier \textbf{resets} periodically; space
lower bounds improve to max over windows of depth d, and product
arguments can be modularized across windows. This mirrors blockchain
confirmations and serializable databases with checkpoints.

\emph{Note: This extension primarily applies to distributed systems and
can be skipped for standard complexity analysis.}

\paragraph{5.6 Practical Considerations: Cut
Selection}\label{56-practical-considerations-cut-selection}

\textbf{Finding optimal cuts} (minimal \(\lambda\) =
\(\Sigma\)\_\{v\(\in\)C\}(R\_v - q\_v) residual): Note that q is
run-dependent (credited at first use), so practical selection combines
the structural profile \(\lambda\)\_base with observed q\_v over the
run.

\begin{itemize}
\tightlist
\item
  \textbf{Layer cuts:} Any layer \(\ell\) separating roots from sinks
\item
  \textbf{Antichains:} Sets hitting all root\(\to\)sink paths
\item
  \textbf{Reconvergence hotspots:} Heavy fan-in/fan-out regions
\item
  \textbf{Algorithms:} Standard max-flow/min-cut (Ford-Fulkerson, Dinic)
  with node capacities R\_v - q\_v
\end{itemize}

\begin{center}\rule{0.5\linewidth}{0.5pt}\end{center}

\paragraph{5.7 Parallelism: Time vs
Work/Bandwidth/Depth}\label{57-parallelism-time-vs-workbandwidthdepth}

\textbf{SCL (model-agnostic).} For any correct computation, across the
min-cut C*,

\(\Sigma\)\_\{v\(\in\)C*\} q\_v + \(\Sigma\)\_\{v\(\in\)C*\}
log\(_2\)(Alt\_v) \(\geq\) \(\Sigma\)\_\{v\(\in\)C*\} R\_v, \(\lambda\)
:= \(\Sigma\)\_\{v\(\in\)C*\}(R\_v - q\_v).

This conservation law holds for all classical models solving L* - the
need for 2\^{}\(\lambda\) distinguishable states is universal for L*.

\textbf{Note:} This section provides an overview. For the formal
round-credit SCL formulation that captures both depth and work
precisely, see §7.3.11.

\textbf{Parallel complexity bounds.} In parallel models, the
instantaneous legitimate progress scales with resources, changing the
complexity landscape. Here \(\beta\)* denotes a per-round upper bound on
legitimate cross-cut credit across the min-cut C* (sum of per-processor
contributions; see §7.3.11):

\textbf{T\_par \(\geq\) max\{2\^{}\(\lambda\)/P,
\(\lambda\)/\(\beta\)*\}}

where P is the number of processors and \(\beta\)* is the per-round
legitimate-credit budget across the min-cut (§7.3.11 provides precise
definitions and proofs).

\textbf{Work law.} Total legitimate processing across all processors is
at least 2\^{}\(\lambda\) seed-consistent worlds (reads or priced
refutations); hence

W := P × T\_par \(\geq\) \(\Omega\)(2\^{}\(\lambda\)) (Work)

independent of parallelism.

\textbf{Model-specific bounds:}

\begin{itemize}
\tightlist
\item
  \textbf{Sequential k-tape TM:} \(\beta\)* = B = O(1), so T \(\geq\)
  2\^{}\(\lambda\)/B (our main result)
\item
  \textbf{PRAM with P processors:} \(\beta\)* = \(\Theta\)(P × B), so
  T\_par \(\geq\) max\{2\^{}\(\lambda\)/P, \(\lambda\)/(P×B)\}
\item
  \textbf{Boolean circuits:} Depth \(\geq\) \(\lambda\)/w\_circ where
  w\_circ is circuit width (gates per layer), Size \(\geq\)
  2\^{}\(\lambda\) (adapter via width/size bridges; see §5.4 and
  Appendix G)
\item
  \textbf{Distributed systems:} Rounds \(\geq\)
  \(\lambda\)/\(\beta\)\_net where \(\beta\)\_net is per-round network
  bandwidth
\end{itemize}

\textbf{Key insight:} Parallelism can reduce depth/time by widening the
instantaneous funnel (increasing \(\beta\)*), but cannot reduce total
work below 2\^{}\(\lambda\). The distinction between sequential and
parallel complexity is fundamental: with exponentially many processors
(e.g., 2\^{}\(\lambda\)), parallel time can be O(1) while work remains
exponential.

\textbf{Dependency-only complexity.} When the bottleneck residual
\(\lambda\) is 0 (or O(log n\_core)), \textbf{multiplicative} costs do
not arise; any remaining cost is \textbf{additive/scheduling}, which is
polynomial on our explicit DAG overlays (depth = poly,
\textbar{}V\textbar{} = poly). Parallel \textbf{span} is captured by D
\(\geq\) \(\lambda\)/\(\beta\)* (Round-credit SCL); with \(\lambda\) = 0
the bound yields \textbf{no super-poly span}.

\textbf{Relation to P vs NC:} The round-credit formulation (§7.3.11)
explains why some P problems resist parallelization despite low
\(\lambda\) - structural dependencies limit \(\beta\)* even with many
processors, forcing sequential depth.

\begin{center}\rule{0.5\linewidth}{0.5pt}\end{center}

\textbf{Section 5 Summary:} All paradigms obey the same conservation law
- whether tracking tree branches, table keys, diagram width, or proof
clauses. The min-cut framework reveals an unavoidable bottleneck for any
classical algorithm solving L*. In sequential models this yields
super-polynomial complexity (quasi-polynomial for QP-sharp, exponential
for flat); in parallel models it yields exponential work bounds and
depth/bandwidth bounds dependent on \(\beta\)*.

\subsection{Part III: The Construction}\label{part-iii-the-construction}

\subsubsection{6. Instance Construction \&
Invariants}\label{6-instance-construction--invariants}

Sections 4 and 5 established the abstract framework: axioms
\textbf{A1-A5} that imply the Semantic Conservation Law, and showed how
this law manifests universally across computational paradigms. But do
instances satisfying A1-A5 actually \emph{exist}? And can we make them
NP-complete?

\textbf{Purpose of Section 6:} We construct the explicit NP-complete
language \textbf{L*} - a blockchain-inspired, cryptography-free
dependency DAG that provably satisfies axioms \textbf{A1-A5}
(Hermeticity, Injectivity, Emergence, Closure, Dependency). This is the
\textbf{existence proof}: we do not just claim instances with these
antagonistic properties exist; we build them deterministically via
witness-preserving reduction from 3-SAT.

\textbf{Construction Approach.} The Structural OWF construction (§9)
uses instance axioms \textbf{A1-A5} (Hermeticity, Injectivity,
Emergence, Closure, Dependency) with per-run properties \textbf{H1-H5}
(disjoint pools, unique-neighbor, Enc injectivity, realizability, no
cross-coupling) for cut composition. Section 6 focuses on the structural
properties needed for these instance axioms.

\textbf{Roadmap:}

\begin{itemize}
\tightlist
\item
  \textbf{§6.1}: Design rationale - why this construction? (five
  antagonisms, three fundamental obstacles)
\item
  \textbf{§6.2}: Construction (one node) - seeds, designated addresses,
  decode schema, gates, FG
\item
  \textbf{§6.3}: Emergence Lemma (provable) - rank(H\_v) = R\_v forces
  R\_v fresh bits at each node
\item
  \textbf{§6.4}: Completeness (rank forcing) - y\_v = H\_v x\_v fully
  determines emergent information
\item
  \textbf{§6.5}: Dependency (no look-ahead) - parents must complete
  before children
\item
  \textbf{§6.6-§6.8}: Sizing, invariants summary, gadget glossary
\end{itemize}

\textbf{From abstract to concrete:} Section 7 will prove that
\textbf{A1-A5} \emph{mathematically imply} SCL (with H1-H5 for cut
composition). Section 6\textquotesingle{}s job is to show that A1-A5 are
\emph{realizable} in an NP-complete language. Together: realizability
(§6) + implication (§7) = structural lower bounds for L*.

Notation (sizes). See §4 (Notation). We write n := n\_core =
\textbar{}Encode(\(\varphi\))\textbar{} unless explicitly stated; n\_tot
:= \(\Sigma\)\_v R\_v = \(\Theta\)(n\_core · log n\_core) under the flat
profile.

\textbf{Overview:} This section constructs explicit instances where the
\textbf{three-dimensional computational trap} (Storage, Resolution,
Elimination; §7.5) becomes \textbf{mathematically unavoidable} for any
classical algorithm. For the Structural OWF construction (§9),
\textbf{every FG-wired instance} exhibits per-instance deterministic
hardness (Theorem 8.A) via these structural invariants.

\paragraph{6.1 Why This Construction?}\label{61-why-this-construction}

\textbf{L* is a blockchain-inspired, cryptography-free dependency DAG -
a hashless proof-of-search (PoSearch) analogue where structure (not
hashes) compels search} (PoSearch: proof of search work as a correctness
requirement; see §1.1 for why search is necessary). The construction
instantiates the framework axioms \textbf{A1-A5} from §4.2 through
specific mechanisms that enforce all three computational dimensions
simultaneously: states become unmergeable (Storage), information becomes
unresolvable (Resolution), and candidates become uneliminable
(Elimination). This section proves the key properties: Emergence (Lemma
6.1), Completeness (Lemma 6.2), Dependency (§6.5), and Injectivity
(encoding Enc §6.2.7), ensuring the Semantic Conservation Law q +
\(\Phi\) \(\geq\) R holds at each node.

\textbf{Hourglass DAG Architecture.} L* uses an hourglass-shaped DAG:
wide (variables) \(\to\) narrow (FG bottleneck) \(\to\) wide (clauses).
All information must flow through the narrow "pinch point" at the
FrontierGate (FG):

\begin{verbatim}
        Source
           \ensuremath{\downarrow}
    +------+------+------+------+
    v\_1 v\_2 v\_3 ... v\_n        \ensuremath{\leftarrow} WIDE: n variable nodes (witness \ensuremath{\alpha} enters here)
    +------+------+------+------+
           \ensuremath{\downarrow}
         +---------+
         | FG |             \ensuremath{\leftarrow} NARROW: bottleneck (R independent bits emerge, A3)
         +---------+
           \ensuremath{\downarrow}
    +------+------+------+------+
    C\_1 C\_2 C\_3 ... C\_m        \ensuremath{\leftarrow} WIDE: m clause nodes (ALL depend on FG)
    +------+------+------+------+
           \ensuremath{\downarrow}
      Reduction tree        \ensuremath{\leftarrow} O(log m) depth combining
\end{verbatim}

The FG bottleneck creates the 2\^{}R search space: wrong \(\alpha\)
\(\to\) wrong variable seeds \(\to\) wrong FG seed \(\to\) ALL clause
seeds wrong \(\to\) garbage decoding. See §6.2.8 for FG wiring details.

\subparagraph{\texorpdfstring{6.1.1 L*\textquotesingle{}s
Five-Dimensional Antagonism (Why \(\lambda\) Remains
Large)}{6.1.1 L*\textquotesingle s Five-Dimensional Antagonism (Why \textbackslash lambda Remains Large)}}\label{611-ls-five-dimensional-antagonism-why-ux3bb-remains-large}

\emph{We call an instance \textbf{antagonistic} when, \textbf{after all
admissible polynomial-time inferences} (our verifier-auditable
overlay/inference library), the min-cut residual}

\(\lambda\)(C*) = \(\Sigma\)\_\{v\(\in\)C*\}(R\_v - q\_v) =
\(\Sigma\)\_\{v\(\in\)C*\} R\_v - \(\Sigma\)\_\{v\(\in\)C*\} q\_v

\emph{remains large (where C} is a min-cut). L* is engineered so that
five \textbf{antagonistic global constraints} jointly keep \(\lambda\)
large - in stark contrast to P problems where global constraints are
\textbf{synergistic} and permit \(\lambda\) = O(log n).*

\textbf{A. Per-node antagonism (factoring-style global constraint).}
Each node enforces a full-rank global requirement y\_v = H\_v x\_v (A3:
Emergence/Rank). This constraint enforces a \textbf{unique preimage} on
the designated coordinates (rank R\_v); rejecting a candidate prunes at
most one world (\(\leq\)1 bit), consistent with SCL. A wrong x\_v yields
no shortcut to the correct x\_v. Thus, unless the designated reads are
paid, the per-node residual R\_v - q\_v remains - feeding directly into
SCL\textquotesingle{}s \(\Phi\)\_v \(\geq\) R\_v - q\_v.

\textbf{B. Cross-node antagonism (cascaded independence).} Requirements
propagate along the DAG while designated address pools are
\textbf{pairwise disjoint} (A1: Hermeticity/Disjoint Atoms + Lemma 6.5;
§6.3.2). Satisfying parents \textbf{creates} rigid, fresh obligations
downstream (A5: Dependency), but those obligations \textbf{do not
coalesce} across parallel branches - so residuals \textbf{add} across
any cut (min-cut \(\Lambda\)).

\textbf{C. Temporal antagonism (lock-in).} Seeds are
\textbf{content-addressed} (A2: Keyedness + Injectivity): early
commitments fix addresses for future reads. A wrong root choice
propagates through descendants; backtracking discards all dependent
work. This prevents "try many, merge later" strategies: \textbf{merging
non-equivalent worlds is incorrect} (collision \(\Rightarrow\) error),
so residual cannot be retroactively collapsed.

\textbf{D. Emergence antagonism (fresh requirements despite knowledge).}
Even with all parent seeds known, R\_v \textbf{fresh, independent} bits
must still \emph{emerge} at v (A3: Emergence). Global inferences cannot
predict those bits without paying legitimate reads; any attempt to
short-circuit just reencounters SCL\textquotesingle{}s per-node bound.

\textbf{E. Parallel antagonism (path isolation).} Different commitment
paths induce \textbf{incompatible futures} via content-addressing;
computations from distinct paths cannot be memoized across seeds. In
sequential time this forces many \textbf{segments}; in parallel models
we claim Work \(\geq\) 2\^{}\(\lambda\) and \textbf{Depth \(\geq\)
\(\lambda\)/\(\beta\)*} (Round-credit SCL §7.3.11) - wall-clock time can
shrink with processors but work remains exponential.

\textbf{What each antagonism structurally blocks:}

\begin{itemize}
\tightlist
\item
  \textbf{DP/memoization} handles \emph{temporal} reuse, but
  \textbf{cannot legally merge} states with different seeds under
  \textbf{path isolation} (E).
\item
  \textbf{Massively parallel search} spreads per-node attempts (A), but
  \textbf{cross-node disjointness} (B) and \textbf{Keyedness} (C)
  prevent sharing; must pay exponential \textbf{work} (and depth unless
  \(\beta\)* is huge).
\item
  \textbf{Backtracking} can revise decisions, but \textbf{temporal
  lock-in} (C) makes revision discard all dependent work; cannot
  amortize residual.
\item
  \textbf{Algebraic/global elimination} excels when constraints cohere
  (e.g., XOR-SAT), but \textbf{emergence} (D) hides fresh bits behind
  designated reads and seed-gated addresses; elimination cannot cross
  seeds without paying Q.
\end{itemize}

\textbf{Information barrier (connecting the antagonisms).} The five
antagonisms combine to create a fundamental information barrier:
\textbf{algorithms need information to make progress; L* provides no
information.} Standard problem instances provide structural clues - unit
constraints, symmetries, frequencies, correlations - that enable
algorithmic shortcuts. L*\textquotesingle{}s seed-locked encoding
(§10.1.1 OAP) systematically removes ALL structural information by
encoding the problem itself as data accessible only through correct seed
computation. Accessing any structural property of \(\varphi\) (unit
clauses, literal polarities, clause structure, variable frequencies)
requires the correct seed chain, which requires knowing the assignment
\(\alpha\), which IS the solution. This circular dependency transforms
the computational barrier from algorithmic ("we haven\textquotesingle{}t
found a fast method") to information-theoretic ("no information exists
to guide faster methods"). Every algorithmic technique - propagation,
learning, pruning, heuristics - requires information L* does not
provide. The antagonisms enforce this barrier: emergence (D) ensures
fresh bits cannot be predicted, hermeticity (A1) ensures reads cannot be
shared, keyedness (C) ensures different guesses cannot be merged, and
dependency (B) ensures local barriers multiply globally. The result:
exponential search is not an algorithmic limitation but an
information-theoretic necessity.

\textbf{Meta-effect.} These antagonisms do not merely add; through SCL
they \textbf{multiply}: residuals \textbf{add across cuts} by
disjointness; taking logs of the product of alternatives gives the
\textbf{multiplication of worlds} (SCL cut form). After all admissible
inference (credited in Q), the \textbf{residual} \(\lambda\) across C*
remains large, forcing exponential \textbf{Work} (and sequential time;
parallel depth via \(\lambda\)/\(\beta\)*).

\emph{Scope note.} This is an \textbf{intuition} subsection. Formal
lower bounds are stated in §§7-9. In parallel models we claim
\textbf{Work/Depth} bounds (Round-credit SCL), not universal time
bounds.

\textbf{Parallel span:} See §7.3.11 for the precise Round-credit SCL: D
\(\geq\) \(\lambda\)/\(\beta\)*, W \(\geq\) 2\^{}\(\lambda\).

\subparagraph{6.1.1.1 From Five Antagonisms to Three Fundamental
Obstacles}\label{6111-from-five-antagonisms-to-three-fundamental-obstacles}

The five antagonisms (A-E) described above create three orthogonal
computational barriers. These are not merely different challenges - they
are \textbf{independent dimensions} of obstruction that cannot be
optimized around:

\textbf{Obstacle 1: Unmergeable States (Space)}

L*\textquotesingle{}s injective seed construction creates
2\^{}\(\lambda\) distinct computational states that cannot be merged
without producing incorrect results.

\begin{itemize}
\tightlist
\item
  \emph{Enforced by}: Antagonisms (C) Temporal + (E) Parallel
\item
  \emph{Mechanism}: Keyedness (§4.4) - different Seed\_v values \(\to\)
  different designated addresses F\_overlay(Seed\_v; j,\(\ell\)) \(\to\)
  merging produces wrong computations (Lemma 7.I; §7.2.1)
\item
  \emph{Result}: Must maintain Alt\_v \(\geq\) 2\^{}(R\_v-q\_v)
  simultaneous artifacts at node v
\end{itemize}

\textbf{Obstacle 2: Unresolvable Information (Time-Forward)}

Fresh bits at each node cannot be inferred or predicted - must be
explicitly read, and reading is slow.

\begin{itemize}
\tightlist
\item
  \emph{Enforced by}: Antagonism (D) Emergence + per-step bandwidth
  limit
\item
  \emph{Mechanism}: R\_v fresh bits at each node (A3; §6.2.5; Lemma 6.1)
  + bandwidth \(\leq\) B = O(1) per step (Lemma 5.5.1)
\item
  \emph{Result}: Single-run lane requires m\_seg \(\geq\)
  2\^{}(\(\rho\)-s) segments (App. C.2), each costing
  \(\Omega\)(n/W\_min) steps (App. C.1.1)
\end{itemize}

\textbf{Obstacle 3: Uneliminable Candidates (Time-Backward)}

Testing wrong candidates eliminates at most one possibility (\(\leq\)1
bit) with no cascade - cannot prune efficiently.

\begin{itemize}
\tightlist
\item
  \emph{Enforced by}: Antagonism (A) Per-node + CDT
\item
  \emph{Mechanism}: Factoring-style constraint y\_v = H\_v x\_v
  (§6.1.1.A above) + no unbacked semantic progress (Lemma
  CDT-1\textquotesingle{}; App. C)
\item
  \emph{Result}: Restart lane requires \ensuremath{\mathbb{E}}{[}tries{]} \(\geq\)
  2\^{}(\(\Delta\)(C*)) independent attempts (Lemma 7.R; App. C.4.2)
\end{itemize}

\textbf{Why These Obstacles Are Independent:}

Overcoming one dimension does not reduce the others:

\begin{itemize}
\tightlist
\item
  Efficient storage does not make reading faster (Resolution) or pruning
  cheaper (Elimination)
\item
  Fast reading (Resolution) does not reduce testing costs (Elimination)
\item
  Efficient testing (Elimination) does not bypass the need to read fresh
  bits (Resolution)
\end{itemize}

This orthogonality is what makes L* maximally incompressible -
algorithms cannot trade off between dimensions or find a "fourth way"
around them.

\textbf{Composition Across the DAG:}

Antagonism (B) Cross-node ensures these local obstacles compose
multiplicatively into global exponential bounds. Because designated
address pools \{U\_v\} are pairwise disjoint (A1: Hermeticity; Lemma
6.5), residuals \textbf{add} across any cut: \(\lambda\)(C) =
\(\Sigma\)\_\{v\(\in\)C\}(R\_v - q\_v). Taking exponentials:
2\^{}(\(\lambda\)(C)) = \(\prod\)\_\{v\(\in\)C\} 2\^{}(R\_v-q\_v).

\textbf{Connection to Other Frameworks:}

These three fundamental obstacles provide the conceptual foundation for
the three-dimensional framework formalized in §7.5. See §7.5 for the
formal structure and conservation law, Appendix C for how the Two Lanes
(Restart vs. Single-Run) fail against different dimensions, and §10.5
for a comprehensive summary contrasting verification (which collapses
all three dimensions via the witness) with search (which faces
exponential barriers in all three).

\begin{center}\rule{0.5\linewidth}{0.5pt}\end{center}

\textbf{Counter-examples \& Guardrails Box}

\textbf{Antagonistic global constraints with small \(\lambda\) (easy
despite conflict):}

\begin{itemize}
\tightlist
\item
  \textbf{XOR-SAT}: Parity constraints form antagonistic global
  requirements, yet Gaussian elimination achieves \(\lambda\) = O(log
  n\_core)
\item
  \textbf{2-SAT}: Many global clauses fight locally, yet
  implication-graph SCCs achieve \(\lambda\) = O(log n)
\item
  \textbf{Maximum Flow}: Competing global capacity constraints, yet
  augmenting paths achieve \(\lambda\) = O(log n)
\end{itemize}

\textbf{Synergistic global constraints with large \(\lambda\) (hard
despite coherence):}

\begin{itemize}
\tightlist
\item
  \textbf{Unique Games (small \(\varepsilon\))}: Global constraints
  nearly consistent, yet \(\lambda\) remains large
\item
  \textbf{Factoring}: One beautiful global constraint n = pq, yet no
  polynomial-time method achieves small \(\lambda\)
\end{itemize}

\textbf{Key insight}: P problems have \textbf{synergistic global
constraints} permitting small \(\lambda\). NP-complete problems have
\textbf{antagonistic global constraints} maintaining large \(\lambda\).
What matters is the \textbf{residual \(\lambda\)}, not surface
appearance.

\begin{center}\rule{0.5\linewidth}{0.5pt}\end{center}

\paragraph{6.2 Construction (one node)}\label{62-construction-one-node}

Fix node v in DAG G=(V,E).

\subparagraph{6.2.1 Public parameters}\label{621-public-parameters}

\begin{itemize}
\tightlist
\item
  \textbf{Branch set.} J\_v with size K\_v := \textbar{}J\_v\textbar{}.
\item
  \textbf{Micro-arity.} \(\kappa\) := \(\Theta\)(log n\_core) denotes
  the number of designated atoms per branch.
\item
  \textbf{Emergence rank.} R\_v \(\leq\) K\_v\(\kappa\). (Flat profile
  chooses R\_v = \(\Theta\)(n\_core/log n\_core).)
\end{itemize}

\subparagraph{6.2.2 Anchors (DI: Distinct/Unique per
branch)}\label{622-anchors-di-distinctunique-per-branch}

Each branch j \(\in\) J\_v is associated with an \textbf{anchor} symbol
(a fixed index drawn from the core). Anchors are used only to
parameterize which designated atoms a branch references, ensuring each
branch touches \textbf{designated new atoms} and prevents global reuse.

Definition (DI). "Distinct/Unique per branch" means each branch j
touches designated atoms unique to that branch, preventing cross-branch
reuse.

\subparagraph{6.2.3 Scatter mapping (UN: Unique-Neighbor/bounded
overlap)}\label{623-scatter-mapping-un-unique-neighborbounded-overlap}

\begin{itemize}
\tightlist
\item
  Create a fresh address range U\_v (disjoint across nodes) that stores
  \textbf{salt words} \{\(\sigma\)\_u\}\_\{u \(\in\) U\_v\}, each
  \(\sigma\)\_u \(\in\) \{0,1\}\^{}(s\_salt) with s\_salt =
  \(\Theta\)(log n\_core).
\item
  Define a \textbf{seed-dependent permutation} \(\pi\)\_v:
  J\_v×{[}\(\kappa\){]} \(\to\) U\_v derived from \textbf{Seed\_v} that
  maps each pair (j,\(\ell\)) to a \textbf{distinct} designated address.
  We write F\_overlay(Seed\_v; j,\(\ell\)) := \(\pi\)\_v(j,\(\ell\)) for
  the overlay address function. Node address ranges U\_v are pairwise
  disjoint. Thus different (j,\(\ell\)) never read the same designated
  cell. This ensures disjoint designated cells give independence and
  avoid amortization.
\item
  \textbf{Pool vs. designated cells.} U\_v is the \textbf{address pool}
  for node v; the \textbf{designated} addresses used by a run are
  u\_v,\_j,\(\ell\) := F\_overlay(Seed\_v; j,\(\ell\)) \(\in\) U\_v (hence
  depend on the actual Seed\_v).
\end{itemize}

Definition (UN). "Unique-Neighbor" indicates that
(j,\(\ell\))\(\mapsto\)u is injective within node v; together with
disjoint U\_v across nodes, designated addresses never overlap.

See Appendix A.1-A.2 for explicit \(\pi\)\_v and salt layouts.

For each (j,\(\ell\)) \(\in\) L\_v := J\_v × {[}\(\kappa\){]}, let
u\_v,\_j,\(\ell\) := F\_overlay(Seed\_v; j,\(\ell\)) =
\(\pi\)\_v(j,\(\ell\)) \(\in\) U\_v be the designated cell holding
\(\sigma\)\_\{u\_v,\_j,\(\ell\)\}.

Clarification (Seed separation vs per-index injectivity). Our
construction does not require the impossible condition that for every
fixed (j,\(\ell\)), F\_overlay(Seed\(_1\);j,\(\ell\)) \(\neq\)
F\_overlay(Seed\(_2\);j,\(\ell\)) for all distinct seeds
Seed\(_1\)\(\neq\)Seed\(_2\) (which would contradict finite pool size).
What is required - and what we prove - is: What we require (summary): \textbullet{}
Injective seed\(\to\)permutation: different seeds induce different
\(\pi\)\_v (via injective key derivation from Seed\_v to round keys; see
Appendix A.1.1). \textbullet{} Pool confinement and bijection (per seed): for each
fixed seed, \(\pi\)\_v is a bijection L\_v\(\to\)U\_v, so designated
addresses lie in U\_v and are unique within v (Unique-Neighbor) while
pools \{U\_v\} are disjoint across nodes (no cross-node aliasing). \textbullet{}
Full-churn separation on designated sets: for any two seeds, the
canonical designated set S(P) moves by a constant fraction, i.e.,
\textbar{}\(\pi\)\_v\^{}(seed)(S(P)) \(\Delta\)
\(\pi\)\_v\^{}(seed\textquotesingle)(S(P))\textbar{} \(\geq\)
c·\textbar{}S(P)\textbar{} (Lemma A.1.F; §A.1.1, §A.9). Together with
Keyedness (artifacts must be seed-consistent), these imply "collision
\(\Rightarrow\) error": using a mismatched seed misaddresses and breaks
correctness (see §4.4, §6.2.7 note on Keyedness; Appendix C.1.1). This
is exactly the property used in the SCL/SMP product factoring and the
OAP bypass arguments (Lemma J.1-Cart; Step 3 of §10.4.1-BYP).

\textbf{Lemma 6.2.3-CUT (Cut-Independence via Disjoint Pools - Forward
Reference).}

Under the disjoint designated address pools U\_v (§6.2.3 above) and
injective Enc (§6.2.7), feasible worlds across any cut C in the DAG
satisfy the Cartesian product bound:

\textbar{}\(\Pi\)\_C(\(\pi\))\textbar{} = \(\prod\)\_\{v\(\in\)C\}
\textbar{}S\_v\textbar{}

where \(\Pi\)\_C(\(\pi\)) denotes feasible world choices across cut C
consistent with transcript prefix \(\pi\), and S\_v denotes the set of
locally feasible assignments to unresolved coordinates at node v that
produce distinct designated addresses on the current seed chain.

\textbf{Intuition}: Because address pools \{U\_v\}\emph{\{v\(\in\)V\}
are pairwise disjoint (U\_v \(\cap\) U\_v\textquotesingle{} =
\(\emptyset\) for v \(\neq\) v\textquotesingle{}) and
F\_overlay(Seed\_v; j, \(\ell\)) maps to U\_v, changing the local choice
at node v only affects addresses in pool U\_v, not in other pools.
Combined with Enc injectivity (different local outputs \(\to\) different
seeds \(\to\) different designated addresses) and no cross-coupling
across the cut (Frontier-Gate horizon ensures post-horizon digests
depend only on pre-horizon data; Lemma J.1-NC), this forces full
Cartesian product structure: every combination of local choices
\{s\_v\}}\{v\(\in\)C\} with s\_v \(\in\) S\_v corresponds to exactly one
distinct feasible world.

\textbf{Implications}: This product structure is crucial for cut
composition in the SCL proof (Theorem 7.2.1 Step 4; §7.2.1) and ensures
that residuals add across cuts: \(\lambda\)(C) =
\(\Sigma\)\_\{v\(\in\)C\}(R\_v - q\_v), yielding \(\Phi\)(C) \(\geq\)
\(\lambda\)(C) by taking logarithms of both sides of the product bound.

\emph{Full bijection proof with both injectivity and surjectivity}:

\begin{itemize}
\tightlist
\item
  Main text: Theorem 7.2.1 Step 4 (§7.2.1)
\item
  Appendix: Theorem J.1-PROD (Hypotheses H1-H5), Lemma J.1-Cart
  (Cartesian factoring), Lemma J.1-NC (no cross-coupling)
\end{itemize}

\subparagraph{6.2.4 Primitive checks (Tiny-AND/parity variant with AAS:
Anti-Algebraic
Salting)}\label{624-primitive-checks-tiny-andparity-variant-with-aas-anti-algebraic-salting}

Let Z\_v(w,x) \(\in\) \{0,1\}\^{}(m\(_0\)) be a fixed O(1)-length bit
vector derived from the base core (e.g., a few witness/base bits or
edge/tour indicators). For public coefficient rows a\_v,\_j,\(\ell\)
\(\in\) \{0,1\}\^{}(m\(_0\)), b\_v,\_j,\(\ell\) \(\in\)
\{0,1\}\^{}(s\_salt) \textbackslash{} \{0\}, define the
\textbf{primitive bit}:

e\_v,\_j,\(\ell\) := \ensuremath{\langle}a\_v,\_j,\(\ell\), Z\_v(w,x)\ensuremath{\rangle} \(\oplus\)
\ensuremath{\langle}b\_v,\_j,\(\ell\), \(\sigma\)\_\{u\_v,\_j,\(\ell\)\}\ensuremath{\rangle} \(\in\) \{0,1\},

where \ensuremath{\langle}·,·\ensuremath{\rangle} is inner product over \ensuremath{\mathbb{F}}\(_2\). The salts \(\sigma\)\_u
\(\in\) \{0,1\}\^{}(s\_salt) are pairwise-distinct published constants
(AAS) that prevent algebraic merging without resolution and ensure each
branch bit represents an \textbf{intermediate state} that must be
computed rather than inferred. This is constant-time to evaluate.

Definition (AAS). "Anti-Algebraic Salting" means salts are explicit
constants embedded in the instance to block algebraic inference without
designated reads; pairwise-distinctness across designated addresses
suffices for our arguments.

See Appendix A.2 (salts) and A.4 (constraints) for explicit choices.

\textbf{Important:} The "independence" here refers only to the R\_v
designated addresses selected by Sel\_v at each node - we use explicit,
pairwise-distinct s\_salt-bit constants. No actual randomness is
involved; independence is a property of the specific subset of addresses
accessed, not the entire salt pool. These can be deterministic bit
patterns (e.g., leading bits of well-known mathematical constants offset
by the address index). During instance construction, all salts are
published constants.

6.2.4.a Witness Distribution Policy (Bounded Coverage)

To make explicit the intended bounded-coverage behavior for core bits
used by primitives, the reduction adopts the following deterministic
policy for forming Z\_v(w,x):

Lemma 6.2.4.a (Witness distribution on bottleneck cuts; bounded reuse).
In Algorithm R, for the canonical bottleneck cut C* (the designated
min-cut), each witness bit w{[}i{]} appears in Z\_v(w,x) for at most c =
O(1) nodes v \(\in\) C*. Consequently, across C*, the total reuse of any
primitive input coordinate is bounded by a universal constant.

Explicit schedule (by construction). Let the nodes on C* be ordered
v\_1, \ldots{}, v\_k by their public ids. Fix m\(_0\) = O(1). Define a
public stride

s := max\{ m\(_0\), \(\lfloor\) n\_core / k \(\rfloor\) \}.

Set, for r = 1..k,

Z\_\{v\_r\}(w,x) := {[} w{[}(r\(-\)1)·s + t mod n\_core{]}
{]}\_\{t=0\}\^{}(m\(_0\)-1).

\begin{itemize}
\tightlist
\item
  Capacity/no-wrap. Under QP-sharp calibration (k = \(\Theta\)(log
  n\_core), m\(_0\) = O(1)), we have k·m\(_0\) \(\leq\) n\_core, hence s
  = \(\lfloor\)n\_core/k\(\rfloor\) \(\geq\) m\(_0\) and the windows do
  not overlap across C*. Therefore c = 1 (each witness bit appears in at
  most one Z\_\{v\_r\}).
\item
  General bound. For any profile, the stride s \(\geq\) m\(_0\)
  guarantees each index participates in at most one window across C*, so
  c \(\leq\) 1 holds universally. If desired, a simpler variant with
  stride s = m\(_0\) yields c \(\leq\) m\(_0\) = O(1).
\end{itemize}

The schedule is deterministic and instance-public; it preserves all
structural properties of §6 and does not affect OAP or FG wiring. (Other
equivalent constant-reuse schedules are acceptable.)

Stack them as a column vector

e\_v := (e\_v,\_j,\(\ell\))\_\{(j,\(\ell\)) \(\in\)
J\_v×{[}\(\kappa\){]}\} \(\in\) \{0,1\}\^{}(K\_v\(\kappa\)).

\subparagraph{6.2.5 Branch bits (linear aggregation of
primitives)}\label{625-branch-bits-linear-aggregation-of-primitives}

Let L\_v := J\_v × {[}\(\kappa\){]} be the primitive-index pool. Choose
a public selector matrix Sel\_v \(\in\) \{0,1\}\^{}(R\_v×K\_v\(\kappa\))
of \textbf{full row rank R\_v} with \textbf{unit row-weight} (each row
has exactly one 1); different rows select different coordinates of e\_v.
This fixes the \textbf{selected index subset} L\_v\^{}(sel) =
\{(j,\(\ell\)) \(\in\) L\_v selected by Sel\_v\},
\textbar{}L\_v\^{}(sel)\textbar{} = R\_v, and, via seed-keyed
addressing, the \textbf{designated address subset} Addr\_v =
\{F\_\{overlay\}(Seed\_v; j,\(\ell\)) : (j,\(\ell\)) \(\in\)
L\_v\^{}(sel)\} \(\subseteq\) U\_v, \textbar{}Addr\_v\textbar{} = R\_v,
so the R\_v selected primitives read R\_v \textbf{distinct} salts.
Define the node\textquotesingle{}s \textbf{forced micro-facts vector}

x\_v := Sel\_v e\_v (mod 2) \(\in\) \{0,1\}\^{}(R\_v).

(If desired, you may also define per-branch scalars x\_v,\_j by grouping
rows of Sel\_v; for the invariants we only need the length-R\_v vector
x\_v.)

\subparagraph{6.2.6 Completeness (node
resolution)}\label{626-completeness-node-resolution}

We present two gate choices, both achieving full rank R\_v:

\textbf{Standard gate (used for EO/DP/Resolution).} Choose H\_v =
I\_\{R\_v\} (the identity matrix), so

y\_v := H\_vx\_v = x\_v \(\in\) \{0,1\}\^{}(R\_v).

\textbf{Order-robust gate (for OBDD).} For stronger OBDD bounds, use an
expander parity gate: fix a d-regular expander Exp\_v on R\_v vertices
with matching decomposition M\(_1\),...,M\_d, and define y\_v as the
edge parities on a \textbf{constant number} of disjoint matchings (e.g.,
M\(_1\)\(\cup\)M\(_2\)). This keeps size \(\Theta\)(R\_v), preserves
order-robust width blow-up for \textbf{every} variable order, and yields
\textbf{rank(H\_v)=R\_v} (details in Appendix B).

Both choices satisfy rank(H\_v) = R\_v, ensuring the Completeness
invariant (Lemma 6.2).

Construction note (parity gate rank). The expander-parity gate uses
\(\Theta\)(R\_v) parity equations drawn from matching edges, but it is
constructed so that exactly R\_v rows are linearly independent. Thus
rank(H\_v)=R\_v holds exactly (no \(\Theta\)(·) slack), ensuring precise
emergence accounting needed for SCL.

\subparagraph{6.2.7 Public addressing \& dependency
(DAG)}\label{627-public-addressing--dependency-dag}

\begin{itemize}
\tightlist
\item
  \textbf{Overlay addressing.} A public function F\_overlay(Seed\_v;
  j,\(\ell\)) maps the node seed Seed\_v and pair (j,\(\ell\)) to the
  designated address u\_v,\_j,\(\ell\) \(\in\) U\_v.
\end{itemize}

\textbf{Remark (naming).} A \textbf{node} is the phase v together with
its data (Seed\_v, U\_v, Sel\_v, H\_v, F\_overlay). The subsets
L\_v\^{}(sel) and A\_v are \emph{induced} by that phase; the node is not
itself a subset.

\begin{itemize}
\tightlist
\item
  \textbf{Seed (DAG).} For node v with parents P(v), define
\end{itemize}

Seed\_v := Enc(v \textbar{}\textbar{}
sort(\{(u,Seed\_u,y\_u):u\(\in\)P(v)\}) \textbar{}\textbar{}
GateDigest\_v),

\textbf{Fan-in bound (complexity hygiene).} We fix max in-degree
\(\Delta\)\_in := O(1) for G. This ensures \textbar{}P(v)\textbar{}
\(\leq\) \(\Delta\)\_in so seed parsing/encoding and per-node
verification remain O(log² n) (see §10.1).

where Enc is a public, efficiently computable injective encoding (e.g.,
length-delimited concatenation). This content-addressed seed models
\textbf{Dependency}: addresses at v are undefined until its
parents\textquotesingle{} y\_u are fixed; addresses are derived from
Seed\_v via F\_overlay.

\textbf{Lemma 6.2.7 (Seed size bound; parseable Enc).} With max
in-degree \(\Delta\)\_in=O(1) and canonical, length-delimited,
domain-separated Enc, for any node v we have \textbar{}Seed\_v\textbar{}
= O(depth(v) · log n\_core). Proof sketch. Each Enc tuple appends: (i)
node id v (O(log n\_core) bits); (ii) \(\leq\)\(\Delta\)\_in parent
triples (u,Seed\_u,y\_u) where u is O(log n\_core) and y\_u has length
R\_u=O(log n\_core) under our sizing; and (iii) GateDigest\_v, a
canonical record with a path id plus a parity bit over S(P(v)), both
O(log n\_core). All fields are length-delimited, so parsing is
unambiguous and linear in size. With \(\Delta\)\_in=O(1) and
depth(v)=O(log n\_core) in the balanced profiles,
\textbar{}Seed\_v\textbar{} = O(log² n\_core); along a solution path of
O(log n\_core) nodes, total seed bits remain O(log² n\_core), keeping
verification polynomial (see §10.1; §A.1.1 on seed-derived round keys).

\subparagraph{6.2.8 Gate-Ledger (Mandatory FG
Wiring)}\label{628-gate-ledger-mandatory-fg-wiring}

Intuition: The ledger binds each step of progress to priced designated
work on the current seed chain, preventing front-loading and enabling
single-run time pricing from Alt to steps.

To make the Frontier-Gate requirement structural and verifier-checkable,
we wire a \textbf{gate-ledger digest} into seeds beyond a published
horizon.

\textbullet{} \textbf{Published gate horizon.} The instance publishes a Boolean map
GREQ\_v \(\in\) \{0,1\} indicating whether node v requires a gate
digest. The map is chosen so that along every root\(\to\)sink path
hitting the bottleneck cut C* (i.e., containing nodes from C*), from a
published horizon onward, every node has GREQ=1. (A canonical choice:
all nodes in the last L\_tail layers intersecting C*, with L\_tail =
\(\Theta\)(n/W\_min) for QP-sharp.)

Calibration remark: With this calibration, each gate-required node v
beyond the horizon has a canonical path P(v) whose designated set
satisfies \textbar{}S(P(v))\textbar{} = \(\Theta\)(n/W\_min), so
producing GateDigest\_v entails evaluating and XORing
\(\Theta\)(n/W\_min) seed-dependent terms on the current seed chain (cf.
Lemma 5.5.1.c).

\textbullet{} \textbf{Canonical path assignment.} Each gate-required node v is
assigned a canonical path P(v) \(\in\) \ensuremath{\mathcal{P}} (Appendix C.1.1) hitting C*
(containing nodes from C*), published as PathOf(v).

\textbullet{} \textbf{Gate digest (binding).} For GREQ\_v=1, define

GateDigest\_v := Enc\_gate(P(v) \textbar{}\textbar{}
XOR\_\{(v\textquotesingle{},(j,\(\ell\)))\(\in\)S(P(v))\}
e\_\{v\textquotesingle{},j,\(\ell\)\}).

(Enc\_gate is a canonical, length-delimited encoder for gate digests;
see §6.2.9 below.)

This XOR is evaluated on the current seed chain (designated addresses
depend on Seeds). For GREQ\_v=0, GateDigest\_v := \(\varepsilon\)
(empty).

\textbf{Architectural note (Discriminator vs Hardness).} The XOR
computes a 1-bit parity over R emergent bits. This parity serves as a
\emph{discriminator}, not the hardness source:

\begin{itemize}
\tightlist
\item
  \textbf{Discriminator role}: Parity witnesses that two configurations
  differ. If parity(cfg\(_1\)) \(\neq\) parity(cfg\(_2\)), then
  cfg\(_1\) \(\neq\) cfg\(_2\). Flipping any single unobserved bit flips
  parity, so incomplete observation always permits finding two configs
  with different parities.
\item
  \textbf{Hardness source}: The 2\^{}R lower bound comes from A2
  injectivity on the FULL R-bit emergent vector: different configs
  \(\to\) different seeds \(\to\) different addresses.
\item
  \textbf{Proof chain}: incomplete obs \(\to\) parity differs
  {[}discriminator{]} \(\to\) configs differ \(\to\) seeds differ
  {[}A2{]} \(\to\) must explore 2\^{}R configs {[}hardness{]}.
\end{itemize}

If parity alone were the hardness source, there would be only 2
possibilities. The R-bit architecture with A2 injectivity enforces
2\^{}R exploration.

Pre-horizon-only inputs. The published index sets S(P(v)) are drawn
exclusively from pre-horizon nodes (those with GREQ=0) along the
canonical paths. Consequently, GateDigest\_v depends only on pre-horizon
designated primitives and never on unresolved coordinates on the
bottleneck cut (see Lemma 7.2.1-NC; Appendix C.1.1).

Unbypassability of GREQ=1 (decode and witness). \textbullet{} Horizon parameter. The
horizon is explicit: choose L\_tail = \(\Theta\)(n/W\_min) tail layers
intersecting C*; all nodes beyond this horizon have GREQ=1. This is
published and verifier-checkable. \textbullet{} Seed-locked decode uses gated seeds.
The seed-locked decode schema \(\Phi\)\textasciitilde{} stores (or derives) its mask
bits exclusively at designated addresses keyed by seeds with GREQ=1
along canonical paths hitting C* (containing nodes from C*). Since
Seed\_v := Enc(v \textbar{}\textbar{} ... \textbar{}\textbar{}
GateDigest\_v) includes GateDigest\_v whenever GREQ\_v=1, decoding
\(\varphi\) via \(\Phi\)\textasciitilde{} necessitates computing those GateDigest\_v
values on the current seed chain. Hence \(\varphi\) cannot be fully
decoded from pre-horizon (GREQ=0) data alone. \textbullet{} Canonical witness
requires all published gate paths. The canonical witness includes
G\_\(\tau\) = \{(P,S(P)) : P \(\in\) \ensuremath{\mathcal{P}}\} and the corresponding digest
bits Dig\_\(\tau\). The verifier checks that the
witness\textquotesingle{}s G\_\(\tau\) equals the published canonical
set (coverage/completeness) and that each digest bit equals the
recomputed parity over S(P) on the current seed chain (see §4.4,
§C.1.1). Thus even given w, a valid canonical witness cannot be produced
without engaging the GREQ=1 nodes beyond the horizon. \textbullet{} Mass of gated
information (counting). Let m be the number of clauses and b :=
\(\lceil\)log\(_2\) n\_core\(\rceil\)+1 the literal-encoding length.
\(\Phi\)\textasciitilde{} one-time-pads all 3m·b literal bits via mask bits
R{[}i,p{]}{[}t{]} that are placed at designated addresses keyed by
GREQ=1 seeds along canonical paths P(i) (see §10.1.1; Appendix C.1.1 for
path construction). Thus \(\Theta\)(m·log n\_core) independent mask bits
reside beyond the horizon; without evaluating these gated seeds, literal
identities remain hidden and \(\varphi\) cannot be recovered from GREQ=0
data. \textbullet{} Decode sub-pools are disjoint. The decode sub-pools
U\_v\^{}(dec) used for \(\Phi\)\textasciitilde{}\textquotesingle s mask bits are
dedicated slices disjoint from the gate-digest pools used for S(P) in
path gates (see §10.1.1). This prevents interference between decode
storage and gate computations and keeps verifier checks modular.

\textbullet{} Information-theoretic independence (structure reveals \(\approx\)0
about \(\varphi\)). Let X\_struct denote the public overlay metadata and
decode schema (G, \{Sel\_v\}, \{H\_v\}, Enc, F\_overlay, GREQ, \{S(P)\},
\(\Phi\)\textasciitilde{}), excluding the decode payload values. Under the sampler \ensuremath{\mathcal{S}}
(§9.2), X\_struct is drawn independently of \(\varphi\) and seed
computation is w-independent (Lemma 10.1.1-Adj). Hence H(\(\varphi\))
\(-\) O(log n) \(\leq\) H(\(\varphi\) \textbar{} X\_struct) \(\leq\)
H(\(\varphi\)). Equivalently, the overlay-as-problem (OAP) structure is
information-theoretically hiding with respect to \(\varphi\); recovering
\(\varphi\) requires the gated decode payload bits beyond the horizon.

\textbf{Lemma 10.1.IT (OAP independence bound).} Under the packaging
distribution \ensuremath{\mathcal{S}} where X\_struct \ensuremath{\perp} \(\varphi\) and seed computation is
w-independent, we have H(\(\varphi\)) \(-\) O(log n) \(\leq\)
H(\(\varphi\) \textbar{} X\_struct) \(\leq\) H(\(\varphi\)). In
particular, any mutual information I(\(\varphi\); X\_struct) is at most
O(log n) from explicit length fields. Thus any nontrivial information
about \(\varphi\) resides in the gated payload bits, not in the public
overlay structure.

\emph{Proof.} By independence, for any measurable S of structures and
\(\Phi\) of formulas, Pr{[}\(\varphi\)\(\in\)\(\Phi\) \textbar{}
X\_struct\(\in\)S{]} = Pr{[}\(\varphi\)\(\in\)\(\Phi\){]}. Entropy
equality follows, with an additive O(log n) for sizes/lengths explicitly
encoded in X\_struct. \(\blacksquare\)

Lemma (Gate inputs are seed-bound; not a function of (w,v) alone). Let v
be beyond the horizon with GREQ\_v=1 and path P=P(v). The digest
GateDigest\_v is the XOR of primitives e\_\{u,j,\(\ell\)\} at addresses
u\_\{u,j,\(\ell\)\}=F\_overlay(Seed\_u; j,\(\ell\)) with u on P and
GREQ\_u=0 (acyclicity). Although the primitives lie at pre-horizon
nodes, their designated addresses depend on the current seed chain
through the seeds at those nodes. Since seeds on P depend on the sorted
parent (Seed,y) tuples (Enc) and, for post-horizon seeds, on GateDigest
values, the addresses in S(P) are bound to the traversal history.
Therefore, GateDigest\_v cannot be computed from (w,v) alone without
reconstructing the seed chain; different seed histories induce different
address selections (constant-fraction movement; §A.1.1), invalidating
any post-hoc synthesis that ignores the chain.

\begin{quote}
\textbf{Acyclicity Lemma.} S(P) includes only indices from nodes with
GREQ=0 (pre-horizon), so GateDigest\_v never appears in any Seed\_u that
is an ancestor of v. \emph{Proof.} By construction (Appendix C.1.1), for
each gated node v the set S(P(v)) is chosen from pre-horizon nodes
(those with GREQ=0). Therefore GateDigest\_v is a function only of
pre-horizon primitives and does not depend on any gated
node\textquotesingle{}s digest. Since Seed\_u for any ancestor u encodes
only sorted parent tuples and (when applicable) its own GateDigest\_u,
and no ancestor\textquotesingle{}s GateDigest depends on any
descendant\textquotesingle{}s digest, dependencies follow the
DAG\textquotesingle{}s topological order. Hence GateDigest\_v cannot
occur in any ancestor seed. \(\blacksquare\)
\end{quote}

\textbullet{} \textbf{Optional prefix-keyed receipts (analysis-only).} Fix
k=\(\lfloor\)\(\tau\)R\_v\(\rfloor\) and let PSeed\_v\^{}(k) := Enc(v
\textbar{}\textbar{} sort(\{(u,Seed\_u,y\_u)\}\_\{u\(\in\)P(v)\})
\textbar{}\textbar{} y\_v{[}1..k{]}). One may derive receipts from
PSeed\_v\^{}(k) to audit that designated accesses correspond to the
current seed chain. These receipts are an optional analytical
instrument: they are not part of the language or witness, and the
verifier does not require or check them. Our time lower bounds do not
rely on receipts.

\textbullet{} \textbf{Seeds carry the digest.} Seed\_v includes GateDigest\_v (see
§6.2.8 for construction). Thus computing addresses at descendants
requires producing GateDigest\_v. By Lemma 5.5.1.c, computing
GateDigest\_v requires evaluating and XORing \textbar{}S(P(v))\textbar{}
= \(\Theta\)(n/W\_min) seed-dependent terms; this is a computational
cost that persists even if salts were pre-scanned and cached.

\textbullet{} \textbf{Seed binding.} Since designated addresses for S(P(v)) depend
on the current seed chain, any change in unresolved bits changes Seeds
and hence the designated addresses - invalidating prior digests. This is
the essence of FG-4 (no cross-seed sharing).

\textbullet{} \textbf{Verifier check.} The NP verifier recomputes Seeds and
GateDigest\_v for each v with GREQ\_v=1 using the explicit schema in
Appendix C.1.1; any mismatch leads to rejection. All checks run in
polynomial time.

\textbf{Remark (Gate overlaps and accounting).} The lower bounds do not
assume disjointness of S(P(v)) across different gate-required nodes v.
Our single-run bound only needs that each non-accepting segment that
progresses the final chain performs at least one priced event: a
\textbf{fresh} GateDigest on the current seed chain, a \textbf{growth of
NF\_C} (ConstraintDigest change), or a \textbf{refutation of
WorldCommit\_C}. (NF\_C is the canonical normal form summarizing
cut-level constraints; see §C.2.a for the full definition and verifier
contract.) When multiple priced events occur along the same chain, any
overlaps in designated checks only reduce total cost; the bifurcation
analysis uses the existence of a priced event per segment to lower-bound
time by \(\Omega\)(n/W\_min) and does not rely on additivity across
events.

\subparagraph{\texorpdfstring{6.2.9 Gate Horizon Budget (Structural
\(\tau\)-Cap)}{6.2.9 Gate Horizon Budget (Structural \textbackslash tau-Cap)}}\label{629-gate-horizon-budget-structural-ux3c4-cap}

To make the single-run bound tight unconditionally, the instance
publishes a \textbf{gate horizon budget} that structurally caps
pre-final agreement across the bottleneck cut.

\textbf{BOX - \(\mu\)-\(\tau\) FG Summary (single-run TMs).} Let
\(\lambda\)\_base := \(\Sigma\)\_\{v\(\in\)C*\} R\_v. With
FG\textquotesingle{}s global allowance \(\tau\)(n) and local allowance
\(\mu\):

\begin{itemize}
\tightlist
\item
  \textbf{Pre-final agreement:} s \(\leq\)
  \(\Theta\)(\(\tau\)(n)·\(\lambda\)\_base) + \(\mu\)·\(\lambda\)\_base.
\item
  \textbf{Segments:} m\_seg \(\geq\) 2\^{}(\(\rho\)-s) with \(\rho\)
  \(\geq\) \(\lambda\)\_base \(-\) s.
\item
  \textbf{Per-segment cost:} \(\Omega\)(n/W\_min) TM steps (priced
  event: gate digest, NF\_C growth, or committed-world refutation).
  Hence T(n) \(\geq\) 2\^{}(\(\rho\)-s)·\(\Omega\)(n/W\_min) =
  2\^{}((1-\(\Theta\)(\(\tau\)+\(\mu\)))\(\lambda\)\_base) /
  poly(n\_core). For \textbf{QP-sharp}
  (\(\lambda\)\_base=\(\Theta\)(log² n), \(\tau\)(n)=\(\Theta\)(log
  n/\(\lambda\)\_base)): T(n) = n\^{}(\(\Theta\)(log n\_core)) (tight up
  to poly factors).
\end{itemize}

Note (asymptotics). The \(\tau\) budget and each
\textbar{}S(P)\textbar{} are polynomially bounded in
\textbar{}x*\textbar{} by construction, so the cost to list or check the
\(\tau\) items themselves is polynomial. The super-polynomial lower
bound arises from the product of (i) the forced number of non-accepting
rollback segments m\_seg \(\geq\) 2\^{}(\(\rho\)-s) (segment counting)
and (ii) the per-segment baseline \(\Omega\)(n/W\_min) beyond the gate
horizon - not from \(\tau\) alone.

\textbf{Example (Computing \(\lambda\)\_base).} Consider a toy
bottleneck cut C* = \{v\(_1\), v\(_2\), v\(_3\)\} with emergence ranks R
= (6, 4, 5). Then:

\begin{itemize}
\tightlist
\item
  \(\lambda\)\_base = \(\Sigma\)\_\{v\(\in\)C*\} R\_v = 6 + 4 + 5 = 15
  (purely structural)
\item
  For QP-sharp profile: set \(\tau\)(n) = \(\Theta\)(log n /
  \(\lambda\)\_base) = \(\Theta\)(log n / 15)
\item
  Gate horizon budget: S(n) = \(\Theta\)(\(\tau\) · \(\lambda\)\_base) =
  \(\Theta\)(log n\_core)
\item
  GREQ assignment: Choose GREQ\_v so that \(\Sigma\)\_\{v\(\in\)C*,
  GREQ\_v=0\} R\_v \(\leq\) S(n)
\end{itemize}

This structurally enforces s \(\leq\) B(n) where s is pre-final
agreement, yielding single-run bound m\_seg \(\geq\)
2\^{}((1-\(\tau\))\(\lambda\)\_base) segments via SC (Appendix C.2).

\textbullet{} \textbf{Global budget (\(\tau\)-cap).} Let C* be the designated
bottleneck cut with base residual \(\lambda\)\_base =
\(\Sigma\)\_\{v\(\in\)C*\} R\_v. The published GREQ map is chosen so
that across C* \(\Sigma\)\_\{v\(\in\)C*, GREQ\_v=0\} R\_v \(\leq\) S(n)
:= \(\Theta\)(\(\tau\)(n) · \(\lambda\)\_base(n)), where \(\tau\)(n) =
\(\Theta\)(log n/\(\lambda\)\_base(n)) for QP-sharp, and similarly
calibrated for other profiles. Thus, at most S(n) cut bits can be
resolved before encountering gate-required nodes along any
root\(\to\)sink path through C*.

\textbullet{} \textbf{Local allowance (\(\mu\)).} Optionally, for each node v we
designate a local allowance \(\mu\)·R\_v of bits attributable before the
gate horizon; this is realized by partitioning selector rows and
publishing which rows are eligible pre-horizon. (This does not affect NP
verification cost.)

Cross-reference (residual after \(\tau\)-cap). By Lemma C.2.1, the
gate-horizon budget implies the run-specific residual obeys \(\rho\)
\(\geq\) \(\lambda\)\_base \(-\) s, where s \(\leq\) B(n) is the
pre-final agreement across C*. Combining with Segment Counting (Appendix
C.2) yields m\_seg \(\geq\) 2\^{}(\(\rho\)-s) \(\geq\)
2\^{}((1-o(1))\(\rho\))/poly(n\_core) under QP-sharp calibration.

\textbullet{} \textbf{Calibration of \(\mu\) (with \(\tau\)).} Since pre-horizon
bits attributable per node are \(\leq\) \(\mu\)·R\_v, across C* we have
s \(\leq\) S(n) + \(\mu\)·\(\sum\)\_\{v\(\in\)C*\}R\_v =
\(\Theta\)(\(\tau\)·\(\lambda\)\_base) + \(\mu\)·\(\lambda\)\_base. With
\(\rho\) \(\geq\) \(\lambda\)\_base \(-\) s (Lemma C.2.1), s \(\leq\)
(\(\tau\)+\(\mu\))·\(\lambda\)\_base \(\leq\)
((\(\tau\)+\(\mu\))/(1\(-\)(\(\tau\)+\(\mu\))))·\(\rho\). Choosing
\(\tau\)+\(\mu\) \(\leq\) 1/2 (e.g., \(\mu\) \(\leq\) \(\tau\)/4 with
\(\tau\) \(\leq\) 1/2) gives s \(\leq\) 2(\(\tau\)+\(\mu\))·\(\rho\) and
thus m\_seg \(\geq\) 2\^{}(\(\rho\)-s) \(\geq\)
2\^{}((1-2(\(\tau\)+\(\mu\)))·\(\rho\)) = 2\^{}(\(\rho\))/poly(n\_core).
This shows \(\mu\) can be tuned alongside \(\tau\) without affecting
tightness.

Calibration examples.

\begin{itemize}
\tightlist
\item
  Exponential profile (\(\lambda\)\_base=\(\Theta\)(n\_core)). Choose
  constants \(\tau\)=1/8 and \(\mu\)=1/16. Then s \(\leq\)
  (\(\tau\)+\(\mu\))·\(\lambda\)\_base = (3/16)·\(\lambda\)\_base and
  \(\rho\)\(-\)s \(\geq\) (13/16)·\(\lambda\)\_base, so Segment Counting
  yields m\_seg \(\geq\) 2\^{}((13/16)·\(\lambda\)\_base) (exponential).
\item
  QP-sharp profile (\(\lambda\)\_base=\(\Theta\)(log² n\_core)). Set
  \(\tau\)(n)=c·log n\_core/\(\lambda\)\_base for a fixed c\(\in\)(0,1)
  and \(\mu\) \(\leq\) \(\tau\)/4. Then s \(\leq\) \(\Theta\)(log
  n\_core) and \(\rho\)\(-\)s \(\geq\)
  (1\(-\)c\(-\)o(1))·\(\lambda\)\_base, so m\_seg \(\geq\)
  2\^{}(\(\Omega\)(\(\lambda\)\_base)) = n\_core\^{}(\(\Omega\)(log
  n\_core)) (quasi-polynomial).
\end{itemize}

\textbullet{} \textbf{Gate count along the final chain (Case B context).} Beyond the
horizon, every node on any root\(\to\)sink path hitting C* (containing
nodes from C*) has GREQ=1. Hence any accepting final chain that
progresses beyond s must traverse \(\Theta\)(\(\rho\)/R\_avg)
gate-required nodes (R\_avg := average rank on C*). Producing t gate
proofs on that chain entails \(\Omega\)(t·\textbar S(P)\textbar) =
\(\Omega\)(t·n/W\_min) designated reads (read-or-x.md §6.2.8),
reinforcing the Case-B pricing of prior gate work.

\textbullet{} \textbf{Structural consequence.} For any run, the \textbf{pre-final
agreement} s (the number of distinct cut bits first revealed before the
last segment) is bounded by s \(\leq\) S(n) = \(\Theta\)(\(\tau\)(n) ·
\(\lambda\)\_base(n)). This follows because beyond the horizon all nodes
on C* have GREQ=1 and require gate digests bound to the current seed
chain for any progress.

This budget is purely \textbf{instance-side} (encoded by the published
GREQ map and selectors) and is fully verifiable in polynomial time by
the NP verifier.

\subparagraph{6.2.10 Gate (rank-forcing) - publication
note}\label{6210-gate-rank-forcing---publication-note}

Publish H\_v\(\in\)\{0,1\}\^{}(R\_v×R\_v) with rank(H\_v)=R\_v. Either
take H\_v=I (identity) or an \textbf{expander-parity} family
(order-robust for OBDD).

\emph{Rank witness.} When using expander-parity, include the public
seed/parameters that define the map; verifying full rank is then
straightforward from the published description and is verifiable in
polynomial time by the NP verifier.

All objects above are explicit and computable in time poly(n\_core);
their total size is polynomial (see Appendix A.7).

\textbf{Canonical Enc (content-addressed; Git/Nix style).} \texttt{Enc}
(aka \textbf{CA-Enc}) is a \textbf{canonical}, length-delimited,
domain-separated encoding; collisions are irrelevant to the proofs (we
only need injectivity), but a canonical format eliminates ambiguity and
matches real content-addressed systems. We similarly use a canonical,
length-delimited encoder \texttt{Enc\_gate} for gate digests.

\textbf{Construction Recipe (node v):}

\begin{enumerate}
\def\labelenumi{\arabic{enumi}.}
\tightlist
\item
  Choose \(\kappa\) = \(\Theta\)(log n\_core), R\_v per flat profile
\item
  Define a seed-dependent permutation \(\pi\)\_v(Seed\_v;·) over D\_v =
  K\_v\(\kappa\) cells and set designated addresses by u\_v,\_j,\(\ell\)
  := F\_overlay(Seed\_v; j,\(\ell\)) = \(\pi\)\_v(Seed\_v; j,\(\ell\))
\item
  Write pairwise-distinct published salt constants
  \{\(\sigma\)\_u\}\_\{u\(\in\)U\_v\} at those designated cells (AAS)
\item
  Define primitives e\_v,\_j,\(\ell\) = \ensuremath{\langle}a\_v,\_j,\(\ell\),Z\_v\ensuremath{\rangle} \(\oplus\)
  \ensuremath{\langle}b\_v,\_j,\(\ell\),\(\sigma\)\_\{u\_v,\_j,\(\ell\)\}\ensuremath{\rangle} (Tiny-AND/parity)
\item
  Pick selector Sel\_v of row rank R\_v with unit row weight; set x\_v =
  Sel\_ve\_v
\end{enumerate}

\textbf{Size bound (explicit).} Fix \(\kappa\) = \(\Theta\)(log n\_core)
and choose K\_v = \(\Theta\)(R\_v/log n) so that
\textbar{}L\_v\textbar{} = K\_v·\(\kappa\) = \(\Theta\)(R\_v) and
\textbar{}L\_v\^{}(sel)\textbar{} = R\_v. Each node needs R\_v salts
(one per selected primitive). With pairwise-disjoint salt pools and s =
\(\Theta\)(log n\_core) bits per salt, the overlay salt mass is
\(\Sigma\)\_v R\_v·s = \(\Theta\)(\(\Sigma\)\_v R\_v·log n) =
\(\Theta\)(n log² n) bits under the flat profile where \(\Sigma\)\_v
R\_v = \(\Theta\)(n log n). Structural metadata adds
\(\Theta\)(\(\Sigma\)\_v R\_v) = \(\Theta\)(n log n) bits. Publishing
GREQ\_v and PathOf(v) adds
O(\textbar{}V\textbar{}·log\textbar{}\ensuremath{\mathcal{P}}\textbar{}) = poly(n\_core) bits.
Hence \textbar{}x*\textbar{} = poly(\textbar{}x\textbar{}) under the
flat profile (\(\Sigma\)\_v R\_v = \(\Theta\)(n log n)) and our
parameter choices. 6. Completeness: H\_v = I\_R\_v so y\_v = x\_v (rank
forcing) 7. Serialize: Use the DAG seed formula: Seed\_v = Enc(v
\textbar{}\textbar{} sort(\{(u,Seed\_u,y\_u):u\(\in\)P(v)\})
\textbar{}\textbar{} GateDigest\_v) (GateDigest\_v empty when GREQ\_v=0)
For chains where P(succ(v))=\{v\}, this reduces to
Enc(...\textbar{}\textbar{}v\textbar{}\textbar{}Seed\_v\textbar{}\textbar{}y\_v),
equivalent to the simpler notation

This works for \textbf{CLIQUE/TSP/3-SAT}: Z\_v(w,x) pulls O(1) core bits
(literal/edge/tour indicators).

\textbf{§6.2 Summary:} Constructed per-node machinery: disjoint address
pools U\_v (A1: Hermeticity), injective seed encoding Seed\_v = Enc(...)
(A2: Injectivity), rank-R\_v selector/check matrices (A3: Emergence,
proved Lemma 6.1; Completeness proved Lemma 6.2), seed-dependent
addressing F\_overlay (Keyedness), DAG dependency structure (A5, §6.5),
and FG wiring (§6.2.8-6.2.9 for per-instance bounds). Together these
realize SCL q+\(\Phi\)\(\geq\)R at each node; invariants summarized in
§6.7.

\textbf{Lemma 6.H5 (No Cross-Coupling; Acyclicity).} For any published
gate path P(v) in the FG wiring, the selector index set S(P) references
only nodes and coordinates \textbf{strictly before the gate horizon}
(upstream in DAG order). Consequently, evaluating GateDigest\_v at a
node v on the designated bottleneck cut C* does not create functional
dependencies between unresolved coordinates across C* - it depends only
on pre-horizon data already determined by the seed chain.

\emph{Proof sketch (full in App. C.1.1/J).} By construction
(§6.2.8-§6.2.10), nodes with GREQ\_v=1 have published S(P) sets
constructed via acyclic path selection: each (u,(j,\(\ell\)))\(\in\)S(P)
satisfies u \(\in\) Anc(C*) (ancestors of C*) and (j,\(\ell\)) indexes
designated rows at pre-horizon positions. The published GREQ map ensures
no forward edges cross the horizon. By DAG structure (A5) and path
acyclicity, GateDigest\_v cannot couple unresolved cut coordinates with
each other - it only reads pre-resolved upstream values. This "no
cross-coupling" property (H5) is essential for cut composition (Appendix
J: Theorem J.1, Lemma J.1-Cart). \(\blacksquare\)

\begin{center}\rule{0.5\linewidth}{0.5pt}\end{center}

\paragraph{6.3 Emergence Lemma
(provable)}\label{63-emergence-lemma-provable}

\textbf{Lemma 6.1 (Emergence / Realizability).}

\textbf{Main Statement (Per-Instance Property). Before any node-v
discovery read (by RWA convention), every x\_v \(\in\) \{0,1\}\^{}(R\_v)
remains realizable} by the construction: for any target value, we can
choose the (unread) designated salts to produce that value. This
realizability property holds for \textbf{every constructed instance} and
is the foundation for per-instance deterministic bounds (Theorem 8.A).

Hypotheses (for readability): Sel\_v has full row rank and unit row
weight; the designated map (j,\(\ell\)) \(\mapsto\) u\_v,\_j,\(\ell\) is
injective within node v; pools \{U\_v\} are disjoint across nodes; salts
\(\sigma\)\_u are pairwise-distinct published constants (AAS).

\emph{Proof (Realizability).} We fixed Sel\_v with full row rank and
unit row-weight, where different rows select different coordinates of
e\_v. Combined with the injective addressing (j,\(\ell\)) \(\mapsto\)
u\_v,\_j,\(\ell\), the R\_v selected primitives read R\_v distinct salts.
For fixed w,x and coefficients, each selected primitive e\_v,\_j,\(\ell\)
is an affine linear form of its \textbf{unread} salt
\(\sigma\)\_\{u\_v,\_j,\(\ell\)\}. Before any discovery read, we can
choose these R\_v unread salts to produce any target x\_v. Full rank of
Sel\_v ensures the mapping is onto \{0,1\}\^{}(R\_v). \(\blacksquare\)

\textbf{Consequence for Per-Instance Bounds.} This realizability
property means L*\textquotesingle{}s structure doesn\textquotesingle{}t
privilege any particular x\_v value before reading - all 2\^{}(R\_v)
possibilities are structurally valid. Therefore, any algorithm that
doesn\textquotesingle{}t read designated salts faces 2\^{}(R\_v)
genuinely distinct cases (Keyedness + Injectivity \(\to\) different x\_v
lead to different seeds \(\to\) different addresses \(\to\)
verifier-detectable). This holds for \textbf{every instance}, not just
distributional expectations.

\textbf{Twin-Input Flip (witness).} If some designated address
u\(\in\)U\_v has \textbf{not} been read before node v, then there exist
two inputs that (i) agree on the entire transcript and all previously
read cells, (ii) differ only at \(\sigma\)\_u, and (iii) induce
\textbf{different} y\_v. Hence a solver that does not read u either errs
on one input or must maintain both distinguishable artifacts (counted by
Alt\_v).

\textbf{Lemma 6.1-ZMI (Emergence via Zero Mutual Information -
strengthened).} Let Tr\_\{\textless{}v\} be any transcript that contains
\textbf{no discovery read} for node v (by the Receiving-Window
Attribution (RWA) convention; prefetch before node v is still charged to
node v and hence excluded from Tr\_\{\textless{}v\}). Then:

\textbf{Main Statement (Per-Instance).} Realizability: \textbar{}\{x\_v
consistent with Tr\_\{\textless{}v\}\}\textbar{} = 2\^{}(R\_v). Before
discovery reads, all 2\^{}(R\_v) values of x\_v remain structurally
valid.

\textbf{Lemma 6.1-RWA (Order-invariance of resolutions under RWA).}
Under the Receiving-Window Attribution convention, the per-node
resolution count q\_v depends only on which node-v facts are
functionally determined by the transcript, not on the wall-clock order
or timing of reads. In particular, moving reads of node-v salts earlier
in time, or precomputing functions of them before the node boundary,
does not reduce q\_v: their first use is still charged to node v.

\emph{Proof (semantic).} Let \ensuremath{\mathbb{F}}\_v be the \(\sigma\)-algebra generated by
the observations that the run ever uses and that depend on node-v salts.
By definition of q\_v, it is the number of independent node-v bits of
x\_v determined by \ensuremath{\mathbb{F}}\_v. Reordering the same observations in time does
not change \ensuremath{\mathbb{F}}\_v, hence does not change q\_v. By RWA, the first use of
any node-v dependent observation is attributed to node v regardless of
when the raw bytes were read, so prefetch cannot reduce q\_v.
Observations that are never used generate no resolutions.
\(\blacksquare\)

\textbf{Lemma 6.1-RZ (Realizability).} For every node v and every
transcript Tr\_\{\textless{}v\} that has not read any node-v salts,
\textbf{all values x\_v \(\in\) \{0,1\}\^{}(R\_v) are realizable} within
the construction.

\emph{Proof.} The node-v coordinates of x\_v are deterministic functions
of the designated salts. For any target u \(\in\) \{0,1\}\^{}(R\_v),
choosing the (unread) salts accordingly (consistently with
Tr\_\{\textless{}v\}) yields x\_v = u. \(\blacksquare\)

\textbf{Language vs. analysis (clarity).} The objects
(Sel\_v,H\_v,Enc,salts,selectors,expander) are \textbf{part of the
instance} and are checked by the NP verifier. The rules "Receiving-Node
Charging," "Keyedness," and the artifact counters
(tries/branches/keys/width/subfunctions/degree) are \textbf{analysis
conventions} to measure search complexity; they do \textbf{not} affect
acceptance.

\subparagraph{\texorpdfstring{6.3.1 Why Having the Input \(\neq\)
Knowing the
Intermediates}{6.3.1 Why Having the Input \textbackslash neq Knowing the Intermediates}}\label{631-why-having-the-input--knowing-the-intermediates}

The location of node v\textquotesingle{}s data depends on Seed\_v =
Enc(v \textbar{}\textbar{} sort(\{(u,Seed\_u,y\_u) : u \(\in\) P(v)\})
\textbar{}\textbar{} GateDigest\_v), which requires computing parent
outputs first. The input contains all salts, but \textbf{which salts
belong to node v} is determined by F\_overlay(Seed\_v; ·) - the data is
content-addressed by the computation itself.

This creates three operational barriers at each node with
2\^{}(R\_v-q\_v) live worlds (mathematically the constraint remains
two-dimensional in (q, \(\Phi\))):

\begin{enumerate}
\def\labelenumi{\arabic{enumi}.}
\tightlist
\item
  \textbf{Resolution dimension} (read/learn q\_v bits) \(\to\) forward
  learning, reduces to 2\^{}(R\_v-q\_v) possibilities
\item
  \textbf{Elimination dimension} (test candidates) \(\to\) backward
  pruning via up to 2\^{}(R\_v-q\_v) attempts
\item
  \textbf{Storage dimension} (maintain states) \(\to\) parallel tracking
  requires artifact size \(\geq\) 2\^{}(R\_v-q\_v)
\end{enumerate}

The Emergence lemma (realizability: all 2\^{}(R\_v) values of x\_v
remain structurally valid before discovery reads; Lemma 6.1-RZ)
formalizes this: until observations depend on x\_v, you have
2\^{}(R\_v-q\_v) indistinguishable worlds.

\subparagraph{\texorpdfstring{6.3.2 UN-Selector coverage \(\to\)
disjoint designated
atoms}{6.3.2 UN-Selector coverage \textbackslash to disjoint designated atoms}}\label{632-un-selector-coverage--disjoint-designated-atoms}

We record the structural disjointness that prevents "one read helps many
obligations".

\textbf{Lemma 6.5 (Disjoint designated atoms).} Because (j,\(\ell\))
\(\mapsto\) u\_v,\_j,\(\ell\) is injective and nodes use disjoint ranges
U\_v, the designated addresses are \textbf{pairwise disjoint}. Thus the
entire family already satisfies the disjointness property (coverage
parameter C=1: a single read can satisfy at most one obligation in this
subfamily).

\emph{Proof.} Immediate from the injective mapping \(\pi\)\_v defined in
\textbf{Appendix A.1} and per-node disjointness of U\_v.
\(\blacksquare\)

\textbf{Corollary (No batching across the disjoint subfamily).} Any
single read/derivation can merge distinguishable artifacts for at most
\textbf{one} obligation in that subfamily.

\textbf{Lemma 6.6 (Instance Completability).} For any partial
specification of salts (corresponding to reads made during any
computational run), we can consistently assign values to all remaining
unread salts to produce a complete, valid instance x* \(\in\) L*. The
construction guarantees:

\begin{enumerate}
\def\labelenumi{\arabic{enumi}.}
\tightlist
\item
  Each node\textquotesingle{}s selected primitives touch distinct salts
  (Lemma 6.5)
\item
  Node pools U\_v are pairwise disjoint (§6.2)
\item
  H\_v has full rank (Completeness, Lemma 6.2)
\end{enumerate}

Therefore, such a completion always exists and the resulting instance x*
has length poly(n\_core).

\emph{Proof:} For unread positions, choose x\_v arbitrarily and compute
y\_v = H\_v x\_v. Injectivity of Enc ensures descendant seeds remain
consistent. The flat sizing ensures \textbar{}x*\textbar{} =
poly(n\_core). \ensuremath{\square}

\textbf{Consequence.} No computational run can "paint itself into a
corner" - whatever salts have been read, the remaining salts can always
be assigned consistently. This property enables per-instance analysis
where every completion is valid.

\begin{center}\rule{0.5\linewidth}{0.5pt}\end{center}

\paragraph{6.4 Completeness (rank
forcing)}\label{64-completeness-rank-forcing}

\textbf{Lemma 6.2 (Rank forcing).} Let y\_v = H\_vx\_v with rank(H\_v) =
R\_v. Any correct procedure that outputs the true y\_v for \textbf{all}
possible x\_v must \textbf{fix} (learn) at least R\_v independent
coordinates of x\_v.

\emph{Proof.} If fewer than R\_v independent coordinates of x\_v are
determined by the transcript, there exist two distinct possibilities
x\_v \(\neq\) x\textquotesingle{}\_v consistent with the transcript.
Since H\_v has full row rank, H\_vx\_v \(\neq\)
H\_vx\textquotesingle{}\_v. Thus no procedure can output the true y\_v
on both possibilities without reading more; contradiction.
\(\blacksquare\)

\textbf{Key:} Since Seed\_\{child\} depends on all coordinates of y\_v,
any uncertainty forces distinguishable artifacts - exactly what SCL
requires.

\textbf{Completeness (verifier).} The verifier recomputes primitives,
forms x\_v = Sel\_v e\_v, computes y\_v = H\_vx\_v, then Seed\_v = Enc(v
\textbar{}\textbar{} sort(\{(u,Seed\_u,y\_u):u\(\in\)P(v)\})
\textbar{}\textbar{} GateDigest\_v). We do not publish y\_v; the
verifier derives it. Polynomial time overall.

\textbf{Soundness.} Any error in primitives or y\_v causes immediate
verification failure.

\begin{center}\rule{0.5\linewidth}{0.5pt}\end{center}

\paragraph{6.5 Dependency (no
look-ahead)}\label{65-dependency-no-look-ahead}

\textbf{Observation (Dependency established by construction).} Before
y\_v is fixed, the node seed Seed\_v = Enc(v \textbar{}\textbar{}
sort(\{(u, Seed\_u, y\_u) : u \(\in\) P(v)\}) \textbar{}\textbar{}
GateDigest\_v) and thus all addresses F\_overlay(Seed\_v; ·) are
\textbf{undefined}. (GateDigest\_v is empty when GREQ\_v=0 and, when
present, ranges only over pre-horizon indices - acyclic; see §6.2.8.)
Any solver that accesses designated cells at v must therefore have
\textbf{resolved} parent outputs and - when GREQ\_v=1 - produced the
required gate digest on the current seed chain.

\emph{Note.} This follows immediately from the definition of Seed\_v as
a function of ancestor outputs - it is how the construction establishes
dependency.

\textbf{Merge and multi-parent nodes.} If a solver wishes to "coalesce"
two distinguishable artifacts (states), it must realize a \textbf{new}
node whose seed is derived from \textbf{both parents} (see §4.4).
Coalescing without such a step is ruled out by Keyedness.

GREQ semantics reference. When GREQ\_v=1, computing GateDigest\_v is
required at v before Seed\_v is defined (see §6.2.8-§6.2.9 for
GREQ/GateDigest semantics).

\begin{center}\rule{0.5\linewidth}{0.5pt}\end{center}

\paragraph{6.6 Sizing: flat NP variant
(expository)}\label{66-sizing-flat-np-variant-expository}

We instantiate a \textbf{flat} profile variant that keeps verification
polynomial while forcing exponential complexity for \emph{search}
without a witness. Our main quasi-polynomial profile (QP-sharp) is
specified in §4.3.

\begin{itemize}
\tightlist
\item
  \textbf{DAG depth.} depth(G) = O(log n\_core).
\item
  \textbf{Emergence budget.} \(\Sigma\)\_v\(\in\)V R\_v = \(\Theta\)(n
  log n). With O(log n\_core) layers, this gives \(\Theta\)(n) rank per
  layer, ensuring min-cut \(\lambda\)\_base = \(\Theta\)(n).
\item
  \textbf{Micro-arity.} \(\kappa\) = \(\Theta\)(log n\_core) so that the
  permutation/selector can pick R\_v disjoint primitive positions.
\item
  \textbf{Salt footprint.} Fix D\_v = K\_v\(\kappa\) = \(\Theta\)(R\_v)
  by choosing K\_v = \(\Theta\)(R\_v/\(\kappa\)). With \(\kappa\) =
  \(\Theta\)(log n\_core), the total salt budget across nodes is
  \(\Sigma\)\_v D\_v = \(\Sigma\)\_v \(\Theta\)(R\_v) = \(\Theta\)(n log
  n).
\item
  \textbf{Completeness sparsity.} Choose Sel\_v and H\_v with O(1) row
  weight (explicit LDPC/MDS-style families); both are public and
  polynomial size. Since we use H\_v = I\_R\_v in the main construction,
  this is trivially satisfied.
\item
  \textbf{Metadata size.} Permutation descriptions, selector seeds,
  Seed\_v, H\_v, addressing F\_overlay, and Enc are all poly(n\_core).
  Salt ranges U\_v occupy \(\Theta\)(R\_v) words each; overall input
  blow-up is polynomial.
\end{itemize}

Under this profile, a one-path verifier resolves all R\_v at each node
and runs in time poly(n\_core), while the bounds in §5 yield
\textbf{exponential} artifacts for EO/DP/OBDD/Resolution when the solver
avoids resolution on many nodes.

\textbf{Why exponential:} Shallow-wide DAG (O(log n\_core) depth,
\(\Theta\)(n) rank per layer) yields \(\lambda\)\_base = \(\Theta\)(n)
\(\to\) 2\^{}\(\Theta\)(n) complexity. Deep DAGs would give only
polynomial bounds.

\begin{center}\rule{0.5\linewidth}{0.5pt}\end{center}

\paragraph{6.7 Instance invariants
summary}\label{67-instance-invariants-summary}

\begin{itemize}
\tightlist
\item
  \textbf{Hermeticity (A1):} Designated pools \{U\_v\} pairwise
  disjoint; no hidden channels (§6.2.3; Lemma 6.5)
\item
  \textbf{Injectivity (A2):} Different histories \(\to\) different seeds
  (Enc injective) (§6.2.7)
\item
  \textbf{Emergence (A3):} All 2\^{}(R\_v) values of x\_v remain
  realizable before discovery reads - fresh coordinates cannot be
  derived without designated reads (Lemma 6.1; §6.3)
\item
  \textbf{Closure (A4):} Seeds deterministically recover ancestors (Enc
  parseable) (§6.2.7; Appendix L)
\item
  \textbf{Dependency (A5):} Seed\_v = Enc(v \textbar{}\textbar{}
  sort(\{(u,Seed\_u,y\_u): u \(\in\) P(v)\}) \textbar{}\textbar{}
  GateDigest\_v) requires parent completion (and gate digest when
  GREQ\_v=1) (§6.5)
\item
  \textbf{Completeness:} rank(H\_v) = R\_v forces complete computation
  (Lemma 6.2; §6.4)
\item
  \textbf{Early-read irrelevance:} RWA charges at first valid use (Lemma
  5.5.1.b)
\item
  \textbf{Designated reads requirement:} Each complete try reads
  \(\geq\) \(\Sigma\)\_v R\_v designated salts (Appendix C.4)
\end{itemize}

\textbf{Combined: With a valid witness, all R\_v bits are resolved
(\(\lambda\) = 0) \(\to\) polynomial verification. Without a witness,
maintaining correctness requires tracking 2\^{}(R\_v-q\_v)
distinguishable states at each unresolved node. These properties hold
for every instance}, enabling per-instance deterministic bounds (Theorem
8.A).

\begin{center}\rule{0.5\linewidth}{0.5pt}\end{center}

\paragraph{6.8 Gadget Glossary (Quick
Reference)}\label{68-gadget-glossary-quick-reference}

\begin{itemize}
\tightlist
\item
  Seed\_v: Content-addressed seed Seed\_v = Enc(v \textbar{}\textbar{}
  sort(\{(u,Seed\_u,y\_u):u\(\in\)P(v)\}) \textbar{}\textbar{}
  GateDigest\_v); binds addresses to parent outputs and (optionally)
  digest. See §6.2.7.
\item
  Enc: Public injective, parseable encoder used for seeds (and digests);
  enables Closure/Recoverability. See §6.2.7; App. L.
\item
  F\_overlay: Public addressing u\_\{v,j,\(\ell\)\} =
  F\_overlay(Seed\_v; j,\(\ell\)); a seed-dependent permutation
  \(\pi\)\_v with Unique-Neighbor within U\_v that makes designated
  reads seed-dependent. See §6.2.3/§6.2.7.
\item
  PathOf(v): The canonical path P(v) assigned to gate-required nodes v
  for computing GateDigest\_v. See §6.2.8; Appendix C.1.1.
\item
  S(P): Published path-gate index set on canonical path P;
  \textbar{}S(P)\textbar{} = \(\Theta\)(n/W\_min). See Appendix A.9;
  Appendix C.1.1.
\item
  U\_v: Disjoint address pool per node; prevents cross-node overlap. See
  §6.2.3.
\item
  \(\pi\)\_v: Explicit seed-dependent permutation with full-churn
  avalanche; drives address-churn. See App. A.1.1; Lemma A.1.\(\Delta\).
\item
  Sel\_v: Full-rank, unit-row-weight selector; fixes R\_v selected
  primitives. See §6.2.5.
\item
  L\_v\^{}(sel): Selected index set for primitives at v;
  \textbar{}L\_v\^{}(sel)\textbar{} = R\_v. See §6.2.5.
\item
  H\_v: Check/Completeness matrix (rank R\_v); standard (identity) or
  expander-parity choice. See §6.2.6.
\item
  Z\_v(w,x): O(1)-bit core fragment from base instance; used in
  primitives. See §6.2.4.
\item
  \(\sigma\)\_u (salts): Published constants in U\_v; Anti-Algebraic
  Salting (AAS) blocks inference without reads. See §6.2.4; App. A.2.
\item
  GateDigest\_v: Parity/XOR over the designated set S(P(v)) along the
  current seed chain; binds progress to designated work. See
  §6.2.8-§6.2.9.
\item
  GREQ\_v: Gate-requirement flag at node v; when 1, computing
  GateDigest\_v is required before Seed\_v is defined. See
  §6.2.8-§6.2.9.
\item
  PathOf(v) / P(v): Canonical path index set used to define gate digests
  at v. See Appendix A.9 / C.1.1.
\item
  e\_\{v,j,\(\ell\)\}: Primitive bit \ensuremath{\langle}a, Z\_v\ensuremath{\rangle} \(\oplus\) \ensuremath{\langle}b,
  \(\sigma\)\_\{u\}\ensuremath{\rangle}; constant-time check unit. See §6.2.4.
\item
  x\_v: Branch/intermediate vector x\_v = Sel\_v e\_v of length R\_v.
  See §6.2.5.
\item
  y\_v: Node output y\_v = H\_v x\_v feeding into children. See §6.2.6.
\item
  GREQ\_v: Published gate-requirement flag; marks nodes requiring
  GateDigest. See §6.2.8.
\item
  PathOf(v): Canonical path assignment P(v) \(\in\) \ensuremath{\mathcal{P}} for gate checks.
  See §6.2.8; App. C.1.1.
\item
  S(P): Published index set for path P with XOR-cancellation;
  \textbar{}S(P)\textbar{} = \(\Theta\)(n/W\_min). See App. A.9; App.
  C.1.1.
\item
  GateDigest\_v: Enc\_gate(P(v) \textbar{}\textbar{}
  \(\oplus\)\emph{\{(v\textquotesingle{},(j,\(\ell\)))\(\in\)S(P(v))\}
  e}\{v\textquotesingle{},j,\(\ell\)\}); wired into Seed\_v when
  GREQ\_v=1. See §6.2.8; App. C.1.1.
\item
  Frontier-Gate (FG): Instance-side throttle (\(\mu\),\(\tau\)
  allowances; gated progress) ensuring tight single-run bounds. See App.
  C.1.
\item
  Keyedness: Seed-consistency of artifacts; forbids cross-seed merges
  (collision \(\Rightarrow\) error). See §4.2; §7.2.1.
\item
  Hermeticity: No hidden channels; only designated reads count. See
  §4.4; §10.
\item
  RWA: Receiving-Window Attribution; charges bits at first valid use.
  See §4.2; §5.5.1.
\item
  Address-churn: Constant-fraction change of S(P) addresses per seed
  change; forces re-work. See Lemma A.1.\(\Delta\).
\item
  Cut-gate proofs: (P,S(P)) parity objects certifying progress; weight
  \(\Theta\)(n/W\_min). See App. C.1.1.
\item
  Segment Counting (SC): m\_seg \(\geq\) 2\^{}(\(\rho\)-s) rollback
  segments in a single run. See App. C.2.
\item
  Min-cut residual \(\lambda\)(A,x): Run-dependent bottleneck
  \(\lambda\) = min\_C \(\Sigma\)\_\{v\(\in\)C\}(R\_v\(-\)q\_v). See
  §7.2.1; App. J.
\item
  Algorithm V (Verifier): Explicit polynomial-time verifier for L*. See
  §10.
\end{itemize}

\subsection{Part IV: Core Technical
Results}\label{part-iv-core-technical-results}

\subsubsection{7. The Semantic Conservation Law (SCL) and Semantic
Necessity}\label{7-the-semantic-conservation-law-scl-and-semantic-necessity}

Section 6 constructed L* and proved it satisfies axioms \textbf{A1-A5}:
Hermeticity, Injectivity, Emergence, Closure, and Dependency. These are
structural properties built into L* by design, independent of any
algorithm. Now comes the critical step: proving these properties
\emph{mathematically force} the Semantic Conservation Law.

\textbf{Purpose of Section 7:} This is the \textbf{core mathematical
proof} of the paper. We rigorously prove that *\emph{A1-A5 \(\to\) SCL*}
(Theorem 7.A, §7.2.1 Consolidated SCL Theorem). This establishes that q
+ \(\Phi\) \(\geq\) R is not an assumption or heuristic - it is a
mathematical necessity arising from L*\textquotesingle{}s structure.
Everything after this section (per-instance bounds, Structural OWF
construction, separation) follows by applying this proven conservation
law.

\textbf{Proof Approach.} The result (P \(\neq\) NP via Structural OWF
construction, §9-§10) uses \textbf{A1-A5} with per-instance
deterministic bounds (§8). Section 7 focuses on the A1-A5 framework with
per-run properties H1-H5.

\textbf{Roadmap:}

\begin{itemize}
\tightlist
\item
  \textbf{§7.0}: Abstract Setup \& Terminology Map - define semantic
  quantities precisely
\item
  \textbf{§7.1}: Framework Summary and Key Results - overview of main
  theorems
\item
  \textbf{§7.2}: \textbf{The Semantic Conservation Law: Formal
  Statement} (Theorem 7.A: q\_v + \(\Phi\)\_v \(\geq\) R\_v)

  \begin{itemize}
  \tightlist
  \item
    \textbf{§7.2.1}: Consolidated SCL Theorem - complete proof
    *\emph{A1-A5 \(\to\) SCL*} (the logical heart)
  \end{itemize}
\item
  \textbf{§7.3}: Extensions and Manifestations - SCL \(\to\) Time
  (Theorem 7.B), paradigm adapters
\item
  \textbf{§7.4-§7.8}: Dependency \& multiplicative growth, operational
  routes, paradigm instantiations, conclusion
\end{itemize}

\textbf{Why this section is critical:} Before §7, we had intuitions
(§1-§2), framework definitions (§4), paradigm manifestations (§5), and a
construction (§6). But we have not \emph{proven} the conservation law
holds. Section 7 is where hand-waving stops and rigorous proof begins.
After §7, we can confidently say: "For L*, the law q + \(\Phi\) \(\geq\)
R is mathematically necessary." Section 8 then proves every FG-wired
instance has per-instance deterministic hardness (Theorem 8.A),
establishing the foundation for Structural OWF construction (§9).

\textbf{Preservation note.} All SCL statements are made on the native
run (DAG semantics) using ledger-defined counters (RWA-credited q,
seed-consistent Alt, min-cut \(\lambda\)). Compiled models are used only
to connect these counters to standard measures (e.g.,
width/subfunctions) via the SNF wrapper; we do not assume simulations
preserve semantic structure by themselves.

\paragraph{7.0 Abstract Setup \& Terminology
Map}\label{70-abstract-setup--terminology-map}

\subparagraph{7.0.1 The Abstract
Framework}\label{701-the-abstract-framework}

Let G = (V, E) be a DAG representing computation\textquotesingle{}s
dependency structure.

\textbf{Core Quantities:}

\textbf{Key Semantic Quantities}

\begin{itemize}
\tightlist
\item
  \textbf{Information requirement (R\_v)}: Bits that must be determined
  at node v
\item
  \textbf{Resolved bits (q\_v)}: Bits learned at node v
  (Receiving-Window Attribution)
\item
  \textbf{Residual uncertainty (R\_v - q\_v)}: Unresolved information
  after processing v
\item
  \textbf{Distinguishable artifacts (Alt\_v)}: Distinct computational
  states at v
\item
  \textbf{Min-cut capacity (\(\lambda\)(A,x))}: min\_\{cuts C\}
  \(\Sigma\)\_\{v\(\in\)C\}(R\_v-q\_v) - Global bottleneck
  (run-dependent)
\end{itemize}

\textbf{Notation sanity (\(\rho\) vs \(\lambda\)).} We write
\textbf{\(\lambda\)(A,x)} for the run-dependent min-cut residual in
general, and \textbf{\(\rho\)} for the same quantity evaluated just
before the last segment in the single-run analysis (§7.3.9). When
context is clear we may write \(\lambda\) for brevity.

\subparagraph{7.0.2 The Conservation
Principle}\label{702-the-conservation-principle}

When computation faces R bits of uncertainty but only resolves q bits,
the remaining (R - q) bits must manifest as distinguishable artifacts.
If Alt \textless{} 2\^{}(R - q), distinct worlds collide \(\to\)
incorrect computation. Thus the SCL: \textbf{q + \(\Phi\) \(\geq\) R}
(where \(\Phi\) = log\(_2\)(Alt)). For the complete proof see §7.2.1
(Consolidated SCL Theorem).

\emph{(Metric note: \(\Phi\) is Hartley (Rényi-0) entropy over feasible
worlds; see §1.6 for Hartley vs Shannon and why SCL is zero-error.)}

RWA and the definition of q. RWA is the operational, order-invariant
implementation of the semantic "functional determination" definition of
q used throughout. In other words, we define q semantically, and we
attribute (and price) it in runs via RWA.

Why RWA is canonical (q attribution).

\begin{itemize}
\tightlist
\item
  Fresh information that shrinks feasible worlds can only enter through
  first valid uses of designated reads; later re-reads do not further
  refine the partition. RWA charges exactly those first-use bits.
\item
  Order-invariant: q\_v depends on which observations are ultimately
  used (\(\sigma\)-algebra \ensuremath{\mathbb{F}}\_v), not when they are read; see Lemma
  6.1-RWA. See also Theorem I.1 (ledger soundness \&
  schedule-invariance).
\item
  Model-neutral: RWA works for any schedule or caching; it is a
  representation-independent accounting of what was necessarily learned.
  Combined with Hermeticity, it aligns per-node q\_v with actual
  designated reads (Lemma 5.5.1.b).
\end{itemize}

\subparagraph{7.0.3 Mapping to L*
Construction}\label{703-mapping-to-l-construction}

\textbf{Abstract SCL to L* Mapping}

\begin{itemize}
\tightlist
\item
  \textbf{Information requirement R\_v}: L* Instance: Emergence rank;
  Note: Bits that emerge at node v
\item
  \textbf{Resolved bits q\_v}: L* Instance: Designated reads; Note:
  Receiving-Window Attribution (RWA)
\item
  \textbf{Dependencies}: L* Instance: Seed chains
  Enc(v\textbar{}\textbar{}parents\textbar{}\textbar{}...); Note:
  Enforced by construction
\item
  \textbf{Distinguishable artifacts}: L* Instance: Branches/keys/width;
  Note: Paradigm-specific
\end{itemize}

\paragraph{7.1 Framework Summary and Key
Results}\label{71-framework-summary-and-key-results}

\textbf{Core requirements:} (1) DAG dependency structure, (2)
Information emergence at nodes, (3) Receiving-Window Attribution, (4)
Distinguishable artifacts cannot merge without resolution.

\textbf{Key bounds:}

\begin{itemize}
\tightlist
\item
  \textbf{Per node:} q + \(\Phi\) \(\geq\) R where \(\Phi\) =
  log\(_2\)(Alt) (Theorem 7.A)
\item
  \textbf{Across cuts:} For any s-t cut C, \(\Sigma\)\_\{v\(\in\)C\}
  log\(_2\) Alt\_v \(\geq\) \(\Sigma\)\_\{v\(\in\)C\}(R\_v\(-\)q\_v);
  minimizing over C gives \(\lambda\)(A,x) = min\_C
  \(\Sigma\)\_\{v\(\in\)C\}(R\_v\(-\)q\_v), and for the minimizing cut
  C*, log\(_2\) Alt(C*) \(\geq\) \(\lambda\)(A,x) (Theorem J.1)
\end{itemize}

\textbf{Complexity:}

\begin{itemize}
\tightlist
\item
  \(\lambda\)(A,x) = \(\Theta\)(log² n) \(\to\) n\^{}(\(\Theta\)(log
  n\_core)) (quasi-polynomial)
\item
  \(\lambda\)(A,x) = \(\Theta\)(n) \(\to\) 2\^{}(\(\Theta\)(n))
  (exponential)
\end{itemize}

\textbf{Theorem 7.B (SCL \(\to\) TM Time; deterministic).} Notation:
\(\Delta\)(C*) := \(\Lambda\)(C*) \(-\) (Q(C*) + log\(_2\) Alt(C*)) is
the across-tries deficit on the bottleneck cut C*; Alt(C*) :=
\(\prod\)\emph{\{v\(\in\)C*\} Alt\_v. Fix any deterministic k-tape
Turing machine A with alphabet \(\Gamma\) and let B :=
k·\(\lceil\)log\(_2\)\textbar{}\(\Gamma\)\textbar{}\(\rceil\) (Lemma
5.5.1). For any input x to L*, let \(\lambda\)(A,x) := min\_C
\(\Sigma\)}\{v\(\in\)C\}(R\_v\(-\)q\_v) be the run-dependent min-cut
residual under RWA/Hermeticity/Keyedness (§4-§6). Exactly one of the
following lanes occurs, and each yields a time lower bound:

\begin{itemize}
\item
  (Restart/brute-force lane) If A does not reuse keyed state across
  resolution-prefix changes, then time(A,x) \(\geq\)
  2\^{}(\(\lambda\)(A,x) - log\(_2\) Alt(C*))/B.
  (Kraft\textquotesingle{}s inequality refines this to time(A,x)
  \(\geq\) 2\^{}(\(\Delta\)(C*))·\(\Delta\)(C*)/B.)
\item
  (Single-run lane with FG) If A persists keyed state, then with
  Frontier-Gate wiring (§6.2.8; App. C.1.1) and the gate-horizon budget
  (\(\tau\)-cap), the run partitions into rollback segments and
  time(A,x) \(\geq\) 2\^{}(\(\rho\)-s) · \(\Omega\)(n/W\_min), where
  \(\rho\) \(\geq\) \(\lambda\)\_base(x) \(-\) s and s \(\leq\)
  \(\Theta\)(\(\tau\)(n)·\(\lambda\)\_base(x)).
\end{itemize}

\textbf{Asymmetry (discovery vs execution).} These lanes exhibit
complementary pricing: restart counts \textbf{discovery attempts}
(exponentially many independent tries sampling distinct resolution
prefixes to find one that works, with each try costing \(\geq\)1/B),
while single-run counts \textbf{execution work} (many rollback segments
once committed to a seed chain, each costing \(\Omega\)(n/W\_min) to
compute gate digests). Restart is brute-force search: state forgotten
between tries prevents learning, forcing 2\^{}(\(\Delta\)(C*))
independent sampling attempts in expectation. This distinction arises
from L*\textquotesingle{}s gate horizon structure (§6.2.9): failed
restart tries may abort at pre-horizon nodes before expensive gates,
while persistent single-run segments must compute post-horizon gates on
each rollback. See Appendix C.4.2 for the detailed restart bound
derivation.

In particular, for QP-sharp profiles (\(\lambda\)\_base=\(\Theta\)(log²
n), W\_min=\(\Theta\)(log n\_core)) we get time \(\geq\)
n\^{}(\(\Omega\)(log n\_core)); for flat profiles
(\(\lambda\)\_base=\(\Theta\)(n), W\_min=\(\Theta\)(1)) we get time
\(\geq\) 2\^{}(\(\Omega\)(n)).

\emph{Proof.} Theorem J.1 together with Lemma J.1-NC yield, for any run,
the cut inequality Q(C)+log\(_2\) Alt(C) \(\geq\) \(\Lambda\)(C) and
thus the residual \(\lambda\)(A,x) = min\_C
\(\Sigma\)\_\{v\(\in\)C\}(R\_v\(-\)q\_v). Lane analysis (Appendix C)
applies.

Restart/brute-force lane: See "Quick Restart Bound" (Appendix C.4.2) for
the self-contained derivation via across-tries inequality (Lemma 7.R),
RWA, and per-step inflow.

Single-run lane: By Segment Counting (Appendix C.2), the number of
non-accepting rollback segments that progress the final chain satisfies
m\_seg \(\geq\) 2\^{}(\(\rho\)-s), where \(\rho\) is the effective
residual and s the pre-final agreement on the designated bottleneck cut.
By the FG per-segment baseline (Appendix C.1.1) and TM digest cost
(Lemma 5.5.1.c) together with address-churn (Lemma A.1.\(\Delta\)), each
such segment requires \(\Omega\)(n/W\_min) TM steps. Acceptance
uniqueness (\textbar{}\ensuremath{\mathcal{U}}\textbar{}\(\leq\)1 at the start of the accepting
segment) follows from digest binding (Lemma C.2.ACC) and also as a
logical consequence (Lemma C.2.ACC-logical). Multiplying completes the
bound. \ensuremath{\square}

Remark (exhaustiveness and hybrids). Lane exhaustiveness is proved in
Appendix C; see Lemma C.EXH (Lane Exhaustiveness). Hybrid strategies are
tracked by the hybrid potential \ensuremath{\mathcal{H}}\_v with q\_v+\ensuremath{\mathcal{H}}\_v \(\geq\) R\_v (Lemma
7.3.6.7).

\textbf{Hypotheses Used (A1-A5).}

\begin{itemize}
\tightlist
\item
  \textbf{A1 (Hermeticity)}: All instance information enters via
  designated payloads on the read-only input (no hidden channels).
\item
  \textbf{A2 (Injectivity)}: Distinct (Seed,y) histories induce distinct
  designated addresses; merging non-equivalent worlds is
  verifier-detectable (+ Keyedness).
\item
  \textbf{A3 (Emergence)}: At node v there are R\_v local coordinates
  that must be realized; feasible choices exist for each coordinate
  pattern (realizability).
\item
  \textbf{A4 (Closure)}: Seeds deterministically recover ancestors (Enc
  parseable); verifier can reconstruct seed chains.
\item
  \textbf{A5 (Dependency)}: Parents complete before children; seeds
  propagate via Enc(v \textbar{}\textbar{} sorted parent tuples
  \textbar{}\textbar{} GateDigest) respecting DAG topological order.
\item
  \textbf{RWA (Attribution mechanism)}: q counts first-use, legitimate
  resolutions; schedule-invariant and model-neutral.
\end{itemize}

For cut composition (Step 4 in §7.2.1), we use per-run hypotheses
\textbf{H1-H5} (disjoint pools, unique-neighbor, Enc injectivity,
realizability, no cross-coupling) detailed in §7.2.1.

\paragraph{7.2 The Semantic Conservation Law: Formal
Statement}\label{72-the-semantic-conservation-law-formal-statement}

\textbf{Lemma 7.I (Injectivity \(\Rightarrow\) Alt\_v lower bound).} For
any node v in an A1-A5 instance, if s\_v := R\_v \(-\) q\_v bits remain
unresolved at v, then Alt\_v \(\geq\) 2\^{}(s\_v).

\emph{Proof.} Step 1 (Unresolved worlds). By the definition of q\_v
(first-use, RWA-credited designated reads), there are exactly
2\^{}(s\_v) feasible worlds consistent with the current transcript
prefix at v.

Step 2 (Distinct seeds via Injectivity). By A2 (Injectivity), the Enc
tuple over sorted parent (Seed,y) values is injective. Distinct
assignments to the s\_v unresolved local coordinates induce distinct Enc
tuples and hence distinct Seed\_v values across the 2\^{}(s\_v) worlds.

Step 3 (Distinct addresses via Keyedness). Keyedness asserts that
designated addresses are computed as F\_overlay(Seed\_v; j,\(\ell\)), so
if Seed\_v differs then for some designated index (j,\(\ell\)) we must
have F\_overlay(Seed\_v; j,\(\ell\)) \(\neq\) F\_overlay(Seed'\_v;
j,\(\ell\)). Thus the 2\^{}(s\_v) worlds induce pairwise distinct
designated access patterns on U\_v.

Step 4 (Hermetic distinguishability). By A1 (Hermeticity), all
legitimate information enters only through designated payloads on
disjoint pools U\_v. Therefore, two worlds that differ on designated
addresses or designated payload values yield different memory
projections on U\_v and are simultaneously distinguishable at v
(operational definition below).

Step 5 (Counting artifacts). The 2\^{}(s\_v) worlds map to 2\^{}(s\_v)
distinct projections on U\_v, hence Alt\_v \(\geq\) 2\^{}(s\_v) and
log\(_2\) Alt\_v \(\geq\) R\_v \(-\) q\_v. \(\blacksquare\)

\emph{(See §6.2.3 for disjoint pools and §6.2.7 for Seed/Enc.)}

\textbf{Falsifiability Note:} Breaking Cartesian factorization across
cuts (showing \textbar{}\(\Pi\)\_C\textbar{} \(\ll\)
\(\prod\)\_\{v\(\in\)C\}\textbar S\_v\textbar{} via witness-coupling
compression) would falsify this theorem. §3.6 addresses this attack
vector via two-level factorization---Cartesian product holds given
witness w; witness and overlay spaces must both be searched.

\textbf{Definition (Seed-Consistent Equivalence Classes).} Two worlds
\(\omega\) \textasciitilde{} \(\omega\)' at node v if they induce
identical future designated addresses and outcomes on the seed chain
through v. Define Alt\_v :=
\textbar{}\(\Pi\)\_v/\textasciitilde\textbar{} as the number of
equivalence classes. By A2 (Injectivity) + Keyedness, a correct
algorithm must realize at least Alt\_v disjoint computational states;
otherwise collision between inequivalent worlds causes errors.

\textbf{Definition 7.2.1.1 (Distinguishable Artifacts).} Two worlds
\(\omega\)\_i, \(\omega\)\_j are simultaneously distinguishable at v iff
\(\exists\) designated address \(\ell\) \(\in\) U\_v such that
Memory\_i{[}\(\ell\){]} \(\neq\) Memory\_j{[}\(\ell\){]} when computed
on the current seed chain. Define Alt\_v := \textbar{}\{memory
projections on U\_v over feasible worlds\}\textbar{}. Under Keyedness,
correctness requires these classes be represented by disjoint
computational states.

\textbf{Proposition 7.2.1.1 (Equivalence of definitions).} On L*, the
operational notion above coincides with seed-consistent equivalence:
\(\omega\) \textasciitilde{} \(\omega\)' iff their memory projections on
U\_v match on all designated reads along the seed chain through v. In
particular, Alt\_v defined operationally equals Alt\_v defined via
seed-consistent classes.

\textbf{Theorem 7.K-Addr (F\_overlay seed-injectivity).} For any L*
instance satisfying A1-A5, the overlay address map F\_overlay(Seed\_v;
j,\(\ell\)) is injective in the Seed\_v argument on the designated pools
U\_v required for correctness.

\emph{Proof (structural necessity).} Suppose not: \(\exists\) Seed
\(\neq\) Seed' with F\_overlay(Seed; j,\(\ell\)) = F\_overlay(Seed';
j,\(\ell\)) = addr for all designated indices needed at v. Then either
(i) both worlds read the same addr with the same value, collapsing
distinct unresolved worlds without resolution (contradicts A3 Emergence
and Lemma 7.I); or (ii) an overwrite semantics is required to separate
worlds (contradicts A1 Hermeticity: read-only, write-disjoint pools); or
(iii) addresses coincide but downstream behavior diverges on some
designated index, re-establishing seed-injectivity by definition. In all
cases, non-injectivity is incompatible with the A1-A5 framework's
correctness obligations. Therefore F\_overlay must be injective in its
seed argument on the designated domain. \(\blacksquare\) \emph{Proof.}
(\(\Rightarrow\)) If projections differ at some designated address,
future outcomes or addresses differ, so \(\omega\) \ensuremath{\not\approx} \(\omega\)'.
(\(\Leftarrow\)) If \(\omega\) \ensuremath{\not\approx} \(\omega\)' under seed-consistency,
then by Keyedness+Injectivity there exists a designated index
(j,\(\ell\)) whose address or outcome diverges; thus memory projections
differ at that \(\ell\). Representation invariance (Lemma 5.3.1) yields
equality of the counts. \(\blacksquare\)

\textbf{Theorem 7.A (Semantic Conservation Law - Abstract):} For any
computation on a DAG G = (V, E) that correctly distinguishes between all
possible outcomes:

\textbf{Universal form:} q + \(\Phi\) \(\geq\) R (where \(\Phi\) =
log\(_2\)(Alt))

This constraint holds at every node v, where:

\begin{itemize}
\tightlist
\item
  q\_v = bits of information resolved at node v (Receiving-Window
  Attribution)
\item
  Alt\_v = number of distinguishable computational artifacts maintained
  at v
\item
  R\_v = information requirement at node v
\end{itemize}

\textbf{Proof:} By Lemma 7.I, with s\_v=R\_v\(-\)q\_v unresolved local
coordinates at v, there are at least 2\^{}(s\_v) future-distinguishable
classes; any sufficient artifact state must realize at least that many
distinct images. Hence Alt\_v \(\geq\) 2\^{}(R\_v-q\_v) and
log\(_2\)(Alt\_v) \(\geq\) R\_v\(-\)q\_v. \(\blacksquare\)

\textbf{For a complete, structured proof:} See §7.2.1 (Consolidated SCL
Theorem) which explicitly shows the derivation: L* construction \(\to\)
\textbf{A1-A5} \(\to\) Distinguishability \(\to\) SCL.

\textbf{Machine-verified proof:} Lean formalization in
\texttt{SCLNode.lean} (theorem \texttt{SCL\_node}), with cut composition
in \texttt{SCLCut.lean} (theorem \texttt{SCL\_cut}). Both theorems are
proven with zero axioms, using only constructive counting via
\texttt{Fintype.card} instances.

\textbf{Dependencies (all proved earlier):} Information-theoretic
counting plus Injectivity (Lemma 7.I); Alt\_v counts
seed-consistent/future-distinguishable classes (Def. above). Keyedness
is used to argue that any correct algorithm must realize at least Alt\_v
disjoint states (no cross-seed merges). See Appendix I for
sufficient-statistic formalization.

\textbf{Theorem 7.A.1 (Semantic Conservation Law - Across Cuts):} For
any cut C through the DAG:

\(\Sigma\)\_\{v\(\in\)C\} log\(_2\)(Alt\_v) \(\geq\) \(\lambda\)(C),
where \(\lambda\)(C) = \(\Sigma\)\_\{v\(\in\)C\}(R\_v - q\_v).

Equivalently, \(\Pi\)\_\{v\(\in\)C\} Alt\_v \(\geq\)
2\^{}(\(\lambda\)(C)).

Micro-example (two-node cut). If C = \{v\(_1\), v\(_2\)\} with
s\_\{v\(_1\)\} = s\_\{v\(_2\)\} = 1 (one unresolved bit each), then
Alt\_\{v\(_1\)\} \(\geq\) 2 and Alt\_\{v\(_2\)\} \(\geq\) 2, so \(\Pi\)
Alt\_v \(\geq\) 4 and \(\Sigma\) log\(_2\) Alt\_v \(\geq\) 2. This is
exactly \(\Sigma\)(R\_v\(-\)q\_v) for the cut.

\textbf{Proof:} For each v\(\in\)C, Theorem 7.A gives Alt\_v \(\geq\)
2\^{}(R\_v-q\_v). Multiplying over v\(\in\)C yields
\(\Pi\)\_\{v\(\in\)C\} Alt\_v \(\geq\) \(\Pi\)\_\{v\(\in\)C\}
2\^{}(R\_v-q\_v) = 2\^{}(\(\Sigma\)\_\{v\(\in\)C\}(R\_v-q\_v)) =
2\^{}(\(\lambda\)(C)). Taking log\(_2\) of both sides gives
\(\Sigma\)\_\{v\(\in\)C\} log\(_2\)(Alt\_v) \(\geq\) \(\lambda\)(C).
\(\blacksquare\)

\textbf{SCL \(\leftrightarrow\) SMP Equivalence:} The Semantic
Conservation Law is the logarithmic form of the Semantic Multiplication
Principle. Under \textbf{Keyedness} (no cross-seed merges), SMP with
Alt\_v counting within-try distinguishable artifacts yields SCL.
Conversely, exponentiating SCL gives the multiplicative refinement bound
of SMP. Both express the same fundamental constraint: unresolved
information forces multiplicative growth of distinguishable
possibilities.

Notation clarity (s\_v vs s). We use s\_v := R\_v \(-\) q\_v to denote
the unresolved-bit count at a specific node v. Separately, we use s
(without subscript) for the pre-final agreement across the designated
bottleneck cut in the single-run analysis (Appendix C.2). These are
distinct quantities.

\textbf{Theorem 7.A0 (Minimal SCL - structural, accounting-free form).}
Hypotheses at a node v:

\begin{enumerate}
\def\labelenumi{\arabic{enumi})}
\tightlist
\item
  Hermeticity (instance-side): the only admissible observations are
  designated payload bits published in x* (no hidden channels).
\item
  Completeness (rank forcing): there are R\_v semantically necessary
  local bits (coordinates of x\_v) such that for any two assignments
  differing on these coordinates, there exist two feasible worlds
  consistent with the run prefix that realize them (realizability).
\item
  Injectivity + Keyedness: distinct assignments to the unresolved
  coordinates induce distinct future designated addresses/seed chains;
  merging non-equivalent worlds causes a verifier-detectable error.
\item
  Closure/Recoverability: the verifier can reconstruct from the sink
  transcript the ancestor (Seed,y) tuples so seed/address differences
  are observable.
\end{enumerate}

Conclusion: Let s\_v be the number of unresolved local coordinates at v
with respect to the current run prefix\textquotesingle{}s used
observations (functional determination, independent of any crediting
schedule). Any correct computation must realize at least 2\^{}(s\_v)
future-distinguishable artifact classes at v. Equivalently, defining
\(\Phi\)\_v := log\(_2\) Alt\_v, we have \(\Phi\)\_v \(\geq\) s\_v and
hence q\_v + \(\Phi\)\_v \(\geq\) R\_v with q\_v := R\_v \(-\) s\_v.

Proof (sketch). Under Completeness/realizability there exist 2\^{}(s\_v)
feasible worlds that agree with the prefix and differ exactly on the
unresolved local coordinates. Injectivity + Keyedness imply these worlds
induce distinct future seed/address traces; Closure makes these
differences verifier-observable. If the computation realized fewer than
2\^{}(s\_v) artifact classes, two such worlds would map to the same
artifact class, implying a later collision detectable by the verifier
(distinct designated addresses/outcomes). Hermeticity rules out external
information to pre-rule such worlds. Therefore Alt\_v \(\geq\)
2\^{}(s\_v). This argument does not appeal to any accounting rule; it
uses only correctness and instance-side structure. \(\blacksquare\)

Metric note. The "resolved" count q\_v here is defined via functional
determination with respect to the set of published payload bits actually
used in the run prefix, not via when those bits are read. Later, RWA is
employed only to price first-use bits in time (Appendix D), not to
derive SCL.

\textbf{Lemma 7.R (Across-tries inequality).} Fix any s-t cut C and a
deterministic run that may restart (or a randomized run with coins fixed
via Yao\textquotesingle{}s principle). Let the run decompose into tries
T\(_1\),...,T\_m, where a try ends exactly when keyed state is discarded
across a change of the resolution prefix. For try i, let q\_i(C) denote
the RWA-counted first-use bits across C inside that try, and let
Alt\_i(C) bound the maximum number of simultaneously distinguishable
artifacts across C within that try. Let \(\Lambda\)(C) :=
\(\Sigma\)\_\{v\(\in\)C\} R\_v.

Deterministic form: \(\Lambda\)(C) \(\leq\) \(\Sigma\)\_\{i=1\}\^{}m
q\_i(C) + \(\Sigma\)\_\{i=1\}\^{}m log\(_2\) Alt\_i(C).

In particular, if Alt\_i(C) \(\leq\) Alt(C) for all i (i.e., Alt(C) is
any uniform upper bound on the per-try artifact peak), then m \(\geq\)
(\(\Lambda\)(C) \(-\) \(\Sigma\)\_i q\_i(C)) / log\(_2\) Alt(C).

Expectation form: For any bound Alt(C) \(\geq\) Alt\_i(C), letting q(C)
:= \ensuremath{\mathbb{E}}{[}\(\Sigma\)\_i q\_i(C){]} and m := \ensuremath{\mathbb{E}}{[}\#tries{]},

q(C) + log\(_2\) Alt(C) + log\(_2\) m \(\geq\) \(\Lambda\)(C).

Equivalently, m \(\geq\) 2\^{}(\(\Lambda\)(C) \(-\) q(C) - log\(_2\)
Alt(C)).

\textbf{Dependencies Map (Where Each Property Is Used):}

\textbf{Property Usage in Core Proofs:}

\begin{itemize}
\tightlist
\item
  Lemma 7.I (Alt\_v lower bound): Uses A2 (Injectivity)

  \begin{itemize}
  \tightlist
  \item
    Distinct histories yield distinct seeds
  \end{itemize}
\item
  SCL \(\leftrightarrow\) SMP equivalence: Uses Keyedness

  \begin{itemize}
  \tightlist
  \item
    No cross-seed merges within tries
  \end{itemize}
\item
  Cut composition (Theorem J.1): Uses disjoint pools + Injectivity

  \begin{itemize}
  \tightlist
  \item
    Cartesian product structure (worst-case lane)
  \end{itemize}
\item
  Time bounds (§5.5, App C): Uses Dependency + RWA

  \begin{itemize}
  \tightlist
  \item
    Information flow through DAG
  \end{itemize}
\end{itemize}

\textbf{Semantic Necessity:} At node v, resolving q\_v bits reduces
2\^{}R\_v initial possibilities to 2\^{}(R\_v-q\_v) remaining worlds;
correctness forces at least 2\^{}(R\_v-q\_v) simultaneously
distinguishable artifacts unless those bits are resolved. Across any cut
C, requirements multiply to 2\^{}(\(\Sigma\)\_\{v\(\in\)C\}(R\_v-q\_v)).
This multiplication is unavoidable - fewer states means collision,
collision means wrong answers.

\textbf{Corollary 7.A.2 (Parametric Complexity):} The run-dependent
min-cut capacity \(\lambda\)(A,x) determines computational complexity:

\textbf{Capacity-Complexity Correspondence}

\begin{itemize}
\tightlist
\item
  \textbf{Capacity \(\lambda\)(A,x) = 0}: Artifacts Required: O(1);
  Typical Complexity: Polynomial (verification)
\item
  \textbf{Capacity \(\lambda\)(A,x) = O(log n)}: Artifacts Required:
  poly(n\_core); Typical Complexity: Polynomial
\item
  \textbf{Capacity \(\lambda\)(A,x) = \(\Theta\)(log² n)}: Artifacts
  Required: n\^{}\(\Theta\)(log n\_core); Typical Complexity:
  Quasi-polynomial
\item
  \textbf{Capacity \(\lambda\)(A,x) = \(\Theta\)(\ensuremath{\sqrt{n}})}: Artifacts
  Required: 2\^{}\(\Theta\)(\ensuremath{\sqrt{n}}); Typical Complexity: Sub-exponential
\item
  \textbf{Capacity \(\lambda\)(A,x) = \(\Theta\)(n)}: Artifacts
  Required: 2\^{}\(\Theta\)(n); Typical Complexity: Exponential
\end{itemize}

Semantic potential. A measure of computational uncertainty

\ensuremath{\Psi}\_v = log\(_2\)(\# of worlds consistent with observations at node v)

\textbf{How it evolves:}

\begin{itemize}
\tightlist
\item
  Node v: \ensuremath{\Psi}\_v bits of uncertainty
\item
  Resolve q\_v bits at node v
\item
  Successor nodes: uncertainty propagates through DAG
\end{itemize}

Why this is unavoidable. L*\textquotesingle{}s structure ensures these
2\^{}(R\_v-q\_v) states correspond to genuinely different computational
paths that lead to different outcomes. For Structural OWF construction
(§9), \textbf{every FG-wired instance} exhibits this property via
per-instance deterministic bounds (Theorem 8.A). Merging them would
cause incorrect results - hence the necessity.

\textbf{Key Insight:} Alt\_v counts \textbf{semantic artifacts}
(simultaneously distinguishable computational objects required for
correctness), not raw machine states or Shannon information.

\textbf{Frontier bound (DAG).} \textbf{FrontierPeak \(\geq\)
2\^{}(\(\lambda\)(A,x)) (by Theorem J.1}, there exists a cut across
which that many distinguishable artifacts must coexist; see also
Corollary 7.1.1).

\textbf{Receiving-Window Attribution (RWA).} We attribute resolutions to
the first node where they are used, independent of read order. See
Appendix I for one auditable implementation; the semantic bounds require
only the attribution principle, not a specific mechanism.

\textbf{DAG-wide bounds via min-cut.} For any s-t cut C in the DAG,
distinguishable artifacts across the cut satisfy
\(\Sigma\)\_\{v\(\in\)C\} log\(_2\)(Alt\_v) \(\geq\) \(\lambda\)(A,x)
where \(\lambda\)(A,x) = min\_C \(\Sigma\)\_\{v\(\in\)C\}(R\_v \(-\)
q\_v). We write Alt(C) := \(\prod\)\_\{v\(\in\)C\} Alt\_v for the total
across cut C; then log\(_2\) Alt(C) \(\geq\) \(\lambda\)(A,x). This
yields the tightest possible bound for DAG-structured computations. For
linear DAGs (chains), this reduces to the path sum \(\Sigma\)\_v(R\_v -
q\_v).

\textbf{Checkpoints (finality nodes; optional).} If a subset F of nodes
is declared \textbf{final} (must resolve fully, q\_v=R\_v), then
distinguishable artifacts multiply only across non-final nodes in cuts.

\begin{center}\rule{0.5\linewidth}{0.5pt}\end{center}

\paragraph{\texorpdfstring{7.2.1 Consolidated SCL Theorem (L*
construction \(\Rightarrow\) A1-A5 \(\Rightarrow\) Distinguishability
\(\Rightarrow\)
SCL)}{7.2.1 Consolidated SCL Theorem (L* construction \textbackslash Rightarrow A1-A5 \textbackslash Rightarrow Distinguishability \textbackslash Rightarrow SCL)}}\label{721-consolidated-scl-theorem-l-construction--a1-a5--distinguishability--scl}

\textbf{Theorem 7.2.1 (Consolidated SCL - L*).} Under A1-A5
(Hermeticity, Injectivity, Emergence, Closure/Recoverability,
Dependency) for the L* overlay, for any run of any solver on input x:

\begin{itemize}
\tightlist
\item
  Per-node: for every node v, q\_v + log\(_2\) Alt\_v \(\geq\) R\_v.
\item
  Across any s-t cut C: Q(C) + log\(_2\) Alt(C) \(\geq\) \(\Lambda\)(C),
  with Alt(C) := \(\prod\)\emph{\{v\(\in\)C\} Alt\_v, Q(C) :=
  \(\sum\)}\{v\(\in\)C\} q\_v, and \(\Lambda\)(C) :=
  \(\sum\)\_\{v\(\in\)C\} R\_v.
\end{itemize}

Equivalently, Alt\_v \(\geq\) 2\^{}(R\_v-q\_v) and Alt(C) \(\geq\)
2\^{}(\(\Lambda\)(C)-Q(C)). In particular, letting \(\lambda\)(A,x) :=
min\_C (\(\Lambda\)(C) \(-\) Q(C)) yields log\(_2\) Alt(C*) \(\geq\)
\(\lambda\)(A,x) for some cut C*.

\emph{Note:} The Structural OWF construction uses H1-H5 (including
disjoint-pool Cartesian factoring via H1) for per-instance deterministic
bounds.

Proof sketch: Structured below; see Steps 1-5.

Hypotheses used for cut factoring (H1-H5, worst-case lane).

\begin{itemize}
\tightlist
\item
  H1) Disjoint designated pools across the cut: U\_v \(\cap\) U\_\{v'\}
  = \(\emptyset\) for v \(\neq\) v' on C (no aliasing across cut nodes;
  §6.2.3).
\item
  H2) Unique-Neighbor mapping: designated addresses encode their pool id
  (v) and a key; an address identifies its unique cut node
  (instance-side property of F\_overlay; Appendix A.1).
\item
  H3) Enc injectivity + Closure/Recoverability: sink seeds encode a
  parsable chain of (Seed,y) tuples reconstructible by the verifier
  (Lemma J.0); different parent tuples yield different child seeds.
\item
  H4) Completeness/Realizability: for each v \(\in\) C every admissible
  local assignment s\_v \(\in\) S\_v extends to at least one feasible
  world consistent with the transcript prefix.
\item
  H5) No cross-coupling across the cut: post-horizon digests do not
  depend on unresolved coordinates on C and there are no instance-side
  constraints simultaneously restricting unresolved coordinates of
  distinct cut nodes (Lemma 7.2.1-NC; Lemma J.1-NC).
\end{itemize}

Under H1-H5, feasible worlds across any fixed cut C factor as a
Cartesian product (Appendix J: Theorem J.1/PROD), yielding the cut-level
SCL stated above.

\textbf{Lemma 7.2.1-NC (No cross-coupling across the cut).} In L* with
FG horizon, for any s-t cut C, every GateDigest wired into any Seed\_u
depends only on designated indices from pre-horizon nodes (GREQ=0).
Consequently, no post-horizon digest or other instance-side constraint
simultaneously restricts unresolved local coordinates from two distinct
nodes v \(\neq\) v' in C. Equivalently, the admissible local settings
\{S\_v\}\_\{v\(\in\)C\} are not coupled by instance structure.

\emph{Proof.} By §6.2.8 (gate horizon), GateDigest\_v includes only
indices S(P(v)) drawn from pre-horizon nodes (GREQ=0). Enc is parseable
and injective and depends only on parent (Seed,y) tuples and
GateDigest\_v; designated addresses carry pool-id = v and the pools
\{U\_v\} are disjoint (H1; §6.2.3). Therefore, digest values never
reference unresolved coordinates on C, and seed propagation introduces
no cross-node algebraic constraints beyond parent determinism. Appendix
J (Lemma J.1-NC) formalizes the no-cross-coupling condition used for
Cartesian factoring. Hence H5 holds. \(\blacksquare\)

Appendix Map (quick refs)

\begin{itemize}
\tightlist
\item
  Appendix J: Theorem J.1 (DAG min-cut/product), Lemma J.1-NC (no
  cross-coupling), Lemma J.1-Cart (Cartesian factoring), Lemma J.1-REAL
  (global realizability), Lemma J.0 (recoverability)
\item
  Appendix C: C.1.1 (Frontier-Gate paths; per-segment baseline), C.2
  (Segment Counting: m\_seg \(\geq\) 2\^{}(\(\rho\)-s)), C.4.2 (Quick
  Restart Bound)
\item
  Appendix I: Theorem I.1 (ledger soundness \& schedule-invariance),
  Lemma I.2.4 (monotonicity), Lemma I.2.5 (Frontier = classes)
\end{itemize}

\textbf{Setup.} Let x \(\in\) L* be an instance produced by the overlay
construction (§6). For every node v on the DAG:

\begin{itemize}
\item
  \textbf{Recall (Feasible world, §4.1)}: A total assignment consistent
  with the fixed transcript prefix and overlay constraints.
\item
  R\_v = required fresh bits (rank/requirement at v)
\item
  q\_v = \textbf{legitimate first-use} bits read at v (RWA accounting)
\item
  Define an equivalence relation on feasible worlds at v: \(\omega\)
  \(\approx\)\_v \(\omega\)' iff they induce the same \textbf{seed
  chain} and therefore the same designated addresses and outcomes
  downstream from v (Keyedness + Recoverability)
\item
  Alt\_v := \textbar{}\(\Pi\)\_v/\(\approx\)\_v\textbar{} and
  \(\Phi\)\_v := log\(_2\)(Alt\_v)

  \begin{quote}
  \textbf{Representation-invariance:} Alt\_v counts equivalence classes
  independent of internal encoding (see Lemma 4.2.2 for
  sufficient-statistic bound; Lemma 5.3.1 for full
  representation-invariant formulation)
  \end{quote}
\item
  For a cut C, write Q(C) := \(\Sigma\)\_\{v\(\in\)C\} q\_v,
  \(\Lambda\)(C) := \(\Sigma\)\_\{v\(\in\)C\} R\_v, \(\Phi\)(C) :=
  \(\Sigma\)\_\{v\(\in\)C\} \(\Phi\)\_v
\end{itemize}

\textbf{Notation.} \(\Lambda\)(C) := \(\Sigma\)\_\{v\(\in\)C\} R\_v
(total requirement); \(\lambda\)(C) := \(\Sigma\)\_\{v\(\in\)C\}(R\_v
\(-\) q\_v) (residual).

\textbf{Restatement (from Theorem 7.2.1).} For every node v,

q\_v + log\(_2\)(Alt\_v) \(\geq\) R\_v

Consequently, for every cut C,

Q(C) + \(\Phi\)(C) \(\geq\) \(\Lambda\)(C)

Equivalently, with Alt(C) := \(\prod\)\_\{v\(\in\)C\} Alt\_v we have
log\(_2\) Alt(C) = \(\Phi\)(C), and \(\Phi\)\_v \(\geq\) R\_v - q\_v and
\(\Phi\)(C) \(\geq\) \(\Lambda\)(C) - Q(C).

\textbf{Proof (structured):}

\textbf{Step 1 - Instance compliance (A1-A5).} The L* overlay satisfies
the axioms by construction:

\begin{itemize}
\tightlist
\item
  \textbf{A1 Hermeticity:} Designated pools; no side channels
\item
  \textbf{A2 Injectivity:} Length-delimited Enc over sorted ancestry and
  local outputs
\item
  \textbf{A3 Emergence:} Rank/requirement R\_v is well-defined and
  auditable
\item
  \textbf{A4 Closure/Recoverability:} Seeds/digests are recomputable by
  the verifier
\item
  \textbf{A5 Dependency:} Seed chaining makes addresses undefined before
  parents fix
\end{itemize}

For cut composition (Step 4), we use \textbf{H1-H5} (worst-case
properties: disjoint pools, unique-neighbor, Enc injectivity,
realizability, no cross-coupling).

\emph{No use of SCL in this step.}

\emph{Note on information measures.} In §5.5 we use Shannon-style
mutual-information only to upper-bound \textbf{per-step fresh-bit
inflow} by B (Lemma 5.5.1) and to connect first-use bits to actual read
steps (Lemma 5.5.1.b). The \textbf{SCL itself} is worst-case, zero-error
(Hartley) counting; no Shannon averaging enters the SCL inequality.

\textbf{Step 2 - Distinguishability (collision \(\Rightarrow\) error).}
("Merging causes errors.") If two worlds at v disagree on any of the
R\_v \(-\) q\_v unresolved bits, Keyedness + Injectivity imply they
induce \textbf{different seeds}, hence \textbf{different designated
addresses} downstream. Specifically, if algorithm state S represents
both world \(\omega\)\(_1\) (with Seed\(_1\)) and world \(\omega\)\(_2\)
(with Seed\(_2\)), then S can compute addresses for at most one seed.
When S computes F\_overlay(Seed\(_1\); j,\(\ell\)) (see Appendix A) but
the actual world is \(\omega\)\(_2\), it reads the wrong memory
location, obtaining incorrect salt \(\sigma\)\_wrong instead of
\(\sigma\)\_correct, leading to wrong output y\_v. Any state merge
before a legitimate read misroutes an address, yielding an incorrect
transcript. Therefore, each unresolved bit doubles the number of
seed-consistent classes that must be kept distinct (formalized as Lemma
7.Misroute below).

\textbf{Lemma 7.Misroute (Collision implies verifiable error).} Let Pref
be any transcript prefix up to node v and suppose two feasible worlds
\(\omega\)\(_1\) \ensuremath{\not\approx}\_v \(\omega\)\(_2\) (inequivalent under
seed-consistent equivalence) are represented by a single machine state S
at the moment immediately before the next designated read on the seed
chain through v. Then there exists a designated index (j,\(\ell\)) such
that F\_overlay(Seed\_v(\(\omega\)\(_1\)); j,\(\ell\)) \(\neq\)
F\_overlay(Seed\_v(\(\omega\)\(_2\)); j,\(\ell\)), and the next
designated read performed by S yields a payload inconsistent with at
least one of \{\(\omega\)\(_1\),\(\omega\)\(_2\)\}. Consequently,
extending Pref with that read produces a transcript that the NP verifier
rejects for at least one of \{\(\omega\)\(_1\),\(\omega\)\(_2\)\}.

\emph{Proof.} By Injectivity and Keyedness, inequivalence
\(\omega\)\(_1\) \ensuremath{\not\approx}\_v \(\omega\)\(_2\) implies Seed\_v(\(\omega\)\(_1\))
\(\neq\) Seed\_v(\(\omega\)\(_2\)) and induces different designated
addresses downstream along the current chain. Hence there exists
(j,\(\ell\)) with F\_overlay differing. Since Hermeticity forbids hidden
channels, the next designated read used to continue along the chain must
target exactly one of these addresses. If S reads using
Seed\_v(\(\omega\)\(_1\)), but the actual world is \(\omega\)\(_2\) (or
vice versa), the read returns \(\sigma\)\_wrong rather than
\(\sigma\)\_correct. Recoverability (Appendix J, Lemma J.0) ensures the
seeds and ancestor tuples are reconstructible from the transcript; the
verifier recomputes the designated address for the actual seed and
detects the mismatch, rejecting the transcript. \(\blacksquare\)

\textbf{Lemma 7.K (Keyedness \(\Rightarrow\) no merge before
resolution).} Let two partial worlds \(\omega\)\(_1\),\(\omega\)\(_2\)
at node v induce different seeds Seed\(_1\)\(\neq\)Seed\(_2\) while some
bits of y\_v remain unresolved. Any solver state S that purports to
represent both \(\omega\)\(_1\) and \(\omega\)\(_2\) must, upon the next
designated access tied to v\textquotesingle{}s chain, compute addresses
with either Seed\(_1\) or Seed\(_2\). In the former case, if the actual
world is \(\omega\)\(_2\), S misaddresses and reads \(\sigma\)\_wrong;
symmetrically if the latter and the world is \(\omega\)\(_1\). Hence a
single state cannot remain correct for both worlds without first
resolving the distinguishing bits. Therefore, different seeds cannot be
merged prior to resolution. \emph{Proof.} Address functions are
seed-dependent (Appendix A); disjoint pools prevent accidental aliasing
(§6.2.3). Using the wrong seed yields a different read, breaking
correctness under Hermeticity. \(\blacksquare\)

\textbf{Step 3 - Per-node counting.} Let s\_v := R\_v - q\_v (unresolved
bits at v). By Step 2, each unresolved bit doubles the number of
inequivalent seed-consistent classes that must be kept distinct
(different values \(\to\) different seeds \(\to\) different designated
addresses). Therefore Alt\_v \(\geq\) 2\^{}(s\_v) = 2\^{}(R\_v-q\_v)
(trivial when R\_v = q\_v, since Alt\_v \(\geq\) 1). Taking logarithms:
q\_v + log\(_2\)(Alt\_v) \(\geq\) R\_v.

H1-H5 (worst-case lane hypotheses): H1 Disjoint designated pools
\{U\_v\}; H2 Keyedness/Injectivity (no cross-seed merges); H3
Closure/Recoverability (parseable Enc); H4 Realizability (Emergence at
v); H5 No cross-coupling across the horizon (gate digests depend only on
pre-horizon indices).

\textbf{Step 4 - Composition across a cut (worst-case proof via H1-H5).}
We first verify L* satisfies the hypotheses H1-H5, then apply Cartesian
factoring.

\textbf{Verification that L* satisfies H1-H5:}

\begin{itemize}
\tightlist
\item
  \textbf{H1 (Disjoint pools)}: §6.2.3 assigns disjoint address ranges
  U\_v to each node v, ensuring U\_v \(\cap\) U\_\{v\textquotesingle{}\}
  = \(\emptyset\) for v \(\neq\) v\textquotesingle{} on any cut C
\item
  \textbf{H2 (Enc injectivity)}: §6.2.7 defines Enc as injective
  length-delimited encoding with sorted ancestry; Seed\_v = Enc(v
  \textbar{}\textbar{} sort(\{(u,Seed\_u,y\_u):u\(\in\)P(v)\})
  \textbar{}\textbar{} GateDigest\_v)
\item
  \textbf{H3 (Closure/Recoverability)}: Lemma J.0 (Appendix J) proves
  sink seeds deterministically parse and recover all ancestor (Seed,y)
  tuples via Enc parseability
\item
  \textbf{H4 (Realizability)}: Lemma 6.1 (§6.3 Emergence/Realizability)
  proves every x\_v \(\in\) \{0,1\}\^{}(R\_v) remains realizable - we
  can choose unread salts to produce any target value; across cut C,
  disjoint pools (H1) + no cross-coupling (H5 below) ensure independent
  extension to feasible worlds
\item
  \textbf{H5 (No cross-coupling)}: Lemma 7.2.1-NC (proven below) shows
  gate digests GateDigest\_v use only pre-horizon indices S(P) with
  GREQ=0, so post-horizon digests don\textquotesingle{}t couple
  unresolved coordinates across distinct v \(\in\) C
\end{itemize}

\textbf{Cut-Factoring Checklist} (prerequisites for Cartesian product
across cut C):

\begin{itemize}
\tightlist
\item
  {[}YES{]} \textbf{H1}: Disjoint designated pools \{U\_v\} for v
  \(\in\) C (§6.2.3)
\item
  {[}YES{]} \textbf{H2}: Enc injective and parseable; Seed\_v uniquely
  determined (§6.2.7)
\item
  {[}YES{]} \textbf{H3}: Closure/Recoverability from sink seeds (Lemma
  J.0)
\item
  {[}YES{]} \textbf{H4}: Realizability - every local assignment extends
  to feasible world (Lemma 6.1)
\item
  {[}YES{]} \textbf{H5}: No cross-coupling - digests use only
  pre-horizon indices (Lemma 7.2.1-NC below)
\end{itemize}

With H1-H5 verified \(\Rightarrow\) \(\Pi\)\_C(\(\pi\)) \(\cong\)
\(\prod\)\_\{v\(\in\)C\} S\_v (Lemma J.1-Cart, Appendix J)

\textbf{Cartesian factoring:} By Lemma J.1-Cart (Appendix J), H1-H5
imply feasible worlds across cut C factor as \(\Pi\)\_C \(\cong\)
\(\prod\)\emph{\{v\(\in\)C\} S\_v, yielding Alt(C) =
\(\prod\)}\{v\(\in\)C\} Alt\_v \(\geq\)
2\^{}(\(\Sigma\)\_\{v\(\in\)C\}(R\_v-q\_v)). Therefore
\(\Sigma\)\_\{v\(\in\)C\} log\(_2\)(Alt\_v) \(\geq\)
\(\Sigma\)\_\{v\(\in\)C\}(R\_v - q\_v), i.e., Q(C) + \(\Phi\)(C)
\(\geq\) \(\Lambda\)(C).

\textbf{Cut Cartesian Product - proof sketch.} Fix a transcript prefix
\(\pi\). For each node v in a cut C, let S\_v denote the set of locally
feasible assignments to the unresolved coordinates at v that are
consistent with \(\pi\) (realizability, H4) and produce distinct
designated addresses on the current seed chain (Keyedness + Injectivity,
H2). Define \(\Pi\)\_C(\(\pi\)) to be the feasible world choices across
the cut consistent with \(\pi\). We show a bijection \(\Pi\)\_C(\(\pi\))
\(\cong\) \(\prod\)\_\{v\(\in\)C\} S\_v:

\begin{itemize}
\item
  \textbf{Injectivity}: Disjoint address pools (H1) ensure that changing
  the local choice at any single v changes only addresses in U\_v; Enc
  is parseable/injective (H3), so two distinct local choices at v cannot
  yield the same future designated addresses. Hence different tuples in
  \(\prod\) S\_v map to different elements of \(\Pi\)\_C(\(\pi\)).
\item
  \textbf{Surjectivity}: Closure/recoverability (H3) ensures that any
  feasible world choice across C consistent with \(\pi\) decomposes into
  local choices \{s\_v\} with s\_v \(\in\) S\_v, because post-horizon
  gate digests depend only on pre-horizon indices and cannot couple
  unresolved coordinates across distinct v \(\in\) C (no cross-coupling,
  H5; Lemma 7.2.1-NC). Thus every element of \(\Pi\)\_C(\(\pi\)) arises
  from some tuple in \(\prod\) S\_v.
\end{itemize}

Therefore \textbar{}\(\Pi\)\_C(\(\pi\))\textbar{} =
\(\prod\)\_\{v\(\in\)C\} \textbar{}S\_v\textbar{}, yielding the
Cartesian product bound claimed. See Appendix J (Theorem J.1, Lemma
J.1-Cart) for the full formalization.

\emph{(Technical note: The Cartesian factoring applies to
overlay-feasible worlds - the combinatorial structure of cut-coordinate
tuples that produce valid overlay configurations. The satisfiability
constraint (whether the decoded \(\varphi\) admits a witness w) is an
independent filter applied afterward and does not introduce coupling
among the cut coordinates. This two-level structure ensures the product
bound holds for the overlay enumeration space, which drives the lower
bound. See Lemma 5.1.NC and Appendix J post-J.1-REAL discussion.)}

\textbf{Lemma 7.2.1-NC (No cross-coupling across the bottleneck cut).}
In the Frontier-Gate (FG) wiring (§6.2.8; Appendix C.1.1), each
GateDigest\_v is computed only from pre-horizon indices S(P(v)) where
GREQ=0, so it does not depend on any unresolved coordinate on the
bottleneck cut C*. Hence post-horizon digests cannot couple free choices
across C*, and the Cartesian-product factoring used above is valid.

\emph{Proof.} By construction in §6.2.8 and Appendix C.1.1, the index
set S(P) for gate digest GateDigest\_v is drawn exclusively from nodes
with GREQ=0 (pre-horizon). The bottleneck cut C* contains nodes with
both GREQ=0 (pre-horizon, can contribute to s) and GREQ=1 (post-horizon,
require gate proofs). Since GateDigest\_v depends only on GREQ=0
indices, and the unresolved coordinates on C* at any prefix are those on
GREQ=1 nodes, GateDigest\_v cannot couple these unresolved coordinates.
Therefore different choices for unresolved bits at distinct
v,v\textquotesingle{} \(\in\) C can be made independently without
creating digest-based constraints across them. \(\blacksquare\)

\textbf{Step 5 - (For later use) Min-cut residual.} Define the
\textbf{run-dependent residual}

\(\lambda\)(A,x) := min\_C (\(\Lambda\)(C) - Q(C))

(Here Q depends on A via RWA; we drop the subscript for brevity,
consistent with §1-§5.)

The theorem itself establishes the cut inequality; tightness/min-cut
optimality and pricing into time are handled in §7.3.

\ensuremath{\square}

\textbf{Where each ingredient lives:}

\begin{itemize}
\tightlist
\item
  \textbf{A1-A5 compliance:} §6 (overlay construction), Lemma 4.A (for
  A1 see also Lemma 4.2.I)
\item
  \textbf{H1-H5 (worst-case cut properties):} Lines 3454-3459
  (Hypotheses section above); formalized in Appendix J
\item
  \textbf{Keyedness/Injectivity lemmas:} Lemma 7.K (local proof), §4.2
  (framework), §6.2.7 (seeds)
\item
  \textbf{Per-node Alt bound:} Lemma 7.I above (see also Lemma 4.2.2 and
  Lemma 5.3.1 for representation-invariance)
\item
  \textbf{Cut composition (worst-case):} Step 4 above via H1-H5, Lemma
  J.1-NC (no cross-coupling), Appendix J (Theorem J.1, Lemma J.1-Cart)
\item
  \textbf{RWA justification:} "Why RWA is canonical" explanation in
  §7.0.2
\item
  \textbf{Min-cut residual \& pricing:} §7.3 (segments/arity), Appendix
  C (time lower bounds)
\end{itemize}

\textbf{Remark (Shannon vs SCL).} SCL is a \textbf{zero-error,
worst-case} conservation law (Hartley/Rényi-0) over feasible worlds,
whereas Shannon is \textbf{average-case}. We use SCL here because it
composes across cuts and feeds \textbf{directly into worst-case time}
via the later pricing lemmas.

\begin{center}\rule{0.5\linewidth}{0.5pt}\end{center}

\paragraph{7.3 Extensions and
Manifestations}\label{73-extensions-and-manifestations}

Notation (\(\Delta\) matches Appendix J). We write \(\Delta\)(C) :=
\(\Lambda\)(C) \(-\) (Q(C) + log\(_2\) Alt(C)) (Eq. (7.3)) for the
residual deficit on a cut C, matching the \(\Delta\) used in Appendix J.
On a designated bottleneck cut C*, \(\Delta\)(C*) controls tries via
\ensuremath{\mathbb{E}}{[}tries{]} \(\geq\) 2\^{}(\(\Delta\)(C*)) and underpins the phase
threshold (Corollary 7.3.PT). To avoid symbol clashes: we denote graph
max in-degree by \(\Delta\)\_in (not \(\Delta\)); \(\Delta\) in lemma
names (e.g., Lemma A.1.\(\Delta\)) is part of the label, not a variable;
and \(\Delta\)(C) is reserved for the across-tries deficit.

Manifestations. Having established the Semantic Conservation Law
(Theorem 7.A), we now show how it appears across computational paradigms
through projection templates.

\subparagraph{7.3.1 Projection Template: Backtracking/Search
Trees}\label{731-projection-template-backtrackingsearch-trees}

This template shows how to apply the abstract SCL to a specific
computational paradigm.

\textbf{Projection Contract for Backtracking:}

\textbf{1. Artifacts Counted (What is Alt\_v):}

\begin{itemize}
\tightlist
\item
  Alt\_v = number of active branches at node v
\item
  Each branch represents a distinct partial solution being explored
\item
  Branches are distinguishable if they make different choices up to node
  v
\end{itemize}

\textbf{2. Receiving-Window Attribution (RWA) (How q\_v is measured):}

\begin{itemize}
\tightlist
\item
  q\_v = bits of information learned through exploring at node v
\item
  Includes: variable assignments made, constraints checked, pruning
  decisions
\item
  Counted at first discovery (not re-counted if same info used later)
\end{itemize}

\textbf{3. Cut Exposure (How cuts reveal complexity):}

\begin{itemize}
\tightlist
\item
  A cut C through the DAG represents a frontier in the search
\item
  Total active branches = \(\Pi\)\_\{v\(\in\)C\} (branches at v)
\item
  When each node has 2\^{}(R\_v - q\_v) possibilities, total =
  2\^{}\(\Sigma\)(R\_v - q\_v) = 2\^{}(\(\lambda\)(A,x)) Example: two
  nodes on C with one unresolved bit each force \(\geq\) 4 live branches
  at the frontier.
\end{itemize}

\textbf{4. Cost Translation (Paradigm-specific bounds):}

\begin{itemize}
\tightlist
\item
  Tree size \(\geq\) total branches explored \(\geq\)
  2\^{}(\(\lambda\)(A,x))
\item
  Time complexity \(\geq\) tree size (must visit each branch)
\item
  Space complexity \(\geq\) max concurrent branches (at widest frontier)
\end{itemize}

\textbf{Application Result:} Backtracking requires tree size \(\geq\)
2\^{}(\(\lambda\)(A,x)) nodes

\textbf{Uses:} Theorem 7.A, Receiving-Window Attribution (RWA), DAG cut
analysis

\subparagraph{7.3.2 Projection Template
Guide}\label{732-projection-template-guide}

To apply SCL to other paradigms, follow this template:

\begin{enumerate}
\def\labelenumi{\arabic{enumi}.}
\tightlist
\item
  \textbf{Define Artifacts:} What are the distinguishable computational
  objects?
\item
  \textbf{Define First-Use:} How do you measure information acquisition?
\item
  \textbf{Analyze Cuts:} How do artifacts compose across parallel nodes?
\item
  \textbf{Translate Costs:} How do artifact counts become time/space?
\end{enumerate}

Example instantiations:

\begin{itemize}
\tightlist
\item
  \textbf{Dynamic Programming:} Alt\_v = distinct table keys;
  composition = Cartesian product
\item
  \textbf{Resolution:} Alt\_v = active clauses; composition = clause
  combinations
\item
  \textbf{OBDD:} Alt\_v = nodes at level; composition = subfunctions
  represented
\item
  \textbf{CDCL:} Alt\_v = learned clauses; composition = clause database
  size
\end{itemize}

Each paradigm mechanically plugs into the same abstract framework.

\subparagraph{7.3.3 Artifact Extractors and Cost Bridges
(formal)}\label{733-artifact-extractors-and-cost-bridges-formal}

We make explicit the artifact extractors E\_v and the bridges from
artifact counts/width/degree to standard resource measures. Throughout,
Lemma 7.I supplies Alt\_v \(\geq\) 2\^{}(R\_v-q\_v).

\begin{itemize}
\item
  Backtracking / explicit search trees

  \begin{itemize}
  \tightlist
  \item
    Extractor E\_v: active partial assignments/choices (branches) up to
    v, keyed by the resolution prefix.
  \item
    Bridge: time \(\geq\) number of visited nodes; maximal concurrent
    branches \(\geq\) width at a frontier. Hence size(tree) \(\geq\)
    Alt(C) and time \(\geq\) 2\^{}(\(\lambda\)(A,x)).
  \end{itemize}
\item
  Dynamic programming (table-based)

  \begin{itemize}
  \tightlist
  \item
    Extractor E\_v: distinct state keys required for correctness
    (including seed-keyed history where applicable).
  \item
    Bridge: memory \(\geq\) \#keys at widest frontier; time \(\geq\)
    total distinct key computations. Hence space/time \(\geq\)
    2\^{}(\(\lambda\)(A,x)) across a min-cut.
  \end{itemize}
\item
  Resolution / CDCL (width \(\to\) size)

  \begin{itemize}
  \tightlist
  \item
    Extractor E\_v: clause profiles sufficient to refute assignments
    differing on unresolved bits (degree-type variant counts width).
  \item
    Bridge: width lower bound \(\Omega\)(\(\sum\) s\_v on a cut) yields
    exponential size by {[}BEN01{]}. Result: size \(\geq\)
    2\^{}(\(\Omega\)(\(\lambda\)(A,x))).
  \end{itemize}
\item
  OBDD / ROBDD (any variable order)

  \begin{itemize}
  \tightlist
  \item
    Extractor E\_v: distinct nodes at some level corresponding to
    unresolved bits at v (seed-consistent labeling makes worlds
    future-distinguishable for any order).
  \item
    Bridge: There exists a level with width \(\geq\)
    2\^{}(\(\Omega\)(\(\lambda\)(A,x))) for appropriate
    expander/Tseitin-style encodings of our gates; see {[}WEG00{]} for
    order-robust lower bounds. Result: size \(\geq\)
    2\^{}(\(\Omega\)(\(\lambda\)(A,x))) for our gate family.
  \end{itemize}
\end{itemize}

\subparagraph{7.3.4 {[}Technical Implementation
Note{]}}\label{734-technical-implementation-note}

Readers may skip to §7.3.5. This subsection records technical details
about single-run versus restart analysis that are not essential for
following the main results.

It is helpful to distinguish \textbf{semantic indistinguishability} (a
property of the \emph{instance}) from \textbf{operational
indistinguishability} (a property of a \emph{run} of the solver).

\textbf{Operational indistinguishability (segments).} In a
\textbf{single run} of a deterministic TM, information may persist. The
same bit, once legitimately revealed, can support many "virtual leaves"
in the semantic tree without being repaid; summing leaf depths would
then \textbf{double-count information}. To respect persistence, we
partition the run into \textbf{rollback segments} (maximal subruns
between rollback keys / resolution prefixes). Within a segment the
resolution prefix across the bottleneck cut (including the cut-level
ConstraintDigest and the functional determinations q on C) is fixed; at
a rollback the prefix changes. By Lemma C.2.3, a non-accepting segment
can eliminate at most \textbf{one} additional world beyond the pruning
caused by newly determined cut information \(\Delta\)q, giving the
\textbf{Segment Counting} bound

m\_seg \(\geq\) 2\^{}(\(\rho\)-s),

where m\_seg is the number of segments, \(\rho\) is the effective
residual on the cut, and s is the number of cut bits that (by the start
of the last segment) already agree with the accepting world. In short:
\textbf{when knowledge persists, we count segments} and then multiply by
a per-segment time baseline.

\textbf{Paths \(\to\) Segments}: The 2\^{}(R-q) distinct computational
paths (semantic possibilities) cannot be compressed - they manifest
operationally as \(\geq\) 2\^{}(\(\rho\)-s) distinct exploration
segments. Each segment represents a committed resolution pattern;
keyedness prevents merging across different patterns without causing
errors.

\textbf{The worlds-to-segments bridge:} Segments are the operational
mechanism through which algorithms physically explore the space of
unresolved worlds. Each segment commits to a distinct prefix of
decisions about those worlds (a specific resolution pattern), and
keyedness prevents segments with different prefixes from sharing
computational work. Thus the 2\^{}(R-q) semantic possibilities (worlds)
manifest as \(\geq\) 2\^{}(\(\lambda\)-s) operationally distinct
exploration segments. The structural necessity of maintaining
distinguishable artifacts for unresolved worlds translates directly into
the operational necessity of exploring distinct segments - the worlds do
not just exist abstractly, they force concrete computational exploration
through segments.

\textbf{Why not use the semantic tree directly in a single run?} The
semantic tree lower-bounds the \emph{sum of leaf depths}. In a single
run, many leaves share the same early revelations; those early bits are
paid \textbf{once}, not once per leaf. Counting worlds would therefore
overestimate cost. Counting \textbf{segments} restores a prefix-free
structure exactly where the key changes (rollbacks), i.e., precisely
where the semantic core forces fresh distinguishing work.

\textbf{Remark (Segment keys and reuse).} Segment boundaries coincide
with changes to the resolution prefix across the bottleneck cut
(seed-chain projection plus ConstraintDigest\_C and WorldCommit\_C). Any
reuse of designated gate objects or committed-world refutations across
segments would require an unchanged prefix, contradicting the definition
of a new segment. Reuse within a segment is permitted and already
accounted for in the per-segment baseline; reuse across segments is
blocked by Keyedness/CDT and the definition of segments.

\textbf{Per-segment model hook.} Our path-based "MV-Rollback" argument
under the QP-sharp profile shows that each \textbf{non-accepting}
segment performs \(\Omega\)(n/log n) designated verification steps (via
producing a fresh gate object, growing NF\_C, or refuting the committed
world); on a k-tape TM, each designated step costs \(\Omega\)(1) time.
Consequently a single run satisfies

time\_TM \(\geq\) \(\Omega\)(2\^{}(\(\rho\)-s) · n/log n).

This bound follows directly from Segment Counting, and becomes
quasi-polynomial whenever s is bounded away from \(\lambda\)(A,x) (see
Appendix C for instance-side mechanisms that enforce such caps in
concrete families).

\textbf{Lane summary.}

\begin{itemize}
\tightlist
\item
  \textbf{Single run (unbounded space):} persistence \(\Rightarrow\)
  \textbf{count segments}; SC × MV-Rollback \(\Rightarrow\) time
  \(\Omega\)(2\^{}(\(\rho\)\(-\)s)·n/W\_min). Concrete caps on s are
  instance-specific; see Appendix C.
\end{itemize}

\textbf{On concrete caps for s.} For specific instance families (e.g.,
L*), one can enforce caps on pre-final agreement s by embedding
structural constraints; see Appendix C for a complete treatment and
tight bounds.

\textbf{Why caching does not help across segments (single run).}

\textbf{Lemma C.2.1 (Bottleneck residual under \(\tau\)-cap).} (alias:
formerly Lemma C.2.0) Let C* be the designated bottleneck cut for the
published profile with base residual \(\lambda\)\_base. Let s be the
pre-final agreement across C* as bounded by the gate-horizon budget
(Section 6.2.10). Then the run-specific min-cut residual \(\rho\)
satisfies

\(\rho\) \(\geq\) \(\lambda\)\_base \(-\) s.

\emph{Proof.} By profile construction, C* attains the minimum base
residual among all cuts, i.e., for every cut C, base\_residual(C)
\(\geq\) \(\lambda\)\_base. Early revelation allowed before the gate
horizon is confined to GREQ=0 nodes on C* (Section 6.2.10), totaling at
most s bits (Gate-Horizon Budget). Any other cut C either (i) also
crosses those GREQ=0 nodes, in which case its residual drops by at most
s, or (ii) does not, in which case it retains base\_residual(C) \(\geq\)
\(\lambda\)\_base. Therefore the run-specific residual on every cut
remains \(\geq\) \(\lambda\)\_base \(-\) s, hence \(\rho\) = min\_C
residual(C) \(\geq\) \(\lambda\)\_base \(-\) s. \ensuremath{\square} \textbf{Q:} If the
solver caches results, can it avoid paying again after a rollback?
\textbf{A:} No. Rollbacks change the resolution key. Different keys mean
different designated addresses (by keyedness). Previous artifacts become
invalid for the new key\textquotesingle{}s addresses. The solver must
re-compute all verifications for the new chain, even if similar
computations were done under a different key. Together these enforce the
per-segment baseline (MV-Rollback) once the key changes; see the
per-segment baseline lemma in \textbf{Appendix C.1}.

\subparagraph{7.3.5 Direct Consequences}\label{735-direct-consequences}

\textbf{Corollary 7.1.1 (FrontierPeak Lower Bound).} For any s-t cut C:
log\(_2\)(frontier across C) \(\geq\) \(\Sigma\)\_\{v\(\in\)C\}(R\_v -
q\_v), hence FrontierPeak \(\geq\) 2\^{}(\(\lambda\)(A,x)) (Eq. (7.1)),
where \(\lambda\)(A,x) is the run-dependent min-cut capacity.

Notation disambiguation (Alt(C*) contexts).

\begin{itemize}
\tightlist
\item
  FrontierPeak/single-run context: Alt(C*, t) counts concurrent
  distinguishable artifacts across the fixed witness cut C* at time t,
  and FrontierPeak := max\_t Alt(C*, t). In runs that maintain keyed
  state, FrontierPeak \(\geq\) 2\^{}(\(\lambda\)(A,x)) over the entire
  run.
\item
  Restart/tries context: Alt\_i(C*) denotes the maximum artifacts
  maintained within a single try i when state is forgotten between
  tries; typically Alt\_i(C*) \(\leq\) poly(n\_core). The across-tries
  inequality (Lemma 7.R) uses this per-try cap Alt\_i(·), not
  FrontierPeak.
\end{itemize}

These are different quantities. In the restart lane the artifacts per
try are small (polynomial), but exponentially many tries are required;
in the single-run lane the run maintains exponentially many concurrent
artifacts, leading to different pricing but the same conservation-law
obligation.

\emph{Follows from Theorem J.1 (DAG Min-Cut Read-or-X; cut composition
of the per-node SCL).} \textbf{Prefetch / reorder:} RWA charges the
\emph{first revealing read} to node v; reordering cannot reduce q\_v.
\textbf{Guess / partial y\_v:} GSS + Keyedness \(\Rightarrow\) multiple
candidate seeds \(\Rightarrow\) \textbf{distinguishable artifacts}.
\textbf{Cross-attempt caching:} caches must be \textbf{keyed} by the
resolution prefix y\_\{\(\leq\)v\}; unkeyed reuse is incorrect on such
instances. \textbf{No look-ahead:} future addresses/artifacts are
determined by current resolutions; without resolving, they remain
undefined for the solver. \textbf{No premature inference:} independence
ensures unresolved coordinates remain undetermined until their own
first-use revelations. \textbf{Order choice (OBDD):} expander-parity
gate gives order-robust width growth. \textbf{Hybrid paradigms:} tracked
by the hybrid potential \ensuremath{\mathcal{H}}\_v with q\_v+\ensuremath{\mathcal{H}}\_v\(\geq\)R\_v; families
\textbf{multiply} across nodes by Dependency. \textbf{Parallel /
wide-word reads:} q\_v counts \textbf{independent bits learned}
(information units), not I/O ops. \textbf{Randomized algorithms:}
Per-instance bounds (Theorem 8.A) apply to any fixed run (coin-fixing).
For Structural OWF construction (§9), every FG-wired instance x*
requires super-polynomial witness-finding time regardless of coins.
\textbf{Compression / hashing:} a single \textbf{unkeyed} artifact
cannot be correct for multiple seeds on our instances; \textbf{keyed}
families are exactly what Alt\_v counts. \textbf{Streaming / one-pass
small memory:} a streaming solver that does not resolve the missing bits
must maintain the live frontier of distinguishable artifacts. By the
\textbf{FrontierPeak bound (Corollary 7.1.1), this requires space
\(\geq\) max\_v \(\prod\)\_i\(\leq\)\_t 2\^{}(R\_i-q\_i); otherwise it
must resolve} and is charged in q\_v. Either way, Read-or-X applies.
\textbf{Incomplete parents (DAG):} advancing at v without all parents
fixed creates multiple candidate successor identities; distinct
possibilities form a frontier of artifacts counted in Alt\_v.

\textbf{Corollary 7.1.2 (Space Lower Bound).} For any FG-wired instance
x* (Structural OWF construction, §9), any correct solver requires space
\(\geq\) max\_t \(\prod\)\_\{i\(\leq\)t\} 2\^{}(R\_i-q\_i) for
witness-finding.

\textbf{Corollary 7.1.3 (Universal Manifestation).} For any classical
search algorithm on DAG computations with information requirements:

\begin{itemize}
\tightlist
\item
  Either \(\lambda\)(A,x) = o(n) (polynomial verification through full
  commitment), OR
\item
  The algorithm exhibits \textbf{exponential} growth in at least one
  semantic measure:

  \begin{itemize}
  \tightlist
  \item
    Backtracking: tree branches \(\geq\) 2\^{}(\(\lambda\)(A,x))
  \item
    Dynamic Programming: distinct keys \(\geq\) 2\^{}(\(\lambda\)(A,x))
  \item
    Resolution/CDCL: proof size \(\geq\) 2\^{}(\(\lambda\)(A,x))
  \item
    OBDD: diagram size \(\geq\) 2\^{}(\(\lambda\)(A,x))
  \item
    Branching Programs: states/subfunctions \(\geq\)
    2\^{}(\(\lambda\)(A,x))
  \item
    Algebraic systems: degree \(\geq\) \(\Omega\)(\(\lambda\)(A,x))
  \end{itemize}
\end{itemize}

\textbf{This is not merely algorithmic limitations - SCL formalizes the
structural necessity that enforces these costs on any correct algorithm;
for L*, the construction makes the obligations explicit and verifiable.}

\begin{center}\rule{0.5\linewidth}{0.5pt}\end{center}

\subparagraph{\texorpdfstring{7.3.6 Why Standard Solver Tricks
Don\textquotesingle{}t
Help}{7.3.6 Why Standard Solver Tricks Don\textquotesingle t Help}}\label{736-why-standard-solver-tricks-dont-help}

\textbf{Common solver optimizations and how SCL constrains them:}

\begin{itemize}
\tightlist
\item
  \textbf{Prefetch/reorder:} RWA charges the \emph{first revealing read}
  to node v; reordering cannot reduce q\_v
\item
  \textbf{Guess/partial y\_v:} Multiple candidate successors \(\to\)
  \textbf{distinguishable artifacts} (by Dependency)
\item
  \textbf{Cross-attempt caching:} Caches must be \textbf{keyed} by the
  resolution prefix; unkeyed reuse is incorrect
\item
  \textbf{Look-ahead:} Future artifacts are determined by unresolved
  outputs; cannot be known before first-use
\item
  \textbf{No premature inference:} Unresolved coordinates remain
  undetermined until their own first-use revelations
\item
  \textbf{Order choice (OBDD):} Width lower bounds are order-robust (see
  Appendix B.1)
\item
  \textbf{Hybrid paradigms:} Tracked by hybrid potential \ensuremath{\mathcal{H}}\_v; families
  multiply across nodes
\item
  \textbf{Parallel reads:} q\_v counts \textbf{independent bits
  learned}, not I/O operations
\item
  \textbf{Randomization:} Per-instance bounds (§8) apply to any fixed
  run via coin-fixing
\item
  \textbf{Compression:} Merging distinct worlds is incorrect; maintained
  families are counted in Alt\_v
\item
  \textbf{Streaming:} Must maintain a frontier of distinguishable
  artifacts or resolve (either way pays)
\item
  \textbf{Incomplete DAG traversal:} Creates multiple successor
  candidates \(\to\) distinguishable artifacts
\end{itemize}

Micro-lemmas (formal; one-line proofs)

Lemma 7.3.6.1 (Prefetch/reorder doesn\textquotesingle{}t reduce q\_v).
Under RWA, q\_v is charged at first valid use and depends only on
functional determination, not wall-clock order (see Lemma 6.1-RWA).
Proof: Reordering the same observations leaves the \(\sigma\)-algebra
unchanged; q\_v is invariant. \(\blacksquare\)

Lemma 7.3.6.2 (Guess/partial y\_v \(\Rightarrow\) distinguishable
artifacts). Different guesses for y\_v induce different seeds for
descendants via Enc; by Lemma 7.I they are future-distinguishable and
counted in Alt\_v unless resolved. Proof: y\_v\(\neq\)y\_v'
\(\Rightarrow\) Seed\_child\(\neq\)Seed\_child' \(\Rightarrow\) distinct
designated addresses. \(\blacksquare\)

Lemma 7.3.6.3 (Cross-attempt caching must be keyed). Any cache reused
across attempts must be keyed by the resolution prefix; unkeyed reuse
conflates \(\neq\)-equivalent worlds and is incorrect. Proof: Keyedness
demands seed-consistency; otherwise different future addresses/outcomes
would share one artifact. \(\blacksquare\)

Lemma 7.3.6.4 (No look-ahead). Before first-use, future designated
addresses depend on unresolved seeds; Hermeticity forbids any other
channel to compute them. Proof: Dependency + Hermeticity implies
addresses are undefined until parents (and digests, if any) are fixed.
\(\blacksquare\)

Lemma 7.3.6.5 (No premature inference). With no node-v discovery reads,
I(x\_v;Tr\_\{\textless{}v\})=0 and all x\_v are realizable (Lemma
6.1-ZMI/RZ). Proof: Any inference would contradict Emergence/zero mutual
information. \(\blacksquare\)

Lemma 7.3.6.6 (OBDD order choice does not help). For expander-FG gates,
width lower bounds are order-robust (Appendix B/G; {[}WEG00{]}), so size
\(\geq\) 2\^{}(\(\Omega\)(\(\lambda\)(A,x))) for any variable order.
Proof: Expander structure forces \(\Omega\)(s) distinct subfunctions at
some level independent of order. \(\blacksquare\)

Lemma 7.3.6.7 (Hybrids tracked by hybrid potential). Hybrid strategies
can\textquotesingle{}t merge \(\neq\)-equivalent worlds; their artifact
families multiply across nodes and are counted by a hybrid potential
\ensuremath{\mathcal{H}}\_v with q\_v+\ensuremath{\mathcal{H}}\_v\(\geq\)R\_v. Proof: Lemma 7.I applies to any
sufficient statistic; composition follows Dependency. \(\blacksquare\)

Lemma 7.3.6.8 (Parallel reads don\textquotesingle{}t change
q\_v\textquotesingle{}s semantics). q\_v counts independent bits
learned, not I/O ops; parallel/wide-word reads reduce steps but not the
resolved bit count (cf. Lemma 5.5.1). Proof: Mutual-information inflow
is bounded per fresh bits, independent of scheduling. \(\blacksquare\)

Lemma 7.3.6.9 (Randomization doesn\textquotesingle{}t help; per-instance
bounds). Per-instance deterministic bounds (Theorem 8.A) apply to any
fixed run. For randomized algorithms, coin-fixing (Yao) reduces to
deterministic analysis: any coin string yields a deterministic run
subject to SCL. Hence per-instance bounds extend to randomized
algorithms. \(\blacksquare\)

Lemma 7.3.6.10 (Compression/hashing doesn\textquotesingle{}t help).
Compressing artifacts cannot merge future-distinguishability classes
without error; keyed families are exactly what Alt\_v counts. Proof:
Lemma 7.I requires injectivity on \(\equiv\)-classes; unkeyed merges are
incorrect. \(\blacksquare\)

Lemma 7.3.6.11 (Streaming/one-pass). Without resolving, a streaming
solver must maintain the live frontier; by FrontierPeak \(\geq\)
2\^{}(\(\lambda\)(A,x)), space blows up unless committing. Proof:
Otherwise it conflates \(\neq\)-equivalent worlds and errs.
\(\blacksquare\)

Lemma 7.3.6.12 (Incomplete DAG traversal). Advancing at v without all
parents fixed induces multiple candidate Seed\_v and successors; these
form a frontier counted in Alt\_v. Proof: Dependency:
Seed\_v=Enc(v\textbar{}\textbar{}parents\textbar{}\textbar{}digest) is
undefined until all parents are fixed. \(\blacksquare\)

\begin{center}\rule{0.5\linewidth}{0.5pt}\end{center}

\subparagraph{7.3.7 The Read-or-Rollbacks
Extension}\label{737-the-read-or-rollbacks-extension}

\textbf{Clarification on "tries":} In a single deterministic TM run, a
"try" denotes a maximal segment between rollback keys (resolution
prefixes); multiple tries occur via internal backtracking within one
computation, not separate machine runs.

\emph{Per-instance bounds.} The per-node SCL requirements apply to any
run on any instance. For Structural OWF construction (§9), \textbf{every
FG-wired instance} has deterministic witness-finding lower bounds
(Theorem 8.A), so the rollback bound holds for all such instances.

\textbf{Rollback key (a.k.a. restart key).} In any run that may
rollback/try a new candidate, define the attempt\textquotesingle{}s
\textbf{rollback key} at node v to be the resolution prefix
y\_\{\(\leq\)v\}. Because successor computations are determined by y\_v
(via Dependency), different prefixes induce different successor
identities and cannot be merged.

\textbf{Scope of a rollback (nearest consistent prefix).} Changing a
resolved choice y\_v after any successor of v has been seeded requires
rolling back to the \emph{shortest transcript prefix} whose frontier is
consistent with the new y\_v. A rollback is therefore \textbf{not
necessarily from scratch}; it is from the nearest \textbf{consistent
cut} above v.

\textbf{Change-of-mind checklist.}

\begin{enumerate}
\def\labelenumi{\arabic{enumi})}
\tightlist
\item
  \textbf{Before first reveal at v:} free to reconsider (no q\_v, no
  Alt\_v yet).
\item
  \textbf{After resolution at v, before any child is seeded:} either \textbullet{}
  \textbf{retry locally} (discard and choose a new y\_v; previously paid
  q\_v is sunk, and a new try begins just before v), or \textbullet{} \textbf{fork}
  (keep both options; Alt\_v \(\leftarrow\) Alt\_v ·
  2\^{}(\(\Delta\)q\_v)).
\item
  \textbf{After any child of v is seeded:} \textbf{cut-rollback} to the
  nearest consistent frontier; prior artifacts keyed to the old
  resolution prefix become invalid (\textbf{Dependency} +
  \textbf{Finality}). Cost appears as \textbf{extra tries}
  (Read-or-Rollbacks) and/or \textbf{extra LOADs/time} (windowed-LOAD).
\end{enumerate}

\textbf{One step or many?} If no child was seeded, the rollback is one
step (just before v); otherwise it drops exactly the invalidated sub-DAG
and no more.

\textbf{Accounting reminder.} We count \textbf{distinct rollback keys}
(resolution prefixes) across attempts; caches must be keyed to the same
prefix. Switching keys invalidates caches and is priced by
\textbf{tries/time}; keeping multiple keys alive is priced by
\textbf{Alt\_v}.

\textbf{Definition:} Let \textbf{tries\_\{\(\leq\)v\}} be the number of
\textbf{distinct} rollback keys y\_\{\(\leq\)v\} the solver ever
instantiates across its attempts (entire run).

\textbf{Persistent caches across attempts.} Any cache or materialized
state reused across attempts must be keyed by the resolution prefix
y\_\{\(\leq\)v\}; unkeyed reuse is incorrect on such instances. Such
keying contributes to \(\Phi\)\_v via the count of distinct keys/states.

\textbf{Lemma 7.C (Worst-case Read-or-Rollbacks).} For any deterministic
solver A and node v, \textbf{either} A resolves q\_v bits at node v,
\textbf{or} there exists an explicit input (consistent with the overlay)
on which

log\(_2\)(tries\_\{\(\leq\)v\}) \(\geq\) R\_v - q\_v.

\emph{Semantic interpretation:} "Tries" denotes \textbf{distinct
computational paths} characterized by their resolution prefixes
y\_\{\(\leq\)v\}. The solver must distinguish among 2\^{}(R\_v - q\_v)
distinguishable artifacts at node v.

\emph{Proof (per-instance).} For any FG-wired instance x*, if A avoids
resolving s\_v = R\_v \(-\) q\_v bits at node v, L*\textquotesingle{}s
structure forces A to explore distinguishable artifacts. By Dependency,
different y\_v induce different successor identities, hence distinct
rollback keys. By Keyedness and Injectivity (A2), this forces A to
realize at least 2\^{}(s\_v) distinct y\_\{\(\leq\)v\}. \(\blacksquare\)

\textbf{Corollary 7.1-RG (Worst-case global tries).} For any
deterministic solver A with resolution profile \{q\_v\}, there exists an
explicit input on which

tries\_\{\(\leq\)D\} \(\geq\) 2\^{}(\(\sum\)\_v (R\_v-q\_v)).

\emph{Proof.} Apply Lemma 7.C at each node and use Dependency: distinct
node-v choices yield distinct successor identities, so rollback keys
\textbf{multiply} along any root-to-leaf evolution. \textbf{The
multiplicative step:} Distinct y\_v values induce distinct successors;
hence distinct prefixes cannot coalesce later. The total number of
realized prefixes across the \textbf{entire} run is the \textbf{product}
of per-node factors. \(\blacksquare\)

\textbf{Distributional variant (via Yao\textquotesingle{}s principle).}
Under an \textbf{independent per-node requirement} distribution (when
available), for any randomized solver and a designated bottleneck cut
C*: \ensuremath{\mathbb{E}}{[}tries\_\{\(\leq\)sink\}{]} \(\geq\) 2\^{}(\(\lambda\)(A,x) -
log\(_2\) Alt(C*)), and hence \textbf{some fixed input} in the support
achieves this bound by Yao\textquotesingle{}s principle. When Alt(C*)
\(\leq\) poly(n\_core), this simplifies to
\ensuremath{\mathbb{E}}{[}tries\_\{\(\leq\)sink\}{]} \(\geq\)
2\^{}(\(\Omega\)(\(\lambda\)(A,x)))/poly(n\_core).

\textbf{Spectrum note.} Intermediate resolution deficits interpolate
smoothly: if \(\lambda\)(A,x) = \(\Theta\)(log² n) then
tries\_\{\(\leq\)sink\} = 2\^{}(\(\Theta\)(log² n)) (quasi-polynomial).

\subparagraph{7.3.8 Segment Counting
Analysis}\label{738-segment-counting-analysis}

Statement (informal). Partition any single run into rollback segments:
maximal subruns during which the resolution prefix across the designated
bottleneck cut C (including the cut-level counters) remains fixed. Let s
be the pre-final agreement (cut bits determined before the gate horizon)
and let \(\rho\) be the effective residual on C. Then the number of
non-accepting rollback segments that progress the final chain satisfies

m\_seg \(\geq\) 2\^{}(\(\rho\) - s).

Intuition. With a fixed prefix, the unresolved coordinates across C
define 2\^{}(\(\rho\)-s) feasible worlds that are pairwise
future-distinguishable (Injectivity + Keyedness). A non-accepting
segment can eliminate at most one such world beyond the pruning already
credited to newly determined cut information \(\Delta\)q during the
segment; otherwise two inequivalent worlds would be merged,
contradicting SCL. Hence at least 2\^{}(\(\rho\)-s) segments are
required to reach acceptance once s bits of agreement are spent.
Acceptance uniqueness at the final segment follows from digest binding
(Lemma C.2.ACC) and from the logical acceptance lemma (Lemma
C.2.ACC-logical). Frontier-Gate supplies the per-segment priced event
beyond the horizon.

For the formal lemma, definitions of the resolution prefix/counters, and
a complete proof, see \textbf{Appendix C.2}.

\subparagraph{7.3.9 Per-Try Cost and Time
Bounds}\label{739-per-try-cost-and-time-bounds}

\textbf{Model Hook.} Through arity-bounded analysis (§5.5.1), the
per-step information limit of B :=
k\(\lceil\)log\(_2\)\textbar{}\(\Gamma\)\textbar{}\(\rceil\) bits
directly yields exponential time lower bounds.

\emph{Notation.} Here \(\rho\) denotes the run-specific residual on the
min-cut just before the last segment; s is the number of cut bits
revealed before that segment (§12.2).

\textbf{Corollary C.1.1 (FG per-segment baseline).} Under FG-3 (Gate
requirement), every non-accepting rollback segment that progresses the
final chain computes a cut-gate digest over \textbar{}S(P)\textbar{} =
\(\Theta\)(n/W\_min) seed-dependent terms and therefore costs
\(\Omega\)(n/W\_min) TM steps (Lemma 5.5.1.c + Lemma A.1.\(\Delta\)).

\emph{No shortcut.} Because indices in S(P) are seed-dependent and
designated pools are disjoint - with \(\Theta\)(\textbar S(P)\textbar)
churn under seed changes (Lemma A.1.\(\Delta\)) - any attempt to produce
GateDigest without evaluating the \(\Theta\)(n/W\_min) selected terms
fails to compute the correct parity for some seed chain. Even with all
salts pre-cached, computing the digest requires evaluating these
addresses and XORing the corresponding cached values (Lemma 5.5.1.c).

\emph{Distributional \(\Rightarrow\) worst-case.} Since \ensuremath{\mathbb{E}}\_s{[}time{]}
\(\geq\) 2\^{}(\(\Omega\)(\(\lambda\))) under the product-salts
distribution, Yao\textquotesingle{}s principle yields a fixed seed s in
support with time \(\geq\) 2\^{}(\(\Omega\)(\(\lambda\))) for each
algorithm.

\textbf{Sharp verify/search threshold (Cor. 7.5.PT).} Under polynomial
space/hardware, \textbf{polynomial time} requires Q(C*) \(\geq\)
\(\Lambda\)(C*) - O(log n). Any \(\omega\)(log n) shortfall forces
super-polynomial time: time \(\geq\) 2\^{}(\(\omega\)(log
n))·\(\omega\)(log n).

\textbf{Per-try costs:} Each complete try requires at least
\(\Sigma\)\_v R\_v first-use revelations attributable along the cut. For
complete per-try analysis, see \textbf{Appendix C.4.3 (Lemma C.4.1)}.

\textbf{Corollary 7.3.PT (Sharp Phase Transition for Polynomial Time).}
(See also Corollary 8.3 for the constant-success threshold in terms of
Q(C*) and \(\lambda\)(A,x).)

\emph{Intuition:} There exists a fundamental discontinuity in the
complexity spectrum - algorithms must choose between polynomial time (by
resolving almost everything) or genuine search capability (requiring
super-polynomial time). There is no middle ground.

Formally, let C* be a min-cut with information requirement
\(\Lambda\):=\(\Lambda\)(C*). For any classical solver using polynomial
space/hardware (maintaining \(\leq\) poly(n\_core) states, launching
\(\leq\) poly(n\_core) parallel tries):

\textbf{The Phase Transition:} If expected time is polynomial, then the
resolution Q(C*) must satisfy:

Q(C*) \(\geq\) \(\Lambda\)(C*) - O(log n)

{[}Near-total resolution required{]}

Equivalently, the \textbf{residual deficit} \(\Delta\)(C*) :=
\(\Lambda\)(C*) - (Q(C*) + log\(_2\) Alt(C*)) must be O(log n\_core).
(We use \(\Delta\) here to avoid confusion with the per-step inflow
budget B in §5.5.1.)

\textbf{Critical Threshold:} Any strategy with Q(C*) \(\leq\)
\(\Lambda\)(C*) - \(\omega\)(log n) requires at least
2\^{}(\(\omega\)(log n)) = super-polynomial expected time. The optimal
is quasi-polynomial complexity 2\^{}(polylog(n)), achieved with
polylog(n) unresolved bits (e.g., \(\Theta\)(log² n) or \(\Theta\)((log
n)\^{}(3/2))).

\emph{Proof.} The across-tries inequality (Lemma 7.R) gives \ensuremath{\mathbb{E}}{[}tries{]}
\(\geq\) 2\^{}(\(\Delta\)(C*)) where \(\Delta\)(C*) := \(\Lambda\)(C*)
\(-\) (Q(C*) + log\(_2\) Alt(C*)). Polynomial resources reduce the
effective deficit by at most O(log n\_core) since:

\begin{itemize}
\tightlist
\item
  Parallel tries: log\(_2\)(poly(n\_core)) = O(log n)
\item
  Stored states: log\(_2\)(poly(n\_core)) = O(log n)
\end{itemize}

Thus \ensuremath{\mathbb{E}}{[}time{]} \(\geq\) 2\^{}(\(\Delta\)(C*)-O(log n\_core)) ·
(\(\Delta\)(C*)-O(log n\_core)) through bit-level analysis (the
information-theoretic accounting that a semantic decision tree with
2\^{}(\(\lambda\)) leaves has average leaf depth \(\geq\) \(\lambda\) by
Kraft\textquotesingle{}s inequality, so total first-use information
\(\geq\) 2\^{}(\(\lambda\)) · \(\lambda\)). For polynomial time, we need
\(\Delta\)(C*) = O(log n\_core). \ensuremath{\square}

\textbf{Optimality of Quasi-Polynomial:} This sharp threshold proves
quasi-polynomial is optimal for maintaining search capability:

\textbf{Resolution-Search Tradeoffs}

\begin{itemize}
\tightlist
\item
  \textbf{Near-verification}: Resolution Saved: O(log n); Complexity:
  Polynomial; Search Type: Not genuine search
\item
  \textbf{Optimal hybrid}: Resolution Saved: polylog(n); Complexity:
  Quasi-polynomial; Search Type: Best possible
\item
  \textbf{Moderate search}: Resolution Saved:
  \(\Theta\)(n\^{}\(\varepsilon\)),
  0\textless{}\(\varepsilon\)\textless1; Complexity: Sub-exponential;
  Search Type: Worse than optimal
\item
  \textbf{Heavy search}: Resolution Saved: \(\Theta\)(n); Complexity:
  Exponential; Search Type: Excessive
\end{itemize}

The quasi-polynomial strategy (leaving polylog(n) bits unresolved)
achieves the optimal balance - the minimum super-polynomial complexity
required for meaningful search.

\textbf{Corollary 7.B.1 (Time lower bounds; model-hooked)} (cf.
\textbf{Thm. C.4.2})

\textbf{Part A: Model-Independent (distributional +
Yao\textquotesingle{}s principle).} There exists a distribution \ensuremath{\mathcal{D}} over
instances such that for \textbf{every} randomized solver A,

\ensuremath{\mathbb{E}}\_\{x\textasciitilde{}\ensuremath{\mathcal{D}}\}{[}time\_A(x){]} \(\geq\)
2\^{}(\(\lambda\)(A,x)) · \(\lambda\)(A,x)

via arity-bounded analysis, where \(\lambda\) is the min-cut value.
Consequently, for \textbf{each} randomized solver A, there exists a
\textbf{fixed} instance x\_A in the support of \ensuremath{\mathcal{D}} with
\ensuremath{\mathbb{E}}{[}time\_A(x\_A){]} \(\geq\) 2\^{}\(\Omega\)(\(\lambda\)).

\emph{(Distributional context.) We do not claim a single instance works
simultaneously for all randomized solvers in this distributional bound.
In contrast, §§8-10 analyze FG-wired instances and prove per-instance
deterministic hardness that applies to any fixed run (see Scope note at
§8 and §9).}

\textbf{Part B: Lower bound via arity-bounded analysis.} Through
arity-bounded information flow (§5.5.1), each k-tape TM step can acquire
at most B bits (with B :=
k\(\lceil\)log\(_2\)\textbar{}\(\Gamma\)\textbar{}\(\rceil\)), directly
yielding time \(\geq\) 2\^{}\(\Omega\)(\(\lambda\)) where \(\lambda\) is
the min-cut residual. Appendix C.4 provides calibrations and mechanisms
achieving tight bounds on concrete instance families.

\textbf{Note:} This bound is model-agnostic. Each first-use read costs
\(\Omega\)(1) time in any standard model (TM head move, RAM word access,
etc.), so the information-theoretic bound translates directly to time.

Keyed cross-attempt reuse does not reduce the 2\^{}(\(\lambda\)(A,x))
factor: any correct reuse must be keyed by distinct prefixes
y\_\{\(\leq\)v\} (or Seed\_\{succ(v)\}), which are exactly what the
semantic counters (tries/keys) count (see §4.4 Keyedness; proofs in
Appendix B).

The bit-level analysis accounts for all paths automatically through the
exploration structure.

\textbf{Streaming corollary (exponential time).} If
\(\lambda\)=\(\Theta\)(n), then any \textbf{classical streaming} solver
(any number of passes p\(\geq\)1 with subexponential memory) satisfies

time\_M \(\geq\) 2\^{}(\(\lambda\)(A,x) - log\(_2\) Alt(C*)) ·
\(\lambda\)(A,x) (if Alt(C*) \(\leq\) poly(n\_core), then \(\geq\)
2\^{}(\(\lambda\)(A,x))/poly(n\_core)·\(\lambda\)(A,x)).

\emph{Proof.} Part A: Distributional tries bound +
Yao\textquotesingle{}s principle. Part B: Bit-level analysis (via
Kraft\textquotesingle{}s inequality on a semantic decision tree with
2\^{}\(\lambda\) leaves and average leaf depth \(\geq\) \(\lambda\))
shows total first-use reads \(\geq\) 2\^{}(\(\lambda\)(A,x)) ·
\(\lambda\)(A,x); each read costs \(\Omega\)(1) time in any standard
model.

\emph{(This converts the semantic "tries" bound to \textbf{time} under
worst-case instances.)}

\begin{center}\rule{0.5\linewidth}{0.5pt}\end{center}

\textbf{Lemma 4.1-H (Read-or-Hybrid; per-node law).} For a
\textbf{hybrid} solver that may simultaneously branch, memoize, maintain
an OBDD, and derive resolution clauses, let \ensuremath{\mathcal{H}}\_v be its hybrid potential
at node v. Then

q\_v + \ensuremath{\mathcal{H}}\_v \(\geq\) R\_v

\emph{Proof.} Each live world induces a \textbf{composite tuple}
\(\tau\)\_v of artifact identifiers (e.g., branch id, state key
sufficient for correctness, OBDD node at the relevant level, and a
clause profile sufficient to separate cases). Correctness requirements
imply \textbf{injectivity}: two distinct worlds cannot share the same
\(\tau\)\_v without resolution. Hence

tuples\_v \(\geq\) A\_v = 2\^{}(R\_v\(-\)q\_v)

The number of realized tuples is at most the product of component
cardinalities, so:

log\(_2\) tuples\_v \(\leq\) log\(_2\) branches\_v + log\(_2\) keys\_v +
log\(_2\) width\_OBDD\_v + log\(_2\) cases\_RES\_v = \ensuremath{\mathcal{H}}\_v

Combining gives: \textbf{q\_v + \ensuremath{\mathcal{H}}\_v \(\geq\) R\_v}. \(\blacksquare\)

\textbf{Lemma 4.1b (Direct-Sum Anti-Compression).} Within a single node
v, the \textbf{R\_v independent bits} of x\_v cannot be compressed
without violating correctness requirements across any paradigm:

\begin{itemize}
\tightlist
\item
  \textbf{EO/backtracking.} Branch information for \textbf{independent}
  ambiguity blocks cannot be encoded with fewer than the product of
  their distinguishable artifacts; one branch step merges at most
  \textbf{one} block.
\item
  \textbf{DP.} Distinct values of the unresolved bits induce
  \textbf{distinct state keys} sufficient for correctness; merging them
  would be incorrect, so one state creation merges at most \textbf{one}
  block.
\item
  \textbf{OBDD.} The expander parity gate ensures that for any variable
  order, width blow-up occurs when avoiding resolution.
\item
  \textbf{Resolution/CDCL.} Excluding k independent distinguishable
  artifacts forces \textbf{width \(\Omega\)(k)}; widths add across
  independent bundles, and size follows via width\(\to\)size tradeoffs.
\end{itemize}

Consequently, \textbf{within node v} the solver\textquotesingle{}s costs
\textbf{add over ambiguity blocks}; combined with \textbf{Dependency},
costs \textbf{multiply across nodes}.

\subparagraph{7.3.10 Direct Time Bound via Arity-Bounded Information
Flow}\label{7310-direct-time-bound-via-arity-bounded-information-flow}

The arity-bounded proof (§5.5.1) directly establishes exponential time
lower bounds for all deterministic k-tape TMs on DAG computations with
these information requirements.

\textbf{Time bounds via arity:} The arity-bounded information flow
(§5.5.1) directly yields exponential time bounds for all deterministic
k-tape TMs on such DAG computations. See \textbf{Appendix C.4.3 (Theorem
C.4.2)} for conversions and tightness mechanisms on concrete instances.

\begin{center}\rule{0.5\linewidth}{0.5pt}\end{center}

\subparagraph{7.3.11 Depth/Span via Round-Credit
SCL}\label{7311-depthspan-via-round-credit-scl}

Here \(\beta\)* denotes an upper bound on the per-round legitimate
cross-cut credit across the bottleneck cut C* (sum over processors
within a round).

The SCL captures both information bottlenecks (\(\lambda\)) and
resolution bandwidth (\(\beta\)*) - the rate at which verifiable
progress can reduce residual across a cut. In parallel models, this
two-dimensional structure determines both depth and span.

\textbf{Round-credit budget (definition).} For a parallel model and cut
C, let \(\Delta\)\_legit,t(C) be the maximum verifier-auditable
reduction in residual across C that a single round t can legally effect.
We count as "legitimate credit" both (i) first-use reads credited by RWA
and (ii) priced structural updates that shrink the feasible-world set
across C (e.g., NF\_C growth or WorldCommit\_C refutation; cf. Appendix
C.2/C.4). Define \(\beta\)\_t(C) \(\geq\) \(\Delta\)\_legit,t(C) as any
model-specific upper bound, and \(\beta\)*(C) := sup\_t \(\beta\)\_t(C).
Write \(\beta\)* := \(\beta\)*(C*) for the min-cut.

\textbf{Depth lower bound (Round-SCL).} For any correct parallel
computation that finishes in D synchronous rounds,

\textbf{D \(\geq\) \(\lceil\)\(\lambda\)/\(\beta\)*\(\rceil\)}

where \(\lambda\) = \(\Lambda\)(C*) \(-\) Q(C*) is the residual on the
min-cut C*.

\emph{Proof sketch.} By definition, each round effects at most
\(\beta\)\_t(C*) legitimate credit across C*. By acceptance purity
(ACC-1; Appendix C.2) and monotone sufficiency of NF\_C (CDT-1/2),
completing a run requires total legitimate credit \(\geq\) \(\lambda\)
across some s-t cut; at the min-cut we need at least \(\lambda\).
Summing over D rounds gives \(\Sigma\)\_\{t\(\leq\)D\} \(\beta\)\_t(C*)
\(\geq\) \(\lambda\) while \(\Sigma\)\_\{t\(\leq\)D\} \(\beta\)\_t(C*)
\(\leq\) D·\(\beta\)*, hence D \(\geq\) \(\lambda\)/\(\beta\)*. \ensuremath{\square}

\textbf{Work lower bound (aggregate processing).} Independently of
parallelism,

W = P · D \(\geq\) \(\Omega\)(2\^{}\(\lambda\))

because at least 2\^{}\(\lambda\) seed-consistent worlds must be
distinguished or refuted in aggregate (Theorem 7.A across cuts; Segment
Counting/FG pricing for refutations, Appendix C.2/C.1.1).

\textbf{Brent-style corollary (time on P processors).} Let T\_P be
parallel time with P processors and per-round credit \(\beta\)*. Then

\textbf{T\_P \(\geq\) max\{2\^{}\(\lambda\)/P, \(\lambda\)/\(\beta\)*\}}

(first term from Work; second from Depth).

\textbf{Specializations:}

\begin{itemize}
\item
  \textbf{Sequential k-tape TM.} \(\beta\)* = B = O(k
  log\textbar{}\(\Gamma\)\textbar) = O(1). Depth and time coincide: T
  \(\geq\) 2\^{}\(\lambda\), recovering our main bounds.
\item
  \textbf{PRAM (idealized).} In one round, at most \(\Theta\)(P·B)
  legitimate bits can be credited across C*. Thus D \(\geq\)
  \(\lambda\)/(\(\Theta\)(P)B) and T\_P \(\geq\)
  max\{2\^{}\(\lambda\)/P, \(\lambda\)/(\(\Theta\)(P)B)\}. With P =
  2\^{}\(\lambda\) processors, depth D can be O(1), but Work remains
  \(\Omega\)(2\^{}\(\lambda\)).
\item
  \textbf{Boolean circuits.} If at most c(C*) wires cross C* per level,
  then \(\beta\)* = c(C*). Thus depth \(\geq\) \(\lambda\)/c(C*) and
  size \(\geq\) 2\^{}\(\lambda\).
\item
  \textbf{Distributed models.} With per-round bandwidth \(\beta\)*
  across the partition, rounds \(\geq\) \(\lambda\)/\(\beta\)*.
\end{itemize}

\textbf{Key insight.} Parallelism widens the instantaneous funnel
(bigger \(\beta\)*), reducing depth/time, but cannot reduce total work
\(\Omega\)(2\^{}\(\lambda\)). SCL remains unchanged; only the pricing
varies with the model.

\textbf{Relation to P-completeness.} This formulation explains why some
problems in P resist parallelization despite low \(\lambda\). For
example, Circuit Value Problem has \(\lambda\) = 0 (deterministic
evaluation) but structural dependencies force \(\beta\)* = O(1) even
with many processors, yielding depth = \(\Theta\)(circuit\_depth). The
round-credit extension thus captures both information complexity
(\(\lambda\)) and dependency complexity (\(\lambda\)/\(\beta\)*).

\textbf{Corollary (Depth via round-credit).} For any correct computation
on our explicit overlay and any parallel model with per-round
legitimate-credit budget \(\beta\)*, the span satisfies

\textbf{D \(\geq\) \(\lceil\)\(\lambda\)/\(\beta\)*\(\rceil\)}

Thus if the min-cut residual is \(\lambda\) = O(log n), then \textbf{D
is polynomial} (and in the sequential model, \textbf{time is
polynomial}). In particular, when \(\lambda\) = 0, span lower bounds
vanish (D \(\geq\) 0) and a trivial upper bound is \textbf{D \(\leq\)
depth(G) \(\leq\) poly(n\_core)} by topological evaluation.

\textbf{§7.3 Summary:} Established SCL \(\to\) paradigm bounds via
projection templates (7.3.1-7.3.3): backtracking (tree size), DP (keys),
OBDD (width), Resolution (proof size), TM (segments). Derived time
bounds: arity-bounded flow (7.3.10: time \(\geq\) 2\^{}\(\lambda\)/B
steps), segment counting (7.3.7-7.3.9: m\_seg \(\geq\) 2\^{}(\(\rho\)-s)
rollbacks; Appendix C.2), depth/span (7.3.11: D \(\geq\)
\(\lambda\)/\(\beta\)*). Confirmed solver tricks don\textquotesingle{}t
compress \(\lambda\) (7.3.6). Used in: §8 (per-instance bounds Theorem
8.A), Appendix C (FG wiring, tight single-run analysis).

\begin{center}\rule{0.5\linewidth}{0.5pt}\end{center}

\paragraph{\texorpdfstring{7.4 Dependency \(\Rightarrow\) multiplicative
growth}{7.4 Dependency \textbackslash Rightarrow multiplicative growth}}\label{74-dependency--multiplicative-growth}

\textbf{Theorem (Multiplication across nodes).} \emph{Refs:} Per-node
SCL (Theorem 7.A) and dependency composition; see Appendix J for min-cut
composition. Under Dependency, the size of the computation artifact
obeys

artifact size \(\geq\) \(\prod\)\emph{\{v on chain\} Alt\_v \(\geq\)
\(\prod\)}\{v on chain\} 2\^{}(R\_v - q\_v) = 2\^{}(\(\sum\)\_\{v on
chain\}(R\_v - q\_v)).

Note: Nodes with R\_v = 0 contribute Alt\_v = 1 and do not affect the
product.

\emph{Proof.} By induction on nodes: Because successor artifacts are
determined by their immediate predecessors (Dependency), for any
root-to-leaf evolution we have Alt\_\{\(\leq\)succ(v)\} \(\geq\)
Alt\_\{\(\leq\)v\} · Alt\_\{succ(v)\}. Distinct y\_v values induce
distinct successors, so distinguishable artifacts from different node-v
branches cannot share successor artifacts. Thus the numbers of live
distinguishable artifacts maintained at each node multiply along any
root-to-leaf evolution, yielding the product bound; the inequality
Alt\_v \(\geq\) 2\^{}(R\_v\(-\)q\_v) is Theorem 7.A. \(\blacksquare\)

\emph{Emphasis.} \textbf{Keyedness \(\Rightarrow\) no sharing across
different seeds \(\Rightarrow\) multiplication (not addition)} along
paths.

\textbf{Hybrid multiplication (corollary).} For hybrids, let tuples\_v
be the number of realized composite tuples at node v. By the same
keyed-dependency argument, tuples across nodes \textbf{multiply} along
any path. Since tuples\_v \(\geq\) 2\^{}(R\_v\(-\)q\_v) per node, the
total number of composite tuples is at least 2\^{}(\(\lambda\)(A,x)).

\textbf{Block-ownership axiom (semantic).} Each node\textquotesingle{}s
ambiguity blocks own disjoint designated components. Thus one operation
can reduce distinguishable artifacts in \textbf{at most one} block.
Combined with Dependency (successor artifacts determined by
predecessors), this forces \textbf{multiplication} across nodes rather
than shared amortization.

\textbf{Corollary 7.4.1 (Exponential under repeated non-resolution).} If
there is a set T of nodes with q\_v \(\leq\) R\_v \(-\) 1 for all v
\(\in\) T, then artifact size \(\geq\) 2\^{}\textbar{}T\textbar{}. In
the flat profile, \(\lambda\)\_base=\(\Theta\)(n) yields artifact size
\(\geq\) 2\^{}\(\Theta\)(n).

\textbf{Corollary 7.4.2 (Additive under verification).} If the solver
\textbf{verifies} every node (q\_v=R\_v for all v), then Alt\_v=1 and
Theorem 4.2 yields a unit product. The overall work reduces to learning
\(\Sigma\)\_v R\_v forced facts (the one-path polynomial verifier).

\begin{center}\rule{0.5\linewidth}{0.5pt}\end{center}

\subparagraph{7.4.1 Adaptive Resolution Example (Parametric
Nature)}\label{741-adaptive-resolution-example-parametric-nature}

\textbf{Example (Adaptive Resolution Strategy).} Consider the strategy
q\_v = R\_v - \ensuremath{\sqrt{\text{depth}}}(v) across D = \(\Theta\)(log n) depth:

\begin{itemize}
\tightlist
\item
  Per node: Resolve R\_v - \ensuremath{\sqrt{\text{depth}}}(v) facts, branch on \ensuremath{\sqrt{\text{depth}}}(v) facts
\item
  Distinct states per node: 2\^{}\ensuremath{\sqrt{\text{depth}}}(v)
\item
  Total residual: \(\lambda\)(A,x) = \(\Theta\)((log n)\^{}1.5)
\item
  Artifact size: 2\^{}(\(\Theta\)((log n)\^{}1.5)) =
  \textbf{quasi-polynomial}
\end{itemize}

This demonstrates the parametric tradeoff: leaving only \ensuremath{\sqrt{\text{depth}}}(v) =
O(\ensuremath{\sqrt{\log}} n) facts unresolved per node yields quasi-polynomial complexity,
exactly as our bound predicts. Sub-exponential artifacts arise from
\(\Theta\)(n\^{}\(\varepsilon\)) total unresolved bits, while
exponential artifacts require \(\Sigma\)(R\_v - q\_v) = \(\Theta\)(n)
unresolved bits, which occurs only in pure search mode.

\paragraph{7.5 Three Operational Routes (two-dimensional
constraint)}\label{75-three-operational-routes-two-dimensional-constraint}

\textbf{Core Insight:} At each node v, any classical algorithm faces
exactly three options for handling 2\^{}(R\_v-q\_v) genuinely distinct
computational states. This is not a limitation - it is a mathematical
consequence of the DAG\textquotesingle{}s information structure.

These three options are operational strategies for satisfying the
two-dimensional constraint q + \(\Phi\) \(\geq\) R:

\textbf{Dimension 1: Storage (Space)}

\begin{itemize}
\tightlist
\item
  \emph{Option}: Maintain multiple simultaneous artifact states
\item
  \emph{Measure}: log\(_2\)(states\_v), where states\_v = number of
  non-mergeable artifact classes
\item
  \emph{L*\textquotesingle{}s enforcement}: Keyedness (§4.4) prevents
  merging artifacts with different seed histories - different Seed\_v
  values produce different designated addresses F\_overlay(Seed\_v;
  j,\(\ell\)), so merging would compute incorrect addresses and violate
  verifiability (Lemma 7.I; §7.2.1)
\item
  \emph{Consequence}: Must maintain \(\geq\) log\(_2\)(2\^{}(R\_v-q\_v))
  = R\_v - q\_v bits of distinguishable state information
\end{itemize}

\textbf{Dimension 2: Resolution (Time-Forward)}

\begin{itemize}
\tightlist
\item
  \emph{Option}: Learn which state is correct by explicit reading
\item
  \emph{Measure}: q\_v\^{}(read) = bits resolved via designated reads at
  node v
\item
  \emph{L*\textquotesingle{}s enforcement}:

  \begin{itemize}
  \tightlist
  \item
    Emergence (A3; §6.2.5): R\_v fresh bits at v cannot be inferred -
    must read explicitly from designated addresses
  \item
    Bandwidth (Lemma 5.5.1): Each step acquires \(\leq\) B = O(1) bits
    of fresh information
  \end{itemize}
\item
  \emph{Consequence}: Reading (R\_v - q\_v) bits requires \(\geq\) (R\_v
  - q\_v)/B sequential steps
\end{itemize}

\textbf{Dimension 3: Elimination (Time-Backward)}

\begin{itemize}
\tightlist
\item
  \emph{Option}: Prune wrong states by testing and rejection
\item
  \emph{Measure}: q\_v\^{}(elim) = bits eliminated via testing wrong
  candidates
\item
  \emph{L*\textquotesingle{}s enforcement}:

  \begin{itemize}
  \tightlist
  \item
    Per-node antagonism (§6.1.1.A): Each wrong candidate eliminates
    \(\leq\)1 bit (factoring-style; no cascade)
  \item
    CDT (Lemma CDT-1\textquotesingle{}; Appendix C): Semantic progress
    requires designated work - no "free" inference
  \end{itemize}
\item
  \emph{Consequence}: Eliminating wrong possibilities among
  2\^{}(R\_v-q\_v) candidates requires testing exponentially many
  (cannot bulk-prune)
\end{itemize}

\textbf{The Conservation Law (Theorem 7.A, refined form):}

For any correct algorithm at node v, the total information accounted for
- whether through explicit reading (q\_v\^{}(read)), elimination via
testing (q\_v\^{}(elim)), or maintaining as stored states
(log\(_2\)(states\_v)) - must satisfy:

(q\_v\^{}(read) + q\_v\^{}(elim)) + log\(_2\)(states\_v) \(\geq\) R\_v

\emph{Intuition}: Total resolved (via reading OR testing) plus
maintained (as distinguishable artifacts) must cover the total required.
The algorithm may freely trade off between reading, testing, and
storing, but cannot escape the total requirement R\_v.

\textbf{Orthogonality of Dimensions:}

These three dimensions attack the problem from independent directions:

\begin{itemize}
\tightlist
\item
  \textbf{Fast resolution} (Dimension 2) does not reduce elimination
  costs (Dimension 3) - high read bandwidth does not help determine
  which of the 2\^{}(R\_v-q\_v) possibilities to pursue; you must still
  explore/test candidates
\item
  \textbf{Fast elimination} (Dimension 3) does not reduce resolution
  costs (Dimension 2) - efficiently pruning wrongs does not tell you
  which remaining candidate is correct; you must still read to resolve
\item
  \textbf{Efficient storage} (Dimension 1) does not reduce either time
  dimension - maintaining fewer artifacts does not make reading or
  testing faster
\end{itemize}

L* enforces exponential cost in ALL THREE dimensions simultaneously,
creating a "triple trap" where optimizing any single dimension leaves
the others exponential.

\textbf{Why No Fourth Option:}

L*\textquotesingle{}s structure ensures that these 2\^{}(R\_v-q\_v)
states lead to different outcomes. For Structural OWF construction (§9),
\textbf{every FG-wired instance} satisfies per-instance deterministic
bounds (Theorem 8.A). Any attempt to bypass all three dimensions
simultaneously would merge distinct computational paths, causing
incorrect results (different seeds \(\to\) wrong designated addresses
\(\to\) verifier detects error; see §1.1). The mathematics does not
allow it - \textbf{A1-A5} properties make this structurally impossible.

\textbf{Lemma 7.T (Three-dimensional factorization).} Within a try at
node v, let branches\_v denote the number of live semantic branches
(root-to-v decision paths consistent with the current transcript), and
let states\_v denote the number of live seed-consistent artifact classes
at v (Keyedness; distinct seeds/keys cannot merge). Then the number of
simultaneously distinguishable artifacts at v satisfies

Alt\_v \(\geq\) branches\_v · states\_v,

and thus by SCL

q\_v + log\(_2\)(branches\_v) + log\(_2\)(states\_v) \(\geq\) R\_v.

\emph{Proof.} Tag each live artifact at v by the pair (branch-id,
key-id), where branch-id identifies the root-to-v choice path within the
current try, and key-id is the seed-consistent class label at v. Two
artifacts that differ on either component are future-distinguishable:
distinct branches differ on some earlier choice; distinct keys induce
distinct future seed chains and addresses (Keyedness), and cannot be
merged without violating injectivity/verifiability. Hence the pairs are
all distinct, yielding at least branches\_v · states\_v artifacts.
Taking logs and applying Theorem 7.A gives the inequality.
\(\blacksquare\)

\textbf{Cross-Reference:}

See §6.1.1.1 for how the five-dimensional antagonisms (A-E) map to these
three fundamental obstacles, and Appendix C (lane analysis) for how the
Two Lanes (Restart vs. Single-Run) represent different strategies that
fail against different dimensions.

\begin{center}\rule{0.5\linewidth}{0.5pt}\end{center}

\textbf{Quick Reference: Three-Dimensional Framework \(\to\) Technical
Mechanisms}

\emph{For readers seeking specific technical details of each dimension:}

\textbf{Dimension 1: Unmergeable States (Space)}

\begin{itemize}
\tightlist
\item
  \textbf{Property}: Keyedness (§4.4) - artifacts correct iff
  seed-consistent
\item
  \textbf{Construction}: Injective Enc (§6.2.7) - different histories
  \(\to\) different seeds
\item
  \textbf{Enforcement}: Content-addressing (§6.2.3) - different seeds
  \(\to\) different addresses F\_overlay(Seed\_v; j,\(\ell\))
\item
  \textbf{SCL Proof}: Theorem 7.A + Lemma 7.I (§7.2.1) - collision
  \(\to\) error
\item
  \textbf{Consequence}: Alt\_v \(\geq\) 2\^{}(R\_v-q\_v) distinguishable
  artifacts at node v
\end{itemize}

\textbf{Dimension 2: Unresolvable Information (Time-Resolution)}

\begin{itemize}
\tightlist
\item
  \textbf{Property}: Emergence (A3; §6.2.5) - R\_v fresh bits at each
  node v
\item
  \textbf{Proof}: Lemma 6.1 (distributional independence until read)
\item
  \textbf{Bandwidth}: Lemma 5.5.1 - \(\leq\) B = O(1) bits per step
\item
  \textbf{Segments}: Segment Counting (Appendix C.2) - m\_seg \(\geq\)
  2\^{}(\(\rho\)-s) rollback segments
\item
  \textbf{FG Enforcement}: Frontier-Gate (Appendix C.1) -
  \(\Omega\)(n/W\_min) steps per segment
\item
  \textbf{Consequence}: Time \(\geq\) 2\^{}(\(\rho\)-s) ·
  \(\Omega\)(n/W\_min) in single-run lane
\end{itemize}

\textbf{Dimension 3: Uneliminable Candidates (Time-Elimination)}

\begin{itemize}
\tightlist
\item
  \textbf{Property}: Per-node antagonism (§6.1.1.A) - reject eliminates
  \(\leq\)1 bit
\item
  \textbf{Construction}: Full-rank y\_v = H\_v x\_v (§6.2.5) -
  factoring-style
\item
  \textbf{CDT}: Lemma CDT-1\textquotesingle{} (Appendix C) - no unbacked
  progress
\item
  \textbf{Tries Bound}: Across-tries inequality (Lemma 7.R; Appendix
  C.4.2)
\item
  \textbf{Extraction}: Yao\textquotesingle{}s principle (Appendix C) -
  distributional \(\to\) worst-case
\item
  \textbf{Consequence}: \ensuremath{\mathbb{E}}{[}tries{]} \(\geq\) 2\^{}(\(\Delta\)(C*))
  \(\to\) Time \(\geq\) 2\^{}(\(\Delta\)(C*))/B in restart lane
\end{itemize}

\textbf{Cross-Cutting Mechanisms:}

\begin{itemize}
\tightlist
\item
  \textbf{Composition}: Cross-node antagonism (B; §6.1.1) + Disjoint
  pools (A1; Lemma 6.5) \(\to\) residuals add across cuts
\item
  \textbf{Lane Exhaustivity}: Lane analysis (Appendix C) \(\to\) all
  strategies fall into restart or single-run
\item
  \textbf{No Bypass}: Theorem 10.4.1-BYP (§10.4.1) \(\to\) canonical
  witness requires overlay engagement
\item
  \textbf{Per-Instance Bounds}: Theorem 8.A \(\to\) every FG-wired
  instance hard (OWF path)
\end{itemize}

\emph{Note}: All mechanisms work in concert. The three dimensions are
independent (overcoming one does not help with others) but enforced by
overlapping mechanisms from the five antagonisms (A-E; §6.1.1).

\begin{center}\rule{0.5\linewidth}{0.5pt}\end{center}

\paragraph{7.6 Paradigm-Specific
Instantiations}\label{76-paradigm-specific-instantiations}

Having rigorously established SCL (§7.2.1), we now instantiate it into
multiple paradigms:

\textbf{Core paradigm manifestations:}

\begin{itemize}
\tightlist
\item
  \textbf{EO / Backtracking:} maintained artifacts = \textbf{branches}
  \(\Rightarrow\) \textbf{Resolution/Elimination or Storage-branches};
  tree size \(\geq\) 2\^{}(\(\lambda\)(A,x)).
\item
  \textbf{DP / Memoization:} maintained artifacts = \textbf{distinct
  keys} \(\Rightarrow\) \textbf{Resolution/Elimination or
  Storage-table}; key count \(\geq\) 2\^{}(\(\lambda\)(A,x)).
\item
  \textbf{Resolution / CDCL:} maintained artifacts = \textbf{cases}
  encoded in \textbf{clause width} \(\Rightarrow\)
  \textbf{Resolution/Elimination or Storage-width} per node; widths
  \textbf{add} and standard width\(\to\)size yields exponential proof
  size.
\item
  \textbf{OBDD (order-robust):} maintained artifacts = \textbf{width} at
  some cut \(\Rightarrow\) \textbf{Resolution/Elimination or
  Storage-width}; widths multiply across nodes (requires
  expander-parity; §6.2.11; {[}HOO06{]}, {[}WEG00{]}).
\end{itemize}

\textbf{Uses:} Theorem 7.A + projection lemmas (§5.2) for each paradigm;
width\(\to\)size (Appendix G) and OBDD width (Appendix B.1) where
applicable.

\textbf{Additional semantic characterizations (§4.4):}

\begin{itemize}
\tightlist
\item
  \textbf{Search with Rollbacks:} maintained artifacts =
  \textbf{distinct rollback keys} (resolution prefixes) \(\Rightarrow\)
  \textbf{read-or-rollbacks}; on worst-case instances, \ensuremath{\mathbb{E}}{[}tries{]}
  \(\geq\) 2\^{}(\(\Delta\)(C*)) (across-tries; Lemma 7.R).
\item
  \textbf{General Branching Programs:} maintained artifacts =
  \textbf{residual subfunctions} \(\Rightarrow\)
  \textbf{read-or-subfunction}; size \(\geq\)
  2\^{}(\(\Omega\)(\(\Sigma\)(R\_v - q\_v))).
\item
  \textbf{Algebraic Proof Systems (PC/CP/SoS):} \textbf{read-or-degree}
  \(\Rightarrow\) degree \(\geq\) \(\Omega\)(\(\Sigma\)(R\_v - q\_v))
  \textbf{unconditionally}; size \(\geq\) g(\(\Sigma\)(R\_v - q\_v))
  requires \textbf{Assumption A} (standard degree\(\to\)size bridge).
\end{itemize}

We also record an additive (strong polynomial) per-path lemma for
\textbf{streaming/oblivious verifiers}: with small workspace S, each
pass adds \(\leq\)S resolutions, so per node \textbf{passes \(\geq\)
\(\lceil\)(R\_v - q\_v)/S\(\rceil\)} and \textbf{passes add} across
nodes.

\begin{center}\rule{0.5\linewidth}{0.5pt}\end{center}

\paragraph{7.7 Robustness to
Randomization}\label{77-robustness-to-randomization}

\textbf{Structural OWF construction (§9):} SCL bounds are
\textbf{per-run} and \textbf{information-theoretic} - they apply to any
deterministic computational run regardless of how it arose. For
randomized algorithms, \textbf{coin-fixing} (Yao\textquotesingle{}s
min-max principle {[}YAO77{]}) establishes that per-instance bounds
extend: any fixed coin string yields a deterministic run subject to SCL;
if instance x* requires time T on every fixed run, then randomization
cannot reduce this. Hence per-instance deterministic bounds (Theorem
8.A) directly imply bounds for randomized algorithms.

\begin{center}\rule{0.5\linewidth}{0.5pt}\end{center}

\subparagraph{7.7.1 Example: How L* Instantiates
SCL}\label{771-example-how-l-instantiates-scl}

To see how a specific problem triggers SCL, consider the language L*
from §6:

\textbf{L* \(\to\) SCL Mapping:}

\textbf{L* Properties to SCL Parameters}

\begin{itemize}
\tightlist
\item
  \textbf{Overlay DAG structure}: SCL Parameter: G = (V, E); Concrete
  Value: Depth O(log n\_core) DAG
\item
  \textbf{Emergence rank at v}: SCL Parameter: R\_v; Concrete Value:
  \(\kappa\) log n bits
\item
  \textbf{GREQ commitments}: SCL Parameter: q\_v; Concrete Value:
  \(\leq\) B per algorithm step
\item
  \textbf{Disjoint salt pools}: SCL Parameter: Independence; Concrete
  Value: No pre-derivation possible
\item
  \textbf{Seed chain dependencies}: SCL Parameter: DAG edges; Concrete
  Value: Enc(v\textbar{}\textbar{}parents\textbar{}\textbar{}...)
  enforces order
\item
  \textbf{Min-cut capacity}: SCL Parameter: \(\lambda\)\_base; Concrete
  Value: \(\Theta\)(log² n) for QP-sharp profile
\end{itemize}

\textbf{Application of Theorem 7.A:}

\begin{enumerate}
\def\labelenumi{\arabic{enumi}.}
\tightlist
\item
  L* creates a DAG with information requirements R\_v at each node
\item
  Algorithms can resolve at most q\_v bits per step at v
\item
  The residual R\_v - q\_v must manifest as distinguishable artifacts
\item
  Min-cut analysis yields \(\lambda\)\_base = \(\Theta\)(log² n)
\item
  Therefore: In L* we can prove that \(\geq\) 2\^{}(\(\Theta\)(log² n))
  = n\^{}(\(\Theta\)(log n\_core)) artifacts are mathematically
  necessary
\end{enumerate}

\textbf{Key Point:} The full L* construction (§6) ensures these
parameters through axioms A1-A5 (Hermeticity, Injectivity, Emergence,
Closure, Dependency). Here we only need that L* satisfies the abstract
SCL interface - the specific mechanisms (gates, selectors, encoding
functions) are implementation details that establish the interface
requirements.

\textbf{Uses:} Theorem 7.A applied to L*\textquotesingle{}s parameter
instantiation.

\begin{center}\rule{0.5\linewidth}{0.5pt}\end{center}

\paragraph{7.8 Conclusion}\label{78-conclusion}

\textbf{What We\textquotesingle{}ve Proven:} The Semantic Conservation
Law (Theorem 7.A, rigorously proven in §7.2.1) is a mathematical theorem
about computation on DAGs with information requirements. When applied to
our constructed NP-complete language L*, it shows every classical
algorithm must satisfy q + \(\Phi\) \(\geq\) R at each node.

Perspective.

\begin{itemize}
\tightlist
\item
  \textbf{Old View:} "Algorithms lack sufficient power to solve
  NP-complete problems efficiently"
\item
  \textbf{New View:} "L*\textquotesingle{}s structure creates
  unavoidable multiplicative requirements"
\end{itemize}

Next steps: §8 proves per-instance deterministic bounds for every
FG-wired instance, §9 packages these bounds into an unconditional
Structural OWF, and §10 applies the classical bridge to conclude P
\(\neq\) NP for classical uniform PPT.

Why this matters.

\begin{enumerate}
\def\labelenumi{\arabic{enumi}.}
\tightlist
\item
  \textbf{It\textquotesingle{}s Proven, Not Conjectured:} For Structural
  OWF construction (§9), \textbf{every FG-wired instance} exhibits
  deterministic hardness (Theorem 8.A).
\item
  \textbf{It\textquotesingle{}s General:} It applies to any classical
  algorithm, not just known paradigms
\item
  \textbf{It\textquotesingle{}s Structural:} The bounds arise from
  L*\textquotesingle{}s correctness requirements, not algorithmic
  limitations
\item
  \textbf{It\textquotesingle{}s Orthogonal to Shannon:} We measure
  semantic necessity, not information entropy
\item
  \textbf{It\textquotesingle{}s Information-Theoretic:} Algorithms need
  information to make progress; L* provides no information. The barrier
  operates on three levels: (a) \textbf{Bootstrapping problem} - must
  have information before you can get information (circular dependency:
  need \(\varphi\) to find \(\alpha\), need \(\alpha\) to compute seeds,
  need seeds to decode \(\varphi\)); (b) \textbf{No structural clues} -
  standard problems provide unit constraints, symmetries, frequencies
  that enable shortcuts, but L*\textquotesingle{}s seed-locked encoding
  removes ALL structural information; (c) \textbf{No useful feedback} -
  wrong guesses yield only "digest mismatch" with zero guidance about
  which variables to change or candidates to eliminate (WC-1: each test
  eliminates \(\leq\)1 of 2\^{}\(\lambda\) possibilities, forcing
  exhaustive search). This triple barrier transforms "we
  haven\textquotesingle{}t found a fast algorithm" into "fast algorithms
  are information-theoretically impossible."
\end{enumerate}

\textbf{The Bottom Line:} When facing 2\^{}(R\_v) possible computational
states at node v in L*, SCL imposes a two-dimensional constraint (q,
\(\Phi\)) that in practice yields three operational routes:
\textbf{Storage} (maintaining distinguishable states),
\textbf{Resolution} (learning correct answers via reading), and
\textbf{Elimination} (pruning wrong candidates via testing).
L*\textquotesingle{}s structure creates exponential barriers across
these routes simultaneously - not because we lack imagination, but
because \textbf{A1-A5 properties make this mathematically unavoidable
(compressing below 2\^{}(R\_v-q\_v) artifacts causes collisions \(\to\)
different seeds map to same artifact \(\to\) wrong addresses \(\to\)
errors; see §1.1). The exponential barrier is not algorithmic but
information-theoretic} - every technique that could enable shortcuts
(propagation, learning, pruning, heuristics) requires structural
information that L* does not provide. This structural inevitability
yields super-polynomial time for producing canonical witnesses. Combined
with Structural OWF construction (§9) \(\to\) classical bridge (§10),
this proves P \(\neq\) NP.

\begin{center}\rule{0.5\linewidth}{0.5pt}\end{center}

\subsubsection{8. Per-Instance Deterministic
Bounds}\label{8-per-instance-deterministic-bounds}

Section 7 proved the Semantic Conservation Law: for L*, A1-A5
mathematically imply q + \(\Phi\) \(\geq\) R at each node, and the
framework translates this into bounds depending on residual \(\lambda\).
Section 8 establishes the \textbf{per-instance deterministic lower
bounds} that form the foundation for the Structural OWF construction
(§9). For the lane dichotomy used in pricing time from \(\lambda\), see
also Appendix C, Lemma C.EXH (Lane Exhaustiveness).

\textbf{Purpose of Section 8:} We prove that \textbf{every FG-wired
instance has a deterministic witness-finding lower bound} that applies
to any fixed computational run (Theorem 8.A). This per-run bound is
schedule/coin-independent - a structural property of the instance, not
algorithmic behavior. Combined with coin-fixing (Yao), this extends to
randomized PPT adversaries, enabling Structural OWF security (§9.4).

Assumptions box (scope for §8):

\begin{itemize}
\tightlist
\item
  Model: deterministic k-tape TMs (randomized handled by coin-fixing);
  Hermeticity/no advice/oracles
\item
  \textbf{FG wiring precondition}: Instance has published S(P) gates and
  GREQ flags per §6.2.8 (Gate-Ledger) and gate-horizon budget S(n) =
  \(\Theta\)(\(\tau\)·\(\lambda\)\_base) per §6.2.9; detailed in
  Appendix C.1.1
\item
  Ledger/SNF: RWA credits first-use designated reads; per-segment
  pricing uses the SNF ledger (Appendix I, C)
\item
  Notation: n := n\_core; W\_min per §4.3; \(\lambda\)\_base per
  profile; \(\rho\),s as in §4.3/Appendix C
\end{itemize}

\textbf{Roadmap:}

\begin{itemize}
\tightlist
\item
  \textbf{§8.0}: Structural hardness (Theorem 8.0.SH) - hardness arises
  from L*\textquotesingle{}s properties (A1-A5), not algorithmic
  deficiencies
\item
  \textbf{§8.1}: \textbf{Per-Instance Deterministic Bounds (Theorem
  8.A)} - every FG-wired instance requires super-polynomial time for
  witness-finding on any fixed run
\item
  \textbf{§8.2}: Connection to Structural OWF construction (§9) - how
  deterministic bounds enable one-wayness via coin-fixing
\item
  \textbf{§8.3}: Solution multiplicity assumption - security holds when
  \#SAT(\(\varphi\)) = 2\^{}\{o(\(\lambda\))\} (satisfied by all
  standard CNF families)
\end{itemize}

\textbf{The proof architecture:} L*\textquotesingle{}s structural
properties (A1-A5) + FG wiring \(\to\) per-node SCL (§7) \(\to\)
cut-level residual \(\lambda\) \(\to\) \textbf{per-instance
deterministic time bound} (any fixed run). No distributional arguments
needed.

\textbf{Why this matters:} Theorem 8.A provides the \textbf{per-instance
hardness} that makes Structural OWF construction possible. Every sampled
instance f(r) has this deterministic bound, so coin-fixing any
successful randomized inverter yields a contradiction (§9.4).

\textbf{Notation (used locally in §8):} \(\rho\) = effective residual at
bottleneck cut; s = pre-final agreement; \(\lambda\)(A,x) =
run-dependent min-cut residual; \(\lambda\)\_base = instance/profile
residual (§4.3); W\_min = profile\textquotesingle{}s minimal window
parameter; C*(x) = run-dependent minimizing cut.

\paragraph{8.0 Structural Hardness is
Universal}\label{80-structural-hardness-is-universal}

\textbf{Core Insight:} The hardness of L* arises from its structural
properties - Hermeticity, Injectivity, Emergence, Closure, and
Dependency (A1-A5) - not from algorithm-specific tailoring. These
properties ensure that the conservation law q + \(\Phi\) \(\geq\) R
holds for \textbf{every algorithm on every FG-wired instance}, yielding
per-run deterministic bounds that apply to any fixed computational trace
(schedule/coin-independent).

\textbf{Theorem 8.0.SH (Structural Hardness).} For the language L*
constructed in §6, the lower bounds arise from L*\textquotesingle{}s
structural properties (A1-A5), not from algorithmic limitations.
Concretely:

\begin{enumerate}
\def\labelenumi{\arabic{enumi})}
\tightlist
\item
  Construction independence. A1-A5 are built into L* by the
  deterministic Karp reduction f: \(\varphi\) \(\to\) x* (§6),
  independent of any solver A.
\item
  Universal SCL. For every algorithm A and run on x*, the Semantic
  Conservation Law holds (Theorem 7.A; §7.2.1), yielding per-node q\_v +
  \(\Phi\)\_v \(\geq\) R\_v via Keyedness (A1 Hermeticity + A2
  Injectivity).
\item
  Artifact necessity. By Keyedness + Injectivity, Alt\_v \(\geq\)
  2\^{}(R\_v-q\_v) at each node (Lemma 7.I), so any correct solver must
  realize exponentially many representation-invariant, seed-consistent
  artifact classes when residual is large.
\item
  Per-run bounds. Artifact requirements apply to \textbf{any fixed
  computational trace} (schedule/coin-independent), translating to time
  via FG wiring (Appendix C), giving profile-tight bounds (e.g.,
  n\^{}(\(\Omega\)(log n\_core)) for QP-sharp; 2\^{}(\(\Omega\)(n)) for
  flat).
\end{enumerate}

Therefore, the lower bounds are a consequence of instance-side
structural obligations (A1-A5 + FG) rather than deficiencies of
algorithms. These \textbf{per-run, information-theoretic bounds} extend
to randomized algorithms via coin-fixing (§9.4). \(\blacksquare\)

\textbf{Implication.} The P \(\neq\) NP separation (§9-§10) is
structural and unconditional: L*\textquotesingle{}s construction forces
super-polynomial witness-finding (per-instance deterministic bounds,
proven in §8.1) while verification remains polynomial. Structural OWF
security against randomized PPT follows via coin-fixing.

\textbf{Lemma 8.0 (Interface satisfaction).} The language L* from §6
satisfies the abstract SCL interface:

\begin{itemize}
\tightlist
\item
  Information requirements R\_v per node (Emergence, Lemma 6.1)
\item
  Receiving-Window Attribution of resolved bits q\_v (RWA)
\item
  No pre-derivation of unresolved coordinates (instance-side:
  Hermeticity + disjoint pools; not a probabilistic assumption)
\item
  DAG dependencies propagating requirements
\end{itemize}

Therefore Alt\_v \(\geq\) 2\^{}(R\_v-q\_v) at each node and Theorems
7.A, 7.A.1 apply to L*. The complete derivation from
L*\textquotesingle{}s construction to SCL is shown in §7.2.1
(Consolidated SCL Theorem). \(\blacksquare\)

Scope note (FG enabled). For §§8-10 we build L* with FG wiring (§6.2.8;
Appendix C.1.1), so the per-instance lower bounds are deterministic (no
probabilistic arguments; bounds apply to any fixed computational run).

\paragraph{8.1 Per-Instance Deterministic Bounds (Foundation for
OWF)}\label{81-per-instance-deterministic-bounds-foundation-for-owf}

This subsection establishes the core theorem underlying the Structural
OWF construction: \textbf{every FG-wired instance has a deterministic
witness-finding lower bound that applies to any fixed computational
run}.

\textbf{Progress Invariant (Segment Counting Mechanism).} A
\textbf{rollback segment} is a maximal contiguous subrun during which
the \textbf{resolution prefix} across the designated bottleneck cut C*
remains fixed - specifically, the tuple (seed-chain tag,
ConstraintDigest\_C, WorldCommit\_C) is constant (detailed definitions
in Appendix C.2.a).

\textbf{Invariant}: At the start of the final segment, there are at most
2\^{}\(\rho\) seed-consistent worlds across C*, where \(\rho\) :=
\(\Sigma\)\_\{v\(\in\)C*\}(R\_v\(-\)q\_v) is the effective residual at
that time (see §7.3.9; Appendix C.2):

\begin{itemize}
\tightlist
\item
  Each \textbf{non-accepting segment} can eliminate:

  \begin{itemize}
  \tightlist
  \item
    At most 2\^{}t worlds via \textbf{functional determination} t (new
    \(\Delta\)q bits within segment), PLUS
  \item
    \textbf{Exactly 1 additional world} via terminal refutation of
    WorldCommit\_C (Lemma C.2.3)
  \end{itemize}
\item
  \textbf{Acceptance requires} \(\leq\)1 world remaining (Lemma C.2.ACC)
\end{itemize}

\textbf{Consequence}: Shrinking from 2\^{}\(\rho\) worlds to 1 accepting
world, when total functional determination across all m\_seg\(-\)1
non-accepting segments is s bits, requires:

2\^{}(\(\rho\)-s) \(-\) (m\_seg\(-\)1) \(\leq\) 1 \(\to\) m\_seg
\(\geq\) 2\^{}(\(\rho\)-s)

\emph{(Proof: Appendix C.2: Segment Counting via CDT and WC lemmas,
including Lemma WC-1)}

\textbf{Theorem 8.A (Per-Instance Deterministic Bounds under FG).} For
\textbf{any L* instance with FG wiring} (§6.2.8-§6.2.9; Appendix C.1.1:
published S(P) gates and GREQ flags with gate-horizon budget) and
\textbf{any fixed computational run} (deterministic TM or randomized TM
with fixed coins), letting \(\rho\) be the effective residual on the
designated bottleneck cut and s the pre-final agreement:

time(x) \(\geq\) 2\^{}(\(\rho\)-s) · \(\Omega\)(n/W\_min)

By the gate-horizon budget (§6.2.9), s \(\leq\)
\(\Theta\)(\(\tau\)(n)·\(\lambda\)\_base(n)) \textbf{deterministically}
(independent of salt values and coins). By Lemma C.2.1, \(\rho\)
\(\geq\) \(\lambda\)\_base \(-\) s. Therefore

\(\rho\) \(-\) s \(\geq\) \(\lambda\)\_base \(-\) 2s \(\geq\) (1 \(-\)
\(\Theta\)(\(\tau\)))·\(\lambda\)\_base.

Instantiating profiles: for QP-sharp, 2\^{}(\(\rho\)-s) \(\geq\)
2\^{}(\(\Omega\)((log n)\^{}2)) = n\^{}(\(\Omega\)(log n)); for flat,
2\^{}(\(\rho\)-s) \(\geq\) 2\^{}(\(\Omega\)(n)). The
\(\Omega\)(n/W\_min) factor multiplies these; if W\_min \(\leq\) n
(typical), this is at least a constant factor.

Here n := n\_core. The machine model is a deterministic k-tape TM;
randomized TMs reduce to fixed-coin deterministic runs by coin-fixing.

\textbf{Crucially}: This bound applies to \textbf{any fixed
computational trace} - whether from deterministic choices or from fixing
random coins. The bound depends only on (x, computational\_trace), not
on how the trace was generated.

\emph{Proof.} By Lemma C.2 (Segment Counting), the number of
non-accepting rollback segments that progress the final chain obeys
m\_seg \(\geq\) 2\^{}(\(\rho\)-s). By Corollary C.1.1 (FG per-segment
baseline), each such segment incurs \(\Omega\)(n/W\_min) TM steps
regardless of pre-scanning or caching. Multiplying yields the claimed
bound. The s bound is deterministic from the published GREQ map
(§6.2.9), independent of salt values or coins. The bound applies to any
fixed computational trace because SCL (§7) is information-theoretic and
schedule-independent. \(\blacksquare\)

\textbf{Machine-verified proof chain:} Lean formalization across
multiple layers: segment counting
(\texttt{Layer3\_InformationBounds/SegmentReduction/SegmentReduction.lean},
theorem \texttt{refutation\_count\_exponential\_bound}), TM execution
semantics
(\texttt{Layer4\_Operational/TimeBridge/TMToExecutionPrefix.lean},
function \texttt{buildRunFromTMTrace}), and profile-specific time bounds
(\texttt{TMAdapterQP.lean} and \texttt{TMAdapterExponential.lean},
theorem \texttt{fg\_first\_commit\_time\_lower\_bound}). Complete proof
chain verified with zero \texttt{sorry} statements.

\textbf{Key Insight:} This \textbf{per-instance deterministic bound} is
the foundation for Structural OWF security. Every instance x* = f(r)
sampled in §9 has FG wiring, so Theorem 8.A applies. Coin-fixing (§9.4)
extends this to randomized PPT adversaries: any successful randomized
inverter yields a successful deterministic run, contradicting the bound.

\textbf{Falsifiability Note:} This theorem is falsifiable via
algorithmic test---exhibiting a uniform PPT algorithm achieving
\(\lambda\)(A,x*) = o(\(\lambda\)\_base) on FG-wired instances would
disprove this result. See §3.6 for explicit falsification criteria and
empirical test specifications.

\paragraph{8.2 Connection to OWF
Construction}\label{82-connection-to-owf-construction}

Theorem 8.A enables the Structural OWF construction (§9) through a
simple but powerful observation:

\textbf{How per-instance bounds yield Structural OWF security:}

\begin{enumerate}
\def\labelenumi{\arabic{enumi}.}
\item
  \textbf{Construction} (§9.2): Define f: r \(\mapsto\) x* where x* =
  Plant(\(\varphi\), r) with FG wiring (length-regular; see
  §9.2/Appendix O)

  \begin{itemize}
  \tightlist
  \item
    Every output x* = f(r) is an FG-wired L* instance
  \item
    Theorem 8.A applies to \textbf{every} x* (not just some fraction)
  \end{itemize}
\item
  \textbf{Per-instance hardness} (Theorem 8.A):

  \begin{itemize}
  \tightlist
  \item
    Every x* = f(r) has witness-finding bound \(\geq\) super-poly
    \textbf{on any fixed run}
  \item
    This includes runs from fixing coins of randomized algorithms
  \end{itemize}
\item
  \textbf{Extractor + Domain Constraint} (§9.3): If \ensuremath{\mathcal{I}} outputs r' \(\in\)
  D(\(\varphi\)) where f(r') = x*, then by Lemma 9.DOM, r'.assignment
  satisfies \(\varphi\); Ext produces the canonical witness W containing
  r'.assignment in poly-time

  \begin{itemize}
  \tightlist
  \item
    Maps successful inversion \(\to\) witness extraction (any valid
    domain element suffices, not necessarily r*)
  \end{itemize}
\item
  \textbf{Contradiction via coin-fixing} (§9.4):

  \begin{itemize}
  \tightlist
  \item
    Suppose a uniform randomized PPT inverter \ensuremath{\mathcal{I}} succeeds with
    probability \(\geq\) 1/poly(n) over (r, coins)
  \item
    By averaging over \ensuremath{\mathcal{I}}\textquotesingle{}s random coins (Yao
    {[}YAO77{]}), \(\exists\) fixed coins c- such that
    Pr\_\{r\(\leftarrow\)D\_n\}{[}\ensuremath{\mathcal{I}}(f(r); c-) inverts{]} \(\geq\)
    1/poly(n)
  \item
    By averaging over r, \(\exists\)r* where \ensuremath{\mathcal{I}}(f(r*); c-) succeeds
    \(\to\) deterministic run on x* = f(r*)
  \item
    Ext(\ensuremath{\mathcal{I}}(x*; c-), x*) produces witness in poly-time
  \item
    \textbf{Contradicts Theorem 8.A} applied to the fixed run with coins
    c-
  \item
    Therefore \ensuremath{\mathcal{I}} cannot succeed \(\to\) f is one-way against classical
    PPT
  \end{itemize}
\end{enumerate}

\textbf{Key architectural point}: We don\textquotesingle{}t need
distributional hardness or worst-case extraction methods. The
\textbf{per-instance deterministic bound} (Theorem 8.A) +
\textbf{coin-fixing} (Yao) + \textbf{extractor} (§9.3) directly yield
Structural OWF security against randomized PPT.

\paragraph{8.3 Solution Multiplicity
Assumption}\label{83-solution-multiplicity-assumption}

\textbf{Remark (Solution Multiplicity).} The Structural OWF security
bound is 2\^{}\(\lambda\) where \(\lambda\) is the
profile\textquotesingle{}s residual parameter. If the planted CNF
\(\varphi\) has K satisfying assignments, the effective security becomes
2\^{}\(\lambda\)/K (since any of the K solutions yields a valid
preimage). For security to hold, we need K = 2\^{}\{o(\(\lambda\))\},
i.e., sub-exponential in \(\lambda\).

\textbf{Concrete analysis for QP-sharp profile (\(\lambda\) = (log
n)²):}

\begin{itemize}
\tightlist
\item
  Security bound: 2\^{}\{(log n)²\}
\item
  If K \(\leq\) n\^{}c solutions (polynomial): K = 2\^{}\{c·log n\}
\item
  Effective security: 2\^{}\{(log n)² - c·log n\} = 2\^{}\{(log n)(log n
  - c)\}
\item
  For n sufficiently large: (log n - c) \textgreater{} 0, so bound
  remains \textbf{super-polynomial}
\end{itemize}

\textbf{Concrete analysis for flat/exponential profile (\(\lambda\) =
n):}

\begin{itemize}
\tightlist
\item
  Security bound: 2\^{}n
\item
  If K \(\leq\) n\^{}c solutions: K = 2\^{}\{c·log n\}
\item
  Effective security: 2\^{}\{n - c·log n\} \(\to\) \textbf{exponential}
  (overwhelming margin)
\end{itemize}

\textbf{Standard CNF families satisfying this bound:}

\begin{itemize}
\tightlist
\item
  Random 3-SAT near threshold: O(1) solutions (trivially satisfies)
\item
  Cryptographic reductions (PRG inversion, hash preimage): unique
  solution by construction
\item
  Planted SAT with unique solution: exactly 1 solution
\item
  k-SAT with bounded solution density: poly(n) solutions
\end{itemize}

\textbf{When this assumption fails:} Only for "almost trivial" CNFs
where K \(\approx\) 2\^{}\(\lambda\) - meaning nearly every assignment
satisfies the formula. Such formulas are computationally uninteresting
for Structural OWF construction (trivial to solve). All standard CNF
families from complexity theory and cryptography have solution counts
far below this threshold.

\textbf{Formal statement:} The security proof assumes \#SAT(\(\varphi\))
\(\leq\) 2\^{}\{o(\(\lambda\))\}, which holds for all non-trivial CNF
families arising from standard reductions. The Lean formalization
includes \texttt{CNFFamily.BoundedSolutions} and
\texttt{CNFFamily.SubexponentialSolutions} definitions encoding this
property (see
\texttt{Layer3\_InformationBounds/Theorems/Quantitative.lean}).

\begin{center}\rule{0.5\linewidth}{0.5pt}\end{center}

\subsection{Part V: The Payoff}\label{part-v-the-payoff}

Notation recap (used throughout Part V and appendices):

\begin{itemize}
\tightlist
\item
  \(\lambda\)(A,x): run-dependent min-cut residual = min\_C
  \(\Sigma\)\_\{v\(\in\)C\}(R\_v \(-\) q\_v)
\item
  \(\lambda\)\_base: profile/instance residual (from §4.3), enforced up
  to o(\(\lambda\)\_base) by FG
\item
  \(\lambda\)(C): residual on a fixed cut C =
  \(\Sigma\)\_\{v\(\in\)C\}(R\_v \(-\) q\_v)
\item
  \(\rho\): effective residual at the bottleneck in a run (\(\rho\)
  \(\approx\) \(\lambda\)(A,x))
\item
  s: pre-final agreement (number of cut bits first revealed before the
  final segment)
\item
  W\_min: profile\textquotesingle{}s minimal window parameter (see §4.3)
  When context is clear, we write \(\lambda\) for \(\lambda\)(A,x).
\end{itemize}

\subsubsection{9. Unconditional One-Way Function
Construction}\label{9-unconditional-one-way-function-construction}

Sections 1-8 established the foundation: the Semantic Conservation Law
(§7) and \textbf{per-instance deterministic bounds} (§8, Theorem 8.A)
showing every FG-wired instance requires super-polynomial
witness-finding time on any fixed computational run.

\textbf{Purpose of Section 9:} We construct an \textbf{unconditional
one-way function} from L*\textquotesingle{}s structural properties,
yielding \textbf{P \(\neq\) NP} for classical uniform PPT via the
classical bridge OWF \(\Rightarrow\) FP \(\neq\) FNP \(\Rightarrow\) P
\(\neq\) NP. This is the \textbf{main result} of the paper - an explicit
Structural OWF construction (not assumption) enabling unconditional
separation.

\paragraph{Structural OWF: Construction
Overview}\label{structural-owf-construction-overview}

The \textbf{Structural OWF} derives hardness from A2 injectivity on
R-bit emergent configurations, with 1-bit parity as discriminator. This
makes the one-wayness information-theoretic rather than computational
(no factoring/discrete-log assumptions).

\textbf{Basic Structural OWF (for P\(\neq\)NP):}

Given a fixed 3-SAT formula \(\varphi\), define f: D(\(\varphi\))
\(\to\) L* where D(\(\varphi\)) contains valid preimages:

\begin{itemize}
\tightlist
\item
  \textbf{Input}: r = (\(\alpha\), gateDigests, salt) where \(\alpha\)
  satisfies \(\varphi\)
\item
  \textbf{Process}: Plant(\(\varphi\), r) --- build DAG overlay, compute
  seed chain with variable seeds from \(\alpha\), apply FrontierGate
  (parity over emergence bits), wire digest into downstream seeds
\item
  \textbf{Output}: x* \(\in\) L* containing overlay structure and
  identity digests (\(\alpha\) is NOT exposed)
\item
  \textbf{Forward}: O(poly(n)) --- seed chain propagation
\item
  \textbf{Backward}: \(\Omega\)(2\^{}n) or \(\Omega\)(n\^{}\{log n\})
  --- must resolve all emergence bits (Theorem 8.A)
\end{itemize}

\textbf{Discriminator vs Hardness Architecture.} The FrontierGate
computes 1-bit parity (XOR) over R emergence bits. This parity is a
\emph{discriminator} that witnesses configuration differences:

\begin{itemize}
\tightlist
\item
  \textbf{Why parity?} Parity has full sensitivity---flipping ANY input
  bit flips the output. This guarantees incomplete observation leaves
  two configs with different parities.
\item
  \textbf{Hardness source}: The 2\^{}R lower bound comes from A2
  injectivity: different R-bit configurations \(\to\) different seeds
  \(\to\) different addresses. Parity witnesses the difference; A2
  enforces the cost.
\item
  \textbf{Key insight}: If parity alone were the hardness source, there
  would be only 2 possibilities (parity 0 vs 1). The 2\^{}R hardness
  comes from needing to identify which of 2\^{}R distinct configurations
  is correct.
\end{itemize}

\textbf{Trapdoor Structural OWF (enables Cryptomania):}

The trapdoor variant starts from the answer and generates the puzzle:

\begin{itemize}
\tightlist
\item
  \textbf{Input}: \(\alpha\) (Alice\textquotesingle{}s secret
  assignment)
\item
  \textbf{Generate \(\varphi\) from \(\alpha\)}: For each variable i,
  add unit clause (x\_i) if \(\alpha\)(i)=1, else (¬x\_i)
\item
  \textbf{Result}: \(\varphi\) where \(\alpha\) is the unique satisfying
  assignment (proven, not assumed)
\item
  \textbf{Public key}: pk = x* = Plant(\(\varphi\), r) --- the
  OAP-encoded instance (\(\varphi\) is seed-locked, not plaintext;
  §10.1.1)
\item
  \textbf{Private key}: sk = \(\alpha\)
\item
  \textbf{Operation}: Anyone can verify witness validity against x*; pk
  defines the OWF challenge
\end{itemize}

With trapdoor (knows \(\alpha\)): O(poly(n)) --- can invert x* to
recover r (\(\alpha\) enables seed computation) Without trapdoor:
\(\Omega\)(2\^{}n) --- must break OAP to find \(\alpha\) (circular:
decode \(\varphi\) \(\to\) need seeds \(\to\) need \(\alpha\) \(\to\)
solve \(\varphi\))

\textbf{OAP Security}: The unit-clause structure is irrelevant to
security because OAP (§10.1.1) seed-locks all formula data. An attacker
sees only \texttt{literal\_i\ =\ enc(actual\_literal\_i)\ \ensuremath{\oplus}\ R\_mask\_i}
where masks reside at seed-dependent addresses. Decoding requires
computing seeds, which requires knowing \(\alpha\)---creating the
circular dependency that forces exponential search.

This places us in Cryptomania --- both private-key (Minicrypt) and
public-key cryptography are possible from L*.

\textbf{Roadmap:}

\begin{itemize}
\tightlist
\item
  \textbf{§9.1}: Parameters and Canonical Decode Order - sizing and
  output structure
\item
  \textbf{§9.2}: Sampler and Total Function f - the planting function f:
  r \(\mapsto\) x* with FG wiring
\item
  \textbf{§9.3}: Extractor and Reduction - poly-time extraction from any
  valid preimage
\item
  \textbf{§9.4}: Security Proof - coin-fixing argument showing f is
  one-way against classical PPT
\end{itemize}

\textbf{The Construction Architecture}: We define a total length-regular
function f: D(\(\varphi\)) \(\to\) L*\_yes by sampling r = (assignment,
gateDigests, structuralSalt) uniformly from the domain D(\(\varphi\))
and outputting x* := f(r), a planted instance with FG wiring and
identity-based digests (non-leaking). Key properties ensure one-wayness:

\begin{itemize}
\tightlist
\item
  \textbf{Every output} x* = f(r) is an FG-wired L* instance with
  per-instance deterministic lower bound (Theorem 8.A); no assignment
  bits are exposed in x*
\item
  \textbf{Extractor}: Given any r' \(\in\) D(\(\varphi\)) where f(r') =
  x*, algorithm Ext produces witness W containing r'.assignment in
  poly-time (§9.3); by domain constraint, r'.assignment satisfies
  \(\varphi\)
\item
  \textbf{Security}: Coin-fixing converts any successful randomized PPT
  inverter to a deterministic run outputting r' \(\in\) D(\(\varphi\));
  by Lemma 9.DOM, r'.assignment satisfies \(\varphi\); Ext produces
  witness in poly-time, contradicting the per-instance bound (§9.4)
\end{itemize}

\textbf{Result}: f is one-way against classical uniform PPT,
establishing P \(\neq\) NP unconditionally (no cryptographic assumptions
required).

\textbf{Universality architecture (constructed, not assumed).}
BuildOverlay embeds Frontier-Gate (FG) wiring during instance
construction, and Plant realizes it in x*. Hence every output x*
inherits FG wiring deterministically. The per-instance lower bound of
Theorem 8.A is therefore an instance-side property (\(\forall\)x*), not
a hypothesis about solvers.

\textbf{Definition (One-Way Function Security). For input size n (where
n := n\_core = core CNF size; §4 Notation), let m(n) denote the preimage
length and D\_n the uniform distribution on \{0,1\}\^{}(m(n)). A
function family \{f\_n: \{0,1\}\^{}(m(n)) \(\to\) \{0,1\}\^{}(poly(n))\}
is one-way} if: Notation: We sometimes write f for f\_n when the input
length n is clear from context.

\begin{enumerate}
\def\labelenumi{\arabic{enumi}.}
\item
  f\_n is computable in polynomial time
\item
  For every uniform PPT inverter \ensuremath{\mathcal{I}}:

  Pr\_\{r\(\leftarrow\)D\_n, coins\}{[} f\_n(\ensuremath{\mathcal{I}}(f\_n(r))) = f\_n(r) {]}
  \(\leq\) negl(n)

  where the probability is over r sampled from D\_n and
  \ensuremath{\mathcal{I}}\textquotesingle{}s internal randomness, and negl(n) denotes a
  negligible function (for all constants c, negl(n) \textless{} 1/n\^{}c
  for sufficiently large n).
\end{enumerate}

\textbf{Our Stronger Property.} Our per-instance deterministic lower
bound (Theorem 8.A) is stronger than standard Structural OWF security:
for \textbf{every} output x* = f(r) and \textbf{any} fixed deterministic
run (including deterministic runs obtained by fixing a randomized
algorithm\textquotesingle{}s coins), witness-finding requires
super-polynomial time. Consequently, no deterministic poly-time run
succeeds on any x*, so any randomized PPT inverter (a distribution over
deterministic runs) has negligible success probability over D\_n. We
present the argument via coin-fixing (§9.4, Appendix O) for alignment
with standard formulations.

\textbf{Boundary (cryptographic).} We do not claim public-key encryption
or key agreement from OWFs; fully black-box constructions are impossible
(Impagliazzo-Rudich \textquotesingle{}89). Non-black-box routes from
OWFs remain open; PKE/KA follow under stronger assumptions (e.g.,
trapdoor permutations, DDH/LWE).

\paragraph{9.1 Parameters, Domain, and Non-Leaking Security
Model}\label{91-parameters-domain-and-non-leaking-security-model}

Fix n\_core and the QP-sharp or flat profile (§4.3). Let
BuildOverlay(n\_core) emit the public metadata (G, \{Sel\_v\}, \{H\_v\},
Enc, F\_overlay, GREQ, \{S(P)\}, \(\Phi\)\textasciitilde{}) as in §6 and §10.1.1. Let
T\_dec be the set of all decode indices (i,p,t) used by \(\Phi\)\textasciitilde{}.
Define a canonical total order \ensuremath{\prec} on T\_dec by lexicographic order on
(i,p,t) composed with the canonical path enumeration used by \(\Phi\)\textasciitilde{}
(seed-independent; §10.1.1). No circularity: decode mask bits are
selected after seeds; decode pools are disjoint from gate-digest pools
(§10.1.1).

\textbf{Non-Leaking Security Model.} The Structural OWF construction
follows a \textbf{domain-constrained, non-leaking} model aligned with
parametrized OWF families (Goldreich 2001, Vol. 2):

\begin{itemize}
\item
  \textbf{Domain constraint}: D(\(\varphi\)) = \{ r \textbar{}
  WellFormedRandomness(\(\varphi\), r) \} where WellFormedRandomness
  requires:

  \begin{enumerate}
  \def\labelenumi{\arabic{enumi}.}
  \tightlist
  \item
    \(\varphi\).satisfies(r.assignment) --- the assignment in r must
    satisfy the CNF
  \item
    r.gateDigests encode the correct parities of emergent configurations
    at FG gates
  \end{enumerate}
\item
  \textbf{Non-leaking property}: The public instance x* contains only
  \textbf{identity-based digests} (XOR of emergent configurations), NOT
  the assignment bits directly. No assignment information is exposed in
  x*.
\item
  \textbf{Inversion success}: An adversary succeeds if it outputs r'
  \(\in\) D(\(\varphi\)) such that plant(\(\varphi\), r') = x*. The
  adversary need not recover the exact planted preimage
  r*---\textbf{any} r' in the domain producing the same output suffices.
\item
  \textbf{Security reduction}: Finding any valid r' \(\in\)
  D(\(\varphi\)) requires finding a satisfying assignment for
  \(\varphi\) (since WellFormedRandomness demands
  \(\varphi\).satisfies(r'.assignment)), which is hard by Theorem 8.A.
\end{itemize}

This model differs from naive Structural OWF constructions that encode
the witness w directly into the output (which would enable trivial
inversion by reading w). Here, the instance reveals only parity
information---insufficient to recover any satisfying assignment without
exponential work.

\textbf{Preimage structure (length-regular).} Let b\_w(n) := n\_core be
the assignment length. Define preimage length m(n) := b\_w(n) + d\_g
where d\_g is the total digest-bit budget for FG gates. For r \(\in\)
\{0,1\}\^{}(m(n)), parse r as:

r = (assignment, gateDigests, structuralSalt)

where assignment \(\in\) \{0,1\}\^{}(b\_w(n)) is a candidate CNF
solution, gateDigests encodes the parity bits for FG gates, and
structuralSalt provides randomness for pool addressing. Seed computation
depends only on gateDigests and metadata---\textbf{never on the
assignment bits directly}.

\textbf{Domain membership (poly-time verifiable).} Given r and
\(\varphi\), checking r \(\in\) D(\(\varphi\)) requires:

\begin{enumerate}
\def\labelenumi{\arabic{enumi}.}
\tightlist
\item
  Verify \(\varphi\).satisfies(r.assignment) --- standard CNF
  evaluation, O(\textbar{}\(\varphi\)\textbar) time
\item
  Verify gateDigests match the emergent parities induced by r.assignment
  --- requires computing the seed chain and checking each FG
  gate\textquotesingle{}s parity, polynomial in
  \textbar{}\(\varphi\)\textbar{}
\end{enumerate}

Both checks are polynomial-time, ensuring the inversion relation is in
FNP (§9.4).

\paragraph{9.2 Sampler \ensuremath{\mathcal{S}} and Total Function
f}\label{92-sampler-ux1d4ae-and-total-function-f}

Sampler \ensuremath{\mathcal{S}}(1\^{}n). Sample r \(\leftarrow\) D(\(\varphi\)) uniformly from
the valid domain and output x* := f(r). In practice, sample assignment w
uniformly, compute the required gateDigests as parities of emergent
configurations under w, and form r = (w, gateDigests, structuralSalt).

\textbf{Definition (f).} On input r = (assignment, gateDigests,
structuralSalt) where r \(\in\) D(\(\varphi\)):

\begin{enumerate}
\def\labelenumi{\arabic{enumi})}
\tightlist
\item
  Compute metadata := BuildOverlay(n\_core).
\item
  Use r.gateDigests directly as the digest values for GREQ=1 nodes. By
  WellFormedRandomness (§9.1), these already encode the correct parities
  of emergent configurations induced by r.assignment.
\item
  Compute seeds in topological order using Enc: GateDigest\_v :=
  r.gateDigests{[}v{]} when GREQ\_v=1; \(\varepsilon\) otherwise.
\item
  Set pool payloads using r.structuralSalt for collision avoidance
  (salted addressing). The decode mask encodes a CNF \(\varphi\) such
  that \(\varphi\)(r.assignment) = 1.
\item
  Emit all published fields as in §6. The instance x* contains only
  identity-based digests---\textbf{no assignment bits are written to
  x*}.
\end{enumerate}

Output x* = (metadata, salts/layout, identity digests) as in §10.

\textbf{Properties.} f is total, polytime, and length-regular (fixed
m(n)); \ensuremath{\mathcal{S}} is efficiently sampleable.

\textbf{Non-leaking verification.} The public instance x* contains:

\begin{itemize}
\tightlist
\item
  Overlay metadata (DAG structure, addressing functions, FG gate
  configuration)
\item
  Parity-based digests (GateDigest\_v = XOR of emergent configurations)
\item
  Structural salts (for pool addressing)
\end{itemize}

Crucially, x* does \textbf{not} contain r.assignment. The assignment is
implicit: any r' \(\in\) D(\(\varphi\)) with plant(\(\varphi\), r') = x*
must have \(\varphi\).satisfies(r'.assignment), but the specific
assignment cannot be read from x*. This is the non-leaking property that
prevents trivial inversion.

\textbf{Lemma 9.LR (Length-Regularity of f).} For each input length n
(core size n\_core), there exists m(n) = b\_w(n) + d\_g + d\_s =
\(\Theta\)(n\_core) + O(poly(n\_core)) such that f\_n: D(\(\varphi\)\_n)
\(\to\) \{0,1\}\^{}(poly(n)) and D\_n is the uniform distribution on the
domain D(\(\varphi\)\_n).

\emph{Proof.} By §9.1, r parses as (assignment, gateDigests,
structuralSalt) with \textbar{}assignment\textbar{} = b\_w(n) =
\(\Theta\)(n\_core), \textbar{}gateDigests\textbar{} = d\_g =
O(poly(n\_core)), and \textbar{}structuralSalt\textbar{} = d\_s = O(1).
By construction (§9.2), f(r) = x* has length polynomial in n\_core.
Hence f is length-regular and D\_n is uniform on the domain.
\(\blacksquare\)

\textbf{FG wiring verification.} BuildOverlay(n\_core) includes the
Frontier-Gate structure from Appendix C.1.1: for each canonical path P
hitting the bottleneck cut C* (containing nodes from C*), the index set
S(P) is published (defined via XOR-cancellation as in App. C.1.1), and
GREQ flags are set for nodes beyond the gate horizon. Step
2\textquotesingle{}s digest plan D assigns parity values to these gates,
and step 5 realizes them via payload adjustments. Therefore every x* =
f(r) has the published S(P) gates wired into seeds via GateDigest\_v at
GREQ=1 nodes, satisfying the FG wiring precondition of Theorem 8.A
(per-instance deterministic bounds).

\textbf{Universality architecture (constructed, not assumed).} The
universality prerequisite for applying the single-run bounds (Theorem
8.A) is enforced by the instance itself: FG wiring is embedded during
BuildOverlay and realized by f's planting steps, so every output x*
inherits it. No algorithmic or distributional assumption is required;
universality is a property of the constructed instances (instance-side),
not a hypothesis about solvers.

\paragraph{9.3 Extractor Ext and
Reduction}\label{93-extractor-ext-and-reduction}

Cross-reference. Canonical output format and checks are specified in
§3.6 (Canonical Forms) and Appendix O.2.1 (Verifier/Extractor
Canonicality Checklist).

\textbf{Extractor Ext(r, x*).} Given any valid domain element r \(\in\)
D(\(\varphi\)) and instance x* = f(r), output canonical witness W = (w,
G\_\(\tau\), Dig\_\(\tau\)) in polynomial time (used in §9.4). The key
property: since r \(\in\) D(\(\varphi\)), we have
\(\varphi\).satisfies(r.assignment) by domain definition, so the
extracted assignment is a valid CNF witness.

\textbf{Input:}

\begin{itemize}
\tightlist
\item
  r = (assignment, gateDigests, structuralSalt) where r \(\in\)
  D(\(\varphi\)) --- by domain constraint, assignment satisfies
  \(\varphi\)
\item
  x* containing overlay metadata (G, Sel\_v, H\_v, Enc, F\_overlay,
  PathOf, S(P), salts)
\end{itemize}

\textbf{Algorithm:}

\begin{enumerate}
\def\labelenumi{\arabic{enumi}.}
\item
  \textbf{Parse x* metadata:} Extract DAG structure G, selector matrices
  \{Sel\_v\}, linear maps \{H\_v\}, address function F\_overlay,
  published path index sets \{S(P)\}, decode schema \(\Phi\)\textasciitilde{}, and all
  designated salts (O(\textbar{}x*\textbar{}) = poly(n) time)
\item
  \textbf{Reconstruct seed chain:} Topologically traverse G from roots
  to sinks:

  \begin{itemize}
  \item
    At roots: Initialize Seed\_root as specified in x*
  \item
    At each node v with parents P(v): a) Retrieve parent outputs y\_u =
    H\_u · x\_u (where x\_u = Sel\_u · e\_u, with e\_u computed from
    salts in x*) b) If GREQ\_v = 1: Compute GateDigest\_v by evaluating
    published pre-horizon indices: GateDigest\_v =
    \(\oplus\)\emph{\{(u,(j,\(\ell\)))\(\in\)S(P(v))\}
    e}\{u,j,\(\ell\)\} where each address is a(u,j,\(\ell\)) =
    F\_overlay(Seed\_u; j,\(\ell\)) and e\_\{u,j,\(\ell\)\} uses
    \(\sigma\)\_\{a(u,j,\(\ell\))\}. c) Form Seed\_v = Enc(v
    \textbar{}\textbar{} sort\{(u, Seed\_u, y\_u) : u \(\in\) P(v)\}
    \textbar{}\textbar{} GateDigest\_v)
  \item
    Total: \(\Theta\)(\(\sum\)\_\{v: GREQ\_v=1\}
    \textbar{}S(P(v))\textbar{}) cell-probe operations (linear in the
    total published index sizes); polynomial in \textbar{}x*\textbar{}
  \end{itemize}
\item
  \textbf{Extract assignment component:} Set W.assignment :=
  r.assignment (the satisfying assignment from the domain element)
\item
  \textbf{Compute gate proof list G\_\(\tau\):} For each published path
  P in x*\textquotesingle{}s gate horizon:

  \begin{itemize}
  \tightlist
  \item
    Record (P, S(P)) where S(P) is the published index set from x*
  \item
    Total: \textbar{}G\_\(\tau\)\textbar{} = O(\(\tau\) · n\_tot /
    R\_avg) = poly(n) entries (§6.2.10)
  \end{itemize}
\item
  \textbf{Compute digest vector Dig\_\(\tau\):} For each path P \(\in\)
  G\_\(\tau\):

  \begin{itemize}
  \tightlist
  \item
    Using final seed chain from step 2, compute addresses
    u\_\{v,j,\(\ell\)\} = F\_overlay(Seed\_v; j,\(\ell\)) for each
    (j,\(\ell\)) \(\in\) S(P)
  \item
    Read salt values \(\sigma\)\_\{u\_\{v,j,\(\ell\)\}\} from x*,
    compute primitives e\_\{v,j,\(\ell\)\}
  \item
    Compute Dig\_\(\tau\){[}P{]} =
    \(\oplus\)\emph{\{(j,\(\ell\))\(\in\)S(P)\} e}\{v,j,\(\ell\)\}
  \item
    Total: O(\(\Sigma\)\_P \textbar{}S(P)\textbar{}) = O(n) XOR
    operations
  \end{itemize}
\item
  \textbf{Return} W = (r.assignment, G\_\(\tau\), Dig\_\(\tau\))
\end{enumerate}

\textbf{Runtime:} Step 1: O(\textbar{}x*\textbar{}); Step 2:
\(\Theta\)(\(\sum\)\emph{\{v: GREQ\_v=1\} \textbar{}S(P(v))\textbar{})
cell-probe operations; Steps 3-6:
O(\(\sum\)}\{P\(\in\)G\_\(\tau\)\}\textbar S(P)\textbar). Total is
linear in published index sizes and thus polynomial in
\textbar{}x*\textbar{}.

In particular, \textbar{}G\_\(\tau\)\textbar{} = O(\(\tau\) · n\_tot /
R\_avg) = poly(n). Using §4.3 (\(\lambda\)\_base = W\_min·r and the rank
budget L·W\_min·r = \(\Theta\)(n\_tot) with R\_avg = \(\Theta\)(r)), we
have n\_tot/R\_avg = O(\(\lambda\)\_base·poly(n)). Hence

\textbar{}G\_\(\tau\)\textbar{} = O(\(\tau\) · \(\lambda\)\_base ·
poly(n)) = poly(n)

and verification runs in

O(\textbar{}G\_\(\tau\)\textbar{} · n/W\_min) = O(\(\tau\) ·
\(\lambda\)\_base · poly(n))

time.

Note (overlay work, not a bypass). Step 2 necessarily evaluates
GateDigest\_v at GREQ\_v=1 nodes by reading designated salts at
seed-bound addresses for the current chain. This is the same overlay
engagement priced in our lower bounds; Ext performs it only after an
inverter supplies a preimage r and is still polynomial in
\textbar{}x*\textbar{}.

\textbf{Lemma 9.Ext (Extractor Correctness).} For any r' \(\in\)
D(\(\varphi\)) with f(r') = x*, the extractor Ext(r', x*) produces a
valid canonical witness W = (w, G\_\(\tau\), Dig\_\(\tau\)) accepted by
Algorithm V in polynomial time. Crucially, r' need not equal the
originally planted r*---any domain element producing x* suffices.

\emph{Proof.}

\begin{enumerate}
\def\labelenumi{\arabic{enumi}.}
\tightlist
\item
  Parse r' = (assignment, gateDigests, structuralSalt) where r' \(\in\)
  D(\(\varphi\)) and f(r') = x*
\item
  By domain constraint (§9.1), \(\varphi\).satisfies(r'.assignment) ---
  this is the key: \textbf{any valid preimage contains a satisfying
  assignment}
\item
  Recompute seed chain from x*\textquotesingle{}s public overlay
  (polynomial-time; deterministic from public data and r'.gateDigests)
\item
  Compute G\_\(\tau\) and Dig\_\(\tau\) via Steps 4-5 above (O(n) XOR
  operations over seed-dependent addresses)
\item
  By WellFormedRandomness, r'.gateDigests encode the correct parities
  matching the seed chain; seeds/digests are deterministically
  reconstructable
\item
  Algorithm V (§10.2) verifies (x*, W) by recomputing seeds/digests and
  checking \(\varphi\)(r'.assignment)=1 \(\to\) accepts in polynomial
  time
\end{enumerate}

\(\blacksquare\)

\textbf{Remark (Observation completeness for planted instances).} The
extractor\textquotesingle{}s correctness relies on unpacking what
"deterministically reconstructable" (proof step 5 above) means for
planted instances. During planting (§9.2), the digest plan D assigns
parity values (step 2), then step 5 realizes them by adjusting payloads
to ensure GateDigest\_v = \(\oplus\)\emph{\{(j,\(\ell\))\(\in\)S(P(v))\}
\(\sigma\)}\{F\_overlay(Seed\_v;j,\(\ell\))\} for each GREQ=1 node v.
This creates a structural invariant: \textbf{the published FG digests
match the actual emergent parities} that arise when computing seeds from
x*\textquotesingle{}s structure.

\textbf{Semantic\(\to\)Operational Bridge.} The connection from
\textbf{semantic correctness} (producing the right witness W) to
\textbf{operational coverage} (the encoder must have visited all 2\^{}R
emergent configurations during execution) requires bridging two
conceptual levels:

\begin{itemize}
\tightlist
\item
  \emph{Semantic level}: Planted instances have exactly one correct
  parity per FG-gated node (determined by the planted assignment). For
  the TM to output the correct parity, it must distinguish which of the
  2\^{}R possible configurations is the correct one.
\item
  \emph{Operational level}: This distinguishing requirement translates
  to the encoder\textquotesingle{}s execution trace---the TM must have
  computationally explored all 2\^{}R values to identify the unique
  correct one.
\end{itemize}

\textbf{Formalization note (Axiomatized Bridge).} The Lean formalization
axiomatizes this connection via
\texttt{parity\_distinguishability\_required\_for\_planted\_correctness}
(TMAxioms.lean): for planted instances with well-formed randomness, if a
TM produces a correct witness, then its encoder must have visited all
2\^{}R emergent configuration values during execution. This axiom
represents the gap between information-theoretic requirements (proven in
Layers 0-3) and operational execution (Layer 4). The
formalization\textquotesingle{}s trust boundary consists of this axiom
plus the Church-Turing thesis.

\textbf{Lemma 9.DOM (Domain-Constrained Inversion).} For any r' \(\in\)
D(\(\varphi\)) with f(r') = x*, the assignment r'.assignment satisfies
\(\varphi\). That is, successful inversion implies SAT-solving.

\emph{Proof.} By definition of the domain D(\(\varphi\)) (§9.1), any r'
\(\in\) D(\(\varphi\)) must satisfy WellFormedRandomness(\(\varphi\),
r'), which requires \(\varphi\).satisfies(r'.assignment). Therefore, any
valid domain element that produces x* contains a satisfying assignment
for \(\varphi\). \(\blacksquare\)

\textbf{Remark (Non-uniqueness of preimages).} Unlike Structural OWF
constructions that encode the witness directly into the output (enabling
unique preimage recovery), our non-leaking model admits \textbf{multiple
valid preimages} for the same x*. If \(\varphi\) has K satisfying
assignments w\(_1\), ..., w\_K, then each can be extended to a valid
domain element r\_i \(\in\) D(\(\varphi\)) with f(r\_i) = x* (provided
the gateDigests are set correctly for each w\_i). The security bound
accounts for this: effective security is 2\^{}\(\lambda\)/K, which
remains super-polynomial when K = 2\^{}\{o(\(\lambda\))\} (see §8.3
Solution Multiplicity).

\textbf{Why non-uniqueness doesn\textquotesingle{}t weaken security.}
The adversary\textquotesingle{}s goal is to find \textbf{any} r' \(\in\)
D(\(\varphi\)) with f(r') = x*. By Lemma 9.DOM, success requires finding
a satisfying assignment for \(\varphi\)---regardless of whether it
equals the planted assignment. This is hard by Theorem 8.A: finding any
satisfying assignment requires super-polynomial time on any fixed run.

\paragraph{9.4 Security (Deterministic FG-based +
Coin-Fixing)}\label{94-security-deterministic-fg-based--coin-fixing}

Scope note (no distributional assumptions). This section relies only on
per-instance deterministic bounds (A1-A5 + FG) and coin-fixing; A6
(Independence) and distributional/average-case machinery are not used.
Standard references on coin-fixing/Yao\textquotesingle{}s min--max
principle: Yao {[}YAO77{]}, Goldreich {[}GOL01{]}, Arora--Barak
{[}AB09{]}.

Every instance x* sampled by \ensuremath{\mathcal{S}} has FG wiring (Appendix C.1.1) by
construction: BuildOverlay includes the FG structure (published S(P)
gates and GREQ flags), and §9.2 steps 2-5 instantiate it. By
\textbf{Theorem 8.A}, for ANY L* instance with FG wiring, the single-run
time satisfies time\_A(x*) \(\geq\) 2\^{}(\(\rho\)-s) ·
\(\Omega\)(n/W\_min) \textbf{on any fixed computational run} (any fixed
coin string). The structural bound from Lemma C.2.1 gives s \(\leq\)
\(\Theta\)(\(\tau\)·\(\lambda\)\_base) \textbf{deterministically
(independent of salt values and coins), yielding \(\rho\) \(-\) s
\(\geq\) (1\(-\)\(\Theta\)(\(\tau\)))·\(\lambda\)\_base and thus time
\(\geq\) 2\^{}(\(\Omega\)(\(\lambda\)\_base))/poly(n\_core) for every
sampled instance, any fixed run}.

\textbf{Theorem 9.DSO (Deterministic strong one-wayness).} For every
output x* in Range(f\_n), no deterministic polynomial-time algorithm A
outputs a valid preimage r \(\in\) D(\(\varphi\)) with f\_n(r) = x*.

\emph{Proof.} Suppose a deterministic poly-time A outputs r \(\in\)
D(\(\varphi\)) with f\_n(r) = x*. By Lemma 9.DOM, r.assignment satisfies
\(\varphi\). By Lemma 9.Ext, Ext(r, x*) produces a canonical witness W
containing r.assignment in polynomial time. This contradicts Theorem
8.A, which lower-bounds witness-finding on any fixed run for every
FG-wired instance x*. Therefore no deterministic poly-time inverter
exists for any x*. \(\blacksquare\)

\textbf{Theorem 9.CF (Coin-fixing reduction; Yao).} Suppose a uniform
PPT inverter \ensuremath{\mathcal{I}} satisfies

Pr\_\{r\(\leftarrow\)D\_n, coins\}{[} \ensuremath{\mathcal{I}}(f\_n(r); coins) \(\in\)
D(\(\varphi\)) \ensuremath{\land} f\_n(\ensuremath{\mathcal{I}}(...)) = f\_n(r) {]} \(\geq\) 1/poly(n)

Then there exists a fixed coin string c- such that

Pr\_\{r\(\leftarrow\)D\_n\}{[} \ensuremath{\mathcal{I}}\_\{c-\}(f\_n(r)) \(\in\) D(\(\varphi\))
\ensuremath{\land} f\_n(\ensuremath{\mathcal{I}}\_\{c-\}(...)) = f\_n(r) {]} \(\geq\) 1/poly(n)

In particular, there exists r* with \ensuremath{\mathcal{I}}\_\{c-\}(f\_n(r*)) = r' where r'
\(\in\) D(\(\varphi\)) and f\_n(r') = f\_n(r*) = x*. By Lemma 9.DOM,
r'.assignment satisfies \(\varphi\). Composing \ensuremath{\mathcal{I}}\_\{c-\} with Ext yields
a deterministic poly-time witness for x*, contradicting Theorem 8.A
(per-instance deterministic lower bound). Therefore f\_n is one-way
against classical PPT.

\textbf{Key insight}: The security argument does \textbf{not} require
preimage uniqueness. By the domain constraint, \textbf{any} valid
preimage r' \(\in\) D(\(\varphi\)) contains a satisfying assignment
(Lemma 9.DOM). The adversary cannot output r' \(\in\) D(\(\varphi\))
without having solved SAT for \(\varphi\), which is hard by Theorem 8.A.
Coin-fixing preserves this: per-instance bounds apply to any fixed
computational trace. Profile-specific bounds give success \(\leq\)
n\^{}\{-\(\Theta\)(log n)\} (QP-sharp) or \(\leq\)
2\^{}\{-\(\Theta\)(n)\} (flat).

\textbf{Remark (why inversion requires SAT-solving).} The
adversary\textquotesingle{}s task is to output r' \(\in\) D(\(\varphi\))
with f(r') = x*. The domain constraint D(\(\varphi\)) = \{ r \textbar{}
WellFormedRandomness(\(\varphi\), r) \} requires
\(\varphi\).satisfies(r.assignment)---there is no way to produce a valid
domain element without knowing a satisfying assignment. Since x*
contains only identity digests (not assignment bits), the adversary
cannot read the assignment from x*; it must solve SAT. This is precisely
what Theorem 8.A rules out in polynomial time.

\textbf{Conclusion.} We have constructed an unconditional one-way
function f from L*\textquotesingle{}s structural properties via
per-instance deterministic bounds (Theorem 8.A), extended to randomized
PPT via coin-fixing. Therefore \textbf{f is one-way against classical
PPT}. Combined with L*\textquotesingle{}s NP-completeness (§10.2-10.3)
and the classical bridge OWF \(\Rightarrow\) FP \(\neq\) FNP
\(\Rightarrow\) P \(\neq\) NP (§10.4), this establishes \textbf{P
\(\neq\) NP for classical uniform PPT} (§10.5).

\begin{center}\rule{0.5\linewidth}{0.5pt}\end{center}

\subsubsection{10. NP-Completeness and Classical
Bridge}\label{10-np-completeness-and-classical-bridge}

\textbf{Formalization note.} This section contains both paper-only and
Lean-verified content. \textbf{§10.1-10.3 (NP-completeness via 3-SAT
reduction)}: Paper exposition only---Lean bypasses this path.
\textbf{§10.4-10.5 (Classical bridge OWF \(\to\) FP\(\neq\)FNP \(\to\)
P\(\neq\)NP, main theorem)}: Proven in BOTH paper and Lean. The Lean
formalization establishes P\(\neq\)NP directly via the OWF route (see
\texttt{Layer5\_Applications/PvsNP/PrimaryPath/ParametricBitstringBridge.lean},
theorem \texttt{fpnefnp\_implies\_not\_peqnp}), which is mathematically
simpler and fully machine-verified with 0 sorries. \textbf{For
authoritative verification, consult the Lean formalization.}

Scope note. All complexity-class claims in this section refer to
classical uniform PPT unless otherwise specified. Model equivalence.
Randomized PPT adversaries reduce to deterministic runs via coin-fixing;
separation is stated for deterministic k-tape TMs.

Section 9 constructed an unconditional one-way function f from
L*\textquotesingle{}s structural properties (per-instance deterministic
bounds, Theorem 8.A), proving f is one-way against classical PPT via
coin-fixing. But the \textbf{classical bridge} OWF \(\Rightarrow\) FP
\(\neq\) FNP \(\Rightarrow\) P \(\neq\) NP requires that the OWF is
constructed from an \textbf{NP-complete} problem. Otherwise, we have
shown a hard-to-invert function exists - not that P \(\neq\) NP.

\textbf{Purpose of Section 10:} We complete the separation by proving
\textbf{L* is NP-complete} and establishing the classical bridge. This
requires:

\begin{enumerate}
\def\labelenumi{\arabic{enumi}.}
\tightlist
\item
  \textbf{L* \(\in\) NP} (Theorem 10.1, §10.2): Polynomial-time
  verification with witness
\item
  \textbf{L* is NP-hard} (Theorem 10.2, §10.3): Witness-preserving
  reduction from 3-SAT
\item
  \textbf{Classical Bridge} (§10.4): OWF \(\Rightarrow\) FP \(\neq\) FNP
  \(\Rightarrow\) P \(\neq\) NP (standard result, applies to our
  constructed OWF)
\end{enumerate}

\textbf{Main Theorem} (§10.5): \textbf{L* is NP-complete} (§10.2-10.3) +
\textbf{f is one-way} (§9) + \textbf{classical bridge} (§10.4) \(\to\)
\textbf{P \(\neq\) NP unconditionally}.

\textbf{Roadmap:}

\begin{itemize}
\tightlist
\item
  \textbf{§10.1}: \textbf{Formal language membership} and reduction from
  3-SAT (seed-locked decode schema \(\Phi\)\textasciitilde{})
\item
  \textbf{§10.2}: \textbf{L* \(\in\) NP} - explicit polynomial-time
  verifier Algorithm V
\item
  \textbf{§10.3}: \textbf{L* is NP-hard} - deterministic
  witness-preserving Karp reduction (Algorithm R)
\item
  \textbf{§10.4}: \textbf{Classical Bridge} - OWF \(\Rightarrow\) FP
  \(\neq\) FNP \(\Rightarrow\) P \(\neq\) NP (Proposition 10.4 connects
  §9\textquotesingle{}s OWF to P \(\neq\) NP)
\item
  \textbf{§10.5}: \textbf{Main Theorem} - P \(\neq\) NP (Theorem 10.5
  combining all results)
\end{itemize}

\textbf{Why this section completes the proof:} Sections 1-9 constructed
a one-way function f from structural properties, but one-wayness alone
doesn\textquotesingle{}t imply P \(\neq\) NP - we must show f is built
from an NP-complete problem. Section 10 closes this gap by proving L* is
NP-complete (membership + hardness) and invoking the classical bridge:
OWF from NP-complete problem \(\Rightarrow\) FP \(\neq\) FNP
\(\Rightarrow\) P \(\neq\) NP. This transforms the Structural OWF
construction into an unconditional complexity separation. The result is
model-specific (deterministic k-tape TMs with constant parameters) but
unconditional within that model - no unproven assumptions, no
cryptographic conjectures, just structural properties of an explicit
language.

\paragraph{10.1 Formal Language Membership (Reduction from
3-SAT)}\label{101-formal-language-membership-reduction-from-3-sat}

Cross-reference. The verifier operates on canonical representations
only; see §3.6 (Canonical Forms) for definitions and Appendix O.2.1 for
the verifier/extractor checklist.

Reader\textquotesingle{}s summary. Two key properties underpin
NP-membership and the no-bypass results:

\begin{itemize}
\tightlist
\item
  Decoding requires overlay engagement (OAP): literal identities are
  recovered by designated reads at seed-dependent addresses; there is no
  standalone CNF table to read first.
\item
  Verification is deterministic and polynomial: the verifier recomputes
  seeds/addresses and unmasks E deterministically from public metadata;
  all required bits are present on the read-only input.
\end{itemize}

V accepts only canonical W; equivalently V\(\circ\)Norm accepts all
valid W.

\textbf{Context}: L* was constructed in §6 with overlay structure
(seeds, DAG, emergence) satisfying properties A1-A5. This section
provides the \textbf{formal language membership predicate} for
NP-completeness and details how instances are generated \textbf{via
reduction from 3-SAT}.

\textbf{Reduction Overview}: Given 3-CNF \(\varphi\) over n\_core
variables, Algorithm R (below) produces instance x* with:

\begin{itemize}
\tightlist
\item
  Seed-locked decode schema \(\Phi\)\textasciitilde{} encoding \(\varphi\) (§10.1.1)
\item
  Overlay DAG metadata from §6 (seeds, gates, addressing)
\end{itemize}

\textbf{Formal Language Membership:}

Notation (decoding). Let \(\varphi\) := Decode\_from\_seeds(\(\Phi\)\textasciitilde{};
Seed\_chain) denote the 3-CNF decoded from schema \(\Phi\)\textasciitilde{} using the
recomputed seed chain. Addresses depend only on seeds/metadata;
structuralSalt is used only for pool addressing (never in seed
computation).

x* \(\in\) L* \(\leftrightarrow\) \(\exists\)w such that:

\begin{itemize}
\tightlist
\item
  (1) Decoded CNF satisfied: \(\varphi\)(w) = 1 (via \(\Phi\)\textasciitilde{} and seed
  chain)
\item
  (2) Overlay constraints hold: all DAG gates from §6 satisfied on (x*,
  w)
\end{itemize}

\textbf{Instance Structure} (from reduction; details in Algorithm R):

\begin{itemize}
\tightlist
\item
  x* = (G, Sel\_v, H\_v, Enc schema, F\_overlay, GREQ, PathOf, S(P),
  salts, public parameters, \(\Phi\)\textasciitilde{})
\item
  DAG from §6: depth O(log n\_core), \(\Sigma\)\_v R\_v =
  \(\Theta\)(n\_core log n\_core)
\item
  All objects poly(n\_core) size
\item
  Encoding: canonical binary format (Appendix D.5), consistent with §3.6
  Canonical Forms; parsing time O(\textbar{}x*\textbar{})
\end{itemize}

\subparagraph{10.1.1 Seed-Locked Decode (OAP) - the overlay is the
problem
(default)}\label{1011-seed-locked-decode-oap---the-overlay-is-the-problem-default}

Remark (digest function Dg). Dg is used only for
canonicalization/compactness in cut-level digests; no cryptographic
collision resistance is assumed or required. Soundness relies on the
verifier\textquotesingle{}s deterministic recomputation of all digest
inputs on the current seed chain.

Context. This construction realizes the structural security principle
introduced in §1.1 : L* achieves "structurally what blockchain
Proof-of-Work achieves cryptographically," compelling work through
problem structure rather than cryptographic assumptions. The seed-locked
decode schema \(\Phi\)\textasciitilde{} is the technical mechanism that binds CNF
recovery to overlay computation, ensuring that any solver must perform
designated work to access the problem definition itself.

Motivation. To prevent a solver from bypassing the overlay by consulting
a standalone CNF table, we adopt a seed-locked decode schema \(\Phi\)\textasciitilde{}
as part of the instance. Under OAP, any implementation must engage the
overlay to even evaluate \(\varphi\)(w).

\textbf{Core Insight (Information Barrier).} All efficient algorithms -
SAT solvers, heuristic search, constraint propagation, learning
algorithms - rely on extracting structural information from the problem
instance to guide search, prune candidates, or make informed decisions.
L* systematically removes ALL sources of structural information through
seed-locked encoding. Standard 3-SAT instances provide clues that enable
algorithmic shortcuts: unit clauses enable unit propagation, pure
literals enable free assignments, clause structure enables subsumption
and resolution, variable frequency guides branching heuristics (VSIDS,
VMTF), conflict analysis enables clause learning (CDCL), and symmetries
enable symmetry breaking. L* blocks every technique: the CNF formula
\(\varphi\) is encoded as E{[}i,p{]} = enc(lit{[}i,p{]}) \(\oplus\)
R{[}i,p{]} where mask bits R{[}i,p{]} reside at seed-dependent addresses
that require correct assignment \(\alpha\) to compute. This creates a
circular dependency: accessing \(\varphi\) requires seeds, computing
seeds requires \(\alpha\), finding \(\alpha\) requires solving
\(\varphi\). \textbf{The result is a perfect information barrier -
algorithms need information to make progress; L* provides no
information.} Accessing ANY structural property of \(\varphi\) (unit
clauses, literal polarities, clause overlaps, variable frequencies)
requires computing the correct seed chain, which requires knowing the
solution. Every query is answered only by "the solution would tell you
that." This transforms computational hardness from "we
haven\textquotesingle{}t found a fast algorithm" into "fast algorithms
are information-theoretically impossible" - the exponential barrier is
not algorithmic but fundamental.

Definition (OAP schema). Replace \(\Phi\)\_enc with a decode schema
\(\Phi\)\textasciitilde{} that specifies how to derive the clause/literal identities
from the current seed chain along canonical paths via Enc and
F\_overlay. Intuitively, \(\Phi\)\textasciitilde{} describes "how to decode \(\varphi\)
from seeds," not \(\varphi\)\textquotesingle s clauses verbatim.

Explicit construction (seed-locked decode; \(\Phi\)\textasciitilde{}). We instantiate
the abstract definition with a concrete, parseable schema that is both
verification-friendly and blocks structural bypass prior to gate-horizon
computation, while keeping the reduction deterministic (Karp).

\begin{itemize}
\item
  Canonical literal encoding. Let b := \(\lceil\)log\(_2\)
  n\_core\(\rceil\) + 1. Encode each literal as a b-bit string lit
  \(\in\) \{0,1\}\^{}b comprising a variable id in {[}n\_core{]} and a
  polarity bit.
\item
  Published artifacts in \(\Phi\)\textasciitilde{}.

  \begin{enumerate}
  \def\labelenumi{\arabic{enumi})}
  \tightlist
  \item
    A masked literal table E \(\in\) (\{0,1\}\^{}b)\^{}(m×3), where for
    each clause i\(\in\){[}m{]} and position p\(\in\)\{1,2,3\}:
    E{[}i,p{]} := enc(lit{[}i,p{]}) \(\oplus\) R{[}i,p{]}. Here enc is
    the fixed canonical encoding and R{[}i,p{]} \(\in\) \{0,1\}\^{}b is
    a per-entry mask defined below.
  \item
    A seed-locked mask recipe \ensuremath{\mathcal{R}} that, for every (i,p,t) with
    t\(\in\){[}b{]}, specifies a designated address selector of the form
    (P(i), sel(i,p,t)). Canonical path assignment: enumerate the
    canonical path family \ensuremath{\mathcal{P}} in a fixed public order (e.g., lexicographic
    by node ids) and define P(i) := \ensuremath{\mathcal{P}}{[}(i\(-\)1) mod
    \textbar{}\ensuremath{\mathcal{P}}\textbar{}{]}. This assignment is
    \(\varphi\)-independent. The selector template sel(i,p,t) is a
    parseable, \(\varphi\)-independent rule that, given the current seed
    chain, yields a designated address u = F\_overlay(Seed\_v;
    j,\(\ell\)) with GREQ\_v=1 along P(i) according to a public
    round-robin over layer positions and local indices. For simplicity
    and tight pricing, each mask bit R{[}i,p{]}{[}t{]} is bound to
    exactly one such designated primitive value e\_u \(\in\) \{0,1\}
    read at u (generalization to small XORs is allowed but not
    required). The mapping (i,p,t) \(\mapsto\) (v,(j,\(\ell\))) used by
    \ensuremath{\mathcal{R}} is injective (no reuse of designated indices) and draws from a
    dedicated decode sub-pool U\_v\^{}(dec) \(\subset\) U\_v that is
    disjoint from the gate-digest pool for S(P), ensuring
    non-interference and independent assignment of mask bits.
  \end{enumerate}
\item
  Mask assignment (instance side; deterministic). The reduction computes
  E deterministically from \(\varphi\) and fixes designated payload bits
  in the dedicated decode sub-pools so that, for every (i,p,t), the
  designated primitive at u\_t := sel(i,p,t)(Seed\_chain) yields a fixed
  bit R{[}i,p{]}{[}t{]} and E{[}i,p{]} = enc(lit{[}i,p{]}) \(\oplus\)
  R{[}i,p{]}. Concretely, for each selector class
  (i,p,t)\(\to\)(v,(j,\(\ell\))), the instance uses a canonical
  "last-bit as adjuster" rule: it sets all but one payload bit in the
  primitive\textquotesingle{}s definition by a fixed parseable recipe
  from published salts, and assigns the remaining adjuster bit so that
  the evaluated primitive equals the required R{[}i,p{]}{[}t{]}.
  Injectivity of (i,p,t)\(\mapsto\)(v,(j,\(\ell\))) and a dedicated
  decode sub-pool U\_v\^{}(dec) ensure no cross-constraints. No
  randomness is used; the mapping is fully deterministic and independent
  of w.

  \textbf{Lemma 10.1.1-Adj (Adjuster uniqueness and polytime
  assignment).} For every mask bit position (i,p,t), the designated
  decode primitive e\_\{u\_t\} is constructed to depend only on
  decode-pool payload bits (not on Z\_v(w,x)). For decode sub-pools we
  set a\_\{v,j,\(\ell\)\}:=0 and fix b\_\{v,j,\(\ell\)\} to the unit
  vector e\(_1\); hence e\_\{u\_t\} = \ensuremath{\langle}b\_\{v,j,\(\ell\)\},
  \(\sigma\)\_\{u\_t\}\ensuremath{\rangle} = \(\sigma\)\_\{u\_t\}{[}1{]}. Fix all payload
  bits of \(\sigma\)\_\{u\_t\} except the adjuster entry
  \(\sigma\)\_\{u\_t\}{[}1{]} by the canonical recipe, and set

  \(\sigma\)\_\{u\_t\}{[}1{]} := R{[}i,p{]}{[}t{]}.

  Thus the linear equation \ensuremath{\langle}b,\(\sigma\)\ensuremath{\rangle} = R{[}i,p{]}{[}t{]} over
  \ensuremath{\mathbb{F}}\(_2\) has a unique solution for the adjuster entry, computable in
  O(1) time. Because (i,p,t)\(\mapsto\)(v,(j,\(\ell\))) is injective and
  uses a dedicated decode sub-pool U\_v\^{}(dec), these assignments are
  independent across all (i,p,t). Therefore assigning all adjuster bits
  takes time \(\Theta\)(m·b) and yields a unique, consistent mask R.
  \(\blacksquare\)
\item
  Decoding algorithm (verifier and solvers). Given the recomputed seed
  chain, for each (i,p):

  \begin{enumerate}
  \def\labelenumi{\arabic{enumi})}
  \tightlist
  \item
    For t=1..b, compute u\_t := sel(i,p,t)(Seed\_chain) (seed-bound
    address) and evaluate e\_\{u\_t\} \(\in\) \{0,1\} via the designated
    primitive.
  \item
    Assemble R{[}i,p{]} := (e\_\{u\_1\}, ..., e\_\{u\_b\}) and output
    enc(lit{[}i,p{]}) := E{[}i,p{]} \(\oplus\) R{[}i,p{]}.
  \item
    Parse enc(lit{[}i,p{]}) to obtain the literal identity.
  \end{enumerate}
\item
  Size/parsability. E has size O(m·b) = O(n\_core·log n\_core); \ensuremath{\mathcal{R}} is
  parseable, \(\varphi\)-independent, and of size O(m·b) references into
  the canonical path family with constant-time selectors. \(\Phi\)\textasciitilde{}
  remains polynomial in \textbar{}\(\varphi\)\textbar.
\end{itemize}

\subparagraph{\texorpdfstring{Algorithm R (Explicit Karp Reduction f:
\(\varphi\) \(\to\)
x*)}{Algorithm R (Explicit Karp Reduction f: \textbackslash varphi \textbackslash to x*)}}\label{algorithm-r-explicit-karp-reduction-f-ux3c6--x}

Input: 3-CNF \(\varphi\) over n\_core variables with m clauses.

Output: Instance x* = O = (G, Sel\_v, H\_v, Enc schema, F\_overlay,
GREQ, PathOf, S(P), salts, public parameters, \(\Phi\)\textasciitilde{}=(E\_lit,\ensuremath{\mathcal{R}})) such
that \(\varphi\)(w)=1 iff (x*,w) is accepted by the base verifier.

Procedure:

\begin{enumerate}
\def\labelenumi{\arabic{enumi})}
\item
  Construct overlay skeleton (graph and parameters).

  \begin{itemize}
  \tightlist
  \item
    Build the canonical DAG G of depth O(log n\_core) with node
    parameters \{R\_v\} so that, by default, \(\Sigma\)\_v R\_v =
    \(\Theta\)(n\_core·log n\_core).
  \item
    Capacity calibration (decode pools). Ensure sufficient decode
    capacity by enforcing \(\Sigma\)\_v D\_v = \(\Sigma\)\_v
    (K\_v·\(\kappa\)) \(\geq\) 3m·b. If m exceeds a constant multiple of
    n\_core, increase per-node K\_v or \(\kappa\), or add a bounded
    number of auxiliary decode-only nodes (GREQ=0, not placed on the
    designated bottleneck cut C*) so that \(\Sigma\)\_v D\_v \(\geq\)
    3m·b holds while keeping all published objects polynomial in
    \textbar{}\(\varphi\)\textbar. This preserves verifier costs and all
    structural properties (disjointness, Keyedness, FG, min-cut
    calibration).
  \item
    Emit canonical Enc schema (length-delimited), F\_overlay (\(\pi\)\_v
    family), and public parameters (\(\kappa\), K\_v, m\(_0\), etc.).
  \item
    Time: poly(n\_core) (each node description and parameter computed in
    O(1)-O(log n\_core)).
  \end{itemize}

  Note (Z\_v indices). The per-node core-bit indices used in Z\_v(w,x)
  are chosen by the bounded-coverage policy of §6.2.4.a so that, for any
  cut C, each witness bit appears in at most c = O(1) Z\_v on C (see
  A.9.1; Lemma C.1.1').
\item
  Selectors and linear maps.

  \begin{itemize}
  \tightlist
  \item
    For each v, choose Sel\_v with full row rank R\_v and unit
    row-weight; choose H\_v with rank(H\_v)=R\_v (Completeness).
  \item
    Time: \(\Sigma\)\_v poly(R\_v) = poly(n\_core·log n\_core).
  \end{itemize}
\item
  Canonical paths and gate index sets.

  \begin{itemize}
  \tightlist
  \item
    Compute the canonical path family \ensuremath{\mathcal{P}} and fix a public total order on
    \ensuremath{\mathcal{P}} (e.g., lexicographic by node ids). For each P\(\in\)\ensuremath{\mathcal{P}}, compute the
    index set S(P) via Appendix A.9 (XOR-cancellation) and C.1.1 rules.
  \item
    Time: O(\(\Sigma\)\_v R\_v) = O(n\_core·log n\_core) (constant
    m\(_0\) patterns; remove O(1) per pattern class).
  \end{itemize}
\item
  Seed-locked decode schema \(\Phi\)\textasciitilde{}.

  \begin{itemize}
  \tightlist
  \item
    Literal encoding: set b := \(\lceil\)log\(_2\) n\_core\(\rceil\)+1.
    Allocate E\_lit \(\in\) (\{0,1\}\^{}b)\^{}(m×3).
  \item
    Mask recipe \ensuremath{\mathcal{R}} (\(\varphi\)-independent template): For each (i,p,t),
    set P(i):=\ensuremath{\mathcal{P}}{[}(i\(-\)1) mod \textbar{}\ensuremath{\mathcal{P}}\textbar{}{]} using the fixed
    public order of \ensuremath{\mathcal{P}}. Define sel(i,p,t) by a public,
    \(\varphi\)-independent round-robin over layer positions along P(i)
    and local indices (j,\(\ell\)), with GREQ=1 along P(i). Ensure
    (i,p,t)\(\mapsto\)(v,(j,\(\ell\))) is injective into decode
    sub-pools U\_v\^{}(dec).
  \item
    Time: O(m·b) to enumerate and encode selectors.
  \end{itemize}
\item
  Decode sub-pool payload assignment (OAP adjusters).

  \begin{itemize}
  \tightlist
  \item
    For each (i,p,t): fix all but one payload bit of
    \(\sigma\)\_\{u\_t\} by canonical recipe; choose an index h with
    b{[}h{]}=1; set \(\sigma\)\_\{u\_t\}{[}h{]} to satisfy e\_\{u\_t\} =
    R{[}i,p{]}{[}t{]} using Lemma 10.1.1-Adj; set E\_lit{[}i,p{]} :=
    enc(lit{[}i,p{]}) \(\oplus\) R{[}i,p{]}.
  \item
    Time: O(1) per bit; total O(m·b).
  \end{itemize}
\item
  Emit disjoint address pools and salts.

  \begin{itemize}
  \tightlist
  \item
    Pool partition (decode/gate separation): For each node v, let d\_v
    denote the number of decode mask indices (i,p,t) assigned to v by
    the round-robin selector template sel from step 4. Partition the
    address pool U\_v = {[}0, D\_v\(-\)1{]} as: U\_v\^{}(dec) := {[}0,
    d\_v\(-\)1{]} (decode sub-pool) and {[}d\_v, D\_v\(-\)1{]} (gate
    sub-pool for S(P) primitives). By injectivity of the mapping
    (i,p,t)\(\mapsto\)(v,(j,\(\ell\))) and the round-robin distribution
    over GREQ=1 nodes along canonical paths, \(\Sigma\)\_v d\_v = 3m·b.
    Capacity condition: enforce \(\Sigma\)\_v D\_v \(\geq\) 3m·b by the
    calibration in step 1; with this, d\_v \(\leq\) D\_v for all v
    (verifiable during reduction construction). This keeps decode and
    gate pools disjoint and ensures non-interference.
  \item
    Emit U\_v and pairwise-distinct salts \(\sigma\)\_u for all
    u\(\in\)U\_v (covering both decode and gate sub-pools), using
    canonical generation.
  \item
    Time/size: O(\(\Sigma\)\_v R\_v·s) with s=\(\Theta\)(log n\_core);
    overall \(\Theta\)(n\_core·log² n\_core) bits to write; time linear
    in output size.
  \end{itemize}
\item
  Publish GREQ and PathOf.

  \begin{itemize}
  \tightlist
  \item
    Set GREQ by the published gate-horizon calibration; publish
    PathOf(v) for canonical paths.
  \item
    Time: O(\textbar{}V\textbar{}) = poly(n\_core).
  \end{itemize}
\end{enumerate}

Complexity analysis:

\begin{itemize}
\tightlist
\item
  Each step runs in time polynomial in n\_core and m. The dominating
  costs are constructing S(P) and writing salts/metadata, both linear in
  the output size \textbar{}x*\textbar{}.
\item
  Size bound: \textbar{}x*\textbar{} = \(\Theta\)(n\_core·log² n\_core)
  bits (salts) + \(\Theta\)(n\_core·log n\_core) (metadata) + O(m·log
  n\_core) (E\_lit,\ensuremath{\mathcal{R}}), hence \textbar{}x*\textbar{} =
  poly(\textbar{}\(\varphi\)\textbar).
\item
  Therefore Algorithm R runs in polynomial time and outputs a
  polynomial-size instance.
\end{itemize}

Correctness:

\begin{itemize}
\tightlist
\item
  By construction, for any w, the verifier decodes \(\varphi\) from
  \(\Phi\)\textasciitilde{} by evaluating mask bits via \ensuremath{\mathcal{R}} on the current seed chain and
  checks overlay constraints; thus \(\varphi\)(w)=1 iff (x*,w) is
  accepted. Witness preservation holds: the same w that satisfies
  decoded \(\varphi\) certifies x* \(\in\) L*. \(\blacksquare\)
\end{itemize}

Non-inferability without post-horizon reads (deterministic Karp). Denote
by Pub the set of all published overlay fields (G, H\_v, Sel\_v, salts
U\_v, Enc, F\_overlay, PathOf, GREQ, S(P), parameters) together with the
schema \(\Phi\)\textasciitilde{} = (E\_lit, \ensuremath{\mathcal{R}}). Let Seeds\_+ be the seed chain values at
nodes with GREQ=1 along canonical paths; these are not published and
require designated computation.

Lemma 10.1.1-NI (Non-inferability for OAP; deterministic Karp). Fix Pub.
There exist two instances x*\_0, x*\_1 with identical Pub but different
decode masks R (hence different \(\varphi\)) such that any
polynomial-time algorithm A that does not evaluate any Seeds\_+-dependent
designated primitive produces identical outputs on x*\_0 and x*\_1. In
particular, no such A can correctly compute any nontrivial structural
functional g(\(\varphi\)) that depends on literal identities across all
instances with this Pub.

Proof. E\_lit = enc(lit) \(\oplus\) R entry-wise. Holding Pub fixed
allows changing only designated payload bits. By injectivity of
(i,p,t)\(\mapsto\)(v,(j,\(\ell\))) into the dedicated decode sub-pools
U\_v\^{}(dec) and Hermeticity, for any choice of mask bits R we can
realize those values by setting the independent adjuster entries (Lemma
10.1.1-Adj) without altering Pub. Hence there exist completions R\textsuperscript{0}, R¹
with R\textsuperscript{0} \(\neq\) R¹ and corresponding \(\varphi\)\textsuperscript{0} \(\neq\) \(\varphi\)¹
that share the same Pub. For any A that does not evaluate
Seeds\_+-dependent designated primitives, Appendix C.1.2 (transcript
indistinguishability without reads) implies A\textquotesingle{}s
transcript on x*\_0 and x*\_1 is identical; therefore
A\textquotesingle{}s output is identical on both. Since g(\(\varphi\)\textsuperscript{0})
\(\neq\) g(\(\varphi\)¹), A cannot compute g correctly on both,
establishing non-inferability. \(\blacksquare\)

Operational corollary. Any algorithm that gains nontrivial information
about literal identities (e.g., variable frequencies, clause overlaps,
backbones) must evaluate Seeds\_+-dependent designated primitives to
determine the corresponding mask bits. By seed-bound addressing and FG
wiring (§6.2.8; Appendix C.1.1), accumulating these determinations
incurs the designated post-horizon work priced by Theorem 10.4.1-BYP.

Mechanism (structural access control). The construction E\_lit =
enc(lit) \(\oplus\) R uses one-time-pad structure but achieves
information hiding via provable access costs rather than secrecy
assumptions (see §1.1):

\begin{itemize}
\tightlist
\item
  R stored at designated addresses u = F\_overlay(Seed\_v; sel) where
  GREQ\_v = 1 (gate-protected)
\item
  Computing Seeds\_+ requires evaluating gate digests GateDigest\_v
  (§6.2.8)
\item
  Each gate digest requires \(\Omega\)(\textbar S(P)\textbar) =
  \(\Omega\)(n/W\_min) designated primitive evaluations (Appendix C.1.1)
\item
  On k-tape TMs: \(\Omega\)(n/W\_min) operations per segment (Lemma
  5.5.1.c), even with pre-cached salts (Lemma C.1.3)
\item
  Unpredictability: skipping any designated term admits two
  transcript-consistent completions with opposite parity (Lemma C.1.2)
\item
  Verifier: computes Seeds\_+ and accesses R in polynomial time; no
  secrets withheld (NP consistency)
\item
  No assumptions: access cost follows from TM mechanics, independent of
  P vs NP or cryptographic conjectures
\end{itemize}

Notes.

\begin{itemize}
\tightlist
\item
  Representation-invariance. The hiding claim is agnostic to how
  literals are encoded (any fixed injective enc works) and to small XOR
  mask recipes; using one primitive per mask bit simplifies accounting
  and maximizes unpredictability by C.1.2.
\item
  Template neutrality. \ensuremath{\mathcal{R}}\textquotesingle{}s structure (path references
  and selectors) is \(\varphi\)-independent; thus the schema leaks no
  counts, overlaps, or frequencies beyond trivial sizes (m, n\_core).
\item
  Compatibility. The construction preserves NP-membership: the verifier
  recomputes Seeds, evaluates the required designated primitives,
  recovers \(\varphi\), and checks \(\varphi\)(w)=1 in polynomial time
  (§10.2).
\item
  Determinism (Karp). No randomness is used in the reduction; mask bits
  R are fixed via designated payload assignments. Lemma 10.1.1-NI uses
  Hermeticity and indistinguishability without designated reads, not
  distributional arguments.
\item
  Visibility. The parameters m and n\_core remain public; these trivial
  size parameters are not considered SAT-useful "structure" within the
  scope of Lemma 10.1.1-NI.
\item
  NP consistency. All mask bits R{[}i,p{]}{[}t{]} reside as designated
  payload on the read-only input; "information-hiding" means
  non-inferability from Pub alone without evaluating those designated
  positions. The verifier deterministically reads them in polynomial
  time; no secrets or cryptographic hardness are assumed.
\end{itemize}

Gated-seed requirement (closes SAT-first bypass). The decode schema
\(\Phi\)\textasciitilde{} explicitly uses seeds for nodes with GREQ\_v = 1 along
canonical paths that cross the designated bottleneck cut C*. Since
Seed\_v := Enc(v \textbar{}\textbar{} sort\{(u,Seed\_u,y\_u)\}
\textbar{}\textbar{} GateDigest\_v) includes GateDigest\_v whenever
GREQ\_v=1, decoding \(\varphi\) under \(\Phi\)\textasciitilde{} necessarily requires
computing those GateDigest\_v values on the current seed chain.

No cycles. By the acyclicity property of §6.2.8 and Appendix C.1.1, each
GateDigest\_v depends only on pre-horizon indices (GREQ=0); therefore,
recomputing these digests to obtain the gated seeds used by \(\Phi\)\textasciitilde{} is
well-founded and follows the DAG order.

Acceptance (OAP). Given (x*, w):

\begin{itemize}
\tightlist
\item
  Recompute the relevant seed chain(s) from ancestor tuples (Seed\_u,
  y\_u) via Enc
\item
  Decode \(\varphi\) := Decode\_from\_seeds(\(\Phi\)\textasciitilde{}; Seed\_chain) and
  check \(\varphi\)(w)=1
\item
  Perform the same overlay consistency checks (selectors, H\_v, seeds,
  and, when GREQ\_v=1, gate digests)
\end{itemize}

Reduction (OAP). The Karp reduction f maps \(\varphi\) to x* by
embedding \(\varphi\) inside \(\Phi\)\textasciitilde{} (with the gated-seed requirement
above) and overlay objects (Sel\_v, H\_v, PathOf, etc.). Size remains
poly(\textbar{}\(\varphi\)\textbar). The reduction depends only on
\(\varphi\) (independent of any solver) and is witness-preserving:
\(\varphi\)(w)=1 iff (x*,w) is accepted. NP-membership and NP-hardness
are preserved; the verifier recomputes any needed GateDigest\_v in
polynomial time (Algorithm V).

\textbf{Scope clarification (what we prove).} The reduction from 3-SAT
to L* adds computational structure - the overlay with seed-locked decode
- that makes L*-search harder than bare 3-SAT search. This is valid
because: (1) L* is a standalone NP-complete language (not "3-SAT with
extra steps"); (2) we prove L* \(\in\) NP via polynomial-time
verification (Algorithm V); (3) we prove L* is NP-hard via
witness-preserving reduction from 3-SAT; (4) we prove search/output for
L* requires super-polynomial time (§9). We are NOT claiming to prove
that bare 3-SAT search (with \(\varphi\) given directly) is
super-polynomial. We ARE proving that L*-search (where \(\varphi\) is
hidden and must be decoded) is super-polynomial, and since L* is
NP-complete, this establishes FP \(\neq\) FNP and hence P \(\neq\) NP.
The hardness comes from L*\textquotesingle{}s combined structure (base
predicate + overlay), not from the base 3-SAT alone. This is the
standard approach for NP-completeness: construct a problem with specific
structure, prove it is NP-complete, then prove lower bounds for that
problem.

Note (ordering only). OAP eliminates the "read CNF first, ignore
overlay" path by requiring overlay engagement for decoding. Together
with disjoint address pools, Keyedness, Completeness, Hermeticity, and
FG, Theorem 10.4.1-BYP rules out any post-hoc shortcut from w to
G\_\(\tau\), establishing that producing the canonical witness
necessarily incurs the overlay work unconditionally.

\textbf{§10.1 Summary:} Established formal reduction from 3-SAT to L*:
Algorithm R maps \(\varphi\) \(\to\) x* deterministically
(witness-preserving, poly-time, poly-size); OAP schema \(\Phi\)\textasciitilde{}
seed-locks decode (requires overlay engagement, closes CNF-bypass);
Non-inferability (Lemma 10.1.1-NI) + gated-seed requirement ensure
\(\varphi\) inaccessible without post-horizon work; reduction preserves
NP-membership (§10.2) and establishes NP-hardness (§10.3).

\begin{center}\rule{0.5\linewidth}{0.5pt}\end{center}

\paragraph{10.2 NP-Membership}\label{102-np-membership}

\textbf{NP-Membership Sanity Check.} L* is a proper NP language: given
(x*, w), a uniform deterministic k-tape TM verifier computes seeds and
designated names via Enc/F\_overlay, reads O(\(\Sigma\)\_v R\_v)
on-input payload bits, and checks all constraints in polynomial time -
no advice or oracles. "Emergence" means those bits are not derivable
before first use, not that they are hidden: every required bit resides
in x*\textquotesingle{}s designated records; "addresses" are parseable
names on a read-only input tape (see Appendix D.5).

\textbf{Theorem 10.1 (NP membership).} L* \(\in\) NP (see Algorithm V
for the explicit verifier). Here n := n\_core and b :=
\(\lceil\)log\(_2\) n\_core\(\rceil\)+1.

\emph{Proof.} The verifier:

\begin{enumerate}
\def\labelenumi{\arabic{enumi}.}
\tightlist
\item
  Recomputes seeds along canonical paths; decodes \(\varphi\) :=
  Decode\_from\_seeds(\(\Phi\)\textasciitilde{}; Seed\_chain) and checks
  \(\varphi\)(w)=1
\item
  For each node v \(\in\) V: computes x\_v=Sel\_v e\_v, y\_v=H\_v x\_v,
  Seed\_v
\item
  For GREQ\_v=1 nodes: recomputes GateDigest\_v by XOR over published
  pre-horizon indices (u,(j,\(\ell\)))\(\in\)S(P(v)) using addresses
  a(u,j,\(\ell\))=F\_overlay(Seed\_u; j,\(\ell\)) (the digest never
  depends on itself; see §6.2.8/Appendix C.1.1)
\end{enumerate}

Path caching yields total work O(\(\Sigma\)\_v R\_v +
\textbar{}\ensuremath{\mathcal{P}}\textbar{}·\textbar{}S(P)\textbar{}) = O(n log n).
\(\blacksquare\)

\emph{Key points:} Salts are explicit constants (no randomness). FG via
GateDigest preserves NP membership. All mask bits referenced by
\(\Phi\)\textasciitilde{} are present as on-input designated payload; decoding
\(\varphi\) is deterministic from x* by Algorithm V-OAP.
"Information-hiding" refers only to the Pub fields prior to any
post-horizon designated reads and does not withhold required information
from the verifier.

Algorithm V-OAP (explicit unmask-and-check for \(\Phi\)\textasciitilde{}). Acceptance
predicate matches §4.4\textquotesingle{}s specification.

Input: instance x* with \(\Phi\)\textasciitilde{}=(E\_lit,\ensuremath{\mathcal{R}}); assignment w

\begin{enumerate}
\def\labelenumi{\arabic{enumi})}
\item
  Parse O; recompute Seeds on canonical paths via Enc; for GREQ\_v=1,
  compute required GateDigest\_v by XOR over published pre-horizon
  indices (u,(j,\(\ell\)))\(\in\)S(P(v)) with
  a(u,j,\(\ell\))=F\_overlay(Seed\_u; j,\(\ell\))
\item
  Decode \(\varphi\) from \(\Phi\)\textasciitilde{} by unmasking E using \ensuremath{\mathcal{R}}:

  \begin{itemize}
  \tightlist
  \item
    For each clause i and position p:

    \begin{itemize}
    \tightlist
    \item
      For t = 1..b: compute u\_t := sel(i,p,t)(Seed\_chain) and evaluate
      e\_\{u\_t\} \(\in\) \{0,1\}
    \item
      Set R{[}i,p{]} := (e\_\{u\_1\}, ..., e\_\{u\_b\}); recover
      enc(lit{[}i,p{]}) := E{[}i,p{]} \(\oplus\) R{[}i,p{]}; parse
      literal
    \end{itemize}
  \end{itemize}
\item
  Evaluate \(\varphi\)(w); if false, reject
\item
  Verify overlay constraints (rank(H\_v)=R\_v;
  seed-consistency/Keyedness; digests for GREQ\_v=1)
\item
  Accept
\end{enumerate}

Complexity: Steps 1-4 run in polynomial time; decoding uses O(m·b)
designated evaluations; all objects are poly-size and parseable (§D.5).

\textbf{Lemma 10.2.1 (Verification Complexity Bound).} The verification
algorithm for L* runs in time O(n²) on inputs of size n.

\emph{Proof.} We analyze each verification step:

\begin{enumerate}
\def\labelenumi{\arabic{enumi}.}
\item
  \textbf{CNF check}: Recompute seeds and decode \(\varphi\) via
  \(\Phi\)\textasciitilde{}; evaluating \(\varphi\)(w) takes
  O(\textbar{}Encode(\(\varphi\))\textbar) = O(n) time.
\item
  \textbf{Per-node computations} (for each v \(\in\) V where
  \textbar{}V\textbar{} = O(n/log n)):

  \begin{itemize}
  \tightlist
  \item
    Computing x\_v = Sel\_v e\_v: O(R\_v) = O(log n\_core) time per node
  \item
    Computing y\_v = H\_v x\_v: O(R\_v²) = O(log² n) for dense H\_v
  \item
    Computing Seed\_v via Enc: O(\textbar{}parents\textbar{}·log n) =
    O(log² n) since bounded degree
  \end{itemize}
\item
  \textbf{GateDigest verification} (for GREQ\_v=1 nodes):

  \begin{itemize}
  \tightlist
  \item
    Evaluate cut-gate digests by batching per canonical path P \(\in\) \ensuremath{\mathcal{P}}
  \item
    For each P, compute GateDigest\_v =
    \(\oplus\)\emph{\{(u,(j,\(\ell\)))\(\in\)S(P)\} e}\{u,j,\(\ell\)\}
    where a(u,j,\(\ell\))=F\_overlay(Seed\_u; j,\(\ell\)) and
    e\_\{u,j,\(\ell\)\} uses \(\sigma\)\_\{a(u,j,\(\ell\))\}
  \item
    Total cost across all required gates: O(\textbar{}\ensuremath{\mathcal{P}}\textbar{} ·
    \textbar{}S(P)\textbar{}) = O(n) (since \textbar{}\ensuremath{\mathcal{P}}\textbar{} =
    \(\Theta\)(W\_min) and \textbar{}S(P)\textbar{} =
    \(\Theta\)(n/W\_min))
  \end{itemize}
\item
  \textbf{Path validation}:

  \begin{itemize}
  \tightlist
  \item
    Verifying PathOf(·) consistency: O(depth·\textbar{}V\textbar{}) =
    O(log n · n/log n) = O(n)
  \end{itemize}
\end{enumerate}

Total complexity: O(n) + O((n/log n)·log² n) + O(n) + O(n) = O(n log n).

With careful implementation avoiding redundant recomputation (path
caching), the bound is O(n log n). Even with a naive implementation,
verification remains O(n²), establishing polynomial-time verification.
\(\blacksquare\)

\emph{Fan-in reminder:} With max in-degree \(\Delta\)\_in = O(1) (set in
§6.2.7), per-node seed parsing and Enc operations are O(log² n),
yielding the stated O(n log n) overall bound with caching.

Since \textbar{}x*\textbar{} = \(\Theta\)(n\_core · log² n\_core), this
bound is polynomial in \textbar{}x*\textbar{}; hence L* \(\in\) NP.

\paragraph{10.3 NP-Hardness (explicit
reduction)}\label{103-np-hardness-explicit-reduction}

\textbf{Theorem 10.2 (NP-hardness via 3-SAT; deterministic Karp
reduction).} L* is NP-hard.

\subparagraph{Algorithm V (Explicit Verifier for
L*)}\label{algorithm-v-explicit-verifier-for-l}

Input: Instance x* = overlay and witness w.

Goal: Accept iff (i) w satisfies the decoded CNF \(\varphi\) via
\(\Phi\)\textasciitilde{}, and (ii) all overlay checks succeed along the DAG.

Procedure (runs in polynomial time):

\begin{enumerate}
\def\labelenumi{\arabic{enumi})}
\tightlist
\item
  Parse overlay objects (G, Sel\_v, H\_v, Enc schema, F\_overlay, GREQ,
  PathOf, S(P), salts, public parameters, \(\Phi\)\textasciitilde{}).
\item
  Recompute seeds; decode \(\varphi\) := Decode\_from\_seeds(\(\Phi\)\textasciitilde{};
  Seed\_chain); check \(\varphi\)(w)=1.
\item
  Topologically order G; initialize seeds at roots as specified.
\item
  For each node v in topological order: a) Recompute parent
  (Seed\_u,y\_u) and form Seed\_v := Enc(v \textbar{}\textbar{}
  sort(\{(u,Seed\_u,y\_u)\}) \textbar{}\textbar{} GateDigest\_v) using
  the published schema; when GREQ\_v=1, compute GateDigest\_v by
  evaluating the published path gates G\_\{C*,P\} (Appendix C.1.1) along
  their S(P) sets: for each (u,(j,\(\ell\))) \(\in\) S(P(v)), compute
  a(u,j,\(\ell\))=F\_overlay(Seed\_u; j,\(\ell\)), read
  \(\sigma\)\_\{a(u,j,\(\ell\))\}, compute e\_\{u,j,\(\ell\)\}, and XOR
  per the schema. b) Compute x\_v := Sel\_v·e\_v (selector applies to
  the local primitive vector e\_v derived from (x*,w)). c) Compute y\_v
  := H\_v·x\_v and record (Seed\_v,y\_v) for children. d)
  Acyclicity/pre-horizon check: verify that any published S(P(v)) used
  in GateDigest\_v draws only from GREQ=0 (pre-horizon) indices and that
  no digest depends (directly or indirectly) on itself; reject on
  violation (GREQ acyclicity).
\item
  At sinks, check that the final conditions hold
  (Closure/Recoverability, Section 6.2.7; Lemma J.0 in Appendix J;
  Appendix L): from recorded (Seed,y) tuples the sink seeds
  deterministically parse the ancestor forest.
\end{enumerate}

Accept iff all checks above succeed; otherwise reject.

Complexity: Each step performs O(\(\Sigma\)\_v R\_v) primitive
evaluations and
O(\(\Sigma\)\_\{P\(\in\)\ensuremath{\mathcal{P}}\}\textbar S(P)\textbar)=\(\Theta\)(n/W\_min)
parity computations per gated node; all data are of size poly(n\_core),
so total time is polynomial in \textbar{}x*\textbar{}. (Note: The
verifier does not enumerate feasible cut-worlds; uniqueness at
acceptance follows logically from digest binding by Lemma C.2.ACC.)
\(\blacksquare\)

\textbf{Lemma 10.R (Reduction correctness - witness-preserving, size).}
Let f be Algorithm R mapping \(\varphi\)\(\mapsto\)x*. Then:

\begin{itemize}
\tightlist
\item
  (Completeness) If \(\varphi\)(w)=1, there exists a canonical witness
  W=(w,G\_\(\tau\),Dig\_\(\tau\)) such that (x*,W) is accepted by
  Algorithm V. In particular, V recomputes Seeds, decodes \(\varphi\)
  via \(\Phi\)\textasciitilde{}, checks \(\varphi\)(w)=1, and verifies all overlay
  constraints including any required gate digests.
\item
  (Soundness) If (x*,W) is accepted by V, then letting w be the
  assignment component of W, the decoded formula satisfies
  \(\varphi\)(w)=1. There are no spurious accepting witnesses:
  acceptance requires both decode correctness and overlay consistency
  along the DAG.
\item
  (Size) \textbar{}x*\textbar{} \(\leq\)
  poly(\textbar{}\(\varphi\)\textbar). Concretely, each node contributes
  O(R\_v) selectors plus O(1) metadata; salts contribute
  \(\Theta\)(n\_core·log² n\_core) bits; all published path/selector
  data and \(\Phi\)\textasciitilde{} are polynomial in \textbar{}\(\varphi\)\textbar.
\end{itemize}

Hence f is a deterministic many-one, witness-preserving reduction from
3-SAT to L*. \(\blacksquare\)

\emph{Proof of NP-hardness.} f is poly-time deterministic and
(\(\varphi\),w)\(\in\)3-SAT \(\leftrightarrow\) (x*,w)\(\in\)L*, hence
L* is NP-hard. \(\blacksquare\)

\paragraph{\texorpdfstring{10.4 Classical Bridge: From OWF to P \(\neq\)
NP}{10.4 Classical Bridge: From OWF to P \textbackslash neq NP}}\label{104-classical-bridge-from-owf-to-p--np}

\textbf{Theorem 10.3 (NP-completeness).} L* is NP-complete.

\emph{Proof.} NP membership (Theorem 10.1): polynomial verification via
explicit deterministic checks. NP-hardness (Theorem 10.2): polynomial
reduction from 3-SAT. \(\blacksquare\)

With L* proven NP-complete and the unconditional Structural OWF
constructed in §9, we now establish the \textbf{classical bridge}
connecting Structural OWF security to P \(\neq\) NP. \emph{(See Theorem
10.4.1-BYP.)}

Important note (what is hard). The Structural OWF security reduction
does not require showing that the decoded CNF \(\varphi\) is itself hard
to solve in isolation. Any efficient inverter \ensuremath{\mathcal{I}} producing a preimage r
from x* yields, via Ext(r, x*), a polynomial-time algorithm that outputs
the canonical witness W from x* alone, contradicting the per-instance
deterministic bound (Theorem 8.A). The hardness resides in producing W
from x* (overlay engagement and digest binding), not in an a priori
assumption about \(\varphi\)'s stand-alone difficulty.

\textbf{Proposition 10.4 (Classical Bridge: OWF \(\Rightarrow\) FP
\(\neq\) FNP \(\Rightarrow\) P \(\neq\) NP).} Let f be a total,
polynomial-time computable function that is one-way against PPT
inverters (as proven for our f in §9.4). Since FP \(\subseteq\) PPT (any
deterministic polynomial-time algorithm is also a randomized
polynomial-time algorithm), f is also one-way against FP inverters.
Define the inversion relation R\_f(y, x) := {[}f(x) = y{]}. Then:

\begin{enumerate}
\def\labelenumi{\arabic{enumi}.}
\tightlist
\item
  \textbf{R\_f \(\in\) P}: Verification runs f and checks equality in
  polynomial time
\item
  \textbf{R\_f is an FNP problem}: Find x such that f(x) = y
\item
  \textbf{If P = NP, then FP = FNP}: Standard complexity theory result
\item
  \textbf{FP = FNP implies polynomial-time inverter}: There exists
  polytime Inv where f(Inv(y)) = y for all y \(\in\) Range(f)
\item
  \textbf{Contradiction}: This contradicts f being one-way on y
  \(\leftarrow\) f(U\_n)
\item
  \textbf{Therefore FP \(\neq\) FNP, hence P \(\neq\) NP}
\end{enumerate}

\emph{Proof.} Standard result in complexity theory (see e.g., Goldreich
{[}GOL01{]} or Arora--Barak {[}AB09{]}). The key observation: if all NP
problems have polynomial-time solutions (P=NP), then all FNP problems
have polynomial-time solutions (FP=FNP), so inverting any poly-time
function would be in FP, contradicting one-wayness against FP inverters.
\(\blacksquare\)

\textbf{Application to our construction:}

\begin{itemize}
\tightlist
\item
  §9 proved: f is one-way against classical PPT (via per-instance bounds
  + coin-fixing)
\item
  Theorem 10.3: L* is NP-complete
\item
  Proposition 10.4: OWF \(\Rightarrow\) FP \(\neq\) FNP \(\Rightarrow\)
  P \(\neq\) NP
\item
  \textbf{Conclusion}: P \(\neq\) NP unconditionally (no cryptographic
  assumptions)
\end{itemize}

\paragraph{10.4.1 Canonical Witness for the Search/Output
Task}\label{1041-canonical-witness-for-the-searchoutput-task}

Definition (Canonical seed-consistent witness). For input x*, a witness
is W = (w, G\_\(\tau\), Dig\_\(\tau\)), where:

\begin{itemize}
\tightlist
\item
  w is an assignment with \(\varphi\)(w)=1 (\(\varphi\) decoded via
  \(\Phi\)\textasciitilde{} with gated-seed requirement);
\item
  G\_\(\tau\) is the published canonical gate list G\_\(\tau\)\^{}(pub)
  := \{ (P, S(P)) : P \(\in\) \ensuremath{\mathcal{P}} \} consisting of all canonical paths P
  in the published family \ensuremath{\mathcal{P}} together with their published index sets
  S(P) (Appendix C.1.1). A valid witness must enumerate exactly these
  items (no omissions/additions), all bound to the final seed chain;
\item
  Dig\_\(\tau\) provides, for each item (P, S(P)) \(\in\) G\_\(\tau\),
  the corresponding cut-gate digest bit G\_\{C*,P\}(x*,w) computed on
  the final seed chain.
\end{itemize}

Construction (from (x*, w)). In polynomial time:

\begin{enumerate}
\def\labelenumi{\arabic{enumi})}
\tightlist
\item
  Recompute seeds along the canonical path(s) from parent tuples
  (Seed\_u, y\_u) via Enc; for nodes with GREQ=1, compute GateDigest\_v
  using the published pre-horizon indices (Appendix C.1.1): for each
  (u,(j,\(\ell\)))\(\in\)S(P(v)), compute
  a(u,j,\(\ell\))=F\_overlay(Seed\_u; j,\(\ell\)), evaluate
  e\_\{u,j,\(\ell\)\} from \(\sigma\)\_\{a(u,j,\(\ell\))\}, and XOR;
\item
  Select the \(\tau\)-budgeted set of paths P and their corresponding
  published subsets S(P) (Appendix C.1);
\item
  For each (P, S(P)) \(\in\) G\_\(\tau\), evaluate the seed-bound
  addresses u\_\{v,j,\(\ell\)\} along the final seed chain and compute
  the parity G\_\{C*,P\}(x*,w) by XORing e\_\{v,j,\(\ell\)\} over S(P);
  set Dig\_\(\tau\){[}P{]} := G\_\{C*,P\}(x*,w);
\item
  Output W = (w, G\_\(\tau\), Dig\_\(\tau\)).
\end{enumerate}

Verifier (canonical format). Given (x*, W):

\begin{enumerate}
\def\labelenumi{\arabic{enumi})}
\tightlist
\item
  Decode \(\varphi\) via \(\Phi\)\textasciitilde{} (using gated seeds as specified in
  §10.1.1; see Algorithm V-OAP for the unmask step) and check
  \(\varphi\)(w)=1; perform overlay-consistency checks as in the
  acceptance predicate;
\item
  Check that the provided gate list G\_\(\tau\) equals the published
  canonical set G\_\(\tau\)\^{}(pub) (keys and S(P) match exactly);
\item
  For each (P, S(P)) \(\in\) G\_\(\tau\)\^{}(pub), recompute seeds along
  P; for each (v,(j,\(\ell\))) \(\in\) S(P), compute u\_\{v,j,\(\ell\)\}
  = F\_overlay(Seed\_v; j,\(\ell\)), evaluate e\_\{v,j,\(\ell\)\}, and
  XOR to recompute the digest bit; check Dig\_\(\tau\){[}P{]} equals the
  recomputed value (Appendix C.1). Each item costs
  O(\textbar{}S(P)\textbar{}) time; total verification is polynomial in
  \textbar{}x*\textbar{}.
\end{enumerate}

Soundness clause. The verifier rejects any witness W for which
G\_\(\tau\) \(\neq\) G\_\(\tau\)\^{}(pub) or for which any entry
Dig\_\(\tau\){[}P{]} fails to match the recomputed digest on the final
seed chain.

Equivalence. For every x* and w with \(\varphi\)(w)=1, the above
construction yields (G\_\(\tau\), Dig\_\(\tau\)) such that (w,
G\_\(\tau\), Dig\_\(\tau\)) is accepted by the canonical verifier.
Conversely, if (w, G\_\(\tau\), Dig\_\(\tau\)) is accepted, then (x*, w)
is accepted by the base verifier. Since G\_\(\tau\) is computable in
polynomial time from (x*, w), deciding x* \(\in\) L* remains many-one
equivalent to 3-SAT.

Definition (Output task F\_\{can\}). Given x* with x* \(\in\) L*,
produce a canonical witness W = (w, G\_\(\tau\), Dig\_\(\tau\)) as
above. Verification is polynomial in \textbar{}x*\textbar{}. Our
search/output lower-bound arguments apply to producing W; see §5-§8 and
Appendix C.1.

\subparagraph{10.4.1-BYP No-Overlay Bypass Theorem (Canonical
Witness)}\label{1041-byp-no-overlay-bypass-theorem-canonical-witness}

\textbf{Theorem 10.4.1-BYP (No overlay bypass for W).} Fix the
deterministic k-tape TM model and an instance x* \(\in\) L* constructed
with Seed-Locked Decode (\(\Phi\)\textasciitilde{}), disjoint designated address pools,
injective and parseable Enc, Keyedness, Hermeticity, Receiving-Window
Attribution (RWA), Completeness (rank forcing), and Frontier-Gate wiring
with GREQ beyond the gate horizon (§6.2.8). For any algorithm that
outputs a valid canonical witness W = (w, G\_\(\tau\), Dig\_\(\tau\)) on
x*, producing the listed digest bits Dig\_\(\tau\) necessarily requires
performing, on the native run, the designated computations of the
overlay along the listed gate paths; in particular, there is no strategy
that first computes w and then synthesizes Dig\_\(\tau\) via cheaper
post-processing without incurring the same asymptotic overlay work. Our
lower bounds quantify over runs that produce W from x* alone (assembling
W when w is provided as extra input remains polynomial by
§10.4.1\textquotesingle{}s construction).

\emph{Proof outline (five steps; full details in Appendix C.1.4):}

\textbf{Step 1 (Verifier contract).} Algorithm V (§10.1)
deterministically recomputes each digest Dig\_\(\tau\){[}P{]} by
evaluating the XOR over S(P) using the current seed chain and comparing
against the supplied value. Any mismatch \(\to\) rejection. Therefore, a
correct output must match the verifier\textquotesingle{}s deterministic
recomputation for all (P, S(P)) \(\in\) G\_\(\tau\).

\textbf{Step 2 (Unpredictability-without-reads).} By Lemma C.1.2
(Unpredictability), for any gate path P with index set S(P), skipping
the designated read of any single term at index (j,\(\ell\)) \(\in\)
S(P) admits at least two completions of the unread payload bits that (a)
keep the transcript prefix fixed through all prior designated reads, and
(b) flip the parity of the final XOR digest while remaining consistent
with the overlay structure. Since the verifier checks the digest
deterministically and only one parity value yields acceptance, an
algorithm cannot correctly predict the digest without evaluating all
\textbar{}S(P)\textbar{} designated terms on the current seed chain.

\textbf{Step 3 (Seed-bound addressing).} The address function
F\_overlay(Seed\_v; j, \(\ell\)) maps each index (j,\(\ell\)) \(\in\)
S(P) to a designated memory location dependent on the current seed
Seed\_v. When the resolution prefix changes (seed chain changes), Lemma
A.1.\(\Delta\) (Address Churn) shows that
\(\Theta\)(\textbar S(P)\textbar) designated addresses change.
Precomputed digests for old seed chains use wrong addresses and are
rejected by the verifier\textquotesingle{}s recomputation check.
Therefore, digest computation must engage the overlay on the current
seed chain.

\textbf{Step 4 (RWA credits and first-use reads).} By Receiving-Window
Attribution (RWA; §4.2), each designated address contributes to the
resolution count q\_v when first read on the current seed chain. By
Hermeticity (A1; Lemma 4.2.I), fresh information enters only via
first-use designated reads. Evaluating \textbar{}S(P)\textbar{} =
\(\Theta\)(n/W\_min) designated primitives requires reading
\(\Theta\)(n/W\_min) designated addresses (each contributing \(\leq\) B
bits), establishing the information-flow requirement.

\textbf{Step 5 (TM step pricing).} On a deterministic k-tape TM with
alphabet \(\Gamma\) and k tapes, each tape operation accesses O(1)
cells. By Lemma 5.5.1.c (TM digest cost) and Lemma D.2.1 (TM Addressing
Cost), even with all salt values pre-cached on the tapes, computing
addresses F\_overlay(Seed\_v; j, \(\ell\)) for \(\Theta\)(n/W\_min)
indices and performing the designated XOR requires \(\Omega\)(n/W\_min)
sequential tape operations. Lemma C.1.1' (Bounded Reuse) ensures that
across the \(\Theta\)(\(\tau\)·n\_tot/R\_avg) listed paths in
G\_\(\tau\), the total reuse is bounded, preventing combinatorial
collapse. Hence time \(\geq\) \(\Omega\)(n/W\_min) per fresh gate
digest.

\textbf{Completeness.} Combining Steps 1-5: Any algorithm outputting
valid W = (w, G\_\(\tau\), Dig\_\(\tau\)) must (i) compute all listed
digests correctly (Step 1), (ii) evaluate designated terms for each
digest (Step 2), (iii) on the current seed chain (Step 3), (iv) via
designated reads (Step 4), (v) costing \(\Omega\)(n/W\_min) TM steps per
digest (Step 5). Together with Lemma C.1.1' and segment counting
(Appendix C.2), this yields the stated asymptotic overlay work
requirement. \ensuremath{\square}

\textbf{Cross-reference.} The anti-bypass checklist below addresses six
common shortcut strategies and explains why each fails under the
construction\textquotesingle{}s guarantees.

\textbf{Anti-bypass checklist (soundness against common shortcuts).}

\begin{itemize}
\tightlist
\item
  (Precomputed tables) Precomputing digests for earlier keys does not
  help: GateDigest\_v binds to the current Seed\_v; address churn (Lemma
  A.1.\(\Delta\)) changes the address set; old digests are invalid under
  the new key and are rejected by the verifier.
\item
  (Linearization tricks) Parity is already linear; the cost is in
  touching the designated terms for the current Seed\_v. Lemma 5.5.1.c
  lower-bounds the required tape operations even with all salts
  pre-cached.
\item
  (Algebraic summaries) Any summary not recomputed on the current seed
  chain fails the recomputation check in the verifier\textquotesingle{}s
  step 3; summaries recomputed on the current chain incur the same
  \(\Omega\)(\textbar S(P)\textbar) probe work.
\item
  (Cross-seed reuse) Keyedness forbids merging artifacts across
  different Seed\_v; reuse without recomputation produces wrong
  addresses/digests and is rejected.
\item
  (Address virtualization) Published S(P) indices are interpreted via
  F\_overlay(Seed\_v;·); virtualization cannot shortcut evaluating the
  seed-bound addresses without reading/designating the same payloads,
  which is exactly what the verifier checks.
\item
  (Guess-and-verify) Guesses succeed with probability
  2\^{}(-\textbar{}S(P)\textbar{}); the verifier deterministically
  recomputes digests, so wrong guesses cause rejection. For single-run
  lower bounds, segment counting forces many fresh digests, making
  guessing infeasible.
\end{itemize}

\paragraph{\texorpdfstring{10.5 Main Theorem: P \(\neq\)
NP}{10.5 Main Theorem: P \textbackslash neq NP}}\label{105-main-theorem-p--np}

\textbf{Theorem 10.5 (Main Result: P \(\neq\) NP for classical uniform
PPT).} P \(\neq\) NP for classical computation.

\emph{Proof.} Combining the results from §§6-10:

\begin{enumerate}
\def\labelenumi{\arabic{enumi}.}
\item
  \textbf{NP-complete language} (§6, §10.1-10.3):

  \begin{itemize}
  \tightlist
  \item
    L* explicitly constructed via deterministic Karp reduction from
    3-SAT (Theorem 10.2)
  \item
    L* \(\in\) NP with polynomial-time verifier (Theorem 10.1)
  \item
    Therefore L* is NP-complete (Theorem 10.3)
  \end{itemize}
\item
  \textbf{Per-instance deterministic bounds} (§7-§8):

  \begin{itemize}
  \tightlist
  \item
    Structural properties A1-A5 \(\to\) Semantic Conservation Law
    (§7.2.1)
  \item
    Every FG-wired instance has witness-finding lower bound on any fixed
    run (Theorem 8.A)
  \item
    Bounds: \(\geq\) n\^{}(\(\Omega\)(log n\_core)) (QP-sharp) or
    \(\geq\) 2\^{}(\(\Omega\)(n\_core)) (flat)
  \end{itemize}
\item
  \textbf{Unconditional Structural OWF construction} (§9):

  \begin{itemize}
  \tightlist
  \item
    Function f: r \(\to\) x*, planted instance with FG wiring (§9.2)
  \item
    Every output x* = f(r) satisfies Theorem 8.A
  \item
    Coin-fixing: randomized PPT inverter \(\to\) deterministic run
    \(\to\) Ext produces witness \(\to\) contradiction (§9.4)
  \item
    Therefore f is one-way against classical PPT
  \end{itemize}
\item
  \textbf{Classical bridge} (§10.4):

  \begin{itemize}
  \tightlist
  \item
    OWF \(\Rightarrow\) FP \(\neq\) FNP \(\Rightarrow\) P \(\neq\) NP
    (Proposition 10.4)
  \item
    Applies to our constructed OWF from step 3
  \end{itemize}
\end{enumerate}

\textbf{Therefore P \(\neq\) NP.} \(\blacksquare\)

\textbf{Machine-verified proof:} Complete formalization in
\texttt{ParametricBitstringBridge.lean}:

\begin{itemize}
\tightlist
\item
  \textbf{Main theorem}: \texttt{fpnefnp\_implies\_not\_peqnp}
\item
  \textbf{Size}: 870 lines, self-contained OWF \(\to\) P\(\neq\)NP proof
  chain
\item
  \textbf{Witnesses}: Explicit bitstring witnesses (fully constructive)
\item
  \textbf{Trust boundary}: Two operational axioms (all standard)
\item
  \textbf{Supporting files}:

  \begin{itemize}
  \tightlist
  \item
    \texttt{ParametricBitstringBridge.lean} - FP\(\neq\)FNP \(\to\)
    P\(\neq\)NP bridge
  \item
    \texttt{OWFQP.lean}, \texttt{StructuralOWFExponential.lean} -
    Structural OWF security proofs
  \end{itemize}
\end{itemize}

\textbf{Scope \& Properties:}

\begin{itemize}
\tightlist
\item
  \textbf{Model}: Classical uniform PPT (probabilistic Turing machines;
  constant tapes/alphabet). Quantum adversaries out of scope.
\item
  \textbf{Unconditional}: No cryptographic assumptions (OWF constructed,
  not assumed)
\item
  \textbf{Explicit}: L* is constructive; deterministic reduction from
  3-SAT
\item
  \textbf{Structural}: Hardness from instance properties (A1-A5), not
  algorithmic limitations
\item
  \textbf{Paradigm-invariant}: SCL framework applies across
  computational models
\end{itemize}

\begin{center}\rule{0.5\linewidth}{0.5pt}\end{center}

\textbf{Why This Works: The Three-Dimensional Separation}

The separation between polynomial-time verification and
super-polynomial-time witness-finding can be understood through
SCL\textquotesingle{}s two-dimensional constraint (q, \(\Phi\)) that
induces three operational obstacles L* enforces simultaneously. A
witness W provides perfect compression (\(\lambda\) = 0) for
verification; without it, search faces unavoidable exponential barriers
across these routes:

\textbf{Dimension 1: Space (Unmergeable States)}

\emph{Verification with witness W:}

\begin{itemize}
\tightlist
\item
  Witness specifies unique seed chain: follow single deterministic path
  through DAG
\item
  No alternatives to track: \(\lambda\) = 0 \(\to\) Alt\_v = 1 at all
  nodes
\item
  Space: polynomial (only current node state)
\end{itemize}

\emph{Search without witness:}

\begin{itemize}
\tightlist
\item
  Must explore 2\^{}\(\lambda\) possible seed chains across bottleneck
  cut C*
\item
  Injective seed construction (A2; §6.2.7): different histories \(\to\)
  different Seed\_v \(\to\) different address permutations \(\pi\)\_v
\item
  Keyedness (§4.4): merging artifacts with different Seed\_v computes
  incorrect addresses \(\to\) verification failure
\item
  Result: Must maintain Alt\_v \(\geq\) 2\^{}(R\_v-q\_v) distinguishable
  artifacts simultaneously (Theorem 7.A; §7.2.1)
\item
  Space: exponential \(\geq\) \(\Omega\)(\(\lambda\)) where \(\lambda\)
  = min\_C \(\Sigma\)\_\{v\(\in\)C\}(R\_v - q\_v)
\end{itemize}

\emph{Why unavoidable:} Mathematical impossibility - not algorithmic
limitation. Different seeds produce provably different computational
outcomes (Lemma 7.I); merging would create detectable errors.

\textbf{Dimension 2: Time-Resolution (Unresolvable Information)}

\emph{Verification with witness W:}

\begin{itemize}
\tightlist
\item
  Witness provides all necessary information directly
\item
  Decode \(\varphi\) from \(\Phi\)\textasciitilde{} using witness-provided seed chain
  (§10.1.1)
\item
  Evaluate \(\varphi\)(w), check gate digests: deterministic
  polynomial-time operations
\item
  No search required: \(\lambda\) = 0 \(\to\) all bits resolved upfront
\end{itemize}

\emph{Search without witness:}

\begin{itemize}
\tightlist
\item
  Emergence (A3; §6.2.5; Lemma 6.1): Each node v contributes R\_v fresh,
  unpredictable bits

  \begin{itemize}
  \tightlist
  \item
    Like unobserved coin flips - cannot be inferred or deduced without
    explicit reads
  \item
    No pattern, no formula, no algebraic shortcut
  \end{itemize}
\item
  Bandwidth limit (Lemma 5.5.1): Each TM step acquires \(\leq\) B =
  k\(\lceil\)log\(_2\)\textbar{}\(\Gamma\)\textbar{}\(\rceil\) = O(1)
  bits

  \begin{itemize}
  \tightlist
  \item
    Reading R\_v - q\_v bits requires \(\geq\) (R\_v - q\_v)/B
    sequential steps
  \end{itemize}
\item
  Dependency (A5): Must process DAG in topological order - cannot "leap
  ahead" via reasoning
\item
  Result: Single-run lane requires m\_seg \(\geq\) 2\^{}(\(\rho\)-s)
  rollback segments (Segment Counting; App. C.2), each costing
  \(\Omega\)(n/W\_min) steps via FG (App. C.1.1)
\item
  Time: \(\geq\) 2\^{}(\(\rho\)-s) · \(\Omega\)(n/W\_min) where \(\rho\)
  \(\geq\) \(\lambda\)\_base - s and s \(\leq\)
  \(\Theta\)(\(\tau\)·\(\lambda\)\_base)
\end{itemize}

\emph{Why unavoidable:} Information-theoretic barrier. Fresh bits are
independent of prior observations (distributional independence; Appendix
J); no amount of cleverness can predict unobserved random bits. Must
physically traverse and read.

\textbf{Dimension 3: Time-Elimination (Uneliminable Candidates)}

\emph{Verification with witness W:}

\begin{itemize}
\tightlist
\item
  No elimination needed: witness provides the correct answer
\item
  Verify correctness: check \(\varphi\)(w) = 1, validate gate digests
\item
  No search through wrong candidates
\end{itemize}

\emph{Search without witness:}

\begin{itemize}
\tightlist
\item
  Per-node antagonism (§6.1.1.A): Testing wrong candidate x\_v
  eliminates exactly one world (\(\leq\)1 bit)

  \begin{itemize}
  \tightlist
  \item
    Factoring-style constraint: y\_v = H\_v x\_v with rank R\_v
  \item
    Each incorrect x\_v yields unique wrong y\_v with no cascade to
    other candidates
  \item
    Contrast with SAT: violating clause eliminates 2\^{}k assignments (k
    = clause width)
  \end{itemize}
\item
  CDT (Lemma CDT-1\textquotesingle{}; App. C): Semantic progress
  requires explicit work

  \begin{itemize}
  \tightlist
  \item
    NF\_C (consequence set) only grows via: (i) increasing \(\Delta\)q
    through reads, or (ii) computing gate digests
  \item
    No "free" inference: conflict learning, unit propagation, algebraic
    elimination all require designated work
  \end{itemize}
\item
  Result: Restart lane requires \ensuremath{\mathbb{E}}{[}tries{]} \(\geq\)
  2\^{}(\(\Delta\)(C*)) independent attempts (Lemma 7.R; App. C.4.2)
\item
  Time: \(\geq\) 2\^{}(\(\Delta\)(C*))/B in expectation;
  Yao\textquotesingle{}s minimax yields explicit instance (App. C)
\end{itemize}

\emph{Why unavoidable:} Structural constraint. Full-rank requirement at
each node prevents bulk elimination. Must test exponentially many
candidates; no algorithmic trick can prune faster than \(\leq\)1 bit per
test.

\textbf{Quantitative Results:}

\begin{enumerate}
\def\labelenumi{\arabic{enumi}.}
\item
  \textbf{Verification (with witness):} Polynomial time

  \begin{itemize}
  \tightlist
  \item
    Theorem 10.1: L* \(\in\) NP via Algorithm V (§10.2)
  \item
    Verifier runs in O(n\_core\^{}c · poly(log n\_core)) for some
    constant c
  \item
    All three dimensions collapse: \(\lambda\) = 0
  \end{itemize}
\item
  \textbf{Witness-finding (without witness):} Super-polynomial time on
  any fixed run

  \begin{itemize}
  \tightlist
  \item
    Per-instance deterministic bounds (§8, Theorem 8.A):

    \begin{itemize}
    \tightlist
    \item
      QP-sharp profile (\(\lambda\)\_base = \(\Theta\)(log² n\_core)):
      time \(\geq\) n\^{}(\(\Omega\)(log n\_core))
    \item
      Flat profile (\(\lambda\)\_base = \(\Theta\)(n\_core)): time
      \(\geq\) 2\^{}(\(\Omega\)(n\_core))
    \end{itemize}
  \item
    All three dimensions exponential: unmergeable states, unresolvable
    info, uneliminable candidates
  \end{itemize}
\item
  \textbf{Separation:} P \(\neq\) NP (Theorem 10.5)

  \begin{itemize}
  \tightlist
  \item
    Via unconditional Structural OWF construction (§9) + classical
    bridge (§10.4)
  \item
    Scope: Classical uniform PPT (coin-fixing extends to randomized;
    quantum out of scope)
  \end{itemize}
\end{enumerate}

\textbf{Why This Works: Maximal Incompressibility}

L* is not merely "a hard problem" - it is \emph{provably maximally
incompressible} across all three dimensions:

\begin{itemize}
\item
  \textbf{Other NP-complete problems} may have exponential search spaces
  (Dimension 1), but often allow:

  \begin{itemize}
  \tightlist
  \item
    Inference shortcuts (fails Dimension 2 barrier partially)
  \item
    Bulk elimination via propagation (fails Dimension 3 barrier
    partially)
  \item
    Examples: XOR-SAT (Gaussian elimination), 2-SAT (implication
    graphs), Horn-SAT (unit propagation)
  \end{itemize}
\item
  \textbf{L* blocks ALL three dimensions SIMULTANEOUSLY}:

  \begin{itemize}
  \tightlist
  \item
    Space: Injective seeds + Keyedness \(\to\) provably unmergeable
    (§7.2.1)
  \item
    Resolution: Emergence + Bandwidth \(\to\) provably unresolvable
    (§6.2.5 + Lemma 5.5.1)
  \item
    Elimination: Per-node + CDT \(\to\) provably uneliminable (§6.1.1.A
    + Lemma CDT-1\textquotesingle{})
  \end{itemize}
\end{itemize}

The proof does not rely on "no algorithm has been found" - it
demonstrates "L*\textquotesingle{}s structure makes polynomial-time
algorithms \emph{mathematically impossible}" via three complementary
enforcement mechanisms (operational routes). This is structural
impossibility, not algorithmic limitation.

\textbf{Cross-Reference to Technical Mechanisms:}

\begin{itemize}
\tightlist
\item
  Five-Dimensional Antagonism (§6.1.1) maps to three obstacles
  (§6.1.1.1)
\item
  Three-Dimensional Framework (§7.5) formalizes Storage, Resolution,
  Elimination
\item
  Two Lanes (Appendix C) show both strategies fail on different
  dimensions (Appendix C)
\item
  Formal proofs: Theorem 7.A (SCL core), Lane analysis (Appendix C),
  Theorem 10.4.1-BYP (no bypass)
\end{itemize}

\begin{center}\rule{0.5\linewidth}{0.5pt}\end{center}

\textbf{Analysis Scope:}

\begin{itemize}
\tightlist
\item
  All results proven for L* specifically (§Note)
\item
  Deterministic k-tape TMs; classical sequential model
\item
  No advice, no oracles, no hidden channels (Hermeticity; A1)
\item
  Per-instance deterministic hardness (Theorem 8.A)
\end{itemize}

\textbf{Cross-references for detailed proofs:}

\begin{itemize}
\tightlist
\item
  Appendix J: DAG Min-Cut and Closure
\item
  Appendix C.1: Frontier-Gate wiring and per-segment baseline
\item
  Appendix C.2: Segment Counting
\item
  Appendix C.4: Restart lane quick bound (across-tries)
\end{itemize}

\begin{center}\rule{0.5\linewidth}{0.5pt}\end{center}

\subsection{Part VI: Context and
Extensions}\label{part-vi-context-and-extensions}

\subsubsection{11. Related Work}\label{11-related-work}

Sections 1-10 established \textbf{P \(\neq\) NP} via unconditional
Structural OWF construction from L*, an explicitly constructed
NP-complete language with engineered structural properties (A1-A5). This
section positions the result within complexity theory literature, maps
where detailed technical content appears, and clarifies how this
approach relates to prior lower bound techniques.

\textbf{Navigation note:} Detailed technical proofs appear in §5
(paradigm adapters), §7-8 (SCL and bounds), §9-10 (OWF and separation).
This section provides a literature roadmap, not duplicate exposition.

\paragraph{11.1 Paradigm Connections: Where SCL Projects to Standard
Measures}\label{111-paradigm-connections-where-scl-projects-to-standard-measures}

The Semantic Conservation Law (Theorem 7.A: q + \(\Phi\) \(\geq\) R) is
paradigm-independent. Below we map how \(\lambda\) (min-cut residual)
manifests in standard complexity metrics across computational models.

\textbf{Backtracking (tree size):}

\begin{itemize}
\tightlist
\item
  \textbf{Projection}: Tree nodes \(\geq\) 2\^{}\(\lambda\) (branches
  correspond to artifacts at unresolved cut nodes)
\item
  \textbf{Detailed analysis}: §5.2 (algorithm classes), §5.3
  (correctness requirements)
\item
  \textbf{For L*}: \(\lambda\) = \(\Theta\)(log² n) \(\to\) tree size
  \(\geq\) n\^{}(\(\Theta\)(log n)) (QP-sharp); \(\lambda\) =
  \(\Theta\)(n) \(\to\) size \(\geq\) 2\^{}(\(\Theta\)(n)) (flat)
\end{itemize}

\textbf{Dynamic Programming (table size):}

\begin{itemize}
\tightlist
\item
  \textbf{Projection}: Distinct memoization keys \(\geq\)
  2\^{}\(\lambda\) (seed-dependent subproblems cannot merge without
  resolution)
\item
  \textbf{Detailed analysis}: §5.2 (DP correctness for L*), §5.3
  (seed-consistency requirement)
\item
  \textbf{For L*}: Keys must be seed-consistent (including Seed\_v
  suffices); this forbids cross-seed merges \(\to\) exponential blowup
  when \(\lambda\) = \(\omega\)(log n)
\end{itemize}

\textbf{OBDD (diagram width):}

\begin{itemize}
\tightlist
\item
  \textbf{Projection}: Width \(\geq\) 2\^{}\(\lambda\) at some layer
  (simultaneously live nodes correspond to artifacts)
\item
  \textbf{Order-robustness}: Requires expander-FG gates (any variable
  ordering hits width bottleneck)
\item
  \textbf{Detailed analysis}: §5.3 (OBDD bounds), §6.2.10
  (expander-parity gate construction following {[}HOO06{]})
\item
  \textbf{Key reference}: {[}WEG00{]} (OBDD bounds), {[}HOO06{]}
  (expander-parity techniques - we use same gadgets)
\end{itemize}

\textbf{Resolution/CDCL (proof size via width):}

\begin{itemize}
\tightlist
\item
  \textbf{Projection}: Clause width \(\geq\) \(\lambda\) \(\to\) proof
  size \(\geq\) 2\^{}(\(\Omega\)(\(\lambda\))) (Ben-Sasson \& Wigderson
  width\(\to\)size bridge)
\item
  \textbf{Detailed analysis}: §5.2 (Resolution bounds), Appendix G
  (width\(\to\)size bridge for L*)
\item
  \textbf{Key reference}: {[}BEN01{]} (our SCL generalizes their width
  bottleneck to paradigm-independent min-cut framework)
\end{itemize}

\textbf{Turing Machines (time):}

\begin{itemize}
\tightlist
\item
  \textbf{Projection}: Time \(\geq\) 2\^{}(\(\rho\)-s) ·
  \(\Omega\)(n/W\_min) where \(\rho\) is algorithm-achieved residual, s
  is FG pre-revelation bound
\item
  \textbf{Key mechanism}: Segment Counting (Appendix C.2) forces
  \(\geq\) 2\^{}(\(\rho\)-s) rollback segments; FG (Appendix C.1.1)
  forces \(\Omega\)(n/W\_min) work per segment
\item
  \textbf{Detailed proof}: §8 (per-instance bounds), Appendix C (FG
  mechanism), Appendix D (TM formalization)
\item
  \textbf{Per-instance deterministic}: Applies to any fixed run; extends
  to randomized via coin-fixing (§9.4)
\end{itemize}

\textbf{Time-Space Tradeoffs:}

\begin{itemize}
\tightlist
\item
  \textbf{Formula}: T \(\geq\) \(\Sigma\)\_W max(0, \(\lambda\)\_W \(-\)
  2S·w) where \(\lambda\)\_W is residual on window W, S is space, w is
  window width
\item
  \textbf{Detailed proof}: Appendix C.4
\item
  \textbf{Connection}: Generalizes pebbling lower bounds via min-cut
  accounting
\end{itemize}

\textbf{Communication Complexity:}

\begin{itemize}
\tightlist
\item
  \textbf{Structural connection}: Min-cut with \(\lambda\) unresolved
  bits \(\to\) \(\Omega\)(\(\lambda\)) communication under appropriate
  partition
\item
  \textbf{Lifting}: Standard gadget composition lifts query bounds to
  communication bounds
\item
  \textbf{Note}: We use this structure conceptually but do not construct
  lifting gadgets here; connection is via min-cut methodology shared
  with communication complexity literature {[}KUS97{]}
\end{itemize}

\textbf{Summary table (cross-reference):}

\begin{itemize}
\item
  \textbf{Backtracking}

  \begin{itemize}
  \tightlist
  \item
    SCL Bound: Tree \(\geq\) 2\^{}\(\lambda\)
  \item
    Detailed Analysis: §5.2, §5.3
  \end{itemize}
\item
  \textbf{DP}

  \begin{itemize}
  \tightlist
  \item
    SCL Bound: Keys \(\geq\) 2\^{}\(\lambda\)
  \item
    Detailed Analysis: §5.2, §5.3
  \end{itemize}
\item
  \textbf{OBDD}

  \begin{itemize}
  \tightlist
  \item
    SCL Bound: Width \(\geq\) 2\^{}\(\lambda\)
  \item
    Detailed Analysis: §5.3, §6.2.10
  \item
    Key Reference: {[}WEG00{]}, {[}HOO06{]}
  \end{itemize}
\item
  \textbf{Resolution}

  \begin{itemize}
  \tightlist
  \item
    SCL Bound: Size \(\geq\) 2\^{}(\(\Omega\)(\(\lambda\)))
  \item
    Detailed Analysis: §5.2, Appendix G
  \item
    Key Reference: {[}BEN01{]}
  \end{itemize}
\item
  \textbf{TM Time}

  \begin{itemize}
  \tightlist
  \item
    SCL Bound: T \(\geq\) 2\^{}(\(\rho\)-s)·\(\Omega\)(n/W\_min)
  \item
    Detailed Analysis: §8, Appendix C
  \end{itemize}
\item
  \textbf{Time-Space}

  \begin{itemize}
  \tightlist
  \item
    SCL Bound: T \(\geq\) \(\Sigma\) max(0,\(\lambda\)\_W\(-\)2S·w)
  \item
    Detailed Analysis: Appendix C.4
  \item
    Key Reference: {[}LEN82{]} pebbling
  \end{itemize}
\end{itemize}

All bounds apply to L* with explicitly constructed overlay (§6).
Extension to other problems requires analogous A1-A5 structure (see
§12.5 for natural problem discussion).

\paragraph{11.2 Prior Lower Bound Techniques: Lineage and
Positioning}\label{112-prior-lower-bound-techniques-lineage-and-positioning}

\textbf{What we build on (technical lineage):}

\begin{enumerate}
\def\labelenumi{\arabic{enumi}.}
\item
  \textbf{Width/size tradeoffs} ({[}BEN01{]} Ben-Sasson \& Wigderson
  2001): Width bottleneck \(\to\) exponential proof size in resolution.
  \emph{Our extension}: SCL shows width is one manifestation of min-cut
  residual \(\lambda\); same bottleneck appears across all paradigms
  with different metrics.
\item
  \textbf{OBDD lower bounds} ({[}WEG00{]} Wegener 2000, {[}HOO06{]}
  Hoory et al. 2006): Variable-ordering techniques and expander-parity
  for order-robust bounds. \emph{Our use}: Same expander gadgets
  (§6.2.10), but SCL provides paradigm-independent explanation -
  expanders force Cartesian factoring (Lemma J.1-Cart) preventing state
  compression.
\item
  \textbf{Algorithm\(\to\)hardness connections} ({[}WIL14{]} Williams
  2014): Algorithms for circuits enable circuit lower bounds. \emph{Our
  connection}: Structural OWF construction (§9) is constructive
  algorithm\(\to\)hardness via Ext extractor - any efficient inverter
  yields efficient witness-finder.
\item
  \textbf{Coin-fixing for distributional hardness} ({[}YAO77{]} Yao
  1977): Minimax principle for randomized algorithms. \emph{Our use}:
  §9.4 applies coin-fixing to extend per-instance deterministic bounds
  to randomized PPT.
\item
  \textbf{Single-principle approach} ({[}HAS86{]} Håstad 1986):
  Consistent bounds across models from unified technique. \emph{Our
  philosophy}: Like Håstad\textquotesingle{}s switching lemma for
  AC\^{}0, SCL provides single conservation law (q + \(\Phi\) \(\geq\)
  R) explaining exponential bounds across all paradigms.
\end{enumerate}

\textbf{How we differ (novelty summary):}

\begin{itemize}
\item
  \textbf{Paradigm unification}: Prior work proved exponential bounds
  within specific models (resolution {[}BEN01{]}, OBDDs {[}WEG00{]},
  etc.). SCL unifies these under one framework - min-cut residual
  \(\lambda\) explains all bounds as different "currency" measurements
  of same information bottleneck.
\item
  \textbf{Instance-side obligations}: Prior techniques analyzed
  algorithmic limitations (circuit depth, proof length, etc.). We focus
  on what \emph{instances} mathematically require via structural
  properties (A1-A5), making bounds model-independent.
\item
  \textbf{Barrier circumvention}: Prior unification attempts (circuit
  complexity measures, natural properties) hit barriers (Natural Proofs,
  Relativization, Algebrization - see §12.6 for detailed analysis). Our
  instance-specific, non-relativizing, combinatorial-counting approach
  operates outside these barriers\textquotesingle{} scope.
\item
  \textbf{Unconditional OWF}: We construct (not assume) one-way function
  from L* (§9), yielding unconditional P \(\neq\) NP for classical
  uniform PPT via classical bridge (§10). This is constructive proof,
  not assumption-based separation.
\end{itemize}

\textbf{Why prior P \(\neq\) NP attempts failed:}

Three decades of barrier theorems ({[}RAZ97{]} Natural Proofs,
{[}BAK75{]} Relativization, {[}AAR09{]} Algebrization) proved broad
classes of techniques cannot separate P from NP. Prior attempts either:

\begin{itemize}
\tightlist
\item
  Hit these barriers (circuit complexity measures are "natural" and
  violate {[}RAZ97{]})
\item
  Had proof gaps (e.g., flawed independence assumptions in prior claimed
  proofs)
\item
  Required unproven conjectures
\end{itemize}

\textbf{Our circumvention} (summary; full analysis in §12.6):

\begin{itemize}
\tightlist
\item
  \textbf{Natural Proofs}: We prove bounds for exponentially sparse
  instance family (density \(\leq\)
  2\^{}(-\(\Omega\)(\(\lambda\)\_base))), violating "largeness"
  requirement - not a natural property
\item
  \textbf{Relativization}: L*\textquotesingle{}s seed-dependency
  structure (Seed\_v via content-addressed Enc) is non-relativizing -
  oracles would violate Hermeticity (A1), changing problem definition
\item
  \textbf{Algebrization}: Artifact-counting (\(\Phi\) = log\(_2\)(Alt))
  is combinatorial, not algebraic - no low-degree polynomial extension
  preserves Cartesian factoring (Lemma J.1-Cart)
\end{itemize}

\textbf{See §12.6} for detailed barrier analysis, formal barrier
statements (Appendix N), and intellectual honesty discussion.

\textbf{Positioning in complexity theory landscape:}

Our unconditional Structural OWF construction places classical uniform
PPT in at least \textbf{Minicrypt} (Impagliazzo\textquotesingle{}s five
worlds) - one-way functions exist, ruling out Algorithmica (P=NP) and
Heuristica (average-case only). Whether we\textquotesingle{}re in
Minicrypt (OWFs but no PKE) or Cryptomania (full crypto) remains open;
standard black-box impossibility ({[}IR89{]} Impagliazzo-Rudich 1989)
applies to OWF\(\to\)PKE.

\paragraph{11.3 Synthesis}\label{113-synthesis}

\textbf{The unifying contribution:}

This work demonstrates that exponential lower bounds across resolution,
OBDDs, dynamic programming, backtracking, and Turing machine time stem
from a single structural bottleneck - the min-cut residual \(\lambda\)
measuring information debt accumulation. The same conservation law (q +
\(\Phi\) \(\geq\) R) explains bounds across all paradigms, measured in
different metrics (tree nodes, DP keys, OBDD width, resolution proof
size, TM time). This unification enables unconditional separation:
constructing an OWF from L*\textquotesingle{}s engineered structure
(A1-A5) yields P \(\neq\) NP for classical uniform PPT via the classical
bridge (§10.4-10.5).

\textbf{How it complements prior work:}

Prior paradigm-specific techniques proved exponential bounds within
individual models using distinct methodologies (width-expansion for
resolution, variable-ordering for OBDDs, pebbling for time-space, etc.).
SCL reveals these as manifestations of the same underlying principle -
Injectivity (A2) + Emergence (A3) + Dependency (A5) force
2\^{}\(\lambda\) distinguishable artifacts at bottleneck cuts, which
translates to exponential cost in any computational currency. This
perspective complements model-specific analyses by providing a unified
information-theoretic foundation.

\textbf{Scope and future directions:}

\begin{itemize}
\tightlist
\item
  \textbf{Proven}: P \(\neq\) NP for classical uniform PPT via L*
  (explicitly constructed NP-complete language)
\item
  \textbf{Model extensions}: Circuits, RAM, parallel models -
  straightforward adapter work (§12.11)
\item
  \textbf{Quantum}: BQP vs NP remains genuinely open; unclear if SCL
  survives superposition (§12.11 research question)
\item
  \textbf{Natural problems}: Extending to 3-SAT, TSP, etc. requires
  constructing analogous overlays satisfying A1-A5 (§12.5 discusses
  barriers; §12.12 proposes \(\lambda\)-program research agenda)
\end{itemize}

\textbf{Where to read more:}

\begin{itemize}
\tightlist
\item
  \textbf{§12}: Discussion and implications (complexity spectrum, scope
  boundaries, barrier analysis, natural problems, future work)
\item
  \textbf{§5}: Paradigm-specific bounds and correctness requirements
  (§5.2-5.3)
\item
  \textbf{Appendices C, G, J}: FG mechanism, width\(\to\)size bridge,
  Cartesian factoring
\item
  \textbf{Appendix N}: Formal barrier statements
\end{itemize}

\paragraph{11.4 Observation-Based Semantics: Theoretical
Foundations}\label{114-observation-based-semantics-theoretical-foundations}

\textbf{Why measure "bits observed" instead of "steps taken"?}

The TM time bound (§8, Appendix C) uses an observation-based model:
computational progress is measured by \textbf{information acquired}
(bits of hidden input revealed) rather than raw state transitions. This
approach has strong theoretical pedigree---multiple fields independently
discovered that measuring information acquired yields valid lower
bounds.

\textbf{Prior Lower Bound Techniques (1970s-1990s):}

Each technique independently discovered that computational lower bounds
arise from information requirements:

\begin{itemize}
\item
  \textbf{Decision Trees} {[}WEG87{]}: For Boolean function f on n bits,
  depth D(f) \(\geq\) log\(_2\)(\# leaves) since each query halves
  possibilities. Our use: \texttt{parity\_requires\_all\_bits} (FG
  digest = parity \(\to\) must read all bits).
\item
  \textbf{Communication Complexity} {[}YAO79{]}, {[}KUS97{]}: For f: X ×
  Y \(\to\) \{0,1\}, CC(f) \(\geq\) log\(_2\)(partition number) since
  each bit exchanged refines the rectangle partition. Our use:
  Coin-fixing (§9.4); min-cut bounds information flow.
\item
  \textbf{Pebbling Games} {[}LEN82{]}, {[}COO76{]}, {[}PIP77{]}: For DAG
  computation, time T × space S \(\geq\) \(\Omega\)(n²) since pebbles
  track live values. Our use: Time-space formula (Appendix C.4).
\item
  \textbf{Branching Programs/OBDD} {[}BAR91{]}, {[}BRY86{]}: For
  read-once BPs, size \(\geq\) 2\^{}(width) since width bounds
  distinguishable states per level. Our use: A4 Closure enforces
  read-once; width bounds follow.
\item
  \textbf{Resolution} {[}BEN01{]}: Width w implies size \(\geq\)
  2\^{}(\(\Omega\)(w)) via the width\(\to\)size theorem. Our use:
  Appendix G adapts this bridge.
\end{itemize}

\textbf{Key references:}

\begin{itemize}
\tightlist
\item
  {[}WEG87{]} Wegener, I. (1987). "The Complexity of Boolean Functions."
  Wiley-Teubner.
\item
  {[}YAO79{]} Yao, A.C. (1979). "Some complexity questions related to
  distributive computing." STOC.
\item
  {[}KUS97{]} Kushilevitz, E. \& Nisan, N. (1997). "Communication
  Complexity." Cambridge.
\item
  {[}LEN82{]} Lengauer, T. \& Tarjan, R.E. (1982). "Asymptotically tight
  bounds on time-space trade-offs in a pebble game." JACM 29(4).
\item
  {[}COO76{]} Cook, S.A. \& Sethi, R. (1976). "Storage requirements for
  deterministic polynomial time recognizable languages." JCSS 13.
\item
  {[}BAR91{]} Barrington, D. \& Straubing, H. (1991). "Superlinear lower
  bounds for bounded-width branching programs." Structure in Complexity
  Theory.
\item
  {[}BRY86{]} Bryant, R.E. (1986). "Graph-based algorithms for Boolean
  function manipulation." IEEE Trans. Computers.
\item
  {[}BEN01{]} Ben-Sasson, E. \& Wigderson, A. (2001). "Short proofs are
  narrow---resolution made simple." JACM 48(2).
\end{itemize}

\textbf{Algorithmic Paradigms (how algorithms encounter SCL-like
barriers):}

These paradigms exhibit tradeoffs structurally similar to SCL, though
the precise correspondence varies:

\begin{itemize}
\item
  \textbf{Backtracking/Search}: Tree size \(\geq\) 2\^{}(depth) when
  branches are forced; analogous to \(\Phi\) (active branches) vs. q
  (branches explored).
\item
  \textbf{Dynamic Programming}: Table entries \(\geq\) 2\^{}(subproblem
  diversity); analogous to \(\Phi\) (table size) vs. q (lookups).
\item
  \textbf{CDCL/SAT Solvers}: Proof size \(\geq\)
  2\^{}(\(\Omega\)(width)) via {[}BEN01{]}. \textbf{Caveat}: Clause
  learning with backjumping is non-monotone in "progress"---learned
  clauses can enable jumps that skip exploration. The SCL analogy
  applies to the underlying resolution proof, not directly to the
  solver\textquotesingle{}s execution trace.
\item
  \textbf{Streaming} {[}AMS99{]}: For S bits of memory, passes P satisfy
  P·S \(\geq\) \(\Omega\)(n) for many problems. \textbf{Loose analogy
  only}: if we set \(\Phi\) = log\(_2\)(S) (memory states) and interpret
  R \(\approx\) n (bits needed to solve), the bound P \(\geq\)
  \(\Omega\)(n/S) = \(\Omega\)(R/2\^{}\(\Phi\)) resembles
  SCL\textquotesingle{}s structure. However,
  streaming\textquotesingle{}s "q" would be total bits read (P·n), not
  per-pass, and the tradeoff P·S \(\geq\) \(\Omega\)(n) differs
  structurally from q + \(\Phi\) \(\geq\) R. This is a heuristic
  parallel, not a formal SCL instance. Reference: {[}AMS99{]} Alon, N.,
  Matias, Y., \& Szegedy, M. (1999). "The space complexity of
  approximating the frequency moments." JCSS 58(1).
\end{itemize}

\textbf{SCL as Structural Parallel:}

The Semantic Conservation Law (SCL: q + \(\Phi\) \(\geq\) R) captures a
pattern common to these techniques. The correspondence is:

\begin{itemize}
\tightlist
\item
  \textbf{Decision trees}: q = queries, \(\Phi\) = log\(_2\)(tree
  nodes), R = log\(_2\)(\# distinguishable inputs)
\item
  \textbf{Pebbling}: q = pebble placements, \(\Phi\) = pebble count, R =
  DAG complexity measure
\item
  \textbf{Branching programs}: q = path length, \(\Phi\) =
  log\(_2\)(width), R = log\(_2\)(\# input classes)
\item
  \textbf{Communication}: q = bits exchanged, \(\Phi\) =
  log\(_2\)(rectangles), R = log\(_2\)(partition number)
\item
  \textbf{Resolution}: q = proof length, \(\Phi\) = clause width, R =
  log\(_2\)(search space) --- \emph{width\(\to\)size bridge {[}BEN01{]}
  relates these indirectly; Appendix G details adaptation}
\item
  \textbf{TM observation}: q = bits observed, \(\Phi\) =
  log\(_2\)(configs visited), R = emergence requirement R\_v ---
  \textbf{{[}FORMALIZED in Lean{]}}
\end{itemize}

\textbf{Important Clarification:}

These are \textbf{structural parallels}, not derived instances. Each
technique has its own proof machinery; we do not claim SCL formally
subsumes them. Rather, SCL articulates the common information-theoretic
intuition: \emph{to distinguish among 2\^{}R possibilities, an algorithm
needs R bits of information, obtained either by queries (q) or
pre-stored in states (\(\Phi\))}.

\textbf{This Work\textquotesingle{}s Contribution:}

\begin{enumerate}
\def\labelenumi{\arabic{enumi}.}
\item
  \textbf{Articulates common structure} across prior techniques via SCL
  (q + \(\Phi\) \(\geq\) R):

  \begin{itemize}
  \tightlist
  \item
    5 lower bound techniques: decision trees, communication complexity,
    pebbling, branching programs, resolution
  \item
    4 algorithmic paradigms: backtracking, DP, CDCL, streaming
  \end{itemize}

  This is a \textbf{conceptual unification}---identifying shared
  intuition, not claiming formal derivation of prior results from SCL.
\item
  \textbf{Formalizes TM observation paradigm:} bits observed = q,
  configs visited = 2\^{}\(\Phi\)

  \begin{itemize}
  \tightlist
  \item
    Bridges SCL to information theory (Shannon, parity lower bounds)
  \item
    Connects abstract bounds \(\to\) concrete TM time complexity
  \item
    Enables unconditional P\(\neq\)NP via Structural OWF construction
  \item
    \textbf{This is the only paradigm with mechanized Lean proofs} (see
    below)
  \end{itemize}
\end{enumerate}

\textbf{Formalization Status:}

The TM observation paradigm is \textbf{fully formalized} in Lean
(\texttt{lean/Layer4\_Operational/TimeBridge/}). The structural
parallels for other techniques are \textbf{conceptual
correspondences}---mathematically motivated by the shared intuition
above, but not mechanically verified. Formalizing these as proper SCL
instances (proving each technique\textquotesingle{}s bounds derive from
SCL) is future work (§12.12/F5b).

\textbf{Lean Implementation} (TM observation paradigm):

\begin{itemize}
\tightlist
\item
  \textbf{q (bits observed)}: \texttt{ExecutionPrefixReal.revealedBits}
  --- Bits read from designated addresses
\item
  \textbf{2\^{}\(\Phi\) (configs)}:
  \texttt{ExecutionPrefixReal.computedConfigs} --- Configurations
  visited during execution
\item
  \textbf{Observation extraction}:
  \texttt{TMToExecutionPrefix.lean::tmExecutionToPrefix} --- Converts TM
  trace \(\to\) observations
\item
  \textbf{Run construction}:
  \texttt{TMToExecutionPrefix.lean::buildRunFromTMTrace} --- Builds
  abstract run for SCL analysis
\item
  \textbf{Time bound (QP)}:
  \texttt{TMAdapterQP.lean::fg\_first\_commit\_time\_lower\_bound} ---
  Derives n\^{}\(\Omega\)(log n) bound
\item
  \textbf{Time bound (Exp)}:
  \texttt{TMAdapterExponential.lean::fg\_first\_commit\_time\_lower\_bound}
  --- Derives 2\^{}\(\Omega\)(n) bound
\end{itemize}

\textbf{Why This Works for TM Lower Bounds:}

Fundamental property of deterministic TMs: \emph{A deterministic machine
cannot branch based on data it hasn\textquotesingle{}t read yet.}

Therefore: \textbf{TM Steps \(\geq\) Bits That Must Be Read \(\geq\) R}

The Lean formalization proves this bridge via
\texttt{observation\_semantics\_all\_configs\_on\_tape}.

\textbf{Summary:}

\begin{itemize}
\tightlist
\item
  \textbf{Observation principle}: Not novel --- established 1970s-80s
  (pebbling, decision trees)
\item
  \textbf{SCL as conceptual framework}: Novel --- articulates why all
  techniques share common structure
\item
  \textbf{TM formalization}: Novel --- mechanically verified
  SCL\(\to\)TM time bound bridge
\item
  \textbf{Application to P\(\neq\)NP}: Novel --- unconditional TM time
  bounds via Structural OWF construction
\end{itemize}

\begin{center}\rule{0.5\linewidth}{0.5pt}\end{center}

\subsubsection{12. Discussion and
Implications}\label{12-discussion-and-implications}

The main proof is complete. Sections 1-10 established \textbf{P \(\neq\)
NP} unconditionally for classical uniform PPT via the Structural OWF
construction. We constructed L* with structural properties (A1-A5)
forcing the Semantic Conservation Law (§7), applied FG-wiring to create
per-instance deterministic lower bounds (§8), built an OWF where every
output inherits hardness (§9), and connected to P \(\neq\) NP via the
classical bridge (§10). Section 11 positioned this work within the
literature.

\textbf{Purpose of Section 12:} We now explore the implications, design
space, scope boundaries, and future directions. This section addresses
key questions: What complexity regimes does the framework support? What
exactly was proven and what wasn\textquotesingle{}t? How does this
approach avoid traditional barriers? Why does the semantic layer matter?
How does this relate to natural NP problems? What future work does this
enable?

\textbf{Roadmap:}

\begin{itemize}
\tightlist
\item
  \textbf{§12.1}: Profile Choices and Tight Bounds - the spectrum from
  quasi-polynomial to exponential
\item
  \textbf{§12.2}: Main Lower Bound Result - formal theorem statement and
  verification vs. search dichotomy
\item
  \textbf{§12.3}: Scope and Boundaries - precisely what is proven, model
  restrictions, quantifier structure
\item
  \textbf{§12.4}: Why a Semantic Layer? - the role of instance-side
  structure
\item
  \textbf{§12.5}: Relationship to Natural NP Problems - L* vs. SAT, TSP,
  and other canonical problems
\item
  \textbf{§12.6}: Why Our Approach Avoids Known Barriers -
  naturalization, algebrization circumvention
\item
  \textbf{§12.7}: The Universal Information Skeleton - expository view
  of the conservation framework
\item
  \textbf{§12.8}: Why the Semantic Core Properties Appear Everywhere -
  A1-A5 as universal structure
\item
  \textbf{§12.9}: Quantifier Structure and the Hardcoding Barrier - why
  uniform restriction enables the proof
\item
  \textbf{§12.10}: Paper Roadmap - navigation guide for different
  reading paths
\item
  \textbf{§12.11}: The Semantic Kernel and Future Directions -
  generalizing the framework
\item
  \textbf{§12.12}: Future Work - a \(\lambda\)-program for simplifying
  the complexity zoo
\end{itemize}

\textbf{Why this section matters:} Understanding what was proven, its
boundaries, and its implications is as important as the proof itself.
This section helps researchers understand the result\textquotesingle{}s
scope, identify future research directions, and see how the semantic
conservation framework might apply beyond P vs NP.

\textbf{Scope Reminder:} All class-separation statements in this paper
are for deterministic k-tape TMs with constant k and
\textbar{}\(\Gamma\)\textbar. The bounds apply to our explicitly
constructed NP-complete language L*, not to arbitrary NP problems.

\paragraph{12.1 Profile Choices and Tight
Bounds}\label{121-profile-choices-and-tight-bounds}

The SCL framework (§7) is parametric: any residual profile
\(\lambda\)\_base = \(\omega\)(log n) satisfying the core properties
(A1-A5) produces super-polynomial lower bounds. This generality enables
a spectrum of concrete instantiations, from quasi-polynomial to
exponential hardness. We demonstrate three representative profiles, each
yielding tight bounds via FG-wiring (§8):

\begin{itemize}
\item
  \textbf{QP-sharp} (\(\lambda\)\_base = \(\Theta\)(log² n)):
  Concentrates residual hierarchically across O(log n) DAG depth. FG
  budget \(\tau\)(n) = \(\Theta\)(log n/log² n) = \(\Theta\)(1/log n)
  yields T(n) = n\^{}(\(\Theta\)(log n)) (quasi-polynomial). This
  "barely super-polynomial" regime demonstrates that even minimal
  hardness suffices for P \(\neq\) NP separation.
\item
  \textbf{Sub-exponential} (\(\lambda\)\_base = \(\Theta\)(\ensuremath{\sqrt{n}})): Spreads
  residual across O(\ensuremath{\sqrt{n}}) nodes. FG budget \(\tau\)(n) = \(\Theta\)(log
  n/\ensuremath{\sqrt{n}}) yields T(n) = 2\^{}(\(\Theta\)(\ensuremath{\sqrt{n}})) (sub-exponential).
  Intermediate regime balancing depth and concentration.
\item
  \textbf{Flat} (\(\lambda\)\_base = \(\Theta\)(n)): Maximizes per-node
  residual, forcing linear bottleneck across the cut. FG budget
  \(\tau\)(n) = \(\Theta\)(log n/n) yields T(n) = 2\^{}(\(\Theta\)(n))
  (exponential). Demonstrates framework can achieve maximal hardness.
\end{itemize}

\textbf{Tightness:} Bounds are tight up to polynomial factors - for each
profile, brute-force search achieves T(n) = O(poly(n) · stated\_bound),
while FG forces T(n) = \(\Omega\)(stated\_bound/poly(n)) for any correct
algorithm. The exponential scaling is exact; polynomial factors capture
algorithmic optimization space.

This paper uses the \textbf{QP-sharp profile} for the Structural OWF
construction (§9) and P \(\neq\) NP proof (§10), demonstrating that the
classical bridge (OWF \(\Rightarrow\) FP \(\neq\) FNP \(\Rightarrow\) P
\(\neq\) NP) activates even with minimal super-polynomial hardness. The
spectrum existence shows the framework\textquotesingle{}s generality:
the same construction technique (seed-locked DAG + FG) scales from
n\^{}(\(\Theta\)(log n)) to 2\^{}(\(\Theta\)(n)) by varying residual
distribution alone, with all instances retaining NP-completeness and
per-instance deterministic lower bounds.

Notation reminder. In this subsection, n denotes n\_core (core CNF
size). A concrete micro-example: if n\_core = 2\^{}k and
\(\lambda\)\_base = k² (QP-sharp), then 2\^{}(\(\lambda\)\_base) =
2\^{}(k²) = (2\^{}k)\^{}k = n\_core\^{}(k) = n\_core\^{}(\(\Theta\)(log
n\_core)).

\paragraph{12.2 Main Lower Bound
Result}\label{122-main-lower-bound-result}

We now state the main theorem formally, showing how the interplay
between instance structure (profile residual \(\lambda\)\_base) and
algorithmic resolution (achieved residual \(\rho\)) determines
computational complexity. This theorem quantifies the
verification-search dichotomy and establishes tight bounds up to
polynomial factors.

\textbf{Setup and Notation:}

Assumptions and sizing. Deterministic k-tape TM model (randomized via
coin-fixing), FG wiring present (§6.2.8-§6.2.9). This subsection uses
sizing n := \textbar{}x*\textbar{} (instance size) rather than n :=
n\_core to align with standard complexity notation. Key quantities:

\begin{itemize}
\item
  \textbf{\(\lambda\)\_base(n)} = W\_min(n) · r(n):
  Profile\textquotesingle{}s baseline residual (instance property from
  §12.1)

  \begin{itemize}
  \tightlist
  \item
    W\_min: minimal window parameter (typically \(\Theta\)(1) for flat,
    \(\Theta\)(log n) for QP-sharp)
  \item
    r(n): rank parameter (depth-dependent emergence)
  \end{itemize}
\item
  \textbf{\(\rho\)} = \(\Sigma\)\_\{v\(\in\)C\}(R\_v - q\_v): Effective
  uncommitted residual at bottleneck cut C for a specific algorithm run.
  This is \textbf{algorithm-dependent} - different algorithms achieve
  different \(\rho\) values on the same instance.
\item
  \textbf{FG mechanism parameters} (§8, Appendix C):

  \begin{itemize}
  \tightlist
  \item
    \(\mu\) = \(\Theta\)(1): cut-size coefficient
  \item
    \(\tau\)(n) = \(\Theta\)(log n/\(\lambda\)\_base(n)): per-profile
    budget parameter
  \item
    s \(\leq\) min\{\(\mu\)·\textbar C\textbar, \(\tau\)(n)·\(\rho\)\}:
    pre-final agreement (revealed bits before last segment)
  \end{itemize}
\item
  m\_seg \(\geq\) 2\^{}(\(\rho\)-s): Number of forced rollback segments
  (Segment Counting, Appendix C.2)
\end{itemize}

\textbf{Why residual determines complexity:} At the bottleneck cut C,
there are 2\^{}\(\rho\) seed-consistent worlds. Any correct algorithm
must distinguish among these to avoid errors. Polynomial-time algorithms
maintain \(\leq\) 2\^{}(O(log n)) states, so when \(\rho\) =
\(\omega\)(log n), the gap forces 2\^{}(\(\rho\)-O(log n))
super-polynomial explorations. (See §2.7: configuration space vs
resource space.)

\textbf{The Verification-Search Dichotomy:}

The residual \(\rho\) determines computational regime:

\begin{itemize}
\item
  \textbf{Verification} (\(\rho\) = 0): All cut nodes fully resolved
  (q\_v = R\_v for all v \(\in\) C). Algorithm has complete information
  - no search needed. Time: polynomial in n.
\item
  \textbf{Near-verification} (0 \textless{} \(\rho\) \(\leq\) O(log n)):
  Minimal unresolved residual, polynomial states suffice. Time: still
  polynomial.
\item
  \textbf{Search} (\(\rho\) \(\geq\) (log n)\^{}(1+\(\delta\)) for any
  \(\delta\) \textgreater{} 0): Substantial unresolved residual forces
  exploration of 2\^{}(\(\rho\)-s) segments. Time: super-polynomial.
\end{itemize}

The sharp cliff occurs at \(\rho\) = \(\omega\)(log n), where polynomial
resources (\(\leq\) 2\^{}(O(log n)) states) become insufficient to cover
the 2\^{}\(\rho\) configuration space.

\textbf{Theorem 12.2 (Tight Bounds for L*).} \emph{(Here n measures
\textbar{}x}\textbar{}, the instance size.)*

\textbf{Upper bound (algorithm-dependent):} For any algorithm that
achieves residual \(\rho\) on instance x* (i.e., resolves information
down to \(\rho\) unresolved bits at the bottleneck cut), there exists a
deterministic procedure deciding L* in time

T(n) = O(2\^{}\(\rho\) · poly(n\_core)) = O(2\^{}\(\rho\) · poly(n))

\emph{Intuition: Enumerate 2\^{}\(\rho\) remaining worlds, verify each
in poly-time.}

\textbf{Lower bound (per-instance, any fixed run):} For any
deterministic k-tape TM operating in the \textbf{search lane} (\(\rho\)
\(\geq\) (log n)\^{}(1+\(\delta\))), the FG mechanism bounds
pre-revelation as s \(\leq\) min\{\(\mu\)·\textbar C\textbar,
\(\tau\)(n)·\(\rho\)\}, forcing time

T(n) \(\geq\) 2\^{}(\(\rho\)-s) · \(\Omega\)(n/W\_min(n)).

\emph{Intuition: Must perform m\_seg \(\geq\) 2\^{}(\(\rho\)-s) rollback
segments, each costing \(\Omega\)(n/W\_min) for parity evaluation.}

\textbf{Corollary 12.2.1 (Universal Tightness via \(\mu\)-\(\tau\) FG).}

With FG parameters \(\mu\) = \(\Theta\)(1) and \(\tau\)(n) =
\(\Theta\)(log n / \(\lambda\)\_base(n)) (from §12.1), the
pre-revelation bound s \(\leq\) \(\tau\)(n)·\(\rho\) typically dominates
(when \(\rho\) is large enough that \(\tau\)(n)·\(\rho\) \textgreater{}
\(\mu\)·\textbar C\textbar). This yields:

T\_lower(n) \(\geq\) 2\^{}(\(\rho\)-\(\tau\)(n)·\(\rho\)) ·
\(\Omega\)(n/W\_min) = 2\^{}((1-\(\tau\)(n))\(\rho\)) ·
\(\Omega\)(n/W\_min) = 2\^{}(\(\rho\))/poly(n\_core) {[}since
\(\tau\)(n) = \(\Theta\)(log n/\(\lambda\)\_base) and \(\lambda\)\_base
= \(\omega\)(log n){]}

Comparing with the upper bound T\_upper(n) \(\leq\) 2\^{}\(\rho\) ·
poly(n\_core), we see the bounds match up to polynomial factors - the
exponential scaling 2\^{}\(\rho\) is exact. This means:

\begin{itemize}
\tightlist
\item
  Any algorithm achieving residual \(\rho\) can be implemented within
  O(2\^{}\(\rho\) · poly(n))
\item
  No algorithm can do better than \(\Omega\)(2\^{}\(\rho\) / poly(n))
  due to FG
\item
  The only optimization space is in the polynomial factors
\end{itemize}

\textbf{Phase Transition (Complete Spectrum):}

\begin{itemize}
\tightlist
\item
  \(\rho\) = 0: Pure verification \(\to\) polynomial time
\item
  \(\rho\) = O(log n): Near-verification \(\to\) still polynomial
\item
  \(\rho\) = \(\omega\)(log n): \textbf{Sharp cliff} \(\to\)
  super-polynomial (sharp phase transition)
\item
  \(\rho\) = \(\Theta\)(log² n): Quasi-polynomial n\^{}(\(\Theta\)(log
  n)) (minimal super-poly, used in §9-10)
\item
  \(\rho\) = \(\Theta\)(\ensuremath{\sqrt{n}}): Sub-exponential 2\^{}(\(\Theta\)(\ensuremath{\sqrt{n}}))
\item
  \(\rho\) = \(\Theta\)(n): Exponential 2\^{}(\(\Theta\)(n)) (maximal
  hardness)
\end{itemize}

\textbf{Proof Sketch.}

\emph{(i) Upper bound.} Given an algorithm with residual \(\rho\) (i.e.,
\(\rho\) unresolved bits at cut C), construct a decision procedure:

\begin{enumerate}
\def\labelenumi{\arabic{enumi}.}
\tightlist
\item
  Enumerate all 2\^{}\(\rho\) feasible assignments to the cut nodes
\item
  For each assignment, run the polynomial-time verifier (Algorithm V,
  §10.2)
\item
  Accept if any assignment verifies
\end{enumerate}

Since \textbar{}x*\textbar{} = O(n\_core log n\_core) and verification
is poly(\textbar{}x*\textbar{}) = poly(n\_core), total time is
2\^{}\(\rho\) · poly(n\_core). \ensuremath{\square}

\emph{(ii) Lower bound (single run).} Combine three ingredients:

\textbullet{} \textbf{FG Pre-Revelation Bound (Appendix C.1):} At the start of the
final segment on each cut-gate path, at most s \(\leq\)
min\{\(\mu\)·\textbar C\textbar, \(\tau\)(n)·\(\rho\)\} bits are
revealed pre-final. Otherwise, the algorithm already paid
\(\Omega\)(n/W\_min) time per gate.

\textbullet{} \textbf{Segment Counting (Appendix C.2):} The number of rollback
segments satisfies m\_seg \(\geq\) 2\^{}(\(\rho\)-s). When
\(\tau\)(n)·\(\rho\) dominates, m\_seg \(\geq\)
2\^{}((1-\(\tau\)(n))\(\rho\)).

\textbullet{} \textbf{Parity Cost (Appendix C.1.1):} Each segment incurs
\(\Omega\)(n/W\_min) time for parity gate evaluation (profile-tight).

Multiplying: T \(\geq\) m · \(\Omega\)(n/W\_min) \(\geq\)
2\^{}(\(\rho\)-s) · \(\Omega\)(n/W\_min). \ensuremath{\square}

\textbf{Key Observations:}

\begin{enumerate}
\def\labelenumi{\arabic{enumi}.}
\tightlist
\item
  \textbf{Role separation:} Instance fixes \(\lambda\)\_base (structural
  property); algorithm\textquotesingle{}s behavior determines \(\rho\)
  (resolution achieved)
\item
  \textbf{Sharp transition:} \(\rho\) = \(\omega\)(log n) \(\to\)
  Segment Counting forces super-polynomial time (polynomial resources
  insufficient)
\item
  \textbf{FG mechanism:} Prevents pre-accumulation of cut information,
  ensuring survival of exponential factor 2\^{}((1-\(\tau\)(n))\(\rho\))
  \(\approx\) 2\^{}\(\rho\)/poly (see §8 and Appendix C for s \(\leq\)
  \(\Theta\)(\(\tau\)·\(\rho\)) and \(\lambda\)(A,x) \(\geq\)
  \(\lambda\)\_base \(-\) o(\(\lambda\)\_base)).
\item
  \textbf{Universality:} Bound applies to any correct algorithm on
  FG-wired instances, regardless of computational strategy
\end{enumerate}

\textbf{The Complete Parametric Spectrum:}

\begin{itemize}
\item
  \textbf{\(\rho\) = 0}

  \begin{itemize}
  \tightlist
  \item
    Complexity Class: Polynomial
  \item
    Regime: Pure verification
  \end{itemize}
\item
  \textbf{\(\rho\) = O(log n)}

  \begin{itemize}
  \tightlist
  \item
    Complexity Class: Polynomial
  \item
    Regime: Near-verification
  \end{itemize}
\item
  \textbf{\(\rho\) = \(\omega\)(log n)}

  \begin{itemize}
  \tightlist
  \item
    Complexity Class: \textbf{Sharp cliff}
  \item
    Regime: \textbf{Super-polynomial threshold}
  \end{itemize}
\item
  \textbf{\(\rho\) = \(\Theta\)(log² n)}

  \begin{itemize}
  \tightlist
  \item
    Complexity Class: n\^{}(\(\Theta\)(log n))
  \item
    Regime: Quasi-poly (minimal super-poly)
  \end{itemize}
\item
  \(\rho\) = \(\Theta\)(n\^{}\(\varepsilon\)),
  0\textless{}\(\varepsilon\)\textless1

  \begin{itemize}
  \tightlist
  \item
    Complexity Class: 2\^{}(\(\Theta\)(n\^{}\(\varepsilon\)))
  \item
    Regime: Sub-exponential
  \end{itemize}
\item
  \textbf{\(\rho\) = \(\Theta\)(n)}

  \begin{itemize}
  \tightlist
  \item
    Complexity Class: 2\^{}(\(\Theta\)(n))
  \item
    Regime: Exponential (maximal)
  \end{itemize}
\end{itemize}

The transition at \(\rho\) = \(\omega\)(log n) is discontinuous:
crossing this threshold jumps from polynomial to super-polynomial with
no intermediate regime.

\textbf{Scope and Extensions:}

\begin{itemize}
\item
  \textbf{Scope:} These bounds apply specifically to L* (our explicitly
  constructed NP-complete language), not to arbitrary NP problems.
  Extending to other problems would require constructing analogous
  overlay structures satisfying A1-A5.
\item
  \textbf{Randomized algorithms:} Covered via Yao\textquotesingle{}s
  minimax principle (coin-fixing, §9.4). The per-run bound extends to
  expected time for randomized algorithms.
\item
  \textbf{Generality:} The theorem template applies to any profile with
  \(\lambda\)\_base(n) = \(\omega\)(log n), yielding super-polynomial
  bounds. The specific regime (QP to exponential) depends on
  \(\lambda\)\_base scaling.
\item
  \textbf{OWF connection:} The QP-sharp profile (\(\lambda\)\_base =
  \(\Theta\)(log² n)) is used in the Structural OWF construction (§9),
  showing even minimal super-polynomial hardness suffices for P \(\neq\)
  NP separation via the classical bridge.
\end{itemize}

\begin{center}\rule{0.5\linewidth}{0.5pt}\end{center}

\paragraph{12.3 Scope and Boundaries}\label{123-scope-and-boundaries}

This subsection clarifies precisely what was proven and what remains
open. Understanding these boundaries is crucial for interpreting the
result correctly and identifying future research directions.

\textbf{What we establish:} For our explicitly constructed NP-complete
language L* (§6), the structural properties A1-A5 mathematically force
the Semantic Conservation Law q + \(\Phi\) \(\geq\) R (where \(\Phi\) =
log\(_2\)(Alt) counts distinguishable artifacts maintained by the
algorithm) on all classical deterministic algorithms. This is not a
heuristic or average-case observation - it\textquotesingle{}s a proven
mathematical necessity arising from L*\textquotesingle{}s engineered
structure. Combined with the FG mechanism (§8), this yields per-instance
deterministic lower bounds enabling the Structural OWF construction (§9)
and unconditional P \(\neq\) NP separation (§10) for classical uniform
PPT.

Below we detail what was proven (\textbf{Results}) and important
limitations (\textbf{Boundaries}).

\textbf{Results (What We Proved):}

\begin{enumerate}
\def\labelenumi{\arabic{enumi}.}
\item
  \textbf{P \(\neq\) NP for Classical Uniform PPT} (Main Result,
  §9-§10): We constructed an unconditional one-way function f from L*
  (an explicitly constructed NP-complete language) and established FP
  \(\neq\) FNP \(\Rightarrow\) P \(\neq\) NP via the classical bridge.
  The separation applies to deterministic k-tape Turing machines with
  constant k (number of tapes) and constant
  \textbar{}\(\Gamma\)\textbar{} (alphabet size), and extends to
  randomized PPT via coin-fixing.
\item
  \textbf{Parametric Complexity Spectrum} (§12.1-12.2): The same
  construction technique yields a full spectrum of lower bounds from
  quasi-polynomial (n\^{}(\(\Theta\)(log n))) to exponential
  (2\^{}(\(\Theta\)(n))) by varying the residual profile
  \(\lambda\)\_base, with tight bounds up to polynomial factors in each
  regime.
\item
  \textbf{Cross-Paradigm Lower Bounds} (§5, §7): The SCL framework
  yields exponential lower bounds across all major computational
  paradigms - backtracking tree size, DP table size, resolution proof
  size, OBDD width - with the same min-cut residual \(\lambda\)
  explaining all bounds. OBDD bounds are order-robust (hold for any
  variable ordering) via expander-FG gates.
\item
  \textbf{Sharp Verification-Search Dichotomy} (§10, Theorem 12.2): L*
  exhibits NP-completeness in its purest form: polynomial-time
  verification with witness (\(\rho\) = 0) vs. super-polynomial search
  without witness (\(\rho\) \(\geq\) \(\omega\)(log n)), with a sharp
  phase transition at \(\rho\) = \(\omega\)(log n) where polynomial
  resources become insufficient.
\end{enumerate}

\textbf{Boundaries (Important Limitations):}

\begin{enumerate}
\def\labelenumi{\arabic{enumi}.}
\item
  \textbf{L*-Specific (Not General NP):} Our lower bounds and Structural
  OWF construction apply specifically to L*, the NP-complete language we
  explicitly constructed in §6. Extending to other NP-complete problems
  would require constructing analogous overlay structures satisfying
  properties A1-A5 (Hermeticity, Injectivity, Emergence, Closure,
  Dependency). We do not claim all NP problems have this structure - L*
  was carefully engineered to exhibit it.
\item
  \textbf{Classical Computation Only:} Results apply to classical
  deterministic k-tape TMs (with constant k and
  \textbar{}\(\Gamma\)\textbar) and extend to classical randomized PPT
  via coin-fixing (Yao\textquotesingle{}s principle). We do NOT claim
  separation for:

  \begin{itemize}
  \tightlist
  \item
    Quantum algorithms (BQP vs NP remains open)
  \item
    Algorithms with oracles or advice (non-uniform models can trivially
    hardcode solutions)
  \item
    Other computational models (circuits, RAMs, parallel models) without
    adapter construction
  \end{itemize}
\item
  \textbf{Uniform Algorithms Only:} The \(\forall\)x* \(\forall\)A
  quantifier structure (every FG-wired instance is hard for every
  uniform algorithm) is essential for the Structural OWF construction.
  Non-uniform algorithms can hardcode instance-specific solutions,
  defeating per-instance hardness. This is why we prove P \(\neq\) NP
  for uniform PPT, not P/poly \(\neq\) NP.
\item
  \textbf{Per-Instance Deterministic (Not Distributional Average-Case):}
  The main result (§8-§10) uses per-instance deterministic lower bounds:
  every FG-wired instance x* requires super-polynomial time on any fixed
  run. This is stronger than average-case hardness.
\item
  \textbf{Algorithm-Dependent Residual:} The bound
  2\^{}(\(\Omega\)(\(\lambda\)(A,x))) depends on the residual
  \(\lambda\)(A,x) achieved by algorithm A on instance x*. Different
  algorithms may achieve different \(\lambda\) values on the same
  instance (e.g., clever resolution vs. naive backtracking). For
  QP-sharp profiles, FG calibration implies \(\lambda\)(A,x) \(\geq\)
  \(\lambda\)\_base \(-\) o(\(\lambda\)\_base) = \(\Theta\)(log² n) for
  all correct algorithms on FG-wired instances (see §8 and Appendix C),
  hence the lower bound is super-polynomial uniformly; only polynomial
  factors vary across algorithms.
\item
  \textbf{Monotonicity (Robustness Property):} If an analysis
  underestimates resolution (q\_v too low), the residual c\_v = R\_v -
  q\_v is overestimated, yielding a weaker but still valid lower bound.
  This monotonicity ensures conservative analyses remain sound, though
  potentially not tight.
\end{enumerate}

\textbf{Implications for Complexity Theory:}

This result demonstrates that P \(\neq\) NP can be proven for specific
computational models using constructive instance engineering, without
relying on unproven conjectures or navigating around traditional
barriers (naturalization, algebrization). The approach:

\begin{itemize}
\tightlist
\item
  \textbf{Shifts focus} from algorithm limitations to instance
  requirements
\item
  \textbf{Articulates common structure} across paradigm-specific lower
  bounds via information-theoretic framework
\item
  \textbf{Opens pathway} to model-specific unconditional separations
\item
  \textbf{Provides template} for other complexity separations via
  structural conservation laws
\end{itemize}

While model-specific, the result is unconditional within its scope - a
rare achievement in complexity theory.

\textbf{Open Questions and Future Work:}

The main open questions fall into three categories:

\begin{enumerate}
\def\labelenumi{\arabic{enumi}.}
\item
  \textbf{Classical model extensions:} The SCL framework is
  paradigm-independent (§5, §7), so extending to circuits, RAMs, or
  parallel models should be straightforward work applying existing
  adapters. We scoped to k-tape TMs for presentation manageability.
  (Detailed discussion in §12.11.)
\item
  \textbf{Quantum computation:} Unlike classical extensions, quantum
  resistance is genuinely uncertain - quantum computation may sidestep
  the SCL bottleneck through superposition. (BQP vs NP remains open; see
  §12.11.)
\item
  \textbf{Natural problems and other separations:} Can natural
  NP-complete problems satisfy A1-A5? Can SCL-style arguments separate
  other complexity classes? (See §12.5 for natural problem
  correspondence, §12.11 for semantic kernel generalization, and §12.12
  for the \(\lambda\)-program vision.)
\end{enumerate}

\textbf{Formalization Trust Boundary:}

The accompanying Lean 4 formalization (90 files, \textasciitilde{}90,000
lines, available at GitHub repository) makes the complete proof chain
machine-checkable with full transparency about the trust boundary. The
formalization reveals that the proof rests on two foundational
assumptions:

\begin{enumerate}
\def\labelenumi{\arabic{enumi}.}
\item
  \textbf{Church-Turing thesis with polynomial simulation} (standard):
  Every uniform polynomial-time algorithm has a polynomial-time Turing
  machine encoding. This is the standard foundational axiom accepted
  throughout complexity theory.
\item
  \textbf{Semantic\(\to\)Operational bridge}: For planted instances with
  well-formed randomness, if a Turing machine produces a correct
  witness, then its encoder must have visited all 2\^{}R emergent
  configuration values during execution.
\end{enumerate}

\textbf{The gap:} Layers 0-3 of the formalization prove
information-theoretic necessity---producing a correct witness requires
distinguishing which of 2\^{}R informationally-distinct configurations
is correct (Lemma C.1.2). However, connecting this semantic requirement
to operational Turing machine execution traces---that the
encoder\textquotesingle{}s state sequence must have represented all
2\^{}R values---involves model-specific details we axiomatize rather
than prove.

\textbf{Trust boundary:} The formalization axiomatizes this connection
explicitly. Together with the Church-Turing thesis, these constitute the
complete trust boundary. The formalization has 0 sorries, confirming no
other gaps exist.

\begin{center}\rule{0.5\linewidth}{0.5pt}\end{center}

\paragraph{12.4 Why a Semantic Layer?}\label{124-why-a-semantic-layer}

\textbf{The Barrier Prior Approaches Created:} Classical lower bounds
(Ben-Sasson \& Wigderson {[}BEN01{]} for resolution, Wegener {[}WEG00{]}
for OBDDs, etc.) proved exponential complexity within specific
computational models. Each bound required distinct proof techniques
tailored to that model\textquotesingle{}s structure. This
paradigm-specificity created a fundamental obstacle: a polynomial-time
algorithm could potentially circumvent any single
model\textquotesingle{}s lower bound by using a different computational
approach. No prior technique could rule out \emph{all} polynomial-time
strategies simultaneously - the requirement for proving P \(\neq\) NP.

\textbf{What the Semantic Layer Solves:} The semantic layer provides a
\textbf{model-independent metric} - distinguishable artifacts Alt\_v at
each node - that measures the same information-theoretic bottleneck
(2\^{}\(\lambda\) unresolved seed-consistent worlds at the min-cut)
regardless of which computational model attempts to solve L*. This
enables an \textbf{adapter architecture}: we prove the conservation law
q\_v + \(\Phi\)\_v \(\geq\) R\_v once for the abstract framework
(§7.2.1), then provide thin model-specific adapters showing how the same
residual \(\lambda\) translates into different metrics (tree nodes for
backtracking, DP keys, resolution clauses, OBDD states, tape head
motions for TMs).

\textbf{Model Portability:} Because artifacts are defined semantically
(each artifact = one distinguishable hypothesis about the
overlay\textquotesingle{}s seed-dependency structure, independent of how
the algorithm internally represents state), the bound \(\lambda\)
\(\geq\) \(\omega\)(log n) survives model changes. For example: A
backtracking algorithm must maintain \(\geq\) 2\^{}\(\lambda\) tree
nodes to track unresolved seeds at the min-cut; a DP algorithm must
maintain \(\geq\) 2\^{}\(\lambda\) distinct memoization keys (since
different seed histories require separate keys for correctness); an OBDD
must maintain \(\geq\) 2\^{}\(\lambda\) simultaneously-live states when
processing variables across the min-cut. All three are manifestations of
the same 2\^{}\(\lambda\) seed-worlds bottleneck, measured in different
computational metrics.

\textbf{Composition and Auditability:} The semantic layer composes
multiplicatively across DAG structure: for any cut C, Alt(C) =
\(\Pi\)\_\{v\(\in\)C\} Alt\_v (since artifacts at different nodes
represent independent seed choices). This compositional property -
combined with RWA (Receiving-Window Attribution, §6) which provides
order-invariant resolution accounting - makes the bottleneck
\textbf{auditable}: we can prove that maintaining \textless{}
2\^{}\(\lambda\) artifacts at the cut necessarily creates resolution
deficits (q\_v + \(\Phi\)\_v \textless{} R\_v for some v \(\in\) C),
which mathematically force correctness errors via Lemma 7.I
(insufficient distinguishability \(\to\) collision \(\to\) misidentified
seeds).

\textbf{The Content-Addressed Encoder:} The mechanism Seed\_v = Enc(v
\textbar{}\textbar{} sorted\{(u, Seed\_u, y\_u)\} \textbar{}\textbar{}
GateDigest\_v) makes artifact-counting model-independent by ensuring
that "distinguishable artifacts" corresponds directly to distinguishable
seeds (via Injectivity A2). Any two computation paths that resolve
different parent seeds \emph{must} track them as separate artifacts
(cannot merge without full resolution), regardless of the
algorithm\textquotesingle{}s internal representation. This seed-artifact
correspondence is what allows the abstract bound (2\^{}\(\lambda\)
distinguishable artifacts required) to translate mechanically into
model-specific bounds (2\^{}\(\lambda\) tree nodes, DP keys, resolution
clauses, etc.).

\textbf{P vs NP-Hard Dichotomy via Semantic Structure:} The semantic
layer clarifies why P and NP-hard problems differ fundamentally:

\begin{itemize}
\item
  \textbf{P problems (polynomial-time solvable):} Algorithms "discharge
  residual capacity early" - they resolve most required bits q\_v
  \(\approx\) R\_v at each node through propagation, learning, or
  structural simplification, keeping the min-cut residual
  \(\lambda\)(C*) \(\leq\) O(log n). For example, in L* verification
  with a witness (Algorithm V, §10.2), the witness provides all overlay
  seeds directly, so q\_v = R\_v at every node \(\to\) \(\lambda\) = 0
  \(\to\) polynomial time. The semantic assumptions (Emergence, Closure,
  Dependency, Injectivity) are satisfied but \emph{no bottleneck forms}
  because resolution proceeds uniformly.
\item
  \textbf{L* without witness (NP-hard):} The instance\textquotesingle{}s
  engineered structure - specifically properties A1-A5 ensuring
  injective seed-chaining with rank-forcing across O(log n) depth -
  creates a persistent min-cut where \(\lambda\) \(\geq\) \(\omega\)(log
  n). No polynomial-time strategy can reduce this residual because doing
  so would require either: (1) violating Injectivity (merging distinct
  seed-histories prematurely \(\to\) correctness errors per Lemma 7.I),
  or (2) discovering enough propagated information to collapse the
  min-cut (which would constitute a polynomial-time witness-finding
  algorithm, contradicting the OWF\textquotesingle{}s per-instance lower
  bound via Ext, §9.3).
\item
  \textbf{Diagnostic principle:} If a language in P appears to force
  large \(\lambda\)(A,x), at least one semantic assumption
  (Emergence/Closure/Dependency/Injectivity) must fail - either the
  instance structure doesn\textquotesingle{}t actually require
  2\^{}\(\lambda\) distinguishable seeds at the bottleneck, or the
  algorithm has access to shortcuts (advice, witnesses, structural
  bypasses) that reduce the effective residual.
\end{itemize}

\textbf{Why This Enables P \(\neq\) NP:} The adapter architecture means
the exponential bottleneck (2\^{}\(\lambda\) seed-worlds at
L*\textquotesingle{}s min-cut, where \(\lambda\) \(\geq\) \(\omega\)(log
n)) cannot be circumvented by switching computational models - each
model must account for the same 2\^{}\(\lambda\) unresolved
possibilities, just in different metrics. Combined with the Structural
OWF construction (every f(r) output is FG-wired \(\to\) per-instance
deterministic witness-finding lower bound, §8-9), this yields an
unconditional separation: one-way functions exist \(\to\) FP \(\neq\)
FNP \(\to\) P \(\neq\) NP via the classical bridge (§10.4-10.5).

\begin{center}\rule{0.5\linewidth}{0.5pt}\end{center}

\paragraph{12.5 Relationship to Natural NP
Problems}\label{125-relationship-to-natural-np-problems}

\textbf{Why This Question Matters:} Natural NP-complete problems like
3-SAT, Graph Coloring, and TSP are widely believed to be hard - their
intractability underpins cryptography, approximation theory, and
parameterized complexity. Yet despite decades of effort, no one has
proven rigorous super-polynomial lower bounds for these problems on
general Turing machines. The obstacle is mathematical, not empirical:
natural problems exhibit the same information-theoretic tradeoffs
(resolve early vs track exponentially many possibilities) that the SCL
formalizes, but their messy structure - variable symmetries, non-uniform
dependencies, entangled constraints - makes rigorous cross-paradigm
analysis intractable. We cannot cleanly define "required bits R\_v at
node v" when variables appear irregularly, or track "distinguishable
artifacts Alt\_v" when symmetries merge partial solutions unpredictably.

\textbf{The Barrier Natural Problems Create:} Natural problems lack the
mathematical scaffolding needed to prove paradigm-invariant lower bounds
rigorously today. Specifically:

\begin{itemize}
\item
  \textbf{3-SAT:} Variables appear in clauses with non-uniform frequency
  and polarity patterns. There\textquotesingle{}s no clean DAG structure
  with well-defined "R\_v fresh bits emerge at node v" - variables are
  global, and dependencies form a tangled hypergraph rather than a
  layered hierarchy. This makes it intractable to cleanly formalize a
  cut residual \(\lambda\) or prove compositional bounds across
  paradigms (obstructs crisp A3/A5 formalization).
\item
  \textbf{Graph Coloring:} Symmetries (automorphisms of the input graph)
  allow algorithms to merge partial colorings that differ only by
  relabeling. This violates the Injectivity property (A2) - distinct
  "configurations" no longer guarantee distinct "seeds" because
  symmetry-equivalent states can be treated as identical. We cannot
  reliably count distinguishable artifacts Alt\_v (violates A2 via
  symmetry quotienting).
\item
  \textbf{TSP (Traveling Salesman):} Global constraints (the tour must
  visit all cities exactly once) create long-range dependencies that
  prevent clean DAG decomposition. There\textquotesingle{}s no localized
  "node v" where exactly R\_v bits emerge - the constraint structure is
  inherently global, precluding the layered min-cut analysis that SCL
  requires (obstructs clean A3/A5 and H1 factoring).
\item
  \textbf{CSPs (Constraint Satisfaction Problems):} Variables
  participate in multiple overlapping constraints, creating redundant
  information paths. A single variable assignment can propagate through
  many constraints simultaneously, violating Hermeticity (A1) and making
  RWA accounting ambiguous (which constraint "receives" the resolved
  information first?) (violates A1; RWA ambiguity).
\end{itemize}

\textbf{How L* Purifies These Obligations:} L* removes the mathematical
obstacles while preserving the fundamental hardness mechanism. It
provides a "cleaned-up" version of NP-completeness where the
conservation law can be proven rigorously, not just observed
empirically.

\textbf{What L* Provides (Mathematical Scaffolding for Proofs):}

\begin{enumerate}
\def\labelenumi{\arabic{enumi}.}
\item
  \textbf{Clean Emergence (A3):} Exactly R\_v fresh bits emerge at each
  node v via rank(H\_v) = R\_v. This makes "required information"
  well-defined and compositional - we can prove
  \(\Sigma\)\_\{v\(\in\)C\} R\_v for any cut C, enabling min-cut
  analysis. Natural problems lack this: in 3-SAT, how many "fresh bits"
  does fixing variable x\_17 provide? The answer depends on clause
  overlap, propagation, and history - there\textquotesingle{}s no clean
  local measure.
\item
  \textbf{Explicit Dependency (A5 + Closure A4):} The seed encoding
  Seed\_v = Enc(v \textbar{}\textbar{} sorted\{(u, Seed\_u, y\_u)\}
  \textbar{}\textbar{} GateDigest\_v) makes information flow explicit
  and recoverable. Every parent\textquotesingle{}s seed and output
  appears in Seed\_v, creating a verifiable dependency chain. This
  enables compositional artifact-counting: Alt(C) =
  \(\Pi\)\_\{v\(\in\)C\} Alt\_v via Cartesian factoring (Lemma J.1-Cart)
  - seed choices at different cut nodes are provably independent due to
  disjoint address pools (A1 Hermeticity) + Dependency (A5) + no
  cross-coupling (formalized as H1-H5 properties in §7.2.1). Natural
  problems lack this structural clarity - dependencies are implicit in
  constraint entanglement.
\item
  \textbf{Hermeticity (A1):} Information flows only through designated
  DAG edges with disjoint address pools \{U\_v\}. This eliminates
  "hidden channels" where algorithms might resolve information without
  our accounting detecting it. For natural problems,
  there\textquotesingle{}s no clean boundary: a 3-SAT
  solver\textquotesingle{}s unit propagation can spread information
  through arbitrary clause paths, making it impossible to audit which
  node "receives" resolved bits first (RWA ambiguity).
\item
  \textbf{Injectivity (A2):} The injective encoder ensures distinct
  parent-tuples \(\to\) distinct seeds \(\to\) distinct artifacts. This
  makes the lower bound Alt\_v \(\geq\) 2\^{}(R\_v-q\_v) provable (Lemma
  7.I) - we can \emph{count} distinguishable artifacts reliably,
  establishing a tight mathematical floor. Natural
  problems\textquotesingle{} symmetries break this: two graph colorings
  differing only by color-label permutation represent the same
  "artifact" from an algorithm\textquotesingle{}s perspective, but we
  cannot formalize which configurations count as "distinct."
\item
  \textbf{Auditable Accounting (RWA, §6):} Receiving-Window Attribution
  provides order-invariant resolution credit. Regardless of which
  sequence an algorithm uses to resolve seeds, RWA assigns clear q\_v
  values based on first-use of designated address reads. This makes the
  conservation law q\_v + \(\Phi\)\_v \(\geq\) R\_v provable for any
  execution order. Natural problems lack designated addresses - we
  cannot reliably track "when" information gets resolved.
\end{enumerate}

\textbf{The Payoff:} These five properties enable rigorous proofs that
natural problem hardness makes impossible. We can prove (not just
conjecture) that maintaining \textless{} 2\^{}\(\lambda\) artifacts at
the min-cut mathematically forces correctness errors (Lemma 7.I), that
this bottleneck survives across all computational paradigms (§5
adapters), and that the FG mechanism creates per-instance deterministic
witness-finding lower bounds (§8). Natural problems exhibit similar
phenomena, but we cannot prove it - the mathematical structure is too
tangled.

\textbf{Empirical Validation (Practice Mirrors Theory):} Despite their
messy mathematical structure, natural problems exhibit the same
resolve-early-vs-track-exponentially tradeoffs that SCL formalizes.
Practitioners observe this empirically when choosing algorithms:

\begin{itemize}
\item
  \textbf{3-SAT solvers:} DPLL (high commitment - eagerly propagate and
  branch, risk backtracking) vs pure backtracking (low commitment -
  enumerate exhaustively) vs CDCL (adaptive - learn clauses to avoid
  redundant exploration). The tradeoff is fundamentally q vs \(\Phi\):
  resolve more via learning (increase q) or maintain more search
  branches (increase \(\Phi\)).
\item
  \textbf{Graph Coloring:} Greedy heuristics (early resolution - commit
  to color assignments quickly, risk failure on hard instances) vs exact
  algorithms (track all partial colorings explicitly, exponential
  space). Again: resolve early (increase q, low \(\Phi\), fast but
  incomplete) or track comprehensively (low q, high \(\Phi\), complete
  but slow).
\item
  \textbf{TSP:} Nearest-neighbor heuristic (O(n²) time, local greedy
  choices) vs Held-Karp dynamic programming (2\^{}n space, exhaustive
  subproblem tracking). The DP algorithm pays exponential space
  (\(\Phi\)) to avoid resolving the global tour structure prematurely.
\end{itemize}

These are not accidents - they reflect the underlying conservation law q
+ \(\Phi\) \(\geq\) R. Natural problems obey it; we just cannot prove it
rigorously due to their tangled structure. (For detailed mapping of
natural NP families to SCL framework, see §12.5.1 below.)

\textbf{L* as a Mathematical Purification:} Like the Ising model in
statistical physics - a simplified lattice system that illuminates phase
transitions and magnetism by stripping away real-material complexities -
L* provides a "toy model" of NP-completeness that reveals the
fundamental mechanism. The Ising model is not realistic iron, but it
enabled rigorous theorems (critical exponents, universality classes)
impossible to prove for real magnets. Similarly, L* is not a natural
problem, but its clean structure enables rigorous proofs
(paradigm-invariant lower bounds, Structural OWF construction, P
\(\neq\) NP) impossible to achieve for 3-SAT or TSP. The insight:
hardness arises from \textbf{information debt accumulation across
dependency chains} - when problem structure forces algorithms to either
resolve information prematurely (risking correctness errors) or track
exponentially many unresolved possibilities (paying exponential cost).
L* makes this tradeoff mathematically explicit and provable.

\subparagraph{12.5.1 Natural NP Families and SCL
Correspondence}\label{1251-natural-np-families-and-scl-correspondence}

Many natural NP-complete families already exhibit the phenomenon our SCL
formalizes: all constraints are published upfront, but
correctness-relevant information becomes actionable only after
prerequisite decisions (dependencies) are satisfied. Below are
representative examples showing how they map to SCL. While they are
close analogues, they do not include our L*-specific wiring (disjoint
atoms + FG) that yields unconditional TM time bounds.

\textbf{Pebbling CNFs (Graph Pebbling Encoded as SAT)} {[}SAV98{]}

\begin{itemize}
\tightlist
\item
  \textbf{Upfront}: The graph and pebbling constraints are explicit
\item
  \textbf{Gated action}: A node can be pebbled only after its
  predecessors - a natural dependency chain (Keyedness analogue via
  frontier configurations)
\item
  \textbf{Known bounds}: Time/space tradeoffs; exponential proof/width
  bounds for certain graphs
\item
  \textbf{SCL lens}: Frontier = distinguishable artifacts; resolving
  prerequisites raises q; otherwise Alt multiplies along cuts
\end{itemize}

\textbf{Tseitin Formulas on Expanders (Parity Constraints)} {[}BEN01,
HOO06{]}

\begin{itemize}
\tightlist
\item
  \textbf{Upfront}: All parity equations are given
\item
  \textbf{Gated action}: Parity at a cut emerges only after enough local
  bits are fixed; expansion forces wide simultaneous tracking
\item
  \textbf{Known bounds}: Width\(\to\)size in Resolution ({[}BEN01{]});
  order-robust OBDD width via expander gadgets
\item
  \textbf{SCL lens}: R on a cut corresponds to independent parities; if
  q is small, width/Alt must be \(\geq\) 2\^{}(\(\Omega\)(\(\lambda\)))
\end{itemize}

\textbf{Hamiltonian Cycle / k-Coloring on High Treewidth Graphs}

\begin{itemize}
\tightlist
\item
  \textbf{Upfront}: Graph and constraints are explicit
\item
  \textbf{Gated action}: Boundary assignments in a decomposition act
  like seeds; DP must keep 2\^{}(\(\Omega\)(treewidth)) distinct keys
\item
  \textbf{Known bounds}: DP lower bounds parameterized by
  treewidth/pathwidth; OBDD/branching-program width bounds tied to
  decomposition width
\item
  \textbf{SCL lens}: R accumulates across a high-width cut; unless q
  resolves those bits, Alt (keys/width) must grow exponentially
\end{itemize}

\textbf{Pointer-Chasing / Indexing CNFs (Query/Communication Analogues)}

\begin{itemize}
\tightlist
\item
  \textbf{Upfront}: Pointers/table entries are explicit
\item
  \textbf{Gated action}: Must follow the dependency chain; later
  positions are locked behind earlier reveals
\item
  \textbf{Known bounds}: Streaming/communication/branching-program lower
  bounds
\item
  \textbf{SCL lens}: Each hop adds R; without early q, Alt multiplies
  layerwise
\end{itemize}

\textbf{CSPs on Expanders (e.g., Latin/Sudoku-style)}

\begin{itemize}
\tightlist
\item
  \textbf{Upfront}: All constraints published; unit propagation emerges
  only after dependencies are met
\item
  \textbf{Known bounds}: Strong for specific proof systems; general TM
  time lower bounds are open
\item
  \textbf{SCL lens}: Same resolve-or-maintain tradeoff;
  L*\textquotesingle{}s overlay makes it verifiable and composable
  across cuts
\end{itemize}

\textbf{What L* Adds Beyond Natural Problems (Two Key Mechanisms):}

The natural problem examples above demonstrate SCL-like phenomena -
pebbling games force frontier tracking, expander-TSP requires wide DP
tables, etc. - but these remain \emph{paradigm-specific} bounds (width
for OBDDs, space for DP, proof size for resolution). We cannot
rigorously translate them into unconditional Turing machine time lower
bounds. L* bridges this gap via two engineered mechanisms:

\begin{enumerate}
\def\labelenumi{\arabic{enumi}.}
\item
  \textbf{Disjoint-atoms + Enc-recoverability (Cartesian factoring):}
  The disjoint address pools \{U\_v\} combined with parseable seed
  encoding Enc guarantee that artifact choices at different cut nodes
  factor independently in the worst case. Formally: for any cut C, the
  number of simultaneously distinguishable configurations is at least
  \(\Pi\)\_\{v\(\in\)C\} 2\^{}(R\_v-q\_v) (Lemma J.1-Cart, Appendix J).
  This \emph{compositional} property ensures that the min-cut bottleneck
  cannot be circumvented by clever caching or dynamic programming - an
  algorithm tracking \textless{} \(\Pi\) 2\^{}(R\_v-q\_v) states must
  miss some seed-tuple, forcing correctness errors. Natural problems
  lack this guarantee: symmetries and overlapping constraints can allow
  subexponential state counts even at high-width cuts.
\item
  \textbf{Frontier-Gate mechanism (FG, per-segment work):} FG wires
  parity-digest requirements (GateDigest\_v) into seeds at GREQ\_v=1
  nodes, forcing \(\Omega\)(n/W\_min) work per distinguishable artifact
  segment (Appendix C.1.1). This work \emph{survives caching} - even if
  an algorithm memoizes intermediate results, it must still compute gate
  digests for each distinct seed-path. Combined with the Cartesian
  factoring above, this elevates the information-theoretic bottleneck
  (2\^{}\(\lambda\) artifacts required) into a time lower bound:
  \(\geq\) 2\^{}\(\lambda\) segments × \(\Omega\)(n/W\_min) per segment
  = super-polynomial time. Natural problems have computational gates,
  but they\textquotesingle{}re not wired into a composable
  seed-dependency structure that enforces per-artifact work.
\end{enumerate}

\textbf{The Result:} These two mechanisms - disjoint-atoms
(compositional artifact-counting) + FG (per-artifact work) - transform
paradigm-specific observations (OBDD width, DP space, resolution size)
into unconditional Turing machine time lower bounds for L*. This is what
enables the Structural OWF construction (§9) and P \(\neq\) NP proof
(§10). Natural problems exhibit the same information-theoretic tradeoffs
empirically, but lack the mathematical scaffolding to prove them
rigorously across all computational models.

\paragraph{12.6 Why Our Approach Avoids Known
Barriers}\label{126-why-our-approach-avoids-known-barriers}

\textbf{The Big Picture}

Three barrier theorems have blocked P vs NP proofs for decades. Each
proves that a \emph{specific category} of proof techniques cannot work.
The key insight is simple: \textbf{barriers block techniques, not
results}. Our proof uses techniques outside all three categories.

\textbf{What each barrier blocks vs. what we use:}

\begin{itemize}
\tightlist
\item
  \textbf{Relativization (BGS 1975):} Blocks black-box arguments that
  ignore algorithm internals \(\to\) We use white-box analysis of
  L*\textquotesingle{}s specific structure
\item
  \textbf{Natural Proofs (RR 1997):} Blocks properties applying to many
  functions \(\to\) We use sparse, instance-specific properties
\item
  \textbf{Algebrization (AW 2009):} Blocks algebraic/polynomial methods
  \(\to\) We use discrete combinatorial counting
\end{itemize}

We\textquotesingle{}re not "cleverly avoiding"
barriers---we\textquotesingle{}re in fundamentally different
mathematical territory. (Note: barrier avoidance is evidence of viable
approach, not an additional theorem; correctness comes from §7-10.)

\begin{center}\rule{0.5\linewidth}{0.5pt}\end{center}

\textbf{Why This Matters}

For over three decades, the complexity theory community has identified
fundamental obstacles---called \textbf{barrier theorems}---that
rigorously prove certain broad classes of proof techniques CANNOT
resolve P vs NP. Natural Proofs (Razborov \& Rudich, 1997) showed that
techniques relying on "natural" properties of Boolean functions cannot
prove super-polynomial circuit lower bounds without breaking
cryptography. Relativization (Baker-Gill-Solovay, 1975) demonstrated
that any proof technique working uniformly for all oracles cannot
separate P from NP. Algebrization (Aaronson-Wigderson, 2009) extended
this to algebraic techniques. These barriers didn\textquotesingle{}t
just suggest difficulty---they mathematically proved whole categories of
approaches are doomed. Numerous promising lower bound programs hit these
walls and stalled.

\textbf{The Central Question:} Given these barriers have blocked decades
of attempts, how does our proof escape them?

\textbf{The Answer:} Our proof operates in a fundamentally different
technical regime that the barrier theorems don\textquotesingle{}t cover.
We\textquotesingle{}re not "cleverly avoiding" known
obstacles---we\textquotesingle{}re using a proof structure
(instance-specific semantic obligations) that lies outside the
barriers\textquotesingle{} scope by design.

\begin{center}\rule{0.5\linewidth}{0.5pt}\end{center}

\subsubsection{1. Relativization Barrier (Baker-Gill-Solovay
1975)}\label{1-relativization-barrier-baker-gill-solovay-1975}

\textbf{Simple Intuition}

Relativization blocks proofs that treat algorithms as "black boxes"
(input \(\to\) output, ignoring internals). Such proofs work the same
whether or not algorithms can query an oracle. But BGS showed: for some
oracles A, P\^{}A = NP\^{}A. So any oracle-independent proof technique
is doomed.

\textbf{Why We Escape (Intuitive)}

Our proof looks \emph{inside} L*\textquotesingle{}s structure:

\begin{itemize}
\tightlist
\item
  Seed chains create sequential dependencies (must compute parent seeds
  before child seeds)
\item
  Hermeticity (A1) controls exactly where information can enter
\item
  Designated addressing forces specific memory access patterns
\end{itemize}

An oracle would let you "skip" these dependencies---query
"what\textquotesingle{}s the answer?" and get it free, bypassing all the
work our lower bound charges for. Our lower-bound accounting
(Hermeticity/RWA/seed-locking) assumes no oracles; we make no claim
about P\^{}O vs NP\^{}O.

\textbf{Technical Details}

\emph{What it is:} BGS showed that any proof technique that
relativizes---meaning it works identically whether or not algorithms
have access to an arbitrary oracle O---cannot separate P from NP. They
constructed oracles where P\^{}A = NP\^{}A (diagonalization oracle) and
P\^{}B \(\neq\) NP\^{}B (random oracle with PSPACE-complete problem),
proving oracle-independent techniques are fundamentally limited. This
killed approaches based on diagonalization, simulation, and other
"oracle-robust" methods.

\emph{Example killed:} Classical diagonalization arguments (used to
separate DTIME hierarchies) fail for P vs NP because they
relativize---the same diagonalization works with oracles, but P\^{}A =
NP\^{}A for some A.

\emph{Why we escape (technical):} Our proof is fundamentally
\textbf{non-relativizing}---it depends critically on
L*\textquotesingle{}s internal structure: seed-dependency chains
(Seed\_v = Enc(v \textbar{}\textbar{} sorted\{parents\}
\textbar{}\textbar{} GateDigest\_v)), disjoint designated address pools
\{U\_v\}, and Hermeticity (A1). An oracle would provide "free answers"
bypassing the seed-bound addressing and designated read accounting
(RWA), destroying the Cartesian factoring (Lemma J.1-Cart) and
conservation law. Specifically: if algorithms could query an oracle
about seed values or gate digests without paying the
\(\Omega\)(n/W\_min) work per segment (FG mechanism), the lower bound
would collapse. This doesn\textquotesingle{}t mean our proof "fails with
oracles"---it means our lower-bound accounting (Hermeticity A1, RWA,
seed-locking) assumes no oracles. We make no claim about P\^{}O vs
NP\^{}O; our result is for the standard uniform model. The proof
technique inherently relies on explicit instance structure that
doesn\textquotesingle{}t relativize, placing it outside
BGS\textquotesingle{}s scope.

\textbf{Addressing a Potential Confusion}

\emph{"By restricting to uniform algorithms (no oracles),
aren\textquotesingle{}t you just avoiding the barrier?"}

This misunderstands BGS\textquotesingle{}s scope. The barrier applies to
\textbf{proof techniques that relativize}---arguments that work
uniformly across all oracle models. BGS showed such techniques fail
because \(\exists\) oracles where P\^{}A = NP\^{}A.

Our proof doesn\textquotesingle{}t claim to relativize. We prove P
\(\neq\) NP for \textbf{uniform PPT} (the standard model, Abstract line
9), not for all oracle variants. This is a valid model choice---the same
model used in the Millennium Prize statement---not "barrier avoidance."
The technical reason we don\textquotesingle{}t extend to P\^{}O vs
NP\^{}O is that oracles would break L*\textquotesingle{}s information
accounting, requiring entirely different constructions. \textbf{BGS
concerns proofs that work for all oracle models; we prove a result for
the specific uniform model.}

\begin{center}\rule{0.5\linewidth}{0.5pt}\end{center}

\subsubsection{2. Natural Proofs Barrier (Razborov-Rudich
1997)}\label{2-natural-proofs-barrier-razborov-rudich-1997}

\textbf{Simple Intuition}

Natural Proofs blocks techniques that identify hard functions by
properties that are (a) \textbf{large}---apply to many functions
(\(\geq\)1/2\^{}O(n) fraction), and (b)
\textbf{constructive}---efficiently testable from the
function\textquotesingle{}s truth table. Such properties would break
cryptography: you could test if a function is a pseudorandom function by
checking if it has the "hardness property."

\textbf{Why We Escape (Intuitive)}

Our hard instances are \textbf{astronomically rare}:

\begin{itemize}
\tightlist
\item
  Only Plant(\(\varphi\), r) outputs are hard
\item
  Density \(\leq\) 2\^{}(-\(\Omega\)(\(\lambda\)))---exponentially
  sparse
\item
  Violates the "largeness" requirement completely
\end{itemize}

Also, you \textbf{can\textquotesingle{}t test} if something has our
hardness property without knowing the secret planting parameters
(\(\varphi\), r, the FG configuration). The hardness property
isn\textquotesingle{}t "efficiently recognizable from truth tables"---it
requires the planted metadata.

We don\textquotesingle{}t say "most random-looking functions are hard"
(which Natural Proofs forbids). We say "these \emph{specific planted
instances} are hard"---totally different.

\textbf{Technical Details}

\emph{What it is:} The Natural Proofs barrier showed that any proof
technique with two properties---(1) \textbf{constructivity} (can
efficiently recognize hard functions) and (2) \textbf{largeness}
(applies to a large fraction of functions, density \(\geq\)
1/poly(2\^{}n))---cannot prove circuit lower bounds for NP without
breaking standard cryptographic assumptions. This killed approaches
based on "natural" circuit complexity measures (like those used for
AC\^{}0, monotone circuits, etc.), which historically relied on
efficiently computable properties that applied broadly.

\emph{Example killed:} Extending Razborov\textquotesingle{}s AC\^{}0
lower bounds (via switching lemma, approximation method) to stronger
circuit classes---these techniques are "natural" and would contradict
pseudorandom functions if they worked for P vs NP.

\emph{Why we escape (technical):} L*\textquotesingle{}s hardness is
\textbf{instance-specific}, not function-wide. Our explicit generator
Plant(\(\varphi\), r) with FG wiring produces hard instances. The
relevant density measure: among all possible outputs of
Plant(\(\varphi\), r) for varying (\(\varphi\), r), the fraction with
min-cut residual \(\geq\) \(\lambda\)\_base is \(\leq\)
2\^{}(-\(\Omega\)(\(\lambda\)\_base))---exponentially sparse over the
planted family (not over all Boolean functions), violating largeness. We
don\textquotesingle{}t claim "most Boolean functions with property X are
hard" (which Natural Proofs forbids); we show an explicit, efficiently
samplable sparse family is hard (see §9.2 for f and §8.1 for
per-instance bounds). The property we exploit (published overlay with
min-cut residual \(\geq\) \(\lambda\)\_base, gated digests wired into
seeds) requires reading explicit instance
metadata---it\textquotesingle{}s not an efficiently computable predicate
of truth tables. Natural Proofs doesn\textquotesingle{}t apply to
sparse, explicit instance families with structural hardness.

\textbf{Addressing a Potential Confusion}

\emph{"L* \(\in\) NP means we can efficiently verify membership (§10.2
verifier V checks witnesses in polynomial time).
Doesn\textquotesingle{}t this mean we can efficiently recognize
instances with the hardness property?"}

This conflates two distinct questions:

\begin{itemize}
\item
  \textbf{Membership testing} asks: given instance x and witness W, can
  we verify x \(\in\) L* efficiently? Answer: yes---L* \(\in\) NP by
  construction.
\item
  \textbf{Hardness property recognition} asks: given a Boolean function
  f (as truth table or circuit), can we efficiently determine whether f
  exhibits the structural hardness we exploit?
\end{itemize}

These differ fundamentally. Natural Proofs concerns the
latter---efficiently recognizing which \textbf{functions} (across all
computational problems) possess a hardness-inducing property. Our
hardness property requires reading the \textbf{planted instance
structure}---specifically, the overlay metadata (\(\varphi\), r) from
Plant(\(\varphi\), r), FG gate configuration, and DAG\textquotesingle{}s
min-cut residual \(\lambda\)\_base. This metadata is part of
L*\textquotesingle{}s explicit construction (§6), not extractable from a
function\textquotesingle{}s truth table or standard circuit encoding. We
prove specific planted instances (from our generator) are hard, not that
"all functions with efficiently checkable property P are hard."

\textbf{Natural Proofs targets function-wide recognizable properties;
our approach uses instance-specific planted structure.} The barrier
doesn\textquotesingle{}t apply to sparse, explicitly constructed
families where hardness depends on engineered metadata unavailable in
generic function representations.

\begin{center}\rule{0.5\linewidth}{0.5pt}\end{center}

\subsubsection{3. Algebrization Barrier (Aaronson-Wigderson
2009)}\label{3-algebrization-barrier-aaronson-wigderson-2009}

\textbf{Simple Intuition}

Algebrization blocks proofs that survive when algorithms can access not
just an oracle A, but also its "low-degree polynomial extension" Ã over
a finite field. Many algebraic techniques (polynomial method, rank
arguments) have this property---and AW showed they cannot separate P
from NP.

\textbf{Why We Escape (Intuitive)}

We count \textbf{discrete objects}, not algebraic degree:

\begin{itemize}
\tightlist
\item
  "How many distinguishable configurations exist?" \(\to\)
  2\^{}\(\lambda\) (an integer)
\item
  Pigeonhole: k configurations need \(\geq\)k states
\item
  Cardinality constraints: exactly 1 feasible world (not "approximately
  1")
\end{itemize}

Algebra deals in continuous quantities (polynomial degree, field
elements). Our proof deals in discrete quantities (cardinality, exact
Boolean equality). There\textquotesingle{}s no meaningful "low-degree
extension" of "number of elements in a set."

Also, our key constraints are \textbf{exact}:

\begin{itemize}
\tightlist
\item
  Parity must equal digest \emph{exactly} (0 or 1, not 0.7)
\item
  Exactly 1 accepting world (not 1.5)
\item
  Variables are Boolean (true/false, not maybe)
\end{itemize}

Algebraic extensions would blur these into approximate/fractional
constraints, breaking everything.

\textbf{Technical Details}

\emph{What it is:} Algebrization extended relativization to algebraic
proof techniques. It showed that proofs which "algebrize"---meaning they
work not just with oracles but with low-degree extensions of oracles
over finite fields---cannot separate P from NP. This killed hopes that
algebraic methods (used successfully in communication complexity,
polynomial method, etc.) could bypass relativization via low-degree
approximations.

\emph{Example killed:} Algebraic circuit lower bounds via polynomial
method, rank arguments, communication complexity lifts---techniques that
represent functions as low-degree polynomials and bound their algebraic
complexity.

\emph{Why we escape (technical):} Our proof counts \textbf{combinatorial
distinguishability} (2\^{}\(\lambda\) seed-consistent worlds at the
min-cut), not algebraic degree. The artifact-counting ledger---tracking
which seed-histories an algorithm can distinguish via designated address
reads under RWA, enforced by Keyedness (different seeds \(\to\)
different addresses \(\to\) cannot merge without errors)---is
fundamentally discrete and non-algebraic (see Proposition N.3.a in
Appendix N for the formal statement). There\textquotesingle{}s no
natural low-degree polynomial representation that preserves the critical
properties: (1) Cartesian factoring across cuts (Lemma J.1-Cart) relies
on disjoint address pools (combinatorial partitioning), not field
structure; (2) FG FG gates compute XOR over designated address contents,
which is degree-1 over GF(2) but the \emph{addressing} (F\_overlay:
Seed\_v × (j,\(\ell\)) \(\to\) addresses) involves hash-like functions
(assumption: F\_overlay lacks low-degree extensions over larger fields);
(3) the min-cut residual \(\lambda\) measures bits of information
(log\(_2\) of distinguishable states), not algebraic degree---collapsing
states algebraically doesn\textquotesingle{}t preserve
distinguishability correctness. Standard algebraic liftings fail to
simulate our counting argument.

\textbf{Addressing a Potential Confusion}

\emph{"L* membership testing is deterministic (§10.2 verifier V runs in
polynomial time given a witness), so can\textquotesingle{}t L* be
represented as polynomials? Wouldn\textquotesingle{}t that make the
proof algebrize?"}

This confuses two separate properties:

\begin{itemize}
\item
  \textbf{The computation} (checking whether x \(\in\) L* given a
  witness) is indeed deterministic and polynomial-representable---L*
  \(\in\) NP via the standard verifier.
\item
  \textbf{The proof technique}, however, relies on discrete structural
  properties of \emph{planted instances} that break under algebraic
  oracle extensions.
\end{itemize}

Specifically:

(i) WellFormedRandomness (§9.2, PlantedInstanceConsistency) requires
digest bits to \textbf{exactly} equal parity of emergent
configurations---a Boolean constraint (bit = 0 or 1) that cannot be
"partially satisfied." Under low-degree extensions over fields, one
could have "fractional parities" (e.g., polynomial evaluations \(\in\)
{[}0,1{]}), violating the exact-match requirement.

(ii) Planted instance uniqueness
(\texttt{planted\_instances\_have\_uniqueness} in
AcceptanceUniqueness.lean) proves exactly \textbf{card = 1} feasible
world at acceptance---a discrete cardinality constraint. Algebraic
extensions could yield "fractional witnesses" (superpositions
representable as polynomials), breaking the discreteness.

(iii) The parity function itself, while linear over GF(2), becomes
degree-2\^{}n as a polynomial over larger fields when composed with the
addressing function F\_overlay---far from "low-degree."

\textbf{The proof technique hinges on discrete, exact-equality
constraints} (digest = parity, card = 1, Boolean satisfaction) that
algebrization would blur into approximate/fractional constraints.
\textbf{Algebrization concerns whether the proof survives algebraic
oracle extensions, not whether the original computation is
deterministic.} Our proof doesn\textquotesingle{}t algebrize because its
correctness requirements are fundamentally Boolean/discrete, even though
L* membership (with witness) is polynomial-time decidable.

\begin{center}\rule{0.5\linewidth}{0.5pt}\end{center}

\subsubsection{Summary: Why Barrier-Avoidance Validates the
Result}\label{summary-why-barrier-avoidance-validates-the-result}

\textbf{How we escape each barrier:}

\begin{itemize}
\tightlist
\item
  \textbf{Relativization:} Blocks oracle-independent arguments \(\to\)
  We use structure-dependent analysis; our accounting assumes no oracles
\item
  \textbf{Natural Proofs:} Blocks large, constructive properties \(\to\)
  Our density \(\leq\) 2\^{}(-\(\Omega\)(\(\lambda\))) over planted
  family; instance-specific hardness
\item
  \textbf{Algebrization:} Blocks algebraic degree bounds \(\to\) We
  count discrete objects with exact Boolean constraints
\end{itemize}

The three barriers don\textquotesingle{}t represent "things to
avoid"---they represent fundamental limitations of broad proof
strategies. Our approach escapes them not by cleverness but by operating
in a different technical space:

\begin{itemize}
\item
  \textbf{Instance-specificity:} We don\textquotesingle{}t prove "all
  NP-complete problems are hard" (which would hit Natural Proofs). We
  construct one explicit NP-complete problem (L*) with engineered
  structure (A1-A5) that mathematically forces super-polynomial
  witness-finding.
\item
  \textbf{Non-relativizing structure:} We don\textquotesingle{}t use
  oracle-independent diagonalization. We rely on L*\textquotesingle{}s
  explicit dependency chains and designated addressing---structure that
  doesn\textquotesingle{}t exist in relativized worlds.
\item
  \textbf{Combinatorial counting:} We don\textquotesingle{}t bound
  algebraic degree. We count discrete distinguishable artifacts forced
  by Injectivity + Keyedness---a combinatorial accounting that
  doesn\textquotesingle{}t lift to algebraic settings.
\end{itemize}

\textbf{Implication for P \(\neq\) NP:} The fact that our proof lies
outside all three known barriers doesn\textquotesingle{}t guarantee
it\textquotesingle{}s correct---mathematical correctness comes from the
formal proofs in §7-10. Rather, barrier-avoidance explains why this
approach isn\textquotesingle{}t immediately doomed by known
impossibility results. The barriers killed previous programs by proving
"no technique in category X can work." Our technique
isn\textquotesingle{}t in any of those categories. This is evidence (not
proof) that the approach occupies unexplored technical territory where
separation results remain mathematically possible.

\textbf{Intellectual honesty (evidence vs meta-theorem):} The analysis
above offers evidence---not a meta-theorem---that our method operates
outside classic barriers. We cannot prove that no future barrier theorem
will ever subsume instance-specific semantic arguments. But the three
major known barriers (Natural Proofs, Relativization, Algebrization)
don\textquotesingle{}t apply to our proof structure. Our contribution is
the mathematical separation (FP \(\neq\) FNP for L* on classical uniform
PPT, §9-10); this barrier discussion contextualizes why standard
impossibility results don\textquotesingle{}t immediately invalidate the
approach.

\begin{center}\rule{0.5\linewidth}{0.5pt}\end{center}

\subsubsection{Formal Barrier Statements (Proofs in Appendix
N)}\label{formal-barrier-statements-proofs-in-appendix-n}

\textbf{Statement N-R (Non-relativizing).} There exists an oracle O such
that the cut-product factoring and FG gating used in our proofs no
longer characterize the solver\textquotesingle{}s obligations (the
oracle answers can violate Hermeticity/seed-bound addressing), while
P\^{}O = NP\^{}O on suitable promise variants. Hence our proof technique
relies on non-relativizing semantic structure of explicit instances
rather than oracle-robust arguments. We prove P \(\neq\) NP for uniform
PPT (standard model), not for all oracle variants P\^{}O vs NP\^{}O.

\emph{Intuition: Oracles bypass L*\textquotesingle{}s designated
addressing, breaking the problem definition. BGS applies to proofs that
work for ALL oracle models; we prove a result for the specific uniform
model---a valid scope choice, not barrier avoidance.}

\textbf{Statement N-Nat (Non-natural).} The property exploited by the
lower bound (existence of a published overlay with min-cut residual
\(\geq\) \(\lambda\)\_base and gated digests) is not large: the density
of functions exhibiting it is \(\leq\)
2\^{}(-\(\Omega\)(\(\lambda\)\_base)) over the relevant family, and the
recognition procedure depends on explicit overlay metadata (not a small,
efficiently computable predicate of truth tables). Therefore the method
evades the Natural Proofs barrier.

\emph{Intuition: Hard instances are exponentially sparse and require
reading planted structure metadata unavailable in generic function
representations. Natural Proofs concerns function-wide hardness
recognition; we prove instance-specific hardness from engineered
structure. Membership testing (L* \(\in\) NP) \(\neq\) hardness property
recognition.}

\textbf{Statement N-Alg (Non-algebrizing).} No bounded-degree low-field
extension captures the artifact-counting ledger (seed-bound addressing +
RWA + Keyedness) so as to simulate the cut-product lower bound; in
particular, the ledger counts combinatorial distinguishability across
dynamic addresses rather than algebraic degree, and standard liftings
fail to preserve these counters. The proof relies on discrete,
exact-equality constraints of planted instances (WellFormedRandomness
with digest = parity, card = 1 uniqueness) that break under algebraic
extensions allowing fractional/approximate values.

\emph{Intuition: We count discrete distinguishable artifacts and require
exact Boolean constraints, not algebraic degree---the proof technique
doesn\textquotesingle{}t algebrize even though L* membership is
deterministic and polynomial-representable.}

\paragraph{12.7 The Universal Information Skeleton
(expository)}\label{127-the-universal-information-skeleton-expository}

\textbf{Why This Abstraction Matters:} Section §5 showed that the
Semantic Conservation Law q + \(\Phi\) \(\geq\) R manifests across
computational paradigms - as tree size (backtracking), table keys (DP),
OBDD width, resolution clauses, TM time. But these appear to be separate
phenomena analyzed through distinct techniques (tree counting, dynamic
programming analysis, proof complexity, etc.). This section reveals a
deeper unity: all these bounds are instances of a single
\textbf{universal information-accounting principle} that transcends
specific computational models. Understanding this abstraction clarifies
why the SCL isn\textquotesingle{}t a model-specific trick -
it\textquotesingle{}s a formalization of fundamental
information-theoretic constraints that appear throughout computer
science.

\textbf{Meta-inequality (informal, unifying view).} For any model of
computation and any point in its execution, the information that has
been \textbf{learned} plus the information implicit in the
\textbf{distinguishable artifacts maintained} must cover the
\textbf{information required} to separate the remaining cases:

\textbf{Learned} + log\(_2\)(\textbf{Distinct States}) \(\geq\)
\textbf{Required}

READ-OR-X makes this accounting explicit and composable across cuts of
the instance DAG.

\textbf{The Three Fundamental Dimensions:} Any computational strategy
for solving a problem must navigate a three-dimensional tradeoff space:

\begin{enumerate}
\def\labelenumi{\arabic{enumi}.}
\item
  \textbf{Resolution (learn information):} Gain new bits through
  computation, reads, communication, propagation. Increases the
  "Learned" term. Example: unit propagation in SAT solvers, comparison
  results in sorting.
\item
  \textbf{Storage (maintain possibilities):} Track multiple unresolved
  hypotheses simultaneously. The log\(_2\)(\#states) term measures how
  many distinct possibilities remain alive. Example: DP table keys, OBDD
  width, search tree branches.
\item
  \textbf{Elimination (prune via reasoning):} Reduce "Required" by
  exploiting structure - symmetries, dominance, pruning rules that show
  certain cases need not be distinguished. Example: alpha-beta pruning,
  symmetry breaking, learned clauses.
\end{enumerate}

Every algorithm balances these three: resolve more (increase q),
maintain more states (increase \(\Phi\)), or reduce requirements
(decrease R). The conservation law q + \(\Phi\) \(\geq\) R makes this
tradeoff mathematically explicit.

(This subsection is \textbf{expository}; it does not introduce new
theorems.)

\textbf{How the Skeleton Manifests Across Fields:}

Below we show how the universal inequality \textbf{Learned +
log\(_2\)(Distinct States) \(\geq\) Required} appears in classical
computational models. Each row identifies what counts as "learned
information," "maintained states," and "required distinctions" in that
model\textquotesingle{}s native terms:

\begin{itemize}
\item
  \textbf{Decision trees}

  \begin{itemize}
  \tightlist
  \item
    Learn (increase q): 1 bit/comparison
  \item
    Maintain (\(\Phi\)): Branches in tree
  \item
    Required (R): log\(_2\)(\#outputs)
  \item
    Inequality: \#comps \(\geq\) log\(_2\)(n!)
  \end{itemize}
\item
  \textbf{Communication}

  \begin{itemize}
  \tightlist
  \item
    Learn (increase q): Bits sent
  \item
    Maintain (\(\Phi\)): Rectangles in partition
  \item
    Required (R): Input entropy
  \item
    Inequality: comm \(\geq\) log\(_2\)(\#rectangles)
  \end{itemize}
\item
  \textbf{Branching programs}

  \begin{itemize}
  \tightlist
  \item
    Learn (increase q): Queries read
  \item
    Maintain (\(\Phi\)): Width per layer
  \item
    Required (R): Function complexity
  \item
    Inequality: L·log\(_2\)(W) \ensuremath{\gtrsim} required
  \end{itemize}
\item
  \textbf{Resolution}

  \begin{itemize}
  \tightlist
  \item
    Learn (increase q): Unit-prop/derivation
  \item
    Maintain (\(\Phi\)): Clause width
  \item
    Required (R): CNF difficulty
  \item
    Inequality: Small width \(\Rightarrow\) large size
  \end{itemize}
\item
  \textbf{READ-OR-X (L*)}

  \begin{itemize}
  \tightlist
  \item
    Learn (increase q): q\_v (first-use reads)
  \item
    Maintain (\(\Phi\)): Alt(C) (artifacts)
  \item
    Required (R): \(\Sigma\) R\_v (emergence)
  \item
    Inequality: Q(C) + log\(_2\) Alt(C) \(\geq\) \(\Lambda\)(C)
  \end{itemize}
\end{itemize}

\textbf{Decision trees (sorting).} Outputs are n! orders; each
comparison yields 1 bit. Thus \#comps \(\geq\) log\(_2\)(n!) \(\approx\)
n log n. Our framework mirrors this: either \textbf{learn} a bit per
comparison (Resolution dimension: raise q) or \textbf{test} candidates
(Elimination dimension) or \textbf{maintain branches} (Storage
dimension: double distinguishable artifacts).

\textbf{Communication (e.g., disjointness). A k-bit protocol partitions
the matrix into at most 2\^{}k rectangles; if distinguishing requires N
rectangles, then 2\^{}k \(\geq\) N \(\Rightarrow\) k \(\geq\)
log\(_2\)(N). In our framework: communicate} (Resolution: increase q) or
\textbf{maintain rectangles} (Storage: distinguishable artifacts).

\textbf{Branching programs. Width W encodes log\(_2\)(W) bits of
distinguishable state per layer; after L layers you separate at most
W\^{}L cases. If the task needs 2\^{}n distinctions, then L \(\geq\)
n/log\(_2\)(W). Our framework: store} width (Storage dimension:
distinguishable artifacts) or \textbf{resolve} more
(Resolution/Elimination dimensions).

\textbf{Resolution (width \(\to\) size).} Width w acts like log\(_2\) of
tracked partial assignments; if w is small, the proof must revisit many
worlds \(\Rightarrow\) large size. In our framework\textquotesingle{}s
terms, small Storage forces exponential cost in Elimination dimension
(many proof steps/artifacts).

\textbf{Why This Unity Holds (Three Information-Theoretic Principles):}

The universal skeleton isn\textquotesingle{}t a coincidence - it follows
from three fundamental principles that hold regardless of computational
model:

\begin{enumerate}
\def\labelenumi{\arabic{enumi}.}
\item
  \textbf{No information creation:} Algorithms cannot manufacture bits
  from nothing. All "learned" information comes from designated sources
  (reads, queries, messages, comparisons, propagations). This bounds the
  rate at which q can grow. For L*, RWA (Receiving-Window Attribution)
  makes this explicit: q\_v credits only first-use designated reads,
  preventing double-counting.
\item
  \textbf{No compression below entropy:} If 2\^{}R distinct cases must
  be separated and you\textquotesingle{}ve only learned q \textless{} R
  bits, the remaining 2\^{}(R-q) possibilities cannot be compressed into
  fewer than 2\^{}(R-q) distinguishable states without losing
  correctness. Formally: maintaining \textless{} 2\^{}(R-q) artifacts
  forces collision - two inequivalent cases map to the same state,
  causing errors. For L*, this is Lemma 7.I (Injectivity \(\to\) Alt\_v
  \(\geq\) 2\^{}(R\_v-q\_v)).
\item
  \textbf{Independent requirements multiply:} If nodes v\(_1\), v\(_2\)
  in a cut each require 2\^{}(R\(_1\)), 2\^{}(R\(_2\)) distinguishable
  states and their requirements are independent (no shared information
  source), the total across the cut is \(\geq\) 2\^{}(R\(_1\)) ×
  2\^{}(R\(_2\)) = 2\^{}(R\(_1\)+R\(_2\)). For L*, this is Cartesian
  factoring (Lemma J.1-Cart): disjoint address pools \{U\_v\} +
  Dependency (A5) + no cross-coupling ensure independence (formalized as
  H1-H5 worst-case properties), yielding Alt(C) = \(\Pi\) Alt\_v.
\end{enumerate}

Micro-example (multiplicativity). If C = \{v\(_1\), v\(_2\)\} with one
unresolved bit at each (s\_\{v\(_1\)\} = s\_\{v\(_2\)\} = 1), then
Alt\_\{v\(_1\)\} \(\geq\) 2 and Alt\_\{v\(_2\)\} \(\geq\) 2, so Alt(C)
\(\geq\) 2·2 = 4 and \(\Phi\)(C) = log\(_2\) Alt(C) \(\geq\) 2. This
matches \(\Sigma\)\_\{v\(\in\)C\} (R\_v \(-\) q\_v) = 2 exactly.

These three principles - formalized rigorously in §7 for L* via A1-A5
properties - explain why the conservation law appears across all
computational models. It\textquotesingle{}s not a model-specific
phenomenon; it\textquotesingle{}s a consequence of \textbf{fundamental
information-theoretic constraints} on distinguishing possibilities.

\textbf{Connection to P \(\neq\) NP:}

This universal skeleton clarifies what our P \(\neq\) NP proof actually
achieves. We\textquotesingle{}re not introducing a novel lower bound
technique specific to Turing machines - we\textquotesingle{}re
formalizing an information-accounting principle that already underlies
lower bounds across computational models. The innovation is showing
that:

\begin{enumerate}
\def\labelenumi{\arabic{enumi}.}
\item
  \textbf{L*\textquotesingle{}s structure (A1-A5) makes the accounting
  rigorous:} Natural problems exhibit the same tradeoffs empirically
  (§12.5), but their tangled structure prevents formal proofs.
  L*\textquotesingle{}s engineered properties (Hermeticity, Injectivity,
  Emergence, Closure, Dependency) make q\_v, R\_v, Alt\_v well-defined
  and provable.
\item
  \textbf{FG mechanism converts information bottleneck \(\to\) time:}
  The universal skeleton alone gives information-theoretic bounds
  (2\^{}\(\lambda\) artifacts required). FG (Frontier-Gate, §8) wires
  computational work (\(\Omega\)(n/W\_min) per segment) into the
  seed-dependency structure, translating the information bottleneck into
  unconditional Turing machine time lower bounds.
\item
  \textbf{Structural OWF construction uses universality:} Because the
  conservation law holds for \emph{any} computational strategy (not just
  specific algorithms), we can prove \(\forall\)x* \(\forall\)A
  (uniform) \(\to\) time \(\geq\) super-poly. This \(\forall\)x*
  \(\forall\)A quantifier structure is what enables the one-way function
  construction (§9) and P \(\neq\) NP via the classical bridge (§10).
\end{enumerate}

The universal skeleton isn\textquotesingle{}t separate from the main
proof - it\textquotesingle{}s the conceptual foundation explaining why
the approach works across all computational models.

\textbf{Caveats (reviewer guidance).} \textbullet{} This section is an
\textbf{interpretive lens}, not a replacement for the original
lower-bound theorems in each field (we do not restate constants). \textbullet{}
Communication mapping uses the deterministic rectangle bound; randomized
protocols fit via the standard distributional/Yao view. \textbullet{} The
BP/Resolution rows are heuristic paraphrases of the formal bridges (see
§5); our theorems rely only on the READ-OR-X adapters and cut law. \textbullet{}
Constants and precise formulations differ between fields; we show
conceptual unity, not identical quantitative bounds.

\paragraph{12.8 Why the Semantic Core Properties Appear
Everywhere}\label{128-why-the-semantic-core-properties-appear-everywhere}

\textbf{Why This Question Matters:} A natural concern: Are
L*\textquotesingle{}s structural properties (A1-A5) artificially
engineered constraints that exist nowhere else, making L* an isolated
curiosity unrelated to real computational problems? This section
addresses that concern by showing these properties reflect fundamental
patterns that appear throughout computer science - in version control
systems, distributed consensus, verified compilation, and more.
Understanding these connections clarifies that L* isn\textquotesingle{}t
contrived; it\textquotesingle{}s a mathematical purification of
constraints that already govern real systems.

Note on analogy limits. Real systems sometimes violate strict
Hermeticity (e.g., side channels, shared caches) or allow reuse across
equivalent histories; our A1-A5 are idealized semantic constraints used
to obtain rigorous theorems. The correspondences below should be read as
structural analogues, not exact formalizations.

\textbf{The Five Core Properties (A1-A5):}

Our framework rests on five instance-side axioms (A1-A5, §4.2, §6) used
in the Structural OWF construction (§9):

\begin{enumerate}
\def\labelenumi{\arabic{enumi}.}
\item
  \textbf{A1 (Hermeticity):} Information flows only through designated
  channels - no hidden dependencies or external shortcuts. All required
  data must be explicitly read from designated locations.
\item
  \textbf{A2 (Injectivity):} Different computational histories produce
  distinguishable artifacts. If two paths make different choices
  (resolve different parent seeds), they cannot merge into the same
  artifact without losing distinguishability.
\item
  \textbf{A3 (Emergence):} Required information doesn\textquotesingle{}t
  exist upfront - it emerges through computation. Specifically, exactly
  R\_v fresh bits emerge at each node v via rank(H\_v) = R\_v.
\item
  \textbf{A4 (Closure/Recoverability):} Past choices are recoverable
  from current state. The seed encoding Seed\_v = Enc(v
  \textbar{}\textbar{} sorted\{parents\} \textbar{}\textbar{}
  GateDigest\_v) is parseable - you can deterministically extract parent
  information from seeds.
\item
  \textbf{A5 (Dependency):} Computation must follow logical dependencies
  - cannot "leap ahead" via reasoning. Children depend on parents; the
  DAG structure enforces topological order.
\end{enumerate}

\textbf{Where These Properties Appear (Real Systems):}

Below we show three well-known systems that exhibit A1-A5 (or close
analogues), demonstrating these aren\textquotesingle{}t artificial
constraints invented for L*.

\textbf{Example 1: Git Version Control}

Git\textquotesingle{}s commit DAG exhibits all five properties (achieved
via cryptographic hashing, whereas L* uses pure combinatorial
structure):

\begin{itemize}
\item
  \textbf{A3 (Emergence):} Commit hash SHA-1(content
  \textbar{}\textbar{} parent\_hashes \textbar{}\textbar{} metadata)
  doesn\textquotesingle{}t exist until computed. You cannot know a
  commit\textquotesingle{}s hash before creating it - the hash emerges
  from the commit operation.
\item
  \textbf{A2 (Injectivity):} Different commit contents \(\to\) different
  hashes (SHA-1 collision resistance). Two commits with different
  parents or different file trees cannot produce the same hash, ensuring
  distinguishability.
\item
  \textbf{A5 (Dependency):} Child commits depend on parent commits - you
  cannot create commit C referencing parent P before P exists. The DAG
  enforces topological ordering.
\item
  \textbf{A1 (Hermeticity):} A commit\textquotesingle{}s hash depends
  only on its content, parent hashes, and metadata - no external
  information channels. Given the commit object, you can verify the hash
  deterministically.
\item
  \textbf{A4 (Closure):} The commit hash is content-addressed - from the
  hash, you can retrieve the commit object, which contains parent
  hashes, allowing you to traverse the entire history. Past choices
  (parent commits) are recoverable.
\end{itemize}

\textbf{Example 2: Blockchain (Bitcoin-style Proof-of-Work)}

Bitcoin\textquotesingle{}s blockchain demonstrates similar patterns:

\begin{itemize}
\item
  \textbf{A3 (Emergence):} Block hash emerges from mining - solving
  Hash(block\_header \textbar{}\textbar{} nonce) \textless{} target. The
  valid hash doesn\textquotesingle{}t exist until computational work
  finds it.
\item
  \textbf{A2 (Injectivity):} Different blocks \(\to\) different hashes
  (collision resistance). Changing any transaction, previous block
  reference, or nonce changes the hash.
\item
  \textbf{A5 (Dependency):} Each block includes the previous
  block\textquotesingle{}s hash in its header, creating a dependency
  chain. You cannot create block n+1 without block n\textquotesingle{}s
  hash.
\item
  \textbf{A1 (Hermeticity):} Block validity depends only on: (1)
  cryptographic proof (hash meets difficulty target), (2) previous block
  reference, (3) transaction validity. No external "shortcuts" can
  create a valid block without doing the work.
\item
  \textbf{A4 (Closure):} Each block header contains the previous hash,
  allowing full chain verification. The entire blockchain history is
  recoverable by traversing backward from any block.
\end{itemize}

\textbf{Example 3: Verified Compilation (CompCert-style)}

Formally verified compilers exhibit weaker versions of these properties:

\begin{itemize}
\item
  \textbf{A3 (Emergence):} Assembly code emerges from compilation - it
  doesn\textquotesingle{}t exist before the transformation. Intermediate
  representations (IR) emerge from parsing and lowering passes.
\item
  \textbf{A2 (Injectivity - weaker):} Different source programs
  generally produce different IR (modulo optimizations). Verified
  compilation ensures semantic preservation, a form of
  distinguishability.
\item
  \textbf{A5 (Dependency):} Compilation phases must follow order:
  parsing \(\to\) type checking \(\to\) IR generation \(\to\)
  optimization \(\to\) code generation. Cannot skip type checking and
  get correct code.
\item
  \textbf{A1 (Hermeticity - weaker):} In pure functional compilation,
  output depends only on input source and compiler flags - no hidden
  state. Verified compilers formalize this via purity guarantees.
\item
  \textbf{A4 (Closure - partial):} Some compilers maintain provenance
  metadata allowing source reconstruction or debugging (though not
  always full recoverability like Git).
\end{itemize}

\textbf{Key Insight: Cryptography vs Structure}

Git and Bitcoin achieve A1-A5 through \textbf{cryptographic primitives}
(SHA-1, SHA-256). L* achieves analogous properties through \textbf{pure
combinatorial structure}:

\begin{itemize}
\tightlist
\item
  \textbf{Emergence:} Via rank-forcing (rank(H\_v) = R\_v), not
  cryptographic emergence
\item
  \textbf{Injectivity:} Via injective encoding Enc(v
  \textbar{}\textbar{} sorted\{parents\}), not hash collision resistance
\item
  \textbf{Dependency:} Via explicit DAG edges, not cryptographic
  chaining
\item
  \textbf{Hermeticity:} Via disjoint address pools \{U\_v\}, not
  cryptographic isolation
\item
  \textbf{Closure:} Via parseable encoding, not cryptographic
  content-addressing
\end{itemize}

This demonstrates that \textbf{information-theoretic structure alone}
(without cryptographic assumptions) can enforce conservation-law
obligations. L* is a "proof-of-work-less blockchain" - the dependency
DAG compels computational effort through mathematical necessity rather
than cryptographic hardness.

\textbf{Connection to P \(\neq\) NP:}

The fact that A1-A5 appear in real systems (Git, blockchain,
compilation) validates that these aren\textquotesingle{}t contrived
mathematical abstractions - they capture fundamental patterns in systems
that manage dependencies, ensure correctness, and prevent shortcuts.
L*\textquotesingle{}s innovation is showing that these structural
properties, when combined with FG (Frontier-Gate mechanism wiring
computational work into dependencies), mathematically force
super-polynomial witness-finding. The ubiquity of A1-A5 in real systems
suggests the conservation law q + \(\Phi\) \(\geq\) R
isn\textquotesingle{}t a special trick - it\textquotesingle{}s a
formalization of constraints that already govern computational systems.
This conceptual connection strengthens confidence that the approach
addresses fundamental complexity-theoretic questions rather than
exploiting artificial loopholes.

\begin{center}\rule{0.5\linewidth}{0.5pt}\end{center}

\paragraph{\texorpdfstring{12.9 Quantifier Structure: Why the Uniform
Restriction Enables P \(\neq\)
NP}{12.9 Quantifier Structure: Why the Uniform Restriction Enables P \textbackslash neq NP}}\label{129-quantifier-structure-why-the-uniform-restriction-enables-p--np}

\textbf{Why This Question Matters (The Hardcoding Barrier):} A
fundamental obstacle in complexity theory: you cannot prove a
\emph{single} instance is hard for \emph{all} algorithms (including
non-uniform). Why? Trivial hardcoding defeats any such claim. A
non-uniform algorithm can include "advice" - literally hardcode "if
input = x, output L(x)" - and solve that specific instance in O(1) time.
This means the quantifier structure "\(\exists\)x such that \(\forall\)
algorithms A (including non-uniform), A(x) fails" is \textbf{impossible
to prove} for any problem. This creates a dilemma: how do you prove P
\(\neq\) NP without claiming a single instance defeats all algorithms?

\textbf{Our Solution: OWF Construction with Uniform Algorithms}

The proof circumvents the hardcoding barrier through careful quantifier
structure:

\textbf{Quantifier structure:} \(\forall\)x*\(\in\)L*\_\{FG\},
\(\forall\) uniform algorithm A \(\to\) time \(\geq\) super-poly

\textbf{Why this works:} The uniform restriction is the key. Uniform
algorithms \textbf{cannot hardcode instance-specific solutions} (no
advice, no oracles). They must work via a fixed finite description that
applies to all input sizes. This enables a stronger claim:
\textbf{every} FG-wired instance is hard for \textbf{every} uniform
algorithm. The \(\forall\)x* \(\forall\)A structure - every instance
hard for every uniform algorithm - is what makes the Structural OWF
construction possible. We prove:

\begin{enumerate}
\def\labelenumi{\arabic{enumi}.}
\tightlist
\item
  L* is NP-complete (§10.1-10.3)
\item
  Function f: r \(\mapsto\) Plant(\(\varphi\), r) with FG wiring outputs
  instances where every x* = f(r) has per-instance deterministic
  witness-finding lower bound (Theorem 8.A)
\item
  For any uniform PPT inverter: coin-fixing gives deterministic run;
  successful inversion \(\to\) poly-time witness via Ext (§9.3) \(\to\)
  contradicts per-instance bound \(\to\) f is one-way (§9.4)
\item
  OWF existence \(\to\) FP \(\neq\) FNP \(\to\) P \(\neq\) NP via
  classical bridge (§10.4-10.5)
\end{enumerate}

\textbf{Why Quantifier Structure Validates the Result:}

The quantifier choices aren\textquotesingle{}t arbitrary -
they\textquotesingle{}re \textbf{necessary to circumvent the hardcoding
barrier} while maintaining mathematical rigor:

\begin{enumerate}
\def\labelenumi{\arabic{enumi}.}
\item
  \textbf{Uniform restriction:} By restricting to uniform algorithms, we
  escape the hardcoding barrier. Uniform algorithms cannot trivially
  solve specific instances via lookup tables. This enables the
  \(\forall\)x* \(\forall\)A claim that makes Structural OWF
  construction possible.
\item
  \textbf{Structural necessity:} The lower bounds arise from
  L*\textquotesingle{}s instance-side properties (A1-A5 + FG), not
  algorithmic deficiencies. The conservation law q + \(\Phi\) \(\geq\) R
  holds for \textbf{any} computation on L* - it\textquotesingle{}s a
  mathematical requirement of the problem structure. This is why we can
  make strong claims: we\textquotesingle{}re not saying "algorithms are
  bad," we\textquotesingle{}re saying "L* mathematically requires
  super-polynomial work."
\end{enumerate}

The paper\textquotesingle{}s innovation is showing that
instance-specific structural obligations (made rigorous via A1-A5)
enable quantifier structures that circumvent known barriers while
proving P \(\neq\) NP.

\paragraph{12.10 Proof Dependencies and Where to
Attack}\label{1210-proof-dependencies-and-where-to-attack}

\textbf{For Basic Verification:} See \textbf{Question-Driven Navigation}
(in the navigation section before §1) \(\to\) "How do I verify this
myself?" for the verification path.

\textbf{This section provides:} Dependency structure (what breaks if X
fails?), attack strategies (how to efficiently refute), and
failure-point ranking (where errors are most likely).

\textbf{The Three Critical Dependencies:}

Theorem 7.A (A1-A5 -\textgreater{} SCL) --------+ +-\textgreater{}
Theorem 8.A (Per-Instance Bounds) FG mechanism (Appendix C.1.1) -----+
\textbar{} v Ext construction (§9.3) ----------\textgreater{} OWF
Security (§9.4) --\textgreater{} P != NP (§10.5) \^{} Cartesian
factoring (Lemma J.1-Cart) -----------+

\textbf{1. Theorem 7.A (§7.2.1): A1-A5 \(\to\) q + \(\Phi\) \(\geq\) R}

\begin{itemize}
\tightlist
\item
  \textbf{Attack strategy:} Find counterexample where A1-A5 hold but q +
  \(\Phi\) \textless{} R, OR show Keyedness (A1+A2 \(\to\) disjoint
  addresses) has error
\item
  \textbf{Sub-dependency:} Lemma 7.I (Injectivity \(\to\) Alt\_v
  \(\geq\) 2\^{}(R\_v-q\_v)) + Lemma J.1-Cart (Cartesian factoring)
\item
  \textbf{If this fails:} Everything collapses
\end{itemize}

\textbf{2. Theorem 8.A: FG \(\to\) Per-Instance Bounds}

\begin{itemize}
\tightlist
\item
  \textbf{Attack strategy:} Solve FG-wired instance in poly-time on some
  fixed run, OR find digest bypass loophole
\item
  \textbf{Sub-dependency:} GateDigest\_v wired into seeds +
  \(\Omega\)(n/W\_min) work per segment (Appendix C.1.1); caching does
  not reduce this per-segment cost because digests are seed-dependent
  (see §7.3.6 and Appendix C.1.1)
\item
  \textbf{If this fails:} Structural OWF security fails \(\to\) no P
  \(\neq\) NP
\end{itemize}

\textbf{3. §9.4: Per-Instance Bounds + Ext \(\to\) OWF}

\begin{itemize}
\tightlist
\item
  \textbf{Attack strategy:} Show Ext (§9.3) fails to extract witness
  from successful inversion, OR break coin-fixing argument
\item
  \textbf{Sub-dependency:} Coin-fixing averaging + Ext extraction
\item
  \textbf{If this fails:} No OWF \(\to\) no P \(\neq\) NP
\end{itemize}

\textbf{Most Likely Failure Points (Ranked by Novelty):}

\begin{enumerate}
\def\labelenumi{\arabic{enumi}.}
\tightlist
\item
  \textbf{Cartesian factoring (Lemma J.1-Cart):} Most novel; if H1-H5
  (disjoint pools, injectivity/Keyedness, closure, realizability, no
  cross-coupling) don\textquotesingle{}t ensure independence, this fails
  \(\to\) Theorem 7.A fails \(\to\) everything collapses
\item
  \textbf{FG bypass:} If there\textquotesingle{}s a way to avoid
  \(\Omega\)(n/W\_min) work per segment \(\to\) Theorem 8.A fails
\item
  \textbf{Ext extraction failure:} If some inversions
  don\textquotesingle{}t yield witnesses \(\to\) Structural OWF security
  fails
\end{enumerate}

\textbf{What NOT to check:} Natural problem hardness (§12.5), barriers
(§12.6), quantifiers (§12.9) - already addressed.

\paragraph{12.11 Future Directions}\label{1211-future-directions}

Our proof of P \(\neq\) NP via the Structural OWF construction from L*
opens natural research directions. We focus on three high-impact areas
with clear feasibility paths, plus concrete open questions arising
directly from the proof.

\textbf{Direction 1: Extensions to Classical Models Beyond Turing
Machines}

Our framework currently proves bounds for classical uniform PPT
(probabilistic polynomial-time Turing machines) via per-instance
deterministic bounds extended to randomized algorithms through
coin-fixing (Yao\textquotesingle{}s technique). Natural extensions
include:

\emph{Immediate targets (feasible within existing framework):}

\begin{itemize}
\item
  \textbf{Streaming/external-memory models:} Adapt RWA (Receiving-Window
  Attribution) to I/O constraints. Target: convert \texttt{Q(C)} and
  \texttt{Alt(C)} obligations to I/O complexity bounds via red-blue
  pebbling-style connections.
\item
  \textbf{Parallel/distributed computation:} Extend cut composition to
  multi-cut separators. Study how \texttt{Alt(C)} composes under
  fork/join and whether attribution prevents world-merging across
  workers without resolution.
\item
  \textbf{Interactive protocols:} Parameterize \texttt{Q} by
  communication rounds. Test whether SCL yields round/communication
  lower bounds when Hermeticity restricts side-channels.
\item
  \textbf{Circuits/RAM:} For circuits, relate \(\Phi\)(C) to width/size
  via standard bridges (e.g., width\(\to\)size) and test whether
  \(\lambda\) constrains depth (span). For RAM/word models, adapt
  per-step inflow B and per-segment pricing to word operations;
  milestone: derive size/depth lower bounds Size \(\geq\)
  2\^{}(\(\Omega\)(\(\lambda\))) under bounded width, Depth \(\geq\)
  \(\Omega\)(\(\lambda\)/\(\beta\)*) under bandwidth \(\beta\)*.
\end{itemize}

\emph{Key feasibility factor:} These models satisfy Dependency (A5),
Hermeticity (A1), and Emergence (A3) constraints structurally - adapter
design is primary challenge, not framework extension.

\emph{Concrete milestone:} Prove streaming lower bound for L* showing
\(\lambda\)\_base bits must cross I/O boundary, yielding pass complexity
\(\geq\) n\^{}(\(\Omega\)(log n))/M for memory M.

\textbf{Direction 2: Quantum Exploration - Does SCL Survive
Superposition?}

Having proven P \(\neq\) NP classically, the natural question is whether
quantum algorithms can circumvent the information-theoretic bottleneck
\(\lambda\).

\emph{Core challenge:} SCL\textquotesingle{}s Cartesian factoring (Lemma
J.1-Cart) assumes \texttt{Alt(C)\ =\ Pi\_\{vinC\}\ Alt\_v} via disjoint
address pools (A1) and classical seed-consistency. Quantum superposition
may violate this: can quantum algorithms maintain \ensuremath{\sqrt{2}}\^{}\(\lambda\))
amplitude instead of 2\^{}\(\lambda\) distinguishable states? We propose
starting in restricted settings (depth-bounded quantum circuits or
quantum query models) to scope the question tightly.

\emph{Research questions:}

\begin{itemize}
\tightlist
\item
  \textbf{Q1 (amplitude vs artifacts):} Does Grover search reduce
  \(\lambda\)\_base from \(\Theta\)(log² n) to \(\Theta\)(log n) by
  amplitude amplification, or does seed-dependency structure prevent
  quadratic speedup?
\item
  \textbf{Q2 (measurement-forcing):} Do A1-A5 properties force
  measurement at DAG nodes, collapsing superposition before
  \(\lambda\)-bottleneck and reducing quantum to classical?
\item
  \textbf{Q3 (quantum OWF):} If quantum algorithms can invert our OWF f:
  r \(\mapsto\) x* efficiently, does this yield BQP-complete
  witness-finding for L*, or does Ext extractor fail in quantum setting?
\end{itemize}

\emph{Why high-impact:} Quantum lower bounds are notoriously difficult.
If SCL extends to quantum (Q2 holds), we gain rare information-theoretic
quantum lower bound. If not (Q1 holds), identifies precise mechanism
where quantum evades classical conservation laws.

\emph{Feasibility path:} Start with bounded-depth quantum circuits on
L*. Analyze whether unitary evolution preserves seed-injectivity or
enables amplitude-based compression of Alt\_v. Leverage recent quantum
information-theoretic tools (e.g., quantum entropy cones).

\emph{Concrete milestone:} Prove either (a) quantum algorithm solving L*
in poly(n) time, identifying SCL breakdown point, or (b) quantum lower
bound \(\geq\) n\^{}(\(\omega\)(1)) by showing measurement-forcing at
bottleneck.

\textbf{Direction 3: Minimality and Tightness - Can the Proof Be
Bypassed?}

Every novel proof invites scrutiny of its assumptions. We identify the
most vulnerable components:

\emph{Open Question 1 (Cartesian factoring tightness):} \textbf{Q:} Is
\texttt{Alt(C)\ =\ Pi\_\{vinC\}\ Alt\_v} tight, or can correlations
reduce artifact count below the product? \textbf{Attack strategy:}
Search for seed structure allowing multiple histories to map to same
artifact without violating Injectivity (A2). If found, would weaken
\(\lambda\) by log factors. \textbf{Why we believe it holds:} Disjoint
address pools (A1) + injective Enc (A2) structurally prevent aliasing.
Formal proof in Lemma J.1-Cart.

\emph{Open Question 2 (FG bypass):} \textbf{Q:} Can algorithm avoid
\(\Omega\)(n/W\_min) work per segment without violating correctness?
\textbf{Attack strategy:} Find preprocessing or memoization strategy
that amortizes parity computations across segments. \textbf{Why we
believe it\textquotesingle{}s hard:} Each segment\textquotesingle{}s
parity depends on unique seed history; no cross-segment reuse without
resolution.

\emph{Open Question 3 (A1-A5 minimality):} \textbf{Q:} Can any of
\{Hermeticity (A1), Injectivity (A2), Emergence (A3), Closure (A4),
Dependency (A5)\} be dropped without breaking the proof?
\textbf{Ablation analysis:}

\begin{itemize}
\tightlist
\item
  Drop A1 (Hermeticity) \(\to\) Allows side-channels bypassing DAG,
  enabling polynomial solutions via hidden communication
\item
  Drop A2 (Injectivity) \(\to\) Distinct histories collapse to same
  seed, breaking Cartesian factoring
\item
  Drop A3 (Emergence) \(\to\) No forced freshness R\_v, trivial
  \(\lambda\)=0
\item
  Drop A4 (Closure) \(\to\) Cannot recover ancestors from seeds, breaks
  composition
\item
  Drop A5 (Dependency) \(\to\) Allows out-of-order computation with
  speculative resolution, breaks cut composition
\end{itemize}

\emph{Concrete milestone:} Formalize ablation theorems showing each
A1-A5 axiom is necessary for at least one paradigm\textquotesingle{}s
lower bound.

\textbf{Concrete Open Questions Arising from the Proof:}

\begin{enumerate}
\def\labelenumi{\arabic{enumi}.}
\item
  \textbf{Can \(\lambda\)\_base be compressed?} Is there polynomial-time
  preprocessing reducing L*\textquotesingle{}s \(\lambda\)\_base below
  \(\Theta\)(log² n) without violating A1-A5? (We conjecture no -
  engineered incompressibility.)
\item
  \textbf{Is Ext extraction universal?} Does every successful OWF
  inversion (any r\textquotesingle{} where f(r\textquotesingle{})=x*)
  yield polynomial-time witness extraction via Ext, or exist
  "non-extractable" inversions? (Current proof: Ext works for any valid
  preimage.)
\item
  \textbf{Uniformity necessity:} Does the \(\forall\)x* \(\forall\)A
  quantifier structure (every instance hard for every uniform algorithm)
  require uniformity restriction, or can it be weakened to larger
  non-uniform classes? (Hardcoding barrier suggests no.)
\item
  \textbf{FNP-completeness of witness-finding:} Is finding canonical
  witness W for FG-wired L* instances FNP-complete? (Would strengthen
  classical bridge beyond current "OWF existence" result.)
\end{enumerate}

\paragraph{\texorpdfstring{12.12 Future Work: Two Paths for
\(\lambda\)-Program
Extensions}{12.12 Future Work: Two Paths for \textbackslash lambda-Program Extensions}}\label{1212-future-work-two-paths-for-ux3bb-program-extensions}

\textbf{What this manuscript claims:} P \(\neq\) NP via one NP-complete
language L* with A1-A5 overlays (§§6-10). Construction: explicit OWF
from L* (§9), per-instance lower bounds (§8), classical bridge (§10).
Foundation: Semantic Conservation Law q + \(\Phi\) \(\geq\) R composes
to \(\lambda\)-bottlenecks.

\(\lambda\)(A,x) = \(\Lambda\)(C*) \(-\) Q(C*), \(\Phi\)(C) = log\(_2\)
Alt(C)

\textbf{What \(\lambda\) measures:} At a high level, \(\lambda\)(A, x)
is the residual incompressible information an algorithm must carry
across a semantic bottleneck, as captured by the SCL inequality q +
\(\Phi\) \(\geq\) R at each node of C*.

\textbf{Extension question:} Can \(\lambda\) extend beyond L* to
organize computational complexity more broadly?

\textbf{Two paths forward:} We identify two conceptually different
extension strategies with distinct goals, trade-offs, and prerequisites.

\begin{center}\rule{0.5\linewidth}{0.5pt}\end{center}

\subsubsection{Path 1 (Primary): Representative-Based Class
Separations}\label{path-1-primary-representative-based-class-separations}

\textbf{Core idea:} For each complexity class C, construct one canonical
complete language L*\_C with A1-A5 overlays (adapting L* template).
Measure \(\lambda\)(L*\_C) directly. Separate classes via comparison:
\(\lambda\)(L*\_C\(_1\)) \(\neq\) \(\lambda\)(L*\_C\(_2\)) suggests
C\(_1\) \(\neq\) C\(_2\).

\textbf{Status:} \textbf{Primary extension strategy} -
engineering-focused, uses proven infrastructure, ready to implement.

\textbf{Method:}

\begin{enumerate}
\def\labelenumi{\arabic{enumi}.}
\tightlist
\item
  \textbf{Choose C-complete problem:} TQBF for PSPACE, circuit-value for
  P, \#SAT for \#P, etc.
\item
  \textbf{Engineer A1-A5 overlay:} Adapt L* layer structure (DAG, pools,
  seeds) to problem
\item
  \textbf{Define \(\lambda\)(L*\_C):} Use the same semantic model as for
  L*\_NP, with \(\lambda\)(A, x) defined as the residual \(\Lambda\)(C*)
  \(-\) Q(C*) at the bottleneck, and \(\lambda\)(L*\_C) obtained by
  aggregating over C* as in §8
\item
  \textbf{Prove bounds:} Lower/upper bounds on \(\lambda\)(L*\_C) via
  SCL framework
\item
  \textbf{Establish completeness:} L*\_C is C-complete under appropriate
  reductions
\end{enumerate}

\textbf{Near-term targets (ordered by feasibility):}

\textbf{T1. NP \(\neq\) PSPACE} (highest priority)

\begin{itemize}
\tightlist
\item
  \textbf{L*\_PSPACE:} TQBF-based construction with quantifier-block
  layers
\item
  \textbf{Expected:} \(\lambda\)(L*\_PSPACE) \textgreater{}
  \(\lambda\)(L*\_NP) due to space reuse challenge
\item
  \textbf{Impact:} Major separation using proven template
\end{itemize}

\textbf{T2. P vs BPP} (randomized variant)

\begin{itemize}
\tightlist
\item
  \textbf{L*\_BPP:} Requires handling promise problems, verifiable tests
\item
  \textbf{Challenge:} BPP lacks clean complete problems under poly-time
  reductions
\item
  \textbf{Impact:} Tests \(\lambda\)-framework limits for probabilistic
  classes
\end{itemize}

\textbf{T3. Counting classes} (\#P hierarchy)

\begin{itemize}
\tightlist
\item
  \textbf{L*\_\#P:} Witness-aligned overlay where Alt(C*) = witness
  count
\item
  \textbf{Challenge:} Requires bijection between seed-equivalence
  classes and witnesses
\item
  \textbf{Impact:} Maps \#P operations to \(\lambda\)-arithmetic
\end{itemize}

\textbf{Infrastructure reuse (engineering estimates):}

\begin{itemize}
\tightlist
\item
  \textbf{From L*\_NP:} Layers 0-4 (SCL framework, construction
  template, information bounds, operational semantics) - roughly 70\%
  reusable
\item
  \textbf{Per-class adaptation:} Layer 2 (problem encoding - about
  30-60\% new work depending on class)
\item
  \textbf{Effort estimate:} L*\_PSPACE \(\approx\) 40\% new, L*\_BPP
  \(\approx\) 60\% new
\end{itemize}

\textbf{Advantages:}

\begin{itemize}
\tightlist
\item
  \textbf{Practical:} Uses proven L* template (ready to extend)
\item
  \textbf{Sufficient:} Achieves primary goal (class separations)
\item
  \textbf{Modular:} Each class independent (parallel development
  possible)
\item
  \textbf{Leverages completeness:} Uses decades of completeness theory
\end{itemize}

\textbf{Limitations:}

\begin{itemize}
\tightlist
\item
  \textbf{Not universal:} Doesn\textquotesingle{}t define \(\lambda\)
  for arbitrary languages
\item
  \textbf{Not characterizations:} Separations only, not "C{[}f{]} = P"
  statements
\item
  \textbf{Representative-dependent:} Robustness across different
  complete problems requires validation
\end{itemize}

\textbf{Escalation criterion:} Pursue Path 2 (universal theory) if Path
1 separations fail or characterizations (not just separations) are
needed.

\begin{center}\rule{0.5\linewidth}{0.5pt}\end{center}

\subsubsection{Path 2 (Aspirational): Universal Overlay
Theory}\label{path-2-aspirational-universal-overlay-theory}

\textbf{Core idea:} Define overlays and \(\lambda\) for \emph{all}
languages, not just canonical representatives. Prove class
characterizations like "C{[}O(log n){]} = P". Make \(\lambda\) a
universal complexity measure (like time/space).

\textbf{Status:} \textbf{Long-term research program} - requires
foundational work on overlay invariance, not ready to implement.

\textbf{Central challenge (F6 - CRITICAL):} Overlay invariance

\begin{itemize}
\tightlist
\item
  \textbf{Problem:} Without invariance, \(\lambda\) becomes
  artifact-dependent (different overlays for same language \(\to\)
  different \(\lambda\) values)
\item
  \textbf{Requirement:} Fix overlay families F that are
  verifier-auditable and closed under reductions
\item
  \textbf{Tasks:}

  \begin{enumerate}
  \def\labelenumi{\arabic{enumi}.}
  \tightlist
  \item
    Axiomatize "reduction-safe overlay" (A1-A5 preserved under
    reductions)
  \item
    Prove: L \(\leq\)\_Karp L\textquotesingle{} and L\textquotesingle{}
    \(\in\) F with bound f \(\to\) L \(\in\) F with bound poly(f)
  \item
    Show such families exist
  \end{enumerate}
\end{itemize}

\textbf{Class functionals (F1):}

\begin{itemize}
\tightlist
\item
  \textbf{Definition:} C{[}f{]} := \{L : \(\Lambda\)*\_L(n) \(\in\)
  O(f(n))\} where \(\Lambda\)*\emph{L(n) =
  sup}\{\textbar{}x\textbar{}=n\} inf\_A \(\lambda\)(A,x)
\item
  \textbf{Conjectures:}

  \begin{itemize}
  \tightlist
  \item
    \textbf{C1 (\(\lambda\)\(\leftrightarrow\)P):} C{[}O(log n){]} = P
    for reduction-stable overlays
  \item
    \textbf{C2 (hierarchy):} C{[}polylog(n){]} \ensuremath{\subsetneq}
    C{[}n\^{}\(\varepsilon\){]} \ensuremath{\subsetneq} C{[}\(\Theta\)(n){]} for 0 \textless{}
    \(\varepsilon\) \textless{} 1
  \item
    \textbf{C3 (probabilistic):} BPP \(\subseteq\) C{[}O(log n){]} with
    verifiable tests
  \end{itemize}
\end{itemize}

\textbf{Advantages if successful:}

\begin{itemize}
\tightlist
\item
  \textbf{Universal:} \(\lambda\) defined for all languages (not just
  representatives)
\item
  \textbf{Characterizations:} "C{[}f{]} = P" statements, not just
  separations
\item
  \textbf{Invariant:} \(\lambda\) becomes fundamental measure (like
  time/space complexity)
\item
  \textbf{Reduction-stable:} Preserves structure under Karp reductions
\end{itemize}

\textbf{Prerequisites (foundational work required):}

\begin{itemize}
\tightlist
\item
  Solve overlay invariance (F6) - no artifacts
\item
  Prove closure under reductions
\item
  Validate across natural NP problems (not just L*)
\item
  \textbf{Timeline:} 4+ years (requires solving open questions)
\end{itemize}

\textbf{Honest assessment:} This is a research program, not accomplished
fact. Success depends on solving F6 (overlay invariance) and empirical
validation showing \(\lambda\) predicts behavior beyond constructed
instances.

\begin{center}\rule{0.5\linewidth}{0.5pt}\end{center}

\subsubsection{Comparison and Decision
Criteria}\label{comparison-and-decision-criteria}

\textbf{Path 1 (Representative) vs Path 2 (Universal):}

\begin{itemize}
\item
  \textbf{Goal}

  \begin{itemize}
  \tightlist
  \item
    Path 1: Class separations
  \item
    Path 2: Class characterizations
  \end{itemize}
\item
  \textbf{Scope}

  \begin{itemize}
  \tightlist
  \item
    Path 1: One L*\_C per class
  \item
    Path 2: All languages
  \end{itemize}
\item
  \textbf{Infrastructure}

  \begin{itemize}
  \tightlist
  \item
    Path 1: Uses proven L* template
  \item
    Path 2: Requires new theory
  \end{itemize}
\item
  \textbf{Readiness}

  \begin{itemize}
  \tightlist
  \item
    Path 1: Ready now
  \item
    Path 2: 4+ years
  \end{itemize}
\item
  \textbf{Prerequisites}

  \begin{itemize}
  \tightlist
  \item
    Path 1: Engineering only
  \item
    Path 2: F6 (overlay invariance)
  \end{itemize}
\item
  \textbf{Advantages}

  \begin{itemize}
  \tightlist
  \item
    Path 1: Practical, modular
  \item
    Path 2: Universal, fundamental
  \end{itemize}
\item
  \textbf{Limitations}

  \begin{itemize}
  \tightlist
  \item
    Path 1: Not universal
  \item
    Path 2: Requires solving open questions
  \end{itemize}
\end{itemize}

\textbf{Recommended strategy:} Focus on Path 1 for 2-4 years. Evidence
from multiple L*\_C representatives (robustness,
\(\lambda\)-differences, structural explanations) will inform whether
Path 2 is necessary.

\textbf{Escalate to Path 2 if:}

\begin{itemize}
\tightlist
\item
  Path 1 separations fail (\(\lambda\)-values too close or
  encoding-dependent)
\item
  Different representatives for same class yield incompatible
  \(\lambda\)
\item
  Need characterizations ("C{[}f{]} = P") beyond separations
\end{itemize}

\textbf{Path 1 suffices if:}

\begin{itemize}
\tightlist
\item
  Robust separations across multiple representatives
\item
  \(\lambda\)-differences are large (not just polynomial factors)
\item
  Structural explanations for \(\lambda\)-values emerge
\end{itemize}

\begin{center}\rule{0.5\linewidth}{0.5pt}\end{center}

\subsubsection{Path 2 Detailed Directions (If
Pursued)}\label{path-2-detailed-directions-if-pursued}

\emph{The following Tier 1-3 directions assume Path 2 (Universal Overlay
Theory) is pursued. They require solving F6 (overlay invariance) first.}

\textbf{Caution (scope of this section):} The Tier 1-3 directions below
assume Path 2 (Universal Overlay Theory) is pursued. Generalizing
\(\lambda\) beyond L* requires constructing A1-A5 overlays for each
target; natural NP languages do not necessarily satisfy these axioms
without augmentation (§12.5). This heavy "overlay-for-every-language"
work belongs only to Path 2, not to the representative-based program
(Path 1).

\paragraph{Tier 1: Foundational Theory (Path 2
prerequisites)}\label{tier-1-foundational-theory-path-2-prerequisites}

\subparagraph{\texorpdfstring{F1. A \(\lambda\)-Spectrum of
Classes}{F1. A \textbackslash lambda-Spectrum of Classes}}\label{f1-a-ux3bb-spectrum-of-classes}

\textbf{Definition (best-achievable residual).} For input x,
\(\lambda\)\^{}*(x) := inf\_A \(\lambda\)(A,x). For a language L,
\(\Lambda\)\^{}*\_L(n) := sup\_\{\textbar{}x\textbar{}=n\}
\(\lambda\)\^{}*(x).

\textbf{Class family.} C{[}f{]} := \{L : \(\Lambda\)\^{}*\_L(n) \(\in\)
O(f(n))\}.

\textbf{{[}YES{]} Proven:} If L \(\in\) C{[}f{]} then L has time
\(\leq\) 2\^{}(O(f(n))) under fixed verifier-auditable overlay with
per-run pricing. L* \(\in\) C{[}\(\Theta\)(log² n){]} (QP-sharp) or
C{[}\(\Theta\)(n){]} (flat).

\textbf{Open problems:}

\begin{enumerate}
\def\labelenumi{\arabic{enumi}.}
\tightlist
\item
  \textbf{(\(\lambda\)\(\leftrightarrow\)P conjecture)} Is C{[}O(log
  n\_core){]} = P for suitably fixed, verifier-auditable overlays?
\item
  \textbf{Closure under Karp reductions.} Identify overlay families
  closed under reductions so C{[}f{]} is reduction-stable within that
  overlay family.
\item
  \textbf{Separations by \(\lambda\).} Show C{[}O(log n\_core){]} \ensuremath{\subsetneq}
  C{[}polylog(n){]} \ensuremath{\subsetneq} C{[}n\^{}\(\varepsilon\){]} \ensuremath{\subsetneq} C{[}\(\Theta\)(n){]}
  for fixed overlays.
\item
  \textbf{Fine-grained separations.} SETH as \(\lambda\)\_SAT(n) = cn
  for specific c \textless{} 1; 3SUM/APSP conjectures as
  \(\lambda\)-barriers within P.
\end{enumerate}

\textbf{First milestone:} Prove \(\lambda\)\(\leftrightarrow\)P for one
natural NP-complete problem beyond L* (e.g., 3-SAT with
verifier-auditable overlay).

\begin{center}\rule{0.5\linewidth}{0.5pt}\end{center}

\subparagraph{F6. Overlay Invariance \& Reduction-Stable Families (!)
CRITICAL}\label{f6-overlay-invariance--reduction-stable-families-warning-critical}

\textbf{Goal.} Fix overlay families F that are
\textbf{verifier-auditable} and \textbf{closed under reductions}, so
\(\lambda\) is not artifact-dependent.

\textbf{Why critical:} Without overlay invariance, entire
\(\lambda\)-program becomes construction-dependent. Different overlays
could yield different \(\lambda\) values for same language, preventing
C{[}f{]} from being well-defined complexity classes.

\textbf{Tasks:}

\begin{enumerate}
\def\labelenumi{\arabic{enumi}.}
\tightlist
\item
  Axiomatize "reduction-safe overlay" (Keyedness, Injectivity,
  Hermeticity preserved under f, witness-preserving).
\item
  Prove: if L \(\leq\)\_\{Karp\} L\textquotesingle{} and
  L\textquotesingle{} \(\in\) F with bound f, then L \(\in\) F with
  bound n \(\mapsto\) f(poly(n\_core)).
\item
  Show overlay families exist that are closed under witness-preserving
  reductions.
\end{enumerate}

\begin{center}\rule{0.5\linewidth}{0.5pt}\end{center}

\subparagraph{F10. Core Conjectures}\label{f10-core-conjectures}

\begin{itemize}
\tightlist
\item
  \textbf{C1 (\(\lambda\)\(\leftrightarrow\)P).} For fixed,
  reduction-stable overlays, C{[}O(log n){]} = P.
\item
  \textbf{C2 (strict \(\lambda\)-spectrum).} For any 0 \textless{}
  \(\varepsilon\) \textless{} 1: C{[}polylog(n){]} \ensuremath{\subsetneq}
  C{[}n\^{}\(\varepsilon\){]} \ensuremath{\subsetneq} C{[}\(\Theta\)(n){]}.
\item
  \textbf{C3 (probabilistic \(\lambda\)).} In overlays with verifiable
  tests, BPP \(\subseteq\) C{[}O(log n){]}.
\item
  \textbf{C4 (space pricing).} There exists a general space-pricing
  theorem making a PSPACE analogue of our time lower bounds.
\end{itemize}

\begin{center}\rule{0.5\linewidth}{0.5pt}\end{center}

\subsubsection{Tier 2: Model Extensions (natural generalizations of
SCL)}\label{tier-2-model-extensions-natural-generalizations-of-scl}

\subparagraph{F2. Randomized Classes via "Verifiable
Tests"}\label{f2-randomized-classes-via-verifiable-tests}

\textbf{Goal.} Formalize *\emph{probabilistic SCL*} when random guesses
are self-verifiable.

\textbf{Tasks:}

\begin{enumerate}
\def\labelenumi{\arabic{enumi}.}
\tightlist
\item
  Define Q\^{}(rand)(C*) := E\_\{coins\}{[}legitimate first-use bits
  resolved{]} in runs with randomness and verifiable checks.
\item
  Prove BPP-style bounds: if E{[}\(\lambda\)(A,x){]} \(\leq\) f(n) and
  tests reject wrong guesses w.h.p., then time 2\^{}(O(f(n))) after
  amplification.
\item
  Map RP/ZPP by one-sided/expected variants of Q\^{}(rand).
\end{enumerate}

\textbf{Note:} Requires verifiable tests so wrong guesses are rejected
w.h.p.; otherwise randomness need not increase Q.

\textbf{Outcome.} A rigorous bridge for "randomness as
\(\lambda\)-reduction" when verification exists.

\begin{center}\rule{0.5\linewidth}{0.5pt}\end{center}

\subparagraph{F3. Space-Time Dual
Pricing}\label{f3-space-time-dual-pricing}

\textbf{Goal.} A \textbf{space-pricing lemma} that complements our time
pricing.

\textbf{Open question:} Why doesn\textquotesingle{}t
Savitch\textquotesingle{}s theorem provide sufficient space-\(\lambda\)
connection? Conjecture: Savitch allows space reuse across branches via
depth-first traversal; \(\lambda\)-pricing may require persistent
storage across simultaneously live worlds (breadth-first bottleneck).

\textbf{Tasks:}

\begin{enumerate}
\def\labelenumi{\arabic{enumi}.}
\tightlist
\item
  Branching programs: prove width W \(\Rightarrow\) space \(\geq\)
  \(\Omega\)(log W) per level; lift to \(\lambda\) across cuts.
\item
  TMs: identify conditions under which maintaining W live worlds across
  L steps forces \(\Omega\)(log W) workspace (beyond Savitch-style
  reuse).
\item
  Characterize languages where \(\lambda\) can be "paid" with space
  reuse (PSPACE-like regimes) vs where \(\lambda\) forces time blowup.
\end{enumerate}

\begin{center}\rule{0.5\linewidth}{0.5pt}\end{center}

\subparagraph{F4. Counting Classes (Witness-Aligned
Overlays)}\label{f4-counting-classes-witness-aligned-overlays}

\textbf{Goal.} Align \(\Phi\)(C*) = log\(_2\)(Alt(C*)) with \textbf{\#
of accepting witnesses}.

\textbf{Critical assumption:} Requires \textbf{witness-aligned overlays}
where seed-consistent classes on a canonical cut biject to
witnesses/paths. This may not hold for natural problems - overlay
engineering may be necessary.

\textbf{Tasks:}

\begin{enumerate}
\def\labelenumi{\arabic{enumi}.}
\tightlist
\item
  Design witness-aligned overlays where seed-consistent classes on a
  canonical cut biject to witnesses/paths.
\item
  In witness-aligned overlays, show: \#P computes Alt(C*), PP decides
  majority over Alt(C*), and \(\oplus\)P decides parity.
\item
  \textbf{Characterize when bijection fails:} Natural problems may have
  witnesses that don\textquotesingle{}t correspond to cut-artifacts.
  Identify structural conditions for alignment vs obstacles.
\end{enumerate}

\begin{center}\rule{0.5\linewidth}{0.5pt}\end{center}

\subparagraph{F5. Relativized SCL
(Oracles)}\label{f5-relativized-scl-oracles}

\textbf{Goal.} Define Q\^{}(O)(C*) that credits \textbf{oracle-learned}
bits and prove (as a target):

Q\^{}(O)(C*) + log\(_2\) Alt(C*) \(\geq\) \(\Lambda\)(C*).

\textbf{Tasks:}

\begin{enumerate}
\def\labelenumi{\arabic{enumi}.}
\tightlist
\item
  Specify \textbf{relativized overlays} (oracle transcripts as
  designated artifacts).
\item
  Show examples where P\^{}(O) collapses \(\lambda\) (easy oracles) and
  where NP\^{}(O) preserves \(\lambda\) (hard oracles).
\item
  Study which classical relativized separations are simply
  \textbf{\(\lambda\)-shifts}.
\item
  Extend to \textbf{interactive oracles}: IP/AM as interactive
  \(\lambda\)-resolution where prover provides Q-hints; MIP as
  distributed oracle \(\lambda\)-verification across multiple provers.
\end{enumerate}

\begin{center}\rule{0.5\linewidth}{0.5pt}\end{center}

\subparagraph{F5b. Formal SCL Embeddings for Classical Lower Bound
Techniques}\label{f5b-formal-scl-embeddings-for-classical-lower-bound-techniques}

\textbf{Goal.} Mechanically verify that classical lower bound techniques
are formal instantiations of SCL.

\textbf{Background (§11.4).} The TM observation paradigm is fully
formalized in Lean (\texttt{lean/Layer4\_Operational/TimeBridge/}). The
correspondences for other techniques---decision trees, pebbling games,
communication complexity, branching programs, resolution, and
algorithmic paradigms (backtracking, DP, CDCL, streaming)---are
currently conceptual mappings, mathematically motivated but not
mechanically verified.

\textbf{Tasks:}

\begin{enumerate}
\def\labelenumi{\arabic{enumi}.}
\tightlist
\item
  \textbf{Decision trees:} Formalize query = q, tree nodes =
  2\^{}\(\Phi\); prove that decision tree depth lower bounds instantiate
  SCL.
\item
  \textbf{Pebbling games:} Formalize placements = q, pebbles = \(\Phi\);
  prove Lengauer-Tarjan time-space tradeoffs are SCL instances.
\item
  \textbf{Communication complexity:} Formalize bits exchanged = q,
  rectangles = 2\^{}\(\Phi\); show rectangle covering bounds follow from
  SCL.
\item
  \textbf{Branching programs/OBDD:} Formalize path length = q, width =
  2\^{}\(\Phi\); prove width lower bounds instantiate SCL.
\item
  \textbf{Resolution:} Formalize proof steps = q, clause width =
  \(\Phi\); verify Ben-Sasson-Wigderson width\(\to\)size relationship
  maps to SCL (note: this mapping is less direct than others).
\item
  \textbf{Algorithmic paradigms:} Show backtracking, DP, CDCL, and
  streaming lower bounds follow SCL structure.
\end{enumerate}

\textbf{Priority:} Decision trees and pebbling (cleanest
correspondences) first; resolution last (most interpretive).

\textbf{Outcome.} If successful, transforms "SCL captures" into "SCL
formally unifies"---strengthening §11.4 from conceptual insight to
mechanized proof. If some embeddings fail, identifies which techniques
are genuine SCL instances vs. analogies.

\begin{center}\rule{0.5\linewidth}{0.5pt}\end{center}

\subparagraph{F6. Overlay Invariance \& Reduction-Stable
Families}\label{f6-overlay-invariance--reduction-stable-families}

\textbf{Goal.} Fix overlay families F that are
\textbf{verifier-auditable} and \textbf{closed under reductions}, so
\(\lambda\) is not artifact-dependent.

\textbf{Tasks:}

\begin{enumerate}
\def\labelenumi{\arabic{enumi}.}
\tightlist
\item
  Axiomatize "reduction-safe overlay" (Keyedness, Injectivity,
  Hermeticity preserved under f, witness-preserving).
\item
  Prove: if L \(\leq\)\_\{Karp\} L\textquotesingle{} and
  L\textquotesingle{} \(\in\) F with bound f, then L \(\in\) F with
  bound n \(\mapsto\) f(poly(n\_core)).
\end{enumerate}

\begin{center}\rule{0.5\linewidth}{0.5pt}\end{center}

\subsubsection{Tier 3: Applied Connections (validation \&
extensions)}\label{tier-3-applied-connections-validation--extensions}

\subparagraph{\texorpdfstring{F7. Algorithmic \(\lambda\)-Reduction
Primitives}{F7. Algorithmic \textbackslash lambda-Reduction Primitives}}\label{f7-algorithmic-ux3bb-reduction-primitives}

\textbf{Goal.} Formalize solver design principles through
\(\lambda\)-lens.

\textbf{Rigorous foundations:}

\begin{enumerate}
\def\labelenumi{\arabic{enumi}.}
\tightlist
\item
  \textbf{\(\Lambda\)-reduction:} structural preprocessing that provably
  lowers \(\Lambda\)(C*) (safe eliminations, kernels, decompositions).
\item
  \textbf{Q-acceleration:} propagation/learning that provably raises
  Q(C*) early (\(\lambda\)-guided branching).
\item
  \textbf{Parameterized complexity:} FPT when \(\lambda\) depends only
  on parameter k; kernelization as k-dependent \(\Lambda\)-reduction;
  W-hierarchy as \(\lambda\)-growth rates in k.
\end{enumerate}

\textbf{Engineering challenges (non-rigorous, practical):} 4.
\textbf{Estimators:} practical proxies \(\Lambda\)\^{}, Q\^{} to predict
residual \(\lambda\)\^{} and steer heuristics - requires empirical
validation (see F8).

\textbf{Open question:} Do existing SAT/CSP heuristics (VSIDS, restarts,
clause learning) implicitly minimize \(\lambda\)? Empirical correlation
study needed.

\begin{center}\rule{0.5\linewidth}{0.5pt}\end{center}

\subparagraph{\texorpdfstring{F8. Empirical
\(\lambda\)-Benchmarks}{F8. Empirical \textbackslash lambda-Benchmarks}}\label{f8-empirical-ux3bb-benchmarks}

\textbf{Goal.} Measure \(\lambda\)\^{} on real instances and correlate with
solver time.

\textbf{Tasks:}

\begin{enumerate}
\def\labelenumi{\arabic{enumi}.}
\tightlist
\item
  Implement auditable overlays for SAT/VC/TSP minors; compute
  \(\Lambda\)\^{} and first-use credits Q\^{}.
\item
  Validate time \(\approx\) poly(n\_core) · 2\^{}(\(\lambda\)\^{}) (up to
  pricing factors) across distributions; ablations for \(\Lambda\)-vs-Q
  contributions.
\end{enumerate}

\begin{center}\rule{0.5\linewidth}{0.5pt}\end{center}

\subparagraph{F12. Extensions to Cryptographic Primitives Beyond
OWFs}\label{f12-extensions-to-cryptographic-primitives-beyond-owfs}

\textbf{{[}YES{]} Proven:} We constructed an explicit one-way function
f: r \(\mapsto\) x* via Plant(\(\varphi\), r) with FG-wiring (§9). Every
output x* has per-instance deterministic witness-finding lower bound,
contradicting poly-time inversion via Ext extractor.

\textbf{Future work:} Extend \(\lambda\)-based characterization to other
cryptographic primitives.

\textbf{Scope clarification:} This does NOT claim \(\lambda\) is "the
source" of cryptographic security - cryptography has well-established
foundations. We explore whether \(\lambda\)-framework provides
\textbf{alternative characterizations} or \textbf{new constructions} for
primitives beyond OWFs.

\textbf{Tasks:}

\begin{enumerate}
\def\labelenumi{\arabic{enumi}.}
\tightlist
\item
  \textbf{Formalize \(\lambda\)-security:} A keyed family
  \{D\_\(\kappa\)\} is \(\lambda\)-secure if for all PPT A:
  Pr{[}\(\lambda\)(A,x) \(\geq\) g(\(\kappa\)){]} \(\geq\) 1 \(-\)
  negl(\(\kappa\)).
\item
  \textbf{Equivalence conjecture:} Does "\(\lambda\)-secure
  \(\leftrightarrow\) one-way function" hold under suitable overlay
  families? (NOT proven - requires showing both directions; our
  construction shows one direction.)
\item
  \textbf{Extensions to other primitives:} Can PRGs, PRFs,
  collision-resistant hash functions be characterized via
  \(\lambda\)-preservation? Our OWF (§9.4) suggests this direction may
  be fruitful.
\item
  \textbf{Hardness amplification:} Does repetition/gap coding raise
  \(\lambda\) \(\to\) \(\Omega\)(t·\(\lambda\)) as expected from
  traditional amplification?
\item
  \textbf{Lattice connections:} Explore whether LWE/LPN can be
  formulated with \(\lambda\)-secure overlays.
\end{enumerate}

\textbf{Outcome.} If successful, provides \(\lambda\)-based design
principles for cryptographic primitives; if unsuccessful, identifies
limitations of \(\lambda\)-framework for cryptographic settings.

\begin{center}\rule{0.5\linewidth}{0.5pt}\end{center}

\subparagraph{F12a. OWF-Enabled Cryptographic Primitives and Structural
Limitations}\label{f12a-owf-enabled-cryptographic-primitives-and-structural-limitations}

\textbf{What the L* OWF enables:} The unconditional Structural OWF
constructed in §9 immediately yields the following cryptographic
primitives via standard black-box reductions:

\begin{itemize}
\tightlist
\item
  \textbf{Pseudorandomness:} PRGs via
  HILL/Håstad-Impagliazzo-Levin-Luby, PRFs via GGM tree construction,
  stream/block-cipher cores
\item
  \textbf{Commitment schemes:} Statistically binding or hiding
  (depending on construction), foundation for many protocols
\item
  \textbf{Digital signatures:} Lamport/Winternitz-style hash-based
  signatures, PRF/MAC-based schemes (Naor-Yung)
\item
  \textbf{Message authentication:} MACs from PRFs, thus from OWFs (via
  PRF construction)
\item
  \textbf{Oblivious transfer / secure computation:} Black-box OT from
  OWFs with interaction (Impagliazzo-Rudich positive direction),
  enabling general MPC
\item
  \textbf{Zero-knowledge proofs:} Computational ZK for NP via commitment
  schemes
\item
  \textbf{Hardness amplification / hashing:} Collision-resistant hashing
  from stronger OWF variants (claw-free, regular, 2-to-1)
\item
  \textbf{Password hashing / proof-of-work:} Salted OWF instantiations;
  memory-hard variants build on OWF baselines
\item
  \textbf{Identification/authentication protocols:} Challenge-response
  protocols built from OWF/PRG pairs
\end{itemize}

\textbf{What the L* OWF does NOT enable (structural limitation):}

\textbf{Public-key encryption (PKE)} and \textbf{non-interactive key
agreement} cannot be constructed from a bare OWF via black-box
techniques. This is NOT a limitation of L* specifically---it is a
fundamental barrier:

\begin{itemize}
\item
  \textbf{Impagliazzo-Rudich barrier (1989):} No black-box construction
  of PKE or key agreement from a generic OWF. There exists an oracle
  relative to which OWFs exist but PKE does not.
\item
  \textbf{Structural reason:} PKE requires a
  \textbf{trapdoor}---asymmetric hardness where inversion is hard for
  everyone \emph{except} the secret key holder. A bare OWF provides
  uniform hardness (hard for everyone). L* has no trapdoor, permutation
  structure, or algebraic group/lattice properties that would enable
  selective inversion.
\item
  \textbf{The verification gap:} NP verification (given witness, check
  correctness) differs from PKE decryption (given ciphertext, recover
  message). Verification checks proposed answers; decryption extracts
  hidden information. The FP \(\neq\) FNP separation proves "finding is
  hard for everyone"---PKE needs "finding is hard for everyone
  \emph{except} the key holder."
\item
  \textbf{No message embedding:} The L*/OWF architecture has no slot for
  a recoverable message. In L*(x, w) \(\to\) y, there is no "plaintext"
  that gets encrypted and can later be decrypted. The witness w is an
  input to the computation, not a key that inverts outputs.
\end{itemize}

\textbf{Future direction (PKE from L*):} Achieving PKE would require
augmenting L* with additional structure:

\begin{enumerate}
\def\labelenumi{\arabic{enumi}.}
\tightlist
\item
  \textbf{Trapdoor permutation variant:} Engineer L* variant where a
  planted trapdoor enables efficient inversion for the trapdoor holder
\item
  \textbf{Lattice/group structure:} Embed L* into algebraic setting
  (LWE-style, group-based) providing homomorphic or trapdoor properties
\item
  \textbf{Non-black-box techniques:} Explore whether
  L*\textquotesingle{}s specific structure (not just OWF property)
  enables PKE construction outside Impagliazzo-Rudich barrier
\end{enumerate}

\textbf{Status:} Speculative. The Impagliazzo-Rudich barrier suggests
this is unlikely without substantial structural augmentation. The L*
OWF\textquotesingle{}s value is in what it \emph{does} enable
unconditionally (symmetric cryptography, commitments, signatures, ZK,
MPC with interaction), not in bypassing known impossibility results.

\begin{center}\rule{0.5\linewidth}{0.5pt}\end{center}

\subparagraph{\texorpdfstring{F13. Approximation and
\(\lambda\)-Relaxation}{F13. Approximation and \textbackslash lambda-Relaxation}}\label{f13-approximation-and-ux3bb-relaxation}

\textbf{Goal.} Extend \(\lambda\)-framework to approximate solutions by
relaxing exactness requirements.

\textbf{Tasks:}

\begin{enumerate}
\def\labelenumi{\arabic{enumi}.}
\tightlist
\item
  Define **approximate SCL**: Q(C) + log\(_2\)
  Alt\^{}((\(\varepsilon\)))(C) \(\geq\)
  \(\Lambda\)\^{}((\(\varepsilon\)))(C) for \(\varepsilon\)-equivalent
  worlds.
\item
  Prove PTAS \(\Rightarrow\) \(\lambda\)\^{}((\(\alpha\))) =
  O(poly(1/\(\alpha\)) log n) for \(\alpha\)-approximation.
\item
  Show APX-hardness \(\Rightarrow\) large \(\lambda\)\^{}((\(\alpha\)))
  via gap-introducing reductions.
\item
  Interpret LP/SDP relaxations as \(\Lambda\)-reduction; rounding as
  Q-acceleration.
\end{enumerate}

\textbf{PCP connection (exploratory):} PCP theorem achieves verification
with O(1) queries through gap amplification. Explore whether this can be
recast as \(\lambda\)-gap amplification - maintaining incompressibility
while reducing verification cost. \textbf{Caveat:} PCP theory is
well-established independently; question is whether \(\lambda\)-lens
provides new insights or constructions, not whether it "explains" PCP.

\textbf{Outcome.} If approximate SCL holds with graceful degradation,
provides unified view of approximation algorithms as
\(\lambda\)-management under relaxed correctness. If \(\lambda\)
degrades catastrophically (e.g., \(\lambda\)\^{}((\(\varepsilon\))) =
O(log n) for all \(\varepsilon\) \textless{} 1), identifies limitation
of framework for approximation settings.

\begin{center}\rule{0.5\linewidth}{0.5pt}\end{center}

\subparagraph{\texorpdfstring{F14. Blockchain Protocol Analysis via
\(\lambda\)-Framework}{F14. Blockchain Protocol Analysis via \textbackslash lambda-Framework}}\label{f14-blockchain-protocol-analysis-via-ux3bb-framework}

\textbf{Context:} Blockchain proof-of-work protocols (Bitcoin, Ethereum
pre-merge) enforce computational work as security mechanism. Mining
requires \textasciitilde{}2\^{}difficulty hash inversions; consensus
selects chain with maximum cumulative work.

\textbf{Research question:} Can \(\lambda\)-framework provide
quantitative security analysis or improved protocol design beyond
existing blockchain theory?

\textbf{Relationship to existing work:} Nakamoto consensus, selfish
mining analysis (Eyal-Sirer), and proof-of-work security are
well-established. We explore whether \(\lambda\)-formalism offers:

\begin{enumerate}
\def\labelenumi{\arabic{enumi}.}
\tightlist
\item
  \textbf{Tighter security bounds} via \(\lambda\)-counting vs
  traditional computational security
\item
  \textbf{Protocol design principles} (\(\lambda\)-optimal: minimize
  honest \(\lambda\), maximize adversarial \(\lambda\))
\item
  \textbf{Formal verification} of consensus properties through
  \(\lambda\)-conservation
\end{enumerate}

\textbf{Tasks:}

\begin{enumerate}
\def\labelenumi{\arabic{enumi}.}
\tightlist
\item
  \textbf{Formalize proof-of-work via \(\lambda\):} Mining maintains
  \(\lambda\) \(\approx\) log\(_2\)(difficulty); difficulty adjustment
  targets expected block time.
\item
  \textbf{Consensus as \(\lambda\)-selection:} Longest chain rule
  selects branch with maximum \(\Sigma\) \(\lambda\)\_block.
\item
  \textbf{Attack analysis:} 51\% attack requires maintaining adversarial
  \(\lambda\)\_rate \textgreater{} honest \(\lambda\)\_rate; selfish
  mining as \(\lambda\)-hoarding strategy - can \(\lambda\)-framework
  provide tighter bounds than existing analysis?
\item
  \textbf{\(\lambda\)-optimal protocol design:} Given security target
  \(\lambda\)\_sec, minimize honest work \(\lambda\)\_honest subject to
  adversarial work \(\geq\) \(\lambda\)\_sec. Does this yield novel
  protocols or recover existing ones?
\item
  \textbf{Proof-of-stake extension (exploratory):} Can virtual stake be
  formalized as \(\lambda\)-capacity? Does slashing correspond to
  \(\lambda\)-penalty?
\end{enumerate}

\textbf{Outcome:} If \(\lambda\)-framework provides tighter security
bounds or novel protocol designs, demonstrates practical value beyond
theoretical lower bounds. If \(\lambda\) merely reparameterizes existing
analysis, identifies limitation of framework for distributed consensus
settings.

\textbf{Scope:} Focus on computational proof-of-work; proof-of-stake is
speculative extension. No claims about market valuations or "explaining"
existing deployments.

\begin{center}\rule{0.5\linewidth}{0.5pt}\end{center}

\subparagraph{\texorpdfstring{F15. Communication Complexity and
Distributed
\(\lambda\)}{F15. Communication Complexity and Distributed \textbackslash lambda}}\label{f15-communication-complexity-and-distributed-ux3bb}

\textbf{Goal.} Extend \(\lambda\)-framework to multi-party settings
where information is distributed across k players.

\textbf{Relationship to existing theory:} Communication complexity has
established lower bounds via information-theoretic arguments (e.g., Yao,
Kushilevitz-Nisan). This direction explores whether
\(\lambda\)-framework provides alternative derivations or strengthens
existing bounds for problems with SCL-compatible structure.

\textbf{Tasks:}

\begin{enumerate}
\def\labelenumi{\arabic{enumi}.}
\tightlist
\item
  Define \textbf{distributed \(\lambda\)} across k parties: partition C*
  into C\_1,...,C\_k; each party sees only their subgraph.
\item
  Prove \textbf{lifting theorems}: decision tree \(\lambda\) lifts to
  communication \(\lambda\) via gadget composition.
\item
  Show \textbf{direct sum/product theorems}: \(\lambda\)(f\^{}n)
  \(\approx\) n·\(\lambda\)(f) for composed problems.
\item
  Analyze \textbf{number-on-forehead model}: shared information reduces
  effective \(\lambda\) exponentially.
\item
  Connect to \textbf{information complexity}: IC(f) as expected
  \(\lambda\) over input distribution.
\end{enumerate}

\textbf{Outcome.} If successful, provides \(\lambda\)-based derivations
of communication lower bounds, potentially extending to settings where
traditional information complexity faces obstacles (e.g., exact vs
approximate communication, non-product distributions).

\begin{center}\rule{0.5\linewidth}{0.5pt}\end{center}

\textbf{Vision.} If successful, this program would provide a unifying
lens for computational complexity: a single measure (\(\lambda\)), a
single law (Q + \(\Phi\) \(\geq\) \(\Lambda\)), and a spectrum C{[}f{]}
to position problems, algorithms, and oracles.

\textbf{Aspirational goal:} Reduce diverse complexity phenomena to one
conservation principle - randomized classes as strategies for managing
\(\lambda\) debt, space-bounded classes as alternative payment methods,
counting classes as operations on forced states, oracle separations as
artificial \(\lambda\)-shifts.

\textbf{Honest assessment:} The broader \(\lambda\)-program outlined in
§12.12 is a research program, not accomplished fact. The results proved
in §§1--10 (SCL for L*, per-instance deterministic bounds, unconditional
Structural OWF construction, and the classical bridge under the
classical uniform PPT model) are established. The ``program'' refers to
generalizations and extensions beyond L* and beyond the classical scope
(e.g., randomized/oracle/quantum models), where additional foundational
work (Tier 1) is required. Empirical validation (F8) will test whether
\(\lambda\) predicts real-world solver behavior or remains primarily a
proof technique for constructed instances.

\textbf{Key insight:} The universality lies in the conservation law
itself, not in requiring one instance to be hard for all algorithms.
Every algorithm faces the same \(\lambda\)-bottleneck on its respective
hard instances - just as energy conservation applies universally even
though different systems store energy differently.

\textbf{§12.12 Summary:} Proposed \(\lambda\)-program explores whether
residual \(\lambda\) can unify complexity theory: Tier 1 (foundational -
C{[}f{]} classes, overlay invariance, \(\lambda\)\(\leftrightarrow\)P
conjecture); Tier 2 (model extensions - randomized/counting/oracle/space
classes via Q/\(\Phi\)/\(\Lambda\) adaptations); Tier 3 (applied -
cryptographic primitives, approximation, blockchain, communication
complexity). Success depends on overlay invariance (F6) and empirical
validation (F8); honest assessment acknowledges this is research
program, not proven fact.

\begin{center}\rule{0.5\linewidth}{0.5pt}\end{center}

\section{Appendices}\label{appendices}

\subsection{Table of Contents}\label{table-of-contents}

\textbf{How to read:} The OWF proof uses Critical Path appendices (A, B,
C, J, O). Model-Specific appendices contain detailed bounds for
particular computational models. Reader Resources include FAQ, glossary,
and supplementary analysis.

\begin{center}\rule{0.5\linewidth}{0.5pt}\end{center}

\subsubsection{Critical Path: OWF Proof
(Essential)}\label{critical-path-owf-proof-essential}

\emph{Required for understanding P \(\neq\) NP via Structural OWF
construction (§9) and classical bridge (§10)}

\begin{itemize}
\tightlist
\item
  \textbf{Appendix A:} Explicit Overlay Construction (L* instance with
  A1-A5 properties)
\item
  \textbf{Appendix B:} Technical Lemmas for Paradigm Bounds (DT, DP,
  OBDD, Resolution adapters)
\item
  \textbf{Appendix C:} Frontier-Gate and Segment Counting (achieving
  tight per-instance bounds)
\item
  \textbf{Appendix J:} DAG Min-Cut Lower Bound (Cartesian factoring, cut
  composition)
\item
  \textbf{Appendix O:} Unconditional OWF Packaging Details (Plant
  construction, Ext extractor)
\end{itemize}

\begin{center}\rule{0.5\linewidth}{0.5pt}\end{center}

\subsubsection{Model-Specific Lower Bounds (Detailed
Proofs)}\label{model-specific-lower-bounds-detailed-proofs}

\emph{In-depth bounds for particular computational models - reference as
needed}

\begin{itemize}
\tightlist
\item
  \textbf{Appendix D:} TM/Word-RAM Time Conversions (k-tape Turing
  machines to time bounds)
\item
  \textbf{Appendix G:} Resolution Width-to-Size on Parity/Expander
  (width\(\to\)size bridge for CDCL)
\item
  \textbf{Appendix H:} Algebraic Proof Systems (Overview)
  (Nullstellensatz, Polynomial Calculus)
\item
  \textbf{Appendix I:} RAM/TM to Layered Branching Programs (reduction
  to BP lower bounds)
\item
  \textbf{Appendix K:} Algebraic Proof Systems via Structural Necessity
  (SCL for algebraic models)
\end{itemize}

\begin{center}\rule{0.5\linewidth}{0.5pt}\end{center}

\subsubsection{Foundations \& Verification (Theoretical
Underpinnings)}\label{foundations--verification-theoretical-underpinnings}

\emph{Foundational lemmas and correctness proofs}

\begin{itemize}
\tightlist
\item
  \textbf{Appendix E:} NP as Selection under Constraints (expository)
\item
  \textbf{Appendix L:} Provenance and Recoverability Proofs (seed
  injectivity, closure properties)
\end{itemize}

\begin{center}\rule{0.5\linewidth}{0.5pt}\end{center}

\subsubsection{Reader Resources (FAQ, Glossary,
Supplementary)}\label{reader-resources-faq-glossary-supplementary}

\emph{Non-essential but helpful for navigation and common questions}

\begin{itemize}
\tightlist
\item
  \textbf{Appendix F:} Reviewer FAQ \& Common Pitfalls (anticipated
  objections, clarifications)
\item
  \textbf{Appendix M:} Glossary of Semantic Terms (R\_v, q\_v,
  \(\Phi\)\_v, \(\lambda\), Alt, etc.)
\item
  \textbf{Appendix N:} Oracle and Barrier Analysis (relativization,
  natural proofs, algebrization)
\end{itemize}

\begin{center}\rule{0.5\linewidth}{0.5pt}\end{center}

\subsection{Appendix A: Explicit Overlay
Construction}\label{appendix-a-explicit-overlay-construction}

\textbf{Purpose:} Explicit construction of L* satisfying A1-A5
properties (referenced from §6, §7.2.1). Proves instance is
polynomial-size with no cryptography - enables SCL (§7) and Structural
OWF construction (§9).

\textbf{Key results:} Addressing (A.1), salts (A.2), selector (A.3),
rank forcing (A.4), seed encoding (A.5) \(\to\) A1-A5 invariants (A.6)
\(\to\) polynomial size (A.7).

\subsubsection{A.0 Symbols \& Domains (quick
reference)}\label{a0-symbols--domains-quick-reference}

\begin{itemize}
\tightlist
\item
  \texttt{J\_v} (branch indices at node v): size \texttt{K\_v}
\item
  \(\kappa\) (micro-arity): designated positions per branch; typically
  \(\Theta\)(log n\_core)
\item
  D\_v := K\_v·\(\kappa\) (address domain size at v); U\_v :=
  \{0,...,D\_v\(-\)1\} (disjoint pool for v)
\item
  G\_v := Z\_\{K\_v\} × Z\_\(\kappa\) (product domain for \(\pi\)\_v
  with coordinates (L,R))
\item
  \(\pi\)\_v: J\_v×{[}\(\kappa\){]}\(\to\)U\_v (seed-dependent
  permutation; bijection on size D\_v)
\item
  L\_v\^{}(sel) (selected indices at v): size R\_v; primitive
  coefficient rows a\_\{v,j,\(\ell\)\} \(\in\) \ensuremath{\mathbb{F}}\(_2\)\^{}(m\(_0\))
  (induced by \texttt{Sel\_v}) with constant m\(_0\) = O(1)
\item
  R\_v (emergence at v); \(\Sigma\)\_v R\_v = \(\Theta\)(n\_core·log
  n\_core) under flat sizing
\item
  S(P) (path-gate index set on canonical path P); per-layer slice
  S\_v(P); size \textbar{}S(P)\textbar{} = \(\Theta\)(n/W\_min)
\item
  \texttt{G\_tau} (published set of path-gate items (P,S(P))) with
  bounded reuse (A.9.1)
\end{itemize}

Ranges: K\_v, \(\kappa\) = \(\Theta\)(log n\_core) for QP-sharp;
\textbar{}V\textbar{} = \(\Theta\)(n\_core/log n\_core); all public
objects are poly(\textbar{}x*\textbar{}).

\subsubsection{\texorpdfstring{A.0.V Verifier Audit (checklist for
S(P)/G\_\(\tau\))}{A.0.V Verifier Audit (checklist for S(P)/G\_\textbackslash tau)}}\label{a0v-verifier-audit-checklist-for-spg_ux3c4}

Given instance x* and a purported (P,S(P)) list with per-node tags, the
verifier checks in time O(\textbar{}G\_\(\tau\)\textbar{} +
\(\Sigma\)\_P \textbar{}S(P)\textbar{}):

\begin{itemize}
\tightlist
\item
  Inclusion: S(P) \(\subseteq\) \ensuremath{\bigcup}\_\{u\(\in\)P\} (\{u\}×L\_u\^{}(sel))
  using node/path labels
\item
  Per-path uniqueness: no duplicate \texttt{(v,(j,l))} within a single
  \texttt{S(P)}
\item
  Cross-path bounded reuse: each \texttt{(v,(j,l))} appears in at most
  \texttt{c\ =\ O(1)} paths (tagged subfamily/labels)
\item
  Near-disjointness: implied by bounded reuse (at most
  \texttt{c'\ =\ O(1)} overlaps)
\item
  Two-axis grid per layer: for each v on P, \texttt{S\_v(P)} encodes an
  interval grid \texttt{I\_L\ x\ I\_R} up to \texttt{O(1)} exceptions
  (as published), enabling affine-avoidance used in A.1.F-AA
\end{itemize}

These checks ensure the assignment rules in A.9/A.9.1 hold and that the
prerequisites for address-churn and XOR-cancellation lemmas are met.

\begin{center}\rule{0.5\linewidth}{0.5pt}\end{center}

\subsection{A.1 Addressing: disjoint "scatter" per
node}\label{a1-addressing-disjoint-scatter-per-node}

\textbf{Disjoint pool allocation.} Fix a topological ordering of DAG
nodes V = \{v\_1, v\_2, ..., v\_\textbar{}V\textbar{}\}. For each node
v\_i, let D\_\{v\_i\} := K\_\{v\_i\}·\(\kappa\) be its pool size. Define
global address offset for each node recursively:

offset(v\_1) := 0 offset(v\_i) := offset(v\_\{i-1\}) + D\_\{v\_\{i-1\}\}
for i \textgreater{} 1

Then assign disjoint address pools:

U\_\{v\_i\} := \{offset(v\_i), offset(v\_i)+1, ...,
offset(v\_i)+D\_\{v\_i\}-1\}

\textbf{Lemma A.1.DISJ (Pool disjointness).} For any two distinct nodes
v\_i \(\neq\) v\_j with i \textless{} j in the topological order:

\begin{itemize}
\tightlist
\item
  U\_\{v\_i\} occupies addresses {[}offset(v\_i), offset(v\_i) +
  D\_\{v\_i\})
\item
  U\_\{v\_j\} occupies addresses {[}offset(v\_j), offset(v\_j) +
  D\_\{v\_j\})
\item
  Since offset(v\_j) \(\geq\) offset(v\_i) + D\_\{v\_i\} by
  construction, we have U\_\{v\_i\} \(\cap\) U\_\{v\_j\} = \(\emptyset\)
\end{itemize}

Therefore all node pools \{U\_v : v \(\in\) V\} are pairwise disjoint.
This establishes H1 (disjoint pools) and A1 (Hermeticity) required for
Cartesian factoring (§7 Step 4, Lemma J.1-Cart).

\textbf{Per-node addressing.} For each node v = v\_i with pool U\_v =
\{offset(v), ..., offset(v)+D\_v-1\}, define a \textbf{public
permutation} \(\pi\)\_v: J\_v×{[}\(\kappa\){]}\(\to\)\{0,1,...,D\_v-1\}
(one-to-one) derived from the seed Seed\_v (e.g., a mixed-radix indexing
followed by a seed-dependent Feistel-style permutation over {[}D\_v{]}
with a \textbf{full-churn avalanche property}: changing any bit of
Seed\_v induces \(\Theta\)(D\_v) output changes, ensuring address churn
across segments; any explicitly computable permutation with this
property suffices). Set the \textbf{designated address}

u\_v,\_j,\(\ell\) := offset(v) + \(\pi\)\_v(j,\(\ell\)) \(\in\) U\_v.

Each pair (j,\(\ell\)) maps to a distinct cell in U\_v, and by Lemma
A.1.DISJ, designated addresses for different nodes never overlap. For
how these logical addresses are represented and accessed on
deterministic k-tape TMs (prefix-free record naming, sequential lookup,
and first-use attribution), see Appendix D.5.

\subsubsection{\texorpdfstring{Lemma A.1.\(\Delta\) (Address-churn under
seed
change).}{Lemma A.1.\textbackslash Delta (Address-churn under seed change).}}\label{lemma-a1ux3b4-address-churn-under-seed-change}

For the explicit \(\pi\)\_v family and the canonical path index set
S(P), for any two seed-chains that differ on any unresolved bit along
the chain to v, the symmetric difference between the designated address
multisets \{u\_\{v,j,\(\ell\)\}\}\_\{(v,(j,\(\ell\)))\(\in\)S(P)\} on
the two chains has size \(\Theta\)(\textbar S(P)\textbar). Consequently,
each new cut-gate digest requires evaluating and XORing
\(\Theta\)(\textbar S(P)\textbar) terms anew. See also §A.1.1 and Lemma
A.1.F (Full-churn), which provide the explicit product-Feistel mixer
guaranteeing a constant-fraction movement for any nonzero seed delta.

\emph{Proof.} Fix a node v and path P with its published index set S(P)
\(\subseteq\) \{v\}×L\_v\^{}(sel). Let \(\pi\)\_v\^{}(seed) denote the
permutation instantiated by Seed\_v on domain J\_v×{[}\(\kappa\){]}
(size D\_v = K\_v·\(\kappa\)). Consider two seed-chains that differ on
at least one unresolved bit along the chain to v; by construction (Enc
with GateDigest binding), this implies Seed\_v \(\neq\)
Seed\textquotesingle{}\_v and therefore \(\pi\)\_v\^{}(seed) \(\neq\)
\(\pi\)\_v\^{}(seed\textquotesingle) as permutations.

We use the explicit 4-round product-Feistel \(\pi\)\_v detailed in
§A.1.1 with invertible affine mixers over Z\_\{K\_v\} and
Z\_\{\(\kappa\)\}. For any input index (j,\(\ell\)), write it as (L,R)
with L\(\in\)Z\_\{K\_v\}, R\(\in\)Z\_\{\(\kappa\)\}. Each round applies
an invertible map (Feistel over product groups) whose parameters
(\(\alpha\)\_r,\(\beta\)\_r,\(\gamma\)\_r,\(\delta\)\_r) are derived
deterministically from Seed\_v (length-delimited inside Seed\_v), and
likewise for Seed\textquotesingle{}\_v.

Define \(\Delta\)\(\pi\) := \{ (j,\(\ell\)) \(\in\) S(P) :
\(\pi\)\_v\^{}(seed)(j,\(\ell\)) \(\neq\)
\(\pi\)\_v\^{}(seed\textquotesingle)(j,\(\ell\)) \}. We prove
\textbar{}\(\Delta\)\(\pi\)\textbar{} \(\geq\)
c·\textbar{}S(P)\textbar{} for a universal constant c\(\in\)(0,1).

\begin{enumerate}
\def\labelenumi{\arabic{enumi})}
\item
  Sensitivity of product-Feistel to key changes. For each round r and
  any nonzero change in the round parameters (equivalently, a change in
  a Seed\_v bit), the composed map over two consecutive rounds yields an
  affine transformation whose output depends nontrivially on both
  inputs. Because \(\alpha\)\_r,\(\gamma\)\_r are invertible in their
  respective rings, a change in any of the parameters produces a
  nonconstant affine perturbation on at least one coordinate of (L,R)
  across the entire input domain. Thus, for any fixed input subset
  T\(\subseteq\)J\_v×{[}\(\kappa\){]} that is not concentrated on the
  zero set of a nontrivial affine form, at most a constant fraction of T
  can map to the same output under both keys.
\item
  Structure of S(P). By Appendix A.9 and C.1.1, S(P) is chosen
  layer-wise along P, drawing \(\Theta\)(R\_u) indices from each node u
  on P with mixed j coverage. In particular, for each layer v the
  per-layer slice S\_v(P) contains a two-axis grid I\_L × I\_R
  \(\subset\) Z\_\{K\_v\}×Z\_\{\(\kappa\)\} up to O(1) exceptions
  (Remark A.3.1). By Lemma A.1.F-AA, such grids intersect the zero set
  of any nontrivial affine constraint over Z\_\{K\_v\}×Z\_\{\(\kappa\)\}
  in at most a (1\(-\)\(\alpha\))-fraction plus O(1) points, for a
  universal \(\alpha\) \(\in\) (0,1). Therefore, for any nonzero
  parameter difference between Seed\_v and Seed\textquotesingle{}\_v, at
  least an \(\alpha\)-fraction of S\_v(P) must move under \(\pi\)\_v, up
  to O(1) slack.
\item
  Constant-fraction movement. Combining (1)-(2), there exists a constant
  c\(\in\)(0,1) (independent of n) such that
  \textbar{}\(\Delta\)\(\pi\)\textbar{} \(\geq\)
  c·\textbar{}S(P)\textbar{}. Consequently, the symmetric difference of
  designated address multisets satisfies
  \textbar{}\{\(\pi\)\_v\^{}(seed)(S(P)) \(\Delta\)
  \(\pi\)\_v\^{}(seed\textquotesingle)(S(P))\}\textbar{} \(\geq\)
  c·\textbar{}S(P)\textbar{} = \(\Theta\)(\textbar S(P)\textbar).
\end{enumerate}

Hence, when the seed chain changes (i.e., Seed\_v changes), any cut-gate
digest that aggregates over S(P) must re-evaluate
\(\Theta\)(\textbar S(P)\textbar) terms; cached addresses/values
computed under the old seed chain do not apply to at least a constant
fraction of the needed locations. This proves the
\(\Theta\)(\textbar S(P)\textbar) address-churn claim. \(\blacksquare\)

\textbf{Takeaway (Lemma A.1.\(\Delta\)):} Seed change \(\Rightarrow\)
\(\Theta\)(\textbar S(P)\textbar) designated addresses change.
Pre-cached digest values on old seed chains cannot be reused; fresh
parity computation over the new address set is required.

\subsubsection{\texorpdfstring{A.1.1 Explicit full-churn permutation
\(\pi\)\_v (construction; full-churn
proof)}{A.1.1 Explicit full-churn permutation \textbackslash pi\_v (construction; full-churn proof)}}\label{a11-explicit-full-churn-permutation-ux3c0_v-construction-full-churn-proof}

We give an explicit, efficiently computable bijection \(\pi\)\_v with
the required full-churn property over the product domain
J\_v×{[}\(\kappa\){]} (cardinality D\_v=K\_v·\(\kappa\)), avoiding
cryptography.

Construction (4-round product-Feistel with affine mixers):

\begin{itemize}
\item
  Domain coordinates: write an index as a pair (L,R) with
  L\(\in\)Z\_\{K\_v\}, R\(\in\)Z\_\{\(\kappa\)\}.
\item
  Partition Seed\_v into small blocks to derive round keys
  (\(\alpha\)\_r,\(\beta\)\_r,\(\gamma\)\_r,\(\delta\)\_r) with
  \(\alpha\)\_r\(\in\)Z\_\{K\_v\}\^{}×,
  \(\gamma\)\_r\(\in\)Z\_\{\(\kappa\)\}\^{}× (invertible),
  \(\beta\)\_r\(\in\)Z\_\{K\_v\}, \(\delta\)\_r\(\in\)Z\_\{\(\kappa\)\}.
  Keys are published constants inside Seed\_v.
\item
  Injective key derivation: the extraction map Seed\_v \(\mapsto\)
  ((\(\alpha\)\_r,\(\beta\)\_r,\(\gamma\)\_r,\(\delta\)\_r))\_\{r=1..4\}
  is injective (length-delimited), so Seed\_v \(\neq\)
  Seed\textquotesingle{}\_v implies distinct round-key tuples and hence
  \(\pi\)\_v\^{}(seed) \(\neq\) \(\pi\)\_v\^{}(seed\textquotesingle).
\item
  Define affine round functions F\_r(R):= (\(\alpha\)\_r·R +
  \(\beta\)\_r) mod K\_v, G\_r(L):= (\(\gamma\)\_r·L + \(\delta\)\_r)
  mod \(\kappa\).
\item
  Rounds (Feistel on product group Z\_\{K\_v\}×Z\_\{\(\kappa\)\}):

  \begin{enumerate}
  \def\labelenumi{\arabic{enumi})}
  \tightlist
  \item
    (L\(_1\),R\(_1\)) = (R\(_0\), L\(_0\) + F\_1(R\(_0\)) mod K\_v)
  \item
    (L\(_2\),R\(_2\)) = (R\(_1\), L\(_1\) + F\_2(R\(_1\)) mod K\_v)
  \item
    (L\(_3\),R\(_3\)) = (R\(_2\) + G\_3(L\(_2\)) mod \(\kappa\),
    L\(_2\))
  \item
    (L\(_4\),R\(_4\)) = (R\(_3\) + G\_4(L\(_3\)) mod \(\kappa\),
    L\(_3\))
  \end{enumerate}
\item
  Output address: u = L\(_4\)·\(\kappa\) + R\(_4\) (mixed-radix
  bijection U\_v\(\cong\)Z\_\{K\_v\}×Z\_\{\(\kappa\)\}).
\end{itemize}

This is a permutation for every Seed\_v because Feistel over groups is
invertible and \(\alpha\)\_r,\(\gamma\)\_r are chosen invertible. All
maps are O(1) arithmetic on O(log n) bits and thus polytime.

\textbf{Lemma A.1.IDX (Indexability and inverse).} For fixed public
parameters (K\_v, \(\kappa\)) and a fixed Seed\_v, the map
(j,\(\ell\))\(\mapsto\)\(\pi\)\_v(j,\(\ell\)) and its inverse
u\(\mapsto\)\(\pi\)\_v\^{}(-1)(u) are computable in O(1) word operations
on O(log n)-bit words. Moreover, the round-key derivation from Seed\_v
is injective and computable in O(1) time per round. Consequently,
designated address evaluation u\_\{v,j,\(\ell\)\}=\(\pi\)\_v(j,\(\ell\))
is uniformly efficient and auditable by the verifier.

\emph{Proof.} Each Feistel round consists of a constant number of
modular affine operations over Z\_\{K\_v\} and Z\_\(\kappa\); the
composition of a constant number of such rounds is O(1). Feistel
networks are involutive up to a fixed swap, so inversion uses the same
round functions in reverse order with the same O(1) cost. The
length-delimited, domain-separated extraction of
(\(\alpha\)\_r,\(\beta\)\_r,\(\gamma\)\_r,\(\delta\)\_r) from Seed\_v is
injective and takes O(1) time on the word-RAM/TM baseline for O(log
n)-bit words. \ensuremath{\square}

\textbf{Lemma A.1.F-AA (Affine-avoidance for two-axis schedules).} Fix a
layer v and write the product address domain as G\_v := Z\_\{K\_v\} ×
Z\_\(\kappa\) with coordinates (L,R). Let S\_v(P) \(\subset\) G\_v be
the per-layer slice of S(P) produced by the published two-axis
stratified schedule (Remark A.3.1; Appendix A.9): for each layer,
S\_v(P) contains a grid I\_L × I\_R (Cartesian product of intervals in
the L and R coordinates), possibly together with O(1) exceptional points
from XOR-cancellation.

Then there exists a universal constant \(\alpha\)\(_0\) \(\in\) (0,1)
(e.g., \(\alpha\)\(_0\) = 1/2) and c\(_0\) = O(1) such that for every
nontrivial affine system over G\_v of the form

(A·L + B·R \(\equiv\) u mod K\_v) and/or (C·L + D·R \(\equiv\) w mod
\(\kappa\)),

not both tautologies on G\_v, we have the avoidance bound

\textbar{}S\_v(P) \(\cap\) Fix\textbar{} \(\leq\) (1 \(-\)
\(\alpha\)\(_0\)) · \textbar{}S\_v(P)\textbar{} + c\(_0\),

where Fix denotes the solution set of the system in G\_v.

Proof. Consider first a single nontrivial congruence over Z\_\{K\_v\}:
A·L + B·R \(\equiv\) u (mod K\_v), with (A,B) \(\neq\) (0,0) mod K\_v.
If A \(\neq\) 0 mod K\_v, then for any fixed R the congruence has at
most gcd(A,K\_v) solutions for L in a full residue class modulo K\_v;
within an interval I\_L, the number of solutions is at most
\textbar{}I\_L\textbar{}/K\_v + O(1). Summing over all
\textbar{}I\_R\textbar{} choices of R gives at most
\textbar{}I\_L\textbar{}\textbar{}I\_R\textbar{}/K\_v +
O(\textbar{}I\_R\textbar{}) solutions in I\_L × I\_R. Normalizing by
\textbar{}I\_L\textbar{}\textbar{}I\_R\textbar{} shows the satisfying
fraction is \(\leq\) 1/K\_v + O(1/\textbar{}I\_L\textbar{}). Since K\_v
\(\geq\) 2 and \textbar{}I\_L\textbar{} \(\geq\) 2 by construction, this
is \(\leq\) 1/2 + o(1). Thus at least a 1/2 \(-\) o(1) fraction of I\_L
× I\_R avoids the congruence; the O(1) boundary slack is absorbed into
c\(_0\).

If instead A = 0 mod K\_v and B \(\neq\) 0 mod K\_v, the same argument
applies with roles of L and R swapped against the \(\kappa\)-congruence
case below. An entirely analogous counting argument holds for a single
nontrivial congruence over Z\_\(\kappa\).

For a system of up to two congruences (one mod K\_v and one mod
\(\kappa\)), the solution set inside the grid is bounded by the product
of the per-axis densities plus O(\textbar{}I\_L\textbar{} +
\textbar{}I\_R\textbar{}), yielding an overall satisfying fraction at
most 1/K\_v + 1/\(\kappa\) + o(1) \(\leq\) 1/2 + o(1). Choosing
\(\alpha\)\(_0\) = 1/2 and absorbing the o(1) term into an absolute
c\(_0\) = O(1) gives the claim.

Therefore any nontrivial affine constraint over G\_v excludes at least
an \(\alpha\)\(_0\)-fraction of the grid up to an additive O(1) slack,
as required. \(\blacksquare\)

Remark. We rely only on the two-axis product structure of the per-layer
slice (grid I\_L × I\_R) and do not require GF(2) vector structure. The
published schedule (Remark A.3.1; Appendix A.9) is chosen so that the
verifier can check this property from the instance.

Lemma A.1.F (Full-churn). For the \(\pi\)\_v above, let \(\pi\) and
\(\pi\)' be the permutations induced by two seeds Seed\_v \(\neq\)
Seed'\_v (differing in at least one round parameter). For every
canonical index set S(P) used in Appendix C.1.1 (as constructed in
§A.9), there exists a universal constant c \(\in\) (0,1) (independent of
n) such that

\textbar{}\(\pi\)(S(P)) \(\Delta\) \(\pi\)'(S(P))\textbar{} \(\geq\) c ·
\textbar{}S(P)\textbar{}.

Proof. Write the domain as the direct product group G := Z\_\{K\_v\} ×
Z\_\{\(\kappa\)\}. The 4-round product-Feistel map with affine mixers is
an affine bijection on G at each fixed key: there exist matrices and
offsets over the product ring such that the overall transformation is of
the form

T(L,R) = (A·L + B·R + a, C·L + D·R + b),

where A,D are units in Z\_\{K\_v\},Z\_\{\(\kappa\)\} respectively and
B,C are linear maps between the factors induced by the round structure
(all operations taken componentwise in the product ring). For two
distinct keys, consider the composition H := T'\^{}(-1) \(\circ\) T,
also an affine bijection on G. Points (L,R) fixed by H satisfy

H(L,R) = (L,R) \(\leftrightarrow\) (A\_H\(-\)I)·L + B\_H·R = u (mod
K\_v), C\_H·L + (D\_H\(-\)I)·R = v (mod \(\kappa\)),

for some affine coefficients A\_H,B\_H,C\_H,D\_H and offsets u,v
determined by the two keys. Since Seed\_v \(\neq\) Seed'\_v and the
round mixers (\(\alpha\)\_r,\(\gamma\)\_r) are invertible, at least one
of the four blocks (A\_H\(-\)I), B\_H, C\_H, (D\_H\(-\)I) is nonzero;
otherwise H would be the identity for all inputs, contradicting Seed\_v
\(\neq\) Seed'\_v. Therefore the fixed-point set Fix(H) is contained in
the solution set of a nontrivial affine system of at most two linear
equations over the product group G and thus is an affine subvariety of
codimension \(\geq\) 1.

Counting in the full domain. A nontrivial affine equation over
Z\_\{K\_v\}×Z\_\{\(\kappa\)\} restricts at least one coordinate
nontrivially. Standard counting on product rings yields
\textbar{}Fix(H)\textbar{} \(\leq\) c\(_0\)·(K\_v + \(\kappa\)) for a
universal constant c\(_0\) (e.g., solve one coordinate freely and the
other from a single congruence; a second independent constraint further
reduces the count). Hence, in the full domain of size D\_v =
K\_v·\(\kappa\), the fraction of fixed points is \(\leq\)
c\(_0\)·(1/K\_v + 1/\(\kappa\)).

Transfer to S(P). By Remark A.3.1 (ordering choice) and §A.9 (selector
construction), each per-layer slice S\_v(P) contains a two-axis grid
I\_L × I\_R. By Lemma A.1.F-AA, these grids are not concentrated on the
zero set of any nontrivial affine constraint over G: there exists a
universal \(\alpha\) \(\in\) (0,1) (e.g., \(\alpha\) = 1/2) and an
absolute constant c\(_1\) such that for every nontrivial affine
constraint over G,

\textbar{}S(P) \(\cap\) Zero(affine)\textbar{} \(\leq\) (1 \(-\)
\(\alpha\))·\textbar S(P)\textbar{} + c\(_1\).

Apply this to the affine system defining Fix(H): either we have a single
nontrivial constraint (codimension 1) or two independent ones
(codimension 2). In both cases, by Lemma A.1.F-AA (affine-avoidance),

\textbar{}S(P) \(\cap\) Fix(H)\textbar{} \(\leq\) (1 \(-\)
\(\alpha\))·\textbar S(P)\textbar{} + c\(_1\).

Therefore at least \(\alpha\)·\textbar S(P)\textbar{} \(-\) c\(_1\)
points of S(P) are moved by H, i.e., differ between \(\pi\) and
\(\pi\)'. Since \(\pi\)(S(P)) \(\Delta\) \(\pi\)'(S(P)) contains all
moved images and S(P) has size growing with n, choose n large enough
that c := \(\alpha\)/2 satisfies \(\alpha\)·\textbar S(P)\textbar{}
\(-\) c\(_1\) \(\geq\) c·\textbar{}S(P)\textbar{}. This c is universal
(independent of v and n) because \(\alpha\) is fixed by the S(P)
construction. Hence

\textbar{}\(\pi\)(S(P)) \(\Delta\) \(\pi\)'(S(P))\textbar{} \(\geq\) c ·
\textbar{}S(P)\textbar{},

establishing the full-churn claim. \(\blacksquare\)

Remark. Any constant-round Feistel with invertible affine round
functions and mixed-radix output suffices; the specific constants
(\(\alpha\)\_r,\(\beta\)\_r,\(\gamma\)\_r,\(\delta\)\_r) are derived
deterministically from Seed\_v. No cryptography or randomness is used;
the property is by construction and published with the instance.

\begin{center}\rule{0.5\linewidth}{0.5pt}\end{center}

\subsection{A.2 Salts and primitive
bits}\label{a2-salts-and-primitive-bits}

\begin{itemize}
\item
  \textbf{Salt length.} Fix s=\(\lceil\)c log n\(\rceil\) for a constant
  c\textgreater{}=2.
\item
  \textbf{Salt array.} For each u\(\in\)U\_v, store an explicit
  bitstring \(\sigma\)\_u\(\in\)\{0,1\}\^{}(s\_salt) as fixed published
  constants.
\item
  \textbf{Core fragment.} Let Z\_v(w,x)\(\in\)\{0,1\}\^{}(m\(_0\)) be a
  fixed O(1)-length bit-vector obtained by constant-time projections
  from the base instance/core (e.g., a few witness bits or edge
  indicators).
\item
  \textbf{Primitive check (linear).} For public coefficient rows
  a\_v,\_j,\(\ell\)\(\in\)\{0,1\}\^{}(m\(_0\)) and
  b\_v,\_j,\(\ell\)\(\in\)\{0,1\}\^{}(s\_salt):
\end{itemize}

e\_v,\_j,\(\ell\) := \ensuremath{\langle}a\_v,\_j,\(\ell\), Z\_v(w,x)\ensuremath{\rangle} \(\oplus\)
\ensuremath{\langle}b\_v,\_j,\(\ell\), \(\sigma\)\_u\_v,\_j,\(\ell\)\ensuremath{\rangle} \(\in\) \{0,1\}.

Each e\_v,\_j,\(\ell\) is computed by reading \textbf{one} designated salt
(constant time).

By injectivity of (j,\(\ell\))\(\to\)u\_v,\_j,\(\ell\), the primitives
\{e\_v,\_j,\(\ell\)\} indexed by distinct pairs (j,\(\ell\)) are
well-defined.

\begin{center}\rule{0.5\linewidth}{0.5pt}\end{center}

\subsection{A.3 Selector Sel\_v and branch vector
x\_v}\label{a3-selector-sel_v-and-branch-vector-x_v}

\begin{itemize}
\tightlist
\item
  \textbf{Selector. Let Sel\_v\(\in\)\{0,1\}\^{}(R\_v×(K\_v\(\kappa\)))
  be a row-selector} with \textbf{full row rank R\_v} and \textbf{unit
  row weight}: each row chooses a \textbf{distinct} primitive
  coordinate. Concretely, fix an ordering of L\_v:=J\_v×{[}\(\kappa\){]}
  and let row r of Sel\_v be the unit vector that selects the r-th
  coordinate of e\_v (for r=1,...,R\_v).
\end{itemize}

Remark A.3.1 (Ordering choice). The ordering is chosen to ensure
layer-wise spread across both product coordinates via a two-axis
stratified schedule: for each layer v we publish a grid I\_L × I\_R
\(\subset\) Z\_\{K\_v\} × Z\_\(\kappa\) (with
\textbar{}I\_L\textbar{},\textbar{}I\_R\textbar{} \(\geq\) 2) that forms
the per-layer slice S\_v(P) up to O(1) parity-adjustment exceptions for
XOR-cancellation. This guarantees the affine-avoidance property required
in Lemma A.1.F-AA. Any ordering yielding a verifiable two-axis grid
suffices; a purely diagonal traversal (i, i mod \(\kappa\)) is not
required and is not used for proofs.

\begin{itemize}
\tightlist
\item
  \textbf{Branch vector.} Set x\_v := Sel\_v e\_v \(\in\)
  \{0,1\}\^{}(R\_v). Thus x\_v is literally the list of the R\_v
  selected primitives.
\end{itemize}

Unit-weight rows + disjoint addressing ensure R\_v distinct salts are
read.

\begin{center}\rule{0.5\linewidth}{0.5pt}\end{center}

\subsection{A.4 Completeness Constraint H\_v (rank
forcing)}\label{a4-completeness-constraint-h_v-rank-forcing}

We present two completeness constraint options:

\textbf{Standard constraint (for EO/DP/Resolution).} Pick H\_v=I\_R\_v
(the identity). Then y\_v := H\_vx\_v = x\_v. Identity has row rank
R\_v; the \textbf{Completeness} lemma becomes trivial: to output y\_v
for all inputs, one must learn all R\_v bits of x\_v.

\textbf{Expander parity gate (for order-robust OBDD).} Fix a d-regular
expander Exp\_v on R\_v vertices. Decompose its edges into matchings
M\(_1\),...,M\_d. Define y\_v as edge parities on
\textbf{M\(_1\)\(\cup\)M\(_2\)}: for e=\{u,v\} in this union,
y\_v(e)=x\_v(u)\(\oplus\)x\_v(v). This has size \(\Theta\)(R\_v) and
\textbf{rank R\_v} (constructed with exactly R\_v linearly independent
rows). See Appendix B for details.

Both achieve rank(H\_v) = R\_v, ensuring rank forcing. All other proofs
(Emergence, Dependency) remain unchanged.

Construction note (parity gate rank). While the parity construction
exposes \(\Theta\)(R\_v) equations overall, the published H\_v includes
exactly R\_v independent rows so that rank(H\_v)=R\_v holds as a precise
equality required by Completeness and SCL.

\begin{center}\rule{0.5\linewidth}{0.5pt}\end{center}

\subsection{A.5 Seed-from-completeness Enc and addressing function
F\_overlay}\label{a5-seed-from-completeness-enc-and-addressing-function-f_overlay}

\begin{itemize}
\tightlist
\item
  \textbf{Unified seed formula (DAG-canonical).} For any node v with
  parents P(v):
\end{itemize}

Seed\_v := Enc( v \textbar{}\textbar{} sort(\{(u, Seed\_u, y\_u) : u
\(\in\) P(v)\}) )

where Enc is a fixed, public injective and parseable encoding from
bitstrings to bitstrings. Seeds are variable-length, parseable byte
strings. All node-v addresses are derived from Seed\_v via F\_overlay.

\textbf{Concrete injective Enc example:} For tuple (v, \{(u\_i,
Seed\_\{u\_i\}, y\_\{u\_i\})\}), encode as:

Enc = \textbar{}v\textbar{} : v : \textbar{}k\textbar{} :
(parent\_tuple\_1)...(parent\_tuple\_k)

where parent\_tuple\_i =
\textbar{}u\_i\textbar{}:u\_i:\textbar{}Seed\_ui\textbar{}:Seed\_ui:\textbar{}y\_ui\textbar{}:y\_ui:

\textbf{Notation:}

\begin{itemize}
\tightlist
\item
  \textbar{}x\textbar{} = length of x in decimal
\item
  \textquotesingle{}:\textquotesingle{} = delimiter (not appearing in
  lengths)
\item
  k = number of parents
\end{itemize}

This achieves injectivity through unique parsing.

In chains where \textbar{}P(v)\textbar{}=1, this reduces to
Enc(v\textbar{}\textbar{}Seed\_parent\textbar{}\textbar{}y\_parent).
Different (Seed\_u,y\_u) yield different Seed\_v, preventing free
coalescing.

\begin{itemize}
\item
  \textbf{Addressing.} Define F\_overlay(Seed\_v; j,\(\ell\)) :=
  \(\pi\)\_v(j,\(\ell\)) where \(\pi\)\_v is the public permutation
  seeded by Seed\_v (A.1). This makes all designated addresses for node
  v \textbf{computable} from Seed\_v.
\item
  \textbf{Hermeticity.} F\_overlay and Enc take only their stated
  inputs.
\end{itemize}

\begin{center}\rule{0.5\linewidth}{0.5pt}\end{center}

\subsection{A.6 Proofs of §6
invariants}\label{a6-proofs-of-6-invariants}

\subsubsection{Emergence (Lemma 6.1)}\label{emergence-lemma-61}

The R\_v coordinates selected by Sel\_v are distinct primitives
depending on distinct salts \(\sigma\)\_\{u\_v,\_j,\(\ell\)\}. Conditioned
on Tr\_\{\textless{}v\}, these salts are unread, independent s\_salt-bit
strings. Each primitive is an independent unbiased bit, so x\_v is
uniform on \{0,1\}\^{}(R\_v) with H\_\(\infty\)(x\_v \textbar{}
Tr\_\{\textless{}v\}) = R\_v.

\subsubsection{Completeness (Lemma 6.2)}\label{completeness-lemma-62}

\textbf{Standard constraint:} With H\_v=I\_R\_v, y\_v=x\_v. Any correct
procedure outputting y\_v for all inputs must learn all R\_v coordinates
of x\_v.

\textbf{Expander constraint:} The published parity matrix H\_v is
constructed so that exactly R\_v rows are linearly independent. To
compute all edge parities correctly, one must learn R\_v independent
bits of x\_v.

Both satisfy the rank-forcing property with rank R\_v.

\subsubsection{Dependency}\label{dependency}

Successor node addresses are computed from seeds including
parent\textquotesingle{}s (Seed\_v, y\_v). Before y\_v is fixed,
successor seeds are undefined, preventing look-ahead.

\subsubsection{Injectivity}\label{injectivity}

Enc is injective by construction: different input tuples map to
different encoded strings through unique parsing. Different parent
states produce different Seed\_v, preventing distinct histories from
merging without resolution.

\begin{center}\rule{0.5\linewidth}{0.5pt}\end{center}

\subsection{A.7 Size/budget (flat
profile)}\label{a7-sizebudget-flat-profile}

Using the flat profile (where n denotes the core input length n\_core):

\begin{itemize}
\tightlist
\item
  \textbf{Depth.} D=O(log n\_core).
\item
  \textbf{Node count.} \textbar{}V\textbar{} = O(n/log n) to achieve
  total emergence \(\Theta\)(n log n) with O(log n\_core) per node.
\item
  \textbf{Emergence budget.} \(\Sigma\)\_v R\_v=\(\Theta\)(n log n) with
  R\_v = O(log² n) per node typically.
\item
  \textbf{Micro-arity.} \(\kappa\)=\(\Theta\)(log n) (gives headroom to
  pick R\_v distinct primitives).
\item
  \textbf{Salt requirements.} Each node v needs R\_v salts (one per
  selected primitive), each s = \(\Theta\)(log n) bits.
\end{itemize}

\textbf{Parameter consistency check:}

\begin{itemize}
\tightlist
\item
  Total nodes: \textbar{}V\textbar{} = O(n/log n)
\item
  Average R\_v per node: \(\Theta\)(log² n)
\item
  Total emergence: \textbar{}V\textbar{} × avg(R\_v) = O(n/log n) ×
  \(\Theta\)(log² n) = \(\Theta\)(n log n) {[}YES{]}
\item
  Total salt bits: \(\Sigma\)\_v R\_v × s = \(\Theta\)(n log n) ×
  \(\Theta\)(log n) = \(\Theta\)(n log² n) {[}YES{]}
\item
  DAG depth × width: O(log n\_core) × O(n/log² n) = O(n/log n) =
  \textbar{}V\textbar{} {[}YES{]}
\end{itemize}

\textbf{Metadata per node:}

\begin{itemize}
\tightlist
\item
  \(\pi\)\_v: permutation seed for addressing; O(log n\_core) bits.
\item
  Seed\_v: parseable encoding of (v, sorted parent tuples); O(log n ×
  degree) bits.
\item
  Sel\_v (selector): R\_v indices in {[}K\_v\(\kappa\){]}; total
  R\_v·O(log n\_core) bits.
\item
  H\_v: identity matrix indicator; O(1) bits.
\end{itemize}

\textbf{Total overlay size:} O(n log² n) for salts + O(n log n) for
metadata = poly(n\_core).

Total overlay size is poly(n\_core).

\begin{center}\rule{0.5\linewidth}{0.5pt}\end{center}

\subsection{A.8 Porting to a concrete NP core (CLIQUE / TSP /
3-SAT)}\label{a8-porting-to-a-concrete-np-core-clique--tsp--3-sat}

\begin{itemize}
\item
  \textbf{CLIQUE.} Z\_v(w,x) can be O(1) bits encoding a handful of
  vertex selections; E\_v,\_j,\(\ell\) is then a salted parity of those
  plus \(\sigma\)\_u\_v,\_j,\(\ell\). The base verifier checks the claimed
  k-clique in polynomial time as usual.
\item
  \textbf{TSP.} Let Z\_v(w,x) pull a few edge-indicator bits of the
  tour; E\_v,\_j,\(\ell\) mixes those with the salt; the base verifier
  checks degree/closure and cost in polytime.
\item
  \textbf{3-SAT.} Z\_v(w,x) can be constant many literal-values under w;
  E\_v,\_j,\(\ell\) mixes them with the salt; the base verifier evaluates
  all clauses in polytime.
\end{itemize}

The overlay is agnostic to the core.

\begin{center}\rule{0.5\linewidth}{0.5pt}\end{center}

\subsection{A.9 Selector Construction for Path Gates (XOR
cancellation)}\label{a9-selector-construction-for-path-gates-xor-cancellation}

We formalize the existence of a large index set S(P) along any canonical
bottleneck path P such that the XOR of the Z-selectors cancels, as used
in Appendix C.1.1.

\textbf{Lemma A.9 (Path-wise XOR cancellation).} Let P be a canonical
root\(\to\)sink path that crosses the bottleneck cut C*. For each node v
on P, let L\_v\^{}(sel) index the R\_v selected primitives with
associated primitive coefficient rows a\_\{v,j,\(\ell\)\} \(\in\)
\ensuremath{\mathbb{F}}\(_2\)\^{}(m\(_0\)). Then there exists a subset S(P) \(\subseteq\)
\ensuremath{\bigcup}\emph{\{v\(\in\)P\} (\{v\}×L\_v\^{}(sel)) of size
\textbar{}S(P)\textbar{} = \(\Sigma\)}\{v\(\in\)P\} R\_v \(-\) O(1) such
that \(\oplus\)\emph{\{(v,(j,\(\ell\)))\(\in\) S(P)\} a}\{v,j,\(\ell\)\}
= 0 \(\in\) F\_2\^{}(m\_0).

\emph{Proof.} Let M := \(\Sigma\)\_\{v\(\in\)P\} R\_v be the total
number of selected rows across P, and for each pattern a \(\in\)
\ensuremath{\mathbb{F}}\(_2\)\^{}(m\(_0\)) let cnt(a) denote its multiplicity in the multiset
\ensuremath{\mathcal{A}} := \{a\_\{v,j,\(\ell\)\} : v\(\in\)P,
(j,\(\ell\))\(\in\)L\_v\^{}(sel)\}. Define S(P) by removing from \ensuremath{\mathcal{A}} at
most one representative for every pattern a with cnt(a) odd. There are
at most 2\^{}(m\(_0\)) such patterns, so \textbar{}S(P)\textbar{}
\(\geq\) M \(-\) 2\^{}(m\(_0\)) = M \(-\) O(1). By construction, every
remaining pattern has even multiplicity, hence their XOR is 0 in
\ensuremath{\mathbb{F}}\(_2\)\^{}(m\(_0\)). Each chosen index (v,(j,\(\ell\))) is unique, and
designated addresses remain distinct because (i) within a node the
selected indices are disjoint, and (ii) across nodes the designated
address pools are disjoint. \(\blacksquare\)

\textbf{Corollary A.9.1.} Since m\(_0\) = O(1), \textbar{}S(P)\textbar{}
= \(\Theta\)(\(\Sigma\)\_\{v\(\in\)P\} R\_v) = \(\Theta\)(n/W\_min)
under the flat sizing, matching the weight used in Appendix C.1.1.

Note (independence from witness coverage). The bounded-reuse property
for gate index sets S(P) in this appendix (c\_S = O(1); see Lemma
C.1.1') is independent of the bounded witness-coverage policy for
Z\_v(w,x) on the bottleneck cut (c\_Z = 1; §6.2.4.a). Both constraints
are enforced by construction and hold simultaneously.

\subsubsection{\texorpdfstring{A.9.1 Assignment Rules for G\_\(\tau\)
(Disjointness and Bounded
Reuse)}{A.9.1 Assignment Rules for G\_\textbackslash tau (Disjointness and Bounded Reuse)}}\label{a91-assignment-rules-for-g_ux3c4-disjointness-and-bounded-reuse}

We fix a deterministic assignment of path-gate index sets G\_\(\tau\) =
\{ (P, S(P)) \} satisfying the following enforceable properties (checked
by the verifier):

\begin{enumerate}
\def\labelenumi{\arabic{enumi})}
\item
  Per-path uniqueness. For any fixed path P, S(P) \(\subseteq\)
  \ensuremath{\bigcup}\_\{v\(\in\)P\} (\{v\}×L\_v\^{}(sel)) and contains no duplicate
  indices: within a node v, each pair (j,\(\ell\)) appears at most once
  in S(P).
\item
  Cross-path bounded reuse. There exists a universal constant c = O(1)
  such that any designated primitive index (v,(j,\(\ell\))) appears in
  S(P) for at most c distinct paths P in G\_\(\tau\). Construction:
  partition L\_v\^{}(sel) into constant-many subfamilies keyed by path
  labels; assign each (j,\(\ell\)) to at most one subfamily; the
  parity-adjustment step of Lemma A.9 may cause at most O(1) carry-over
  collisions per \ensuremath{\mathbb{F}}\(_2\)\^{}(m\(_0\)) pattern. Thus total reuse is
  bounded by c.
\item
  Cross-path near-disjointness. For any two distinct paths P \(\neq\)
  P', \textbar{}S(P) \(\cap\) S(P')\textbar{} \(\leq\) c' where c' =
  O(1). This follows from (2) and from the fact that parity cancellation
  in Lemma A.9 removes at most one representative per pattern class
  (2\^{}(m\(_0\)) = O(1) classes).
\item
  Verifier-checkable. The assignment carries path labels and per-node
  subfamily tags in x* so the verifier can confirm (i) S(P)
  \(\subseteq\) \ensuremath{\bigcup}\_\{v\(\in\)P\} (\{v\}×L\_v\^{}(sel)); (ii) no
  duplicates within S(P); and (iii) each (v,(j,\(\ell\))) is used across
  at most c paths.
\end{enumerate}

These rules imply \textbar{}\ensuremath{\bigcup}\emph{\{(P,S(P))\(\in\)G\_\(\tau\)\}
S(P)\textbar{} = \(\Theta\)(\(\sum\)}\{(P,S(P))\(\in\)G\_\(\tau\)\}
\textbar{}S(P)\textbar{}), as used in Lemma C.1.1'.

\textbf{Affine-avoidance (per-node slice).} For each node v on P, the
slice S\_v(P) := S(P) \(\cap\) (\{v\}×L\_v\^{}(sel)) is layer-wise mixed
in j and is not concentrated on the zero set of any nontrivial affine
form over Z\_\{K\_v\}×Z\_\{\(\kappa\)\}: there exist universal constants
\(\alpha\)\(\in\)(0,1) and c\(_1\)=O(1) such that for every nontrivial
affine constraint F, \textbar{}S\_v(P) \(\cap\) \{F=0\}\textbar{}
\(\leq\) (1\(-\)\(\alpha\))·\textbar S\_v(P)\textbar{} + c\(_1\). (Used
in Lemma A.1.F.)

\subsection{A.10 Design alternatives}\label{a10-design-alternatives}

\textbf{Flexibility:} Primitive checks can use AND/OR/XOR. Any full-rank
linear map works for H\_v. The key requirement is maintaining A1-A5
properties (Hermeticity, Injectivity, Emergence, Closure, Dependency).

\subsection{A.11 Summary}\label{a11-summary}

This overlay enforces A1-A5 properties:

\begin{enumerate}
\def\labelenumi{\arabic{enumi}.}
\tightlist
\item
  \textbf{A1 (Hermeticity)}: Disjoint address pools \{U\_v\} \(\to\) no
  hidden channels
\item
  \textbf{A2 (Injectivity)}: Distinct histories \(\to\) distinct seeds
  (Enc injective)
\item
  \textbf{A3 (Emergence)}: R\_v bits unknown until discovered (rank
  forcing)
\item
  \textbf{A4 (Closure)}: Seeds parseable \(\to\) ancestors recoverable
\item
  \textbf{A5 (Dependency)}: Future addresses depend on current outputs
  (DAG structure)
\end{enumerate}

Consequently:

\begin{itemize}
\tightlist
\item
  §7.2.1: A1-A5 \(\to\) SCL (q + \(\Phi\) \(\geq\) R)
\item
  §5: SCL \(\to\) paradigm-specific lower bounds
\item
  §10: Polynomial-size, witness-preserving reduction \(\to\)
  NP-completeness
\end{itemize}

Algorithms must either resolve hidden bits or maintain exponentially
many possibilities.

\begin{center}\rule{0.5\linewidth}{0.5pt}\end{center}

\emph{End of Appendix A.}

\begin{center}\rule{0.5\linewidth}{0.5pt}\end{center}

\subsection{Appendix B: Technical Lemmas for Paradigm
Bounds}\label{appendix-b-technical-lemmas-for-paradigm-bounds}

\textbf{Purpose:} Establish OBDD width and branching-program subfunction
lower bounds \textbf{under the expander-parity gate} (referenced from
§5.1, Appendix A.4). \emph{Order-robustness holds only for this gate}
(no variable order avoids the blow-up). These lemmas specialize the
general paradigm arguments in §5.

\textbf{Why expander-parity:} Expander structure ensures crossing edges
at any cut \(\to\) independent parities \(\to\) 2\^{}(\(\Omega\)(s\_v))
distinguishable subfunctions regardless of variable order.

Notation note (rank exactness). Throughout this appendix,
"expander-parity gate" refers to the published completeness matrix H\_v
constructed with exactly R\_v linearly independent rows
(rank(H\_v)=R\_v). Any \(\Theta\)(R\_v) statements concern construction
size (e.g., number of parity equations available), not rank; we select
an R\_v-row independent subsystem for H\_v.

\emph{Note: General paradigm proofs are in §5. These lemmas handle the
special expander construction.}

\subsubsection{B.1 OBDD cut lemma for the expander-parity gate
(order-robust)}\label{b1-obdd-cut-lemma-for-the-expander-parity-gate-order-robust}

\textbf{Result:} For any variable order \(\pi\), OBDD width \(\geq\)
2\^{}(\(\Omega\)(s\_v)) when s\_v uncommitted bits remain (referenced
from §5.4.1, §7.3.3).

\textbf{Context:} OBDD width equals the number of distinct residual
states at a level. Expander parities force width \(\geq\)
2\^{}(\(\Omega\)(s\_v)) when s\_v = R\_v \(-\) q\_v bits remain
unresolved.

\textbf{Lemma B.1 (OBDD cut \(\Rightarrow\) width, complete).} Fix node
v, the expander-parity gate on M\(_1\)\(\cup\)M\(_2\), and any variable
order \(\pi\). Let U\_v\(\subseteq\){[}R\_v{]} be the set of
s\_v=R\_v-q\_v \textbf{uncommitted} coordinates of x\_v. Along \(\pi\),
there exists a level \(\ell\) with
\textbar{}U\_v\(\cap\)read\_\(\ell\)\textbar{}\(\in\){[}s\_v/3,2s\_v/3{]}.
For this level:

\begin{enumerate}
\def\labelenumi{\arabic{enumi})}
\item
  \textbf{Many crossing parities.} By edge-expansion and the two
  disjoint matchings M\(_1\),M\(_2\), at least k=\(\Omega\)(s\_v) edges
  of M\(_1\)\(\cup\)M\(_2\) have one endpoint in
  U\_v\(\cap\)read\_\(\ell\) and the other in
  U\_v\(\setminus\)read\_\(\ell\). Moreover, a subset of
  k\textquotesingle{}=\(\Omega\)(s\_v) of these edges are
  \textbf{pairwise disjoint} (matching edges).
\item
  \textbf{Independence.} The parity outputs p\_e = x(u)\(\oplus\)x(v)
  for e in this disjoint subset depend on \textbf{disjoint} pairs of
  variables, hence are mutually independent given read\_\(\ell\).
\item
  \textbf{Subfunction explosion \(\Rightarrow\) width.} Consider the
  sub-OBDD below level \(\ell\). For each assignment \(\alpha\) to the
  \textbf{unread} coordinates U\_v\(\setminus\)read\_\(\ell\), let
  F\_\(\alpha\) be the induced \textbf{residual subfunction} on the
  suffix variables. Two assignments \(\alpha\),\(\beta\) that differ on
  the vector (p\_e)\_e yield \textbf{different} residual subfunctions:
  because the gate includes these parities verbatim in y\_v, and
  \textbf{Dependency} keys all successor nodes addressing by Seed\_v :=
  Enc(v \textbar{}\textbar{} sort(\{(u,Seed\_u,y\_u) : u \(\in\)
  P(v)\})), the suffix instance differs on the locations reached. Hence
  there are at least 2\^{}(k\textquotesingle{}) =
  2\^{}(\(\Omega\)(s\_v)) distinct residual subfunctions at level
  \(\ell\). In an OBDD, the number of nodes at a level is at least the
  number of distinct residual subfunctions computed by that level.
  Therefore width\_\(\ell\) \(\geq\) 2\^{}(\(\Omega\)(s\_v)).
\end{enumerate}

\emph{Proof details.} For (1), the expander property with expansion
factor c\textgreater{}0 guarantees that any subset S\(\subseteq\)V with
\textbar{}S\textbar{}\(\leq\)\textbar V\textbar/2 has at least
c\textbar{}S\textbar{} neighbors outside S. Setting
S=U\_v\(\cap\)read\_\(\ell\) with
\textbar{}S\textbar{}\(\approx\)s\_v/2, we get c·s\_v/2 edges crossing.
Since M\(_1\),M\(_2\) are disjoint matchings, at least half belong to
one matching, giving
k\textquotesingle{}\textgreater{}=c·s\_v/4=\(\Omega\)(s\_v) disjoint
crossing edges. For (2), matching edges share no vertices, so their
parities are independent. For (3), different parity vectors lead to
different y\_v values, hence different Seed\_\{succ(v)\} by injectivity
of Enc, causing divergent suffix computations. \(\blacksquare\)

\subsubsection{B.2 BP residual-subfunction lemma (order-free branching
programs)}\label{b2-bp-residual-subfunction-lemma-order-free-branching-programs}

\textbf{Result:} Every branching program has some cut with \(\geq\)
2\^{}(\(\Omega\)(s\_v)) distinct residual subfunctions when s\_v bits
remain unresolved (referenced from §7.3.3).

\textbf{Context:} Branching programs can read variables in any order yet
must still realize 2\^{}(\(\Omega\)(s\_v)) residual subfunctions at some
cut when s\_v bits remain unresolved.

\textbf{Lemma B.2 (BP cut \(\Rightarrow\) subfunctions, complete).}
Under the same expander-parity gate and notation as above, every
(layered) branching program has some cut L during node v with at least
2\^{}(\(\Omega\)(s\_v)) \textbf{distinct residual subfunctions} on
suffix variables.

\emph{Proof.} Traverse the program\textquotesingle{}s read-order and
stop at the first cut L after \(\lfloor\)s\_v/2\(\rfloor\) uncommitted
bits are read. As in Lemma B.1, take
k\textquotesingle{}=\(\Omega\)(s\_v) \textbf{disjoint crossing edges}
from the expander. These edges have one endpoint among the
\(\approx\)s\_v/2 read uncommitted bits and the other among the
\(\approx\)s\_v/2 unread uncommitted bits.

Assignments to the unread half of U\_v toggle the k\textquotesingle{}
independent parities arbitrarily. Specifically, for two choices
\(\alpha\),\(\beta\) of unread bits that differ on the parity vector,
the induced \textbf{suffix instances} (future nodes) differ because y\_v
differs and thus Seed\_\{succ(v)\} differs (Dependency). Therefore the
residual functions F\_\(\alpha\),F\_\(\beta\) on suffix variables are
\textbf{distinct}.

Since we can independently set each of the k\textquotesingle{} parities
to 0 or 1 by choosing appropriate values for the unread bits (using
independence from disjoint edges), we obtain 2\^{}(k\textquotesingle{})
distinct parity vectors, hence 2\^{}(k\textquotesingle{}) distinct
residual subfunctions at cut L. With k\textquotesingle{} =
\(\Omega\)(s\_v), this gives the required 2\^{}(\(\Omega\)(s\_v)) bound.
\(\blacksquare\)

\subsubsection{B.3 Summary}\label{b3-summary}

\textbf{Result:} The expander-parity construction forces exponential
growth in OBDDs (width) and branching programs (residual subfunctions)
when s\_v = R\_v \(-\) q\_v bits remain unresolved. The bounds are
order-robust.

\textbf{Interpretation:} The structure (Emergence, Dependency,
Injectivity) prevents merging unresolved cases; 2\^{}(\(\Omega\)(s\_v))
distinguishable artifacts must be represented, in line with SCL (§7).

\textbf{Used in:} §8 (per-instance deterministic bounds), §9.4
(Structural OWF security proof).

\begin{center}\rule{0.5\linewidth}{0.5pt}\end{center}

\subsection{Appendix C: Achieving Tight Bounds via Frontier-Gate and
Segment
Counting}\label{appendix-c-achieving-tight-bounds-via-frontier-gate-and-segment-counting}

\textbf{Purpose:} Proves tight per-instance deterministic lower bounds
via two complementary mechanisms (referenced from §8.A, §9.4). FG
(Frontier-Gate) provides per-segment cost \(\Omega\)(n/W\_min); SC
(Segment Counting) proves m\_seg \(\geq\) 2\^{}(\(\rho\)-s) segments
required. Together: time \(\geq\)
2\^{}(\(\rho\)-s)·\(\Omega\)(n/W\_min), achieving n\^{}(\(\Theta\)(log
n)) (QP-sharp) or 2\^{}(\(\Theta\)(n)) (flat).

\textbf{Key results:} Theorem C.2.T establishes min-cut \(\lambda\)(A,x)
is canonical and tight up to subpoly factors. Used in Structural OWF
security proof (§9.4) to contradict poly-time inversion.

Notation alignment. This appendix writes cap(C) :=
\(\Sigma\)\_\{v\(\in\)C\}(R\_v\(-\)q\_v) for the residual on a cut C.
For consistency with §2.1, read \(\lambda\)(C) \(\equiv\) cap(C) and
\(\lambda\)(A,x) := min\_C \(\lambda\)(C).

\textbf{Observation:} FG and SC are construction choices that create
L*\textquotesingle{}s hardness properties. The frontier size
(\(\Theta\)(n log n) bits) arises from the \textbf{A1-A5} framework; FG
was added to control when information can be acquired, preventing
front-loading that would weaken bounds.

\subsubsection{C.0 Lane Exhaustiveness (formal
statement)}\label{c0-lane-exhaustiveness-formal-statement}

\textbf{Lemma C.EXH (Lane Exhaustiveness).} For any native deterministic
run on x* (or any randomized run with fixed coins), exactly one of the
following holds:

\begin{itemize}
\tightlist
\item
  (Restart lane) The run does not persist keyed state across changes in
  the resolution prefix on the designated bottleneck cut C*; expected
  number of independent tries satisfies \ensuremath{\mathbb{E}}{[}tries{]} \(\geq\)
  2\^{}(\(\Delta\)(C*)) with \(\Delta\)(C*) = \(\Lambda\)(C*) \(-\)
  (Q(C*) + log\(_2\) Alt(C*)) (Lemma 7.R; see Theorem 7.B and Appendix
  C.4.2), and each try costs \(\geq\) 1/B steps.
\item
  (Single-run lane) The run persists keyed state; the execution
  partitions into non-accepting rollback segments that progress the
  final chain, with m\_seg \(\geq\) 2\^{}(\(\rho\)-s) (Segment Counting;
  Appendix C.2), and, under FG wiring, each such segment costs
  \(\Omega\)(n/W\_min) steps (Appendix C.1.1), yielding time \(\geq\)
  2\^{}(\(\rho\)-s)·\(\Omega\)(n/W\_min).
\end{itemize}

\emph{Proof sketch.} Define the resolution prefix tag on C* by the tuple
(seed-chain tag, ConstraintDigest\_C, WorldCommit\_C) as in Appendix
C.2.a. If the run discards keyed state whenever this tag changes,
attempts are independent and priced by the across-tries inequality
(Theorem J.1, item 2), giving the restart lane bound. Otherwise, keyed
state persists across tag changes; then any suffix after a tag change
must either increase q on C* or grow NF\_C or refute WorldCommit\_C to
progress (CDT/WC), and each non-accepting suffix constitutes a rollback
segment. Segment Counting yields m\_seg \(\geq\) 2\^{}(\(\rho\)-s). With
FG, each segment computes a priced digest of weight
\(\Theta\)(n/W\_min), giving \(\Omega\)(n/W\_min) steps per segment.
Exactly one of these behaviors applies to any fixed run, establishing
exhaustiveness. \ensuremath{\square}

\subsubsection{C.1 Frontier-Gate (FG)
Mechanism}\label{c1-frontier-gate-fg-mechanism}

\textbf{Result:} FG enforces \(\Omega\)(n/W\_min) computational cost per
rollback segment by wiring GateDigest\_v into seeds (referenced from
§8.A, §9.4). Every non-accepting segment that progresses the final chain
must evaluate a identity digest over \(\Theta\)(n/W\_min) seed-dependent
addresses - this cost is unavoidable even with pre-scanned salts (Lemma
5.5.1.c + address churn Lemma A.1.\(\Delta\)).

\textbf{Note:} FG is implemented by wiring GateDigest into seeds, so
progress requires producing a seed-chain-bound digest. This enforces
controlled acquisition of information at designated points.

\textbf{Technical Reality:} The Frontier-Gate is deliberately built into
L*\textquotesingle{}s instance construction, controlling WHEN the
\(\Theta\)(n log n) bits can be acquired. While arity-bounded analysis
shows algorithms must acquire these bits at \(\leq\)B bits/step, we
designed FG into L* to prevent front-loading - creating the tightness we
seek.

\textbf{FG \(\leftrightarrow\) NP-Verification Compatibility:} FG is
\textbf{instance-side} and preserves NP membership: the verifier replays
cut-gate checks in polytime; no model restriction is assumed - only
Receiving-Window Attribution (RWA) and keyedness. Concretely, FG is
wired into seeds via the GateDigest\_v field (§6.2.8); computing Seed\_v
for gate-required nodes mandates producing a gate digest bound to the
current seed chain. This keeps L* in NP while enforcing conservation-law
obligations on any solver.

\textbf{Definitions (digest weight):}

\begin{enumerate}
\def\labelenumi{\arabic{enumi})}
\tightlist
\item
  \textbf{Arity parameter (for per-try baseline only).} For a k-tape TM
  with alphabet \(\Gamma\), let B :=
  k\(\lceil\)log\(_2\)\textbar{}\(\Gamma\)\textbar{}\(\rceil\). Lemma
  5.5.1 bounds fresh information inflow per step by B; we use this only
  for the per-try baseline.
\item
  \textbf{Cut-gate computational weight.} A \textbf{cut-gate proof}
  G\_C(y\_\(\leq\)C) has \textbf{computational weight} w(G\_C) :=
  \textbar{}S(P)\textbar{} = \(\Theta\)(n/W\_min(n)), the number of
  seed-dependent terms that must be evaluated and XORed to compute its
  digest. By Lemma 5.5.1.c, verifying G\_C requires \(\Omega\)(w(G\_C))
  TM operations even if all salts were pre-scanned and cached.
\end{enumerate}

\textbf{FG - the four clauses:}

\begin{itemize}
\tightlist
\item
  \textbf{FG-1 (Local allowance \(\mu\)).} Before the last receiving
  window, the solver may match up to a \(\mu\)-fraction of bits across
  the bottleneck cut C \textbf{without} producing cut-gate proofs.
\item
  \textbf{FG-2 (Global allowance \(\tau\)).} Across the whole cut, the
  solver may pre-learn at most
  \(\Theta\)(\(\tau\)(n)·\(\lambda\)\_base(n)) \textbf{fresh bits}
  without cut-gate proofs. (The calibration \(\tau\)(n) = \(\Theta\)(log
  n/\(\lambda\)\_base(n)) yields QP-sharp tightness.) \emph{(QP-sharp
  link):} With \(\tau\)(n)=\(\Theta\)(log n/\(\lambda\)\_base(n)), the
  remaining residual \(\rho\) - s forces 2\^{}(\(\rho\)-s) segments;
  combined with per-segment probe work \(\Omega\)(n/W\_min) this yields
  the stated quasi-polynomial lower bound in §9.
\item
  \textbf{FG-3 (Gate requirement).} After exhausting FG-1 and FG-2,
  \emph{any further progress along the final chain} must either (i)
  \textbf{reveal a new cut bit}, or (ii) \textbf{evaluate a cut-gate}
  G\_C(y\_\(\leq\)C) whose computational weight is w(G\_C) =
  \(\Theta\)(n/W\_min).
\item
  \textbf{FG-4 (Seed binding).} All artifacts used to certify (i) or
  (ii) are \textbf{bound} to the current seed chain; when the chain
  changes, prior artifacts are invalid (no cross-seed sharing).
\end{itemize}

\textbf{Immediate time pricing (TM computation):} By Lemma 5.5.1.c and
Lemma A.1.\(\Delta\), computing a cut-gate digest with weight
w(G\_C)=\(\Theta\)(n/W\_min) requires \(\Omega\)(w(G\_C)) =
\(\Omega\)(n/W\_min) TM operations, independent of when salts were read
or cached.

\textbf{Lemma (FG time baseline).} Under FG, at the start of the last
segment: (A) \emph{No prior cut-gate proofs} \(\to\) L* limits s
\(\leq\) min\{\(\mu\)·\textbar C\textbar, \(\tau\)(n)·\(\rho\)\};
\textbf{or} (B) \emph{Prior cut-gate proofs} \(\to\) L* already forced
\(\Omega\)(n/W\_min) time per proof

Total time \(\geq\) \(\Omega\)(2\^{}(\(\rho\)-s)·(n/W\_min)) with the s
bound in (A), or time already paid as in (B).

\textbf{Lemma (Per-segment baseline under FG, TM form).} In any
\textbf{single run}, once the \(\mu\) and \(\tau\) allowances are
exhausted, \textbf{every non-accepting rollback segment} that progresses
the final chain must either (i) reveal a new cut bit or (ii) evaluate a
\textbf{cut-gate} G\_C(y\_\(\leq\)C). By \textbf{FG-3}, each such gate
evaluation requires computing a digest over \textbar{}S(P)\textbar{} =
\(\Theta\)(n/W\_min) terms. By Lemma 5.5.1.c and Lemma A.1.\(\Delta\),
this computation costs \(\Omega\)(n/W\_min) TM operations regardless of
when salts were read; for constant k,\textbar{}\(\Gamma\)\textbar, this
is \(\Omega\)(n/W\_min) steps.

FG was designed to manage the same frontier that arity-bounded analysis
measures. By throttling pre-accumulation to
\(\tau\)(n)·\(\lambda\)\_base = \(\Theta\)(log n) bits, FG transforms
what would be loose bounds into tight quasi-polynomial bounds under
QP-sharp profiles.

\textbf{Lemma C.1 (FG Tightness).} The FG mechanism with \(\tau\)(n) =
\(\Theta\)(log n/\(\lambda\)\_base) ensures unconditional tight bounds
for single-run TMs:

\begin{itemize}
\tightlist
\item
  \textbf{With FG wired into seeds (this work):} Instance-side throttle
  caps s \(\leq\) \(\tau\)·\(\rho\) unconditionally, ensuring m\_seg
  \(\geq\) 2\^{}((1-\(\tau\))\(\rho\)) = 2\^{}\(\rho\)/poly(n\_core),
  achieving tight n\^{}(\(\Theta\)(log n\_core)) unconditionally for
  QP-sharp
\end{itemize}

Since FG is built into L* itself (not a model restriction), the bounds
with FG are unconditional - they apply to ALL single-run k-tape TMs
regardless of their strategy.

\textbf{Interpretation:} By incorporating FG into L*\textquotesingle{}s
design, we achieve the parametric spectrum:

\begin{itemize}
\tightlist
\item
  \textbf{QP-sharp} (\(\lambda\)\_base = \(\Theta\)(log² n)): tight
  n\^{}(\(\Theta\)(log n\_core)) bounds
\item
  \textbf{Exponential} (\(\lambda\)\_base = \(\Theta\)(n)): tight
  2\^{}(\(\Theta\)(n)) bounds
\item
  \textbf{\ensuremath{\sqrt{n}} profiles}: tight 2\^{}(\(\Theta\)(\ensuremath{\sqrt{n}})) bounds
\end{itemize}

All via the same arity-bounded core proof, with FG providing the
tightness calibration.

\paragraph{C.1.1 Cut-Gate Proofs: Explicit, Verifier-Checkable
Schema}\label{c11-cut-gate-proofs-explicit-verifier-checkable-schema}

We make the cut-gate mechanism fully explicit so that any claimed
cut-gate check is a deterministic, polynomial-time computation by the NP
verifier.

\textbf{Canonical bottleneck cut.} Fix a designated bottleneck cut C*
(the min-cut under the construction profile). Let \ensuremath{\mathcal{P}} be a canonical
family of root\(\to\)sink paths that cross C* and cover all layers
(e.g., choose W\_min vertex-disjoint paths via a fixed layerwise
matching; any canonical, poly-time computable family suffices).

\textbf{Gate selection (published).} For each path P \(\in\) \ensuremath{\mathcal{P}}, the
instance publishes a fixed index set S(P) consisting of
\(\Theta\)(n/W\_min) designated primitive indices of the form
(v,(j,\(\ell\))) with v on P and (j,\(\ell\)) \(\in\) L\_v\^{}(sel).
S(P) is computed canonically from the published overlay objects (Enc,
F\_overlay, Sel\_v) and consists solely of selected primitives so that
each element maps to a designated address u\_\{v,j,\(\ell\)\} =
F\_overlay(Seed\_v; j,\(\ell\)).

\textbullet{} \textbf{Pre-horizon constraint (acyclicity).} S(P) uses only
primitives from nodes u on P with GREQ\_u = 0 (the published pre-horizon
region). Consequently, all addresses u\_\{u,j,\(\ell\)\} needed to
evaluate G\_\{C*,P\} are defined without requiring any GateDigest.

\textbf{Gate function (published).} The path gate is a parity over
selected primitives. The index set S(P) is chosen so that the XOR of
their coefficient rows cancels the Z-terms: XOR\_\{(v,(j,\(\ell\)))
\(\in\) S(P)\} a\_\{v,j,\(\ell\)\} = 0 (over \ensuremath{\mathbb{F}}\(_2\)\^{}(m\(_0\))). Thus
the gate depends only on salts via the b·\(\sigma\) parts. The verifier
computes the gate value by evaluating

G\_\{C*,P\}(x*,w) := XOR\_\{(v,(j,\(\ell\))) \(\in\) S(P)\}
e\_\{v,j,\(\ell\)\}(x*,w).

(See Appendix A.9 for a formal construction of S(P) achieving XOR
cancellation.)

Here e\_\{v,j,\(\ell\)\} is the Tiny-AND/parity primitive of §6.2.4. No
target bit is published; the verifier recomputes the XOR directly from
(x*,w) by evaluating the designated primitives along S(P).

\textbf{Cut-gate proof object.} A cut-gate proof for path P is the pair

G\_C := (P, S(P)),

with P \(\in\) \ensuremath{\mathcal{P}} and S(P) as published in x*. No additional non-uniform
data is required in a witness; the verifier recomputes
e\_\{v,j,\(\ell\)\} from (x*,w) and uses their XOR to form the digest.

\emph{Construction note.} Across a path P, the multiset of selected rows
\{(a\_\{v,j,\(\ell\)\}) : v on P, (j,\(\ell\)) \(\in\) L\_v\^{}(sel)\}
has total size M = \(\Sigma\)\_\{v\(\in\)P\} R\_v =
\(\Theta\)(n/W\_min). Since a\_\{v,j,\(\ell\)\} \(\in\)
\ensuremath{\mathbb{F}}\(_2\)\^{}(m\(_0\)) with m\(_0\) = O(1), there are at most
2\^{}(m\(_0\))=O(1) distinct patterns. Removing at most one row per
pattern with odd multiplicity yields a subset S(P)' of size M \(-\) O(1)
whose XOR is 0. This preserves \textbar{}S(P)\textbar{} =
\(\Theta\)(n/W\_min). See Appendix A.9 for a formal lemma.

\textbf{Computational weight (formal).} The computational weight of G\_C
is

w(G\_C) := \textbar{}S(P)\textbar{},

the number of designated primitives whose values must be evaluated and
XORed along P to compute the digest. By Lemma 5.5.1.c, any TM requires
\(\Omega\)(w(G\_C)) tape operations to compute the parity, even if all
\(\sigma\)\_\{u\_\{v,j,\(\ell\)\}\} were pre-scanned and cached. By
construction, \textbar{}S(P)\textbar{} = \(\Theta\)(n/W\_min).

\textbf{Seed binding (addresses bound to seed chain).} Since
u\_\{v,j,\(\ell\)\} = F\_overlay(Seed\_v; j,\(\ell\)) and Seed\_v depend
on parent (Seed\_u,y\_u) tuples, the addresses in S(P) are bound to the
current seed chain along P. Any change in the chain (different
unresolved bits) changes the designated addresses and invalidates prior
artifacts - exactly FG-4\textquotesingle{}s requirement.

\textbf{Address churn under rollback (computational).} Because
F\_overlay depends on the current seed chain and \(\pi\)\_v has an
avalanche property, after a rollback the selected address set for S(P)
changes by \(\Theta\)(\textbar S(P)\textbar) in general. Therefore
computing the new digest still requires evaluating and XORing
\(\Theta\)(\textbar S(P)\textbar) terms; by Lemma 5.5.1.c, this costs
\(\Omega\)(\textbar S(P)\textbar) TM operations even if salts were
pre-scanned.

\textbf{Verifier procedure (poly-time).} Given (x*,w) and path P:

\begin{enumerate}
\def\labelenumi{\arabic{enumi})}
\tightlist
\item
  Recompute seeds along P via Enc and parents\textquotesingle{}
  (Seed\_u,y\_u);
\item
  For each (v,(j,\(\ell\))) \(\in\) S(P), compute u\_\{v,j,\(\ell\)\} =
  F\_overlay(Seed\_v; j,\(\ell\)) and read
  \(\sigma\)\_\{u\_\{v,j,\(\ell\)\}\};
\item
  Compute e\_\{v,j,\(\ell\)\} = \ensuremath{\langle}a\_\{v,j,\(\ell\)\}, Z\_v(w,x)\ensuremath{\rangle}
  \(\oplus\) \ensuremath{\langle}b\_\{v,j,\(\ell\)\}, \(\sigma\)\_\{u\_\{v,j,\(\ell\)\}\}\ensuremath{\rangle};
\item
  XOR them to obtain G\_\{C*,P\}(x*,w).
\end{enumerate}

This costs O(\textbar{}S(P)\textbar{}) designated reads and
O(\textbar{}S(P)\textbar{}) time beyond seed recomputation, i.e.,
poly(\textbar{}x*\textbar{}).

\textbf{Gate horizon budget (structural; auditable).} The \(\tau\)-cap
is encoded by the published GREQ map (Section 6.2.10): across the
bottleneck cut, nodes with GREQ=0 contribute at most
S(n)=\(\Theta\)(\(\tau\)(n)·\(\lambda\)\_base(n)) bits. This
structurally caps pre-final agreement s. Any claimed cut-gate proofs are
recomputable by the verifier using the schema above.

\textbf{Corollary C.1.1 (Per-segment baseline; mandatory).} With
GateDigest\_v wired into seeds for all v in the published gate horizon
(GREQ\_v=1), any non-accepting rollback segment that progresses the
final chain must produce at least one fresh GateDigest on the current
seed chain. By the schema above and Lemma 5.5.1.c + Lemma
A.1.\(\Delta\), computing this digest requires \(\Omega\)(n/W\_min) TM
operations regardless of when salts were read, thus each segment costs
\(\Omega\)(n/W\_min) TM steps.

\textbf{Lemma C.1.1' (Bounded reuse across G\_\(\tau\)).} In the
published cut-gate schema G\_\(\tau\) = \{ (P, S(P)) \}, each designated
primitive index (v,(j,\(\ell\))) appears in at most c entries across
G\_\(\tau\), where c = O(1) is a fixed construction constant.
Consequently,

\textbar{}\ensuremath{\bigcup}\emph{\{(P,S(P))\(\in\)G\_\(\tau\)\} S(P)\textbar{} =
\(\Theta\)(\(\sum\)}\{(P,S(P))\(\in\)G\_\(\tau\)\}
\textbar{}S(P)\textbar{}).

\emph{Proof (via Appendix A.9).} The schema assigns indices by path and
depth with explicit disjointness rules (see Appendix A.9): (i) along
distinct published paths, selected index multisets are disjoint except
for at most O(1) parity-adjustment collisions per pattern; (ii) per node
v, selector families L\_v\^{}(sel) are partitioned by (j,\(\ell\)) so a
given pair is assigned to at most one path at a time; and (iii) the
XOR-cancellation procedure removes at most one row per constant-size
pattern class (O(1) patterns since m\(_0\) = O(1)). Thus each index can
be reused across at most c = O(1) entries in G\_\(\tau\), proving the
\(\Theta\)(·) relation. The verifier enforces these assignment rules
when checking (P, S(P)). \(\blacksquare\)

\textbf{Lemma C.1.2 (Unpredictability without reads).} If an algorithm
computes G\_\{C*,P\} while skipping any salt in S(P), then there exist
two completions of the unread salt values that are both consistent with
the transcript but flip the parity of G\_\{C*,P\}. Hence the
algorithm\textquotesingle{}s output would be wrong on one completion.

\emph{Proof.} By selector cancellation,
\(\oplus\)\emph{\{(v,(j,\(\ell\))) \(\in\) S(P)\} a}\{v,j,\(\ell\)\} =
0, so the gate depends only on salts. For any unread
\(\sigma\)\_\{u\_\{v,j,\(\ell\)\}\}, the two values \(\sigma\)=0 and
\(\sigma\)=1 are both consistent with the transcript but yield different
parities. Since addresses depend on the current seed chain
(seed-binding), the algorithm cannot predict which cells to read without
computing the chain. \ensuremath{\square}

\textbf{Lemma C.1.3 (No asymptotic benefit from pre-scanning).} While a
TM can pre-scan all published salts in O(\(\Sigma\)\_v R\_v) = O(n log
n) time at the start, this does not reduce the asymptotic complexity. By
Lemma 5.5.1.c, computing each gate digest GateDigest\_v still requires
\(\Omega\)(\textbar S(P)\textbar) = \(\Omega\)(n/W\_min) operations per
segment, even with all salts cached. The seed-dependent addressing
(\(\pi\)\_v with avalanche property) ensures that each new seed chain
requires computing a fresh set of \(\Theta\)(n/W\_min) addresses and
XORing the corresponding cached values. Thus pre-scanning shifts work
timing but does not eliminate the \(\Omega\)(n/W\_min) computational
cost per segment on the final chain.

\paragraph{C.1.4 No-Overlay Bypass (formal proof for Theorem
10.4.1-BYP)}\label{c14-no-overlay-bypass-formal-proof-for-theorem-1041-byp}

We prove that any algorithm that outputs a canonical witness W = (w,
G\_\(\tau\), Dig\_\(\tau\)) must realize, on the native run, the listed
overlay computations beyond the gate horizon; consequently, there is no
asymptotically cheaper "compute w, then synthesize Dig\_\(\tau\)"
strategy.

Setup. Fix an instance x* \(\in\) L* constructed with: Seed-Locked
Decode (\(\Phi\)\textasciitilde{}) for CNF decoding; disjoint designated address pools
\{U\_v\}; injective, parseable Enc; Keyedness; Hermeticity; RWA;
Completeness (rank(H\_v)=R\_v); and Frontier-Gate wiring with GREQ on
the published gate horizon (§6.2.8). Let \(\pi\) be any accepting run
that outputs W = (w, G\_\(\tau\), Dig\_\(\tau\)).

Step 1 (Verifier contract). By §10.4.1, the canonical verifier
recomputes each (P, S(P)) \(\in\) G\_\(\tau\) by regenerating seeds
along P and, for each (v,(j,\(\ell\))) \(\in\) S(P), computing the
seed-bound address u\_\{v,j,\(\ell\)\} and evaluating the designated
primitive e\_\{v,j,\(\ell\)\} before XORing to form the digest (Appendix
C.1.1). Acceptance requires the produced digests match the recomputed
values.

Step 2 (Necessity of evaluating all listed primitives). By Lemma C.1.2
(Unpredictability without reads) and Completeness, skipping any
designated term in S(P) admits two completions of the skipped payload
bits that keep the transcript fixed but flip the parity of the digest,
causing a verification failure. Hence to ensure acceptance, the run must
(functionally) determine every term in each listed S(P).

Step 3 (Seed-bound addressing forbids write-in shortcuts). Injective
Enc, Keyedness, and disjoint \{U\_v\} ensure that altering unresolved
coordinates changes downstream seeds and addresses, and artifacts keyed
by different seeds cannot merge. Therefore, computing a correct digest
for the published P and S(P) requires evaluating exactly the designated
payloads named by the current seed chain; any alternative path would
produce addresses inconsistent with the verifier\textquotesingle{}s
recomputation.

Step 4 (RWA credit for first-use). By Hermeticity and RWA, determining
each designated term\textquotesingle{}s value for the first time is
credited as a first-use designated read attributable to its node v,
independent of prefetch order (Lemma 6.1-RWA; §D.5). Across all entries
of G\_\(\tau\), the total number of such first-use credits is
\(\Omega\)(\textbar\ensuremath{\bigcup}\emph{\{(P,S(P))\(\in\)G\_\(\tau\)\}
S(P)\textbar{}). Moreover, by the bounded-reuse property of the
published schema (Lemma C.1.1'), we have \textbar{}\ensuremath{\bigcup} S(P)\textbar{} =
\(\Theta\)(\(\sum\)}\{(P,S(P))\(\in\)G\_\(\tau\)\}
\textbar{}S(P)\textbar{}).

Step 5 (Steps from first-use credits). By Lemma D.2.1 (RWA first-use
reads \(\to\) TM steps), if the run accrues F such first-use credits, it
takes at least F/B steps on a k-tape TM (B :=
k·\(\lceil\)log\(_2\)\textbar{}\(\Gamma\)\textbar{}\(\rceil\)),
independent of scheduling or caching.

Conclusion. Combining Steps 2-5, any accepting run that outputs W must
realize the overlay computations beyond the gate horizon and pay
\(\Omega\)(\textbar\ensuremath{\bigcup} S(P)\textbar{}) first-use credits, implying a step
lower bound of \(\Omega\)(\textbar\ensuremath{\bigcup} S(P)\textbar{}/B). By Lemma C.1.1',
this equals \(\Omega\)(\(\sum\)\textbar S(P)\textbar/B) up to constants.
This rules out an asymptotically cheaper "post-hoc synthesis" from w
alone and establishes Theorem 10.4.1-BYP. \(\blacksquare\)

Scope note (native runs vs post-hoc). This appendix establishes
necessity on native runs that produce W from x* alone. It does not
assert a super-polynomial cost to assemble W when w is provided as extra
input; §10.4.1 explicitly gives a polynomial-time construction of W from
(x*, w), consistent with NP-membership.

\subsubsection{C.2 Segment Counting (SC)}\label{c2-segment-counting-sc}

\textbf{Result:} SC proves m\_seg \(\geq\) 2\^{}(\(\rho\)-s) rollback
segments required to reach acceptance (Lemma C.2, Theorem C.2.T).
Combined with FG\textquotesingle{}s per-segment cost, yields time
\(\geq\) 2\^{}(\(\rho\)-s)·\(\Omega\)(n/W\_min).

\textbf{Structure:} C.2.a defines CDT/WC mechanisms (segment boundaries,
pricing), \textbf{Lemma WC-1} ("+1 only" refutation bound) and
\textbf{Lemma C.2.3} (one extra elimination per segment) establish the
per-segment world reduction, \textbf{Lemma C.2} proves main counting
lemma m\_seg \(\geq\) 2\^{}(\(\rho\)-s), and C.2.T establishes tightness
and min-cut canonicality. SC exposes the multiplicative structure: when
worlds remain indistinguishable, segments expand by a factor \(\approx\)
2\^{}(\(\rho\)-s).

\paragraph{C.2.a Constraint-Digest Tagging (CDT) and World-Commit
(WC)}\label{c2a-constraint-digest-tagging-cdt-and-world-commit-wc}

\textbf{Achievement:} CDT/WC mechanisms ensure (1) segment boundaries
are well-defined via ConstraintDigest\_C and WorldCommit\_C tags, (2)
each non-accepting segment pays designated work (\(\Omega\)(n/W\_min)
per CDT-3), (3) acceptance requires uniqueness
\textbar{}\ensuremath{\mathcal{U}}(\(\pi\))\textbar{} \(\leq\) 1 (ACC-1). Together these enable
the m\_seg \(\geq\) 2\^{}(\(\rho\)-s) segment count (Lemma C.2).

We strengthen the rollback key so that semantic pruning across the
bottleneck cut C is reflected in the tag and properly priced.

\begin{itemize}
\tightlist
\item
  Constraint set over C at prefix \(\pi\): Cons\_C(\(\pi\)) := all
  semantic consequences over cut bits that are derivable from the
  transcript prefix \(\pi\) (public overlay + revealed designated bits),
  restricted to the published, NP-verifiable constraint family (the same
  family used by cut-gates; Appendix C.1.1).
\item
  Canonical normal form: NF\_C(·) maps any finite constraint set over C
  to a unique, duplicate-free, order-invariant basis that forms a
  conservative summary of cut-feasibility constraints
  (over-approximating the feasible-worlds set). The verifier recomputes
  NF\_C in polynomial time from \(\pi\).
\item
  ConstraintDigest\_C(\(\pi\)):= H(NF\_C(Cons\_C(\(\pi\)))). This digest
  is included in the resolution prefix across C.
\item
  WorldCommit\_C(\(\pi\)):= H(CommitSelector(\(\pi\))\textbar\_C), a
  canonical choice of a currently feasible cut-world under
  NF\_C(Cons\_C(\(\pi\))) and revealed bits. H is used only as a compact
  representation; soundness relies on the verifier\textquotesingle{}s
  deterministic recomputation of \(\omega\)* :=
  CommitSelector(\(\pi\))\textbar\_C from NF\_C and q, not on collision
  resistance. Any certificate is checked against this recomputed
  \(\omega\)*. For non-accepting segments we require that, if no new cut
  information is added, the segment must end by refuting this committed
  world (verifier-checkable).
\end{itemize}

Required properties (design axioms):

\begin{itemize}
\tightlist
\item
  CDT-1 (Sufficient summary). For every \(\pi\), NF\_C(Cons\_C(\(\pi\)))
  together with fixed cut bits from q forms a sufficient summary of
  cut-level information such that any suffix of the run that does not
  change NF\_C and does not increase q on C cannot reduce the
  feasible-worlds set across C. (We do not require NF\_C to enumerate
  all logical consequences; only that it is strong enough to block
  read-free pruning within a segment.)
\item
  CDT-2 (Monotone minimality). If NF\_C grows between \(\pi\) and
  \(\pi\)\textquotesingle, at least one previously feasible world is
  excluded (logical strengthening). Conversely, if NF\_C is unchanged
  and q on C is unchanged, feasible worlds across C are unchanged.
\item
  CDT-3 (Charge/cost coupling). Any growth of NF\_C across C or any
  refutation of WorldCommit\_C must be backed by designated work: it
  requires producing a verified cut-gate object or a published cut-gate
  digest (or published equivalent) of weight \(\Theta\)(n/W\_min) on the
  \textbf{current seed chain}, costing \(\Omega\)(n/W\_min) TM steps
  (per §C.1.1 and Lemma 5.5.1.c).
\item
  WC-1 (Refutation bound). If NF\_C and q on C do not change within a
  segment, at most one additional world (the committed one) may be
  refuted at the segment\textquotesingle{}s terminal action. Any other
  refutation would either change NF\_C or q and thus end the segment
  earlier.
\item
  ACC-1 (Acceptance purity; alias: Lemma C.2.ACC-mech). The accept
  action does not introduce new cut-level information: it neither
  increases q on C nor grows NF\_C. Therefore, at the start of the
  accepting segment, the feasible-worlds set across C has size at most 1
  (otherwise acceptance would require new designated artifacts,
  contradicting purity; operationally, \(\leq\)1 world at the start of
  the accepting segment). (See also Lemma C.2.ACC-logical for a logical
  derivation of ACC-1 and Lemma C.2.ACC for the digest-binding
  implication.)
\end{itemize}

\textbf{Lemma C.2.ACC-logical (ACC-1 as a logical consequence).} Fix any
instance x* with the published FG wiring and canonical witness schema,
and a transcript prefix \(\pi\) at the beginning of the accepting
segment. Suppose \textbar{}\ensuremath{\mathcal{U}}(\(\pi\))\textbar{} \(\geq\) 2 across the
designated bottleneck cut C (i.e., there are at least two cut-worlds
consistent with NF\_C(Cons\_C(\(\pi\))) and the revealed \(\Delta\)q on
C). Then no run can output a valid canonical witness W = (w,
G\_\(\tau\), Dig\_\(\tau\)) that passes Algorithm V on x* without either
(i) increasing \(\Delta\)q on C (functionally determining additional cut
bits) or (ii) producing a new priced digest (which ends the segment).
Hence, acceptance with \textbar{}\ensuremath{\mathcal{U}}(\(\pi\))\textbar{} \(\geq\) 2 is
impossible; therefore ACC-1 holds as a logical consequence of
correctness requirements.

\emph{Proof.} Under Hermeticity (A1), the transcript is the only source
of fresh information: cut-level consequences can arise only via (i)
first-use designated reads (counted in \(\Delta\)q) or (ii) priced gate
computations (Appendix C.1.1), which, by CDT-1', are precisely the
events that can grow NF\_C or reduce \textbar{}\ensuremath{\mathcal{U}}\textbar{} within a
segment. By Realizability (H4), every world in \ensuremath{\mathcal{U}}(\(\pi\)) remains
feasible absent new designated artifacts. If
\textbar{}\ensuremath{\mathcal{U}}(\(\pi\))\textbar{} \(\geq\) 2 at the start of the accepting
segment and neither \(\Delta\)q increases nor a priced digest is
produced during the segment, then no new evidence distinguishes the
worlds; any choice among them is a guess. But the canonical witness
includes the vector of digests Dig\_\(\tau\) bound to the current seed
chain. By Lemma C.1.2 (Unpredictability) and address-churn (Lemma
A.1.\(\Delta\)), for any incorrect choice of the cut-world the
recomputed digests on the current chain disagree on at least one
published path P; Algorithm V deterministically recomputes each digest
and rejects. Thus, without new designated evidence (\(\Delta\)q or a
priced digest), a correct witness cannot be produced while
\textbar{}\ensuremath{\mathcal{U}}(\(\pi\))\textbar{} \(\geq\) 2. Therefore acceptance requires
\textbar{}\ensuremath{\mathcal{U}}(\(\pi\))\textbar{} \(\leq\) 1 at the start of the accepting
segment. \ensuremath{\square}

Constraint family and normalization (scope over C).

\begin{itemize}
\item
  We restrict cut-level constraints admitted into NF\_C to two canonical
  forms:

  \begin{enumerate}
  \def\labelenumi{\arabic{enumi})}
  \tightlist
  \item
    Bit determinations (unit equalities) that increase q on C; and
  \item
    UnitRefute(\(\omega\)): a unit exclusion of a single cut-world
    \(\omega\) (fixed assignment to unresolved cut bits). General
    overlay/gate constraints may exist globally but are not admitted
    into NF\_C unless they reduce to (1) or (2) over C.
  \end{enumerate}
\item
  Normalization NF\_C eliminates duplicates and keeps only these
  canonical forms with a deterministic order.
\end{itemize}

\textbf{Algorithm NF\_C (Explicit Normalization).}

Input: Cons\_C(\(\pi\)) - finite constraint set over bottleneck cut C
derivable from prefix \(\pi\) Output: NF\_C(Cons\_C(\(\pi\))) -
canonical normal form (unique, duplicate-free, order-invariant basis)

Procedure:

\begin{enumerate}
\def\labelenumi{\arabic{enumi})}
\item
  Initialize two sets: BitDet := \(\emptyset\), Refuted := \(\emptyset\)
\item
  Filter to canonical forms: For each constraint c \(\in\)
  Cons\_C(\(\pi\)): a) If c is a unit equality (v\_i = b for v\_i
  \(\in\) C, b \(\in\) \{0,1\}): Add (i, b) to BitDet // bit
  determination b) If c is a canonical UnitRefute(\(\omega\)) over C
  (unit exclusion of a single cut-world \(\omega\)): Add w to Refuted //
  w is a fixed assignment to unresolved cut bits c) Otherwise: discard c
  (not admitted to NF\_C)
\item
  Remove duplicates: BitDet := unique(BitDet) // remove duplicate bit
  determinations Refuted := unique(Refuted) // remove duplicate world
  refutations
\item
  Check consistency (polynomial-time satisfiability): Let w\_partial :=
  assignment from BitDet If w\_partial is self-contradictory: return
  \(\perp\) (inconsistent) Remove from Refuted any w that contradicts
  w\_partial (already excluded)
\item
  Canonical ordering: Sort BitDet lexicographically by (index i, bit b)
  Sort Refuted lexicographically by world assignment
\item
  Output canonical basis: return NF := (BitDet\_sorted, Refuted\_sorted)
\end{enumerate}

Complexity: O(\textbar{}Cons\_C\textbar{}) filtering +
O(\textbar{}Cons\_C\textbar{} log \textbar{}Cons\_C\textbar{}) sorting =
O(\textbar{}Cons\_C\textbar{} log \textbar{}Cons\_C\textbar{}) =
polynomial

Properties:

\begin{itemize}
\tightlist
\item
  Unique: lexicographic ordering ensures uniqueness
\item
  Duplicate-free: step 3 eliminates duplicates
\item
  Order-invariant: output independent of input order (canonical sort)
\item
  Polynomial-time: all operations poly in \textbar{}Cons\_C\textbar{}
\item
  Conservative summary: the feasible-worlds set under NF\_C(Cons)
  \(\cup\) \{q\} contains the true feasible-worlds set under Cons;
  discarding non-canonical consequences only weakens pruning and
  strengthens the m-bound.
\end{itemize}

Verifier usage (Appendix C.2.b): The NP verifier recomputes
NF\_C(Cons\_C(\(\pi\))) by:

\begin{itemize}
\tightlist
\item
  Collecting constraints from transcripted designated reads
\item
  Collecting constraints from verified gate objects
\item
  Running Algorithm NF\_C above
\item
  Computing ConstraintDigest\_C := H(NF\_C(Cons\_C(\(\pi\))))
\item
  Checking consistency with resolution prefix
\end{itemize}

All steps polynomial-time; normalization deterministic. \(\blacksquare\)

\textbf{Lemma CDT-1\textquotesingle{} (No unbacked cut consequences).}
For the published constraint family and normalization NF\_C defined
above, if a cut-level consequence becomes derivable from transcript
prefix \(\pi\) without additional designated artifacts being revealed or
produced, then it must be either (i) a bit determination already
reflected as \(\Delta\)q on C, or (ii) a UnitRefute(\(\omega\)) for some
world \(\omega\) already present in NF\_C. Hence NF\_C can only grow
when (i) \(\Delta\)q on C increases, or (ii) a new designated artifact
is produced (gate/digest) that supports a new UnitRefute or bit
determination, implying a tag change and designated cost.

\emph{Proof.} We prove this by analyzing the three sources from which
constraints in NF\_C can arise:

\textbf{Source 1: Overlay-only constraints.} The published overlay (G,
H\_v, Sel\_v, Enc, F\_overlay, GREQ, published parameters) is fixed
before the run begins and is a deterministic function of the base
instance \(\varphi\) (see §10 reduction). All constraints derivable
purely from the overlay structure - without reference to any designated
salt values or algorithm-computed outputs - are constant across all
possible runs and are captured in the initial constraint set
NF\_C(\(\emptyset\)) at prefix \(\pi\) = \(\emptyset\). By
Closure/Recoverability and Realizability (Appendix J; Lemma
"Closure/Recoverability"), any cut-world consistent with
NF\_C(\(\emptyset\)) admits a transcript-consistent completion absent
new designated artifacts; hence overlay-only reasoning cannot further
shrink the feasible set mid-run.

Since these overlay consequences are already present in NF\_C from the
start, they cannot constitute "new" entries that grow NF\_C during the
run. {[}YES{]}

\textbf{Source 2: Designated salt revelations.} When the algorithm reads
a designated cell \(\sigma\)\_u for the first time (a first-use read
under RWA), this reveals one primitive value e\_\{v,j,\(\ell\)\} =
\ensuremath{\langle}a\_\{v,j,\(\ell\)\}, Z\_v(w,x)\ensuremath{\rangle} \(\oplus\) \ensuremath{\langle}b\_\{v,j,\(\ell\)\},
\(\sigma\)\_u\ensuremath{\rangle}. By the construction in Appendix A.2, each such read
reveals exactly one bit of information about the selected coordinates on
cut C (via the full-rank selector Sel\_v).

The verifier tracks these revelations via RWA (Lemma 5.5.1.b) and
credits them as functional determinations \(\Delta\)q on C. Each unit of
\(\Delta\)q corresponds to fixing one coordinate value on C, which may
imply:

\begin{itemize}
\tightlist
\item
  Adding one unit equality (bit determination) to NF\_C, OR
\item
  Resolving a prior constraint to exclude specific worlds (captured as
  \(\Delta\)q effect)
\end{itemize}

Since the verifier explicitly counts \(\Delta\)q and the normalization
NF\_C only admits unit equalities from designated reads, any constraint
from this source is already accounted for in the \(\Delta\)q term. New
designated reads \(\to\) \(\Delta\)q increases \(\to\) already priced.
{[}YES{]}

\textbf{Source 3: Algorithm-computed gate digests.} When GREQ\_v = 1,
the algorithm must compute GateDigest\_v = Enc\_gate(P(v)
\textbar{}\textbar{}
XOR\_\{(v\textquotesingle{},(j,\(\ell\)))\(\in\)S(P(v))\}
e\_\{v\textquotesingle{},j,\(\ell\)\}). This digest is a deterministic
function of:

\begin{itemize}
\tightlist
\item
  The seed chain (Seed\_u, y\_u along ancestors)
\item
  The designated salt values \(\sigma\)\_u for u \(\in\) designated
  addresses in S(P(v))
\item
  The overlay functions (F\_overlay, Enc\_gate)
\end{itemize}

The gate digest itself does not directly enter NF\_C as a constraint
unless it \textbf{implies a cut-level consequence}:

\textbf{Case 3a:} GateDigest\_v determines Seed\_child for a child node
on cut C, which via Dependency propagates to determine y\_child
coordinates. This determination is captured as \(\Delta\)q on C (via the
full-rank H\_child and Sel\_child). Already priced via \(\Delta\)q.
{[}YES{]}

\textbf{Case 3b:} GateDigest\_v, combined with existing NF\_C and
revealed bits, proves that a specific cut-world assignment w is
\textbf{inconsistent}. Canonical mechanisms include:

\begin{itemize}
\tightlist
\item
  Parity contradictions: for a published gate parity over S(P(v)), the
  digest fixes
  \(\oplus\)\emph{\{(v\textquotesingle{},j,\(\ell\))\(\in\)S(P(v))\}
  e}\{v\textquotesingle{},j,\(\ell\)\}; any w whose induced cut-bits
  force the opposite parity is infeasible.
\item
  Rank/linear constraints: with rank(H\_u)=R\_u, linear consequences
  from gate outputs imply linear relations on cut bits; any w violating
  these relations is infeasible.
\item
  Seed/consistency propagation: a digest fixing a child seed/value
  forces a cut assignment incompatible with w. The verifier checks such
  a proof object and admits it as UnitRefute(\(\omega\)) into NF\_C.
\end{itemize}

Crucially, producing such a proof requires:

\begin{itemize}
\tightlist
\item
  The algorithm to have computed GateDigest\_v (designated work
  \(\Omega\)(\textbar S(P)\textbar) by Lemma 5.5.1.c)
\item
  The digest to reference designated artifacts (the XOR over S(P))
\item
  The verifier to check the proof (polynomial time, but the algorithm
  must exhibit it)
\end{itemize}

If GateDigest\_v is \textbf{fresh} (differs from the previous value on
the current seed chain), then by Lemma C.2.2 the seed-chain tag changes,
ending the segment and pricing the digest computation. If GateDigest\_v
is unchanged but leads to a new UnitRefute(\(\omega\)), this refutation
updates NF\_C \(\to\) ConstraintDigest\_C changes \(\to\) tag changes
\(\to\) segment ends. {[}YES{]}

\textbf{Case 3c:} GateDigest\_v is unchanged and implies no new
consequences beyond what \(\Delta\)q and prior NF\_C already capture.
Then NF\_C does not grow from this digest. {[}YES{]}

\textbf{Synthesis:} Every path by which NF\_C can grow requires either:

\begin{itemize}
\tightlist
\item
  New designated artifacts (Source 2 or 3) \(\to\) priced via
  \(\Delta\)q or gate computation
\item
  Consequences already in NF\_C (Sources 1 or prior entries) \(\to\) not
  new growth
\end{itemize}

Therefore, if NF\_C grows during a segment without \(\Delta\)q
increasing, it must be due to a fresh gate digest (Case 3b), which is
priced and ends the segment. If NF\_C does not grow and \(\Delta\)q does
not increase, then no new cut-level consequences can be derived, and the
segment must either accept (if \textbar{}\ensuremath{\mathcal{U}}\textbar{} = 1) or refute
WorldCommit\_C (by WC-1), ending the segment. \(\blacksquare\)

Verifier check. The NP verifier:

\begin{itemize}
\tightlist
\item
  Recomputes NF\_C(Cons\_C(\(\pi\))) from transcripted designated reads
  and gate objects; checks ConstraintDigest\_C consistency.
\item
  Uses CommitSelector(\(\pi\))\textbar\_C := lexicographically least
  assignment to unresolved cut bits consistent with
  NF\_C(Cons\_C(\(\pi\))) and revealed bits; sets WorldCommit\_C :=
  H(CommitSelector(\(\pi\))\textbar\_C).
\item
  Verifies that any non-accepting segment that does not change NF\_C or
  q ends with a valid Refute(WorldCommit\_C) certificate. The
  certificate is a standardized object referencing at least one
  designated gate/digest artifact and proving that the committed world
  contradicts these artifacts given NF\_C and revealed bits. The
  certificate is admitted into NF\_C only as a UnitRefute for that
  world, excluding exactly one world over C. The verifier checks all
  conditions in polynomial time.
\end{itemize}

Lemma index (C.2):

\begin{itemize}
\tightlist
\item
  C.2 - Segment counting (main)
\item
  C.2.1 - Budget cap on s (bottleneck residual under \(\tau\)-cap)
\item
  C.2.2 - Gate/Constraint update implies tag change
\item
  C.2.3 - One extra elimination per segment beyond \(\Delta\)q
\item
  C.2.T - FG tightness and min-cut canonicality (capstone; combines
  necessity + sufficiency)
\end{itemize}

\textbf{Definitions (rollback semantics).}

\begin{itemize}
\tightlist
\item
  Seed-chain tag: the tuple of (Seed\_u,y\_u) along the active ancestor
  chain (Appendix I), extended with (ConstraintDigest\_C,
  WorldCommit\_C) for the designated bottleneck cut C.
\item
  Resolution prefix across C: the projection of the tag to nodes of C,
  including ConstraintDigest\_C and WorldCommit\_C.
\item
  Rollback segment: a maximal contiguous subrun during which the
  resolution prefix is constant. Segment boundaries occur exactly when
  this projection changes (Keyedness + CDT).
\end{itemize}

\textbf{NF\_C scope note.} The feasible-world count for SC is computed
vs. NF\_C \(\cup\) \{q\}, where NF\_C admits only (i) unit equalities
(bit determinations) and (ii) UnitRefute(\(\omega\)) operations. This is
an over-approximation of the true semantic feasible set - the actual set
of consistent worlds may be smaller due to additional constraints not
captured in NF\_C. Using an over-approximation only strengthens the
m-bound, as it conservatively counts more worlds as feasible, making the
required number of segments m\_seg \(\geq\) 2\^{}(\(\rho\)-s) a lower
bound.

\textbf{Lemma C.2.2 (Gate/Constraint update implies tag change).} If
during a subrun the machine computes any fresh GateDigest\_v whose value
differs from its previous value on the current chain, or if
NF\_C(Cons\_C(\(\pi\))) over the bottleneck cut C changes, then the
resolution prefix across C changes, hence the seed-chain tag changes.
Therefore any subrun adding cut-relevant information (either by new bit
determination or constraint growth) ends the current rollback segment.

\textbf{Lemma WC-1 (WorldCommit refutation excludes exactly one world
--- "+1 only").} Let \(\pi\) be a transcript prefix with resolution
prefix across cut C including WorldCommit\_C(\(\pi\)) = H(\(\omega\)*)
where \(\omega\)* = CommitSelector(\(\pi\))\textbar\_C (the
lexicographically least world consistent with NF\_C(Cons\_C(\(\pi\)))
and revealed bits q). Consider a non-accepting segment endpoint where
the algorithm produces a valid Refute(WorldCommit\_C) certificate. Then:

\begin{enumerate}
\def\labelenumi{\arabic{enumi}.}
\item
  \textbf{Exactly one world excluded:} The refutation eliminates world
  \(\omega\)* and no other worlds from the feasible set \ensuremath{\mathcal{U}}(\(\pi\)).
\item
  \textbf{NF\_C update:} The certificate is admitted into NF\_C as a
  UnitRefute(\(\omega\)*) entry, which updates ConstraintDigest\_C and
  ends the segment.
\item
  \textbf{Verifier-checkable:} The NP verifier checks that the
  certificate references designated artifacts (gate digests, revealed
  bits) and that these artifacts contradict \(\omega\)* given NF\_C and
  q.
\end{enumerate}

\emph{Proof.} The verifier deterministically recomputes the target world
\(\omega\)* = CommitSelector(\(\pi\))\textbar\_C from the current NF\_C
and q (Appendix C.2.b). WorldCommit\_C stores H(\(\omega\)*) only as a
compact representation; the verifier checks equality WorldCommit\_C =
H(\(\omega\)*) and then verifies that the certificate contradicts the
recomputed \(\omega\)*. The Refute(WorldCommit\_C) certificate is a
standardized object that:

\begin{itemize}
\tightlist
\item
  References at least one designated gate/digest artifact or revealed
  bit pattern, and
\item
  Proves that \(\omega\)* contradicts these artifacts given the current
  NF\_C(Cons\_C(\(\pi\))) and q.
\end{itemize}

Because the verifier fixes the target world by deterministic
recomputation (independent of any claimed preimage), a valid certificate
can only refute that unique \(\omega\)*. Any other world \(\omega\)
\(\neq\) \(\omega\)* is not considered by the verifier in this check and
therefore is not excluded. The verifier then:

\begin{itemize}
\tightlist
\item
  Recomputes \(\omega\)* from current NF\_C and q
\item
  Verifies WorldCommit\_C = H(\(\omega\)*)
\item
  Checks that referenced artifacts contradict \(\omega\)*
\item
  Admits UnitRefute(\(\omega\)*) into NF\_C, excluding exactly this one
  world
\end{itemize}

Therefore the refutation excludes exactly one world from \ensuremath{\mathcal{U}}(\(\pi\)),
namely \(\omega\)*, and the NF\_C update (adding
UnitRefute(\(\omega\)*)) changes ConstraintDigest\_C, forcing a new
segment by Lemma C.2.2. \(\blacksquare\)

\textbf{Lemma C.2.3 (One extra elimination per segment beyond
\(\Delta\)q).} Consider a non-accepting rollback segment T. Let t be the
increase in functionally determined cut information during T
(\(\Delta\)q across C; not just RWA reads). Then, relative to the
segment start, the number of feasible worlds across C decreases by at
most a factor 2\^{}t and at most one additional world (the committed
one) is excluded by the terminal refutation of T.

\emph{Proof.} Fix C and let \ensuremath{\mathcal{U}}(\(\pi\)) be the feasible-worlds set at
prefix \(\pi\). Let \(\pi\)\_in and \(\pi\)\_out mark
T\textquotesingle{}s endpoints, and let \(\pi\)\_last be the point after
the last \(\Delta\)q event on C within T. By SCL, each unit of
\(\Delta\)q halves \textbar{}\ensuremath{\mathcal{U}}\textbar{}; thus
\textbar{}\ensuremath{\mathcal{U}}(\(\pi\)\_last)\textbar{} =
\textbar{}\ensuremath{\mathcal{U}}(\(\pi\)\_in)\textbar/2\^{}t. By CDT-1', with NF\_C and q
fixed on C between \(\pi\)\_last and the terminal action, the feasible
set cannot shrink in the suffix: \ensuremath{\mathcal{U}}(\(\pi\)\_out\^{}\(-\)) =
\ensuremath{\mathcal{U}}(\(\pi\)\_last). By WC-1, the terminal action refutes exactly one
committed world (WorldCommit\_C), excluding at most one additional world
at that boundary (which also updates NF\_C as a UnitRefute and ends the
segment). Hence \textbar{}\ensuremath{\mathcal{U}}(\(\pi\)\_out)\textbar{} \(\geq\)
\textbar{}\ensuremath{\mathcal{U}}(\(\pi\)\_in)\textbar/2\^{}t \(-\) 1. \(\blacksquare\)

\textbf{Lemma C.2 (Segment counting).} Let C be an s\(\to\)t cut with
effective residual \(\rho\) := \(\Sigma\)\_\{v\(\in\)C\}(R\_v \(-\)
q\_v). Run a deterministic solver A on x* and partition the run into
rollback segments T\(_1\),...,T\_m.

Define: s := \(\Delta\)q\_C accumulated in T\(_1\),...,T\_\{m-1\}
(functional determination across C)

\textbf{Then:} m\_seg \(\geq\) 2\^{}(\(\rho\) - s)

\emph{Proof.} Let s\_i be the cumulative increase in functionally
determined cut information (\(\Delta\)q on C) by the end of T\_i, so
s\_\{m-1\}=s. Initially \textbar{}\ensuremath{\mathcal{U}}(\(\pi\)\_0)\textbar{} =
2\^{}(\(\rho\)). Across the first m\_seg\(-\)1 non-accepting segments,
C.2.3 implies that each segment can exclude at most one world beyond its
\(\Delta\)q contribution. Therefore, after accounting for \(\Delta\)q
totaling s and at most (m\_seg\(-\)1) additional refutations, the number
of worlds remaining at the start of T\_m is at least 2\^{}(\(\rho\)-s)
\(-\) (m\_seg\(-\)1). By Lemma C.2.ACC-mech (Acceptance purity), the
accept action does not introduce new cut-level information; hence the
accepting segment can begin only when at most one cut-world remains
consistent with the current prefix. Therefore 2\^{}(\(\rho\)-s) \(-\)
(m\_seg\(-\)1) \(\leq\) 1, which rearranges to m\_seg \(\geq\)
2\^{}(\(\rho\)-s). By Lemma C.2.2 (Gate/Constraint update
\(\Rightarrow\) tag change), CDT-1', and WC-1, each non-accepting
segment corresponds to either \(\Delta\)q growth, NF\_C growth, or a
single committed-world refutation, all of which are priced per CDT-3.
\(\blacksquare\)

\textbf{Lemma C.2.ACC-mech (Acceptance purity - mechanical check).} The
verifier enforces that acceptance is valid only when either: (i) all cut
bits are functionally determined in NF\_C (i.e.,
\textbar{}\ensuremath{\mathcal{U}}(\(\pi\))\textbar{} = 1 through functional determination
alone), or (ii) the committed world WorldCommit\_C equals the realized
cut assignment and is the unique world consistent with NF\_C \(\cup\)
\{q\}. Specifically, on accept the verifier recomputes
CommitSelector(\(\pi\))\textbar\_C and verifies that this committed
world matches the actual cut values being accepted. This ensures the
accepting segment introduces no new cut-level information beyond what
was already determined or committed.

\textbf{Lemma C.2.ACC (Acceptance implies uniqueness - logical form).}
Suppose an accepting run outputs a canonical witness W = (w,
G\_\(\tau\), Dig\_\(\tau\)) and passes Algorithm V. Then at the start of
the accepting segment, at most one cut-world across C* is consistent
with the transcript prefix \(\pi\). In particular, acceptance implies
\textbar{}\ensuremath{\mathcal{U}}(\(\pi\))\textbar{} \(\leq\) 1 even if NF\_C is used only
analytically.

\emph{Proof.} Let \(\omega\)* be the realized cut assignment on the
final seed chain. By keyedness and Enc injectivity, different cut-worlds
induce different seed chains beyond the bottleneck. For any distinct
\(\omega\)' \(\neq\) \(\omega\)*, the per-path digest family G\_\(\tau\)
includes, by construction, paths whose S(P) sets are evaluated on the
current seed chain (Appendix C.1.1). By Lemma C.1.2 (Unpredictability
without reads) and address-churn (Lemma A.1.\(\Delta\)), the vector of
digests Dig\_\(\tau\) computed on the chain for \(\omega\)* differs from
that for \(\omega\)' on at least one published path P. Algorithm V
recomputes all digests deterministically on the current seed chain and
checks equality with Dig\_\(\tau\); therefore any W consistent with
\(\omega\)' would be rejected. Hence acceptance forces \(\omega\)* to be
the unique cut-world consistent with \(\pi\). \ensuremath{\square}

\emph{Corollary (extreme case).} If s = 0, then m\_seg \(\geq\)
2\^{}(\(\rho\)).

\textbf{Lemma C.2.1 (Budget cap on s).} With the gate-horizon budget
(Section 6.2.10), the pre-final agreement satisfies s \(\leq\) S(n) =
\(\Theta\)(\(\tau\)(n)·\(\lambda\)\_base(n)). Moreover, the effective
residual at the bottleneck obeys \(\rho\) \(\geq\) \(\lambda\)\_base
\(-\) s. Thus for \(\tau\)(n) \(\leq\) 1/2, m\_seg \(\geq\)
2\^{}(\(\rho\) - s) \(\geq\) 2\^{}((1-2\(\tau\))\(\rho\)) =
2\^{}(\(\rho\))/poly(n) (QP-sharp). Combining with the FG time baseline
(Appendix C.1.1) yields unconditional tight single-run bounds.

\emph{Proof.} The published GREQ map marks nodes on C* with GREQ=0
(pre-horizon) and GREQ=1 (gated). By FG-2 the total allowance across C*
before the last receiving window is
\(\Theta\)(\(\tau\)·\(\lambda\)\_base); hence s \(\leq\) S(n). By Lemma
C.2.1 we have \(\rho\) \(\geq\) \(\lambda\)\_base \(-\) s. Apply Lemma
C.2 with this s and the per-segment baseline from Corollary C.1.1 to
obtain the stated inequalities. \(\blacksquare\)

\begin{center}\rule{0.5\linewidth}{0.5pt}\end{center}

\textbf{Theorem C.2.T (FG Tightness and Min-Cut Canonicality).}

Fix the L* construction with FG wiring. Let \(\lambda\)\_base be the
profile residual on the designated bottleneck cut C*, W\_min the
profile\textquotesingle{}s window parameter, \(\rho\) the run-dependent
residual at the start of the last segment, and s the pre-final agreement
across C*. Then for every deterministic k-tape TM A:

\begin{enumerate}
\def\labelenumi{\arabic{enumi}.}
\item
  \textbf{Single-run lower bound:} T\_A(x) \(\geq\)
  2\^{}(\(\rho\)-s)·\(\Omega\)(n/W\_min). In particular, FG enforces s
  \(\leq\) \(\Theta\)(\(\tau\)(n)·\(\lambda\)\_base) and \(\rho\)
  \(\geq\) \(\lambda\)\_base - s, so T\_A(x) \(\geq\)
  2\^{}((1-\(\Theta\)(\(\tau\)))\(\lambda\)\_base)·\(\Omega\)(n/W\_min).
\item
  \textbf{Achievability; tightness up to subpoly: There exist calibrated
  instance/solver pairs whose complexity is
  2\^{}(\(\lambda\)\_base+o(\(\lambda\)\_base)) (QP-sharp:
  n\^{}(\(\Theta\)(log n\_core))), hence the min-cut aggregator
  \(\lambda\)(A,x) = min\_C \(\Sigma\)\_\{v\(\in\)C\}(R\_v-q\_v) is
  canonical and tight up to subpolynomial factors}.
\end{enumerate}

\emph{Proof.}

\textbf{Step 1 - Segment multiplicity.} Let the run be partitioned into
rollback segments T\(_1\),...,T\_m. Write \(\rho\) =
\(\Sigma\)\_\{v\(\in\)C\}(R\_v\(-\)q\_v) for the cut residual (the
effective residual on the chosen cut). Define s := \(\Delta\)q\_C
accumulated before the last segment.

\textbf{Segment Counting (Lemma C.2)} shows: m\_seg \(\geq\)
2\^{}(\(\rho\)-s). Intuition: before acceptance, each non-accepting
segment can eliminate \textbf{at most one} extra world on top of
whatever functional determination it paid; thus to shrink 2\^{}\(\rho\)
worlds down to a single accepting world when only s bits were
functionally determined, you need at least 2\^{}(\(\rho\)-s) such
segments.

\textbf{Step 2 - Per-segment time baseline under FG.} FG wires a
\textbf{cut-gate} that must be recomputed on the \textbf{current seed
chain} when progressing the final chain past the gate horizon. For every
non-accepting rollback segment that progresses the chain, FG forces the
evaluation of a identity digest over
\textbar{}S(P)\textbar{}=\(\Theta\)(n/W\_min) seed-dependent addresses;
on a k-tape TM this costs \(\Omega\)(n/W\_min) operations even if salts
are pre-cached, by (i) \textbf{parity has full sensitivity} (Lemma
5.5.1.c) and (ii) \textbf{address churn} (Lemma A.1.\(\Delta\)) moves a
\(\Theta\)-fraction of designated addresses whenever the seed chain
changes. Hence each such segment costs \(\Omega\)(n/W\_min) time.

\textbf{Step 3 - Multiply segment count by baseline.} Combining Steps
1-2 yields the one-run bound T\_A(x) \(\geq\) m·\(\Omega\)(n/W\_min)
\(\geq\) 2\^{}(\(\rho\)-s)·\(\Omega\)(n/W\_min). This is exactly Theorem
8.A\textquotesingle{}s single-run FG inequality.

\textbf{Step 4 - Capping pre-final agreement and relating \(\rho\) to
\(\lambda\)\_base.} FG\textquotesingle{}s \textbf{gate-horizon budget}
(Lemma C.2.1) structurally limits pre-final agreement: for the published
\(\tau\)-calibration we have s \(\leq\)
\(\Theta\)(\(\tau\)(n)·\(\lambda\)\_base(n)), \(\rho\) \(\geq\)
\(\lambda\)\_base - s. Thus 2\^{}(\(\rho\)-s) \(\geq\)
2\^{}((1-\(\Theta\)(\(\tau\)))\(\lambda\)\_base) (and in the QP-sharp
profile with \(\tau\)(n)=\(\Theta\)(log n/\(\lambda\)\_base), this gives
2\^{}(\(\lambda\)\_base)/poly(n\_core)).

Putting Steps 1-4 together proves (1).

\textbf{Step 5 - Achievability (sufficiency) and canonical tightness of
min-cut.}

\textbf{Construction of calibrated achievers.} Consider the
exhaustive-search algorithm A\_exhaust that: (i) reads all overlay
parameters (Enc, F\_overlay, H\_v, Sel\_v) and salt pools \{U\_v\} in
O(n log n) time; (ii) at the designated bottleneck cut C* with residual
\(\lambda\)\_base, systematically enumerates all
2\^{}(\(\lambda\)\_base) possible assignments to the unresolved
coordinates across C*; (iii) for each candidate assignment, checks
consistency with overlay constraints (rank conditions via Completeness,
seed propagation via Enc, gate digests via published schema in Appendix
C.1.1) in poly(n\_core) time; (iv) upon finding a consistent assignment,
follows the unique determined path to the sink via seed propagation and
verifies the acceptance predicate (by Closure/Recoverability, Appendix
J: filling cut choices determines a unique global completion consistent
with the fixed prefix). Total time: T\_exhaust(n) =
O(2\^{}(\(\lambda\)\_base) · poly(n\_core)).

Since \(\lambda\)\_base = min\_C \(\Sigma\)\_\{v\(\in\)C\}(R\_v-q\_v) is
the \textbf{minimum} over all cuts, any algorithm must cross some cut
with residual \(\geq\) \(\lambda\)\_base. Therefore A\_exhaust is
optimal up to poly(n\_core) factors, establishing T(n) \(\leq\)
2\^{}(\(\lambda\)\_base+o(\(\lambda\)\_base))·poly(n\_core).

\textbf{For QP-sharp:} \(\lambda\)\_base = \(\Theta\)(log² n)
\(\Rightarrow\) T\_exhaust(n) = n\^{}(\(\Theta\)(log n\_core)).
\textbf{For flat:} \(\lambda\)\_base = \(\Theta\)(n) \(\Rightarrow\)
T\_exhaust(n) = 2\^{}(\(\Theta\)(n)).

\textbf{Canonical tightness.} By the Sharpness axiom (Def. 5.1.3.1), any
admissible aggregator A\textquotesingle{} must satisfy
A\textquotesingle{} \(\leq\) \(\lambda\)+o(\(\lambda\)\_base) (otherwise
the achiever A\_exhaust would contradict the claimed bound
2\^{}(A\textquotesingle{}+o(A\textquotesingle{}))). Conversely, SCL +
cut composition (Theorem 7.A, Theorem J.1-PROD) make \(\lambda\)
necessary, whence \(\lambda\) \(\leq\) A\textquotesingle{}. Therefore
\(\lambda\)(A,x) \(\leq\) A\textquotesingle{} \(\leq\)
\(\lambda\)(A,x)+o(\(\lambda\)\_base), i.e., min-cut is
\textbf{canonical and tight up to subpolynomial factors}.
\(\blacksquare\)

\textbf{Corollaries (standard profiles):}

\textbullet{} \textbf{QP-sharp profile.} \(\lambda\)\_base=\(\Theta\)(log² n),
W\_min=\(\Theta\)(log n\_core), \(\tau\)(n)=\(\Theta\)(log
n/\(\lambda\)\_base). Then T(n) \(\geq\) n\^{}(\(\Omega\)(log n\_core))
and calibrated achievers attain T(n) \(\leq\) n\^{}(\(\Theta\)(log
n\_core)) (tight up to poly factors).

\textbullet{} \textbf{Flat (exponential) profile.} \(\lambda\)\_base=\(\Theta\)(n),
W\_min=\(\Theta\)(1). Then T(n) \(\geq\) 2\^{}(\(\Theta\)(n)), and
achievers match within poly factors.

\begin{center}\rule{0.5\linewidth}{0.5pt}\end{center}

\subsubsection{\texorpdfstring{C.3 How FG and SC Interact:
L*\textquotesingle{}s Structure Forces
Tightness}{C.3 How FG and SC Interact: L*\textquotesingle s Structure Forces Tightness}}\label{c3-how-fg-and-sc-interact-ls-structure-forces-tightness}

\textbf{Note:} FG and SC are complementary views of how
L*\textquotesingle{}s structure yields tight bounds: SC captures
multiplicative growth; FG prevents front-loading.

\textbf{The Mathematical Symphony}:

\begin{enumerate}
\def\labelenumi{\arabic{enumi}.}
\tightlist
\item
  \textbf{SC reveals}: L*\textquotesingle{}s structure FORCES m\_seg
  \(\geq\) 2\^{}(\(\rho\)-s) segments (not algorithmic choice)
\item
  \textbf{FG by design}: We built FG into L* to LIMIT s \(\leq\)
  \(\tau\)·\(\rho\) (structural throttling)
\item
  \textbf{Resulting bound}: m\_seg \(\geq\) 2\^{}((1-\(\tau\))\(\rho\))
  = 2\^{}\(\rho\)/poly(n\_core) unconditionally
\end{enumerate}

\textbf{Implication:} L* does not permit front-loading: gate digests are
wired into seeds (Keyedness + RWA), making premature accumulation
ineffective and unverifiable. The bounds are tight because this
structure is embedded in the language:

\begin{itemize}
\tightlist
\item
  \textbf{QP-sharp} (\(\lambda\)\_base = \(\Theta\)(log² n)): L* FORCES
  n\^{}(\(\Theta\)(log n\_core)) complexity
\item
  \textbf{Exponential} (\(\lambda\)\_base = \(\Theta\)(n)): L* FORCES
  2\^{}(\(\Theta\)(n))/poly(n\_core) complexity
\item
  \textbf{All profiles}: L* creates corresponding tight necessities
\end{itemize}

\textbf{Conclusion:} Because we built FG into L*\textquotesingle{}s
construction, the bounds are unconditional (model-independent) for this
specific language.

\subsubsection{\texorpdfstring{C.4 Time Conversion: How
L*\textquotesingle{}s Structure Manifests Across
Models}{C.4 Time Conversion: How L*\textquotesingle s Structure Manifests Across Models}}\label{c4-time-conversion-how-ls-structure-manifests-across-models}

\textbf{Result:} Converts SCL bottlenecks (\(\lambda\)(A,x) residual, m
segments) into time bounds across k-tape TMs, RAMs, and branching
programs via model-specific operational facts (per-step inflow B, LOAD
costs, probe work).

\textbf{Clarification:} Time bounds reflect how L*\textquotesingle{}s
semantic requirements manifest in each computational framework; the same
conservation-law obligation appears differently across models.

\paragraph{C.4.1 Universal Structural
Requirements}\label{c41-universal-structural-requirements}

\textbf{L*\textquotesingle{}s semantic necessity is absolute:} The
artifact bounds (branches, states, width) emerge from
L*\textquotesingle{}s structure - independent of any computational
model.

\textbf{How structure creates time bounds:} L*\textquotesingle{}s
requirements translate to time through fundamental model properties:

\begin{itemize}
\tightlist
\item
  \textbf{k-tape TMs}: L* FORCES acquisition of \(\Theta\)(n log n)
  bits; physics limits to B :=
  k\(\lceil\)log\(_2\)\textbar{}\(\Gamma\)\textbar{}\(\rceil\) bits/step
  \(\to\) unavoidable time
\item
  \textbf{Word-RAM}: L*\textquotesingle{}s information requirements
  divided by word size \(\to\) structural time necessity
\item
  \textbf{Branching programs}: Each LOAD L* forces costs time \(\to\)
  multiplication is unavoidable
\item
  \textbf{Key truth}: These are not bandwidth limits - they are how
  L*\textquotesingle{}s structure manifests
\end{itemize}

\paragraph{C.4.2 Time Consequences for Specific
Models}\label{c42-time-consequences-for-specific-models}

\textbf{Generic pricing.} If a designated read/LOAD costs at least a
constant, one accepting try costs \(\Omega\)(\(\Sigma\)\_v
(R\_v\(-\)q\_v)).

\textbf{TM pricing (restart lane).} In the restart/brute-force lane
where each try is independent, every distinct try (exploring a distinct
resolution prefix) requires \(\geq\)1 first-use designated read (RWA),
so with B :=
k\(\lceil\)log\(_2\)\textbar{}\(\Gamma\)\textbar{}\(\rceil\), total time
\(\geq\) 2\^{}(\(\lambda\)(A,x) - log\(_2\) Alt(C*))/B.

\textbf{Quick Restart Bound (self-contained).} This lane can be stated
and used in a minimal, airtight form:

\begin{itemize}
\tightlist
\item
  Seed-locked decode forces designated-read engagement to evaluate
  \(\varphi\)(w); there is no CNF-first shortcut (Scope note;
  seed-locked decode \(\Phi\)\textasciitilde{}).
\item
  Fresh information is admitted only via first-use designated reads and
  is schedule-agnostic (RWA monotonicity/subadditivity).
\item
  Under forgetting across resolution-prefix changes at the bottleneck
  cut C*, the across-tries inequality applies: Q(C*) + log\(_2\) Alt(C*)
  + log\(_2\) \ensuremath{\mathbb{E}}{[}tries{]} \(\geq\) \(\Lambda\)(C*). Hence \ensuremath{\mathbb{E}}{[}tries{]}
  \(\geq\) 2\^{}(\(\Lambda\)(C*)-(Q(C*)+log\(_2\) Alt(C*))) =
  2\^{}(\(\lambda\)(A,x) - log\(_2\) Alt(C*)).
\item
  Each \textbf{distinct try} (exploring a distinct resolution prefix
  across C*) must include at least one first-use designated read to
  establish that prefix; otherwise it deterministically replicates a
  prior try and does not count toward \ensuremath{\mathbb{E}}{[}tries{]}. Therefore, on a
  k-tape TM, steps \(\geq\) \ensuremath{\mathbb{E}}{[}tries{]}/B \(\geq\)
  2\^{}(\(\lambda\)(A,x) - log\(_2\) Alt(C*))/B.
\item
  \textbf{Early failure detection.} Most bits at C* lie beyond the gate
  horizon (GREQ\_v=1; see §6.2.9 gate horizon budget:
  \(\Sigma\)\_\{v\(\in\)C*, GREQ\_v=0\} R\_v \(\leq\)
  \(\Theta\)(\(\tau\) \(\lambda\)\_base) with \(\tau\) = o(1)).
  Computing seeds for GREQ\_v=1 nodes requires GateDigest\_v \(\to\)
  \(\Omega\)(n/W\_min) steps (Lemma 5.5.1.c). However, \textbf{failed
  tries may abort at pre-horizon nodes} (GREQ\_v=0) when random guesses
  hit seed-mismatch or early constraint violations - before incurring
  the full gate computation cost. Our lower bound charges only the
  guaranteed minimal per-try cost (\(\geq\)1 first-use read at inflow
  B). If a failed try crosses the horizon and computes any
  GateDigest\_v, that only increases its cost and does not weaken the
  bound; in such regimes the \textbf{single-run} analysis (Appendix C.2)
  becomes the tighter characterization. This separation makes explicit
  that the \textbf{across-tries bound prices discovery} (exponentially
  many attempts to find the right seed chain), while the
  \textbf{single-run bound prices execution} (linearly many gates once
  committed).
\end{itemize}

This argument applies equally to decision-style attempts on x*:
seed-locked decode is mandatory before evaluating \(\exists\)w:
\(\varphi\)(w)=1, so designated-read accounting and the across-tries
lower bound still govern runtime in the restart/brute-force lane.

\textbf{TM pricing (single-run lane).} With GateDigest\_v wired into
seeds (GREQ\_v=1 on the gate horizon) and the CDT/WC extensions
(ConstraintDigest\_C and WorldCommit\_C) included in the resolution
prefix across the bottleneck cut, every non-accepting segment that
progresses the final chain must either (i) produce a fresh gate digest,
or (ii) grow NF\_C (ConstraintDigest change), or (iii) refute
WorldCommit\_C. By CDT-3 each such event requires designated work
costing \(\Omega\)(n/W\_min) TM steps (probe-work pricing, not arity;
Lemma 5.5.1.c + Lemma A.1.\(\Delta\)). The gate horizon budget (§6.2.9)
enforces s \(\leq\) B(n) = \(\Theta\)(\(\tau\)(n)·\(\lambda\)\_base(n))
along the bottleneck; since \(\rho\) \(\geq\) \(\lambda\)\_base \(-\) s,
this also implies s = O(\(\tau\)\(\rho\)). Combining the per-segment
baseline with Segment Counting (Appendix C.2) yields tight
\textbf{single-run} bounds.

\paragraph{C.4.3 Per-Try Cost Analysis}\label{c43-per-try-cost-analysis}

\textbf{Lemma C.4.1 (Per-try must-touch).} In any complete try, the
solver reads at least \(\Sigma\)\_v R\_v \textbf{designated salt words}
(one per selected primitive) across the D nodes; designated address
ranges are disjoint across nodes.

\emph{Proof sketch.} By construction, each of the R\_v selected
primitives depends on a \textbf{distinct} unread salt at node v; by
Dependency, computing y\_v (hence Seed\_\{succ(v)\}) requires fixing all
R\_v coordinates. Summing over nodes gives \(\Sigma\)\_v R\_v reads.
\(\Delta\)

\textbf{Theorem C.4.2 (L*\textquotesingle{}s Structure Creates Time
Complexity).} For any deterministic k-tape TM with alphabet \(\Gamma\)
attempting to solve L*:

\begin{enumerate}
\def\labelenumi{\arabic{enumi}.}
\tightlist
\item
  \textbf{Physics meets structure:} B :=
  k\(\lceil\)log\(_2\)\textbar{}\(\Gamma\)\textbar{}\(\rceil\) bits/step
  (information flow limit)
\item
  \textbf{In L*, we can prove:} \(\Sigma\)\_v R\_v = \(\Theta\)(n log n)
  bits must be acquired for correctness
\item
  \textbf{Mathematical necessity:} Steps \(\geq\) \(\Theta\)(n log n)/B
  = \(\Omega\)(n log n) when B = O(1)
\end{enumerate}

\textbf{How L*\textquotesingle{}s structure manifests through different
algorithmic behaviors:} The restart/forgetting and single-run
consequences are captured immediately above by the restart pricing (with
the Quick Restart Bound) and the single-run per-segment baseline; we
omit redundant summaries here.

\paragraph{C.4.4 Model-Specific Time
Bounds}\label{c44-model-specific-time-bounds}

\textbf{Corollary C.4.3 (Worst-case extraction from the uniform bound).}
For every deterministic A and infinitely many n, there exists a seed s
\(\in\) \ensuremath{\mathbf{G}}\_n such that time\_M(A, G(n,s)) \(\geq\)
2\^{}\(\Omega\)(\(\lambda\)(A,x; n)) · \(\lambda\)(A,x; n).

\emph{Proof.} By the averaging principle: if \ensuremath{\mathbb{E}}{[}time{]} \(\geq\) T,
then \(\exists\)s with time(s) \(\geq\) T. Apply this to the uniform
bound; choose an infinite subsequence of n with growing \(\lambda\)(A,x;
n).

\textbf{Streaming/Windowed Models:} A windowed load lemma (Appendix C.4)
translates cut residuals into LOAD lower bounds within node windows,
giving conservative time bounds T \(\geq\) \(\Sigma\)\_W max(0,
\(\lambda\)\_W - 2S·w) in RAM/branching-program/cell-probe models.

\textbf{Time from queries/proofs:} Our query\(\to\)time and
LOAD\(\to\)time conversions follow established practice in decision-tree
and cell-probe literatures. Similarly, proof size\(\to\)runtime is
standard for DPLL/CDCL solvers implementing Resolution.

\textbf{Key insight:} The time bounds depend fundamentally on the
semantic necessity established in the main text. The conversions here
are standard applications of model-specific operational facts.

\emph{Notes.} The \textbf{disjoint-atoms} hypothesis (Lemma 6.5, §6.3.2)
is the only place we need "no cross-help" between obligations; it is
already proved from the unique-neighbor mapping and per-node disjoint
U\_v ranges in §6.2.3. Recoverability/closure (Axiom A4, §6.2.7; Lemma
J.0) ensures that filling cut choices determines a unique global
completion consistent with the fixed transcript prefix (App. J preface).

\begin{center}\rule{0.5\linewidth}{0.5pt}\end{center}

\subsection{Appendix D: TM/Word-RAM Time
Conversions}\label{appendix-d-tmword-ram-time-conversions}

\textbf{Purpose:} Converts SCL\textquotesingle{}s semantic bottlenecks
(\(\lambda\)(A,x) residual, 2\^{}(\(\rho\)-s) segments) into concrete
TM/RAM time bounds via per-step inflow coupling (referenced from §8.A,
§9.4).

\textbf{Key result:} Lemma D.2.1 (RWA first-use reads \(\to\) TM steps):
F first-use bits require \(\geq\) F/B steps where B =
k\(\lceil\)log\(_2\)\textbar{}\(\Gamma\)\textbar{}\(\rceil\). Combined
with 2\^{}\(\Delta\) tries (restart lane) or m segments (single-run
lane), yields time \(\geq\) 2\^{}\(\Omega\)(\(\lambda\))/B.

\textbf{Structure:} D.1 shows verification is poly-time with full
commitment. D.2-D.4 convert exponential artifact requirements into time.
D.5 proves k-tape TMs realize address-pool abstraction.

This appendix sketches how the semantic bounds translate into time on
standard models. The key idea is to couple the per-step information
inflow with the required first-use information (and/or required
artifacts) to obtain time lower bounds.

\subsubsection{D.1 Verification is Polynomial: Full Commitment Enables
Efficiency}\label{d1-verification-is-polynomial-full-commitment-enables-efficiency}

\textbf{Theorem D.1 (Verification Complexity).} \emph{When q\_v = R\_v
(full commitment), verification requires polynomial time on all standard
models:}

\begin{itemize}
\tightlist
\item
  \textbf{One-path verifier:} Commits R\_v per node \(\to\) costs add to
  \(\Sigma\)\_v R\_v = \(\Theta\)(n log n)
\item
  \textbf{Streaming:} Passes add along paths \(\to\) strongly polynomial
  per path
\item
  \textbf{RAM:} Direct reads of R\_v designated salts \(\to\) O(R\_v)
  time per node
\end{itemize}

Full commitment avoids the multiplicative explosion: the
conservation-law obligation is satisfied additively.

\subsubsection{D.2 RAM and the source of
exponentiality}\label{d2-ram-and-the-source-of-exponentiality}

\textbf{Result:} Exponentiality arises from 2\^{}(\(\Delta\)(C*)) tries
structurally required by L*\textquotesingle{}s instance structure across
the bottleneck cut (restart/brute-force lane), not from per-try
verification cost (which is polynomial). Combined with Lemma D.2.1,
yields time \(\geq\) 2\^{}(\(\Delta\)(C*))/B.

\paragraph{\texorpdfstring{D.2.1 Lemma (RWA first-use reads \(\to\) TM
steps)}{D.2.1 Lemma (RWA first-use reads \textbackslash to TM steps)}}\label{d21-lemma-rwa-first-use-reads--tm-steps}

Let A be a deterministic k-tape TM with alphabet \(\Gamma\) and let B :=
k·\(\lceil\)log\(_2\)\textbar{}\(\Gamma\)\textbar{}\(\rceil\). Consider
any accepting run \(\pi\) on input x* where Receiving-Window Attribution
(RWA) credits F first-use designated payload bits (i.e., the set of
bit-positions whose first valid use occurs in \(\pi\)). Then the number
of TM steps T in \(\pi\) satisfies T \(\geq\) F / B.

Proof. In each TM step, at most one symbol per tape is read, importing
at most \(\lceil\)log\(_2\)\textbar{}\(\Gamma\)\textbar{}\(\rceil\)
fresh bits per tape, i.e., at most B fresh bits total. By definition of
RWA, each first-use designated payload bit contributes one unit of fresh
information that must be imported in some step of \(\pi\). Therefore,
importing F such bits requires at least F/B steps. This bound is
schedule-independent (reordering/prefetch does not reduce RWA credits;
see Lemma 6.1-RWA and §D.5). \(\blacksquare\)

Observation. On an unrestricted adaptive RAM, a single verification try
can be polynomial by reading the R\_v words per node that a full
verifier would read.

Where exponentiality arises.

\begin{itemize}
\tightlist
\item
  The 2\^{}(\(\Delta\)(C*)) tries are structurally required by
  L*\textquotesingle{}s instance structure across a bottleneck cut
  (restart/brute-force lane; across-tries inequality)
\item
  Each try incurs at least one first-use designated read (RWA)
\item
  Total time \(\geq\) 2\^{}(\(\Delta\)(C*))/B. If a solver chooses to
  run full verification per try, the per-try cost can be poly(n\_core),
  yielding \(\geq\) 2\^{}(\(\Delta\)(C*))·poly(n\_core). Our lower bound
  uses only the \(\Omega\)(1/B) per-try minimum.
\end{itemize}

This does not contradict the paradigm-specific results.
EO/DP/Resolution/OBDD manifest exponential artifacts
(trees/tables/proofs/diagrams) because their correctness obligations
appear as many simultaneous states within a run; when expressed as
restarts, the same obligation appears as many tries.

\subsubsection{D.3 Semantic vs. machine: structure determines
complexity}\label{d3-semantic-vs-machine-structure-determines-complexity}

We count semantics, not machine steps.

\begin{itemize}
\tightlist
\item
  Computation artifacts: branches, keys, clause width, OBDD nodes
\item
  Distinguishable states: the 2\^{}(\(\lambda\)(A,x)) possibilities L*
  forces algorithms to track
\item
  Structural requirements: what correctness requires, independent of
  hardware
\end{itemize}

Exponential behavior arises from the instance structure - the need to
distinguish 2\^{}(\(\lambda\)(A,x)) states - not from bandwidth or cache
effects.

\subsubsection{D.4 Model hooks: time from semantic
bounds}\label{d4-model-hooks-time-from-semantic-bounds}

\textbf{Result:} Converts SCL bottlenecks into time via model-specific
per-step inflow bounds: B bits/step for k-tape TMs, t·w bits for
cell-probe with w-bit cells. Two lanes: restart (time \(\geq\)
2\^{}\(\Delta\)/B) and single-run (time \(\geq\)
2\^{}(\(\rho\)-s)·\(\Omega\)(n/W\_min)).

\paragraph{D.4.1 Primitive-touch lower bounds from Alt/tries
(tape-level)}\label{d41-primitive-touch-lower-bounds-from-alttries-tape-level}

We record two tape-level conversions that do not rely on any engagement
heuristic, only on semantic counters and instance-side multiplicativity.

\begin{enumerate}
\def\labelenumi{\arabic{enumi})}
\item
  Restart/brute-force lane. If \ensuremath{\mathbb{E}}{[}tries{]} \(\geq\) 2\^{}(\(\Delta\))
  for some residual deficit \(\Delta\) across a cut (Appendix J;
  across-tries inequality), then with B :=
  k\(\lceil\)log\(_2\)\textbar{}\(\Gamma\)\textbar{}\(\rceil\), total
  steps T satisfy \ensuremath{\mathbb{E}}{[}T{]} \(\geq\) \ensuremath{\mathbb{E}}{[}tries{]}/B \(\geq\)
  2\^{}(\(\Delta\))/B because each try consumes at least one first-use
  designated read (RWA) and per-step inflow \(\leq\) B (Lemma D.2.1).
  Yao/minimax fixes coins to give a deterministic input achieving this
  expectation for the coin-fixed solver.
\item
  Single-run lane (structural multiplicativity). If Segment Counting
  (Appendix C.2) yields m\_seg \(\geq\) 2\^{}(\(\rho\)-s) non-accepting
  rollback segments that progress the final chain (purely from SCL + cut
  composition), and each such segment performs at least one priced event
  (fresh gate digest, NF\_C growth, or committed-world refutation), then
  by Appendix C.1.1 and CDT-3 each segment costs \(\Omega\)(n/W\_min) TM
  steps. Thus steps T \(\geq\) m\_seg·\(\Omega\)(n/W\_min) \(\geq\)
  2\^{}(\(\rho\)-s)·\(\Omega\)(n/W\_min). Under the gate-horizon budget
  s \(\leq\) \(\Theta\)(\(\tau\)·\(\lambda\)\_base) (Appendix C.2.1),
  this gives T \(\geq\)
  2\^{}(\(\Omega\)(\(\lambda\)\_base))/poly(n\_core). This conversion
  uses only (i) semantic multiplicativity and (ii) instance-side segment
  multiplicity; it is independent of scheduling/prefetch (RWA
  invariance).
\end{enumerate}

Examples of conversions (informal sketches):

\begin{itemize}
\tightlist
\item
  k-tape TMs. Each step acquires at most B =
  k\(\lceil\)log\(_2\)\textbar{}\(\Gamma\)\textbar{}\(\rceil\) fresh
  bits (§5.5.1). Combined with the 2\^{}(\(\lambda\)(A,x) - log\(_2\)
  Alt(C*)) required tries (restart/brute-force lane; Lemma 7.R), this
  yields expected time \ensuremath{\gtrsim} 2\^{}(\(\lambda\)(A,x) - log\(_2\) Alt(C*)) ·
  poly(n\_core). With FG (Appendix C), single-run bounds match profile
  exponents up to poly factors.
\item
  Cell-probe. With w-bit cells and t probes per step, first-use bits per
  step \(\leq\) t·w; time lower bounds scale accordingly with the
  required first-use information and/or the number of required attempts.
\item
  I/O models. With block size B, required LOADs translate to I/Os via
  the windowed-load accounting (Appendix C.4), giving conservative time
  bounds.
\end{itemize}

These are manifestations of the same conservation-law requirement:
distinguishing 2\^{}(\(\lambda\)(A,x)) cases within a run, or achieving
success across restarts with 2\^{}(\(\lambda\)(A,x) - log\(_2\) Alt(C*))
tries, converts (via per-step inflow) into time.

\subsubsection{D.5 Address pools on a k-tape TM (record layout + 2-tape
sketch)}\label{d5-address-pools-on-a-k-tape-tm-record-layout--2-tape-sketch}

\textbf{Result:} k-tape TMs realize the address-pool abstraction via
parseable record names \texttt{(p,u)} where p = Enc(v) is the pool
identifier. Preserves A1 (Hermeticity via designated payloads), A2
(Injectivity via injective Enc and \(\pi\)\_v), and RWA first-use
accounting. Hence SCL applies verbatim to k-tape TMs: if q falls short
of R at a bottleneck cut, the algorithm must maintain 2\^{}(R-q)
artifacts or pay time to resolve them.

Turing machines lack random access in the RAM sense; they have
sequential tapes. In our framework, an "address" is a self-delimiting
string that names a record on a designated tape. Disjoint pools U\_v are
realized by making the pool identifier part of that name and by using an
injective, parseable encoding (Enc) so names are unambiguous and
composable. Hermeticity (A1) is enforced by requiring that all instance
information the algorithm may import appears as payloads of these named
records; Injectivity (A2) ensures distinct histories yield distinct
names; RWA defines when a read is credited as first-use information.

Record format (prefix-free, parseable; consistent with §A.5 Enc)

\begin{itemize}
\item
  Each record is a concatenation of length-delimited fields; lengths are
  in decimal and separated by a reserved delimiter \texttt{:}. A
  concrete layout is:

  \texttt{\textbar{}p\textbar{}:p:\textbar{}u\textbar{}:u:\textbar{}len\textbar{}:len:payload}

  where \texttt{p\ =\ Enc(v)} (pool id),
  \texttt{u\ in\ U\_v\ =\ \{0,1,...,D\_v-1\}} is the designated address
  (we store \texttt{u} in decimal), and \texttt{len} is the payload
  length in bits (or symbols over \(\Gamma\)). The \texttt{payload} is
  the content of the designated cell (e.g., a salt \(\sigma\)\_u or
  other overlay data as specified in §A.2-A.3). Records are concatenated
  back-to-back on a designated input tape M. Because Enc is injective
  and parseable (§A.5), the tuple \texttt{(p,u)} is uniquely recoverable
  by a left-to-right scan.
\end{itemize}

Naming and designated addresses

\begin{itemize}
\tightlist
\item
  For node v and index pair (j,\(\ell\)), the public addressing function
  is F\_overlay(Seed\_v; j,\(\ell\)) := u = \(\pi\)\_v(j,\(\ell\))
  (§A.1). The canonical record name (logical address) is then the pair
  \texttt{(p,u)} with \texttt{p=Enc(v)}. Pools are the subsets with
  identical \texttt{p} values; "disjoint pools" means that different v
  yield disjoint name prefixes and hence disjoint record sets.
\end{itemize}

-2-tape realization (deterministic; extends to k tapes trivially)

\begin{itemize}
\tightlist
\item
  Tape M (read-only input): Holds the multiset of records for all pools,
  in any fixed canonical order (e.g., lexicographic by \texttt{(p,u)}).
  This tape contains the only importable instance bits, satisfying
  Hermeticity (A1). The machine never writes to M.
\item
  Tape W (work): The algorithm computes seeds and target indices here.
  Given v and (j,\(\ell\)), it computes \texttt{Seed\_v} (§A.5), then
  u=\(\pi\)\_v(j,\(\ell\)) (§A.1), and forms the target key
  \texttt{(p,u)} with \texttt{p=Enc(v)}.
\item
  Lookup(key): Starting from the current head position on M, parse
  successive record headers by reading
  \texttt{\textbar{}p\textbar{}:p:\textbar{}u\textbar{}:u:\textbar{}len\textbar{}:len:};
  when \texttt{(p,u)} matches the target, the next symbol entered is the
  first symbol of \texttt{payload}. For batch retrieval of many keys,
  use a single streaming pass (see §D.5.1) rather than resetting to the
  tape start per key.

  \begin{itemize}
  \tightlist
  \item
    RWA (first-use credit): Credit is assigned when the machine first
    reads a payload symbol that is used to functionally determine a
    counted bit along the current seed chain (q\_v aggregates these
    first-use bits; §5.5, §A.6). Merely scanning headers or unread parts
    of a payload accrues no credit; re-reads do not accrue additional
    credit (RWA box in §4/A.6). This matches the "first valid use" rule
    and ensures every credited bit corresponds to at least one
    designated read in the appropriate pool.
  \item
    Injectivity (A2): Because Enc and \(\pi\)\_v are injective, distinct
    world histories (distinct parent tuples or seeds) induce distinct
    logical names \texttt{(p,u)} and thus cannot be legally merged
    without actually resolving bits.
  \end{itemize}
\end{itemize}

Why this suffices on TMs (intuitive but precise)

\begin{itemize}
\tightlist
\item
  Logical vs physical addressing. The "address" here is a name, not a
  RAM pointer. A TM resolves names by sequential parsing;
  prefix-freeness ensures unique matching. Random access is unnecessary
  for the semantics: SCL counts how much fresh, legitimate information
  is imported, not how fast a head reaches a cell.
\item
  Per-step inflow bound. On a k-tape TM, at most B =
  k\(\lceil\)log\(_2\)\textbar{}\(\Gamma\)\textbar{}\(\rceil\) fresh
  bits can be imported per step (§5.5). Sequential scans can only
  increase time, never reduce the obligation q + \(\Phi\) \(\geq\) R;
  hence our lower bounds are conservative on TMs.
\item
  Disjoint-pool factoring. Since pools are identified by
  \texttt{p=Enc(v)} and \texttt{(p,u)} are unique, worlds factor across
  pools as required for cut composition. The min-cut residual
  \(\lambda\) then prices into time exactly as in §7 and §9.
\end{itemize}

Practical notes and variants

\begin{itemize}
\tightlist
\item
  Canonical order. Storing records on M sorted by \texttt{(p,u)} makes
  lookup time predictable; any fixed canonical order suffices for
  correctness and for the lower bounds.
\item
  Visited marks (for RWA). To avoid mutating M, the TM can keep a
  visited-set log of \texttt{(p,u)} on W; first visit is credited,
  subsequent visits are checked against this log and not credited.
\item
  Address-churn enforces recomputation. When \texttt{Seed\_v} changes
  (different world), pi\_v changes (§A.1.\(\Delta\)), so a constant
  fraction of designated \texttt{(p,u)} change; caches built under a
  previous seed chain provably miss \(\Theta\)(\textbar S(P)\textbar)
  locations, preserving the semantic cost.
\item
  Indices and skips. The TM may build auxiliary indices on W (e.g., a
  sorted map from \texttt{(p,u)} to tape positions) to accelerate
  lookups. These indices are internal artifacts and import no instance
  bits; they cannot reduce the semantic obligation q + \(\Phi\) \(\geq\)
  R, only the scanning overhead.
\end{itemize}

Bottom line. Even without hardware "addresses," k-tape TMs realize the
address-pool abstraction via parseable record names. Hermeticity
constrains all legitimate information to these payloads; Injectivity
prevents legal merging; RWA pins down first-use credit. The Semantic
Conservation Law then applies verbatim: if q falls short of R at a
bottleneck cut, the algorithm must maintain 2\^{}(R\(-\)q)
simultaneously distinguishable artifacts or pay the time to resolve
them.

\paragraph{D.5.1 Scanning overhead: conservative lower bound and a
streaming
bound}\label{d51-scanning-overhead-conservative-lower-bound-and-a-streaming-bound}

\begin{itemize}
\tightlist
\item
  Conservative lower bound (used in main theorems). We price time solely
  by first-use information inflow: each step imports \(\leq\) B =
  k\(\lceil\)log\(_2\)\textbar{}\(\Gamma\)\textbar{}\(\rceil\) fresh
  bits, so learning \(\lambda\) bits requires \(\geq\) \(\lambda\)/B
  steps. Combined with the across-tries inequality, this yields the
  restart-lane bound time \(\geq\) 2\^{}(\(\Delta\)(C*))/B (Appendix
  D.2). We deliberately do not add head-movement costs here; thus the
  bound is conservative on TMs.
\item
  Streaming retrieval (operational upper bound on overhead). Given a
  finite set S of target keys (p,u) for a segment, a 2-tape TM can
  compute S on W and then scan M once left-to-right, reading the first
  payload symbol at each encountered key. This takes
  O(\textbar{}M\textbar{} + \textbar{}S\textbar{}) head moves beyond the
  cost of reading payload symbols, not
  O(\textbar{}M\textbar{}·\textbar{}S\textbar{}). Here
  \textbar{}M\textbar{} is the input length \textbar{}x*\textbar{} (see
  §1.7 "Polynomial baselines"), i.e., \textbar{}M\textbar{} =
  \(\Theta\)(n\_core · log² n\_core). Hence per segment, scanning
  overhead is at most linear in input size plus hits.
\item
  Strengthened single-run bound (already quantified elsewhere). Our
  FG/segment analysis (Appendix C) supplies a per-segment baseline
  independent of scanning: each segment must process
  \(\Omega\)(n/W\_min) designated primitives (GateDigest parity +
  address-churn; Cor. C.1.1), yielding time \(\geq\)
  2\^{}(\(\rho\)-s)·\(\Omega\)(n/W\_min) in the single-run lane (see §9
  and Appendix D.4). Sequential scans can only increase this.
\item
  Takeaway. The paper\textquotesingle{}s time lower bounds do not rely
  on pessimistic head-movement; they use (i) an information-inflow
  baseline (restart lane) and (ii) explicit per-segment work (FG lane).
  Any additional scanning needed by a TM is a multiplicative polynomial
  factor that only strengthens the lower bounds and never weakens them.
\end{itemize}

\begin{center}\rule{0.5\linewidth}{0.5pt}\end{center}

\subsection{Appendix E: NP as Selection under Constraints
(expository)}\label{appendix-e-np-as-selection-under-constraints-expository}

This appendix offers an expository lens: many NP problems can be viewed
as selection under constraints. The aim is intuition, not new formal
theorems. The mapping highlights why verification is additive while
search may need to manage many possibilities when avoiding commitment.

\textbf{Why this matters:} Provides conceptual bridge from familiar NP
problems (3-SAT, CLIQUE, TSP) to L*\textquotesingle{}s structural
framework (emergence, dependency, completeness), helping readers
understand how L* isolates and makes explicit the features that create
the verification/search gap.

\subsubsection{E.1 Subset Selection: Choose or Track 2\^{}m
Possibilities}\label{e1-subset-selection-choose-or-track-2m-possibilities}

\begin{itemize}
\item
  \textbf{3-SAT:} Select assignment z \(\in\) \{0,1\}\^{}m

  \begin{itemize}
  \tightlist
  \item
    Local constraints: each clause tests 3 literals
  \item
    \textbf{Heuristic mapping:} Avoiding early commitments can require
    tracking 2\^{}(m \(-\) committed) assignments
  \end{itemize}
\item
  \textbf{CLIQUE:} Select k vertices via indicator z

  \begin{itemize}
  \tightlist
  \item
    Global constraint: all pairs must have edges
  \item
    \textbf{Heuristic mapping:} Partial selection admits many compatible
    extensions
  \end{itemize}
\end{itemize}

\subsubsection{E.2 Ordering: n! Permutations Force Multiplicative
Growth}\label{e2-ordering-n-permutations-force-multiplicative-growth}

\begin{itemize}
\tightlist
\item
  \textbf{Hamiltonian Path/TSP:} Select permutation \(\pi\) of n
  vertices

  \begin{itemize}
  \tightlist
  \item
    Constraints: valid edges and subtour elimination
  \item
    \textbf{Heuristic mapping:} Partial paths admit many compatible
    extensions (factorial growth in worst case)
  \end{itemize}
\end{itemize}

\subsubsection{E.3 Labeling: Combinatorial Choices Create Structural
Requirements}\label{e3-labeling-combinatorial-choices-create-structural-requirements}

\begin{itemize}
\item
  \textbf{3-COLORING:} Assign \(\ell\)(v) \(\in\) \{1,2,3\} to each
  vertex

  \begin{itemize}
  \tightlist
  \item
    Local constraints: adjacent vertices differ
  \item
    \textbf{Heuristic mapping:} Partial colorings admit many valid
    completions
  \end{itemize}
\item
  \textbf{Exact 3-Cover:} Select disjoint 3-sets covering all elements

  \begin{itemize}
  \tightlist
  \item
    Global constraint: each element covered exactly once
  \item
    \textbf{Heuristic mapping:} Partial covers admit many compatible
    extensions
  \end{itemize}
\end{itemize}

\subsubsection{E.4 The overlay: making structure
explicit}\label{e4-the-overlay-making-structure-explicit}

\textbf{How L* makes the structure explicit:}

\begin{itemize}
\item
  \textbf{Primitive checks:} Each e\_v,\_j,\(\ell\) tests O(1) bits
  \(\to\) verification is additive (SCL: full commitment makes q\_v =
  R\_v so costs add)
\item
  \textbf{Selection matrix:} Sel\_v exposes R\_v independent
  requirements (SCL: identifies the "R" term - required fresh bits at
  node v)
\item
  \textbf{Completeness:} y\_v = H\_vx\_v requires full resolution (no
  partial knowledge) (SCL: correctness demands q\_v reach R\_v;
  otherwise residual R\_v\(-\)q\_v remains)
\item
  \textbf{Dependency:} Seed\_v chains encode sequential dependencies
  (SCL: dependencies compose across cuts - Q(C), \(\Lambda\)(C),
  \(\lambda\)(C) = \(\Lambda\)\(-\)Q)
\end{itemize}

Note. L* does not claim to represent all NP problems; it isolates
structural features (emergence, completeness, dependency) that also
appear in many NP settings.

\textbf{Connection to verification/search gap:} These features explain
why L* admits poly-time verification (with full witness, primitive
checks compose additively) but forces exponential search (without
commitment, dependency chains require maintaining 2\^{}\(\lambda\)
distinguishable seed-consistent worlds across the bottleneck cut, per
SCL in §7).

\subsubsection{E.5 Verification vs. search
(expository)}\label{e5-verification-vs-search-expository}

\begin{itemize}
\tightlist
\item
  \textbf{Verification (full commitment):} Small tests compose
  additively \(\to\) polynomial time
\item
  \textbf{Search (avoiding commitment):} Often requires managing many
  possibilities; under strong structural dependencies (as in L*), this
  becomes exponential
\end{itemize}

This section is heuristic and aims to connect familiar NP problems with
the structural lens developed for L*.

\begin{center}\rule{0.5\linewidth}{0.5pt}\end{center}

\subsection{Appendix F: Reviewer FAQ \& Common
Pitfalls}\label{appendix-f-reviewer-faq--common-pitfalls}

This FAQ addresses common clarifications and distinctions. Answers
emphasize conservation-law obligations (what correctness imposes) rather
than algorithmic limitations.

\textbf{Organization:} F.1 (Core Conceptual) - quantifier structures,
instance-specific dependencies, configuration-space diversity; F.2
(Technical Distinctions) - artifact accounting, information-theoretic
bounds, TM inflow limits; F.3 (Scope and Universality) - algorithm
coverage, hybrid solvers, algebraic systems; F.4 (Proof Structure) -
where to find proofs in main text.

Notation reminder: \(\lambda\)(A,x), \(\Delta\)(C*), \(\rho\) and s
follow §2.1/§8; see also Appendix C for segment counting and FG.

\subsubsection{F.1 Core Conceptual
Questions}\label{f1-core-conceptual-questions}

\textbf{Topics:} Quantifier structure (\(\forall\)x* \(\forall\)A
uniform), instance-specific dependencies and why shortcuts
don\textquotesingle{}t transfer, configuration-space diversity
(2\^{}\(\lambda\) seed chains vs. poly address pools), completeness/salt
assumptions, randomized solvers.

\textbf{Q: Why not a universal RAM/TM exponential?} \textbf{A:} L* shows
that exponential behavior arises from conservation principles, not
machine limitations (restart vs single-run lanes; §7.3, Appendix C).
While a RAM can verify each node efficiently (additive cost), L*
rigorously proves that \ensuremath{\gtrsim} 2\^{}(\(\lambda\)(A,x)) distinct explorations
are mathematically necessary through its carefully crafted structure.

\textbf{Q: Do you claim a single worst-case instance hard for all
algorithms?} \textbf{A:} For FG-wired instances, we prove
\(\forall\)x*\(\in\)L*\_\{FG\}, \(\forall\) \textbf{uniform} algorithm A
\(\to\) super-poly time. This claim ("every instance hard for every
uniform algorithm") is valid because uniform algorithms cannot hardcode
instance-specific solutions. This enables the Structural OWF
construction (§9) and yields P \(\neq\) NP via the classical bridge. The
uniform restriction is what allows the stronger "every instance"
quantifier and enables the Structural OWF construction. See §12.9 for
full details on the quantifier structure.

\textbf{Q: Does Completeness=I\_R\_v weaken generality?} A: No; any
full-rank H\_v works (Lemma 6.2). Identity keeps rank-forcing
transparent.

\textbf{Q: Are salts cryptographic?} A: No; they are published explicit
constants (AAS; Appendix A.2). Emergence is information-theoretic.

\textbf{Q: What if R\_v=0 for some nodes?} A: Then Alt\_v=1; those nodes
contribute neither commits nor multiplication, consistent with our
theorems.

\textbf{Q: Randomized solvers?} A: Distributional (Yao) arguments apply;
some fixed input attains the bound (Yao {[}YAO77{]}, §9.4 coin-fixing).

\textbf{Q: Why can\textquotesingle{}t we find a shortcut in a different
run and apply it to the current instance?} \textbf{A:} Large \(\lambda\)
isn\textquotesingle{}t a runtime limitation - it\textquotesingle{}s a
structural property that persists even with unlimited pre-computation.
Here\textquotesingle{}s why:

\begin{enumerate}
\def\labelenumi{\arabic{enumi}.}
\item
  \textbf{Each instance has unique constraint structure:} Every L*
  instance has different randomly-chosen salts, creating a unique
  constraint topology. Techniques for instance A (which depend on
  A\textquotesingle{}s specific salts) provide no help for instance B
  (which has different salts).
\item
  \textbf{Structure emerges only during computation:} The constraints at
  node v depend on Seed\_v = Enc(parent\_results). These seeds
  don\textquotesingle{}t exist until parent computations complete, so
  the constraint structure literally doesn\textquotesingle{}t exist to
  analyze until you\textquotesingle{}re already computing it.
\item
  \textbf{No universal shortcut exists:} A universal method that works
  for ALL possible salt/seed combinations would need to represent
  2\^{}\(\lambda\) distinct states in fewer states. This violates the
  pigeonhole principle - mathematically impossible for maintaining
  correctness.
\item
  \textbf{Instance-specific dependencies:} Even memorizing all
  2\^{}(\(\Theta\)(log² n)) possible seed\(\to\)result mappings
  doesn\textquotesingle{}t help, as this memorization itself requires
  n\^{}(\(\Theta\)(log n\_core)) space - exactly the complexity
  we\textquotesingle{}re proving!
\end{enumerate}

The conservation law q + \(\Phi\) \(\geq\) R captures a mathematical
necessity: you cannot represent 2\^{}\(\lambda\) distinct constraint
configurations in fewer than 2\^{}\(\lambda\) states without collision.
This isn\textquotesingle{}t about finding a clever strategy -
it\textquotesingle{}s about the mathematical impossibility of
representing instance-specific, dynamically-revealed constraint
structures below their information-theoretic minimum.

\textbf{Q: With only polynomial address pools, how can complexity be
exponential?} \textbf{A:} Exponentiality emerges from
\textbf{configuration-space diversity}, not resource size. The
exponential requirement comes from the number of DISTINCT PERMUTATIONS
(configuration space: 2\^{}\(\lambda\) seed chains \(\to\)
2\^{}\(\lambda\) distinct address-selection patterns), not the number of
AVAILABLE ADDRESSES (resource space: poly(n\_core) cells total).

Think of it like a piano: 88 keys (polynomial) but infinitely many songs
(exponential patterns over those keys). L* has polynomial-sized address
pools U\_v, but 2\^{}\(\lambda\) distinct seed-dependent permutations
determining which addresses to read for each seed chain. You cannot
reuse artifacts across different seeds because each seed induces a
substantially different address pattern (full-churn property, Lemma
A.1.F).

\textbf{See §2.7 for the complete explanation with the piano analogy,
technical details, and why this architectural insight enables the P
\(\neq\) NP proof.} Additional formal details in §6.2.3 (seed-dependent
permutations) and Appendix A.1.1 (explicit construction).

\subsubsection{F.2 Technical
Distinctions}\label{f2-technical-distinctions}

\textbf{Topics:} Artifact accounting (semantic obligations vs. machine
steps), SCL (q + \(\Phi\) \(\geq\) R) preventing shortcuts,
seed-dependent addressing, try-and-verify analysis (rollback keys,
across-tries inequality), information-theoretic time bounds (B bits/step
inflow limits), distributional bounds and Yao/minimax.

\textbf{Q: Why is this different from typical circuit/communication
lower bounds?} \textbf{A:} We count semantic obligations
(distinguishable states/artifacts) rather than machine steps.
Exponential behavior comes from the need to distinguish
2\^{}(\(\lambda\)(A,x)) cases, not from I/O or circuit depth per se.

\textbf{Q: Can a clever algorithm avoid both reading and maintaining
distinguishable artifacts?} \textbf{A:} No. The Semantic Conservation
Law formalizes that either information is resolved or unresolved cases
are maintained; otherwise correctness fails on some instances.

\textbf{Q: But the algorithm has the entire input -
can\textquotesingle{}t it just deduce the intermediate values?}
\textbf{A:} The input contains all bytes, but which bytes belong to node
v is determined by Seed\_v, which depends on parent outputs. The
addresses for v are undefined until parents are computed - building that
dependency chain is the computation.

\textbf{Q: "Try-and-verify" seems polynomial per try - does that evade
the bound?} \textbf{A:} No. Tries are counted via rollback keys
(resolution prefixes). With insufficient commitment, \ensuremath{\mathbb{E}}{[}tries{]}
\(\geq\) 2\^{}(\(\Delta\)(C*)) (across-tries; restart lane), and
per-step inflow limits (e.g., §5.5.1) convert this into time.

\textbf{Q: How do you get time bounds without per-try costs?} A: Our
arity-bounded analysis (§5.5.1) uses information theory. Each TM step
can acquire at most B =
k\(\lceil\)log\(_2\)\textbar{}\(\Gamma\)\textbar{}\(\rceil\) bits. With
\(\Sigma\)\_v R\_v = \(\Theta\)(n log n) total information needed and
L*\textquotesingle{}s structure requiring
2\^{}(\(\Omega\)(\(\lambda\)(A,x))) explorations on FG-wired instances,
time \(\geq\) 2\^{}(\(\Omega\)(\(\lambda\)(A,x))) up to polynomial
factors. For QP-sharp (\(\lambda\)\_base = \(\Theta\)(log² n)), this
gives n\^{}(\(\Omega\)(log n\_core)); for exponential profiles
(\(\lambda\)\_base = \(\Theta\)(n)), it gives 2\^{}(\(\Omega\)(n)).

\textbf{Q: Is this a time lower bound for every RAM/TM?} \textbf{A:} For
L*, distributional bounds give \ensuremath{\mathbb{E}}{[}tries{]} \(\geq\)
2\^{}(\(\lambda\)(A,x) - log\(_2\) Alt(C*)) (in particular, \(\geq\)
2\^{}(\(\lambda\)(A,x))/poly(n\_core) if Alt(C*) \(\leq\)
poly(n\_core)); per-step inflow bounds give time \(\geq\)
2\^{}(\(\lambda\)(A,x) - log\(_2\) Alt(C*))/poly(n\_core). There exists
a fixed explicit instance achieving this (Yao/minimax).
Profile-dependent exponents follow (e.g., QP-sharp yields
n\^{}(\(\Omega\)(log n\_core))).

\subsubsection{F.3 Scope and
Universality}\label{f3-scope-and-universality}

\textbf{Topics:} Coverage of classical algorithms, hybrid solvers
(composite artifact tuples, hybrid potential \ensuremath{\mathcal{H}}\_v), unrestricted
branching programs (residual subfunctions, seed-separation), algebraic
proof systems (PC/CP/SoS degree/size bounds, Assumption A).

\textbf{Q: Does this cover all classical search algorithms on L*?}
\textbf{A:} Yes. Any correct classical solver faces the same tradeoff:
resolve information (additive) or maintain exponentially many states
(multiplicative) on our instances.

\textbf{Q: What about \textbf{hybrid} solvers that mix branching, DP,
OBDD, and resolution?} A: We count the \textbf{composite tuple of
artifact identifiers at each node; correctness requirements imply
injectivity from worlds to tuples, so the number of tuples is \(\geq\)
2\^{}(R\_v\(-\)q\_v). Defining the hybrid potential} \ensuremath{\mathcal{H}}\_v :=
log\(_2\)(branches\_v)+log\(_2\)(keys\_v)+log\(_2\)(width\_OBDD\_v)+log\(_2\)(cases\_RES\_v)
gives the per-node inequality \textbf{q\_v + \ensuremath{\mathcal{H}}\_v \(\geq\) R\_v} and
multiplicative growth across cuts (by dependency). This yields the same
exponential tradeoff curve for hybrids.

\textbf{Q: What about unrestricted branching programs (no fixed order,
read-many)?} A: We use \textbf{residual subfunctions as the semantic
artifact. For the order-robust parity gate, every node forces
2\^{}(\(\Omega\)(R\_v-q\_v)) distinct residual functions at some cut
(Lemma B.2). Injective dependency implies seed-separation}, so these
subfunctions cannot be merged across nodes; the program size is thus
\(\geq\) 2\^{}(\(\Omega\)(\(\lambda\)(A,x))).

\textbf{Q: And stronger algebraic systems (PC/CP/SoS)?} A: We use
\textbf{degree} (or rank/width surrogates) as the semantic artifact.
Each node with s\_v unresolved bits forces degree \(\Omega\)(s\_v) (see
Appendix H, Read-or-Degree) - this is \textbf{unconditional}. For
\textbf{size} bounds, we need the degree\(\to\)size tradeoff (Assumption
A); once established for our parity/expander family (see Appendix H),
size bounds become unconditional as well.

\subsubsection{F.4 Proof Structure and
Location}\label{f4-proof-structure-and-location}

\textbf{Q: Where are the proofs?} \textbf{A:} Proofs are organized as
follows:

\begin{itemize}
\tightlist
\item
  §7: Semantic Conservation Law (core theorem about unavoidable
  requirements)
\item
  §8: Per-instance deterministic bounds (§8.A)
\item
  Appendix B: Paradigm-specific manifestations
\item
  Appendix H: Algebraic system extensions Each shows how
  L*\textquotesingle{}s structure imposes the stated requirements.
\end{itemize}

\textbf{Q: Does any of this rely on bit-level syntactic tricks?}
\textbf{A:} No. We count distinguishable states/artifacts (semantic),
not syntactic machine artifacts. Emergence, Completeness, and Dependency
are instance-side properties.

\begin{itemize}
\tightlist
\item
\end{itemize}

\textbf{Q: Aren\textquotesingle{}t "dependency/charging/keyedness"
non-standard for NP?} \textbf{A:} They are analysis conventions to track
where and when information is first used and how reuse remains correct
(keyed). The instance remains standard NP with polynomial verification.

\begin{center}\rule{0.5\linewidth}{0.5pt}\end{center}

\subsection{Appendix G: Resolution Width-to-Size on
Parity/Expander}\label{appendix-g-resolution-width-to-size-on-parityexpander}

This appendix records the standard width-to-size phenomenon for
Resolution under the expander-parity gate used in our construction, and
how it instantiates the read-or-maintain tradeoff.

\textbf{Purpose:} Establishes Resolution-specific manifestation of SCL
(§7): uncommitted bits force width \(\geq\) s \(\to\) size \(\geq\)
2\^{}(\(\Omega\)(s)) via Ben-Sasson-Wigderson 2001 bridge. Referenced
from Corollary 5.3.2 part 4 (Resolution paradigm bounds). Part of
complete paradigm coverage demonstrating SCL\textquotesingle{}s
universality across proof systems.

\textbf{Theorem G.1 (Width-to-size under expander-parity).} Hypotheses:
expander-parity/Tseitin encoding at node v (as in Corollary 5.3.2 part
4; Appendix B), width measured on the induced parity CNF fragment.

For L* with s = R\_v \(-\) q\_v uncommitted bits at node v under this
encoding, any Resolution refutation requires width \(\geq\) s and size
\(\geq\) 2\^{}(\(\Omega\)(s)) clauses.

\subsubsection{G.1 The Structural
Requirement}\label{g1-the-structural-requirement}

\textbf{Setup:} After committing q\_v bits at node v, we have a CNF F
encoding parity constraints on s = R\_v - q\_v uncommitted bits.
Resolution must refute all 2\^{}s possible assignments.

\textbf{Theorem G.2 (Width forces size).} If a Resolution refutation
must use width \(\geq\) s, then its size is \(\geq\)
2\^{}(\(\Omega\)(s)) (Ben-Sasson-Wigderson).

\subsubsection{G.2 Why Width is
Unavoidable}\label{g2-why-width-is-unavoidable}

\textbf{Lemma G.3 (Width necessity under parity).} For parity
constraints on s uncommitted bits, any Resolution refutation requires
width \(\geq\) s.

\emph{Proof sketch.} Clauses of width \textless{} s cannot distinguish
all 2\^{}s assignments; the property is preserved under Resolution
inferences. Therefore reaching contradiction requires width s.

\textbf{Why parity forces width:} The expander-parity structure over s
bits creates 2\^{}s mutually distinguishable assignments. Any clause
resolving fewer than s literals cannot exclude enough assignments to
make progress - distinguishing all 2\^{}s worlds requires width at least
s, directly instantiating the read-or-maintain tradeoff (resolve bits or
maintain exponential proof artifacts).

\subsubsection{G.3 From Width to Exponential
Size}\label{g3-from-width-to-exponential-size}

\textbf{Width\(\to\)size.} Once width s is necessary, standard
width\(\to\)size tradeoffs imply size \(\geq\) 2\^{}(\(\Omega\)(s)).

\textbf{Ben-Sasson-Wigderson 2001 {[}BEN01{]}:} Any Resolution
refutation with width w and n variables requires size \(\geq\)
2\^{}(\(\Omega\)(w²/n)). On our parity fragments, s =
\(\Theta\)(\#vars), so this yields size \(\geq\) 2\^{}(\(\Omega\)(s)) as
claimed in Theorem G.1.

\subsubsection{G.4 Application to L*}\label{g4-application-to-l}

\textbf{Instantiation for L*.}

\begin{enumerate}
\def\labelenumi{\arabic{enumi}.}
\tightlist
\item
  At each node v: s\_v = R\_v - q\_v uncommitted bits
\item
  Expander-parity constraints force width \(\Omega\)(s\_v)
\item
  Across nodes on a bottleneck cut C: width contributions \textbf{add}
  to give global width \(\Omega\)(\(\Sigma\)\_\{v\(\in\)C\} s\_v) =
  \(\Omega\)(\(\lambda\)(A,x))
\item
  Apply the Ben-Sasson-Wigderson width\(\to\)size bridge to obtain size
  \(\geq\) 2\^{}(\(\Omega\)(\(\Sigma\)\_\{v\(\in\)C\} s\_v)) =
  2\^{}(\(\Omega\)(\(\lambda\)(A,x)))
\end{enumerate}

The multiplicative growth intuition is captured by Alt and FrontierPeak
(product over nodes/cuts); in Resolution we first transfer that product
to a single global width parameter before applying the standard
width\(\to\)size bridge.

These bounds align with the SCL: skipping commitments (q\_v \textless{}
R\_v) manifests as increased width and size.

\subsubsection{G.5 Why This Result is
Fundamental}\label{g5-why-this-result-is-fundamental}

\textbf{Remarks.}

\begin{enumerate}
\def\labelenumi{\arabic{enumi}.}
\tightlist
\item
  The argument is semantic: it counts distinguishable assignments
  induced by the parity structure.
\item
  The same read-or-maintain dichotomy appears here as in other
  paradigms.
\end{enumerate}

\begin{center}\rule{0.5\linewidth}{0.5pt}\end{center}

\subsection{Appendix H: Algebraic Proof Systems
(Overview)}\label{appendix-h-algebraic-proof-systems-overview}

This appendix summarizes how the parity/expander structure induces
degree lower bounds in algebraic proof systems (Polynomial Calculus,
Cutting Planes, Sum-of-Squares), and how these align with the
read-or-maintain tradeoff.

\textbf{Purpose:} Extends paradigm coverage to algebraic proof systems
(beyond combinatorial systems like Resolution), demonstrating
SCL\textquotesingle{}s universality (§7): uncommitted bits force degree
\(\geq\) \(\Omega\)(s) unconditionally; degree\(\to\)size bridge yields
exponential size conditionally via Assumption A. Referenced from §7.6
(paradigm-specific instantiations). Shows read-or-maintain principle
holds across proof system families (combinatorial + algebraic).

\textbf{Assumption A (degree\(\to\)size bridge; system-specific).} For a
given algebraic proof system S over a field \ensuremath{\mathbb{F}} and a family of
tautologies T\_n, there exists a monotone function g such that any
S-proof of T\_n with maximal polynomial degree D has size \(\geq\) g(D).
Typical instantiations include exponential tradeoffs g(D) =
2\^{}(\(\Omega\)(D)) under standard restrictions: \(-\) Polynomial
Calculus (PC): degree\(\to\)size under bounded coefficient growth
{[}CEI96, IMP99{]} \(-\) Sum-of-Squares (SoS): degree\(\to\)rank/size
via moment matrices {[}GRI01, LAU03, SCH08{]} \(-\) Cutting Planes (CP):
degree/rank\(\to\)size for specific constraint families {[}PUD97{]}

Whenever size bounds are claimed for S, we label reliance on Assumption
A explicitly and state the applicable bridge.

\textbf{Theorem H.1 (Read-or-degree).} Hypotheses:
expander-parity/Tseitin encoding at node v (as in §7.6; Appendix B) over
a field \ensuremath{\mathbb{F}} with char(\ensuremath{\mathbb{F}}) \(\neq\) 2.

For L* with s = R\_v \(-\) q\_v uncommitted bits, any algebraic
refutation requires degree \(\geq\) \(\Omega\)(s).

\subsubsection{H.1 The Conservation Principle Manifested Algebraically
in
L*}\label{h1-the-conservation-principle-manifested-algebraically-in-l}

\textbf{Degree as a structural measure.} Algebraic systems differ
syntactically from Resolution, but the same structural choice appears:

\begin{itemize}
\tightlist
\item
  Low degree corresponds to high commitment (many variables fixed)
\item
  High degree reflects tracking many unresolved possibilities
\end{itemize}

\subsubsection{H.2 Degree Lower Bounds: The Mathematical
Inevitability}\label{h2-degree-lower-bounds-the-mathematical-inevitability}

\textbf{Lemma H.2 (Read-or-degree).} If s\_v = R\_v \(-\) q\_v bits
remain uncommitted at node v, any algebraic derivation requires degree
\(\geq\) \(\Omega\)(s\_v).

Independence source. At node v, the expander-parity construction induces
\(\Theta\)(s\_v) independent XOR constraints over disjoint variable sets
(Appendix B), so degree lower bounds add in the bottleneck form.

\emph{Proof sketch.} The s\_v uncommitted bits induce \(\Theta\)(s\_v)
independent XOR constraints; representing these necessitates degree
proportional to the number of independent constraints.

\textbf{Why XOR forces degree:} Each independent XOR constraint over k
variables requires a polynomial of degree \(\geq\) k to represent
(linear functions cannot encode parity over k bits). With
\(\Theta\)(s\_v) independent constraints, the algebraic representation
must track all 2\^{}(s\_v) distinguishable assignments, forcing degree
\(\geq\) \(\Omega\)(s\_v) - directly instantiating the read-or-maintain
tradeoff in algebraic form (resolve variables or maintain high-degree
polynomials).

\subsubsection{H.3 Composition Across the
DAG}\label{h3-composition-across-the-dag}

\textbf{Composition across the DAG (bottleneck form).}

\begin{itemize}
\tightlist
\item
  Variable sets at distinct nodes on a bottleneck cut C are disjoint.
\item
  Along C with \(\Sigma\)\_\{v\(\in\)C\} s\_v =
  \(\Omega\)(\(\lambda\)(A,x)), any refutation must simultaneously
  encode \(\Omega\)(\(\lambda\)(A,x)) independent XOR constraints across
  disjoint variable sets.
\item
  Consequently the \textbf{maximal} degree in the derivation satisfies D
  \(\geq\) \(\Omega\)(\(\lambda\)(A,x)).
\end{itemize}

Size lower bounds from degree apply only via a system-specific bridge
g(D) (Assumption A).

\subsubsection{H.4 From Degree to Size: The Final
Manifestation}\label{h4-from-degree-to-size-the-final-manifestation}

\textbf{From degree to size (conditional).} Under system-specific
degree\(\to\)size tradeoffs (Assumption A), degree D implies size
\(\geq\) 2\^{}(\(\Omega\)(D)). Instantiating with D =
\(\Omega\)(\(\lambda\)(A,x)) yields size \(\geq\)
2\^{}(\(\Omega\)(\(\lambda\)(A,x))), matching other paradigms.

\textbf{Known bridges:} Polynomial Calculus has degree\(\to\)size
tradeoffs under restricted coefficient growth {[}IMP99{]}.
Sum-of-Squares has degree\(\to\)rank/size bounds via moment matrix
analysis {[}GRI01, LAU03{]}. Cutting Planes degree\(\to\)size results
exist for specific constraint families. These bridges make size bounds
conditional on system-specific assumptions, unlike
Resolution\textquotesingle{}s unconditional width\(\to\)size
{[}BEN01{]}.

\subsubsection{H.5 System-Specific
Manifestations}\label{h5-system-specific-manifestations}

\textbf{System sketches.}

\begin{itemize}
\tightlist
\item
  \textbf{Polynomial Calculus {[}CEI96, IMP99{]}:} XOR constraints
  induce high degree; independent constraints drive degree growth.
  Degree lower bounds unconditional; size bounds conditional on
  coefficient growth restrictions.
\item
  \textbf{Cutting Planes {[}PUD97{]}:} Rank/degree measures reflect
  unavoidable information aggregation; size bounds follow from
  system-specific tradeoffs. Application to parity requires careful
  encoding.
\item
  \textbf{Sum-of-Squares {[}GRI01, LAU03, SCH08{]}:} Moment matrices
  capture correlations; degree lower bounds imply rank/size lower bounds
  via standard bridges (when available). Parity constraints force high
  SoS degree unconditionally.
\end{itemize}

\subsubsection{H.6 The Ultimate
Unification}\label{h6-the-ultimate-unification}

\textbf{Conclusion.} The parity/expander structure induces degree lower
bounds that align with the read-or-maintain view. Under
degree\(\to\)size bridges, this yields exponential size, consistent with
other paradigms.

\begin{center}\rule{0.5\linewidth}{0.5pt}\end{center}

\subsection{Appendix I: RAM/TM to Layered Branching
Programs}\label{appendix-i-ramtm-to-layered-branching-programs}

This appendix records the standard simulation of RAMs/TMs by layered
branching programs (tableau-style), which we use as a model hook to
relate machine computations to width/subfunction measures.

\textbf{Purpose:} Establishes SNF (Semantic Normal Form) as foundational
analysis framework making semantic obligations (q, Alt, \(\lambda\))
auditable from native TM/RAM transcripts. Bridges native runs to
paradigm-specific measures (width/subfunctions for OBDD/BP §5.4.1, keys
for DP §5.2, proof size for Resolution §5.3). Used throughout §5
paradigm bounds, §7 SCL proofs, §8 per-instance bounds. Enables uniform
lower bound arguments across computational models without model-specific
structural assumptions.

Scope note. Classical runs only; randomized solvers treated via fixed
coins (Yao). Layered branching programs preserve read-order along
levels.

\textbf{I.Main (Semantic\(\to\)Model Transfer).} Hypotheses:
deterministic k-tape TM (or randomized with fixed coins), layered BP
simulation (Lemma I.1), SNF ledger (RWA for q, Keyedness/Injectivity for
Alt), designated cuts.

Guarantees: compiling a native run of time T yields a layered BP of size
poly(T). For any cut C and any prefix \(\pi\) of the native run, ledger
counters satisfy monotonicity (Lemma I.2.4) and coincide with
model-native measures at designated levels (Lemma I.2.5): Q(C;\(\pi\))
and Alt(C;\(\pi\)) transfer to width/subfunctions. Therefore, SCL bounds
proved on native runs transfer to width/subfunction bounds in the
compiled BP.

\textbf{Lemma I.1 (Compilation to layered BP).} Any RAM/TM computation
on inputs of length n and time T can be simulated by a layered branching
program of size poly(T) (e.g., O(T·polylog T)), preserving read-order
along levels with constant fan-out.

\textbf{Use.} This allows width/subfunction lower bounds to apply to
general classical solvers via the compiled form, aligning with the
paradigm adapters in §5. It is a tool for connecting semantic bounds to
model-specific measures, not a universal claim about all models.

\textbf{Preservation of semantic counters (explicit).} We do not assume
tableau-style simulations preserve information-theoretic structure by
themselves. The SCL is proved directly on the native run with:

\begin{itemize}
\tightlist
\item
  \textbf{RWA} for first-use attribution of q,
\item
  \textbf{Keyedness + Injectivity} for distinguishability (Alt) across
  cuts, and
\item
  model-specific pricing (e.g., \textbf{probe-work} in Appendix C, or
  \textbf{B-bits/step} in §5.5) for time.
\end{itemize}

The SNF ledger records first-use events and frontier classes from the
transcript and is used to relate semantic counters to model-native
measures after compilation. By \textbf{Lemma I.2.4 (Monotonicity)},
compilation cannot reduce ledger-defined counters; \textbf{Lemma I.2.5
(Frontier equivalence)} equates frontier classes with width/subfunctions
at designated cuts. Thus lower bounds transfer without assuming
simulation preserves SCL semantics.

\textbf{Lemma 7.Ledger-Mon (restated).} For any transcript prefixes
\(\pi\) \ensuremath{\preceq} \(\pi\)' of a native TM run on x*, the SNF ledger counters
satisfy Q(C;\(\pi\)) \(\leq\) Q(C;\(\pi\)') and Alt(C;\(\pi\)) \(\leq\)
Alt(C;\(\pi\)') for every cut C. Moreover, compiling the run to a
layered branching program (Lemma I.1) and replaying the transcript
cannot decrease these counters at any designated cut. Consequently,
ledger-based SCL bounds proved on native runs are preserved under
compilation to model-native measures (width/subfunctions).

\begin{center}\rule{0.5\linewidth}{0.5pt}\end{center}

\subsubsection{I.2 Semantic Normal Form: Revealing Hidden
Structure}\label{i2-semantic-normal-form-revealing-hidden-structure}

\textbf{Making Structural Necessity Visible:} Every algorithm has an
implicit structure that our SNF transformation makes explicit. This
doesn\textquotesingle{}t change what the algorithm does - it REVEALS
what the algorithm must necessarily be doing to solve L*.

Clarification. SNF is purely observational: it records first-use events,
tags, and frontier classes from the transcript without enforcing
seed-consistent computation or altering acceptance. We use SNF only to
audit transcripts and relate runs to model-native measures; it is not a
forcing transformation.

\textbf{Lemma I.2 (SNF Reveals Necessity).} \emph{Every classical solver
has a normalized form that makes its conservation-law obligations
explicit:}

\textbf{The SNF Properties (What Structure Must Exist):}

\begin{enumerate}
\def\labelenumi{\arabic{enumi}.}
\tightlist
\item
  \textbf{Information Attribution:} Every bit learned is charged to
  where it\textquotesingle{}s STRUCTURALLY REQUIRED
\item
  \textbf{Dependency Tracking:} The algorithm MUST track dependencies
  (revealed by tags)
\item
  \textbf{State Distinction:} Merging is only safe when states are
  MATHEMATICALLY EQUIVALENT
\item
  \textbf{Balanced Cuts:} Width/subfunctions measured where structure
  FORCES distinction
\item
  \textbf{Monotone Requirements:} SNF never decreases structural
  counters - it reveals what was always necessary:

  \begin{itemize}
  \tightlist
  \item
    branches\^{}(SNF) \(\geq\) branches (forced exploration)
  \item
    keys\^{}(SNF) \(\geq\) keys (forced state tracking)
  \item
    width\^{}(SNF) \(\geq\) width (forced parallel states)
  \end{itemize}
\item
  \textbf{Efficiency Preserved:} Polynomial overhead - structure was
  always there
\end{enumerate}

\textbf{Why SNF helps.}

\begin{itemize}
\tightlist
\item
  Information flow is recorded at first valid use (attribution)
\item
  Dependencies are explicit (seed-derived addressing)
\item
  Unsafe merges are ruled out by keyedness/correctness
\item
  Cuts localize where unresolved bits accumulate
\item
  Counters are monotone under SNF instrumentation
\end{itemize}

\textbf{Bridges to paradigm-specific measures:} SNF ledger counters (q,
Alt, frontier classes) translate directly to model-native artifacts -
width/subfunctions for OBDD/BP (§5.4.1, Lemma I.2.5), keys for DP
(§5.2), clause width for Resolution (§5.3), branches for exhaustive
exploration (§5.4). Lemma I.2.4 (monotonicity) ensures SNF never
underestimates native measures; Theorem I.1 proves q\_RWA = q\_sem
(schedule-invariant attribution).

\paragraph{I.2.1 Formal Wrapper and
Guarantees}\label{i21-formal-wrapper-and-guarantees}

\textbf{Definition I.2.1 (SNF wrapper).} Given a classical solver A,
define A\^{}(SNF) that executes A unchanged and, after each step t,
appends to a ledger L\_t: (i) first-use events (addresses whose earliest
revealing access occurred at t), (ii) the current seed-chain tag (the
tuple of (Seed,y) along the active path), and (iii) the frontier
partition at designated cuts (equivalence classes under \(\equiv\) from
Lemma 4.2.2). A\^{}(SNF) neither modifies A\textquotesingle{}s control
flow nor its inputs.

Note (segment analysis across a bottleneck cut). For Segment Counting in
Appendix C.2, the "resolution prefix across C" is the projection of the
seed-chain tag to the designated bottleneck cut C augmented with the
cut-level fields (ConstraintDigest\_C, WorldCommit\_C) defined in
Appendix C.2.a. This extension is purely analytical and NP-verifiable;
it does not alter the native computation nor the global seed encoding.

\textbf{Lemma I.2.2 (Correctness and overhead).} For all inputs x,
Acc(A\^{}(SNF),x) = Acc(A,x). Moreover, time(A\^{}(SNF),x) \(\leq\)
poly(time(A,x)). \emph{Proof.} The wrapper is observational: it records
the transcript events and computes tags from the public overlay and the
transcript using RWA and Hermeticity, without altering
A\textquotesingle{}s state, inputs, or schedule. Computing tags and
classes is polynomial in the transcript prefix. In particular,
A\^{}(SNF) never evaluates any GateDigest\_v nor any primitive
e\_\{v,j,\(\ell\)\} that was not already revealed by A; the ledger is a
pure observer of A\textquotesingle{}s designated reads and public
metadata. \(\blacksquare\)

\textbf{Lemma I.2.3 (No future-peek).} L\_t is a function of the
transcript up to t and public overlay only. In particular, first-use
events are determined by RWA; seed-chain tags depend on ancestor
(Seed,y) via Enc; frontier classes depend on the current tag and
observed bits. \emph{Proof.} Hermeticity precludes hidden channels; RWA
charges only first revealing accesses; Enc is deterministic/parseable.
\(\blacksquare\)

\textbf{Lemma I.2.4 (Monotonicity of counters).} Let C be any semantic
counter measured from the ledger (branches, keys, width/subfunctions).
Then C\^{}(SNF)(t) \(\geq\) C\^{}(native)(t) for all t.

\textbf{Theorem I.1 (Ledger Soundness \& Schedule-Invariance).} Let
q\_sem(C;\(\pi\)) denote the number of node-v bits functionally
determined across cut C by transcript prefix \(\pi\) (semantic
definition), and let q\_RWA(C;\(\pi\)) denote the sum of first-use bits
attributed by the SNF ledger using Receiving-Window Attribution. Then
for any run and any prefix \(\pi\),

\begin{enumerate}
\def\labelenumi{\arabic{enumi})}
\tightlist
\item
  q\_RWA(C;\(\pi\)) = q\_sem(C;\(\pi\)) (schedule-invariant
  attribution), and
\item
  The per-step legitimate fresh-bit inflow is bounded by B =
  k\(\lceil\)log\(_2\)\textbar{}\(\Gamma\)\textbar{}\(\rceil\) on k-tape
  TMs (information-flow cap), so any increase of q\_RWA by \(\Delta\)
  requires at least \(\Delta\)/B steps in any try.
\end{enumerate}

\textbf{Proof (Expanded).} We prove both claims with careful attention
to the equivalence between semantic and operational definitions.

\textbf{Part (1): q\_RWA(C;\(\pi\)) = q\_sem(C;\(\pi\)) - Semantic and
Operational Equivalence}

We establish equality by proving both directions: q\_RWA \(\leq\) q\_sem
and q\_sem \(\leq\) q\_RWA.

\emph{Direction 1: q\_RWA \(\leq\) q\_sem (RWA does not overcount)}

Every bit credited by RWA corresponds to an actual designated read from
an address pool (by RWA operational definition, §2.1.3). By Hermeticity
(Axiom A1), all fresh instance information enters exclusively through
designated addresses; there are no hidden channels. Therefore:

\begin{itemize}
\tightlist
\item
  Each RWA-credited bit was observed in the transcript (read from a
  designated address)
\item
  Observed bits are by definition functionally determined by the
  transcript
\item
  Therefore q\_RWA \(\leq\) q\_sem (RWA counts only observed bits) \ensuremath{\square}
\end{itemize}

\emph{Direction 2: q\_sem \(\leq\) q\_RWA (RWA does not undercount)}

Suppose a bit b at node v is functionally determined by transcript
prefix \(\pi\). We show b must have been RWA-credited.

\textbf{Case 2.1}: b was directly read from designated address u =
F\_overlay(Seed\_v; j,\(\ell\))

\begin{itemize}
\tightlist
\item
  By RWA operational definition, this read is credited at first use
\item
  Therefore b is counted in q\_RWA {[}YES{]}
\end{itemize}

\textbf{Case 2.2}: b was logically deduced from other bits \{b\(_1\),
..., b\_k\} without direct read

\begin{itemize}
\item
  By Hermeticity (A1) and Lemma 4.2.I (overlay independence), the only
  sources of fresh information are: (a) Designated reads (RWA-credited
  by Case 2.1), OR (b) Public overlay structure (instance-independent,
  cannot determine instance-specific bits)
\item
  Any logical deduction uses observations from (a) or (b)
\item
  Since b is instance-specific and functionally determined, it must
  depend on (a)
\item
  By Lemma 6.1-ZMI (Emergence via Zero Mutual Information), no node-v
  bit can be functionally determined before a node-v discovery read;
  thus any determination of b requires first-use reads attributable to v
\item
  Therefore: the observations enabling the deduction of b were
  RWA-credited
\item
  By Constraint-Digest Tagging (CDT-3, Appendix C.2.a), growing the
  constraint set NF\_C to include implications of \{b\(_1\), ..., b\_k\}
  requires: \(\to\) New designated reads (Case 2.1, already counted), OR
  \(\to\) New gate evaluations costing \(\Omega\)(n/W\_min) (which
  themselves require designated reads per Lemma 5.5.1.c)
\item
  In all sub-cases, the functional determination of b is backed by
  RWA-credited observations
\item
  Therefore b\textquotesingle{}s determination does not exceed RWA
  credits {[}YES{]}
\end{itemize}

\textbf{Case 2.3}: b deduced from public overlay alone (hypothetical)

\begin{itemize}
\tightlist
\item
  Public overlay O is fixed before the run (Lemma 4.2.I) and is
  instance-independent
\item
  By Hermeticity, O does not encode designated payload bits
\item
  Therefore O cannot functionally determine instance-specific bits
\item
  Contradiction: b cannot be functionally determined this way {[}NO{]}
\end{itemize}

Combining Cases 2.1-2.3: Every functionally determined bit is backed by
RWA credits, so q\_sem \(\leq\) q\_RWA \ensuremath{\square}

\emph{Combining both directions: q\_RWA = q\_sem}

Since q\_RWA \(\leq\) q\_sem (Direction 1) and q\_sem \(\leq\) q\_RWA
(Direction 2), we have q\_RWA = q\_sem \ensuremath{\square}

\emph{Schedule Invariance}

Both q\_sem and q\_RWA are defined in terms of \textbf{which
observations} are made, not \textbf{when} they occur:

\begin{itemize}
\tightlist
\item
  q\_sem: Functionally determined = fixed by the set of observations in
  \(\pi\) (order-independent)
\item
  q\_RWA: Credited at first use; Lemma 6.1-RWA proves this is
  order-invariant
\end{itemize}

By Lemma I.2.3, first-use events are functions of the transcript prefix
and public overlay only (no future-peeking). Therefore, reordering
operations in the transcript (while preserving which bits are read) does
not change the set of first-use events, hence does not change q\_RWA.
Similarly, functional determination depends only on which bits are
observed, not their temporal order. Therefore both q\_sem and q\_RWA are
schedule-invariant, and their equality is preserved under read
reordering \ensuremath{\square}

\textbf{Part (2): Per-Step Inflow Bound - Converting Credits to Steps}

By Lemma 5.5.1 (per-step information inflow, TM), in any step a k-tape
TM with alphabet \(\Gamma\) can observe at most B :=
k\(\lceil\)log\(_2\)\textbar{}\(\Gamma\)\textbar{}\(\rceil\) fresh bits
from the instance. By Hermeticity (A1), designated reads are the only
channel for fresh instance information. By RWA definition (§2.1.3), each
credited bit corresponds to at least one actual designated read.
Therefore:

\begin{itemize}
\tightlist
\item
  Accruing Q RWA-credited bits requires reading Q designated symbols
\item
  Reading Q designated symbols requires at least \(\lceil\)Q/B\(\rceil\)
  steps (at most B symbols per step)
\item
  Therefore any increase of q\_RWA by \(\Delta\) requires at least
  \(\Delta\)/B steps
\end{itemize}

The bound is schedule-independent: prefetching or reordering reads does
not reduce the number of designated symbols that must be read (Lemma
6.1-RWA); it only changes when they are read. The step count depends on
how many designated symbols must be read, not on scheduling \ensuremath{\square}

\textbf{Conclusion}: Both claims are established. The semantic
definition (functional determination) and operational definition (RWA
credits) yield identical counts, and these counts price directly into
time via the per-step inflow bound. \(\blacksquare\) \emph{Proof.} The
ledger refines the state partition by distinguishing states that differ
on unresolved bits (Lemma 4.2.2). Refinement cannot reduce counts of
distinguishable classes; compiled width/subfunction measures (Lemma I.1)
are monotone under refinement. \(\blacksquare\)

\textbf{Lemma I.2.5 (Frontier = classes).} At any designated cut C and
time t, the number of frontier classes recorded in the ledger equals the
number of \(\equiv\)-classes \(\Pi\)\_t/\(\equiv\) that intersect the
cut. Therefore \textbar{}frontier\_t(C)\textbar{} \(\geq\)
2\^{}(\(\Sigma\)\_\{v\(\in\)C\}(R\_v-q\_v)) (Theorem J.1). \emph{Proof.}
By definition, two partial worlds that differ on unresolved bits across
C are in distinct \(\equiv\)-classes (Lemma 4.2.2), and conversely
classes that agree on all resolved bits across C induce identical
designated addresses at descendants. Theorem J.1 gives the lower bound.
\(\blacksquare\)

\textbf{Corollary I.2.6 (Transfer of bounds).} Any lower bound proved on
SNF counters applies to A up to polynomial overhead (by Lemma I.2.2 and
Lemma I.2.4).

\subsubsection{I.3 Summary}\label{i3-summary}

SNF exposes the semantic obligations already present in a run and makes
them auditable (first-use, keyedness, cuts). We use it as an analysis
normal form; it does not change accept/reject behavior.

\textbf{Used in:} §5 paradigm adapters (OBDD/BP width §5.4.1, DP keys
§5.2, Resolution width §5.3, exhaustive exploration §5.4), §7 SCL proofs
(Theorem 7.A per-node SCL, Theorem 7.B lane dichotomy), §8 per-instance
deterministic bounds (§8.A via FG + segment counting). SNF provides
uniform framework for proving lower bounds across computational models
without model-specific structural assumptions.

\begin{center}\rule{0.5\linewidth}{0.5pt}\end{center}

\subsection{Appendix J: DAG Min-Cut Lower
Bound}\label{appendix-j-dag-min-cut-lower-bound}

\textbf{Lemma J.0 (Closure).} The sink seeds deterministically parse and
recover every ancestor\textquotesingle{}s (Seed,y). Therefore any valid
sink induces a complete dependency \textbf{closure} containing a
parsable image of all cut nodes. This injectivity is the hinge for the
min-cut bound on distinguishable artifacts.

We prove that the DAG dependency bound extends from max-path to min-cut
characterization; often tighter on parallel DAGs.

\textbf{Purpose:} Establishes min-cut characterization for DAG lower
bounds, enabling tighter bounds on parallel DAG structures than max-path
approach. Proves Alt(C) \(\geq\) 2\^{}(cap(C)) via Cartesian factoring
(Theorem J.1-PROD, Lemma J.1-Cart under H1-H5 hypotheses: disjoint
pools, Enc injectivity, completeness, unique-neighbor, no
cross-coupling). Used in §7.2.1 consolidated SCL proof (cut composition
Theorem 7.A.1), §8.A per-instance deterministic bounds, Appendix C.2
segment counting. Min-cut captures parallel bottlenecks max-path misses;
coincides with max-path on simple chains.

\subsubsection{J.1 Setting}\label{j1-setting}

Let G = (V,E) be a finite DAG with roots S \(\subseteq\) V (in-degree 0)
and sinks T \(\subseteq\) V (out-degree 0). Each node v \(\in\) V
represents a node with:

\begin{itemize}
\tightlist
\item
  \textbf{Emergence rank} R\_v \(\geq\) 0 and \textbf{commitments} q\_v
  \(\in\) {[}0, R\_v{]}
\item
  \textbf{Full-rank gate} H\_v so y\_v = H\_vx\_v with
  \textbar{}y\_v\textbar{} = R\_v
\item
  \textbf{Parents} P(v) = \{u : (u,v) \(\in\) E\}
\end{itemize}

Seeds follow canonical content-addressed dependency (with gate wiring
for designated nodes):

\textbf{Seed\_v = Enc(v \textbar{}\textbar{} sort(\{(u, Seed\_u, y\_u) :
u \(\in\) P(v)\}) \textbar{}\textbar{} GateDigest\_v)}

where GateDigest\_v is empty when GREQ\_v=0 and equals the
seed-chain-bound identity digest when GREQ\_v=1 (see §6.2.8, Appendix
C.1.1).

where Enc is injective and parseable with lexicographically sorted
parents.

Define \textbf{residual capacity} c(v) := R\_v \(-\) q\_v \(\geq\) 0.

An \textbf{s-t cut} C \(\subseteq\) V is a node set intersecting every
root-to-sink path. Its capacity is cap(C) := \(\Sigma\)\_v\(\in\)C c(v).
Let \(\lambda\)(A,x) := min\{cap(C) : C is an s-t cut\} be the
run-dependent min-cut value.

\subsubsection{J.2 Main Theorem}\label{j2-main-theorem}

Notation (this appendix).

\begin{itemize}
\tightlist
\item
  \(\Lambda\)(C) := \(\Sigma\)\_\{v\(\in\)C\} R\_v (total requirement on
  cut C)
\item
  Q(C) := \(\Sigma\)\_\{v\(\in\)C\} q\_v (total resolved bits on C)
\item
  cap(C) := \(\Lambda\)(C) \(-\) Q(C) = \(\Sigma\)\_\{v\(\in\)C\}(R\_v
  \(-\) q\_v) (residual capacity of C)
\item
  Alt(C) := \(\prod\)\_\{v\(\in\)C\} Alt\_v (simultaneously
  distinguishable artifacts across C) (Equivalence: For consistency with
  §2.1, we also write \(\lambda\)(C) \(\equiv\) cap(C) and
  \(\lambda\)(A,x) := min\_C \(\lambda\)(C).)
\end{itemize}

\textbf{Theorem J.1 (DAG Min-Cut Read-or-X).} Under Emergence,
Completeness, canonical DAG dependency, Hermeticity, Keyedness, and RWA,
for any solver respecting parent dependencies:

\begin{enumerate}
\def\labelenumi{\arabic{enumi}.}
\item
  \textbf{(Counting/width paradigms; within-try bound)} For every s-t
  cut C, if Alt\_v counts within-try simultaneously distinguishable
  artifacts at node v (e.g., DP keys/OBDD nodes/Resolution width
  encodings), then \textbf{log\(_2\) Alt(C) \(\geq\) \(\Sigma\)\_\{v
  \(\in\) C\}(R\_v - q\_v) = cap(C)} and hence log\(_2\) Alt(C) \(\geq\)
  \(\lambda\)(A,x).
\item
  \textbf{(Serial/query paradigms; across-tries bound)} For every s-t
  cut C and any single try, \textbf{Q(C) + log\(_2\) Alt(C) \(\geq\)
  \(\Lambda\)(C)} and any residual deficit \(\Delta\) = \(\Lambda\)(C)
  \(-\) (Q(C) + log\(_2\) Alt(C)) \textgreater{} 0 forces

  \ensuremath{\mathbb{E}}{[}tries{]} \(\geq\) 2\^{}\(\Delta\).
\end{enumerate}

Here Alt(C) counts equivalence classes of simultaneously distinguishable
artifacts modulo representation (compression/coordinate changes); see
§4.2.x for the sufficient-statistic formulation.

Reducing Alt(C) requires paying commitments q\_v to decrease residual
capacities.

\textbf{Lemma J.1-Cart (Cut Cartesian Factorization).} Let C be an s-t
cut. Under the following instance-side hypotheses: H1 Disjoint
designated pools: U\_v \(\cap\) U\_\{v\textquotesingle{}\} =
\(\emptyset\) for all distinct v,v\textquotesingle{} \(\in\) C (per-node
disjoint address ranges). H2 Enc injectivity and parseability: Enc is
injective and length-delimited; Seed\_v = Enc(v \textbar{}\textbar{}
sort(\{(u,Seed\_u,y\_u):u\(\in\)P(v)\}) \textbar{}\textbar{}
GateDigest\_v). H3 Completeness: rank(H\_v) = R\_v so y\_v fully
determines the emergent bits at v. H4 Unique-Neighbor/no cross-coupling
across C: each designated atom used for v \(\in\) C depends only on
(Seed\_u,y\_u) from ancestors of v and does not share designated atoms
with v\textquotesingle{} \(\in\) C. H5 Realizability (worst-case lane):
every assignment to the unresolved bits across C extends to a valid
instance consistent with the fixed transcript prefix.

Then the feasible worlds consistent with a fixed transcript prefix
factor across the cut as a Cartesian product: \(\Pi\)\_C \(\cong\)
\(\prod\)\emph{\{v\(\in\)C\} S\_v, where S\_v is the set of feasible
settings of the unresolved coordinates at v. In particular,
multiplicativity across C is justified without probabilistic
independence, and for any family of local lower bounds Alt\_v \(\geq\)
2\^{}(R\_v-q\_v), we obtain Alt(C) = \(\prod\)}\{v\(\in\)C\} Alt\_v
\(\geq\) 2\^{}(\(\Sigma\)\_\{v\(\in\)C\}(R\_v-q\_v)).

\emph{Proof.} Disjointness (H1,H4) ensures that designated atoms
(addresses/bits) contributing to different v \(\in\) C are disjoint, so
choices for one v do not constrain choices for v\textquotesingle{}.
Injectivity/parseability (H2) and Completeness (H3) ensure seeds and
ancestor tuples are uniquely determined by local unresolved bits and the
transcript; thus local choices compose consistently. Realizability (H5)
guarantees that for any tuple of local choices there exists a completion
consistent with the transcript. Therefore, feasible worlds across the
cut factor as a Cartesian product, and multiplicativity follows.
\(\blacksquare\)

\textbf{Corollary J.1.d (Deterministic single-run form).} For
deterministic TMs, interpret "tries" as non-accepting rollback segments
that progress the final chain (Appendix C.2). Then any residual deficit
\(\Delta\) across a cut forces at least 2\^{}\(\Delta\) such segments.
With Frontier-Gate (§6.2.8), each segment requires \(\Omega\)(n/W\_min)
TM steps by Lemma 5.5.1.c + Lemma A.1.\(\Delta\) (computing a cut-gate
digest over \(\Theta\)(n/W\_min) terms), yielding the single-run time
lower bound used in Theorem 7.B.

\subsubsection{J.3 Proof}\label{j3-proof}

We prove (1); (2) follows by minimizing over C.

\textbf{Step 0: Realizability.} As in the chain case: Every setting of
the c(v) unresolved bits extends to a valid instance (Completeness via
Emergence)

\textbf{Step 1: Ancestor decodability from sink seeds.} Let Seed\_T =
(Seed\_v)\_t\(\in\)T be the tuple of sink seeds. Since Enc is parseable
and includes sorted parent triples (u, Seed\_u, \textbf{y\_u}), a
deterministic parser recovers the entire ancestor forest from Seed\_T.

\textbf{Critical point:} Our Enc stores (u, Seed\_u, \textbf{y\_u})
tuples verbatim; thus sink seeds contain a \textbf{fully explicit
(hashless) concatenation DAG} that deterministically recovers every
ancestor y\_v. This is essential for the injectivity argument.

\textbf{Recoverability Lemma:} For any node v that is an ancestor of
some sink, y\_v is a deterministic function of Seed\_T.

Since C is an s-t cut, every v \(\in\) C lies on some root-to-sink path,
hence is recoverable from Seed\_T.

\textbf{Step 2: Injectivity across the cut.} Let Xc be the tuple of
unresolved bits across C (\textbar{}Xc\textbar{} = cap(C)). Consider:

F: \{0,1\}\^{}(cap(C)) \(\to\) \{Seed\_T\}, F(x\_C) := Seed\_T when cut
bits are x\_C

\textbf{Lemma J.1-INJ (Cut Injectivity).} F is injective.

\emph{Proof:} If x\_c \(\neq\) x\textquotesingle{}\emph{c, then for some
v* \(\in\) C, y}\{v*\} \(\neq\) y\textquotesingle{}\emph{\{v*\}. By Step
1 and Lemma J.0 (recoverability), y}\{v*\} is deterministically
reconstructible from the \textbf{sink seeds}, so F(x\_c) \(\neq\)
F(x\textquotesingle{}\_c). \(\blacksquare\)

Therefore \textbar{}Im(F)\textbar{} \(\geq\) 2\^{}(cap(C)).

Cut-Factoring Checklist (at point of use). Product factoring across a
cut uses the following instance-side prerequisites:

\begin{itemize}
\tightlist
\item
  Disjoint designated address pools \{U\_v\} for v\(\in\)C (no
  cross-node aliasing on the cut; §6.2.3/Unique-Neighbor)
\item
  Injective, parseable Enc and recoverability from sink seeds (Lemma
  J.0)
\item
  Keyedness: artifacts must be seed-consistent (no merging across
  different seeds)
\item
  Hermeticity + RWA: no hidden channels; first-use attribution is sound
  These ensure feasible choices at different cut nodes are independent
  coordinates, yielding a Cartesian product.
\end{itemize}

\textbf{Theorem J.1-PROD (Cut Product Theorem - minimal hypotheses).} On
L* instances, fix a transcript prefix \(\pi\) and an s-t cut C. For each
v \(\in\) C let S\_v \(\subseteq\) \{0,1\}\^{}(R\_v-q\_v) be the set of
admissible assignments to v\textquotesingle{}s unresolved coordinates
consistent with \(\pi\). Let \(\Pi\)\_C(\(\pi\)) denote the set of
feasible worlds across C. Assume only:

H1) Disjoint designated pools: U\_v \(\cap\) U\_\{v'\} = \(\emptyset\)
for v \(\neq\) v' on C (no aliasing across cut nodes; §6.2.3).

H2) Unique-Neighbor mapping: designated addresses are of the parseable
form (pool-id, key), with pool-id = v; thus an address identifies its
unique cut node (instance-side property of F\_overlay; Appendix A.1).

H3) Enc injectivity + Closure: the sink seeds encode a parsable chain of
(Seed,y) tuples that is reconstructible by the verifier (Lemma J.0), and
different parent tuples yield different child Seed (Enc injective).

H4) Completeness/realizability: for each v \(\in\) C every assignment in
S\_v extends to at least one feasible world consistent with \(\pi\).

H5) No cross-coupling across the cut: post-horizon digests do not depend
on unresolved coordinates on C (Lemma 7.2.1-NC) and there are no
additional instance-side constraints that simultaneously restrict
unresolved coordinates of distinct nodes in C (Lemma J.1-NC below). This
prevents hidden correlations among choices across different v \(\in\) C.

Then \(\Pi\)\_C(\(\pi\)) \(\cong\) \(\prod\)\emph{\{v\(\in\)C\} S\_v via
an explicit bijection; in particular,
\textbar{}\(\Pi\)\_C(\(\pi\))\textbar{} = \(\prod\)}\{v\(\in\)C\}
\textbar{}S\_v\textbar{} = 2\^{}(\(\Sigma\)\_\{v\(\in\)C\}(R\_v-q\_v)).

\emph{Proof.} Define F: \(\prod\)\emph{\{v\(\in\)C\} S\_v \(\to\)
\(\Pi\)\_C(\(\pi\)) by, for each tuple s = (s\_v)}\{v\(\in\)C\}: set
each v\textquotesingle{}s unresolved local coordinates to s\_v; compute
y\_v = H\_v x\_v; propagate seeds via Enc to complete the transcript up
to the sink. By H4 realizability, the result is a feasible world; thus F
is well-defined.

Injectivity: If two tuples differ at v*, Completeness implies y\_\{v*\}
differs; by Enc injectivity, Seed\_\{v*\} differs; by H2/H1 the
designated addresses reachable from Seed\_\{v*\} include addresses in
pool U\_\{v*\} that cannot alias with any other pool\textquotesingle{}s
addresses. Therefore the induced sink seed differs (Lemma J.1-INJ), so
worlds are distinct.

Surjectivity: Let \(\omega\) \(\in\) \(\Pi\)\_C(\(\pi\)). By Closure
(H3) (Recoverability via Enc parseability; see §6.2.7 and Lemma J.0 in
Appendix J), parse the sink seed to recover all ancestor (Seed,y); by
H2/H1, for each v \(\in\) C recover exactly the unresolved local
coordinates s\_v \(\in\) S\_v from addresses in U\_v. Let s :=
(s\_v)\_\{v\(\in\)C\}. By Lemma J.1-REAL (global realizability), there
exists a feasible world realizing s; replaying Enc\textquotesingle{}s
propagation yields \(\omega\). Hence F is bijective. H5 (and Lemma
J.1-NC) ensure no hidden inter-v constraints couple the choices, so
every combination in \(\prod\) S\_v yields a feasible world.
\(\blacksquare\)

\textbf{Lemma J.1-REAL (Global realizability across a cut).} Fix an s-t
cut C and a transcript prefix \(\pi\). For each v \(\in\) C let S\_v be
the set of admissible local assignments consistent with \(\pi\) as in
Theorem J.1-PROD. Under Emergence, Completeness (rank(H\_v)=R\_v),
disjoint designated pools \{U\_v\}, and injective, parseable Enc, the
following holds:

For every tuple s = (s\_v)\_\{v\(\in\)C\} with s\_v \(\in\) S\_v for all
v, there exists a feasible world \(\omega\) \(\in\) \(\Pi\)\_C(\(\pi\))
that realizes all local unresolved coordinates at the cut according to
s.

\emph{Proof.} Work in topological order from the roots toward the sinks.
Keep all symbols already fixed by \(\pi\) unchanged. For each node u not
in C whose values are not fixed by \(\pi\) but are required to compute
descendant seeds, choose salts in its designated pool U\_u to realize
the values demanded by \(\pi\) (possible by Completeness/realizability
and disjoint pools). For each v \(\in\) C, choose salts in U\_v so that
x\_v encodes s\_v; since U\_v are pairwise disjoint across v \(\in\) C
(H1) and seed computation uses only parent (Seed,y) tuples via Enc (H3),
these choices are mutually independent and do not conflict. Gate digests
that are wired into seeds (when GREQ\_v=1) depend only on pre-horizon
indices (Lemma 7.2.1-NC). Given the prefix \(\pi\) and the tuple s =
(s\_v)\_\{v\(\in\)C\}, all such pre-horizon quantities are
deterministically determined by ancestor values (including any cut nodes
with GREQ=0 whose values are fixed by s and \(\pi\)) and thus introduce
no additional constraints that couple unresolved coordinates across
distinct nodes of C. Propagating seeds via Enc yields a well-formed sink
consistent with \(\pi\) and with the chosen (s\_v) on C. Therefore a
feasible world realizing s exists. \(\blacksquare\)

Interpretation (overlay vs satisfiability levels). The Cartesian
factoring statement concerns overlay-feasible worlds across C (level 1).
Any further satisfiability filtering (e.g., decoded \(\varphi\)
constraints) is handled separately by the verifier and does not
retroactively couple the unresolved coordinates across C in the overlay
semantics; the counting bound applies to the overlay level irrespective
of whether a given combination ultimately yields acceptance.

\textbf{Lemma J.1-NC (No cross-coupling across a cut).} Under the L*
overlay (disjoint address pools, Unique-Neighbor addressing, injective
Enc, Completeness, and FG wiring with the published horizon), there are
no instance-side constraints that simultaneously restrict unresolved
coordinates of two distinct nodes in a fixed s-t cut C beyond those
induced by parent\(\to\)child seed propagation and pre-horizon data. In
particular, post-horizon digests do not depend on unresolved coordinates
on C (Lemma J.1-NC), and all designated addresses that witness local
constraints at a node v \(\in\) C lie in U\_v and are determined solely
by Seed\_v and pre-horizon data. Consequently, for any transcript prefix
\(\pi\) the admissible choice sets \{S\_v\}\_v\(\in\)C are free of
hidden correlations across different nodes in C.

\emph{Proof.} By construction (Appendix C.1.1 and §6.2.8), any
GateDigest\_v wired into Seed\_v depends only on indices S(P(v)) drawn
from pre-horizon nodes (GREQ=0). Thus digest values do not involve
unresolved coordinates on C (Lemma 7.2.1-NC). Designated addresses are
of the form u\_\{v,j,\(\ell\)\} = F\_overlay(Seed\_v; j,\(\ell\)) and,
by Unique-Neighbor addressing, carry pool-id = v; hence reads that
constrain node-v coordinates always reside in U\_v and cannot alias with
U\_\{v'\} for v' \(\neq\) v (H1-H2). Enc is injective and parseable, so
seed propagation introduces no additional algebraic constraints among
unresolved coordinates across different cut nodes beyond the
deterministic dependence of child seeds on parent (Seed,y). Finally,
Completeness ensures that for each v \(\in\) C and any admissible local
setting s\_v \(\in\) S\_v there exist salts in U\_v realizing it,
independent of choices at other cut nodes. Therefore, there is no
instance-side mechanism that enforces simultaneous relations among
unresolved coordinates \{s\_v\}\_v\(\in\)C beyond pre-horizon
information, and the admissible sets \{S\_v\} are not coupled across
nodes in C. \(\blacksquare\)

\textbf{Corollary (Lemma J.1-Cart - alias).} Under the hypotheses of
Theorem J.1-PROD, the feasible worlds across a cut factor as a Cartesian
product: \textbar{}\(\Pi\)\_C(\(\pi\))\textbar{} =
\(\prod\)\_\{v\(\in\)C\} \textbar{}S\_v\textbar{}.

\textbf{Step 3: Three operational barriers (two-dimensional
constraint).} L*\textquotesingle{}s Emergence property ensures
unresolved bits are unknown pre-discovery. To handle 2\^{}(cap(C))
possibilities (cap(C) \(\equiv\) \(\lambda\)(C) by §2.4.2), mathematics
- not algorithms - forces the two-dimensional constraint q + \(\Phi\)
\(\geq\) R to be met via three operational barriers:

\begin{itemize}
\tightlist
\item
  \textbf{Resolution dimension}: Increase \(\Sigma\)\_v\(\in\)C q\_v via
  reading designated information (forward learning), or
\item
  \textbf{Elimination dimension}: Increase \(\Sigma\)\_\{v\(\in\)C\}
  q\_v via testing candidates across multiple tries (backward pruning),
  which (by SMP and \textbf{Yao\textquotesingle{}s principle}; see §9.4)
  entails \ensuremath{\mathbb{E}}{[}tries{]} \(\geq\) 2\^{}(\(\Delta\)) when \(\Delta\) =
  \(\Lambda\)(C) \(-\) (Q(C) + log\(_2\) Alt(C)) \textgreater{} 0, or
\item
  \textbf{Storage dimension}: Maintain 2\^{}(cap(C)) within-try
  distinguishable artifacts (parallel state tracking in counting/width
  paradigms).
\end{itemize}

L*\textquotesingle{}s structure creates the necessity in both forms:

Alt(C) \(\geq\) \textbar{}Im(F)\textbar{} \(\geq\)
2\^{}(\(\Sigma\)\_\{v\(\in\)C\} c(v)) = 2\^{}(cap(C)) (simultaneous
representation) or Q(C) + log\(_2\) Alt(C) + log\(_2\) \ensuremath{\mathbb{E}}{[}tries{]}
\(\geq\) \(\Lambda\)(C) (across tries). \(\blacksquare\)

\subsubsection{J.4 Remarks}\label{j4-remarks}

\begin{itemize}
\tightlist
\item
  The proof uses only existing semantic guards plus canonical parseable
  Enc
\item
  Min-cut and max-path are \textbf{incomparable} in general DAGs:
  min-cut often dominates on parallel structure; max-path captures long
  chains. They coincide on simple chains; otherwise either can be
  tighter.
\item
  For general DAGs with parallel branches, min-cut typically yields
  tighter bounds
\item
  The bound is model-agnostic and lifts through Yao\textquotesingle{}s
  principle exactly as chain proofs do
\end{itemize}

\textbf{Used in:} §7.2.1 consolidated SCL proof (Theorem 7.A.1 cut
composition: \(\Sigma\)\_\{v\(\in\)C\} log\(_2\)(Alt\_v) \(\geq\)
\(\lambda\)(C) via Theorem J.1-PROD Cartesian factoring), §8.A
per-instance deterministic bounds (min-cut residual \(\lambda\)(A,x)
determines time complexity via FG+SC), Appendix C.2 segment counting
(cap(C) = \(\Sigma\)\_\{v\(\in\)C\}(R\_v-q\_v) determines segment
multiplicity m\_seg \(\geq\) 2\^{}(\(\rho\)-s)). Theorem J.1 provides
both within-try bounds (counting/width paradigms) and across-tries
bounds (serial/query paradigms).

\begin{center}\rule{0.5\linewidth}{0.5pt}\end{center}

\subsection{Appendix K: Algebraic Proof Systems via Structural
Necessity}\label{appendix-k-algebraic-proof-systems-via-structural-necessity}

\textbf{Observation:} Algebraic proof systems - Cutting Planes,
Polynomial Calculus, Sum-of-Squares - face the same conservation-law
necessity from L*. The expander-parity gadget creates obligations that
any approach must account for; bounds are stated via standard
degree/size bridges where applicable.

\textbf{Purpose:} Proves structural necessity via explicit theorems for
specific algebraic systems (CP rank \(\geq\) \(\lambda\), PC degree
\(\geq\) \(\lambda\), SoS rounds \(\geq\) \(\lambda\)). Complements
Appendix H\textquotesingle{}s degree bounds overview and Assumption A
discussion with focused structural necessity proofs connecting to
Theorem J.1 min-cut residual. Used in §7.6 algebraic paradigm bounds.
Key distinction from H: K provides explicit necessity proofs (Theorems
K.1-K.3) showing rank/degree/rounds \(\geq\) min\_C \(\lambda\)(C); H
provides broader coverage with conditional degree\(\to\)size bridges via
Assumption A.

\subsubsection{K.1 Cutting Planes \& Polynomial
Calculus}\label{k1-cutting-planes--polynomial-calculus}

L*\textquotesingle{}s expander-parity gadget (§6.2.6) embeds independent
parity constraints at each cut C. For residual capacity \(\lambda\)(C) =
\(\Sigma\)\_\{v\(\in\)C\}(R\_v - q\_v), the gadget ensures
rank\_\(\mathcal{F}\)\(_2\)(C) = \(\Theta\)(\(\lambda\)(C)).

\textbf{Notation.} rank\_\(\mathcal{F}\)\(_2\)(C) denotes the rank over
the field \(\mathcal{F}\)\(_2\) = GF(2) of the system of XOR constraints
induced by the expander-FG gates at nodes in cut C. Since the expander
structure ensures \(\Theta\)(\(\lambda\)(C)) independent parity
constraints (Lemma B.1), we have rank\_\(\mathcal{F}\)\(_2\)(C) =
\(\Theta\)(\(\lambda\)(C)).

\textbf{Notation (rank\_\(\mathcal{F}\)\(_2\) and
\(\Theta\)-tightness).} We use rank\_\(\mathcal{F}\)\(_2\)(C) =
\(\Theta\)(\(\lambda\)(C)) intentionally: rank\_\(\mathcal{F}\)\(_2\)(C)
measures the GF(2) rank of the instance-fixed XOR constraint system at
cut C, while \(\lambda\)(C) = \(\Sigma\)\_\{v\(\in\)C\}(R\_v \(-\) q\_v)
is run-dependent via q\_v. The constructed system satisfies
c\(_1\)·\(\lambda\)(C) \(\leq\) rank\_\(\mathcal{F}\)\(_2\)(C) \(\leq\)
c\(_2\)·\(\lambda\)(C) for fixed constants c\(_1\),c\(_2\)
\textgreater{} 0 (from expansion and disjointness), making
\(\Theta\)(\(\lambda\)(C)) the precise characterization. Constant-factor
tightness suffices for Theorems K.1-K.3 (algebraic proof-system lower
bounds are robust to constant factors in rank/degree).

\textbf{Theorem K.1 (CP Structural Necessity).} Any Cutting Planes
refutation of L* requires rank \(\geq\) min\_C
rank\_\(\mathcal{F}\)\(_2\)(C) = \(\Theta\)(min\_C \(\lambda\)(C)).

\emph{Proof sketch}: Each of the \(\lambda\)(C) independent parities
creates a mod-2 constraint that Cutting Planes must maintain. Rank below
\(\lambda\)(C) cannot capture all constraints - feasibility persists. \ensuremath{\square}

\textbf{Theorem K.2 (PC Degree Necessity).} Over GF(2), any Polynomial
Calculus refutation requires degree \(\geq\) min\_C
rank\_\(\mathcal{F}\)\(_2\)(C) = \(\Theta\)(min\_C \(\lambda\)(C)).

\emph{Proof sketch}: The \(\lambda\)(C) independent XORs force degree
\(\geq\) \(\lambda\)(C) to generate the required polynomial conflicts.
L*\textquotesingle{}s structure makes this unavoidable. \ensuremath{\square}

\textbf{Key insight}: The bound min\_C \(\lambda\)(C) is EXACTLY the
min-cut residual from Theorem J.1 - different proof systems manifest the
SAME conservation-law obligation.

\subsubsection{K.2 Sum-of-Squares
Hierarchy}\label{k2-sum-of-squares-hierarchy}

\textbf{Theorem K.3 (SOS Rank Necessity).} For r = min\_C
rank\_\(\mathcal{F}\)\(_2\)(C), any k-round Sum-of-Squares refutation
with k \textless{} r cannot refute L*. Hence SOS rank \(\geq\) r.

\emph{Proof sketch}: L*\textquotesingle{}s XOR constraints require
r-wise correlations in any valid pseudo-distribution. With k \textless{}
r rounds, the SOS hierarchy cannot express these correlations - a
consistent pseudo-distribution exists, preventing refutation. \ensuremath{\square}

\textbf{Structural truth}: These bounds aren\textquotesingle{}t
limitations of algebraic methods - they\textquotesingle{}re mathematical
consequences of L*\textquotesingle{}s semantic structure forcing
rank/degree/rounds \(\geq\) min-cut residual capacity.

\textbf{Used in:} §7.6 algebraic paradigm bounds. Theorems K.1-K.3
provide explicit structural necessity proofs complementing Appendix
H\textquotesingle{}s overview - H covers degree bounds and Assumption A
for degree\(\to\)size bridges; K proves necessity via
rank\_\(\mathcal{F}\)\(_2\)(C) = \(\Theta\)(\(\lambda\)(C)) connecting
algebraic measures directly to min-cut residual (Theorem J.1).

\begin{center}\rule{0.5\linewidth}{0.5pt}\end{center}

\subsection{Appendix L: Provenance and Recoverability
Proofs}\label{appendix-l-provenance-and-recoverability-proofs}

\textbf{Purpose:} Proves Seed\_v encoding enables complete ancestor
recoverability from run transcripts (Theorem L.1), a structural
foundation for cut injectivity (Lemma J.1-INJ). Without provenance,
distinct seed-histories could collapse to identical sink states,
breaking SCL\textquotesingle{}s Alt(C*) \(\geq\) 2\^{}\(\lambda\) bound
- L*\textquotesingle{}s structure prevents this collapse.

\textbf{Observation:} L*\textquotesingle{}s dependency Seed\_v = Enc(v
\textbar{}\textbar{} sort\{(u, Seed\_u, y\_u)\} \textbar{}\textbar{}
GateDigest\_v) yields traceability of ancestor states via seeds. This
enables closure (Lemma J.0) and supports min-cut bounds.

\textbf{Theorem L.1 (Structural Provenance).} Given sink seeds Seed\_T
from a run\textquotesingle{}s transcript, L*\textquotesingle{}s
canonical encoding functionally recovers any ancestor node
v\textquotesingle{}s contribution (Seed\_v, y\_v). The y\_v values are
not published in x*; they are reconstructible from the transcript via
Enc\textquotesingle{}s parseable structure.

\textbf{Notation.} Let T denote the set of sink nodes (out-degree 0).
Seed\_T := (Seed\_t)\_\{t\(\in\)T\} denotes the tuple of all sink seeds.
For any specific sink t \(\in\) T, Seed\_t denotes its individual seed.

\emph{Proof sketch}:

\begin{enumerate}
\def\labelenumi{\arabic{enumi}.}
\tightlist
\item
  \textbf{Injectivity}: Enc is injective and length-delimited \(\to\)
  unique parsing of Seed\_t yields \{(u, Seed\_u, y\_u)\}
\item
  \textbf{Recursion}: Each parent tuple contains its full Seed\_u
  \(\to\) recursive parsing builds complete DAG
\item
  \textbf{Determinism}: Sorted parent lists + canonical encoding \(\to\)
  unique reconstruction path to any v
\item
  \textbf{Witness}: Concatenated parent-tuples along path from sink to v
  form O(depth·deg\textsuperscript{+}) proof
\end{enumerate}

\textbf{Key insight}: This isn\textquotesingle{}t an algorithmic
property - it\textquotesingle{}s mathematical necessity from
L*\textquotesingle{}s structure.

\textbf{Significance}: Provenance enables the cut injectivity (Lemma
J.1-INJ) proving Alt(C*) \(\geq\) 2\^{}(\(\lambda\)(A,x)) on a min-cut.
Without complete recoverability from the transcript, distinct cut
configurations could map to identical sink seeds, breaking the bound.
L*\textquotesingle{}s structure prevents this collapse.

\textbf{Complexity}: Witness size O(deg\textsuperscript{+}·depth); verification is
\textbf{linear in the witness size} (no hashing/Merkleization).

\textbf{Used in:} §7.2.1 (Consolidated SCL Theorem), Appendix J (min-cut
characterization). Key distinction from J: L proves the provenance
mechanism (Theorem L.1: Enc structure \(\to\) ancestor recoverability);
J applies provenance to min-cut framework (Theorem J.1 Cartesian
factoring, Lemma J.1-INJ cut injectivity).

\begin{center}\rule{0.5\linewidth}{0.5pt}\end{center}

\subsection{Appendix M: Glossary of Semantic
Terms}\label{appendix-m-glossary-of-semantic-terms}

\textbf{Purpose:} Quick reference for semantic terminology used
throughout §1-10 and appendices. Emphasizes conservation-law
requirements (what L*\textquotesingle{}s structure forces) over
algorithmic strategies (what algorithms choose). For detailed proofs and
mechanisms, see main sections and construction appendices.

\textbf{Note:} The terms refer to structural necessities (requirements
for correctness), not algorithmic strategies.

\textbf{Distinguishable artifacts maintained (at node v):} The number of
live partial worlds that L*\textquotesingle{}s semantic law FORCES - not
allows - the solver to maintain (Alt\_v classes). Manifests as branches
(EO), keys (DP), width (OBDD/Resolution). This isn\textquotesingle{}t
algorithmic choice; it\textquotesingle{}s a conservation-law obligation.

\textbf{Correctness requirement:} The mathematical impossibility of
merging distinguishable artifacts without acquiring distinguishing
information. L* creates this necessity through its semantic structure -
algorithms have no escape.

\textbf{Resolution:} Learning one fresh bit at node v (counts toward
q\_v). Each resolution reduces uncertainty but is a payment
L*\textquotesingle{}s structure demands, not a clever algorithmic move.

\textbf{Completeness:} L*\textquotesingle{}s enforcement that y\_v =
H\_vx\_v with rank(H\_v) = R\_v, ensuring intermediate states are fully
resolved before use. This models how real computations (compilation,
proofs) REQUIRE complete intermediates - a conservation-law obligation,
not design choice.

\textbf{Dependency:} L*\textquotesingle{}s DAG structure where Seed\_v =
Enc(v\textbar{}\textbar{}sort(\{(u,Seed\_u,y\_u):u\(\in\)P(v)\})
\textbar{}\textbar{} GateDigest\_v) creates unavoidable chains. Like git
commits or build dependencies, this isn\textquotesingle{}t convention -
it\textquotesingle{}s mathematical requirement.

\textbf{Emergence rank (R\_v):} The R\_v bits that MUST emerge at node v
- not given, but necessarily discovered through computation.
L*\textquotesingle{}s structure makes these bits unknowable until
computed.

\textbf{Node window:} A computation partition where
L*\textquotesingle{}s structure enforces min-cut measure \(\lambda\)\_W.
Not an analytical convenience but how L* creates unavoidable bottlenecks
(§5.5.1).

\textbf{LOAD (windowed accounting):} Each read L*\textquotesingle{}s
structure FORCES during computation. Re-loads count because
L*\textquotesingle{}s dependency chains make information non-reusable
across different seeds - a conservation-law obligation, not
inefficiency.

\textbf{Three-Dimensional Computational Framework:} The mathematical
structure L* enforces via q + \(\Phi\) \(\geq\) R (where \(\Phi\) =
log\(_2\)(Alt)). Three independent dimensions face exponential barriers:
\textbf{Storage (maintaining 2\^{}\(\lambda\) distinguishable states),
Resolution} (learning R\_v correct bits via reading), and
\textbf{Elimination (pruning 2\^{}(R\_v-q\_v) wrong candidates via
testing). Not algorithmic limitations but conservation-law requirements
- A1-A5} properties make all three exponentially hard simultaneously.

\textbf{Resolution (semantic):} L*\textquotesingle{}s requirement that
distinct cases remain separate for correctness. This
isn\textquotesingle{}t the Resolution proof system but the deeper
principle: L*\textquotesingle{}s structure forbids false merging.

\textbf{Residual subfunction (BP/OBDD):} After fixing a prefix, the
remaining function. L*\textquotesingle{}s structure ensures distinct
subfunctions can\textquotesingle{}t collapse - forcing width \(\geq\)
2\^{}(\(\lambda\)(A,x)) mathematically, not algorithmically.

\textbf{Semantic Conservation Law (SCL):} The INESCAPABLE principle q +
\(\Phi\) \(\geq\) R (where \(\Phi\) = log\(_2\)(Alt)).
L*\textquotesingle{}s semantic requirements R cannot vanish - they MUST
be satisfied through information (q) or multiplication (\(\Phi\)). This
isn\textquotesingle{}t what algorithms can\textquotesingle{}t do;
it\textquotesingle{}s what mathematics demands. Cut form:
\(\Sigma\)\_\{v\(\in\)C\} q\_v + log\(_2\) Alt(C) \(\geq\)
\(\Sigma\)\_\{v\(\in\)C\} R\_v (subscripts used only for summation).

\textbf{Semantic Multiplication Principle (SMP):} The mathematical truth
that learning only q of R bits forces 2\^{}(R-q) multiplicative growth.
Not algorithmic failure but a conservation-law obligation: Independent
Choices × Resolution \(\Rightarrow\) Multiplication. L* makes this
unavoidable.

\textbf{Usage:} Consult when encountering unfamiliar terms in main
sections (§1-10) or appendices. For detailed mechanisms: Appendix A
(A1-A5 properties), Appendix I (RWA accounting), Appendix J (min-cut
framework), Appendix L (provenance). For paradigm-specific
manifestations: Appendices C-K.

\begin{center}\rule{0.5\linewidth}{0.5pt}\end{center}

\subsection{Appendix N: Oracle and Barrier
Analysis}\label{appendix-n-oracle-and-barrier-analysis}

\textbf{Purpose:} Clarifies how this work relates to classical proof
barriers (Natural Proofs, Relativization, Algebrization). The
unconditional Structural OWF construction (§9) uses non-natural
properties (exponentially sparse), non-relativizing structure
(content-addressed dependencies), and non-algebrizing counting
(combinatorial min-cut). This appendix provides theoretical context; not
required for main theorems §7-10.

This appendix explains how our approach relates to classical proof
barriers and why the specific claims here do not conflict with them.

\subsubsection{N.1 Natural Proofs (sparsity and constructivity in the
truth-table
sense)}\label{n1-natural-proofs-sparsity-and-constructivity-in-the-truth-table-sense}

Let m denote the number of input variables of a Boolean function
F:\{0,1\}\^{}m\(\to\)\{0,1\}. In Razborov-Rudich, a property P on
functions is called "natural" if (i) P is large (holds for a
2\^{}(-O(1)) fraction of all 2\^{}(2\^{}m) functions), and (ii) P is
constructive: given the full truth table of F (length 2\^{}m), P(F) is
decidable in time poly(2\^{}m).

Definition (Overlay-compatibility property P\_m). P\_m(F)=1 iff there
exists an explicit overlaid instance x* of length poly(m) whose L*
acceptance behavior (over its DAG) matches the behavior induced by F on
the designated probes.

\textbf{Proposition N.1 (Non-naturalness of P\_m; informal).} There
exists c\textgreater{}0 such that for all sufficiently large m:

\begin{itemize}
\tightlist
\item
  (Sparsity) The density of \{F : P\_m(F)=1\} among all Boolean
  functions on m bits is at most 2\^{}(-c·2\^{}m).
\item
  (Constructivity not needed) Our lower bounds do not require P\_m to be
  constructive; sparsity alone precludes naturalness.
\end{itemize}

\emph{Sketch.} For fixed m, there are at most 2\^{}(poly(m)) explicit
overlays x* of length poly(m), and therefore at most 2\^{}(poly(m))
distinct induced behaviors on the designated probes. In contrast, there
are 2\^{}(2\^{}m) Boolean functions on m variables; hence the density
bound. Since non-naturalness already follows from sparsity, we do not
rely on constructivity. (Under standard PRF assumptions, combinatorial
properties that are both large and constructive cannot yield strong
circuit lower bounds; our property is exponentially sparse.)
\(\blacksquare\)

Note. This proposition provides context on barriers and is not used in
the proofs of the main theorems in §§7-10.

Implication. The property used here is not "natural" in the
Razborov-Rudich sense (fails largeness), so the barrier does not apply.
Note (distributional density). This claim concerns worst-case density
among all Boolean functions on m variables. We do not claim
non-naturalness under arbitrary induced distributions over functions; if
needed, such distributional claims should be analyzed separately.

\subsubsection{N.2 Relativization (Capacity-Respecting Oracles;
CRO-stability; non-relativizing by
construction)}\label{n2-relativization-capacity-respecting-oracles-cro-stability-non-relativizing-by-construction}

Two key observations:

\begin{itemize}
\item
  \textbf{Non-relativizing by construction}: L*\textquotesingle{}s
  content-addressed dependency Seed\_v = Enc(v \textbar{}\textbar{}
  sort(\{(u,Seed\_u,y\_u):u\(\in\)P(v)\}) \textbar{}\textbar{}
  GateDigest\_v) ties designated addresses to the
  instance\textquotesingle{}s internal transcript (parent (Seed,y)
  tuples). An arbitrary oracle could reveal unread information and thus
  change the problem; hence we do not claim oracle-robustness.
\item
  \textbf{CRO-stability}: Under Capacity-Respecting Oracles - formally,
  oracles O whose answers at step t are functions only of the current
  transcript prefix and public overlay, and which do not increase the
  mutual information with unresolved bits beyond the RWA-charged
  designated reads - L*\textquotesingle{}s structural necessities
  persist.
\end{itemize}

\textbf{Lemma N.2 (CRO-stability).} Under CRO oracles, Emergence,
Completeness, Dependency, and Theorem J.1 (min-cut) hold unchanged.

\emph{Proof.} By definition, a CRO cannot inject fresh information
beyond first-use reads (RWA) and cannot short-circuit dependency
(addresses remain functions of the current seed chain). Independence and
injectivity arguments across cuts are semantic and remain valid.
\(\blacksquare\)

Implication. The argument is non-relativizing in the Baker-Gill-Solovay
sense (we do not claim statements for arbitrary oracles), but it is
stable under oracles that respect the capacity/accounting model.

\textbf{Lemma N.2.b (CRO mutual-information budget).} Let X\_t denote
the unresolved bits at step t and O\_t the oracle\textquotesingle{}s
answer at step t. For a CRO, for every t, I(X\_t; O\_t \textbar{}
transcript\_\{\(\leq\) t\}) = 0, and the total mutual information gain
per step satisfies I(X\_t; (read symbols at t, O\_t) \textbar{}
transcript\_\{\textless{} t\}) \(\leq\) B, where B is the arity
parameter from Lemma 5.5.1.

\emph{Proof.} By definition of CRO, conditioned on the transcript
prefix, O\_t is a function of the prefix and public overlay only,
independent of X\_t; hence the first claim. The second follows by adding
the oracle output to the per-step read tuple and applying Lemma 5.5.1.
\(\blacksquare\)

\subsubsection{N.3 Algebrization (non-algebrization witness and
scope)}\label{n3-algebrization-non-algebrization-witness-and-scope}

\textbf{Witness N.3 (Degree growth under parity constraints; informal).}
The expander-parity overlay embeds \(\Omega\)(\(\lambda\)(C))
independent XOR constraints across the bottleneck cut. Any algebraic
representation that tracks behavior across the cut must have degree
\(\geq\) \(\Omega\)(\(\lambda\)(C)) over \ensuremath{\mathbb{F}}\(_2\).

\emph{Rationale.} Independent parity constraints force degree growth in
standard algebraic systems (e.g., Polynomial Calculus); this is
structural, not model-specific. Where degree\(\to\)size bridges are
known, we state them explicitly in the corresponding appendices.

Implication. Our counting argument is combinatorial and does not
algebrize. We do not claim unconditional degree\(\to\)size consequences
beyond the labeled bridges.

\textbf{Proposition N.3.a (Non-algebrization scope; informal).} Our
lower bounds are inherently non-algebrizing in the Aaronson-Wigderson
sense. Specifically, if one extends the model with oracle access to a
low-degree extension of the designated overlay, an oracle O\_LDE that on
input (addr, seed-chain prefix) returns E(addr, prefix) over a field \ensuremath{\mathbb{F}}
with deg(E) = polylog(n), then:

\begin{itemize}
\tightlist
\item
  Batched interpolation/sum-check style queries can aggregate designated
  bits across disjoint pools on a cut, effectively exceeding the
  per-step legitimate inflow B = O(1) enforced by RWA/Hermeticity.
\item
  Such access can collapse the residual \(\lambda\) priced by
  SCL\textquotesingle{}s designated-read accounting, invalidating the
  per-run pricing lemmas used for our time bounds.
\end{itemize}

Therefore, our SCL-based lower bounds do not claim to hold under
algebrized access; the technique is non-algebrizing by design. Our
claims are confined to classical, uniform models with Hermetic,
designated reads (no algebraic extension oracles).

\emph{Sketch.} Define O\_LDE for the overlay\textquotesingle{}s
designated address map A(s,·) so that for any fixed seed-chain prefix s,
the Boolean vector of designated bits on a pool P is embedded into a
low-degree polynomial E\_P over \ensuremath{\mathbb{F}}. With poly(n) evaluations of E\_P at
chosen points and standard interpolation, a solver can recover
seed-dependent linear aggregates over P that our model would require
\(\Omega\)(\textbar P\textbar/B) designated reads to obtain. This
violates the RWA-based per-step inflow cap and demonstrates that the
conservation-law pricing does not transfer verbatim to algebrized
models.

\subsubsection{N.4 Scope recap}\label{n4-scope-recap}

Our approach deliberately uses techniques outside classical barrier
scopes, enabling the unconditional Structural OWF construction (§9):

\begin{itemize}
\tightlist
\item
  Natural proofs: properties used here are exponentially
  sparse/non-constructive; barrier inapplicable.
\item
  Relativization: arguments are non-relativizing by construction;
  CRO-stability holds.
\item
  Algebrization: counting is non-algebrizing; algebraic size bounds rely
  on labeled bridges.
\end{itemize}

\textbf{Usage:} This appendix provides theoretical context for readers
familiar with complexity barriers. Not required for understanding main
proof (§7-10). See §9 (Structural OWF construction), §10 (classical
bridge), Appendix O (OWF packaging details).

\subsection{Appendix O: Unconditional OWF Packaging
Details}\label{appendix-o-unconditional-owf-packaging-details}

\textbf{Purpose:} Centralizes OWF packaging details from §9 (sampler,
function definition, Ext, coin-fixing security). Recaps §9.2-§9.4
construction and security argument; no new claims - provides complete
reference for OWF path (primary proof route).

O.1 Sampler and Function Definition (recap of §9.2).

\begin{itemize}
\tightlist
\item
  Sampler \ensuremath{\mathcal{S}}: draws a base CNF \(\varphi\), salts, and overlay metadata;
  samples r \(\in\) D(\(\varphi\)) uniformly.
\item
  Domain D(\(\varphi\)) = \{ r \textbar{}
  WellFormedRandomness(\(\varphi\), r) \} where WellFormedRandomness
  requires \(\varphi\).satisfies(r.assignment) and correct identity
  digests.
\item
  Preimage r := (assignment, gateDigests, structuralSalt). The planting
  procedure creates x* with identity-based digests---\textbf{no
  assignment bits are written to x*}.
\item
  Function family f\_n: D(\(\varphi\)\_n) \(\to\) \{0,1\}\^{}(poly(n))
  with f(r) := Plant(\(\varphi\), r) and FG wiring (GREQ beyond horizon;
  Appendix C.1.1). Every output x* inherits the deterministic single-run
  lower bound (Theorem 8.A).
\item
  Input distribution D\_n: the uniform distribution on the domain
  D(\(\varphi\)\_n). Security is measured over r\(\leftarrow\)D\_n and
  the inverter\textquotesingle{}s coins.
\end{itemize}

O.2 Canonical Witness and Ext (recap of §9.3 and §10.1.1).

\begin{itemize}
\tightlist
\item
  Verifier Algorithm V recomputes seeds, designated addresses, and gate
  digests from x* and checks W = (w, G\_\(\tau\), Dig\_\(\tau\)) in
  polytime (§10.1).
\item
  Extraction Ext(r', x*): Given any r' \(\in\) D(\(\varphi\)) with f(r')
  = x*, parses seed chain from x*, and outputs canonical witness W
  containing r'.assignment in polynomial time
  (poly(\textbar{}x*\textbar{})). By domain constraint, r'.assignment
  satisfies \(\varphi\) (§9.3; §10.1.1).
\item
  Domain constraint note (Lemma 9.DOM): Any valid preimage r' \(\in\)
  D(\(\varphi\)) contains a satisfying assignment. Unlike models
  encoding w in the output, preimages are \textbf{not unique}---multiple
  satisfying assignments can produce valid domain elements. Security
  relies on the SAT-hardness of finding any valid domain element, not on
  preimage uniqueness.
\end{itemize}

O.2.1 Verifier/Extractor Canonicality Checklist

\begin{itemize}
\tightlist
\item
  Seeds and DAG:

  \begin{itemize}
  \tightlist
  \item
    Recompute every \texttt{Seed\_v} from canonical \texttt{Enc} fields;
    verify parent lists are in the prescribed sorted order and that
    recomputed seeds match transcript values.
  \item
    Enforce fixed versioning, endianness, and field widths; reject
    duplicate or ambiguous encodings that would admit multiple byte
    strings for the same structure.
  \end{itemize}
\item
  Decode schema \texttt{Phi\_tilde}:

  \begin{itemize}
  \tightlist
  \item
    Verify the global slot order and published gate parity values.
  \item
    Domain membership: verify r.assignment satisfies \(\varphi\) and
    r.gateDigests match emergent parities.
  \end{itemize}
\item
  Witness serialization:

  \begin{itemize}
  \tightlist
  \item
    Check gates are topologically ordered; parent references use
    canonical IDs; field layouts match the canonical specification.
  \item
    Require that re-encoding \texttt{W} reproduces the exact canonical
    byte string.
  \end{itemize}
\item
  Determinism and complexity:

  \begin{itemize}
  \tightlist
  \item
    \texttt{V} must decide "is\_canonical(W, x*)" in polynomial time;
    \texttt{Ext} must deterministically return a canonical witness W for
    any valid domain element r' \(\in\) D(\(\varphi\)).
  \item
    The extracted witness W contains r'.assignment, which satisfies
    \(\varphi\) by domain constraint (Lemma 9.DOM).
  \item
    Reject valid-but-noncanonical witnesses and any representation that
    relies on alternative ordering or ambiguous parsing.
  \end{itemize}
\end{itemize}

O.3 Security via Coin-Fixing (recap of §9.4).

\begin{itemize}
\item
  Deterministic bound. By Theorem 8.A, any fixed computational run on x*
  requires \(\geq\) 2\^{}(\(\rho\)-s) · \(\Omega\)(n/W\_min) steps;
  \(\rho\) \ensuremath{\gtrsim} \(\lambda\)\_base and s \(\leq\)
  \(\Theta\)(\(\tau\)·\(\lambda\)\_base) by FG (App. C.2).
\item
  Domain-constrained inversion. Suppose a uniform PPT inverter \ensuremath{\mathcal{I}}
  satisfies Pr\_\{r\(\leftarrow\)D\_n, coins\}{[} \ensuremath{\mathcal{I}}(f(r)) \(\in\)
  D(\(\varphi\)) \ensuremath{\land} f(\ensuremath{\mathcal{I}}(...)) = f(r) {]} \(\geq\) 1/poly(n). By
  coin-fixing (averaging over the inverter\textquotesingle{}s coins),
  \(\exists\) coins c- such that Pr\_\{r\(\leftarrow\)D\_n\}{[}
  \ensuremath{\mathcal{I}}\_\{c-\}(f(r)) \(\in\) D(\(\varphi\)) \ensuremath{\land} f(\ensuremath{\mathcal{I}}\_\{c-\}(...)) = f(r) {]}
  \(\geq\) 1/poly(n). In particular, \(\exists\) r* with
  \ensuremath{\mathcal{I}}\_\{c-\}(f(r*)) = r' \(\in\) D(\(\varphi\)). By Lemma 9.DOM,
  r'.assignment satisfies \(\varphi\). Compose \ensuremath{\mathcal{I}}\_\{c-\} with Ext to
  obtain a deterministic poly-time witness-finder on x* = f(r*),
  contradicting Theorem 8.A.

  \begin{itemize}
  \item
    Key insight: Security does NOT require preimage uniqueness. The
    domain constraint D(\(\varphi\)) ensures any valid output contains a
    satisfying assignment. The adversary cannot produce r' \(\in\)
    D(\(\varphi\)) without solving SAT---precisely what Theorem 8.A
    forbids in polynomial time.
  \item
    Equivalent view. Since Theorem 8.A forbids any deterministic
    poly-time run from finding a satisfying assignment for any x*, any
    randomized PPT inverter (a distribution over deterministic runs) has
    negligible success probability over D\_n. Thus f is one-way against
    classical PPT. The classical bridge (§10.4) then yields FP \(\neq\)
    FNP \(\Rightarrow\) P \(\neq\) NP for the stated model scope.
  \end{itemize}
\end{itemize}

O.5 Non-Leaking Planting Encoder (identity-based digests).

\begin{itemize}
\tightlist
\item
  Domain structure. D(\(\varphi\)) = \{ r \textbar{}
  WellFormedRandomness(\(\varphi\), r) \} where r = (assignment,
  gateDigests, structuralSalt).
\item
  Non-leaking property. The planted instance x* contains only
  identity-based digests (XOR of emergent configurations at FG gates),
  NOT the assignment bits directly. No information about r.assignment is
  exposed in x*.
\item
  Parity encoding. For each FG gate v, the digest GateDigest\_v = XOR of
  emergent configuration bits. By WellFormedRandomness, r.gateDigests
  must equal these parities when computed from r.assignment.
\item
  Domain membership (poly-time). Verify: (1)
  \(\varphi\).satisfies(r.assignment) ---
  O(\textbar{}\(\varphi\)\textbar) CNF evaluation; (2) r.gateDigests
  match emergent parities --- polynomial seed chain computation.
\item
  Security via domain constraint. Any valid preimage r' \(\in\)
  D(\(\varphi\)) contains a satisfying assignment (Lemma 9.DOM). Since
  x* reveals only parities---insufficient to determine any satisfying
  assignment---inversion requires SAT-solving. This is hard by Theorem
  8.A.
\end{itemize}

\textbf{Contrast with leaking models.} A naive OWF encoding w directly
into x* would enable trivial inversion (just read w). Our non-leaking
model avoids this: the adversary must find a satisfying assignment to
produce any valid domain element, which is provably hard.

O.4 Scope notes.

\begin{itemize}
\tightlist
\item
  Model: classical uniform PPT; quantum out of scope (§1.7, §12.12).
\item
  Correctness: relies on Keyedness, Hermeticity, Completeness,
  Parseability, FG wiring (§6; App. C.1.1).
\item
  Decoding: OAP removes "CNF-first" bypass; Theorem 10.4.1-BYP rules out
  post-hoc synthesis of Dig\_\(\tau\) from w alone.
\end{itemize}

\textbf{Usage:} Primary OWF proof path (§9 construction \(\to\) §10
classical bridge \(\to\) P \(\neq\) NP). Provides complete packaging
details for readers verifying coin-fixing argument (O.3) or Ext
correctness (O.2). For underlying mechanisms: Appendix A (A1-A5
properties), Appendix C (FG wiring, lane dichotomy).

\begin{center}\rule{0.5\linewidth}{0.5pt}\end{center}

\subsection{Appendix Z: Parameter and Complexity
Cheat-Sheet}\label{appendix-z-parameter-and-complexity-cheat-sheet}

Symbols (general)

\begin{itemize}
\tightlist
\item
  n: instance length \textbar{}x*\textbar{}
\item
  \(\omega\): feasible world variable (assignment in \(\Pi\))
\item
  C*: designated bottleneck cut
\item
  R\_v: emergence rank at node v; q\_v: functional determination at v;
  \(\lambda\)\_base := \(\Sigma\)\_\{v\(\in\)C*\} R\_v; \(\rho\) :=
  \(\Sigma\)\_\{v\(\in\)C*\} (R\_v \(-\) q\_v)
\item
  s: \(\Delta\)q\_C accumulated before the final segment along the
  accepting chain
\item
  \(\tau\)(n), \(\mu\): gate-horizon/global and local allowances; S(n)
  := \(\Theta\)(\(\tau\)(n)·\(\lambda\)\_base) +
  \(\mu\)·\(\lambda\)\_base bounds s
\item
  W\_min(n): profile width parameter; \textbar{}S(P)\textbar{} =
  \(\Theta\)(n/W\_min) per cut-gate digest
\item
  Alt(C) := \(\prod\)\_\{v\(\in\)C\} Alt\_v; \(\Phi\)(C) := log\(_2\)
  Alt(C)
\item
  B := k\(\lceil\)log\(_2\)\textbar{}\(\Gamma\)\textbar{}\(\rceil\)
  fresh bits/step on a k-tape TM
\item
  \(\eta\): randomness for Plant(\(\Phi\)\textasciitilde{}, w; \(\eta\)) in OWF
  packaging (§9)
\item
  s\_salt: salt word length for \(\sigma\)\_u \(\in\)
  \{0,1\}\^{}(s\_salt)
\end{itemize}

Core equalities/inequalities

\begin{itemize}
\tightlist
\item
  SCL (cut form): Q(C) + \(\Phi\)(C) \(\geq\) \(\Lambda\)(C) where
  \(\Lambda\)(C) := \(\Sigma\)\_\{v\(\in\)C\} R\_v and Q(C) :=
  \(\Sigma\)\_\{v\(\in\)C\} q\_v
\item
  Across-tries (Appendix J): Q(C) + log\(_2\) Alt(C) + log\(_2\)
  \ensuremath{\mathbb{E}}{[}tries{]} \(\geq\) \(\Lambda\)(C)
\item
  Segment Counting (Appendix C.2): m\_seg \(\geq\) 2\^{}(\(\rho\) - s)
  non-accepting rollback segments (with CDT/WC)
\item
  Per-segment cost (TM): \(\Omega\)(n/W\_min) steps per non-accepting
  segment (priced event)
\end{itemize}

Time conversions (sequential k-tape TM)

\begin{itemize}
\tightlist
\item
  Restart lane: \ensuremath{\mathbb{E}}{[}tries{]} \(\geq\) 2\^{}(\(\Lambda\)(C*) - (Q(C*) +
  log\(_2\) Alt(C*))) and steps \(\geq\) \ensuremath{\mathbb{E}}{[}tries{]}/B
\item
  Single-run lane: steps \(\geq\) m\_seg·\(\Omega\)(n/W\_min) with
  m\_seg \(\geq\) 2\^{}(\(\rho\) - s)
\end{itemize}

Verification (NP membership)

\begin{itemize}
\tightlist
\item
  Verifier recomputes Seeds and GateDigest\_v; recomputes NF\_C and
  WorldCommit\_C at boundaries; checks ConstraintDigest\_C and
  UnitRefute certificates; all in polynomial time
\item
  Complexity bound: O(n log n) with caching, O(n²) naive (see Lemma
  10.2.1)
\end{itemize}

Typical QP-sharp calibration (illustrative)

\begin{itemize}
\tightlist
\item
  \(\lambda\)\_base(n) = \(\Theta\)(log² n); choose \(\tau\)(n) =
  \(\Theta\)(log n/\(\lambda\)\_base) = \(\Theta\)(1/log n); S(n) =
  \(\Theta\)(log n)
\item
  W\_min(n) = \(\Theta\)(log n); \textbar{}S(P)\textbar{} =
  \(\Theta\)(n/log n)
\item
  Bounds: m\_seg \(\geq\) 2\^{}((1-\(\Theta\)(\(\tau\)))\(\rho\)) and
  per-segment cost \(\Omega\)(n/log n) \(\Rightarrow\) time \(\geq\)
  n\^{}(\(\Theta\)(log n\_core)) up to poly factors
\end{itemize}

Scope notes

\begin{itemize}
\tightlist
\item
  All lower bounds are for L* (construction-specific), not general NP;
  Keyedness and disjoint pools are instance-side properties
\item
  Randomized solvers: distributional bounds + Yao \(\Rightarrow\) fixed
  worst-case instances per solver
\end{itemize}

\begin{center}\rule{0.5\linewidth}{0.5pt}\end{center}

\subsection{References}\label{references}

\subsubsection{Core Papers}\label{core-papers}

\begin{itemize}
\tightlist
\item
  \textbf{{[}AAR09{]}} Aaronson \& Wigderson 2009. Algebrization: A New
  Barrier in Complexity Theory. ACM Trans. Comp. Theory 1(1).
\item
  \textbf{{[}BAK75{]}} Baker, Gill \& Solovay 1975. Relativizations of
  the P =? NP Question. SIAM J. Computing 4(4).
\item
  \textbf{{[}BEA03{]}} Beame, Kautz \& Sabharwal 2003. Clause Learning.
  J. Artificial Intelligence Research 22.
\item
  \textbf{{[}BEN01{]}} Ben-Sasson \& Wigderson 2001. Short proofs are
  narrow - resolution made simple. JACM 48(2).
\item
  \textbf{{[}BRY86{]}} Bryant 1986. Graph-Based Algorithms for Boolean
  Function Manipulation. IEEE Trans. Computers C-35(8).
\item
  \textbf{{[}COO73{]}} Cook \& Reckhow 1973. Time bounded random access
  machines. J. Comp. Sys. Sci. 7(4).
\item
  \textbf{{[}RAZ97{]}} Razborov \& Rudich 1997. Natural Proofs. J. Comp.
  Sys. Sci. 55(1).
\item
  \textbf{{[}WEG00{]}} Wegener 2000. Branching Programs and Binary
  Decision Diagrams. SIAM Monographs.
\item
  \textbf{{[}WIL14{]}} Williams 2014. Algorithms for Circuits and
  Circuits for Algorithms. CACM 57(9).
\item
  \textbf{{[}YAO77{]}} Yao 1977. Probabilistic computations.
  FOCS\textquotesingle{}77.
\end{itemize}

\subsubsection{Foundational Complexity}\label{foundational-complexity}

\begin{itemize}
\tightlist
\item
  \textbf{{[}TUR36{]}} Turing 1936. On computable numbers, with an
  application to the Entscheidungsproblem. Proc. London Math. Soc.
  2(42):230-265.
\item
  \textbf{{[}COO71{]}} Cook 1971. The complexity of theorem-proving
  procedures. STOC\textquotesingle{}71.
\item
  \textbf{{[}KAR72{]}} Karp 1972. Reducibility among combinatorial
  problems. Complexity of Computer Computations.
\item
  \textbf{{[}LEV73{]}} Levin 1973. Universal sequential search problems.
  Problems of Information Transmission 9(3).
\item
  \textbf{{[}IL90{]}} Impagliazzo \& Levin 1990. No better ways to
  generate hard NP instances than picking uniformly at random.
  FOCS\textquotesingle{}90.
\item
  \textbf{{[}IMP95{]}} Impagliazzo 1995. A personal view of average-case
  complexity. Structure in Complexity Theory Conference, pp. 134-147.
\item
  \textbf{{[}GAR79{]}} Garey \& Johnson 1979. Computers and
  Intractability. W.H. Freeman.
\item
  \textbf{{[}PAP94{]}} Papadimitriou 1994. Computational Complexity.
  Addison-Wesley.
\item
  \textbf{{[}HAR65{]}} Hartmanis \& Stearns 1965. On the computational
  complexity of algorithms. Trans. AMS 117.
\item
  \textbf{{[}STE65{]}} Stearns, Hartmanis \& Lewis 1965. Hierarchies of
  memory limited computations. Switching Theory\textquotesingle{}65.
\end{itemize}

\subsubsection{Cryptography}\label{cryptography}

\begin{itemize}
\tightlist
\item
  \textbf{{[}DH76{]}} Diffie \& Hellman 1976. New directions in
  cryptography. IEEE Trans. Information Theory 22(6):644-654.
\item
  \textbf{{[}YAO82{]}} Yao 1982. Theory and applications of trapdoor
  functions. FOCS\textquotesingle{}82, pp. 80-91.
\item
  \textbf{{[}GOL01{]}} Goldreich 2001. Foundations of Cryptography, Vol.
  1: Basic Tools. Cambridge University Press.
\item
  \textbf{{[}IR89{]}} Impagliazzo \& Rudich 1989. Limits on the provable
  consequences of one-way permutations. STOC\textquotesingle{}89.
\end{itemize}

\subsubsection{Information Theory}\label{information-theory}

\begin{itemize}
\tightlist
\item
  \textbf{{[}SHA48{]}} Shannon 1948. A mathematical theory of
  communication. Bell Sys. Tech. J. 27(3-4).
\item
  \textbf{{[}HAR28{]}} Hartley 1928. Transmission of information. Bell
  Sys. Tech. J. 7(3).
\item
  \textbf{{[}REN61{]}} Rényi 1961. On measures of entropy and
  information. Proc. 4th Berkeley Symposium on Math. Stat. and
  Probability 1:547-561.
\item
  \textbf{{[}COV06{]}} Cover \& Thomas 2006. Elements of Information
  Theory, 2nd Ed. Wiley.
\end{itemize}

\subsubsection{Communication \& Pebbling}\label{communication--pebbling}

\begin{itemize}
\tightlist
\item
  \textbf{{[}YAO79{]}} Yao 1979. Some complexity questions related to
  distributive computing. STOC\textquotesingle{}79.
\item
  \textbf{{[}KUS97{]}} Kushilevitz \& Nisan 1997. Communication
  Complexity. Cambridge Univ. Press.
\item
  \textbf{{[}COO76{]}} Cook \& Sethi 1976. Storage Requirements for
  Deterministic Polynomial Time Recognizers. J. Comp. Sys. Sci. 13(1).
\item
  \textbf{{[}LEN82{]}} Lengauer \& Tarjan 1982. Asymptotically tight
  bounds on time-space trade-offs in a pebble game. JACM 29(4).
\item
  \textbf{{[}PIP77{]}} Pippenger 1977. Pebbling. IBM Symp. Math.
  Foundations of CS.
\end{itemize}

\subsubsection{Expanders \& Spectral
Theory}\label{expanders--spectral-theory}

\begin{itemize}
\tightlist
\item
  \textbf{{[}HOO06{]}} Hoory, Linial \& Wigderson 2006. Expander graphs
  and their applications. Bull. AMS 43(4).
\item
  \textbf{{[}LUB88{]}} Lubotzky, Phillips \& Sarnak 1988. Ramanujan
  graphs. Combinatorica 8(3).
\item
  \textbf{{[}REI02{]}} Reingold, Vadhan \& Wigderson 2002. Entropy
  waves, the zig-zag graph product. Ann. Math. 155(1).
\end{itemize}

\subsubsection{Circuit \& Proof
Complexity}\label{circuit--proof-complexity}

\begin{itemize}
\tightlist
\item
  \textbf{{[}CEI96{]}} Clegg, Edmonds \& Impagliazzo 1996. Using the
  Groebner basis algorithm to find proofs of unsatisfiability.
  STOC\textquotesingle{}96.
\item
  \textbf{{[}GRI01{]}} Grigoriev 2001. Linear lower bound on degrees of
  Positivstellensatz calculus proofs for the parity. Theoretical
  Computer Science 259(1-2).
\item
  \textbf{{[}HAS86{]}} Håstad 1986. Almost optimal lower bounds for
  small depth circuits. STOC\textquotesingle{}86.
\item
  \textbf{{[}IMP99{]}} Impagliazzo, Pudlák \& Sgall 1999. Lower bounds
  for the polynomial calculus and the Gröbner basis algorithm.
  Computational Complexity 8(2).
\item
  \textbf{{[}LAU03{]}} Laurent 2003. A comparison of the Sherali-Adams,
  Lovász-Schrijver, and Lasserre relaxations. Mathematics of Operations
  Research 28(3).
\item
  \textbf{{[}PUD97{]}} Pudlák 1997. Lower bounds for resolution and
  cutting plane proofs and monotone computations. J. Symbolic Logic
  62(3).
\item
  \textbf{{[}RAZ85{]}} Razborov 1985. Lower bounds on monotone
  complexity. Soviet Math. Doklady 31.
\item
  \textbf{{[}SCH08{]}} Schoenebeck 2008. Linear level Lasserre lower
  bounds for certain k-CSPs. FOCS\textquotesingle{}08.
\item
  \textbf{{[}SMO87{]}} Smolensky 1987. Algebraic methods in Boolean
  circuit complexity. STOC\textquotesingle{}87.
\item
  \textbf{{[}TSE68{]}} Tseitin 1968. On the complexity of derivation in
  propositional calculus. Studies in Constructive Mathematics and
  Mathematical Logic, Part 2:115-125.
\item
  \textbf{{[}NS94{]}} Nisan \& Szegedy 1994. On the degree of Boolean
  functions as real polynomials. Computational Complexity 4(4):301-313.
\end{itemize}

\subsubsection{SAT Solving \&
Tractability}\label{sat-solving--tractability}

\begin{itemize}
\tightlist
\item
  \textbf{{[}DAV62{]}} Davis, Putnam, Logemann \& Loveland 1962. A
  machine program for theorem-proving. CACM 5(7).
\item
  \textbf{{[}MSS01{]}} Moskewicz, Madigan, Zhao, Zhang \& Malik 2001.
  Chaff: Engineering an Efficient SAT Solver. DAC\textquotesingle{}01,
  pp. 530-535.
\item
  \textbf{{[}APT79{]}} Aspvall, Plass \& Tarjan 1979. A linear-time
  algorithm for testing the truth of certain quantified boolean
  formulas. Information Processing Letters 8(3):121-123.
\end{itemize}

\subsubsection{Network Flow}\label{network-flow}

\begin{itemize}
\tightlist
\item
  \textbf{{[}FF56{]}} Ford \& Fulkerson 1956. Maximal flow through a
  network. Canadian J. Mathematics 8:399-404.
\item
  \textbf{{[}DIN70{]}} Dinic 1970. Algorithm for solution of a problem
  of maximum flow in networks with power estimation. Soviet Math.
  Doklady 11:1277-1280.
\end{itemize}

\subsubsection{Algorithms \&
Approximation}\label{algorithms--approximation}

\begin{itemize}
\tightlist
\item
  \textbf{{[}HEL62{]}} Held \& Karp 1962. Dynamic programming approach
  to sequencing. J. SIAM 10(1).
\item
  \textbf{{[}VAZ01{]}} Vazirani 2001. Approximation Algorithms.
  Springer.
\end{itemize}

\subsubsection{Quantum Computation}\label{quantum-computation}

\begin{itemize}
\tightlist
\item
  \textbf{{[}SHO97{]}} Shor 1997. Polynomial-time algorithms for prime
  factorization and discrete logarithms on a quantum computer. SIAM J.
  Computing 26(5):1484-1509.
\end{itemize}

\subsubsection{Probability \&
Concentration}\label{probability--concentration}

\begin{itemize}
\tightlist
\item
  \textbf{{[}MAR06{]}} Markov 1906. Extension of limit theorems to
  chain-connected variables. (Russian; English in Selected Works).
\item
  \textbf{{[}PAL32{]}} Paley \& Zygmund 1932. A note on analytic
  functions in the unit circle. Math. Proc. Cambridge Phil. Soc. 28(3).
\end{itemize}

\subsubsection{Textbooks}\label{textbooks}

\begin{itemize}
\tightlist
\item
  \textbf{{[}SIP13{]}} Sipser 2013. Introduction to the Theory of
  Computation, 3rd Ed. Cengage Learning.
\item
  \textbf{{[}AB09{]}} Arora \& Barak 2009. Computational Complexity: A
  Modern Approach. Cambridge University Press.
\end{itemize}

\end{document}
